\documentclass[11pt]{amsart}
\usepackage{amsmath,amssymb,amsthm,mathtools,mathrsfs}
\usepackage[margin=1.08in]{geometry}
\usepackage{booktabs}
\usepackage{array}
\usepackage{longtable}
\usepackage{tabularx}
\usepackage[hypertexnames=false]{hyperref}
\hypersetup{hidelinks}
\usepackage[final,expansion=false]{microtype}
\usepackage{url}
\setlength{\parskip}{4pt plus 1pt}
\emergencystretch=2.5em

\newtheorem{theorem}{Theorem}[section]
\newtheorem{proposition}[theorem]{Proposition}
\newtheorem{lemma}[theorem]{Lemma}
\newtheorem{corollary}[theorem]{Corollary}
\theoremstyle{definition}
\newtheorem{definition}[theorem]{Definition}
\newtheorem{problem}[theorem]{Problem}
\theoremstyle{remark}
\newtheorem{remark}[theorem]{Remark}

\newcommand{\C}{\mathbb C}
\newcommand{\R}{\mathbb R}
\newcommand{\N}{\mathbb N}
\newcommand{\Z}{\mathbb Z}
\newcommand{\Q}{\mathbb Q}
\newcommand{\Rel}{\operatorname{Re}}
\newcommand{\Iml}{\operatorname{Im}}
\newcommand{\supp}{\operatorname{supp}}
\newcommand{\TV}{\operatorname{TV}}
\newcommand{\Kp}{\mathcal K^{\mathrm{prim}}_{2T}}
\newcommand{\Kw}{\mathcal K_{2T}}
\newcommand{\PT}{\mathcal P_T}
\newcommand{\IT}{I_T}
\newcommand{\lamP}{\lambda^{\mathrm{prim}}_T}
\newcommand{\lamW}{\lambda^{\mathrm W}_T}
\newcommand{\lamA}{\lambda_{\mathrm A}(T)}
\newcommand{\sT}{\sigma_T}
\newcommand{\gh}[1]{\widehat g_{#1}}
\newcommand{\nuh}{\widehat\nu}
\newcommand{\wh}{\widehat w}
\newcommand{\Fh}{\widehat F}
\newcommand{\cF}{c_F}
\newcommand{\Lev}{\mathfrak E}
\newcommand{\Wcl}{\mathcal W_T}
\newcommand{\NT}{\mathcal N_T}
\newcommand{\Aconst}{\mathsf A_T}
\newcommand{\Ainf}{\mathcal A_\infty}
\newcommand{\PiT}{\Pi_T}
\newcommand{\mT}{\mathfrak m_T}
\newcommand{\laminf}{\lambda^{\infty}_{\min}}
\newcommand{\laminfP}{\lambda^{\infty,\mathrm P}_{\min}}
\newcolumntype{Y}{>{\raggedright\arraybackslash}X}

\begin{document}

\title[Weil positivity on a finite window]
{Weil Positivity on a Finite Window:\\
Duality, Certificates and a Finite Reduction}

\author{David Alejandro Trejo Pizzo}
\thanks{Independent Researcher. \texttt{dtrejopizzo@gmail.com}. The certificate data and verification
scripts of the computer-assisted results accompany the paper.}

\date{\today}

\begin{abstract}
For $T>0$ let $\mathfrak a_T$ be the truncated Weil quadratic form on $L^2(-T,T)$, $\PT$ the
subspace on which the two Tate moments vanish, $\lamA$ the bottom of $\mathfrak a_T$ on $\PT$, and
$\lamW$ the bottom of the complete form $W_T=\mathfrak a_T+2\Rel(A\bar B)$ on $L^2(-T,T)$. We study
the window positivity problem $\lamA\ge0$ for every $T>0$. A certified computation in ball
arithmetic gives $\lamA\ge2.5\cdot10^{-12}$ and $\lamW\ge10^{-18}$ for all $T\le\frac12\log5$. For
general $T$ we develop the dual side of the problem. The archimedean multiplier is a L\'evy exponent
with an explicit completely monotone density; the bottom is bounded below by the value of a
semi-infinite linear program whose dual variables are measures supported outside $[-2T,2T]$; strong
duality holds with finitely atomic dual sequences, the dual value is never attained, and every
certificate of the complete form is, up to an arbitrarily small uniform correction, the explicit formula
evaluated on modulated flat-top test functions. An explicit certificate with $437$ atoms gives
$\lambda_{\mathrm A}(\frac12\log3)\ge8.8577\cdot10^{-5}$. The function $T\mapsto\lamA$ is continuous and
nonincreasing; at one window the complete and the primitive bottoms differ by a diagonal $2\times2$
polar capacity; for every $T$ the sign of $\lamA$ is the sign of an explicit symmetric matrix of size
$2k-2$ indexed by piecewise Tate exponentials; and a first failure of positivity, if any, occurs at a
window $T_*>\frac12\log5$ carrying a primitive null vector of definite parity (by a cited, unrefereed
certificate of Mori, whose computational part we have reproduced, at a window $T_*>\frac12\log15$). We also prove limits of
the method: ceilings and class impossibilities for lattice certificates, the exact decay
$\lamA\le4\cdot10^9Te^{-15.449T}$ of the margin, a doubly exponential envelope threshold, and the
failure at explicit windows of every argument that uses only properties of the primes stable under
finite modification of the weights. We further describe the window form as a truncated Wiener--Hopf operator (a trace-class pencil whose
inertia crossings carry the sign, an exact Schur coupling and a Birman--Schwinger determinant), prove that it
sees the symbol only on the Nyquist mesh $\pi k/2T$ and splits into two sectors, develop a theory of
witnesses of negativity, and prove barrier theorems for linear, determinantal, moment, family, pattern and
local arguments. We show that combs positive at every window exist, smooth and atomic, that nonnegative combs on $\log\N$ cannot
have lower logarithmic density above one, and that positivity for all windows is equivalent to the positive definiteness of one explicit function
of one real variable. The remaining problem is stated precisely.
\end{abstract}

\maketitle
\tableofcontents

\section{Introduction}\label{sec:intro}

\subsection{The problem}
Throughout, $T>0$, $\IT=(-T,T)$, and for $F\in L^2(\IT)$ (extended by zero) we write
$\Fh(\tau)=\int F(t)e^{i\tau t}\,dt$ and $\cF(u)=\int F(t)\overline{F(t-u)}\,dt$ for the
autocorrelation, so that $\cF$ is continuous, supported in $[-2T,2T]$, and
$\widehat{\cF}=|\Fh|^2$ on $\R$. Put
\begin{equation}\label{eq:symbol}
m(\tau)=\Rel\psi\!\left(\tfrac14+\tfrac{i\tau}{2}\right)-\log\pi,\qquad
\sT(\tau)=m(\tau)-\sum_{n\in\NT}v_n\cos(\tau\log n),\qquad v_n=\frac{2\Lambda(n)}{\sqrt n},
\end{equation}
where $\psi$ is the digamma function, $\NT=\{n\ge2:\ \log n<2T\}$ is the set of \emph{contacts} and
$\Aconst:=\sum_{n\in\NT}v_n$ is the total mass of the prime comb. The truncated Weil form is
\begin{equation}\label{eq:form}
\mathfrak a_T(F)=\frac1{2\pi}\int_\R\sT(\tau)\,|\Fh(\tau)|^2\,d\tau
\end{equation}
on the form domain $\mathcal D_T=\{F\in L^2(\IT):\int\log(2+|\tau|)|\Fh|^2<\infty\}$. The two \emph{Tate
moments} are $A(F)=\int F(t)e^{t/2}dt$ and $B(F)=\int F(t)e^{-t/2}dt$, the \emph{primitive subspace} is
$\PT=\{F\in L^2(\IT):A(F)=B(F)=0\}$, and the \emph{complete Weil form} on the window is
\begin{equation}\label{eq:W}
W_T(F)=\mathfrak a_T(F)+2\,\Rel\big(A(F)\overline{B(F)}\big).
\end{equation}
By the explicit formula (\cite{Weil52}, \cite[Thm.~8.4]{Paper45}; Lemma~\ref{lem:roughEF} below for
tests that are not smooth), $W_T(F)=\sum_\rho\Phi_F(\rho)\Phi_F(1-\rho)$ with
$\Phi_F(s)=\int F(t)e^{(s-\frac12)t}dt$, the sum over the nontrivial zeros of $\zeta$. We write
\[
\lamA=\inf\{\mathfrak a_T(F):F\in\PT\cap\mathcal D_T,\ \|F\|=1\},\qquad
\lamW=\inf\{W_T(F):F\in\mathcal D_T,\ \|F\|=1\};
\]
$\lamW\le\lamA$ because the two forms agree on $\PT$. The object of this paper is the following
statement, which we call \emph{window positivity}:
\begin{equation}\label{eq:target}
\lamA\ \ge\ 0\qquad\text{for every }T>0 ,
\end{equation}
together with its complete version $\lamW\ge0$ for every $T>0$. Since $T\mapsto\lamA$ is nonincreasing
(Lemma~\ref{lem:monoT}), \eqref{eq:target} is equivalent to $\lim_{T\to\infty}\lamA\ge0$; it is a
statement about every window at once, and every window carries a finite amount of arithmetic data, the
vector $(\Lambda(n))_{n<e^{2T}}$ (Proposition~\ref{prop:finite}).

Positivity of the complete form is classical for $T\le\frac12\log2$, where the comb is empty
\cite{Yoshida1992,Bombieri2000}, and was certified by Zhu for $T\le0.8$ \cite[Thm.~1.2]{Zhu2026}. In a preprint under review, Liu certifies the
complete form on windows of half-width $T\le1$ with $\lamW\ge2^{-151}$ and $T\le17/16$ with
$\lamW\ge2^{-49162}$ \cite[Thms.~A, B]{Liu2026}; his half-width is our $T$ (support of $F$ in $\IT$, comb
over $\log n<2T$). This result is cited, not used in any proof here. A
third bottom, $\lamP$, the value of a linear program over autocorrelations (Definition~\ref{def:PD}),
satisfies $\lamW\le\lamP\le\lamA$ (Theorem~\ref{thm:A}, Remark~\ref{rem:Kprim}); it is the quantity
that dual certificates bound from below.

\subsection{Results}
\emph{Positivity on an initial range.} A certified computation in ball arithmetic, built on a
frequency-truncated lower form, a Legendre section of $140$ modes per parity and a Schur and Haynsworth
inertia count, gives
\[
\lamA\ \ge\ 2.5\cdot10^{-12},\qquad \lamW\ \ge\ 10^{-18}\qquad\text{for every }0<T\le\tfrac12\log5
\]
(Theorems~\ref{thm:log5A} and \ref{thm:log5W}). The window $\frac12\log5$ is the last one with contact
set $\{2,3,4\}$. Liu \cite{Liu2026} reports a certified complete-form bound up to $T=17/16$, past the
contacts $5,7,8$, with constant $2^{-49162}$. The largest certified window we know of is due to Mori
\cite{Mori15}: $\lamW\ge2^{-2904}$ for every $T\le\frac12\log15=1.35402\ldots$ (contacts up to $13$), in
an unrefereed deposit that builds on earlier papers of the same series at the windows $\frac12\log10$ and
$\frac12\log12$. We have independently reproduced its computational part (the author's checker at
$1280$ bits, finite margins $2.3979\cdot10^{-70}$ even and $3.9458\cdot10^{-66}$ odd, and independent
spot checks of scalars and matrix entries); we have not re-checked its analytic reduction lemmas, and we
label the result as cited. By monotonicity and $\lamW\le\lamP\le\lamA$ it gives $\lamA>0$ on
$(0,\frac12\log15]$. Those constants collapse doubly exponentially in $T$: it is a record window,
not a method that extends to all $T$. For comparison, converged Galerkin computations (numerical, not certified; Galerkin
values are upper bounds) give $\lambda_{\mathrm A}(\frac12\log3)\approx9.448\cdot10^{-4}$,
$\lambda_{\mathrm A}(\log2)\approx9.83\cdot10^{-8}$ and $\lambda_{\mathrm A}(\frac12\log5)\approx3.10\cdot10^{-12}$
(Appendix~\ref{app:num}). Independently, at $T=\frac12\log3$ an explicit dual feasible point with $437$
atoms proves $\lambda_{\mathrm A}(\frac12\log3)\ge\lambda^{\mathrm{prim}}_{\frac12\log3}\ge8.8577\cdot10^{-5}$
(Theorem~\ref{thm:C}).

\emph{Duality.} The archimedean multiplier is a L\'evy exponent:
$m(\tau)=C_0+\int_0^\infty(1-\cos\tau u)K(u)\,du$ with $C_0=\psi(\frac14)-\log\pi$ and the completely
monotone density $K(u)=\frac12/\sinh(u/2)+\frac12/\cosh(u/2)$; equivalently $m-C_0=\varphi(\tau^2)$ with
$\varphi$ a complete Bernstein function (Theorem~\ref{thm:S}). The bottom $\lamP$ is the value of a
semi-infinite linear program whose dual is: maximize $\inf_\tau[\sT+\nuh+\beta_0\gh0+\beta_1\gh1]$ over
finite even measures $\nu$ supported in $\{|u|\ge2T\}$ and $\beta\in\R^2$. Strong duality holds with
finitely atomic $\nu$ (Theorem~\ref{thm:D1}), so positivity with a margin is always provable by a finite
certificate; the dual is never attained (Theorem~\ref{thm:O}); exponential columns exhaust it
(Theorem~\ref{thm:E}); and $\beta=(2,0)$ is the unique multiplier for which a certificate bounds the
complete form (Proposition~\ref{prop:P}).

\emph{The dual variable is the explicit formula.} For an even weight $w$ with $w\equiv1$ on $[-2T,2T]$
and sufficient decay, an explicit measure $\nu_w$ outside the window is dual feasible and
$\sT+\nuh_w+2\gh0=f_w$ with $f_w(\tau)=\sum_\rho\wh(\tau-\gamma_\rho)$ (Theorem~\ref{thm:Z}); every
certificate of the complete form is of this kind up to an arbitrarily small uniform correction, and
$\sup_w\inf f_w=\lamW$ (Theorem~\ref{thm:Zexh}). A flat-top weight never has a nonnegative Fourier
transform (Lemma~\ref{lem:Rrig}), and no finite set of zeros carries a flat-top certificate
(Theorem~\ref{thm:finZ}). At each window the certificates are the absolutely continuous solutions of a
Kre\u{\i}n extension problem (Theorems~\ref{thm:KEP1}--\ref{thm:KEPdich}).

\emph{Structure at one window.} The function $T\mapsto\lamA$ is nonincreasing and continuous
(Lemmas~\ref{lem:monoT} and \ref{lem:contA}). The complete and the primitive form at one window are related
by a diagonal $2\times2$ matrix $\mathfrak G_T$, the polar capacity: $\lamW>0$ if and only if $\lamA>0$
and $\mathfrak G_T\succ0$ (Theorem~\ref{thm:parityschur}, Proposition~\ref{prop:polarcap}), and the
condition at one window controls all smaller ones (Proposition~\ref{prop:onewindow}). The archimedean form
alone is positive on $\PT$ up to a threshold $T^*_{\infty,\mathrm P}\in[0.58,0.60]$ and its negative index
is computed exactly for $T\le0.90$ outside two short intervals (Theorem~\ref{thm:index},
Corollary~\ref{cor:Tarch}); beyond the threshold the primes supply positivity.

\emph{A finite reduction.} Split $\IT$ into $k$ equal blocks and let $\mathcal V_k\subset\PT$ be the
functions satisfying both Tate conditions on every block. The archimedean form is coercive on
$\mathcal V_k$ with an explicit constant, $\mathfrak m_T\ge(\log\frac kT-C_B(T/k))\|\cdot\|^2$
(Theorem~\ref{thm:coer-B}), so for $k\ge k_0(T)$ the form $\mathfrak a_T$ is coercive there, and
\[
\lamA>0\iff S_{T,k}\succ0,\qquad \lamA\ge0\iff S_{T,k}\succeq0,
\]
where $S_{T,k}$ is the Schur complement of $\mathfrak a_T$ on the $(2k-2)$-dimensional space of piecewise
Tate exponentials (Theorem~\ref{thm:coer-schur}). Thus \eqref{eq:target} at each window is the positive
semidefiniteness of an explicit finite matrix; $k=25$ suffices at $T=\frac12\log5$.

\emph{First contact.} If \eqref{eq:target} fails, then by monotonicity, continuity and the certified range
there is a first window $T_*>\frac12\log5$ (and $T_*>17/16$ if \cite[Thm.~B]{Liu2026} is granted) with $\lambda_{\mathrm A}(T_*)=0$, and the kernel of the
primitive operator there contains a nonzero vector of definite parity (Lemma~\ref{lem:firstcontact}). For
such a vector we prove a finite Poisson--Chebyshev operator identity combining Poisson summation with
$\Lambda*1=\log$ (Theorem~\ref{thm:PCfinite}), a faithful observation by the odd poles of the Gamma factor,
an exact balance identity, and a strict contraction of the Euler channel (Proposition~\ref{prop:GEgap}).

\emph{Operator structure.} Splitting $\sT=\sigma_+-\sigma_-$, the primitive form is a positive form minus a
trace-class one, and $\lamA>0$ if and only if the largest eigenvalue $\beta_1$ of the associated pencil is
$<1$; the number of negative directions is the number of $\beta_j>1$, the Birman--Schwinger determinant
$\det(1-sB_T)$ is entire of genus $0$, and $1-\beta_1$ equals $\lamA$ up to explicit factors of order one
(Theorems~\ref{thm:wh-pencil}, \ref{thm:wh-det}, \ref{thm:wh-detinherit}). The complement of the prolate
functions of the negative set is coercive, and the sign is the sign of an explicit coupling quotient
$\rho_T<1$ (Theorems~\ref{thm:wh-LW}, \ref{thm:wh-schur}). By Poisson summation the window form reads the
symbol only on the Nyquist mesh: $\mathfrak a_T(F)=\frac1{4T}\sum_k|\Fh(\pi k/2T)|^2s_T(\pi k/2T)$ with $s_T$ the
symbol whose L\'evy measure is truncated to $[0,2T]$ (Theorem~\ref{thm:nyq}); the two halves of the mesh give
two sectors related by a unitary half-shift, and a two-sector Schur reduction yields an explicit sufficient
certificate without matrix inversion (Theorems~\ref{thm:nyq-sectors}, \ref{thm:nyq-schur}).

\emph{Witnesses.} A witness at $T$ is a primitive $F$ with $\mathfrak a_T(F)<0$. The windows carrying a
witness form an interval $(T_c,\infty)$ with $T_c>\frac12\log5$ ($T_c>17/16$ granting \cite{Liu2026}), a witness persists at all larger windows,
and every asymptotic incompatibility is therefore a statement at a fixed window
(Theorems~\ref{thm:wit-Tc}, \ref{thm:wit-asym}). Positivity implies, for every $\gamma>0$, an inequality
$(E_\gamma)$ for a twisted Chebyshev sum (Proposition~\ref{prop:wit-Egamma}), and the window form is, on the mesh,
a weighted coherence of the prime comb (Proposition~\ref{prop:wit-coh}).

\emph{Combs and one profile.} The Tate conditions turn the main term of the prime comb into a positive energy, so that
positivity is the inequality $\mathcal F_T\le\mathfrak a^\sharp$ between the Chebyshev fluctuation and an archimedean form with
symbol $\Rel\psi(\frac54+\frac{i\tau}2)-\log\pi$ (Proposition~\ref{prop:sharp}). Among combs, the smooth comb $2\kappa e^{u/2}du$ is
positive at every window if and only if $\kappa\ge\kappa_\star$, with $\kappa_\star\ge\frac14|C_0|>1$, and purely atomic combs positive
at every window exist (Theorems~\ref{thm:kappa}, \ref{thm:atomic}). A nonnegative comb on $\log\N$ positive at every window has
lower logarithmic density at most $1$, which the primes attain (Theorem~\ref{thm:ceiling}), and $\Lambda$ is positive on every test space that does not
see near-coincidences of integers (Proposition~\ref{prop:nearfree}). Positivity is equivalent to the positive definiteness of a single
explicit function of one real variable (Theorem~\ref{thm:oneprofile}), and to the positive semidefiniteness of a family of matrices
indexed by integers (Proposition~\ref{prop:AM}); its $2\times2$ case is a two-sided band for the local Chebyshev fluctuation at
every scale, which Poisson combs violate with probability $1-O(e^{-cT\log T})$ for every $c$ (Theorems~\ref{thm:band},
\ref{thm:poisson}).

\subsection{Barriers}
Several results limit what any proof of \eqref{eq:target} can look like. The exchange between escaped
primal mass and value is exactly logarithmic (Theorem~\ref{thm:Vp}), and no pointwise envelope can lower
the verification threshold $2\pi e^{\Aconst}$, doubly exponential in $T$ (Proposition~\ref{prop:env}).
Lattice certificates have an explicit ceiling, one small Toeplitz eigenvalue (Proposition~\ref{prop:TK});
the ceilings converge to $\lamP$ (Theorem~\ref{thm:R}), but no lattice of order $M\le8$ certifies at
$\frac12\log3$ and none of order $M\le24$ at $T=0.68$ or $T=\log2$ (Theorems~\ref{thm:N}, \ref{thm:Np}),
and sampling a flat-top certificate on a lattice costs exactly $\Theta(1/M)$
(Theorem~\ref{thm:latedge}). No partial prime comb of mass below $-C_0$ is positive on a large window
(Theorem~\ref{thm:SC}). The margin decays, $\lamA\le4\cdot10^9Te^{-15.449T}$ for $T\ge1$
(Corollary~\ref{cor:margin1}). Modifications of $\Lambda$ that preserve its asymptotics, the zeros and the
pole of the associated Dirichlet series violate positivity at explicit windows (Theorem~\ref{thm:finmod}), so
no argument that uses only such properties of the primes decides any window $T>0.75$
(Corollary~\ref{cor:void}); the same modifications break the per-window implication from primitive to
complete positivity (Corollary~\ref{cor:convfail}). Every bounded remainder after subtracting a positive
archimedean part has infinite negative index once a contact is present (Theorem~\ref{thm:inertia-nogo}).
Arguments that use only an Euler product, statistics in progressions, relative form bounds, monotone M\"obius iterations or
sums of squares of Dirichlet polynomials fail by explicit falsifiers as well (\S\ref{subsec:bar-arith}).

In \S\S\ref{sec:inertia}, \ref{sec:wh} and \ref{sec:barriers} these limits are made systematic. Every
bounded remainder after subtracting a positive part has infinite negative index (\S\ref{sec:inertia}); every
determinant, moment or trace criterion for the pencil has margin at most a fixed multiple of $1-\beta_1$, and
so of $\lamA$ (Lemma~\ref{lem:wh-T}, Theorem~\ref{thm:wh-detinherit}); bounds by norms in a fixed frame lose a
constant that does not decay (Lemma~\ref{lem:half}, Theorem~\ref{thm:posext}); no functional of the bottoms on a
bounded range of windows decides the target (Theorem~\ref{thm:famblind}); no list of window observables does
(Corollary~\ref{cor:pattern}); and no criterion local on the negative set of the symbol does
(Theorem~\ref{thm:nolocal}). Each barrier is a theorem about a declared class of arguments, proved by an explicit
falsifier, and each section says what its class does not exclude.

\subsection{What remains open}
Apart from the certified range, none of the results gives the sign of $\mathfrak a_T$ at a window
$T>\frac12\log5$. Beyond $\frac12\log5$ the certificates we know of are \cite{Liu2026}
($T\le17/16$, preprint under review, cited) and \cite{Mori15} ($T\le\frac12\log15$, unrefereed deposit,
cited; computational part reproduced by us); the problem is open for $T>\frac12\log15$ in any case. The open problem, in two equivalent forms proved in this paper, is: \emph{for every
$T>\frac12\log5$ and some (equivalently every) $k\ge k_0(T)$, the matrix $S_{T,k(T)}$ of
Theorem~\ref{thm:coer-schur} is positive semidefinite}; and, on the arithmetic side, \emph{a pointwise bound with
sign for the twisted Chebyshev sums of the prime comb at the Nyquist points $\pi k/2T$, weighted by the samples
of primitive test functions} (Proposition~\ref{prop:wit-coh}, Remark~\ref{rem:wit-cheb}). The barriers show that
the second cannot come from mean-square or absolute-value estimates. A third form is the positive semidefiniteness of the
integer-lattice matrices of Proposition~\ref{prop:AM}, where the sign is carried by the near-coincidence part. The precise
list is in \S\ref{sec:open}.

\subsection{Organization}
Sections~\ref{sec:levy}--\ref{sec:Z} develop the archimedean multiplier, the primal and dual programs, the tail
lemmas with an explicit certificate at $\frac12\log3$, and the identification of the dual variable with the
explicit formula, including the Kre\u{\i}n extension problem at one window. Section~\ref{sec:class} treats lattice
certificates, \S\ref{sec:archstr} the archimedean form and the negative set of the symbol, \S\ref{sec:negative}
obstructions to general methods, and \S\S\ref{sec:inertia}--\ref{sec:semilocal} inertia and the semilocal
decomposition at the first prime. Section~\ref{sec:complete} studies the complete form at one window,
\S\ref{sec:log5} certifies positivity up to $\frac12\log5$, \S\ref{sec:coer} proves the finite reduction,
\S\ref{sec:wh} the Wiener--Hopf structure, and \S\ref{sec:nyq} the Nyquist identity. Section~\ref{sec:firstcontact}
analyses a first contact, \S\ref{sec:wit} witnesses, \S\ref{sec:combs} combs and translates of one profile,
\S\ref{sec:barriers} barriers, \S\S\ref{sec:shape}--\ref{sec:spec} the shape of a proof, \S\ref{sec:ext} the exterior form and the
leakage lemma, and \S\ref{sec:open} lists the
open problems. The appendices describe the computer verifications and collect numerical values.

\subsection{Related work}\label{subsec:related}
Positivity of the Weil functional on a fixed window goes back to Yoshida~\cite{Yoshida1992}, who proved
strict positivity of the Hermitian form at $a=\tfrac12\log 2$ by an approximate finite-matrix computation.
Bombieri~\cite{Bombieri2000} gave the variational theory: the infimum on the unit sphere of $L^2(E)$ is
attained \cite[Thm.~3]{Bombieri2000}, the window minima are continuous and decreasing
\cite[Thm.~5]{Bombieri2000}, and the number of negative eigenvalues equals the number of off-line conjugate
pairs \cite[Thm.~8]{Bombieri2000}. Suzuki's work on screw functions
\cite{Suzuki2023jlms,Suzuki2023screwline,Suzuki2025hilbert} recasts the form through Kre\u{\i}n's class
$\mathcal G_a$; Theorem~\ref{thm:finZ} is a finite-rank counterpart of that picture. Groskin~\cite{Groskin2026}
proves an exact finite Guinand--Weil dictionary, and Zhu~\cite{Zhu2026} reduces the window problem to one
finite matrix by a pointwise envelope, certifies $Q(f)\ge 8.9\cdot 10^{-18}\|f\|_2^2$ for
$\supp f\subseteq[-0.8,0.8]$, and proves a doubly exponential barrier for pointwise-envelope certificates
\cite[Thm.~1.4]{Zhu2026}. Liu~\cite{Liu2026} combines a source-exact block-Schur reduction with
tail-compensation bounds and claims complete-form positivity for half-width $T\le1$ (constant $2^{-151}$)
and $T\le17/16$ (constant $2^{-49162}$); the preprint is unrefereed and its certificate is not yet
publicly reproducible. Mori \cite{Mori15} certifies $\lamW\ge2^{-2904}$ for $T\le\frac12\log15$
(his ``window of width $\log\mu$'' is our $e^{2T}=\mu$, and his form is our complete form $W_T$ with the
pole term $2\Rel(A\bar B)$) by a bounded comparison, a weighted Schur bound, an out-of-band tail and a
Legendre projection of dimension $1664$ certified in ball arithmetic, building on his earlier windows
$\frac12\log10$ and $\frac12\log12$; the deposit is unrefereed, its code and data are public, and we
reproduced the certificate check but not the analytic lemmas. The archimedean form alone was studied by Connes and Consani
\cite{ConnesConsani2021}. The Delsarte and Cohn--Elkies linear programming method \cite{CohnElkies2003} is
the structural analogue of Theorem~\ref{thm:D1}. What is new here is the dual formulation and its
identification with the explicit formula, the structural results at one window, the finite reduction, the
first-contact analysis, and an independent certification up to $\frac12\log5$ with constants
$2.5\cdot10^{-12}$ (primitive) and $10^{-18}$ (complete); larger ranges are $T\le17/16$, with constant
$2^{-49162}$, claimed in \cite{Liu2026}, and $T\le\frac12\log15$, with constant $2^{-2904}$, in
\cite{Mori15} (cited; computational part reproduced by us), which is the current certified frontier.

\subsection{Conventions}\label{subsec:conv}
$\widehat h(\tau)=\int_\R h(u)e^{i\tau u}du$; for even $h$ this is $\int h\cos(\tau u)$. For a finite even
measure $\nu$, $\nuh(\tau)=\int\cos(\tau u)\,d\nu(u)$, and for a probability measure $\mu$ on $\R$,
$\check\mu(u)=\int\cos(\tau u)\,d\mu(\tau)$. The letter $\tau$ always denotes a frequency, and $u$ a
variable dual to it. The right shift is $(V_aF)(t)=F(t-a)$, the reflection is $(JF)(t)=F(-t)$, and $Q_T$
denotes restriction to $\IT$ followed by extension by zero. We write $\mT(F)=\frac1{2\pi}\int m|\Fh|^2$
for the archimedean form, $\PiT(F)=-\sum_{n\in\NT}v_n\Rel\cF(\log n)$ for the prime form, so that
$\mathfrak a_T=\mT+\PiT$, and $\Ainf=\mT+2\Rel(A\bar B)$ for the complete archimedean form. The functions
\begin{equation}\label{eq:g}
g_0(u)=\cosh\tfrac u2\,\mathbf 1_{|u|<2T},\qquad g_1(u)=u\sinh\tfrac u2\,\mathbf 1_{|u|<2T},
\end{equation}
have transforms $\gh0,\gh1$ (Lemma~\ref{lem:gbounds}). For a parity label
$\mathrm{par}\in\{\mathrm{ev},\mathrm{odd}\}$, $\chi_{\mathrm{ev}}=\cosh\frac t2\,\mathbf 1_{\IT}$ and
$\chi_{\mathrm{odd}}=\sinh\frac t2\,\mathbf 1_{\IT}$. Letters not listed here are local to the proof or
the subsection in which they occur.

Statements whose proof includes a computer verification in interval or ball arithmetic say so; the
methods are in Appendix~\ref{app:encl}. Galerkin eigenvalues are upper bounds for bottoms and are quoted
only as numerical values, never as inputs to a proof (Appendix~\ref{app:num}). The explicit formula is used
in the form \cite[Thm.~8.4]{Paper45}, i.e.\ Weil's formula \cite{Weil52}, \cite[Thm.~5.12]{IK}, in the
normalization above.


\section{Preliminaries: the archimedean multiplier}\label{sec:levy}

\begin{lemma}[the L\'evy density]\label{lem:K}
For $u>0$ put $K(u)=\dfrac{2e^{3u/2}}{e^{2u}-1}=\dfrac{1}{2\sinh(u/2)}+\dfrac{1}{2\cosh(u/2)}
=2\sum_{k\ge0}e^{-(2k+\frac12)u}$. Then $K$ is positive, completely monotone,
$K(u)=u^{-1}+O(1)$ as $u\downarrow0$, $K(u)=2e^{-u/2}(1+O(e^{-u}))$ as $u\to\infty$, and
\begin{equation}\label{eq:m}
\begin{split}
m(\tau)&=\Rel\psi\!\left(\tfrac14+\tfrac{i\tau}2\right)-\log\pi
= C_0+\int_0^\infty(1-\cos\tau u)\,K(u)\,du,\\
C_0&=\psi(\tfrac14)-\log\pi=-5.3721834\ldots
\end{split}
\end{equation}
\end{lemma}
\begin{proof}
The three expressions for $K$ agree: $\frac{1}{2\sinh(u/2)}+\frac1{2\cosh(u/2)}
=\frac{\cosh(u/2)+\sinh(u/2)}{\sinh u}=\frac{2e^{u/2}}{e^u-e^{-u}}=\frac{2e^{3u/2}}{e^{2u}-1}$, and
$\frac{2e^{3u/2}}{e^{2u}-1}=2e^{-u/2}\sum_{k\ge0}e^{-2ku}$. Complete monotonicity is that of a
convergent series of decaying exponentials with positive coefficients. For \eqref{eq:m} use
$\psi(z)=-\gamma+\int_0^\infty\frac{e^{-t}-e^{-zt}}{1-e^{-t}}dt$ ($\Rel z>0$): with $z=\frac14+\frac{i\tau}2$,
\[
\Rel\psi(z)-\psi(\tfrac14)=\int_0^\infty\frac{e^{-t/4}\big(1-\cos(\tau t/2)\big)}{1-e^{-t}}\,dt
=\int_0^\infty\frac{2e^{-u/2}(1-\cos\tau u)}{1-e^{-2u}}\,du\qquad(t=2u),
\]
and $\frac{2e^{-u/2}}{1-e^{-2u}}=\frac{2e^{3u/2}}{e^{2u}-1}=K(u)$. The integral converges since
$1-\cos\tau u=O(u^2)$ at $0$ and $K=O(e^{-u/2})$ at $\infty$.
\end{proof}

\begin{theorem}[complete Bernstein structure]\label{thm:S}
Put $a_k=k+\frac14$ and $\varphi(x)=\sum_{k\ge0}\dfrac{x}{a_k(x+4a_k^2)}$ for $x\ge0$. Then
\emph{(i)} $\Rel\psi(\frac14+\frac{i\tau}2)-\psi(\frac14)=\varphi(\tau^2)$ and $\varphi$ is a
complete Bernstein function; \emph{(ii)} the L\'evy measure of the negative-definite function
$\tau\mapsto\varphi(\tau^2)$ on $\R$ is $\frac12K(|u|)du$, which is completely monotone on each
half-line, so the associated L\'evy process is a pure-jump subordinate Brownian motion
$B_{S_t}$ with a special subordinator $S$ of Laplace exponent $\varphi$; in particular it has no
Gaussian component; \emph{(iii)} $\int_{2T}^\infty K(u)\,du=\log\coth\frac T2+\pi-2\arctan e^{T}$.
\end{theorem}
\begin{proof}
(i) From $\psi(z)=-\gamma+\sum_{k\ge0}\big[\frac1{k+1}-\frac1{k+z}\big]$,
$\Rel\big[\frac1{a_k}-\frac1{a_k+i\tau/2}\big]=\frac1{a_k}-\frac{a_k}{a_k^2+\tau^2/4}
=\frac{\tau^2}{a_k(4a_k^2+\tau^2)}$, which is the stated series at $x=\tau^2$. Each summand
$x\mapsto a_k^{-1}\,x/(x+4a_k^2)$ is a complete Bernstein function (Stieltjes representation with
a point mass at $4a_k^2$), the summands are $O(k^{-3})$ locally uniformly, and the complete
Bernstein functions form a convex cone closed under pointwise limits
\cite[Ch.~6--7]{SSV2012}. (ii) By Lemma~\ref{lem:K}, $\varphi(\tau^2)=\int_\R(1-\cos\tau u)\,
\frac12K(|u|)du$, so the L\'evy measure is $\frac12K(|u|)du$; it is completely monotone on
$(0,\infty)$ and $\varphi(x)/x\to0$, so there is no Gaussian part; the representation as a
subordinate Brownian motion with a special subordinator is \cite[Thm.~13.10, Cor.~13.12]{SSV2012}.
(iii) $\int\frac{du}{2\sinh(u/2)}=\log\tanh\frac u4$ and $\int\frac{du}{2\cosh(u/2)}=2\arctan e^{u/2}$.
\end{proof}

\begin{theorem}[the form as a L\'evy energy]\label{thm:A}
For every $F\in L^2(\IT)$, with $\Lev(\Phi):=\int_0^\infty(\Phi(0)-\Rel\Phi(u))K(u)\,du\in[0,\infty]$
(for complex $F$ the autocorrelation is Hermitian and only its even part $\Rel\cF$ enters; the
prime term below is likewise $\Rel\cF(\log n)$, and from \S\ref{sec:duality} on all measures are even),
\begin{equation}\label{eq:Aform}
\mathfrak a_T(F)=C_0\|F\|^2+\Lev(\cF)-\sum_{n\in\NT}v_n\cF(\log n),
\end{equation}
both sides being finite exactly when $\int\log(2+|\tau|)|\Fh|^2<\infty$. Consequently, with
$\Kw=\{\cF:F\in L^2(\IT),\|F\|=1\}$ and the primitive class
\begin{equation}\label{eq:Kprim}
\Kp=\Big\{\Phi\in\Kw:\ \int\Phi(u)\cosh\tfrac u2\,du=0,\ \int\Phi(u)\,u\sinh\tfrac u2\,du=0\Big\},
\end{equation}
one has
\begin{align*}
\lamP&:=\inf_{\Phi\in\Kp}\Big[C_0+\Lev(\Phi)-\sum_{n\in\NT}v_n\Phi(\log n)\Big]=C_0-\Lambda_T,\\
\Lambda_T&:=\sup_{\Phi\in\Kp}\Big[\sum_{n\in\NT}v_n\Phi(\log n)-\Lev(\Phi)\Big],
\end{align*}
and $\lamP\le\lamA$.
\end{theorem}
\begin{proof}
By \eqref{eq:m} and Tonelli (both $1-\cos\tau u$ and $K$ are nonnegative),
$\frac1{2\pi}\int m|\Fh|^2=C_0\|F\|^2+\int_0^\infty K(u)\Big[\frac1{2\pi}\int(1-\cos\tau u)|\Fh|^2d\tau\Big]du$,
and $\frac1{2\pi}\int\cos(\tau u)|\Fh(\tau)|^2d\tau=\cF(u)$ by Plancherel; likewise for the
prime terms. Finiteness: $\int_0^\infty(1-\cos\tau u)K(u)du=\log|\tau|+O(1)$ as $|\tau|\to\infty$
by Lemma~\ref{lem:K}, since $\int_0^1(1-\cos\tau u)\frac{du}u=\log|\tau|+O(1)$. For the last
statements: the two integrals in \eqref{eq:Kprim} are $\Rel(A\bar B)$ and $\Rel(A'\bar B-A\bar B')$
with $A'=\int tF(t)e^{t/2}dt$, $B'=\int tF(t)e^{-t/2}dt$ (expand $\cosh\frac{t-s}2$ and
$(t-s)\sinh\frac{t-s}2$ in exponentials), so $\{\cF:F\in\PT,\|F\|=1\}\subseteq\Kp$, whence
$\lamP\le\lamA$. In terms of $\widehat\Phi$, an entire function of exponential type $2T$ with
$\widehat\Phi(i/2)=\overline{\widehat\Phi(-i/2)}$, the two conditions say that
$\Rel\widehat\Phi(i/2)=0$ and $\Rel\big[i(\widehat\Phi'(i/2)-\widehat\Phi'(-i/2))\big]=0$; when
$F\in\PT$ one has the stronger $\widehat\Phi(\pm i/2)=\widehat\Phi'(\pm i/2)=0$, i.e.\ $\widehat\Phi$
vanishes to order $\ge2$ at $\pm i/2$ (Remark~\ref{rem:Kprim}).
\end{proof}

\begin{remark}\label{rem:Kprim}
$\Kp$ is slightly larger than the set of primitive autocorrelations: for real $F$ its two
conditions read $AB=0$ and $A'B-AB'=0$, i.e.\ $F\in\PT$, or $A=A'=0$, or $B=B'=0$. Every
certificate below bounds $\lamP$, hence $\lamA$; the equality $\lamP=\lamA$ is not needed
anywhere and is not asserted. Since $W_T(F)=\mathfrak a_T(F)$ whenever $A\bar B$ is purely
imaginary, also $\lamW\le\lamP$. Whether $\lamP=\lamW$, i.e.\ whether the optimal multipliers of the
primitive problem are the polar ones, is characterized in Theorem~\ref{thm:mult}.
\end{remark}

\begin{corollary}[monotonicity; upper bound]\label{cor:mono}
With the contact set held fixed, $T\mapsto\Lambda_T$ is nondecreasing and $T\mapsto\lamP$
nonincreasing. Moreover
$\lamA\le C_0+\Aconst+\varphi\big((3\pi/2T)^2\big)$, which equals $1.291$ at $T=\frac12\log3$ and
$2.327$ at $T=\log 2$.
\end{corollary}
\begin{proof}
Monotonicity: $\Kp$ grows with $T$ and $\Lev$ does not depend on $T$. Upper bound: the form
$\Lev(\cF)$ is the Dirichlet form of the L\'evy process of Theorem~\ref{thm:S} killed on
leaving $\IT$; its eigenvalues are bounded by those of the subordinate killed Brownian motion,
$\lambda_n\le\varphi((n\pi/2T)^2)$ \cite[Prop.~2.3]{SongVondracek2003}, the Brownian motion being
normalized with characteristic exponent $\xi^2$; $\PT$ has codimension
$2$, so the bottom of $\Lev$ on $\PT$ is at most $\lambda_3$; and $-\sum v_n\cF(\log n)\le\Aconst\|F\|^2$.
\end{proof}

The following elementary bounds on $m$ and on the prime comb are used throughout.

\begin{lemma}[elementary bounds]\label{lem:mbounds}
\emph{(i)} $m$ is increasing on $[0,\infty)$. \emph{(ii)} For every $\tau>0$,
\[
\log\frac{\tau}{2\pi}-\frac1\tau-\frac1{2\tau^2}\ \le\ m(\tau)\ \le\ \log\frac{\tau}{2\pi}+\frac1\tau+\frac12\log\Big(1+\frac1{4\tau^2}\Big),
\]
and the lower bound, which we denote by $m_{\mathrm{lo}}(\tau)$, is increasing in $\tau$. \emph{(iii)}
\[
-\sum_{n\in\NT}v_n\cos(\tau\log n)\ge-\Aconst,\qquad\text{with}\quad
\sup_\tau\sum_{n\in\NT}v_n\cos(\tau\log n)=\Aconst.
\]
\end{lemma}
\begin{proof}
(i) In $\Rel\psi(\frac14+iy)=-\gamma+\sum_{n\ge0}\big[\frac1{n+1}-\frac{n+\frac14}{(n+\frac14)^2+y^2}\big]$
every summand increases with $|y|$. (ii) Use $\psi(z)=\log z-\frac1{2z}-\int_0^\infty g(t)e^{-zt}dt$
for $\Rel z>0$, $g(t)=\frac1{e^t-1}-\frac1t+\frac12$ \cite[5.9.13]{DLMF}, at $z=\frac14+iy$, $y=\tau/2$.
Here $g(0^+)=0$, $g'(t)=t^{-2}-\frac14\sinh^{-2}(t/2)\in(0,\infty)$ and $g(\infty)=\frac12$, so one
integration by parts gives $|\int_0^\infty g e^{-zt}dt|=|z|^{-1}|\int_0^\infty g'e^{-zt}dt|\le\frac1{2|z|}\le\frac1{2y}$.
Also $\Rel\log z=\frac12\log(\frac1{16}+y^2)\in[\log y,\log y+\frac12\log(1+\frac1{16y^2})]$ and
$\Rel\frac1{2z}=\frac{1/8}{1/16+y^2}\in[0,\frac1{8y^2}]$. Substituting $y=\tau/2$ and $-\log\pi$
gives (ii); the derivative of the lower bound is $\tau^{-1}+\tau^{-2}+\tau^{-3}>0$. (iii) The
first is trivial. For the second, the numbers $\log p$, $p$ prime, are linearly independent over
$\Q$, so by Kronecker's theorem the vector $(\tau\log p)_{p<e^{2T}}$ is dense modulo $2\pi$; at
$\tau$ where all $\cos(\tau\log p)>1-\varepsilon$ also $\cos(\tau k\log p)>1-k^2\varepsilon$, so the
sum over all $n\in\NT$ exceeds $\Aconst-O(\varepsilon)$.
\end{proof}

The hypothesis that the contact set is fixed can be removed from Corollary~\ref{cor:mono}, for
all three bottoms at once.

\begin{lemma}[monotonicity across contact thresholds]\label{lem:monoT}
The functions $T\mapsto\lamA$, $T\mapsto\lamP$ and $T\mapsto\lamW$ are nonincreasing on $(0,\infty)$,
including at the points $T=\frac12\log n$ where a contact enters.
\end{lemma}
\begin{proof}
Let $T<T'$ and extend $F\in L^2(\IT)$ by zero to $I_{T'}$. The norm, the Tate moments $A(F),B(F)$ and
the autocorrelation $\cF$ are unchanged, and $\cF$ vanishes on $|u|\ge2T$. A contact
$n\in\mathcal N_{T'}\setminus\NT$ has $\log n\ge2T$, so $\cF(\log n)=0$, and by \eqref{eq:Aform}
$\mathfrak a_{T'}(F)=\mathfrak a_T(F)$. Hence $\PT\subset\mathcal P_{T'}$ isometrically with the same
form values, which gives the statement for $\lamA$; the same computation with $W_T=\mathfrak a_T+2\Rel(A\bar B)$
gives it for $\lamW$, and for $\lamP$ the class $\Kp$ grows with $T$, $\Lev$ is $T$-free and the new
contacts contribute $\Phi(\log n)=0$.
\end{proof}

For the complete form, continuity is \cite[Thm.~5]{Bombieri2000}:
Bombieri proves that the even and odd minima $\mu_\pm$ of the complete functional
on $L^2([M^{-1},M])$ are continuous and decreasing. Here $W_T$ splits into its values on the even and
odd parts of $F$: the symbol is even, and $\Rel(A\bar B)=|C|^2-|S|^2$ with $C=\int F\cosh\frac t2$, which
sees only the even part, and $S=\int F\sinh\frac t2$, which sees only the odd part. So, after the change of
variables $x=e^t$, $\lamW=\min(\mu_+,\mu_-)$ at $M=e^T$. A direct proof in the present normalization follows.

\begin{lemma}[continuity of the complete bottom]\label{lem:contT}
The function $T\mapsto\lamW$ is continuous on $(0,\infty)$.
\end{lemma}
\begin{proof}
Fix $T$ and let $r\in[\frac12,2]$. The dilation $U_rF(t)=r^{-1/2}F(t/r)$ is unitary from $L^2(\IT)$ onto
$L^2(I_{rT})$, with $\widehat{U_rF}(\tau)=r^{1/2}\Fh(r\tau)$ and $c_{U_rF}(u)=\cF(u/r)$. It maps the form
domain onto the form domain, since $\log(2+|s|/r)\le\log(2+|s|)+\log2$. Hence
$\lambda^{\mathrm W}_{rT}=\inf\{W_{rT}(U_rF):F\text{ in the form domain},\ \|F\|=1\}$. By \eqref{eq:form}
and \eqref{eq:W},
\[
W_{rT}(U_rF)=\frac1{2\pi}\int_\R m(s/r)|\Fh(s)|^2ds-\sum_{2\le n<e^{4T}}v_n\Rel\cF(\tfrac{\log n}r)
+2\Rel\big(A(U_rF)\overline{B(U_rF)}\big),
\]
where the comb may be summed over $n<e^{4T}$ for both windows, because $\cF$ vanishes on $|u|\ge2T$. This
also takes care of contacts that enter or leave. We compare the three terms with those of $W_T(F)$.
\emph{(i) Archimedean.} $m$ is even and $m'(x)=-\frac12\Iml\psi'(\frac14+\frac{ix}2)$ with
$\psi'(z)=O(1/|z|)$ on $\Rel z=\frac14$. So $M_1:=\sup_x|xm'(x)|<\infty$ and
$|m(s/r)-m(s)|\le M_1|\log r|$, and the difference of the archimedean terms is at most
$M_1|\log r|\,\|F\|^2$.
\emph{(ii) Comb.} $|\Rel\cF(u)-\Rel\cF(u')|\le\frac1{2\pi}\int\min(2,|\tau||u-u'|)|\Fh|^2d\tau$. Put
$\delta=4T|1/r-1|$ and $\|F\|_Q^2:=\frac1{2\pi}\int\log(2+|\tau|)|\Fh|^2d\tau$. Splitting at
$|\tau|=\delta^{-1/2}$, the difference of the comb terms is at most
$S_T\big(\delta^{1/2}\|F\|^2+2\|F\|_Q^2/\log(2+\delta^{-1/2})\big)$, where $S_T=\sum_{n<e^{4T}}v_n$.
\emph{(iii) Polar.} $A(U_rF)=r^{1/2}\int F(s)e^{rs/2}ds$, so
$|A(U_rF)-A(F)|\le\|F\|\,\|r^{1/2}e^{rs/2}-e^{s/2}\|_{L^2(\IT)}$, and likewise for $B$. Both moments are
bounded by $C_T\|F\|$ uniformly in $r\in[\frac12,2]$. So the difference of the polar terms is at most
$\delta_P(r)\|F\|^2$ with $\delta_P(r)\to0$ as $r\to1$.
By Lemma~\ref{lem:mbounds}(ii) there is $C_1$ with $\log(2+|\tau|)\le m(\tau)+C_1$. By
Lemma~\ref{lem:mbounds}(iii) and $|2A\bar B|\le4\sinh T\,\|F\|^2$ (Cauchy--Schwarz), this gives
$\|F\|_Q^2\le W_T(F)+C'\|F\|^2$ with $C'=C_1+\Aconst+4\sinh T$. Collecting (i)--(iii), there is
$\eta(r)\to0$ as $r\to1$ with
\[
|W_{rT}(U_rF)-W_T(F)|\le\eta(r)\big(W_T(F)+C'\|F\|^2\big)
\]
for all $F$ in the form domain, and the right side is nonnegative. Take infima over unit $F$. This gives
$(1-\eta)\lamW-\eta C'\le\lambda^{\mathrm W}_{rT}\le(1+\eta)\lamW+\eta C'$, so
$|\lambda^{\mathrm W}_{rT}-\lamW|\le\eta(r)(|\lamW|+C')\to0$.
\end{proof}

The same proof with the comb deleted gives continuity of the bottom $\laminf$ of $\Ainf$
(\S\ref{sec:archstr}). It does not apply directly to $\lamA$, because $U_r$ does not preserve $\PT$;
the correction of the two Tate moments is the content of Lemma~\ref{lem:contA} below.

\begin{lemma}[compactness]\label{lem:compact}
For every $V\in\R$ the set $\{F\in\mathcal D_T:\|F\|=1,\ \mathfrak a_T(F)\le V\}$ is relatively compact in
$L^2(\IT)$. Consequently the form $\mathfrak a_T$ restricted to $\PT\cap\mathcal D_T$ is the form of a
self-adjoint operator $H_T$ on $\PT$ with compact resolvent, $\lamA=\min\operatorname{spec}H_T$ is attained,
and the bottoms of $\mathfrak a_T$ and $W_T$ on the unit sphere of $L^2(\IT)$, and on any closed subspace, are
attained.
\end{lemma}
\begin{proof}
By Lemma~\ref{lem:mbounds}, $\sT\ge m-\Aconst$ and $m-C_0\ge0$, so for $\|F\|=1$ and
$\mathfrak a_T(F)\le V$ one has $\frac1{2\pi}\int(m-C_0)|\Fh|^2\le V+\Aconst-C_0$, hence
$\frac1{2\pi}\int_{|\tau|>R}|\Fh|^2\le(V+\Aconst-C_0)/(m(R)-C_0)\to0$ as $R\to\infty$, uniformly. For each
$R$ the map $F\mapsto\Fh\mathbf 1_{[-R,R]}$ is Hilbert--Schmidt from $L^2(\IT)$ to $L^2(\R)$ (its kernel
$e^{i\tau t}$ is square integrable on $[-R,R]\times\IT$), hence compact. A diagonal argument gives a
subsequence along which $\Fh_n$ converges in $L^2([-R,R])$ for every $R$, and with the uniform tail bound
$\Fh_n$ is Cauchy in $L^2(\R)$. The form $\mathfrak a_T$ is closed and bounded below by $C_0-\Aconst$, since its
form norm is equivalent to $\int\log(2+|\tau|)|\Fh|^2$ (Lemma~\ref{lem:mbounds}(ii)); $\PT$ is closed, and
the polar term $2\Rel(A\bar B)$ is continuous on $L^2(\IT)$. Lower semicontinuity of closed forms gives
attainment, and relative compactness of the sublevel sets is compactness of the resolvent.
\end{proof}

On $\PT$ the operator $H_T$ commutes with $J$, since $\sT$ is even and $A(JF)=B(F)$; so its eigenspaces
split into even and odd parts.

\begin{lemma}[continuity of the primitive bottom]\label{lem:contA}
The function $T\mapsto\lamA$ is continuous on $(0,\infty)$.
\end{lemma}

\begin{proof}
Fix $T$ and consider windows $T'\in[T/2,2T]$ and factors $r\in[\frac12,2]$ with $rT'\in[T/2,2T]$; all
constants depend only on $T$. Put $M_0=\mathsf A_{2T}-C_0$, so that $\mathfrak a_S+M_0\|\cdot\|^2\ge0$
for every $S\le2T$ (Lemma~\ref{lem:mbounds}).

\emph{Step 1.} Steps (i) and (ii) of the proof of Lemma~\ref{lem:contT} do not involve the polar term.
Taking $\delta=8T|1/r-1|$ and summing the comb over $n<e^{8T}$, which covers every window in
$[T/2,2T]$ and every contact that enters or leaves, they give $\eta(r)\to0$ as $r\to1$ with
$|\mathfrak a_{rT'}(U_rF)-\mathfrak a_{T'}(F)|\le\eta(r)\big(\mathfrak a_{T'}(F)+C\|F\|^2\big)$ for all
$F$ in the form domain of $\mathfrak a_{T'}$, uniformly in $T'$.

\emph{Step 2.} If $F\in\mathcal P_{T'}$, then $A(U_rF)=r^{1/2}\int F(s)\big(e^{rs/2}-e^{s/2}\big)ds$,
so $|A(U_rF)|\le\varepsilon(r)\|F\|$ with $\varepsilon(r)=r^{1/2}\|e^{rs/2}-e^{s/2}\|_{L^2(I_{2T})}\to0$,
and likewise for $B$.

\emph{Step 3.} Fix real $\phi_1,\phi_2\in C_c^\infty(I_{T/4})$ for which the matrix
$\big(A(\phi_j),B(\phi_j)\big)_{j}$ is invertible. By zero extension (proof of Lemma~\ref{lem:monoT}),
$\mathfrak a_S(\phi_i,\phi_j)$ does not depend on $S\ge T/4$. For $u=U_rF$ with $F\in\mathcal P_{T'}$,
$\|F\|=1$, let $G=c_1\phi_1+c_2\phi_2$ have the same Tate moments as $u$. Then
$v:=u-G\in\mathcal P_{rT'}$, $\|G\|\le C_2\varepsilon(r)$ and $\mathfrak a_{rT'}(G)\le C_3\varepsilon(r)^2$.

\emph{Step 4.} By the Cauchy--Schwarz inequality for the nonnegative form
$\mathfrak a_{rT'}+M_0\|\cdot\|^2$,
$|\mathfrak a_{rT'}(u,G)|\le(\mathfrak a_{rT'}(u)+M_0)^{1/2}(\mathfrak a_{rT'}(G)+M_0\|G\|^2)^{1/2}+M_0\|G\|$.
Hence, if $\mathfrak a_{T'}(F)\le V$ for a fixed $V$, then
$\mathfrak a_{rT'}(v)\le\mathfrak a_{T'}(F)+\omega(r)$ with $\omega(r)\to0$ uniformly in such $F$. Since
also $\mathfrak a_{rT'}(v)\ge-M_0\|v\|^2$ and $\big|\|v\|-1\big|\le C_2\varepsilon(r)$, the Rayleigh
quotient gives $\lambda_{\mathrm A}(rT')\le\mathfrak a_{T'}(F)+\omega'(r)$ with
$\omega'(r)\to0$.

\emph{Step 5.} By Lemma~\ref{lem:monoT}, $\lamA$ is bounded on $[T/2,2T]$ between $C_0-\mathsf A_{2T}$
and $\lambda_{\mathrm A}(T/2)$; take $V$ one more than the latter. Applying Step 4 to
near-minimizers at $T$ gives $\limsup_{r\to1}\lambda_{\mathrm A}(rT)\le\lamA$, and applying
it to near-minimizers at $rT$ with the factor $1/r$ gives the reverse inequality for the $\liminf$.
\end{proof}


\section{Primal and dual programs}\label{sec:duality}

\subsection{Primal and dual}
Fix $T$ and the weights $v=(v_n)_{n\in\NT}$ (any nonnegative weights may be used; $v_n=2\Lambda(n)/\sqrt n$
is the case of $\zeta$).

\begin{definition}\label{def:PD}
\textbf{(P)} $\ \lamP=\inf\big\{\int_\R\sT\,d\mu:\ \mu$ an even probability measure on $\R$,
$\check\mu(u)=0$ for $|u|\ge2T$, $\int\gh0\,d\mu=\int\gh1\,d\mu=0\big\}$.

\textbf{(D)} $\ D_T=\sup\big\{\inf_{\tau\in\R}\big[\sT(\tau)+\nuh(\tau)+\beta_0\gh0(\tau)+\beta_1\gh1(\tau)\big]:\
\nu$ a finite even signed Borel measure with $\supp\nu\subseteq\{|u|\ge2T\}$, $\beta\in\R^2\big\}$.
A triple $(\nu,\beta_0,\beta_1)$ with $f:=\sT+\nuh+\beta\cdot\widehat g\ge c$ on $\R$ is a
\emph{certificate of value $c$}.
\end{definition}

\begin{proposition}[the primal is the eigenvalue problem; weak duality]\label{prop:weak}
The value of \textup{(P)} is the $\lamP$ of Theorem~\ref{thm:A}. Every certificate of value $c$
satisfies $c\le\lamP\le\lamA$.
\end{proposition}
\begin{proof}
If $F$ is a unit vector, $\mu_F=|\Fh|^2d\tau/2\pi$ is feasible ($\check\mu_F=\cF$) with
$\int\sT d\mu_F=\mathfrak a_T(F)$, and $\int\gh j\,d\mu_F=\int g_j\cF$. Conversely if $\mu$ is
feasible, $\Phi=\check\mu$ is continuous, positive definite, $\Phi(0)=1$, supported in $[-2T,2T]$,
so by the Boas--Kac theorem \cite{BoasKac1945} $\Phi=\cF$ for some $F\in L^2(\IT)$, $\|F\|=1$, and
$\Phi\in\Kp$; thus the feasible set of (P) is exactly $\Kp$ read through $\mu=|\Fh|^2d\tau/2\pi$.
Weak duality: for feasible $\mu$, $\int\nuh\,d\mu=\int\check\mu\,d\nu=0$ because $\check\mu=\cF$
vanishes identically on the closed set $\{|u|\ge2T\}$ ($\cF(\pm2T)=0$ by continuity, so an atom
of $\nu$ at $\pm2T$ is harmless), and $\int\gh j\,d\mu=0$; hence
$\int\sT d\mu=\int f\,d\mu\ge c$.
\end{proof}

\begin{lemma}[bounds on the multipliers]\label{lem:gbounds}
$\gh0(\tau)=2\int_0^{2T}\cosh\frac u2\cos\tau u\,du$ and $\gh1(\tau)=2\int_0^{2T}u\sinh\frac u2\cos\tau u\,du$
are even, real, bounded, tend to $0$ at infinity, and for $\tau>0$
\[
|\gh0(\tau)|\le k_0(\tau):=\frac{2\cosh T}{\tau}+\frac{\sinh T}{\tau^2},\qquad
|\gh1(\tau)|\le k_1(\tau):=\frac{8T\sinh T}{\tau},
\]
both bounds decreasing in $\tau$. Moreover $\gh0(\tau)=\frac{2\cosh T\sin2T\tau}{\tau}+O(\tau^{-2})$,
and for any $g=y_0g_0+y_1g_1$, $\widehat g(\tau)=\frac{2g(2T^-)\sin2T\tau}{\tau}+\frac{2g'(2T^-)\cos2T\tau}{\tau^2}+O(\tau^{-3})$.
\end{lemma}
\begin{proof}
Integrate by parts: $\gh0=\frac{2\cosh T\sin2T\tau}\tau-\frac1\tau\int_0^{2T}\sinh\frac u2\sin\tau u\,du$,
and once more, $\frac1\tau\int_0^{2T}\sinh\frac u2\sin\tau u\,du=-\frac{\sinh T\cos2T\tau}{\tau^2}+\frac{\gh0(\tau)}{4\tau^2}$;
solving, $|\gh0|(1+\frac1{4\tau^2})\le\frac{2\cosh T}\tau+\frac{\sinh T}{\tau^2}$. For $\gh1$, one
integration by parts gives $\gh1=\frac{4T\sinh T\sin2T\tau}{\tau}-\frac2\tau\int_0^{2T}\big(\sinh\frac u2+\frac u2\cosh\frac u2\big)\sin\tau u\,du$,
and the integrand is the derivative of $u\sinh\frac u2$, whose integral over $[0,2T]$ is $2T\sinh T$.
The asymptotic expansion is two integrations by parts for a function real-analytic on $[0,2T]$
with $g'(0)=0$.
\end{proof}

\subsection{Strong duality}

\begin{theorem}[strong duality; finitely atomic dual sequences]\label{thm:D1}
$D_T=\lamP$. Moreover, for every $\varepsilon>0$ there are finitely many atoms
$u_1,\dots,u_N\ge2T$, reals $a_1,\dots,a_N$ and $\beta\in\R^2$ such that
$\sT(\tau)+\sum_j2a_j\cos(u_j\tau)+\beta_0\gh0(\tau)+\beta_1\gh1(\tau)>\lamP-\varepsilon$ for all $\tau\in\R$.
In particular, for every $T$ with $\lamP>0$, $\lamA>0$ is provable by a finite certificate.
\end{theorem}
\begin{proof}
\emph{(i) Finite relaxations.} For a finite set $U=\{u_1<\dots<u_N\}\subset[2T,\infty)$ let
$P(U)=\inf\{\int\sT d\mu:\ \mu\ \text{even probability},\ \check\mu(u_i)=0\ \forall i,\ \int\gh j d\mu=0\}$
and $D(U)$ the supremum in (D) restricted to $\nu$ supported on $\pm U$. Then $P(U)\le\lamP$
(the feasible set of (P) is contained in that of $P(U)$) and $D(U)\le D_T\le\lamP$.

\emph{(ii) $P(U)=D(U)$.} Write $f_0\equiv1$, $f_i(\tau)=\cos(u_i\tau)$ ($1\le i\le N$),
$f_{N+1}=\gh0$, $f_{N+2}=\gh1$, $b=(1,0,\dots,0)\in\R^{N+3}$, and let
$C=\{\sum_kt_kf(\tau_k):\ t_k\ge0\}$ be the convex cone generated by the moment curve
$\tau\mapsto f(\tau)=(f_0,\dots,f_{N+2})(\tau)$. The value function
$w(b')=\inf\{\int\sT d\mu:\ \mu\ge0,\ \int f\,d\mu=b'\}$ is convex, $>-\infty$ everywhere
($\sT\ge-\Aconst-|C_0|$ and $\int f_0d\mu=b'_0$ bounds the mass) and finite on $C$ (finite atomic
$\mu$ have finite value). Its Legendre biconjugate at $b$ is the Lagrangian dual
$w^{**}(b)=\sup\{y_0:\ \sT-\sum_ky_kf_k\ge0\text{ on }\R\}$, which is $D(U)$ after
$a_i=-y_i/2$, $\beta_j=-y_{N+1+j}$. A proper convex function agrees with its biconjugate at every
interior point of its domain \cite[Thm.~7.4, Cor.~12.1.1]{Rockafellar}; so it suffices to show
$b\in\operatorname{int}C$.

Suppose not. Since $C$ is a convex cone containing $0$, there is $y\ne0$ with
$\langle y,z\rangle\ge\langle y,b\rangle$ for all $z\in C$; taking $z=0$ gives $y_0\le0$, and
taking $z=tf(\tau)$, $t\to\infty$, gives
$G(\tau):=y_0+\sum_iy_i\cos(u_i\tau)+y_{N+1}\gh0(\tau)+y_{N+2}\gh1(\tau)\ge0$ for all $\tau$.
The Ces\`aro mean $X^{-1}\int_0^XG\to y_0$ (each $\cos(u_i\tau)$ has mean $0$ because $u_i>0$,
and $\gh j\to0$ by Lemma~\ref{lem:gbounds}); as $G\ge0$ this forces $y_0=0$, and then
$X^{-1}\int_0^X|G|\to0$. Testing against the bounded functions $\cos(u_i\tau)$ gives
$y_i/2=\lim X^{-1}\int_0^XG\cos(u_i\tau)\,d\tau=0$ for each $i$. There remains
$\widehat g\ge0$ with $g=y_{N+1}g_0+y_{N+2}g_1$. By the expansion in Lemma~\ref{lem:gbounds},
$\widehat g\ge0$ forces $g(2T^-)=0$ and then $g'(2T^-)=0$; but
$g(2T^-)=y_{N+1}\cosh T+2T\,y_{N+2}\sinh T$ and $g'(2T^-)=\frac12y_{N+1}\sinh T+y_{N+2}(\sinh T+T\cosh T)$,
whose determinant is $\sinh T\cosh T+T>0$. So $y=0$, a contradiction; hence $b\in\operatorname{int}C$.

\emph{(iii) The limit.} Choose finite sets $U_1\subset U_2\subset\cdots$ with $\bigcup U_n$ dense in
$[2T,\infty)$; $P(U_n)$ is nondecreasing and $\le\lamP$. Pick feasible $\mu_n$ with
$\int\sT d\mu_n\le P(U_n)+1/n$. Since $\sT\to+\infty$ and is bounded below, $(\mu_n)$ is tight;
let $\mu$ be a weak-$*$ limit point, a probability measure. For $u\in U_N$ and $n\ge N$,
$\check\mu_n(u)=0$ and $\cos(u\,\cdot)$ is bounded continuous, so $\check\mu(u)=0$; $\check\mu$
is continuous, so $\check\mu\equiv0$ on $[2T,\infty)$; likewise $\int\gh jd\mu=0$. Thus $\mu$ is
feasible for (P), and by lower semicontinuity of $\mu\mapsto\int\sT d\mu$ under weak convergence
($\sT$ continuous and bounded below), $\lamP\le\int\sT d\mu\le\liminf\int\sT d\mu_n=\lim P(U_n)\le\lamP$.
Hence $D_T\ge\sup_nD(U_n)=\sup_nP(U_n)=\lamP\ge D_T$, and the atomic statement follows.
\end{proof}

\begin{remark}\label{rem:D1complete}
Theorem~\ref{thm:D1} says that, at a fixed window, $\lamP\ge0$ holds if and only if for every
$\varepsilon>0$ there are an atomic $\nu$ on $\{|u|\ge2T\}$ and $\beta\in\R^2$ with
$\sT+\nuh+\beta\cdot\widehat g\ge-\varepsilon$ on $\R$; and $\lamP\ge0$ implies $\lamA\ge0$. The same
proof, with the two $\widehat g$-constraints deleted, gives strong duality for the complete form:
$\lamW=\sup\{\inf_\tau[\sT+2\gh0+\nuh]\}$ over the same class of $\nu$ (Proposition~\ref{prop:P}
explains the constant $2$).
\end{remark}

\subsection{Non-attainment}

\begin{theorem}[non-attainment]\label{thm:O}
Let $\nu$ be any finite even signed measure with $\supp\nu\subseteq\{|u|\ge2T\}$ and let
$h_0,\dots,h_J$ be finitely many continuous functions on $[-2T,2T]$ that annihilate $\Kp$
(for instance $g_0,g_1$). Then for every $\beta\in\R^{J+1}$
\[
\inf_{\tau\in\R}\Big[\sT(\tau)+\nuh(\tau)+\sum_j\beta_j\widehat h_j(\tau)\Big]<\lamP\quad\text{strictly.}
\]
The statement holds for every finite weight family, including $v=0$.
\end{theorem}
\begin{proof}
A minimizer $\Phi_\star=c_{F_\star}\in\Kp$ of the primal exists: the sublevel sets of
$\mathfrak a_T$ on the unit sphere of $L^2(\IT)$ are compact (Lemma~\ref{lem:compact}),
and, in the language of measures, the feasible set of (P) is tight and weak-$*$ closed and $\mu\mapsto\int\sT d\mu$ is
weak-$*$ lower semicontinuous, exactly as in step (iii) of the proof of Theorem~\ref{thm:D1}. Put $\varrho=\sT+\nuh+\sum\beta_j\widehat h_j-\lamP$, a
continuous function ($\nu$ finite). By the annihilation properties,
$\frac1{2\pi}\int\varrho|\Fh_\star|^2=\mathfrak a_T(F_\star)-\lamP=0$. If $\varrho\ge0$ then $\varrho=0$ a.e.\ on
$\{\Fh_\star\ne0\}$, whose complement is the discrete zero set of a nonzero entire function of
exponential type; so $\varrho\equiv0$ on $\R$, and by Fourier inversion, as distributions in $u$,
\[
\check\sT+\nu+\sum_j\beta_jh_j=\lamP\,\delta_0 .
\]
Restrict to the open set $0<|u|<2T$, where $\nu=0$: by \eqref{eq:m}, $\check m=-\frac12K(|u|)$
there, and the prime part of $\check\sT$ is $-\sum_{n\in\NT}\frac{v_n}2(\delta_{\log n}+\delta_{-\log n})$;
hence $\sum_j\beta_jh_j(u)=\frac12K(|u|)+\sum_{n\in\NT}\frac{v_n}2(\delta_{\log n}+\delta_{-\log n})$
on $0<|u|<2T$. The left side is a bounded continuous function; the right side is unbounded as
$u\to0$ ($K(u)\sim1/u$) and carries atoms as soon as $\NT\neq\emptyset$. Contradiction.
\end{proof}

\begin{remark}\label{rem:OvsB}
Theorem~\ref{thm:O} forbids attainment, not certificates: whenever $\lamP>0$ any dual value in
$(0,\lamP)$ suffices, and Theorem~\ref{thm:D1} says such values are reached by finite atomic
$\nu$. Every maximizing sequence of the dual must therefore degenerate (its total variation is
unbounded; Corollary~\ref{cor:T2} below quantifies the cost). On the primal side the infimum is
attained \cite[Thm.~3]{Bombieri2000}. The proof shows the reason for the asymmetry: the
singularity $K(u)\sim1/u$ of the archimedean L\'evy density at $u=0$ and the prime atoms lie inside
the window, where the dual measure is not allowed to have support.
\end{remark}

\subsection{The Euler--Lagrange equation}

\begin{theorem}[finite-interval Wiener--Hopf equation]\label{thm:W}
Let $\mathsf S_T$ denote the convolution operator on $\R$ with symbol $\sT$,
$\widehat{\mathsf S_TF}=\sT\Fh$, acting on $L^2(\IT)$ extended by zero and taking values in tempered
distributions. A unit vector $F\in\PT\cap\mathcal D_T$ is a critical point of $\mathfrak a_T$ on the unit
sphere of $\PT$ with critical value $\lambda$ if and only if there are $\alpha,\beta\in\C$ (real for real
$F$) with
\[
(\mathsf S_T-\lambda)F=\alpha\cosh\tfrac t2+\beta\sinh\tfrac t2\qquad\text{on }\IT,
\]
in the sense of distributions on $\IT$. The symbol $\sT-\lambda$ is real, even, tends to
$+\infty$, and changes sign for every $\lambda>C_0-\Aconst$; in that range it is not elliptic, and no
Wiener--Hopf index is defined for it.
\end{theorem}
\begin{proof}
For $G\in C_c^\infty(\IT)$, $\langle(\mathsf S_T-\lambda)F,G\rangle=\frac1{2\pi}\int(\sT-\lambda)\Fh\,\overline{\widehat G}$
is the polarization of $\mathfrak a_T-\lambda\|\cdot\|^2$. Criticality on the unit sphere of $\PT$
with value $\lambda$ means that this vanishes for all $G$ in $\PT\cap\mathcal D_T$; since
$C_c^\infty(\IT)\cap\PT$ is dense in $\PT\cap\mathcal D_T$ for the form norm (mollify and correct the two
Tate moments by a fixed pair of $C_c^\infty$ functions), it is equivalent to require it for
$G\in C_c^\infty(\IT)$ with $\int Ge^{t/2}=\int Ge^{-t/2}=0$. This is a subspace of codimension $2$
of $C_c^\infty(\IT)$, and a distribution on $\IT$ that annihilates the common kernel of two continuous
linear functionals is a linear combination of them; the functionals are integration against
$e^{\pm t/2}$, i.e.\ against $\cosh\frac t2\pm\sinh\frac t2$. Conversely, if $(\mathsf S_T-\lambda)F$
is such a combination on $\IT$, it annihilates every $G$ in that subspace, so $F$ is critical. Finally
$\sT(0)=C_0-\Aconst$ and $\sT\to+\infty$ (Lemma~\ref{lem:mbounds}).
\end{proof}

\subsection{The polar multiplier}

\begin{lemma}[moment interior]\label{lem:W}
For $F\in L^2(\IT)$ real put $Q_0(F)=\int\cF\cosh\frac u2=A B$ and
$Q_1(F)=\int\cF\,u\sinh\frac u2=A'B-AB'$. The image of the set of unit vectors under
$F\mapsto(Q_0,Q_1)$ contains a neighborhood of $0$ in $\R^2$, realised by $C_c^\infty(\IT)$
functions; equivalently, $0$ is an interior point of the (convex) moment image
$\mathfrak m(\mu)=(\int\gh0d\mu,\int\gh1d\mu)$ of the feasible set of (P) without the two
constraints, and the interior is realised by measures of finite value.
\end{lemma}
\begin{proof}
The formulas for $Q_0,Q_1$ follow from expanding $\cosh\frac{t-s}2$ and $(t-s)\sinh\frac{t-s}2$
in exponentials. The functions $e^{t/2},e^{-t/2},te^{t/2},te^{-t/2}$ are linearly independent
on $\IT$, and a nontrivial combination is real-analytic, so it cannot annihilate
$C_c^\infty(\IT)$: the linear map $F\mapsto(A,B,A',B')$ is onto $\R^4$ from $C_c^\infty(\IT)$.
If $0$ were not interior, some $(t_0,t_1)\ne0$ would satisfy $t_0Q_0(F)+t_1Q_1(F)\ge0$ for all
$F$ (normalization only rescales $Q_j$ by a positive factor). Take $(A,B,A',B')=(1,0,0,\mp1)$:
$Q_0=0$, $Q_1=\pm1$, so $t_1=0$; take $(1,\pm1,0,0)$: $Q_1=0$, $Q_0=\pm1$, so $t_0=0$. A finite
convex combination of the four witnesses covers a ball, and each has finite value because
$\Fh$ is Schwartz. For the image of the unit sphere itself: a function with
$(A,B,A',B')=(1,x,0,-y)$ has $(Q_0,Q_1)=(x,y)$; adding an element of the kernel of the moment map
does not change $(Q_0,Q_1)$ and makes the norm any number above the minimal one, and
$(Q_0,Q_1)(F/\|F\|)=(Q_0,Q_1)(F)/\|F\|^2$. Since the minimal norm is bounded on $x^2+y^2=1$, every
point $t(x,y)$ with $0<t\le t_0$, $t_0>0$ uniform, is attained on the unit sphere, and $(0,0)$ is
attained on $\PT$.
\end{proof}

\begin{proposition}[the polar multiplier]\label{prop:P}
For every $F\in L^2(\IT)$ in the form domain,
$\frac1{2\pi}\int\gh0|\Fh|^2=\Rel\big(A(F)\overline{B(F)}\big)$ and hence
\[
\frac1{2\pi}\int\big(\sT+2\gh0\big)|\Fh|^2\,d\tau=W_T(F);
\]
if moreover $F$ is real and $\sum_\rho|\Phi_F(\rho)\Phi_F(1-\rho)|<\infty$ (in particular if
$F\in C_c^\infty(\IT)$ is real), then $W_T(F)=\sum_\rho\Phi_F(\rho)\Phi_F(1-\rho)$
(Lemma~\ref{lem:roughEF} below).
Consequently a certificate with $(\beta_0,\beta_1)=(2,0)$ and value $c$ proves $W_T\ge c$ on all
of $L^2(\IT)$, and $\beta_0=2$ is the unique value of $\beta_0$ with this property.
\end{proposition}
\begin{proof}
$\frac1{2\pi}\int\gh0|\Fh|^2=\int g_0\cF=\iint\cosh\frac{t-s}2F(t)\overline{F(s)}$ and
$\cosh\frac{t-s}2=\frac12e^{t/2}e^{-s/2}+\frac12e^{-t/2}e^{s/2}$ gives $\frac12(A\bar B+B\bar A)$.
The identification of $W_T$ with the zero sum is the explicit formula for the autocorrelation
$\cF$, \cite{Weil52,Bombieri2000}, \cite[Thm.~8.4]{Paper45}, whose polar term is
$2\Phi_F(0)\Phi_F(1)=2\Rel(A\bar B)$ for the real forms considered here. Uniqueness: for
$\beta_0\ne2$, $\frac1{2\pi}\int(\sT+\beta_0\gh0)|\Fh|^2=W_T(F)+(\beta_0-2)\Rel(A\bar B)$, and
$\Rel(A\bar B)$ takes both signs on the unit sphere of $L^2(\IT)$ (Lemma~\ref{lem:W}, witnesses
$(1,\pm1,0,0)$), so no lower bound of the left side bounds $W_T$; the same argument with
$\Rel(A'\bar B-A\bar B')$ forces $\beta_1=0$.
\end{proof}

\cite[Thm.~8.4]{Paper45} is stated for real $F\in C_c^\infty$. The first equality of
Proposition~\ref{prop:P} holds on the whole form domain and is the only one the certificates use;
wherever the zero sum is used for a rough test (Theorem~\ref{thm:margin}), it enters through the
following lemma.

\begin{lemma}[the explicit formula for rough tests]\label{lem:roughEF}
Let $F\in L^2(\IT)$ be real, with \emph{(a)} $\int\log(2+|\tau|)|\Fh|^2<\infty$ and \emph{(b)}
$\sum_\rho|\Phi_F(\rho)\Phi_F(1-\rho)|<\infty$. Then
\[
\sum_\rho\Phi_F(\rho)\Phi_F(1-\rho)=W_T(F)=2\Phi_F(0)\Phi_F(1)+\mathfrak a_T(F),
\]
the sum being absolutely convergent. In particular $\mathfrak a_T(F)=\sum_\rho\Phi_F(\rho)\Phi_F(1-\rho)$
for $F\in\PT$, and a zero $\rho=\frac12+i\gamma$ on the critical line contributes $|\Fh(\gamma)|^2$.
\end{lemma}
\begin{proof}
For real $F$, $\Phi_F(1)=A(F)$ and $\Phi_F(0)=B(F)$, so $2\Phi_F(0)\Phi_F(1)=2\Rel(A\bar B)$. Let
$\eta_\varepsilon\ge0$, $\eta_\varepsilon\in C_c^\infty(-\varepsilon,\varepsilon)$, $\int\eta_\varepsilon=1$,
and $F_\varepsilon=F*\eta_\varepsilon\in C_c^\infty$, supported in $[-T-\varepsilon,T+\varepsilon]$.
Apply \cite[Thm.~8.4]{Paper45} to $F_\varepsilon$ at the window $T+\varepsilon$. Zero side:
$\Phi_{F_\varepsilon}=\Phi_F\Phi_{\eta_\varepsilon}$, $|\Phi_{\eta_\varepsilon}(s)|\le e^{\varepsilon|\Rel s-\frac12|}\le e^{\varepsilon/2}$
in the critical strip and $\Phi_{\eta_\varepsilon}\to1$ pointwise, so by (b) and dominated convergence
the zero sum (which for $F_\varepsilon$ converges absolutely, so that the symmetric order is
immaterial) tends to $\sum_\rho\Phi_F(\rho)\Phi_F(1-\rho)$. Polar term:
$\Phi_{F_\varepsilon}(0)\Phi_{F_\varepsilon}(1)\to\Phi_F(0)\Phi_F(1)$. Archimedean term:
$|\Fh_\varepsilon|^2=|\Fh|^2|\widehat{\eta_\varepsilon}|^2\le|\Fh|^2$ and $|m(\tau)|\le|C_0|+\log(1+|\tau|)$
(Lemma~\ref{lem:mbounds}(i),(ii); $m\le0$ for $|\tau|\le\tau_0$), so by (a) and dominated convergence
$\frac1{2\pi}\int m|\Fh_\varepsilon|^2\to\frac1{2\pi}\int m|\Fh|^2$. Prime term: $c_{F_\varepsilon}\to\cF$
uniformly, the sum runs over the finitely many $n<e^{2T+2\varepsilon}$, and for
$e^{2T}\le n<e^{2T+2\varepsilon}$ one has $\cF(\log n)=0$; so it tends to $-\sum_{n\in\NT}v_n\cF(\log n)$.
\end{proof}


\section{Tail lemmas and an explicit certificate}\label{sec:tails}

A certificate is an inequality $f\ge c$ on all of $\R$. The lemmas of this section reduce it to a
bounded interval; they rest on Lemmas~\ref{lem:mbounds} and \ref{lem:gbounds}.

\subsection{Concentration of the near-optimal feasible measures}\label{subsec:conc}

The program (P) has a minimizer, but no dual optimum is attained (Theorem~\ref{thm:O}); a reduction of
the dual value to a bounded interval cannot be routed through an extremal certificate, and has to be
uniform over the near-optimal feasible set. Write
\[
\mathcal F_\varepsilon=\Big\{\mu\ \text{feasible for (P)}:\ \int\sT\,d\mu\le\lamP+\varepsilon\Big\},
\qquad B_*=\Aconst-C_0,\qquad G(\varepsilon)=\lamP+\varepsilon+B_*,
\]
which is non-empty for every $\varepsilon>0$.

\begin{theorem}[concentration of near-optimal feasible measures]\label{thm:P3}
Let $\varepsilon>0$ and $\eta\in(0,1)$.
\emph{(i)} Every $\mu\in\mathcal F_\varepsilon$ satisfies $\int(m-C_0)\,d\mu\le G(\varepsilon)$ and,
for every $R>0$, $\mu(\{|\tau|\ge R\})\le G(\varepsilon)/(m(R)-C_0)$.
\emph{(ii)} Consequently $\mu(\{|\tau|\ge R\})\le\eta$ for every $\mu\in\mathcal F_\varepsilon$ and
every $R\ge\max\{2,R(\varepsilon,\eta)\}$, where
\[
R(\varepsilon,\eta):=2\pi\exp\Big(1+C_0+\frac{G(\varepsilon)}\eta\Big).
\]
\emph{(iii)} Let $(\nu,\beta)$ be a certificate of value $c$, put
$E_R:=\Aconst+\|\nu\|_{\TV}+|\beta_0|k_0(R)+|\beta_1|k_1(R)-c$, and let
$R\ge\max\{2,R(\varepsilon,\eta)\}$. Then, for every $\mu\in\mathcal F_\varepsilon$,
\[
0\le\lamP-c\le\int_{[-R,R]}(f-c)\,d\mu+\varepsilon_{\mathrm{tail}}(\mu,R),\qquad
\varepsilon_{\mathrm{tail}}(\mu,R)\le\int_{|\tau|>R}(m-C_0)\,d\mu+(E_R+C_0)^+\,\eta ,
\]
the first tail integral being $\le G(\varepsilon)$ always, and $\to0$ as $R\to\infty$ for each fixed
$\mu$.
\end{theorem}
\begin{proof}
(i) $m-C_0\ge0$ by Lemma~\ref{lem:mbounds}(i), and $\int\sT d\mu\ge\int m\,d\mu-\Aconst$ by
Lemma~\ref{lem:mbounds}(iii); hence $\int(m-C_0)d\mu\le\int\sT d\mu+\Aconst-C_0\le G(\varepsilon)$.
Chebyshev's inequality on $\{|\tau|\ge R\}$, where $m-C_0\ge m(R)-C_0$, gives the second bound. (ii) By
Lemma~\ref{lem:mbounds}(ii), $m(R)\ge\log(R/2\pi)-1$ for $R\ge2$, so $m(R)-C_0\ge G(\varepsilon)/\eta$
as soon as $\log(R/2\pi)\ge1+C_0+G(\varepsilon)/\eta$. (iii) By Proposition~\ref{prop:weak},
$\int f\,d\mu=\int\sT d\mu$ for feasible $\mu$, so $\int(f-c)d\mu=\int\sT d\mu-c\in[\lamP-c,\lamP+\varepsilon-c]$;
both pieces of the split at $R$ are nonnegative because $f\ge c$. On $|\tau|>R$,
$f-c\le m+E_R=(m-C_0)+(E_R+C_0)\le(m-C_0)+(E_R+C_0)^+$, by Lemma~\ref{lem:mbounds}(iii),
$|\nuh|\le\|\nu\|_{\TV}$ and Lemma~\ref{lem:gbounds} ($k_0,k_1$ decreasing); apply (i) to the first
summand and (ii) to the second. For fixed $\mu$, $\int(m-C_0)d\mu<\infty$, and the tail integral tends to
$0$ by dominated convergence.
\end{proof}

The positive part in (iii) is needed: $C_0<0$, so $E_R+C_0$ can be negative for certificates of
small total variation, although it is positive for the certificate of Theorem~\ref{thm:C}
($\|\nu\|_{\TV}=90.75$).

\begin{corollary}[verification on a bounded interval]\label{cor:AP}
Suppose $(\nu,\beta)$ satisfies $f\ge c$ \emph{only} on $[-R,R]$ with
$R\ge\max\{2,R(\varepsilon,\eta)\}$, and put
\[
\mathsf V:=\Aconst+\|\nu\|_{\TV}+|\beta_0|\,\gh0(0)+|\beta_1|\,\gh1(0)-C_0 ,
\]
so that $f\ge-\mathsf V$ on all of $\R$. Then $c\le(\lamP+\varepsilon+\mathsf V\eta)/(1-\eta)$ for
every $\varepsilon>0$.
\end{corollary}
\begin{proof}
Since $g_0,g_1\ge0$ on $[0,2T]$, $|\gh j|\le\gh j(0)$, and $m\ge C_0$; this gives $f\ge-\mathsf V$.
For $c<0$ there is nothing to prove. For $c\ge0$ take $\mu\in\mathcal F_\varepsilon$: then
$\lamP+\varepsilon\ge\int\sT d\mu=\int f\,d\mu\ge c\,\mu([-R,R])-\mathsf V\mu(|\tau|>R)\ge c(1-\eta)-\mathsf V\eta$
by Theorem~\ref{thm:P3}(ii).
\end{proof}

\begin{remark}\label{rem:P3exp}
The exponential dependence of $R(\varepsilon,\eta)$ on $1/\eta$ is forced and is not an artifact of
Chebyshev's inequality. Theorem~\ref{thm:Vp}(ii) below constructs, for every $\delta>0$, feasible measures
$\mu'_{\eta,R}$ with mass $\ge\eta(1-o(1))$ beyond $R-\sqrt R$ and value
$\le\lamP+\delta+\eta(\log R+C_1)+o(\eta)$; such a $\mu'$ lies in $\mathcal F_\varepsilon$ as soon as
$\eta(\log R+C_1)+o(\eta)\le\varepsilon-\delta$. So the radius at which mass $\eta$ is forced to have
left satisfies $\log R\asymp1/\eta$, exactly the shape of (ii); we do not know the optimal constant,
which lies between $\varepsilon$ and $G(\varepsilon)$. The radius $R(\varepsilon,\eta)$ has the same
shape as the envelope threshold of Proposition~\ref{prop:env}, with $B_*/\eta$ in place of $\Aconst$.
\end{remark}

\subsection{Lattice atoms: the graded periodic tail}

\begin{lemma}[graded periodic tail]\label{lem:F1}
Let $d=2T/M$, $u_j=jd$ for $M\le j\le J$ (so $u_M=2T$), $\nu=\sum_ja_j(\delta_{u_j}+\delta_{-u_j})$,
$P=2\pi/d$, so that $\nuh(\tau)=\sum_j2a_j\cos(u_j\tau)$ is $P$-periodic. Let
$f=\sT+\nuh+\beta_0\gh0+\beta_1\gh1$. Fix $\tau_1>0$, $N\ge1$, and for $x\in[\tau_1,\tau_1+P)$ put
\begin{align*}
Y_k(x)&=m(x+kP)-\sum_{n\in\NT}v_n\cos\big((x+kP)\log n\big)-|\beta_0|k_0(x+kP)-|\beta_1|k_1(x+kP)\ \ (k<N),\\
Y_N(x)&=m(x+NP)-\Aconst-|\beta_0|k_0(x+NP)-|\beta_1|k_1(x+NP).
\end{align*}
If $\nuh(x)+Y_k(x)\ge c$ for all $x\in[\tau_1,\tau_1+P)$ and all $0\le k\le N$, then $f(\tau)\ge c$
for all $\tau\ge\tau_1$.
\end{lemma}
\begin{proof}
Every $\tau\ge\tau_1$ is uniquely $x+kP$ with $x\in[\tau_1,\tau_1+P)$, $k\ge0$, and
$\nuh(\tau)=\nuh(x)$ since $u_jP=2\pi j$. For $k<N$ the $k$-th row is $f(x+kP)$ with the exact
prime term and the two multiplier terms replaced by the lower bounds of Lemma~\ref{lem:gbounds}.
For $k\ge N$: $m(x+kP)\ge m(x+NP)$, $k_i(x+kP)\le k_i(x+NP)$, and the prime term is $\ge-\Aconst$.
\end{proof}

The prime term is never required to be periodic: it is evaluated exactly on the first $N$
copies of one period of $\nuh$ and replaced by its infimum beyond. The lemma therefore holds
for any finite set of contacts; the incommensurability of $\log2$ and $\log3$ plays no role.
It is a sufficient criterion only (the rows are conservative in $\beta$ by
$|\beta_0|k_0+|\beta_1|k_1$).

\subsection{Exponential columns: the monotone tail}

\begin{proposition}[exponential columns]\label{prop:T0}
For $s>0$ and an integer $p\ge0$ let $d\nu_{s,p}(u)=\frac12(|u|-2T)^pe^{-s(|u|-2T)}\mathbf 1_{|u|\ge2T}du$,
a positive finite measure of mass $p!/s^{p+1}$ supported outside the window. Then
\[
\nuh_{s,p}(\tau)=\Rel\Big[e^{2iT\tau}\frac{p!}{(s-i\tau)^{p+1}}\Big],\qquad
|\nuh_{s,p}(\tau)|\le\operatorname{cap}_{s,p}(\tau):=\frac{p!}{(s^2+\tau^2)^{(p+1)/2}},
\]
and $\operatorname{cap}_{s,p}$ is decreasing in $\tau\ge0$.
\end{proposition}
\begin{proof}
$\nuh_{s,p}(\tau)=\int_{2T}^\infty(u-2T)^pe^{-s(u-2T)}\cos\tau u\,du=\Rel\big[e^{2iT\tau}\int_0^\infty t^pe^{-(s-i\tau)t}dt\big]$.
\end{proof}

\begin{lemma}[monotone tail]\label{lem:T}
Let $\nu=\sum_{k=1}^Kc_k\nu_{s_k,p_k}+\sum_{j=1}^Ja_j(\delta_{u_j}+\delta_{-u_j})$ with $u_j\ge2T$, and
$f=\sT+\nuh+\beta_0\gh0+\beta_1\gh1$. Define for $\tau>0$
\[
L(\tau)=\log\frac\tau{2\pi}-\frac1\tau-\frac1{2\tau^2}-\Aconst-\sum_k|c_k|\operatorname{cap}_{s_k,p_k}(\tau)
-2\sum_j|a_j|-|\beta_0|k_0(\tau)-|\beta_1|k_1(\tau).
\]
Then $f\ge L$ on $(0,\infty)$, $L$ is nondecreasing, and $L(\tau_1)\ge c$ for one $\tau_1$ implies
$f\ge c$ on $[\tau_1,\infty)$.
\end{lemma}
\begin{proof}
Lemma~\ref{lem:mbounds}(ii),(iii), Proposition~\ref{prop:T0}, the trivial bound on the atoms and
Lemma~\ref{lem:gbounds}; every term of $L$ is nondecreasing.
\end{proof}

\begin{corollary}[cost of the tail]\label{cor:T2}
For an atom-free $\nu$ with $A:=\sum_k|c_k|p_k!$ and $A':=|\beta_0|(2\cosh T+\sinh T)+|\beta_1|8T\sinh T$,
the tail inequality $L(\tau_1)\ge c$ holds for every $\tau_1\ge\max\{A+A'+2,\ 2\pi e^{c+\Aconst+1}\}$,
so the tail costs linearly in the $\ell^1$ size of the certificate. For atoms it costs
exponentially, $\tau_1\approx2\pi e^{\Aconst+2\sum|a_j|}$, which is why Lemma~\ref{lem:F1} and not
Lemma~\ref{lem:T} is used for lattice certificates. Neither cost is uniform over the dual class:
by Theorem~\ref{thm:O}, $\|\nu\|\to\infty$ along every maximizing sequence.
\end{corollary}
\begin{proof}
For $\tau\ge1$, $\operatorname{cap}_{s,p}(\tau)\le p!/\tau$ and $|\beta_i|k_i\le A'/\tau$, so
$L(\tau)\ge\log\frac\tau{2\pi}-\Aconst-\frac{A+A'+2}{\tau}$.
\end{proof}

\begin{lemma}[hybrid tail]\label{lem:hybrid}
Let $\nu=\nu_{\mathrm{lat}}+\sum_\ell c_\ell\nu_{s_\ell,p_\ell}$ with $\nu_{\mathrm{lat}}$ as in Lemma~\ref{lem:F1}.
For $x\in[\tau_1,\tau_1+P)$ define $H_k(x)=\nuh_{\mathrm{lat}}(x)+Y_k(x)+\sum_\ell c_\ell\nuh_{s_\ell,p_\ell}(x+kP)$
for $0\le k<N$, and $H_N(x)=\nuh_{\mathrm{lat}}(x)+Y_N(x)-\sum_\ell|c_\ell|\operatorname{cap}_{s_\ell,p_\ell}(x+NP)$.
If $H_k\ge c$ on $[\tau_1,\tau_1+P)$ for all $0\le k\le N$ then $f\ge c$ on $[\tau_1,\infty)$.
\end{lemma}
\begin{proof}
As in Lemma~\ref{lem:F1}; on the rows $k<N$ the column term is evaluated exactly, and for
$k\ge N$ it is bounded by $-\sum_\ell|c_\ell|\operatorname{cap}(x+kP)\ge-\sum_\ell|c_\ell|\operatorname{cap}(x+NP)$
since $\operatorname{cap}$ decreases.
\end{proof}

\subsection{Exhaustion by exponential columns}

\begin{theorem}[exponential columns exhaust the dual]\label{thm:E}
Let $\mathcal E$ be the real span of $\{\nu_{s,p}:s>0,p\ge0\}$. Then
$\sup\{\inf_\tau[\sT+\nuh+\beta\cdot\widehat g]:\ \nu\in\mathcal E,\beta\in\R^2\}=\lamP$, and the same
holds with $\mathcal E$ replaced by the absolutely continuous measures with bounded density and
compact support in $\{|u|\ge2T\}$.
\end{theorem}
\begin{proof}
``$\le$'' is weak duality. For ``$\ge$'' fix $\varepsilon>0$ and, by Theorem~\ref{thm:D1}, an atomic
$\nu_0=\sum_{j\le J}a_j(\delta_{u_j}+\delta_{-u_j})$ and $\beta$ with
$c_0:=\inf[\sT+\nuh_0+\beta\cdot\widehat g]>\lamP-\varepsilon/2$. By Lemma~\ref{lem:T} there is
$\tau_1<\infty$ with $\sT+\nuh_0+\beta\cdot\widehat g\ge c_0-\varepsilon/8$ on $[\tau_1,\infty)$,
and the lemma uses only $\sum|a_j|$, so the same bound holds for every signed measure of the
same total variation supported in $\{|u|\ge2T\}$. Replace each atom $a_j\delta_{u_j}$ by
$a_j$ times the uniform probability density on $(u_j,u_j+\eta)$: since
$|\cos u\tau-\cos u_j\tau|\le|\tau|\eta$, the transform changes by at most $2\sum|a_j|\tau_1\eta$
on $[0,\tau_1]$, which is $<\varepsilon/8$ for small $\eta$; the new $\nu_1$ has bounded density
and compact support and value $\ge c_0-\varepsilon/4$. Finally $\{t^pe^{-t}\}_{p\ge0}$ is total in
$L^1(0,\infty)$ (if $g\in L^\infty$ annihilates them, its Laplace transform is holomorphic on
$\Rel z>0$ with all derivatives vanishing at $z=1$, hence $g=0$), so there is $\nu_2\in\mathcal E$
with $\|\nu_2-\nu_1\|_{\TV}<\varepsilon/8$, and $\|\nuh_2-\nuh_1\|_\infty\le\|\nu_2-\nu_1\|_{\TV}$ on all of $\R$.
\end{proof}


\begin{proposition}[the canonical column; band-limitation]\label{prop:M}
Let $d\nu_K=\frac12K(|u|)\mathbf 1_{|u|\ge2T}du$, a positive measure of mass
$c_K=\int_{2T}^\infty K$ (Theorem~\ref{thm:S}(iii)). Then for all $\tau$
\[
\sT(\tau)+\nuh_K(\tau)=\big(C_0+c_K\big)+\int_0^{2T}(1-\cos\tau u)K(u)\,du-\sum_{n\in\NT}v_n\cos(\tau\log n),
\]
a function whose distributional Fourier transform is supported in $[-2T,2T]$ (it still grows
like $\log|\tau|$: the singularity of $K$ at $u=0$ lies inside the window). Moreover
$d\nu_K=2e^{-T}\big(1+O(e^{-2|u|})\big)\,d\nu_{\frac12,0}$:
the archimedean L\'evy density selects the scale $s=\frac12$ among the exponential columns.
\end{proposition}
\begin{proof}
By \eqref{eq:m}, split $\int_0^\infty(1-\cos\tau u)K$ at $2T$ and use
$\nuh_K(\tau)=\int_{2T}^\infty K(u)\cos\tau u\,du$; the tail contributes $c_K-\nuh_K(\tau)$. For the support
statement let $\psi$ be a Schwartz function whose Fourier transform vanishes on a neighbourhood of
$[-2T,2T]$. Then $\int\psi=0$, $\int\psi(\tau)\cos(\tau\log n)d\tau=0$ for $n\in\NT$, and, by Fubini
(justified by $|1-\cos\tau u|K(u)\le\min(2,\tau^2u^2)K(u)$),
$\int\psi(\tau)\int_0^{2T}(1-\cos\tau u)K(u)\,du\,d\tau=\int_0^{2T}K(u)\big(\widehat\psi(0)-\Rel\widehat\psi(u)\big)du=0$.
Finally, by Lemma~\ref{lem:K}, $\frac12K(u)=e^{-u/2}(1+O(e^{-2u}))$, while the density of $\nu_{\frac12,0}$
is $\frac12e^{T}e^{-u/2}$ for $u\ge2T$.
\end{proof}

By Theorem~\ref{thm:O}, $(\nu_K,0,0)$ is not a certificate of value $\lamP$. The identity says that the
outside-window L\'evy tail removes from the symbol every frequency outside the window; apart from the
$\log|\tau|$ growth generated by $K$ near $u=0$, what remains is carried by frequencies in $[0,2T]$.
Section~\ref{sec:Z} identifies this remainder.

\subsection{An explicit certificate at \texorpdfstring{$T=\tfrac12\log3$}{T = (1/2) log 3}}\label{sec:cert}

At $T=\frac12\log3$ the window is $[-\log3,\log3]$, the contact set is $\NT=\{2\}$, and
$v_2=2\log2/\sqrt2=0.98025814\ldots=\Aconst$. Positivity of the complete form here is known by
\cite[Thm.~1.2]{Zhu2026}, and positivity of both forms also follows from Theorems~\ref{thm:log5A} and
\ref{thm:log5W} below. The point of the following theorem is the method: a proof by an explicit dual
feasible point, with a margin about $9.4\%$ of the numerical value $\lambda_{\mathrm A}(\frac12\log3)\approx9.448\cdot10^{-4}$.

\begin{theorem}[explicit certificate]\label{thm:C}
Let $T=\frac12\log3$, $d=2T/24=\frac{\log3}{24}$,
$u_j=jd$ for $24\le j\le460$ (so $u_{24}=\log3$ and $u_{460}=21.0567\ldots$), let
$a_{24},\dots,a_{460}$ be the $437$ rational numbers given as exact decimal strings in the
certificate data accompanying this paper, and let
$\beta_0=2.00256170925026\ldots$, $\beta_1=-0.00965443404485\ldots$. Put
$\nu=\sum_ja_j(\delta_{u_j}+\delta_{-u_j})$, so $\|\nu\|_{\TV}=2\sum|a_j|=90.75$, $\nuh(0)=1.7337$.
Then
\[
\sT(\tau)+\nuh(\tau)+\beta_0\gh0(\tau)+\beta_1\gh1(\tau)\ \ge\ c:=8.8577\cdot10^{-5}\qquad\text{for every }\tau\in\R,
\]
and therefore $\lamA\ge\lamP\ge8.8577\cdot10^{-5}>0$. The proof includes a computer verification
in interval arithmetic (Appendix~\ref{app:encl}).
\end{theorem}

\begin{proof}
The function $f$ is even, so it suffices to take $\tau\ge0$. Split $[0,\infty)$ as $[0,\tau_1]$
with $\tau_1=30$, where $f$ itself is bounded below, and $[\tau_1,\infty)$, where
Lemma~\ref{lem:F1} with $N=2$ and $P=2\pi/d=137.2608\ldots$ reduces the inequality to the three
rows $\nuh+Y_k\ge c$, $k=0,1,2$, on the single period $[30,167.2608]$.

Each of the four functions ($f$ on $[0,\tau_1]$ and the three rows) is covered by the exhaustive
second-order procedure of Proposition~\ref{prop:TAY} (Appendix~\ref{app:encl}). The second-derivative
bound is $M_2=6264.92$ for $f$, computed from the exact decimals; the three rows admit the smaller
constant $6262.99$, and the larger one is used for all four. The initial cell width is $0.02$, the
minimal width $w_{\min}=10^{-9}$; there is no cap on the number of cells and no cell is discarded.
The cover uses $2\,378\,593$ evaluations, leaves no unresolved interval, and consists of $20$
segments tiling the four intervals exactly.

The target in the cover is $8.8578\cdot10^{-5}$ minus a floating-point allowance
$\varepsilon_{\mathrm{fp}}=10^{-9}$. This allowance exceeds the a priori error of the double-precision
evaluator: argument rounding in the $437$ cosines contributes at most
$2\sum_j|a_j|\cdot u_{460}\cdot(\tau_1+NP)\cdot2^{-53}<10^{-10}$, the Taylor remainder at the minimal
width at most $M_2w_{\min}^2/8<10^{-15}$, and the summation of the $445$ terms together with the
evaluation of $\psi$ less than $10^{-12}$.

The constants $H_0,H_1$, $\psi'(\frac14)$, $k_0,k_1$ and the values of $m$ at the four minimizers
are enclosed in interval arithmetic at $40$ digits, $m$ through the recurrence
$\Rel\psi(z)=\Rel\psi(z+L)-\sum_{j<L}\Rel(z+j)^{-1}$ and the asymptotic series with Stieltjes'
remainder bound. Writing $h_k=\nuh+Y_k$ for the rows, the four minima $f(17.645)$, $h_0(37.815)$, $h_1$, $h_2(58.817)$
are strictly positive intervals with lower ends $\ge8.8578\cdot10^{-5}$; the smallest is
$h_0(37.8153)=8.85786\cdot10^{-5}$. The certificate data and the verification statistics accompany
the paper, and an independent replay from the decimal strings reproduces all
figures.
\end{proof}

\begin{remark}[how the certificate was found]
The atoms were placed on the lattice $u_j=j\cdot2T/M$ because Lemma~\ref{lem:F1} then controls the
tail on one period; the resolution near the window edge, rather than the range of the atoms, turned
out to be the decisive parameter. A linear program in $(a_j,\beta,c)$ with cutting planes on a
$\tau$-grid produced an approximate profile; since values on a grid are unreliable at the required
accuracy, the profile was re-optimized in $128$ blocks on a grid $10^3$ times finer. The result has three visible features:
\begin{enumerate}
\item $\beta_0=2.0026$ and $\beta_1\approx0$: the optimizer restores the two polar terms of the
explicit formula that the primitive class removes (Proposition~\ref{prop:P});
\item the weights $a_j$ alternate in sign at $70\%$ of consecutive lattice sites, decay smoothly away
from the window edge (mean $|a_j|$ from $0.27$ near $u=2T$ to a few hundredths near the last atom $u_{460}\approx21$), and show no
arithmetic structure (no $\log2$-spacing, no equal amplitudes);
\item the contact set $\{f-c<10^{-5}\}$ is a comb: a band inside $[0,13]$, just beyond the last
sign change of $\sT$ at $\tau=10.53$, plus one phase per tail row. This is the complementary
slackness behind Theorem~\ref{thm:O}: the contact set must grow as $c\uparrow\lamP$ and can never be
all of $\R$.
\end{enumerate}
\end{remark}

Zhu \cite[Thm.~1.2]{Zhu2026} combines a pointwise envelope of $\sT$ with a certified Cholesky
factorization of one $200\times200$ Legendre block; Theorem~\ref{thm:C} is instead an explicit dual
feasible point, whose ingredients (polar multiplier, outside-window measure, contact comb) are
interpretable, and whose existence for every $c<\lamP$ is guaranteed by Theorem~\ref{thm:D1}.


\section{The dual variable is the explicit formula}\label{sec:Z}

So far every certificate has been a list of numbers produced by an optimizer. This section
identifies these numbers. Up to an arbitrarily small uniform correction, every certificate of the
complete Weil form is the Weil explicit formula evaluated on a family of modulated test functions,
and its value at $\tau$ is a smoothing of the zero measure at $\tau$.

\subsection{Flat-top weights}
\begin{definition}\label{def:W}
$\Wcl$ is the set of even measurable $w:\R\to\R$ with $w\equiv1$ on $[-2T,2T]$ and
$|w(u)|\le Ce^{-(\frac12+\delta)|u|}$ for $|u|\ge2T$, for some $C,\delta>0$. $\Wcl^{\mathrm{reg}}\subset\Wcl$
consists of the $w\in C^\infty(\R)$ with compact support. For $w\in\Wcl$ define the even measure
\begin{equation}\label{eq:nuw}
d\nu_w=\Big[\tfrac12\big(1-w(u)\big)K(|u|)+2w(u)\cosh\tfrac u2\Big]\mathbf 1_{|u|\ge2T}\,du
-\sum_{n\notin\NT}\tfrac{v_n}2\,w(\log n)\,\big(\delta_{\log n}+\delta_{-\log n}\big),
\end{equation}
and $f_w:=\sT+\nuh_w+2\gh0$.
\end{definition}
The measure $\nu_w$ is finite ($K$ is integrable on $[2T,\infty)$, $|w|\cosh\frac u2$ is
integrable by the decay condition, and $\sum_nv_n|w(\log n)|\le2C\sum\Lambda(n)n^{-1-\delta}$) and
supported in $\{|u|\ge2T\}$. Hence $(\nu_w,2,0)$ is a dual feasible point, and by
Propositions~\ref{prop:weak} and \ref{prop:P}, $W_T\ge\inf f_w$ on $L^2(\IT)$.

\begin{theorem}\label{thm:Z}
\emph{(i)} For every $w\in\Wcl$ and $\tau\in\R$,
\begin{equation}\label{eq:fw}
\begin{split}
f_w(\tau)=C_0+\int_0^\infty\big(1-w(u)\cos\tau u\big)K(u)\,du&-\sum_{n\ge2}v_n\,w(\log n)\cos(\tau\log n)\\
&+2\int_\R w(u)\cosh\tfrac u2\cos(\tau u)\,du,
\end{split}
\end{equation}
which is the right-hand side of the Weil explicit formula for the even test function
$h_\tau(u)=w(u)\cos(\tau u)$.
\emph{(ii)} For $w\in\Wcl^{\mathrm{reg}}$, with $\wh(t)=\int w(u)e^{itu}du$ (an even entire function
of exponential type),
\begin{equation}\label{eq:Z}
f_w(\tau)=\sum_\rho\wh(\tau-\gamma_\rho),\qquad\rho=\tfrac12+i\gamma_\rho,
\end{equation}
the sum running over the non-trivial zeros of $\zeta$ with multiplicity, and converging absolutely.
A zero on the critical line contributes $\wh(\tau-\gamma_\rho)$ with $\gamma_\rho$ real, so on
the zeros of the critical line $f_w$ is the smoothing of the zero-counting measure by the kernel $\wh$.
\end{theorem}
\begin{proof}
(i) By \eqref{eq:m}, $m(\tau)+\int_{2T}^\infty(1-w)K\cos\tau u\,du=C_0+\int_0^{2T}(1-\cos\tau u)K
+\int_{2T}^\infty K-\int_{2T}^\infty wK\cos\tau u=C_0+\int_0^\infty(1-w\cos\tau u)K$, using $w=1$ on
$[0,2T]$. The prime terms inside and outside the window combine into $-\sum_{n\ge2}v_nw(\log n)\cos(\tau\log n)$,
again because $w(\log n)=1$ for $n\in\NT$, and $2\gh0+2\int_{|u|\ge2T}w\cosh\frac u2\cos\tau u=2\int_\R w\cosh\frac u2\cos\tau u$.
(ii) For even $h\in C_c^\infty(\R)$ the Weil explicit formula \cite{Weil52}, \cite[Thm.~5.12]{IK}, in the
normalization of this paper, reads
\[
\sum_\rho\widehat h(\gamma_\rho)=2\int_\R h(u)\cosh\tfrac u2\,du-\sum_{n\ge2}v_n\,h(\log n)+\frac1{2\pi}\int_\R m(t)\,\widehat h(t)\,dt .
\]
Apply it to $h_\tau=w\cos(\tau\,\cdot)$, whose transform is $\frac12[\wh(t-\tau)+\wh(t+\tau)]$. The
zero sum is $\sum_\rho\wh(\tau-\gamma_\rho)$, since $\wh$ is even and the multiset $\{\gamma_\rho\}$ is
invariant under $\gamma\mapsto-\gamma$ ($\rho\mapsto1-\rho$). The polar and prime terms are the
last two terms of \eqref{eq:fw}. For the archimedean term insert \eqref{eq:m}: by Fourier
inversion $\frac1{2\pi}\int\wh(t\mp\tau)\,dt=w(0)=1$ and
$\frac1{2\pi}\int\wh(t\mp\tau)\cos(tu)\,dt=w(u)\cos(\tau u)$, and Fubini is justified by
$\int|\wh(t-\tau)|\int_0^\infty\min(2,t^2u^2)K(u)\,du\,dt<\infty$; the result is
$C_0+\int_0^\infty(1-w(u)\cos\tau u)K(u)\,du$. Absolute convergence of the zero sum holds because $\wh$ is
Schwartz on horizontal lines of the strip $|\Iml z|\le\frac12$, uniformly, and $N(t)\ll t\log t$.
\end{proof}


\begin{corollary}[certificates are statements about zeros]\label{cor:zeros}
If $w\in\Wcl^{\mathrm{reg}}$ and the elementary function \eqref{eq:fw} satisfies $f_w\ge c$ on $\R$,
then $\sum_\rho\wh(\tau-\gamma_\rho)\ge c$ for every real $\tau$, the sum including any zeros off the
critical line through $\wh$ at complex arguments. By Theorem~\ref{thm:finZ} below, this inequality
cannot be verified on any finite list of zeros: every truncation of the sum to finitely many
ordinates takes negative values.
\end{corollary}
\begin{proof}
Immediate from Theorem~\ref{thm:Z}(ii), the zero sum being absolutely convergent.
\end{proof}

\begin{remark}\label{rem:Zcheck}
Numerically, both sides of \eqref{eq:fw}--\eqref{eq:Z} agree for $C^2$ flat-top weights at eight
windows in $[\frac12\log3,1.20]$, the zero side computed from the first $10^5$ ordinates of Odlyzko's
table and the prime side from the elementary right-hand side of \eqref{eq:fw}, to between
$2.6\cdot10^{-7}$ and $2.8\cdot10^{-6}$ (double precision, not certified). At $T=\frac12\log3$ the
truncations of the zero side to the first $10$, $10^2$ and $10^3$ ordinates leave residuals
$4.0\cdot10^{-4}$, $6.0\cdot10^{-6}$ and $1.2\cdot10^{-7}$, of the orders that Lemma~\ref{lem:smooth}(iii)
predicts.
\end{remark}

\begin{theorem}[exhaustion by flat-top weights]\label{thm:Zexh}
Let $T\ge\frac12\log2$. \emph{(i)} $\sup_{w\in\Wcl^{\mathrm{reg}}}\inf_{\tau\in\R}f_w(\tau)=\lamW$.
\emph{(ii)} $\sup\{\inf_\tau[f_w+(\beta_0-2)\gh0+\beta_1\gh1]:\ w\in\Wcl^{\mathrm{reg}},\beta\in\R^2\}=\lamP$.
\end{theorem}
Thus every certificate of the complete Weil form is a flat-top explicit-formula certificate up to
an arbitrarily small uniform correction, and every certificate of the primitive form is one plus a
correction in the two primitivity multipliers.
\begin{proof}
``$\le$'' is weak duality in both cases. For ``$\ge$'' in (i): the proof of Theorem~\ref{thm:D1}
with the two $\widehat g$-constraints removed gives $\lamW=\sup\inf[\sT+2\gh0+\nuh]$ over finite
atomic $\nu$ outside the window, and the argument of Theorem~\ref{thm:E} (with $\beta=(2,0)$
fixed) replaces the atoms by a bounded even density $h$ with compact support in $\{|u|\ge2T\}$,
losing at most $\varepsilon/2$. It remains to realize $h\,du$ as $\nu_w$ up to total variation
$\varepsilon/2$, with $w\in\Wcl^{\mathrm{reg}}$; the transform then moves by at most $\varepsilon/2$
uniformly. On $u\ge2T$ the denominator $d_*(u)=2\cosh\frac u2-\frac12K(u)$ is positive and increasing
($d_*(0.3)>0$ and $d_*'>0$ on $(0,\infty)$), because $2T\ge\log2>0.3$; put
$w_h(u)=\big(h(u)-\tfrac12K(u)\big)/d_*(u)$ for $u\ge2T$, extended evenly and by $1$ on the window.
Then $|w_h|\ll e^{-|u|}$ and the absolutely continuous part of $\nu_{w_h}$ is exactly $h\,du$, while
its atomic part is $-\sum_{n\notin\NT}\frac{v_n}2w_h(\log n)(\delta_{\pm\log n})$, of total variation
$\ll\sum_n\Lambda(n)n^{-3/2}<\infty$. Now modify, in this order: (a) replace $h$ on $[2T,2T+\eta]$
by a smooth function matching $d_*(u)+\frac12K(u)$ to infinite order at $2T$ (so that $w_h\equiv1$
smoothly across the window edge); (b) mollify $h$ and cut $w_h$ off smoothly beyond a large $U_1$
(where $w_h\ll e^{-u}$; the a.c.\ part of $\nu_w$ becomes $\frac12K$ there, of total variation
$\ll e^{-U_1/2}$, and all atoms with $\log n>U_1$ disappear because $w=0$ there); (c) for the finitely
many $n\notin\NT$ with $\log n\le U_1$, replace $h$ by $\frac12K$ on disjoint intervals of small total
length around $\log n$ (so that $w_h$ vanishes there and the remaining atoms disappear). Each step
changes $\nu_{w_h}$ by an arbitrarily small total variation, and the result is in
$\Wcl^{\mathrm{reg}}$. (ii) is the same argument applied to the certificates of Theorem~\ref{thm:E}
with their own $\beta$.
\end{proof}

\subsection{Rigidity, smoothness, and finite zero sets}

\begin{lemma}[rigidity of flat-top weights]\label{lem:Rrig}
Let $w\in L^1(\R)$ be even and continuous with $w\equiv1$ on $[-2T,2T]$ for some $T>0$. Then
$\wh\ge0$ on $\R$ is impossible. Equivalently, if an even continuous $w\in L^1$ with $\wh\ge0$
satisfies $w\equiv1$ on an interval around the origin, then $w\equiv1$ on $\R$, which contradicts
integrability.
\end{lemma}
\begin{proof}
Suppose $w\in L^1$ is continuous with $\wh\ge0$. Then $\wh\in L^1$ and the inversion formula
$w(u)=\frac1{2\pi}\int\wh(t)e^{-itu}dt$ holds pointwise: this is the classical summability argument
(Gauss--Weierstrass or Fej\'er means converge to $w$ uniformly on compacta, and monotone convergence
applied to the nonnegative $\wh$ against the approximate identity gives $\int\wh=2\pi w(0)<\infty$;
see \cite[Ch.~VI]{Katznelson}). Hence $\wh\,dt/2\pi$ is a finite positive measure whose Fourier
transform is $w$, so by Bochner's theorem $w$ is positive definite, with $w(0)=1$. For a positive
definite $\phi$ the nonnegativity of the $3\times3$ Gram matrix at the points $0,-u,a$ gives the
standard inequality
\[
|\phi(u+a)-\phi(u)|^2\ \le\ 2\phi(0)\big(\phi(0)-\Rel\phi(a)\big)\qquad(u,a\in\R).
\]
Take any $a$ with $0<|a|\le2T$: then $w(a)=1=w(0)$, the right side vanishes, and $w(u+a)=w(u)$ for
every $u$. So $w$ is periodic with every period in $(0,2T]$; being continuous, it is constant, equal
to $w(0)=1$. A constant function is not in $L^1(\R)$.
\end{proof}

\begin{corollary}[positivity is never termwise]\label{cor:nontermwise}
For every even continuous $w\in\Wcl$ the kernel $\wh$ takes negative values. Consequently the
inequality $f_w=\sum_\rho\wh(\tau-\gamma_\rho)\ge0$ for $w\in\Wcl^{\mathrm{reg}}$ always expresses
cancellation across the zeros, never the termwise positivity of a kernel: there is no flat-top $w$
with $\wh\ge0$.
\end{corollary}
\begin{proof}
Lemma~\ref{lem:Rrig} applied to $w$, which lies in $L^1$ by the decay hypothesis of
Definition~\ref{def:W}, is even and equals $1$ on the window. A merely measurable $w$ with $\wh\ge0$
agrees almost everywhere with a continuous function, so nothing is lost.
\end{proof}

The following elementary lemma extends Theorem~\ref{thm:Z}(ii), stated for $w\in C_c^\infty$, to the weights $w\in C^k_c$ that arise in finite computations, and bounds the contribution of far zeros.

\begin{lemma}[smoothness and far zeros]\label{lem:smooth}
Let $w\in\Wcl$ be supported in $[-U,U]$, and let $k\ge0$ be an integer such that $w^{(k)}$ is
continuous on $\R$ (in particular across the window edge) and $w^{(k+1)}$ is of bounded variation,
with total variation $V=V(w^{(k+1)})$. Then:
\emph{(i)} for every $z\in\C$, $|\wh(z)|\le V\,e^{U|\Iml z|}|z|^{-k-2}$; in particular
$\wh(t)=O(|t|^{-k-2})$ on $\R$.
\emph{(ii)} The zero sum $\sum_\rho\wh(\tau-\gamma_\rho)$ converges absolutely and
Theorem~\ref{thm:Z}(ii) holds verbatim for such $w$; thus \eqref{eq:Z} is valid on
$\Wcl\cap C_c^k$ for every $k\ge0$, not only on $\Wcl^{\mathrm{reg}}$.
\emph{(iii)} Uniformly for $|\tau|\le H/2$,
\[
\sum_{|\gamma_\rho|>H}\big|\wh(\tau-\gamma_\rho)\big|\ \ll_{U}\ V\,H^{-k-1}\log H .
\]
\emph{(iv)} A zero $\rho=\frac12+i\gamma_\rho$ off the critical line, necessarily with
$|\Iml\gamma_\rho|\le\frac12$, contributes to the sum at most $e^{U/2}$ times the bound that
\emph{(i)} gives for a zero on the line with the same real ordinate.
\end{lemma}
\begin{proof}
(i) Integrating by parts $k+2$ times, the boundary terms vanishing because $w^{(j)}$ is continuous
and compactly supported for $j\le k$,
$\wh(z)=\int w(u)e^{izu}du=(iz)^{-(k+2)}\int e^{izu}\,dw^{(k+1)}(u)$, and $|e^{izu}|\le e^{U|\Iml z|}$
on $[-U,U]$. (ii) By (i) with $k\ge0$ we have $\wh(t)\ll V(1+|t|)^{-2}$ on $\R$ and
$|\wh(t+iy)|\ll Ve^{U/2}(1+|t|)^{-2}$ for $|y|\le\frac12$, while the number of zeros with
$|\gamma_\rho|\in[t,t+1]$ is $O(\log(2+t))$ by $N(t)\ll t\log t$; hence the sum converges absolutely,
and every step of the proof of Theorem~\ref{thm:Z}(ii) (the explicit formula for even $h$ of bounded
variation with compact support, and the Fubini justification, which used only
$\int|\wh(t-\tau)|\int_0^\infty\min(2,t^2u^2)K(u)du\,dt<\infty$) applies unchanged.
(iii) With $|\tau|\le H/2$ and $|\gamma_\rho|>H$ one has $|\tau-\gamma_\rho|\ge|\gamma_\rho|/2$,
so by (i) and partial summation against $N(t)=\frac t{2\pi}\log\frac t{2\pi e}+O(\log t)$,
$\sum_{|\gamma_\rho|>H}|\wh(\tau-\gamma_\rho)|\ll Ve^{U/2}\int_H^\infty t^{-k-2}\log t\,dt
\ll_UV H^{-k-1}\log H$. (iv) is the factor $e^{U|\Iml z|}$ of (i) at $|\Iml z|\le\frac12$.
\end{proof}

For $k=-1$, that is, a weight with a jump at the edge such as the box, the same computation gives
only $\wh(t)=O(|t|^{-1})$; the zero sum then converges conditionally, and (iii) becomes a statement
about oscillation rather than an absolute bound. For this reason the flat-top weights used in
practice are continuous across $u=\pm2T$.

\begin{lemma}[finite zero sets and positive definiteness]\label{lem:fdict}
Let $T>0$, let $Z\subset\R\setminus\{0\}$ be a finite symmetric multiset, and let $w:\R\to\R$ be
even, continuous and integrable with $w\equiv1$ on $[-2T,2T]$. Put
$g_Z(u)=\sum_{\gamma\in Z}e^{i\gamma u}=\sum_{\gamma\in Z}\cos(\gamma u)$ and $\phi:=w\,g_Z$. Then
$\phi$ is continuous and integrable, $\widehat\phi(\tau)=\sum_{\gamma\in Z}\wh(\tau-\gamma)$,
$\phi(0)=|Z|$, $\phi\equiv g_Z$ on $[-2T,2T]$, and
\[
\sum_{\gamma\in Z}\wh(\cdot-\gamma)\ \ge\ 0\ \text{ on }\R
\qquad\Longleftrightarrow\qquad
\phi\ \text{ is positive definite.}
\]
\end{lemma}
\begin{proof}
Continuity, $|\phi|\le|Z|\,|w|\in L^1$ and the values on the window are immediate from $w\equiv1$
there; $\widehat\phi=\wh*\mu_Z$ with $\mu_Z=\sum_{\gamma\in Z}\delta_\gamma$ finite is the
convolution theorem. ($\Leftarrow$) If $\phi$ is positive definite, its spectral measure is
nonnegative, and $\phi\in L^1\cap C$ forces that measure to be $\widehat\phi\,d\tau/2\pi$, so
$\widehat\phi\ge0$. ($\Rightarrow$) If $\widehat\phi\ge0$, the summability argument of
Lemma~\ref{lem:Rrig} (Gauss--Weierstrass means, then monotone convergence against the nonnegative
$\widehat\phi$) gives $\int\widehat\phi=2\pi\phi(0)<\infty$, so $\widehat\phi\,d\tau/2\pi$ is a
finite positive measure with Fourier transform $\phi$, and $\phi$ is positive definite by Bochner.
\end{proof}

\begin{theorem}[finite zero sets are determinate]\label{thm:finDet}
Let $T>0$, let $Z\subset\R\setminus\{0\}$ be a finite symmetric multiset with $|Z|=\ell\ge2$
(multiplicities allowed), and put $g_Z(u)=\sum_{\gamma\in Z}\cos(\gamma u)$. Let $\nu$ be a
finite positive symmetric measure on $\R$ with $\nuh=g_Z$ on $(-2T,2T)$. Then
$\nu=\mu_Z=\sum_{\gamma\in Z}\delta_\gamma$.
\end{theorem}
No decay hypothesis on $\nu$ and no product structure is assumed: the extension problem for a
finite zero set is determinate.
\begin{proof}
Since $\nu$ is finite, $\phi:=\nuh$ is continuous, bounded, real, even and positive definite.

\emph{Step 1 (a lattice inside the window).} Let $w_1,\dots,w_r$ be the distinct elements of
$Z$, with multiplicities $m_1,\dots,m_r\ge1$, so $\sum_jm_j=\ell$ and $r\le\ell$. Put $z_j=e^{iw_ja}$ for
a step $a>0$, so that $g_Z(ka)=\sum_{j\le r}m_jz_j^{\,k}$ for all $k\in\Z$. Since $Z$ is symmetric
with multiplicity, the list $(z_j,m_j)$ is closed under conjugation. Each condition $z_j=z_{j'}$
($j\ne j'$) excludes a discrete set of $a$, so let $E\subset(0,2T/\ell)$ be that exceptional discrete
set and fix $a\in(0,2T/\ell)\setminus E$, so that $z_1,\dots,z_r$ are pairwise distinct. As $r\le\ell$ we
have $ra\le\ell a<2T$, hence $\phi(ka)=g_Z(ka)$ for $|k|\le r$.

\emph{Step 2 (a rank-$r$ Toeplitz block).} Put $c_k:=\phi(ka)\in\R$ ($\phi$ is real and even) and
$A:=[c_{i-j}]_{0\le i,j\le r}$. By Step 1, $A=V^*DV$ with $V_{j,i}=z_j^{\,i}$ and
$D=\operatorname{diag}(m_1,\dots,m_r)$, $m_j>0$; $V$ has rank $r$ because the $z_j$ are distinct, so
$\operatorname{rank}A=r$ and
\[
\ker A=\Big\{v\in\C^{r+1}:\ \textstyle\sum_{i=0}^{r}v_iz_j^{\,i}=0\ (j\le r)\Big\}=\C\cdot v,
\]
$v$ being the coefficient vector of $P(z)=\prod_{j\le r}(z-z_j)$. The root set is closed under
conjugation, so $v\in\R^{r+1}$, $v_r=1$ and $|v_0|=\prod_j|z_j|=1\ne0$.

\emph{Step 3 (the recursion forced by positive definiteness).} Fix $m>r$. The Gram matrix of $\phi$
at $\{0,a,\dots,ra,ma\}$ is positive semidefinite; in block form it is
$M=\begin{psmallmatrix}A&b\\b^*&c_0\end{psmallmatrix}$ with $b_i=\phi(ia-ma)=c_{m-i}$. For
$v\in\ker A$ and $x=(v,t)$ one has $x^*Mx=2\Rel(\bar t\,b^*v)+|t|^2c_0\ge0$ for every $t\in\C$;
choosing the phase of $t$ adversely and letting $|t|\to0$ gives $b^*v=0$, i.e.
\begin{equation}\label{eq:frec}
\sum_{i=0}^{r}v_i\,c_{m-i}=0 .
\end{equation}
The sequence $k\mapsto g_Z(ka)$ satisfies the same relation: $\sum_iv_iz^{m-i}=z^{m-r}z^rP(1/z)$
vanishes at $z=z_j$ because $1/z_j=\bar z_j$ is again a root of $P$; multiplying by $m_j$ and
summing over $j$ gives $\sum_iv_ig_Z((m-i)a)=0$.

\emph{Step 4 (induction).} By Step 1, $c_k=g_Z(ka)$ for $0\le k\le r$. If $c_k=g_Z(ka)$ for all
$0\le k\le m-1$ with $m>r$, then in \eqref{eq:frec} every term with $i\ge1$ equals the
corresponding term for $g_Z$; subtracting the two relations and using $v_0\ne0$ gives
$c_m=g_Z(ma)$. Hence $\phi(ma)=g_Z(ma)$ for every integer $m\ge0$ and every
$a\in(0,2T/\ell)\setminus E$.

\emph{Step 5 (a variable step).} Fix an integer $m\ge1$. By Step 4 the identity
$\phi(ma)=g_Z(ma)$ holds for every $a\in(0,2T/\ell)$ outside the discrete set $E$; both sides are
continuous in $a$, so it holds for every $a\in(0,2T/\ell)$. Since
$\bigcup_{m\ge1}m\cdot(0,2T/\ell)=(0,\infty)$ and both functions are even, $\phi=g_Z$ on $\R$. The
spectral measure of a continuous positive definite function is unique (uniqueness of
Fourier--Stieltjes transforms), so $\nu=\mu_Z$.
\end{proof}

\begin{theorem}[no finite zero set carries a flat-top certificate]\label{thm:finZ}
Let $T>0$, let $Z\subset\R\setminus\{0\}$ be a finite symmetric multiset with $|Z|=\ell\ge2$
(multiplicities allowed), and let $w:\R\to\R$ be even, continuous and integrable with $w\equiv1$ on
$[-2T,2T]$ and $w(u)\to0$ as $|u|\to\infty$. Then
\[
f_{w,Z}(\tau):=\sum_{\gamma\in Z}\wh(\tau-\gamma)
\]
takes negative values on $\R$. In particular, for every continuous $w\in\Wcl$ and every
$H\ge\gamma_1=14.13\ldots$ such that all zeros with $|\Iml\rho|\le H$ lie on the critical line (by
\cite{PlattTrudgian2021}, every $H\le3\cdot10^{12}$), the truncation
$\sum_{|\gamma_\rho|\le H}\wh(\tau-\gamma_\rho)$ of a flat-top certificate is not nonnegative.
\end{theorem}
\begin{proof}
Suppose $f_{w,Z}\ge0$. By Lemma~\ref{lem:fdict} the function $\phi=w\,g_Z$ is continuous and
positive definite with $\phi(0)=\ell$, and $\widehat\phi=f_{w,Z}\in L^1$ with $\int f_{w,Z}=2\pi\ell$, so
$\nu_w:=f_{w,Z}\,d\tau/2\pi$ is a finite positive symmetric measure with $\nuh_w=\phi$, which equals
$g_Z$ on $(-2T,2T)$ because $w\equiv1$ there. By Theorem~\ref{thm:finDet}, $\nu_w=\mu_Z$, i.e.
$w\,g_Z=g_Z$ on $\R$. But along a lattice $\{ma\}$ as in Step~1 of that proof, the $z_j$ are
distinct points of the unit circle, so the Ces\`aro averages of $(z_j\bar z_{j'})^m$ vanish for
$j\ne j'$ (a geometric sum) and
\[
\lim_{M\to\infty}\frac1M\sum_{m=1}^{M}|g_Z(ma)|^2=\sum_{j\le r}m_j^2\ \ge\ r\ \ge\ 1 ,
\]
whence $\limsup_m|g_Z(ma)|\ge1$, while $|w(ma)g_Z(ma)|\le\ell|w(ma)|\to0$, a contradiction. Hence
$f_{w,Z}\ge0$ is impossible; being continuous with $\int_\R f_{w,Z}=2\pi\ell>0$, it takes negative
values.
\end{proof}


The proof of Theorem~\ref{thm:finZ} does not locate the negative values: it contradicts positive
definiteness along a lattice in the variable $u$ and says nothing about $\tau$. Under additional
smoothness hypotheses, integration by parts locates them.

\begin{proposition}[the location of the negativity]\label{prop:loc}
Let $w\in\Wcl$ be supported in $[-U,U]$, let $k\ge0$, assume that $w^{(k)}$ is continuous on $\R$ and
that $w^{(k+1)}$ is of bounded variation on $\R$, absolutely continuous on $\R\setminus\{\pm U\}$, with
a jump $J\ne0$ at $u=U$, i.e.\ $J=w^{(k+1)}(U^+)-w^{(k+1)}(U^-)$. Let $Z\subset\R\setminus\{0\}$ be a
finite symmetric multiset with $S:=\sum_{\gamma\in Z}\cos(U\gamma)\ne0$. Put $c_k(x)=(-1)^{k/2+1}\cos x$
for even $k$ and $c_k(x)=(-1)^{(k-1)/2}\sin x$ for odd $k$. Then, as $t\to+\infty$,
\[
\wh(t)=\frac{2J\,c_k(Ut)}{t^{k+2}}+o\big(t^{-k-2}\big),\qquad
f_{w,Z}(\tau)=\frac{2JS\,c_k(U\tau)}{\tau^{k+2}}+o\big(\tau^{-k-2}\big),
\]
so $f_{w,Z}(\tau)<0$ for arbitrarily large $\tau$, at the rate $\tau^{-k-2}$ and the frequency $U$.
\end{proposition}
\begin{proof}
Integrating by parts $k+2$ times as in Lemma~\ref{lem:smooth}(i), the boundary terms vanishing because
$w^{(j)}$ is continuous and compactly supported for $j\le k$,
$\wh(t)=(-1/it)^{k+2}\int e^{itu}\,dw^{(k+1)}(u)$. The measure $dw^{(k+1)}$ is an $L^1$ density plus
atoms at $\pm U$; the density contributes $o(1)$ to the integral by the Riemann--Lebesgue lemma. Since
$w^{(k+1)}$ has parity $(-1)^{k+1}$, its jump at $-U$ is $(-1)^kJ$, so the atoms contribute
$J(e^{itU}+(-1)^ke^{-itU})$, that is $2J\cos(Ut)$ for even $k$ and $2iJ\sin(Ut)$ for odd $k$. Multiplying
by $(-1/it)^{k+2}=(-1)^k i^{-(k+2)}t^{-(k+2)}$ gives the first formula. For the second, sum over the
finitely many $\gamma\in Z$, use $(\tau-\gamma)^{-k-2}=\tau^{-k-2}(1+O(1/\tau))$, and
$\sum_\gamma c_k(U(\tau-\gamma))=S\,c_k(U\tau)$, because $Z$ is symmetric, so that
$\sum_\gamma\sin(U\gamma)=0$. Since $JS\ne0$ and $c_k(U\tau)$ equals $\pm1$ on two arithmetic
progressions of step $2\pi/U$, the leading term is $\le-2|JS|\tau^{-k-2}$ on one of them.
\end{proof}

Theorem~\ref{thm:finZ} gives the negativity under minimal hypotheses; Proposition~\ref{prop:loc} gives
its location and size under stronger ones. Neither statement contains the other.

\begin{corollary}[finite sets of zeros do not suffice]\label{cor:infinite}
Let $T>0$. Every flat-top certificate $f_w\ge0$, $w\in\Wcl^{\mathrm{reg}}$, is a statement about the complete,
infinite zero multiset of $\zeta$: no truncation of $\sum_\rho\wh(\tau-\gamma_\rho)$ to finitely many
zeros is nonnegative, for any window, any admissible weight and any truncation height $H\ge\gamma_1$ up
to which the zeros are known to lie on the critical line. A finite
table of zeros can evaluate $f_w$ to any accuracy on a bounded range (Lemma~\ref{lem:smooth}(iii)),
but cannot establish $f_w\ge0$. Quantitatively, if $f_w\ge c>0$ for some $w\in\Wcl\cap C_c^k$
supported in $[-U,U]$ with $V=V(w^{(k+1)})$, and $H$ is large enough that
$C_UVH^{-k-1}\log H\le c/2$, then the truncation $f_{w,H}$ is $\ge c/2$ on
$|\tau|\le H/2$ while being negative somewhere; hence all of its negativity lies at
$|\tau|\gtrsim H/2$.
\end{corollary}
\begin{proof}
The first assertion is Theorem~\ref{thm:finZ} applied to the multiset
$Z=\{\gamma_\rho:|\gamma_\rho|\le H\}$, which is finite, symmetric, free of $0$, real by the choice of
$H$, and has at least two elements since $H\ge\gamma_1$. The
evaluation statement and the quantitative one follow from Lemma~\ref{lem:smooth}(iii), which bounds
$|f_w-f_{w,H}|$ by $C_UVH^{-k-1}\log H$ uniformly for $|\tau|\le H/2$.
\end{proof}


In the opposite direction, a finite zero set does admit a dual certificate, but only a trivial one.

\begin{proposition}[finite zero sets]\label{prop:finuniv}
Let $Z\subset\R\setminus\{0\}$ be a finite symmetric multiset, $\mu_Z=\sum_{\gamma\in Z}\delta_\gamma$,
$g_Z=\widehat{\mu_Z}$, and $W_Z(h):=\sum_{\gamma\in Z}\widehat h(\gamma)$ for even Schwartz $h$. Then:
\emph{(i)} $W_Z(F\star\widetilde F)=\sum_{\gamma\in Z}|\Fh(\gamma)|^2\ge0$ for every Schwartz $F$, where
$\widetilde F(u)=\overline{F(-u)}$;
\emph{(ii)} $\mu_Z$ is the unique symmetric positive tempered measure $\mu$ with
$W_Z(h)=\int\widehat h\,d\mu$ for all even Schwartz $h$;
\emph{(iii)} $g_Z$ is continuous and positive definite, with spectral measure $\mu_Z$.
\end{proposition}
\begin{proof}
(i) $\widehat{F\star\widetilde F}=|\Fh|^2$. (ii) $W_Z(h)=\int\widehat h\,d\mu_Z$ by definition; if $\mu$ is
another such measure, then $\int\widehat h\,d\mu=\int\widehat h\,d\mu_Z$ for every even Schwartz $h$, and
both measures are symmetric, so they agree on all Schwartz functions (the odd part of a test function
integrates to $0$ against both); hence $\mu=\mu_Z$. (iii) is Bochner's theorem for the finite positive
measure $\mu_Z$, and uniqueness of Fourier--Stieltjes transforms.
\end{proof}

For a finite zero set the dual certificate is therefore the spectrum itself: positivity of $W_Z$ is
immediate and the representing measure contains no arithmetic information, while by
Theorem~\ref{thm:finZ} the one-sided smoothing $\sum_{\gamma\in Z}\wh(\tau-\gamma)\ge0$ is never
available for a flat-top weight.

The mechanism is the following. A finite $Z$ makes the Toeplitz forms $[g_Z((i-j)a)]$ of finite rank;
beyond that rank, positive semidefiniteness of any extension propagates the linear recursion of $P$
indefinitely, so the extension is unique and does not decay (Theorem~\ref{thm:finDet}), whereas a
flat-top weight must decay. For the zero set of $\zeta$ no such $P$ exists and the rank is never finite.

\subsection{The decay tail of a flat-top certificate}

Proposition~\ref{prop:env} bounds the tail of any certificate at $2\pi e^{\Aconst+\|\nu\|_{\TV}}$,
and Lemma~\ref{lem:T} at a threshold exponential in the atomic mass. Both ignore the decay of
$\nuh_w$. For a $C^2$ flat-top weight this decay is available. The following theorem extends
Lemma~\ref{lem:T}: the single term $2\sum_j|a_j|$ is split into an atomic part, which does not decay,
and an absolutely continuous part, which decays like $\tau^{-2}$. It is also the smooth analogue of
Lemma~\ref{lem:F1}, with one essential difference recorded in Remark~\ref{rem:Dgrade}.

\begin{theorem}[decay tail bound for flat-top weights]\label{thm:Dtail}
Let $T>0$ and let $w\in\Wcl\cap C^2(\R)$ (so $w$ is even, $w\equiv1$ on $[-2T,2T]$, and $w'$, $w''$
match across $u=\pm2T$; in particular $w'(2T)=w''(2T)=0$), with $w''$ of bounded variation, and
assume either \emph{(a)} $\supp w\subset[-U,U]$, or \emph{(b)} $|w^{(j)}(u)|\le Ce^{-s|u|}$ for
$j\le2$ with $s>\frac12$. Put $\Omega(u)=w(|u|)\cosh\frac u2$ and $R=(1-w)K$ on $[2T,\infty)$, and
\[
D_c(w)=4\int_0^\infty\Big|\big(w(u)\cosh\tfrac u2\big)''\Big|\,du,\qquad
D_K(w)=\int_{2T}^\infty\big|\big((1-w)K\big)''(u)\big|\,du,
\]
\[
D(w)=D_c(w)+D_K(w),\qquad E_{\mathrm{at}}(w)=\sum_{\log n>2T}v_n\,|w(\log n)| ,
\]
all four finite. Then, with $\Psi_w(\tau)=\int_{2T}^\infty(1-w)K\cos\tau u\,du$,
\begin{equation}\label{eq:Dsplit}
\nuh_w(\tau)+2\gh0(\tau)=2\widehat\Omega(\tau)+\Psi_w(\tau)-\sum_{\log n>2T}v_n\,w(\log n)\cos(\tau\log n),
\end{equation}
and
\[
|2\widehat\Omega(\tau)|\le D_c(w)\,\tau^{-2},\qquad |\Psi_w(\tau)|\le D_K(w)\,\tau^{-2}.
\]
Consequently, for $f=f_w+(\beta_0-2)\gh0+\beta_1\gh1$ and every $\tau>0$,
\[
f(\tau)\ \ge\ L_w(\tau):=\log\frac\tau{2\pi}-\frac1\tau-\frac1{2\tau^2}-\Aconst-E_{\mathrm{at}}(w)
-|\beta_0-2|k_0(\tau)-|\beta_1|k_1(\tau)-D(w)\,\tau^{-2},
\]
$L_w$ is nondecreasing on $(0,\infty)$, and positivity of $f$ on $\R$ reduces to the compact
interval $[0,\tau_{\mathrm{stop}}(w)]$, $\tau_{\mathrm{stop}}(w)=\min\{\tau:L_w(\tau)\ge0\}$.
\end{theorem}
\begin{proof}
\emph{The identity \eqref{eq:Dsplit}, and why the $\cosh$ part must be regrouped with $2\gh0$.}
By \eqref{eq:nuw} the absolutely continuous part of $\nu_w$ contributes
$\Psi_w(\tau)+4\int_{2T}^\infty w\cosh\frac u2\cos\tau u\,du$ and its atomic part contributes the
last sum of \eqref{eq:Dsplit}. Since $\gh0(\tau)=2\int_0^{2T}\cosh\frac u2\cos\tau u\,du$ and
$w\equiv1$ on the window, $4\int_{2T}^\infty w\cosh\frac u2\cos\tau u\,du+2\gh0(\tau)
=4\int_0^\infty w\cosh\frac u2\cos\tau u\,du=2\widehat\Omega(\tau)$. The regrouping is necessary: the
density of $\nu_w$ at $u=2T^+$ equals $2\cosh T\ne0$, so $\nu_w$ has a jump at the window edge and
one integration by parts leaves the irreducible term $-4\cosh T\sin(2T\tau)/\tau$. Thus $\nuh_w$
alone is only $O(\tau^{-1})$, and no bound $|\nuh_w|\le D\tau^{-2}$ can hold. By
Lemma~\ref{lem:gbounds}, $2\gh0(\tau)=4\cosh T\sin(2T\tau)/\tau+O(\tau^{-2})$ cancels that term
exactly, with the same coefficient and the opposite sign; this is visible at once in
$2\widehat\Omega$, since $\Omega$ has no edge.

\emph{The two bounds.} $\Omega$ is even and, by the hypotheses, globally $C^2$ on $\R$ with $\Omega''$ of
bounded variation: on the window $\Omega=\cosh\frac u2$, and $w,w',w''$ match across $\pm2T$. In case
(a) $\Omega$ is compactly supported; in case (b) $|\Omega^{(j)}|\le C'e^{-(s-\frac12)|u|}$ for $j\le2$ is
integrable, which is exactly where $s>\frac12$ is needed. Two integrations by parts, all boundary
terms vanishing, give $\widehat\Omega(\tau)=-\tau^{-2}\int_\R e^{i\tau u}\,d\Omega'(u)$, so
$|2\widehat\Omega|\le2V(\Omega')\tau^{-2}=4\int_0^\infty|\Omega''|\,du\cdot\tau^{-2}=D_c(w)\tau^{-2}$ by evenness.
For $R$, extend by $0$ to $(-2T,2T)$ and evenly: $R(2T)=(1-w(2T))K(2T)=0$ and
$R'(2T)=-w'(2T)K(2T)+(1-w(2T))K'(2T)=0$ because $w\equiv1$ on the window and $w\in C^1$, so the
extension is $C^1$ on $\R$ with $R'$ of bounded variation, $V(R')=2\int_{2T}^\infty|R''|$ (in case
(a), $R=K$ beyond $U$, and $K,K',K''$ are integrable at infinity; in case (b),
$|R''|\le(|w''|+2|w'|+|1-w|)\,(K+|K'|+|K''|)$ with $K+|K'|+|K''|\le Ce^{-u/2}$ on $[2T,\infty)$). As $\Psi_w=\frac12\widehat R$,
two integrations by parts give $|\Psi_w|\le\frac12V(R')\tau^{-2}=D_K(w)\tau^{-2}$.

\emph{The atoms carry no decay}, and are bounded only by their mass:
\[
\Big|\sum_{\log n>2T}v_nw(\log n)\cos(\tau\log n)\Big|\le E_{\mathrm{at}}(w),
\]
the sum being finite in case
(a) and convergent in case (b) because $\sum_n\Lambda(n)n^{-\frac12-s}<\infty$ for $s>\frac12$,
the same threshold as above.

\emph{The lower bound.} $\sT(\tau)\ge\log\frac\tau{2\pi}-\frac1\tau-\frac1{2\tau^2}-\Aconst$ by
Lemma~\ref{lem:mbounds}(ii),(iii), the two multiplier terms are bounded by Lemma~\ref{lem:gbounds},
and the three pieces of \eqref{eq:Dsplit} by the above. Each term of $L_w$ is nondecreasing, hence so
is $L_w$, and $L_w(\tau_{\mathrm{stop}})\ge0$ gives $f\ge0$ on $[\tau_{\mathrm{stop}},\infty)$.
\end{proof}

\begin{remark}\label{rem:Dgrade}
The exponent $\tau^{-2}$ is the one Lemma~\ref{lem:smooth}(i) gives at $k=0$ after the two extra
derivatives, and it cannot be improved without $w\in C^3$: integrating by parts once more, a jump of
$w'''$ at the support endpoint contributes a genuine $\tau^{-2}$ to $\widehat\Omega$. The bound on
the atoms is sharp in the worst case by Lemma~\ref{lem:mbounds}(iii). This is where the method
differs from Lemma~\ref{lem:F1}: a lattice $\nuh$ is $P$-periodic, so the comb can be evaluated
exactly on $N$ periods and replaced by its infimum beyond, whereas a flat-top $\nuh_w$ is not
periodic and its comb can only be bounded by its mass. Theorem~\ref{thm:Dtail} handles the absolutely
continuous part by decay. We do not know how to combine the two techniques; the remaining
exponential dependence on the atomic mass comes from this gap.
\end{remark}

\begin{corollary}[comb-vanishing weights]\label{cor:comb}
Assume in addition $\supp w\subset[-U,U]$ and $w(\log n)=0$ for every prime power $n$ with
$2T<\log n\le U$ (a finite set of linear conditions). Then $E_{\mathrm{at}}(w)=0$, and for
$D(w)\le D_0$,
\[
\tau_{\mathrm{stop}}(w)\ \le\ \max\Big\{2\pi e^{\Aconst+1+|\beta_0-2|k_0(1)+|\beta_1|k_1(1)},\ \sqrt{D_0}+2\Big\},
\]
since for $\tau$ beyond the first entry one has $\log\frac\tau{2\pi}-\Aconst-|\beta_0-2|k_0-|\beta_1|k_1\ge1$
and then $L_w(\tau)\ge0$ as soon as $\tau^2\ge D_0+O(\tau)$. In particular
$\tau_{\mathrm{stop}}=O(D(w)^{1/2})$ as $D(w)\to\infty$: the verification interval is polynomial,
not exponential, in the size of the weight, compared with $2\pi e^{\Aconst+E_{\mathrm{at}}(w)}$ in
general and $2\pi e^{\Aconst+\|\nu\|_{\TV}}$ in Proposition~\ref{prop:env}.
\end{corollary}


\begin{corollary}[bounded profiles]\label{cor:bprof}
Let $w$ satisfy the hypotheses of Theorem~\ref{thm:Dtail}(a), with $0\le w\le1$ and $\supp w\subset[-U,U]$,
and put $\kappa_0(T)=K(2T)$, $\kappa_1(T)=|K'(2T)|$. Then
\[
D_c(w)\le4\cosh\tfrac U2\Big(\|w''\|_1+\|w'\|_1+\tfrac14\|w\|_1\Big),\qquad
D_K(w)\le\kappa_0(T)\|w''\|_1+2\kappa_1(T)\|w'\|_1+2\kappa_1(T),
\]
so $D(w)=O(e^{U/2})$ with a constant depending on $\|w\|_{W^{2,1}}$ and $T$. Boundedness, however,
controls only $D$: it bounds $E_{\mathrm{at}}(w)$ by nothing better than $\sum_{2T<\log n\le U}v_n$, which
grows like $4e^{U/2}$. Only the comb conditions of Corollary~\ref{cor:comb} control the atoms.
\end{corollary}
\begin{proof}
On $[0,U]$, $(w\cosh\frac u2)''=w''\cosh\frac u2+w'\sinh\frac u2+\frac14w\cosh\frac u2$, and
$|\sinh|\le\cosh\le\cosh\frac U2$ there; integrating over $[0,\infty)$ gives the bound for $D_c$. On
$[2T,\infty)$, $((1-w)K)''=-w''K-2w'K'+(1-w)K''$. By Lemma~\ref{lem:K}, $K$ is completely monotone, so
$K$ and $|K'|$ are decreasing and $K''\ge0$; hence $|((1-w)K)''|\le\kappa_0|w''|+2\kappa_1|w'|+K''$
(using $0\le1-w\le1$), and $\int_{2T}^\infty K''=-K'(2T)=\kappa_1$. For the atoms, $0\le w\le1$ gives
$E_{\mathrm{at}}(w)\le\sum_{2T<\log n\le U}v_n$, with equality when $w=1$ at those points, and
$\sum_{n\le x}2\Lambda(n)n^{-1/2}=4\sqrt x\,(1+o(1))$ by the prime number theorem and partial summation.
\end{proof}

\begin{remark}\label{rem:comblp}
Numerically, the comb-vanishing class of Corollary~\ref{cor:comb} does not produce a certificate in the
following finite family: compactly supported $C^2$ cubic splines on a lattice of $48$ cells with $U=6$,
matched to $1$ across the window edge, with the comb equalities imposed exactly and $\beta$ free. A
linear program over $\tau\in[0,200]$ returns negative minima at $T=\frac12\log3$, $0.58$, $0.62$ and
$\log2$, of sizes about $10^{-2}$ to $10^{-1}$ (double precision, not certified), and for the
resulting weights $\tau_{\mathrm{stop}}<200$, so the failure concerns all of $\R$. The binding minimum
sits at $\tau\approx4$--$8$, whereas without the comb conditions it lies near $\tau\approx43$. One
structural observation: consecutive prime powers near $e^u$ are spaced $\approx e^{-u}$ on the
logarithmic scale, so once the lattice step exceeds that spacing every added cell is consumed by the
conditions inside it, and comb vanishing forces an effectively bounded support.
\end{remark}

\subsection{The Kre\u{\i}n extension problem at one window}\label{subsec:KEP}
The dual variable of Section~\ref{sec:Z} is a positive measure prescribed only on the window. This
is an extension problem for a positive definite datum given on an interval, in the sense of
Kre\u{\i}n \cite{Krein1940}, \cite[Ch.~2]{Sasvari1994}. In this subsection we make the identification
at one window. The classical theory concerns continuous positive
definite functions, whereas the datum below is a distribution with atoms; we therefore import
nothing from it, and every statement is proved from Theorems~\ref{thm:Z} and \ref{thm:Zexh} and
elementary Fourier analysis.

\begin{definition}[the window datum, and $\mathrm{KEP}(T)$]\label{def:KEP}
Let $g$ be the even distribution defined on even $h\in C_c^\infty(\R)$ by
\begin{equation}\label{eq:gdatum}
\langle g,h\rangle=C_0\,h(0)+\int_0^\infty\big(h(0)-h(u)\big)K(u)\,du-\sum_{n\ge2}v_n\,h(\log n)
+2\int_\R h(u)\cosh\tfrac u2\,du,
\end{equation}
that is, the right-hand side of the explicit formula in the normalization of the proof of
Theorem~\ref{thm:Z}(ii): the L\'evy pair $C_0\delta_0-\frac12\mathrm{Pf}\,K(|u|)$, whose Fourier
transform is the archimedean multiplier $m$ of \eqref{eq:m}; the polar term $2\cosh\frac u2$; and
the prime atoms $-\frac{v_n}2\big(\delta_{\log n}+\delta_{-\log n}\big)$, $n\ge2$. For $T>0$ put
\[
\begin{aligned}
\mathfrak S(T)=\Big\{\nu\ \text{a positive symmetric tempered measure on }\R:\ \\
\langle\nuh-g,h\rangle=0\ \text{for every even }h\in C_c^\infty(-2T,2T)\Big\}.
\end{aligned}
\]
The \emph{Kre\u{\i}n extension problem} $\mathrm{KEP}(T)$ is \emph{solvable} if
$\mathfrak S(T)\ne\emptyset$, \emph{determinate} if $\#\mathfrak S(T)=1$, and \emph{indeterminate}
if $\#\mathfrak S(T)\ge2$.
\end{definition}
Symmetry is imposed because even test functions
determine only the even part of a representing measure, so without it every problem would be
indeterminate for a trivial reason. We define ``the datum is positive definite on the window'' to
mean solvability; by the Schwartz--Bochner theorem this is equivalent to saying that the restriction
of $g$ to $(-2T,2T)$ extends to a positive definite tempered distribution on $\R$.

\begin{lemma}[one negative square in the polar term]\label{lem:polar}
At any $n\ge2$ distinct points $t_1,\dots,t_n$ the Gram matrix of the kernel
$(s,t)\mapsto2\cosh\frac{s-t}2$ is $ab^{\mathsf T}+ba^{\mathsf T}$ with $a_i=e^{t_i/2}$ and
$b_i=e^{-t_i/2}$; it has rank $2$ and signature $(1,1)$. Hence $2\cosh\frac u2$ is positive definite
on no interval, and the solvability of $\mathrm{KEP}(T)$ is automatic at no window, however small.
\end{lemma}
\begin{proof}
$2\cosh\frac{t_i-t_j}2=e^{t_i/2}e^{-t_j/2}+e^{-t_i/2}e^{t_j/2}$, which is the stated rank-two matrix.
It vanishes on $\{a,b\}^\perp$ and on $\mathrm{span}\{a,b\}$ has eigenvalues $a\cdot b\pm|a||b|$;
$a\cdot b=n$ while $|a||b|>n$ strictly by Cauchy--Schwarz, the vectors $a,b$ not being parallel when
the $t_i$ are distinct and $n\ge2$. So one eigenvalue is positive and one negative.
\end{proof}
The datum \eqref{eq:gdatum} is therefore the sum of a positive definite part ($m\ge0$), a part with
exactly one negative square (the pole), and a negative part (the prime atoms). This is the reason a
polar multiplier $\beta_0$ must be carried at all (Proposition~\ref{prop:P}).

\begin{theorem}[a certificate is an absolutely continuous solution]\label{thm:KEP1}
Let $T\ge\frac12\log2$ and let $w\in\Wcl^{\mathrm{reg}}$ satisfy $f_w\ge c$ on $\R$.
\emph{(i)} If $c\ge0$ then $\nu_w:=f_w(\tau)\,d\tau/2\pi$ belongs to $\mathfrak S(T)$; in particular
$\mathrm{KEP}(T)$ is solvable, by an absolutely continuous solution.
\emph{(ii)} If $c>0$ then $\mathfrak S(T)$ contains a continuum of pairwise distinct absolutely
continuous measures; in particular $\mathrm{KEP}(T)$ is indeterminate.
\end{theorem}
\begin{proof}
(i) For $w\in\Wcl^{\mathrm{reg}}$ the product $w\cdot g$ is a tempered distribution and
Theorem~\ref{thm:Z} says $\widehat{w\cdot g}=f_w$, a continuous function with
$f_w(\tau)=\log|\tau|+O(1)$ (Lemma~\ref{lem:mbounds}, Theorem~\ref{thm:Dtail}), hence tempered;
Fourier inversion gives $\nuh_w=w\cdot g$, which agrees with $g$ on $(-2T,2T)$ because $w\equiv1$
there. Positivity of $\nu_w$ is $f_w\ge0$, and symmetry is the evenness of $w$ and $g$.
(ii) Fix an even $\chi\in C_c^\infty(\R)$, $\chi\ge0$, $\chi\not\equiv0$, with
$\supp\chi\subset\{|u|\ge\max(2T,\log2)\}$. For $\varepsilon\in\R$ put $w_\varepsilon=w+\varepsilon\chi
\in\Wcl^{\mathrm{reg}}$; then $f_{w_\varepsilon}=f_w+\varepsilon r_\chi$ with
\[
r_\chi(\tau)=2\int\chi(u)\cosh\tfrac u2\cos\tau u\,du-\sum_{n\ge2}v_n\chi(\log n)\cos(\tau\log n)
-\tfrac12\int_\R\chi(u)K(u)\cos\tau u\,du ,
\]
a bounded continuous function, say $|r_\chi|\le R$ ($\chi$ is compactly supported away from the
singularity of $K$, and $\sum_nv_n|\chi(\log n)|<\infty$). For $|\varepsilon|<c/R$ one has
$f_{w_\varepsilon}>0$, so $\nu_{w_\varepsilon}\in\mathfrak S(T)$ by (i). They are pairwise distinct:
$\nu_{w_\varepsilon}-\nu_{w_{\varepsilon'}}=\frac{\varepsilon-\varepsilon'}{2\pi}r_\chi\,d\tau$, and
$r_\chi\equiv0$ would force $\chi\cdot g=0$, whereas on $\supp\chi\subset\{|u|\ge\log2\}$ the
absolutely continuous part of $g$ is $\chi(u)d_*(u)$ with $d_*(u)=2\cosh\frac u2-\frac12K(u)>0$
(positivity and monotonicity of $d_*$ on $[0.3,\infty)$ are established in the proof of
Theorem~\ref{thm:Zexh}), which is nonnegative and not identically zero and which the singular prime
atoms cannot cancel.
\end{proof}

\begin{theorem}[solvability implies Weil positivity at that window]\label{thm:KEP2}
If $\mathfrak S(T)\ne\emptyset$ then $\lamW\ge0$.
\end{theorem}
\begin{proof}
Let $\nu\in\mathfrak S(T)$ and $F\in C_c^\infty(\IT)$, and put $h=F\star\widetilde F$ with
$\widetilde F(u)=\overline{F(-u)}$. Then $h$ is even, $\supp h\subset(-2T,2T)$ and
$\widehat h=|\Fh|^2\ge0$ is rapidly decreasing, so the pairing against the positive tempered measure
$\nu$ is legitimate. Since $h$ vanishes off the window, only the contacts $n\in\NT$ appear in
\eqref{eq:gdatum}, and comparison with \eqref{eq:Aform} and \eqref{eq:W} gives
$\langle g,h\rangle=\mathfrak a_T(F)+2\Rel\big(A(F)\overline{B(F)}\big)=W_T(F)$. By
Definition~\ref{def:KEP}, $W_T(F)=\int\widehat h\,d\nu=\int|\Fh|^2d\nu\ge0$. Finally
$C_c^\infty(\IT)$ is a form core: for $F$ in the form domain, the dilates $F_\lambda(t)=F(t/\lambda)$,
$\lambda\uparrow1$, are supported in $\lambda\IT$ and converge to $F$ in $L^2$ with
$\Fh_\lambda(\tau)=\lambda\Fh(\lambda\tau)$, so that
$\int\log(2+|\tau|)|\Fh_\lambda-\Fh|^2\to0$ (the weight is doubling), and mollifying $F_\lambda$ at a
scale below $(1-\lambda)T$ keeps the support inside $\IT$ and converges in the same norm. Hence
$W_T\ge0$ on the whole form domain and $\lamW\ge0$.
\end{proof}

\begin{theorem}[the dichotomy at one window]\label{thm:KEPdich}
For every $T\ge\frac12\log2$,
\[
\lamW>0\ \Longrightarrow\ \mathrm{KEP}(T)\ \text{is indeterminate}\ \Longrightarrow\ \lamW\ge0 ,
\]
so certificates and indeterminacy coincide except on the set $\{T:\lamW=0\}$.
\end{theorem}
\begin{proof}
If $\lamW>0$, Theorem~\ref{thm:Zexh}(i) gives $w\in\Wcl^{\mathrm{reg}}$ with $\inf f_w>0$, and
Theorem~\ref{thm:KEP1}(ii) applies; the second implication is Theorem~\ref{thm:KEP2}.
\end{proof}

Off the set $\{T:\lamW=0\}$ the solutions of the extension problem and the certificates therefore
determine each other. A Nevanlinna-type parametrization of $\mathfrak S(T)$ in the indeterminate case
would describe the set of certificates, but it presupposes one known solution.


\section{Lattice certificates and their ceilings}\label{sec:class}

This section bounds from above what a given class of dual objects can certify.

\subsection{A class-dependent bound}
\begin{problem}[class-dependent value bound]\label{prob:A}
For the lattice class of step $d=2T/M$ and range $U$, with the tail rows of Lemma~\ref{lem:F1},
find an explicit $D(T,d,U)\to0$ as $d\to0$, $U\to\infty$, such that the best certificate value
$c(T,d,U)$ satisfies $c\ge\lamP-D(T,d,U)$. Equivalently: a quantitative one-sided approximation
theorem for the function $-\sT-\beta\cdot\widehat g$ on $[0,\tau_{\mathrm{env}}]$ by $P$-periodic
trigonometric polynomials whose low modes $|j|<M$ vanish, in the norm $L^1(|\Fh_\star|^2d\tau/2\pi)$
of the primal minimizer.
\end{problem}
A solution would decide each window by a computation of predictable size and would turn
Theorem~\ref{thm:C} into a construction. By the identity $\lamP-c=\int(f-c)\,d\mu_\star$
(Proposition~\ref{prop:weak}) the deficit is the $\mu_\star$-average of the slack, so the
approximation is needed only where $|\Fh_\star|^2$ lives. Theorem~\ref{thm:Vp} shows that $D$ must
depend on the class, not on a radius alone. If Zhu's conjectural law $-\log\lamP\asymp e^{2T}$ holds,
the effort in Problem~\ref{prob:A} grows at least like $e^{2T}$: finite for each $T$, never uniform.

The converse of Problem~\ref{prob:A}, an upper bound on the value of a class, is a theorem.

\begin{proposition}[test-kernel ceiling for lattice classes]\label{prop:TK}
Let $d=2T/M$, $P=2\pi/d$, and let $\nu$ be any finite even measure supported on the lattice
$\{\pm jd:j\ge M\}$ (any range, any number of atoms). Put $S_\beta=\sT+\beta_0\gh0+\beta_1\gh1$ and,
for $\tau_0\in\R$, let $C_M(\tau_0,\beta)$ be the $M\times M$ Hermitian Toeplitz matrix with entries
$c_{j-k}$, $c_\ell=P^{-1}\int_{\tau_0}^{\tau_0+P}S_\beta(\tau)e^{i\ell d\tau}\,d\tau$. If
$\sT+\nuh+\beta\cdot\widehat g\ge c$ on $\R$ then
\[
c\ \le\ \lambda_{\min}\big(C_M(\tau_0,\beta)\big)\ \text{ for every }\tau_0,\qquad\text{hence}\qquad
c\ \le\ \operatorname{Ceil}(M,T):=\sup_{\beta\in\R^2}\inf_{\tau_0}\lambda_{\min}\big(C_M(\tau_0,\beta)\big),
\]
and $\beta\mapsto\inf_{\tau_0}\lambda_{\min}(C_M(\tau_0,\beta))$ is concave. The ceiling depends on
the class only through $M$.
\end{proposition}
\begin{proof}
$\nuh(\tau)=\sum_{j\ge M}2a_j\cos(jd\tau)$ is $P$-periodic with Fourier modes $|j|\ge M$ only; in
particular it has mean zero over any period. Let $q(z)=\sum_{k=0}^{M-1}q_kz^k$ and
$\vartheta(\tau)=|q(e^{id\tau})|^2\ge0$, a trigonometric polynomial in $d\tau$ of degree $\le M-1$. Then
$\int_{\tau_0}^{\tau_0+P}\nuh\vartheta\,d\tau=0$, so integrating $f\ge c$ against $\vartheta$ gives
$c\int\vartheta\le\int S_\beta\vartheta$. Now $\int_{\tau_0}^{\tau_0+P}\vartheta=P\sum|q_k|^2$ and
$\int_{\tau_0}^{\tau_0+P}S_\beta\vartheta=P\sum_{j,k}\bar q_kq_jc_{j-k}$, so
$c\le\langle C_Mq,q\rangle/\|q\|^2$ for every $q\ne0$, i.e.\ $c\le\lambda_{\min}(C_M)$; by the
Fejér--Riesz theorem every nonnegative trigonometric polynomial of degree $\le M-1$ is of the form
$|q|^2$, so nothing is lost. For fixed $q$ and $\tau_0$ the Rayleigh quotient is affine in $\beta$
($S_\beta$ is), so $\inf_{q,\tau_0}$ is concave in $\beta$.
\end{proof}

The ceiling is the smallest eigenvalue of a small matrix and needs no optimizer; like
Theorem~\ref{thm:Vp}, it depends on the class and not on a radius. Numerically, at
$T=\frac12\log3$ it is negative for $M=8$ and positive for $M=12,24$, and the certificate of
Theorem~\ref{thm:C} realizes a fraction of about $15\%$ of its own class ceiling
(double-precision values, Appendix~\ref{app:num}). The negative sign at $M\le8$ survives an interval enclosure; this is
Theorem~\ref{thm:N} below. Since the lattices $2T/2^kM_0$ are nested with dense union,
Theorem~\ref{thm:D1} gives $\sup_M\operatorname{Ceil}(M,T)\ge\lamP$; that the ceilings in fact
converge to $\lamP$, so that the test-kernel family is not lossy in the limit, is
Theorem~\ref{thm:R} and Corollary~\ref{cor:R3} below; the ceilings exceed $\lamP$ by at most $O(1/\log M)$
(Corollary~\ref{cor:Rrate}). We do not know whether the family is lossy at fixed $M$; this is the tractable half of
Problem~\ref{prob:A}.


\subsection{Lattices of small order at $T=\tfrac12\log3$}\label{subsec:N}
The ceiling of Proposition~\ref{prop:TK} is a supremum over $\beta\in\R^2$, and turning a numerical
ceiling into a theorem requires control of $\beta$. The following a priori bound provides it: the
two primitivity multipliers of a certificate cannot be large, with a constant determined by the
geometry of the moment image of Lemma~\ref{lem:W}.

\begin{theorem}[a priori bound on the multipliers]\label{thm:Bbeta}
Fix $T$ and nonnegative weights $v$. Suppose there are $r_0>0$, $V^\star<\infty$ and finitely many
unit vectors $F_1,\dots,F_N\in H^1_0(\IT)$ such that, writing $w_i=(Q_0(F_i),Q_1(F_i))$ in the
notation of Lemma~\ref{lem:W} and $V_i=\mathfrak a_T(F_i)$,
\[
\{w\in\R^2:|w|\le r_0\}\subseteq\operatorname{conv}\{w_1,\dots,w_N\},\qquad V_i\le V^\star\ \ (1\le i\le N).
\]
Then every certificate $(\nu,\beta)$ of value $c$ satisfies $\|\beta\|_2\le(V^\star-c)/r_0$.
\end{theorem}
\begin{proof}
For a unit $F\in H^1_0(\IT)$ the measure $\mu_F=|\Fh|^2d\tau/2\pi$ is an even probability measure with
$\check\mu_F=\cF$ supported in $[-2T,2T]$ and $\cF(\pm2T)=0$, so $\int\nuh\,d\mu_F=\int\cF\,d\nu=0$ for
every measure $\nu$ supported in $\{|u|\ge2T\}$, exactly as in the proof of
Proposition~\ref{prop:weak}; moreover $\int\gh j\,d\mu_F=\int g_j\cF=Q_j(F)$ and
$\int\sT d\mu_F=\mathfrak a_T(F)$. Integrating $f\ge c$ against $\mu_F$ gives
$\mathfrak a_T(F)+\beta\cdot w(F)\ge c$. The relaxed feasible set (even probability measures whose
cosine transform vanishes off the window, with the two moment constraints dropped) is
convex, and both $\mu\mapsto\int\sT d\mu$ and $\mu\mapsto\mathfrak m(\mu)$ are affine; hence for
$w=\sum\lambda_iw_i$ in the convex hull the mixture $\mu=\sum\lambda_i\mu_{F_i}$ is admissible, has
moment $w$ and value at most $V^\star$, so $\beta\cdot w\ge c-V^\star$ for every $|w|\le r_0$.
Applying this to $w$ and $-w$ gives $|\beta\cdot w|\le V^\star-c$ on the disc of radius $r_0$, and
taking the supremum there, $r_0\|\beta\|_2\le V^\star-c$.
\end{proof}

\begin{lemma}[value of a witness]\label{lem:val}
For a unit $F\in H^1_0(\IT)$ with $L:=\|F'\|_2$ and every $u_0>0$,
\[
\mathfrak a_T(F)\ \le\ C_0+\Aconst+\frac{L^2}2\int_0^{u_0}u^2K(u)\,du+2\int_{u_0}^\infty K(u)\,du,
\]
the right side being minimized at $u_0=2/L$; the second integral is
$\log\coth\frac{u_0}2+\pi-2\arctan e^{u_0}$ by Theorem~\ref{thm:S}(iii).
\end{lemma}
\begin{proof}
By Theorem~\ref{thm:A}, $\mathfrak a_T(F)=C_0+\Lev(\cF)-\sum v_n\cF(\log n)\le C_0+\Aconst+
\int_0^\infty(1-\Rel\cF(u))K(u)du$, since $|\cF|\le\|F\|^2=1$. Now
$1-\Rel\cF(u)=\frac1{2\pi}\int(1-\cos\tau u)|\Fh|^2d\tau\le\min\{2,\frac{u^2}2\|F'\|^2\}$, the first
bound from $|\cF|\le1$ and the second from $1-\cos x\le x^2/2$ and Plancherel; split the integral at
$u_0$. The derivative in $u_0$ is $K(u_0)(\frac{L^2}2u_0^2-2)$, which vanishes and changes sign at
$u_0=2/L$.
\end{proof}

At $T=\frac12\log3$ we take as witnesses the unit vectors of the four-dimensional span of
$\psi_k(t)=(T^2-t^2)t^k$, $k=0,1,2,3$, inside $H^1_0(\IT)$: all the quantities entering
$Q_0=AB$, $Q_1=A'B-AB'$ and Lemma~\ref{lem:val} are elementary integrals of a polynomial against
$e^{\pm t/2}$. Taking, for each of $48$ directions, the top generalized eigenvector of the
corresponding combination of the two moment forms relative to the Gram matrix, computing the convex
hull of the $48$ moment vectors and the distance from the origin to its boundary gives
\[
r_0=8.4515\cdot10^{-5},\qquad V^\star=3.817495,\qquad
\|\beta\|_\infty\le\frac{V^\star}{r_0}=45169.4
\]
\[
\text{for every certificate of value }c\ge0 .
\]
The size of this constant is structural rather than a defect of the witnesses: on the window
$g_0(u)-\frac14g_1(u)=\cosh\frac u2-\frac u4\sinh\frac u2=1-\frac{u^4}{384}+O(u^6)$, so
$Q_0(F)-\frac14Q_1(F)=|\int F|^2+O(T^4)$ and the moment image is a thin sliver lying essentially on
one side of the line $Q_0=\frac14Q_1$; hence $r_0=O(T^4)$ for any witness family, and the bound is
$O(T^{-4})$. It is far too large to restrict a search, and good enough for the qualitative
statements below, because the coefficients it multiplies are of size $10^{-9}$.

\begin{theorem}[class impossibility at $T=\frac12\log3$]\label{thm:N}
Let $T=\frac12\log3$, $\NT=\{2\}$, $M\in\{1,\dots,8\}$,
$d=2T/M$, and let $\nu$ be any finite even signed measure supported on the lattice
$\{\pm jd:j\ge M\}$ (any range, any number of atoms), and $\beta\in\R^2$ arbitrary. Then
\[
\inf_{\tau\in\R}\big[\sT(\tau)+\nuh(\tau)+\beta_0\gh0(\tau)+\beta_1\gh1(\tau)\big]\ <\ 0 .
\]
Consequently no dual certificate in the sense of Definition~\ref{def:PD} supported on a lattice of
step $2T/M$ with $M\le8$ has nonnegative value, and none proves $\lamP>0$ at this window, for any
$\beta$, any range, any atom count and whichever of the tail lemmas of \S\ref{sec:tails} is used to
verify the inequality. The proof includes a computer verification in interval arithmetic
(Appendix~\ref{app:encl}).
\end{theorem}
\begin{proof}
Suppose $f:=\sT+\nuh+\beta\cdot\widehat g\ge c$ on $\R$ with $c\ge0$.

\emph{Step 1: $\beta$ is bounded.} By Theorem~\ref{thm:Bbeta} and the constants above,
$\|\beta\|_\infty\le45169.4$. All that is used below is $\|\beta\|_\infty\le B$ for some
$B\le9.7\cdot10^5$; the computed bound is about $4.5\cdot10^4$, so the hypothesis holds with a
factor $21$ to spare, and no sharp value of $V^\star/r_0$ is needed.

\emph{Step 2: one test kernel.} Fix $M$ and a pair $(\tau_0,q)$ with $\tau_0\in\Q$ and
$q\in\C^M$, and put $\vartheta(\tau)=|\sum_{k<M}q_ke^{ikd\tau}|^2\ge0$, $\mathcal I_0=[\tau_0,\tau_0+P]$,
$P=2\pi/d=\pi M/T$. As in the proof of Proposition~\ref{prop:TK}, $\nuh$ is $P$-periodic with Fourier
modes in $\{|j|\ge M\}$ while $\vartheta$ has modes in $\{|j|\le M-1\}$, so $\int_{\mathcal I_0}\nuh\vartheta=0$;
integrating $f\ge c$ against $\vartheta\ge0$ over $\mathcal I_0$ and dividing by $\int_{\mathcal I_0}\vartheta=P\|q\|^2$ (exact,
by orthogonality over a full period) gives
\[
c\ \le\ a+b_0\beta_0+b_1\beta_1,\qquad a=\frac{\int_{\mathcal I_0}\sT\vartheta}{P\|q\|^2},\quad
b_j=\frac{\int_{\mathcal I_0}\gh j\vartheta}{P\|q\|^2},
\]
valid for every $\beta$. This is Proposition~\ref{prop:TK} at a single $(\tau_0,\vartheta)$: no
infimum over $\tau_0$, no eigenvalue and no concave maximization are needed, which is what makes the
enclosure cheap. Any pair $(\tau_0,q)$ gives a valid bound, so the way the pair was found (in
practice, as a double-precision eigenvector) plays no role in the proof.

\emph{Step 3: the enclosure at $M=8$.} With $\tau_0=-27.223$ and the vector $q$ contained in the
certificate data accompanying this paper, interval arithmetic (Appendix~\ref{app:encl}) gives
\[
a\in[-6.009422\cdot10^{-3},\,-6.006218\cdot10^{-3}],\qquad |b_0|+|b_1|\le6.18\cdot10^{-9},
\]
so $c\le-6.006218\cdot10^{-3}+6.18\cdot10^{-9}B\le-5.70\cdot10^{-3}<0$ for every $B\le9.7\cdot10^5$,
contradicting $c\ge0$.

\emph{Step 4: the remaining $M$.} The same four enclosures with their own pairs $(\tau_0,q)$ give
$c\le-9.038987\cdot10^{-2}$ $(M=3)$, $-6.234677\cdot10^{-2}$ $(M=5)$, $-2.175677\cdot10^{-2}$ $(M=6)$
and $-4.242754\cdot10^{-2}$ $(M=7)$, all negative, the allowance $(|b_0|+|b_1|)B$ already included.
For $M\in\{1,2,4\}$ no computation is needed: $M\mid8$, so every multiple of $2T/M$ that is $\ge2T$ is
a multiple of $2T/8$ that is $\ge2T$, i.e.\ the class at step $2T/M$ is contained in the class at step
$2T/8$, which Step~3 already excludes.
\end{proof}

\begin{remark}
Theorem~\ref{thm:N} is not a statement about $\lamP(\frac12\log3)$, which is positive by
Theorem~\ref{thm:C}, a certificate at $M=24$; it says that the class $M\le8$ cannot detect this
positivity.
\end{remark}

\subsection{The same impossibility at the two-prime windows}\label{subsec:Np}

The enclosure of Theorem~\ref{thm:N} is limited by the size of $|a|$, which at $\frac12\log3$ is of
order $10^{-3}$. At the two-prime windows the ceilings are four orders of magnitude smaller, and
reaching them needs a quadrature rule of high order and a test kernel designed against the
$\beta$-box rather than merely against $\sT$. Both are elementary, and they give the same kind of
statement at windows where no certificate is known.

\begin{theorem}[class impossibility at the two-prime windows]\label{thm:Np}
Let $T\in\{0.68,\log2\}$, so $\NT=\{2,3\}$ and
$\Aconst=2.24883$. Let $M\in\{1,\dots,24\}$, $d=2T/M$, let $\nu$ be any finite even signed measure
supported on the lattice $\{\pm jd:j\ge M\}$ (any range, any number of atoms, any signs), and
$\beta\in\R^2$ arbitrary. Then
\[
\inf_{\tau\in\R}\big[\sT(\tau)+\nuh(\tau)+\beta_0\gh0(\tau)+\beta_1\gh1(\tau)\big]\ <\ 0 .
\]
Consequently no dual certificate in the sense of Definition~\ref{def:PD} supported on a lattice of
step $2T/M$ with $M\le24$ has nonnegative value at either window, and none proves $\lamP>0$ there,
for any $\beta$, any range, any atom count and whichever tail lemma of \S\ref{sec:tails} is used to
verify the inequality. In particular
\[
\operatorname{Ceil}(24,0.68)\le-7.493129\cdot10^{-6},\qquad
\operatorname{Ceil}(24,\log2)\le-6.584323\cdot10^{-7}.
\]
The proof includes a computer verification in interval arithmetic (Appendix~\ref{app:encl}).
\end{theorem}
\begin{proof}
The three steps are those of Theorem~\ref{thm:N}. \emph{Step 1.} Theorem~\ref{thm:Bbeta}, with its
witnesses and moment polygon recomputed at these two windows, gives $r_0=2.4309\cdot10^{-4}$,
$V^\star=4.595697$ at $T=0.68$ and $r_0=2.6720\cdot10^{-4}$, $V^\star=4.550671$ at $T=\log2$, hence
$\|\beta\|_\infty\le B$ for every certificate of value $c\ge0$ with $B=1.89\cdot10^4$, resp.
$1.70\cdot10^4$. \emph{Step 2.} For a pair $(\tau_0,q)$ put $\vartheta(\tau)=|\sum_{k<M}q_ke^{ikd\tau}|^2$,
$\mathcal I_0=[\tau_0,\tau_0+P]$, $P=\pi M/T$; as in Proposition~\ref{prop:TK}, $\int_{\mathcal I_0}\nuh\vartheta=0$ and
$c\le a+b_0\beta_0+b_1\beta_1$ with $a=(P\|q\|^2)^{-1}\int_{\mathcal I_0}\sT\vartheta$ and
$b_j=(P\|q\|^2)^{-1}\int_{\mathcal I_0}\gh j\vartheta$, so $c\le a+(|b_0|+|b_1|)B$. \emph{Step 3.} Sixteen enclosures
per window, at $M\in\{6,7,9,11,13,14,\dots,24\}$, are all negative; the two displayed values are the
upper endpoints at $M=24$, the whole $\beta$-allowance included. For the remaining
$M\in\{1,2,3,4,5,8,10,12\}$ no computation is needed: the class at step $2T/M$ is contained in the
class at step $2T/(kM)$ whenever $k\in\N$, and every such $M$ divides $20$ or $24$, both of which are
enclosed.
\end{proof}

\begin{remark}
Theorem~\ref{thm:Np} is an upper bound on what a class can certify; the sign of $\lamP(0.68)$ and of
$\lamP(\log2)$ is not decided by the methods of this paper. It shows that the class $M=24$, which
carries the certificate of Theorem~\ref{thm:C} at $\frac12\log3$, contains no certificate of
nonnegative value at either of these windows, and neither does any lattice class of order $M\le24$.
\end{remark}

\subsection{The ceilings converge}\label{subsec:R}
Proposition~\ref{prop:TK} bounds a lattice class from above by one eigenvalue. The question left
open after it is whether this bound is asymptotically sharp, i.e.\ whether the test kernels see all
of the primal problem. They do, in the limit; the mechanism is Poisson summation used in the
direction opposite to the one that makes a lattice $\nuh$ periodic.

\begin{lemma}[the periodization of a primal measure is a test kernel]\label{lem:R1}
Let $\mu$ be feasible for \textup{(P)}, $M\in\N$, $d=2T/M$, $P=2\pi/d=\pi M/T$, and let $\rho$ be the
periodization of $\mu$ on $\R/P\Z$, $\rho(E)=\sum_{k\in\Z}\mu(E+kP)$. Then $\rho$ is a probability
measure with Fourier coefficients $\widehat\rho(j)=\check\mu(jd)$, which vanish for $|j|\ge M$. Hence
$\rho$ is absolutely continuous with a nonnegative trigonometric polynomial density of degree
$\le M-1$ in $\theta=d\tau$, and by the Fej\'er--Riesz theorem
$d\rho=P^{-1}|q(e^{id\tau})|^2d\tau$ for some polynomial $q$ of degree $\le M-1$: $\rho$ is a test
kernel of order $M$ in the sense of Proposition~\ref{prop:TK}.
\end{lemma}
\begin{proof}
$\rho$ is a finite positive measure of mass $\mu(\R)=1$, and
$\widehat\rho(j)=\int_{[\tau_0,\tau_0+P)}e^{-ijd\tau}d\rho=\int_\R e^{-ijd\tau}d\mu=\check\mu(jd)$,
because $e^{-ijd(\tau+kP)}=e^{-ijd\tau}$ ($dP=2\pi$) and $\mu$ is even. Since $|jd|\ge2T$ exactly when
$|j|\ge M$, and feasibility is the closed condition $\check\mu=0$ on $\{|u|\ge2T\}$ (so the boundary
index $j=M$ is included), only finitely many coefficients are nonzero; a measure with finitely many
nonzero Fourier coefficients is a trigonometric polynomial, here nonnegative, and Fej\'er--Riesz
factorizes it.
\end{proof}

\begin{theorem}[the ceilings converge to the primal value]\label{thm:R}
For $\beta\in\R^2$ write $\operatorname{Ceil}(M,T;\beta)=\inf_{\tau_0}\lambda_{\min}(C_M(\tau_0,\beta))$
and $S_\beta=\sT+\beta\cdot\widehat g$, and put
\[
K_\beta:=2\Aconst+2|\beta_0|\gh0(0)+2|\beta_1|\gh1(0),
\]
which does not depend on $M$. Then for every $\mu$ feasible for \textup{(P)} and every $M\ge1$,
with $P/2=\pi M/2T$,
\[
\operatorname{Ceil}(M,T;\beta)\ \le\ \int\sT\,d\mu+K_\beta\,\mu\big(\{|\tau|>P/2\}\big)
\ \le\ \int\sT\,d\mu+K_\beta\,\frac{\int\sT\,d\mu+\Aconst-C_0}{m(P/2)-C_0}.
\]
Consequently $\limsup_{M\to\infty}\operatorname{Ceil}(M,T;\beta)\le\lamP$ for
every fixed $\beta$, and uniformly for $\beta$ in a bounded box.
\end{theorem}
\begin{proof}
Take $\tau_0=-P/2$, $\mathcal I_0=[-P/2,P/2]$, and let $\pi_0(\tau)=\tau-kP\in\mathcal I_0$ be the folding map. By
Lemma~\ref{lem:R1} and the Fej\'er--Riesz description of the matrix $C_M$ in the proof of
Proposition~\ref{prop:TK}, the infimum over test kernels is at most its value at $P\rho$, so
$\operatorname{Ceil}(M,T;\beta)\le\int_{\mathcal I_0}S_\beta\,d\rho=\int_\R S_\beta(\pi_0\tau)\,d\mu(\tau)$. Now
$\int S_\beta d\mu=\int\sT d\mu$ because $\mu$ has both moments zero, and $\pi_0\tau=\tau$ on
$|\tau|\le P/2$, so
\[
\int S_\beta\circ\pi_0\,d\mu-\int\sT\,d\mu=\int_{|\tau|>P/2}\big(S_\beta\circ\pi_0-S_\beta\big)d\mu
\le\Big(\sup_{\mathcal I_0}S_\beta-\inf_{|\tau|\ge P/2}S_\beta\Big)\,\mu(|\tau|>P/2).
\]
Since $m$ is increasing on $[0,\infty)$ (Lemma~\ref{lem:mbounds}(i)), $|\sum v_n\cos(\tau\log n)|\le\Aconst$ and
$|\gh j|\le\gh j(0)$, one has $\sup_{\mathcal I_0}S_\beta\le m(P/2)+\Aconst+|\beta_0|\gh0(0)+|\beta_1|\gh1(0)$ and
$\inf_{|\tau|\ge P/2}S_\beta\ge m(P/2)-\Aconst-|\beta_0|\gh0(0)-|\beta_1|\gh1(0)$; the difference is $K_\beta$.
This is the first inequality. For the second, $m-C_0\ge0$ everywhere and $m-C_0\ge m(P/2)-C_0>0$ on
$|\tau|\ge P/2$, while $\int(m-C_0)\,d\mu=\int\sT\,d\mu-C_0+\sum_{n\in\NT}v_n\int\cos(\tau\log n)\,d\mu\le\int\sT\,d\mu-C_0+\Aconst$;
Chebyshev's inequality gives $\mu(|\tau|>P/2)\le(\int\sT d\mu+\Aconst-C_0)/(m(P/2)-C_0)$. For the limit, fix
$\delta>0$ and $\mu$ with $\int\sT d\mu\le\lamP+\delta$; since $m(P/2)\to\infty$ as $M\to\infty$
(Lemma~\ref{lem:mbounds}(ii)) and $K_\beta$ is bounded for $\beta$ in a box, the right side tends to
$\lamP+\delta$, uniformly in the box.
\end{proof}

\begin{corollary}[rate of convergence of the ceilings]\label{cor:Rrate}
Let $\mu_\star$ be a minimizer of \textup{(P)}, which exists (proof of Theorem~\ref{thm:O}), so that
$\int\sT d\mu_\star=\lamP$. Then for every $\beta\in\R^2$ and every $M\ge4T$,
\[
\operatorname{Ceil}(M,T;\beta)-\lamP\ \le\ K_\beta\,\mu_\star\big(\{|\tau|>\pi M/2T\}\big)
\ \le\ \frac{K_\beta\,(\lamP+\Aconst-C_0)}{\log\frac{M}{4T}-\frac{2T}{\pi M}-\frac{2T^2}{\pi^2M^2}-C_0},
\]
the denominator being at least $5$. In particular $\operatorname{Ceil}(M,T;\beta)-\lamP=O_{\beta,T}(1/\log M)$,
and with the box $\|\beta\|_\infty\le B$ of Corollary~\ref{cor:R3},
$\operatorname{Ceil}_B(M,T)-\lamP=O_{B,T}(1/\log M)$. If the minimizer satisfies a tail bound
$\mu_\star(|\tau|>R)\le C_aR^{-a}$ for some $a>0$, the rate improves to $O(M^{-a})$.
\end{corollary}
\begin{proof}
Apply Theorem~\ref{thm:R} with $\mu=\mu_\star$ and the lower bound $m(\tau)\ge\log\frac\tau{2\pi}-\frac1\tau-\frac1{2\tau^2}$
of Lemma~\ref{lem:mbounds}(ii) at $\tau=\pi M/2T$. For $M\ge4T$ the logarithm is $\ge0$ and the two
correction terms are at most $\frac1{2\pi}+\frac1{8\pi^2}<0.18$, while $-C_0>5.37$. The last sentence is the
first inequality with $R=\pi M/2T$.
\end{proof}
The rate is logarithmic because the only information used about $\mu_\star$ is that it lies in the form domain,
$\int(m-C_0)\,d\mu_\star<\infty$; a tail bound for the primal minimizer beyond the form domain is not known. The
rate concerns the excess of the ceiling over $\lamP$; the deficit $\lamP-\operatorname{Ceil}_B(M,T)$, when
positive, measures the loss of the lattice class at fixed $M$ and is the subject of Problem~\ref{prob:A}.

\begin{corollary}[the test-kernel family is not lossy in the limit]\label{cor:R3}
Suppose $\lamP>0$, and let $B=V^\star/r_0$ be the a priori bound of Theorem~\ref{thm:Bbeta}, so that
every certificate of value $c\ge0$ has $\|\beta\|_\infty\le B$. Put
$\operatorname{Ceil}_B(M,T)=\sup_{\|\beta\|_\infty\le B}\operatorname{Ceil}(M,T;\beta)$. Then
$\operatorname{Ceil}_B(M,T)\to\lamP$ along the nested lattices $2T/2^kM_0$ with growing range. In
particular the subfamily of test kernels $\{|q|^2\}$ used in Proposition~\ref{prop:TK} is not lossy in
the limit: its optimal value is the same $\lamP$ as that of the full relaxed-feasible family.
\end{corollary}
\begin{proof}
$\limsup_M\operatorname{Ceil}_B\le\lamP$ is Theorem~\ref{thm:R}, $K_\beta$ being bounded uniformly on
the box. Conversely, by Theorem~\ref{thm:D1} there are, for every $\varepsilon>0$, lattice
certificates of value $c>\lamP-\varepsilon$ on the nested lattices (whose union is dense in
$[2T,\infty)$); eventually $c\ge0$, so by Theorem~\ref{thm:Bbeta} their $\beta$ lies in the box, and
Proposition~\ref{prop:TK} gives $c\le\operatorname{Ceil}_B(M,T)$.
\end{proof}
The ceilings exceed $\lamP$ by at most $O(1/\log M)$ (Corollary~\ref{cor:Rrate}); whether the family is lossy at
fixed $M$, i.e.\ by how much the ceilings fall short of $\lamP$, we do not know.


\subsection{Sampling a flat-top certificate}\label{subsec:lat}
Problem~\ref{prob:A} asks for the loss incurred by restricting the dual to a lattice class. For the
flat-top certificates of \S\ref{sec:Z} the loss of the sampled certificate can be computed
exactly. It is an edge effect of order $1/M$, not a smoothness effect, and it is carried entirely by
the polar block: the density of $\nu_w$ in \eqref{eq:nuw} is $\frac12(1-w)K+2w\cosh\frac u2$, whose
second summand does not vanish at $u=2T^+$ (it equals $2\cosh T$ there, the jump already met in
the proof of Theorem~\ref{thm:Dtail}), so no smoothness of $w$ can make the sampling error
$O(M^{-p})$.

\begin{theorem}[the edge cost of lattice sampling]\label{thm:latedge}
Let $p\ge2$ and let $w\in\Wcl$ be even with $w^{(p-1)}$ locally absolutely continuous and
$w^{(p)}\in \mathrm{BV}\cap L^1$, $\supp w\subseteq[-U,U]$ with $C^{p-1}$ contact with $0$ at $\pm U$,
and $w(\log n)=0$ for every $n$ with $\log n>2T$ \textup{(}comb-vanishing, finitely many linear
conditions\textup{)}. No sign condition and no bound $0\le w\le1$ is assumed. The boundary case
$\log n=2T$ is not excluded; it is treated by carrying the atom it produces \textup{(}see
$\nu_{\mathrm{lat}}$ below\textup{)}, which is legitimate because Definition~\ref{def:W} forces
$w(2T)=1$ there, so that atom is a fixed multiple of $\delta_{\pm2T}$ and can be reproduced exactly.
Put
$\Omega(u)=w(|u|)\cosh\frac u2$, $R=\frac12(1-w)K\mathbf 1_{|u|\ge2T}$ and
$G=R+2\Omega\mathbf 1_{|u|\ge2T}$, which by comb-vanishing is the whole absolutely continuous density
of $\nu_w$, the only surviving atom being the boundary one at $\pm2T$ when $2T\in\{\log n\}$; put
$V_0=\|G\|_{L^1}$, $V_p=\|R^{(p)}\|_{L^1}+2\|\Omega^{(p)}\|_{L^1}$ and $c_w=\inf_\R f_w$. For $M\in\N$ set
$d=2T/M$, $P=2\pi/d=\pi M/T$ and
\[
\nu_{\mathrm{lat}}=\sum_{|j|\ge M}d\,\kappa_j\,G(jd)\,\delta_{jd}
-\sum_{n:\,\log n=2T}\tfrac{v_n}2\big(\delta_{2T}+\delta_{-2T}\big),\qquad
\kappa_{\pm M}=\tfrac12,\quad\kappa_j=1\ (|j|>M),
\]
an even signed measure with countably many atoms and finite total variation
$\|\nu_{\mathrm{lat}}\|_{\TV}=V_0+O(d)$. Then $(\nu_{\mathrm{lat}},\beta=(2,0))$ is dual feasible and
\[
c_{\mathrm{lat}}(M)\ \ge\ \min\Big\{c_w-E_{\mathrm{edge}}(M)-\varepsilon_p(M),\
m\big(\tfrac{\pi M}{2T}\big)-\Aconst-\|\nu_{\mathrm{lat}}\|_{\TV}-2\gh0(0)\Big\},
\]
\begin{gather*}
E_{\mathrm{edge}}(M)=\frac{8\cosh T}{P}+\frac{2\pi^2\sinh T}{P^2}+\frac{16.9\cosh T}{P^3}
=\frac4\pi\cosh T\cdot d+O(d^2),\\
\varepsilon_p(M)=2\big(1-2^{-p}\big)\zeta(p)V_p\Big(\frac{\pi M}{2T}\Big)^{-p}.
\end{gather*}
Moreover $\sup_\tau|f_{\mathrm{lat}}-f_w|=\frac4\pi\cosh T\cdot d\,(1+O(d))$: the rate $M^{-1}$ and
its constant are attained, and no smoothness of $w$ improves them.
\end{theorem}
\begin{proof}
Feasibility: the atoms lie at $|u|=|j|d\ge Md=2T$, $(D)$ imposes no moment equality on $\nu$
(Proposition~\ref{prop:weak}), Definition~\ref{def:PD} allows $\nu$ to be signed, the boundary atom
at $\pm2T$ is harmless because $\cF(\pm2T)=0$, and $\beta$ is free. The atoms with $|j|<M$ are absent
because $G$ vanishes on the window. Finiteness of the total variation is
$\|\nu_{\mathrm{lat}}\|_{\TV}=d\sum_{|j|\ge M}\kappa_j|G(jd)|<\infty$, since $G\in L^1\cap
\mathrm{BV}$ on $[2T,\infty)$; the number of atoms is countably infinite, as $G=\frac12K>0$ beyond
$U$.

Value: since $w(\log n)=0$ for $\log n>2T$, $\nu_w=G(u)\,du-\sum_{\log n=2T}\frac{v_n}2(\delta_{2T}+\delta_{-2T})$,
and the second sum, present exactly when $2T\in\{\log n:n\ge2\}$ (as at the window $T=\frac12\log3$,
where it is $-\frac{v_3}2(\delta_{2T}+\delta_{-2T})$), has been put verbatim into
$\nu_{\mathrm{lat}}$, so it cancels in $f_{\mathrm{lat}}-f_w$ and may be ignored below.
Theorem~\ref{thm:Dtail}
gives $\nuh_w+2\gh0=2\widehat\Omega+\Psi_w$ with $\Psi_w=\widehat R$. Because $2T=Md$ is a lattice point,
the samples of $G$ on $|j|\ge M$ split into those of $R$ and those of $2\Omega$, so
\[
\nuh_{\mathrm{lat}}(\tau)=\Big[d\sum_{|j|\ge M}\kappa_jR(jd)e^{ijd\tau}\Big]
+\Big[d\sum_{j\in\Z}2\Omega(jd)e^{ijd\tau}-T_d(\tau)\Big],
\]
where $T_d$ is the end-point-corrected trapezoidal sum, on the nodes of the window, of
$h=2\cosh\frac u2\mathbf 1_{[-2T,2T]}$, whose transform is $\widehat h=2\gh0$. Both $R$ and $2\Omega$ are
$C^{p-1}$ with $p$-th derivative in $L^1$, so Poisson summation applies absolutely to each and the
two bracketed defects are bounded on $|\tau|\le P/2$ by
$V_p\sum_{k\ne0}|\tau+kP|^{-p}\le\varepsilon_p(M)$. Hence
$f_{\mathrm{lat}}-f_w=[2\gh0-T_d]+O(\varepsilon_p)$: the whole non-smooth cost is the quadrature
error of the polar block. Poisson summation for the compactly supported function $h$ of bounded
variation, with its symmetric values at the two jumps, gives $T_d=\sum_k\widehat h(\tau+kP)$, and
since $2TP=2\pi M$ every alias carries the same phase:
\[
2\gh0-T_d=-4\Big[\tfrac12\sinh T\cos2T\tau\sum_{k\ne0}\frac1{\xi_k^2+\frac14}
+\cosh T\sin2T\tau\sum_{k\ne0}\frac{\xi_k}{\xi_k^2+\frac14}\Big],\qquad\xi_k=\tau+kP .
\]
The first sum is at most $\pi^2/P^2$. For the second write
$\xi/(\xi^2+\frac14)=\xi^{-1}-\frac14\xi^{-1}(\xi^2+\frac14)^{-1}$ and use
$\sum_{k\ne0}\xi_k^{-1}=P^{-1}\big[\pi\cot\pi z-z^{-1}\big]$ with $z=\tau/P$, a function
regular at $z=0$ (the subtracted $z^{-1}$ is exactly the omitted $k=0$ term), odd and of
modulus at most $2/P$ on $|z|\le\frac12$; the correction is at most $4.21/P^3$. This is
$E_{\mathrm{edge}}$, uniformly in $\tau$, and it is attained where $\sin2T\tau=\pm1$. For
$|\tau|>P/2$ use $\sT\ge m(P/2)-\Aconst$ (Lemma~\ref{lem:mbounds}),
$|\nuh_{\mathrm{lat}}|\le\|\nu_{\mathrm{lat}}\|_{\TV}$ and $|2\gh0|\le2\gh0(0)$.
\end{proof}

\begin{corollary}\label{cor:ALdef}
At fixed $w$ and for $M\ge M_0(w)$ (large enough that the second branch of the minimum in
Theorem~\ref{thm:latedge} does not bind, which happens for all large $M$ since $m(\pi M/2T)\to\infty$),
the lattice sampling costs
$c_w-c_{\mathrm{lat}}(M)\le\frac8\pi T\cosh T\cdot M^{-1}+O(M^{-2})+O_p(V_pM^{-p})$, and the
$M^{-1}$ is not removable: it is the alias of the jump of $2\cosh\frac u2\mathbf 1_{|u|<2T}$ at the
window edge, an object independent of $w$. Without the end-point correction
\textup{(}$\kappa_{\pm M}=1$\textup{)} the sharp constant is
$\cosh T\sqrt{4+16/\pi^2}\,d=2.3708\cosh T\cdot d$, so the boundary atom of weight $d\cosh T$ at
$\pm2T$ improves the constant by the factor $1.862$ and nothing improves the order.
\end{corollary}

Two remarks on Corollary~\ref{cor:ALdef}. First, for sampled certificates the whole lattice cost
sits at $u=\pm2T$, the same place as the non-attainment of Theorem~\ref{thm:O}; the smooth blocks $R$
and $\Omega$ contribute only $O_p(M^{-p})$. We do not know whether the lattice class, with weights
optimized rather than sampled, also loses $\Theta(M^{-1})$; this is the remaining half of
Problem~\ref{prob:A}. Second, up to $O_p(M^{-p})$ one has $f_{\mathrm{lat}}-f_w=2\gh0-T_d$, the
discrepancy between the polar multiplier and its Riemann sum over the window; $T_d$ is a
$P$-periodic trigonometric sum of degree $M$ and $2\gh0$ is not a trigonometric polynomial, so the
difference is not one either. Heuristically, this is why Proposition~\ref{prop:TK}, whose ceiling
rests on orthogonality to every mode $|j|<M$, need not bind flat-top certificates; neither
statement implies this.

\section{The archimedean form}\label{sec:archstr}

With the primes deleted, the Weil form is an explicit convolution operator. This section studies
its thresholds, its negative index, and the number of even modes that the prime form can compete
with.

\subsection{Thresholds and nodal structure}\label{subsec:arch}

Put, for $F\in L^2(\IT)$,
\[
\Ainf(F)=2A(F)B(F)+\frac1{2\pi}\int_\R m(\tau)|\Fh(\tau)|^2\,d\tau
=2A(F)B(F)+C_0\|F\|^2+\Lev(\cF),
\]
the complete Weil form with the prime comb deleted (the second equality is Theorem~\ref{thm:A} with
$\NT=\emptyset$), and write $\PiT(F)=-\sum_{n\in\NT}v_n\Rel\cF(\log n)$ for the prime form (the real part is redundant for real $F$), so that
$W_T=\Ainf+\PiT$. Let $\laminf(T)$ be the bottom of $\Ainf$ on the unit sphere of $L^2(\IT)$ and
$\laminfP(T)$ its bottom on the unit sphere of $\PT$, where the pole term vanishes identically and
$\Ainf=C_0\|F\|^2+\Lev(\cF)$.

\begin{proposition}[two archimedean thresholds]\label{prop:arch}
\emph{(i)} $\Ainf$ does not depend on $T$ and the feasible sets increase with $T$, so both
$\laminf$ and $\laminfP$ are nonincreasing (for $\laminfP$ strictly, Lemma~\ref{lem:strict}).
\emph{(ii)} $\laminf(T)\le\laminfP(T)\le C_0+\varphi\big((3\pi/2T)^2\big)$; the first inequality is
inclusion of feasible sets and the second is the proof of Corollary~\ref{cor:mono} with the comb deleted. Hence, with
$T_0(0)=0.74921$ the root of $\varphi((3\pi/2T)^2)=-C_0$, both are nonpositive at $T=T_0(0)$ and negative for $T>T_0(0)$ (at $T_0(0)$
itself, as for every $T\ge0.60$, $\laminf\le\laminfP<0$ by Corollary~\ref{cor:Tarch}), and $\laminfP(T)\downarrow C_0$ as
$T\to\infty$, since $\Lev\ge0$ and $\varphi((3\pi/2T)^2)\downarrow0$.
\emph{(iii)} For $2T<\log2$ the contact set is empty, so $W_T=\Ainf$ and $\laminf(T)\ge0$ by the
known positivity of the complete form at autocorrelation support $\le\log2$
\textup{\cite{Yoshida1992,Bombieri2000}} (and, certified, \cite[Thm.~1.2]{Zhu2026}).
\emph{(iv)} Consequently the thresholds
$T^*_\infty=\sup\{T:\laminf(T)>0\}$ and $T^*_{\infty,\mathrm P}=\sup\{T:\laminfP(T)>0\}$ exist, and
\[
T^*_\infty\ \le\ T^*_{\infty,\mathrm P}\ \le\ 0.74921 ;
\]
for $T^*_{\infty,\mathrm P}$ this is sharpened to $0.58\le T^*_{\infty,\mathrm P}\le0.60$ by
Corollary~\ref{cor:Tarch}, and for $T^*_\infty$ to $T^*_\infty\le0.4608<0.58\le T^*_{\infty,\mathrm P}$ by
Proposition~\ref{prop:archneg}.
\end{proposition}
\begin{proof}
(i) $K$, $C_0$ and the two Tate moments are $T$-free, and an infimum over a growing set is
nonincreasing. (ii) $\PT\subset L^2(\IT)$, so the restricted infimum is the larger one; the proof of
Corollary~\ref{cor:mono} bounds the bottom of $\Lev$ on the codimension-two subspace $\PT$ by
$\lambda_3\le\varphi((3\pi/2T)^2)$ and never mentions the primes, and the comb contributes $0$.
(iii) and (iv) are immediate; a nonincreasing function changes sign once, and the thresholds are
suprema. No continuity is used. For $\laminf$ it holds by the proof of Lemma~\ref{lem:contT}, and for
$\laminfP$ it is not asserted.
\end{proof}
\begin{remark}
In (iii) the strict inequality $\laminf>0$ is not re-proved here, and that is the only reason
(iv) is stated for $T^*_\infty$ as an upper bound: strict positivity at $T=\frac12\log2$,
available in \cite{Yoshida1992}, would add the lower bound $\frac12\log2\le T^*_\infty$. For
$T^*_{\infty,\mathrm P}$ a two-sided bound is proved below (Corollary~\ref{cor:Tarch}).
\end{remark}
\begin{remark}
The archimedean form by itself was studied by Connes and Consani \cite{ConnesConsani2021}, who prove
its positivity on a window by a Toeplitz/prolate-spheroidal compression and relate the archimedean
place to the trace formula. Proposition~\ref{prop:arch} looks at the same form from the other side:
it locates the widths beyond which the archimedean form alone stops being positive, first on
$L^2(I_T)$ and then on $\PT$, so that past $T^*_{\infty,\mathrm P}$ the prime form $\PiT$ has to
supply positivity rather than merely be dominated.
\end{remark}

The bottom eigenfunction of the archimedean form on $\PT$ is never of one sign, and this is forced
by the two Tate conditions alone.

\begin{proposition}[nodal structure of $\PT$]\label{prop:nodes}
Recall that $\mT(F)=C_0\|F\|^2+\Lev(\cF)$ is the archimedean form without the pole term on
$L^2(\IT)$, so that $\Ainf=\mT$ on $\PT$. Then:
\emph{(i)} the bottom of $\mT$ on all of $L^2(\IT)$ is a simple eigenvalue whose
eigenfunction can be taken a.e.\ strictly positive;
\emph{(ii)} say that $F\in L^2(\IT)$ has \emph{at most one sign change} if for some
$t_0\in[-T,T]$ one has $F\ge0$ a.e.\ on $(-T,t_0)$ and $F\le0$ a.e.\ on $(t_0,T)$, or the same
with the two inequalities exchanged (the endpoints $t_0=\pm T$ being the case of one sign).
Then every nonzero $F\in\PT$ has at least two sign changes in this sense; in particular no
nonzero element of $\PT$ is of one sign, and an even element has at least one nodal pair in
$(0,T)$;
\emph{(iii)} consequently the bottom eigenfunction of $\mT$ on $\PT$ is not of one sign,
and the Sturm index of $\Ainf|\PT$ is shifted by one nodal pair relative to $(i)$.
\end{proposition}
\begin{proof}
(i) For $F\in L^2(\IT)$ extended by zero, Theorem~\ref{thm:A} and Fubini give the pure-jump
Dirichlet form with killing
\[
\begin{aligned}
\Lev(\cF)&=\frac14\iint_{\IT\times\IT}K(|x-y|)\,|F(x)-F(y)|^2\,dx\,dy
+\frac12\int_{\IT}\kappa(x)|F(x)|^2\,dx,\\
\kappa(x)&=\int_{|y|>T}K(|x-y|)\,dy<\infty,
\end{aligned}
\]
obtained by splitting $2\Lev(\cF)=\int_\R K(|u|)\big(\|F\|^2-\Rel\cF(u)\big)du
=\frac12\iint_{\R^2}K(|x-y|)|\widetilde F(x)-\widetilde F(y)|^2$ into the square $\IT\times\IT$
and the two strips on which the zero extension $\widetilde F$ vanishes; the divergence of $K$ at
the origin is absorbed by $|F(x)-F(y)|^2$ on the form domain, and no $\delta$-kernel exists,
$\int_0^{2T}K=+\infty$. The form is therefore a regular Dirichlet form of pure-jump type with
nonnegative jump density $K$ and killing $\kappa$, up to the additive constant $C_0\|F\|^2$,
which only shifts the spectrum. By the Beurling--Deny criteria the semigroup $e^{-t\mT}$ is
positivity preserving, and since $K>0$ everywhere it is irreducible, hence positivity improving;
Perron--Frobenius then gives a simple, a.e.\ strictly positive ground state, provided the bottom
is an eigenvalue. It is: $m$ is increasing (Lemma~\ref{lem:mbounds}(i)) and $m(\tau)\to\infty$, so
$\int_{|\tau|>R}|\Fh|^2\le2\pi\,(m(R)-C_0)^{-1}\Lev(\cF)$ uniformly on form-bounded sets, while all
competitors are supported in the fixed compact $\IT$; by the Riesz--Kolmogorov criterion the form
domain embeds compactly in $L^2(\IT)$ and the resolvent is compact.
(ii) The pair $\{e^{t/2},e^{-t/2}\}$ is a Chebyshev system of order $2$ on $(-T,T)$: a nonzero
combination $c_1e^{t/2}+c_2e^{-t/2}$ vanishes only where $e^{t}=-c_2/c_1$, i.e.\ at most once. Let
$0\ne F\in\PT$ and suppose $F$ had at most one sign change, at $t_0$. Choose
$G=c_1e^{t/2}+c_2e^{-t/2}$ vanishing exactly at $t_0$ (any nonzero $G$ of the family if
$t_0=\pm T$), with the sign of $G$ matched to that of $F$ on each side. Then $FG\ge0$ a.e.\ and
$FG$ is not a.e.\ zero, so $\int FG\ne0$; but $F\in\PT$ means $\langle F,e^{t/2}\rangle=\langle
F,e^{-t/2}\rangle=0$, whence $\int FG=0$, a contradiction. For even $F$ the two sign changes are
symmetric, so at least one lies in $(0,T)$.
(iii) Immediate from (ii).
\end{proof}

\emph{Parity sectors.} For even $F$ the two Tate moments coincide,
$A(F)=B(F)=\int F(t)\cosh\frac t2\,dt$, and for odd $F$ one has $A(F)=-B(F)=\int F(t)\sinh\frac t2\,dt$.
Write $\PT^{\mathrm{ev}}$ and $\PT^{\mathrm{odd}}$ for the even and the odd elements of $\PT$
($\PT^{\mathrm{ev}}$ is the space of Theorem~\ref{thm:arena} below), $\lambda_{\mathrm{ev}}(T)$ and
$\lambda_{\mathrm{odd}}(T)$ for the bottoms of $\Ainf=\mT$ on their unit spheres (so
$\lambda_{\mathrm{ev}}=\lambda^{\mathrm{ev}}_1$ in the notation of Theorem~\ref{thm:arena}),
$n_{\mathrm{ev}}(T)$ and $n_{\mathrm{odd}}(T)$ for the numbers of negative eigenvalues of $\Ainf$ on them,
and $T^*_{\mathrm{ev}}=\sup\{T:\lambda_{\mathrm{ev}}(T)>0\}$, $T^*_{\mathrm{odd}}=\sup\{T:\lambda_{\mathrm{odd}}(T)>0\}$.
By Lemma~\ref{lem:parity} below, $\laminfP=\min(\lambda_{\mathrm{ev}},\lambda_{\mathrm{odd}})$
and $n_-(\Ainf|\PT)=n_{\mathrm{ev}}+n_{\mathrm{odd}}$. For a parity label $\mathrm{par}\in\{\mathrm{ev},\mathrm{odd}\}$
we write $\PT^{\mathrm{par}}$, $\lambda_{\mathrm{par}}$, $n_{\mathrm{par}}$, $T^*_{\mathrm{par}}$.

\begin{lemma}[parity splitting]\label{lem:parity}
$\PT=\PT^{\mathrm{ev}}\oplus\PT^{\mathrm{odd}}$, the sum being orthogonal both in
$L^2(\IT)$ and for $\mT$, which equals $\Ainf$ on $\PT$. Hence
\[
n_-(\Ainf|\PT)=n_{\mathrm{ev}}(T)+n_{\mathrm{odd}}(T),\qquad
\laminfP(T)=\min\big(\lambda_{\mathrm{ev}}(T),\lambda_{\mathrm{odd}}(T)\big).
\]
Here $n_-$ counts negative eigenvalues with multiplicity; for the closed,
lower-bounded form $\mT$, which has compact resolvent on each of these closed subspaces (proof of
Proposition~\ref{prop:nodes}(i)), this is the largest dimension of a subspace of the form domain on
which the form is negative definite.
\end{lemma}
\begin{proof}
$A(JF)=B(F)$ and $B(JF)=A(F)$, so for $F\in\PT$ both $F_{\mathrm e}=\frac12(F+JF)$ and
$F_{\mathrm o}=\frac12(F-JF)$ have vanishing Tate moments: $\PT$ is $J$-invariant and splits. The
parts are $L^2$-orthogonal, and
$\mT(F_{\mathrm e},F_{\mathrm o})=\frac1{2\pi}\int m\,\widehat{F_{\mathrm e}}\,\overline{\widehat{F_{\mathrm o}}}=0$,
because $m$ is even, $\widehat{F_{\mathrm e}}$ is even and $\widehat{F_{\mathrm o}}$ is odd. Negative
indices add over sums that are orthogonal for the norm and for the form, and bottoms take the minimum.
\end{proof}
The definition of $\Ainf$ writes the pole term as $2A(F)B(F)$, i.e.\ on real $F$, while $L^2(\IT)$ is
complex. Nothing depends on the choice: for $F=u+iv$ with $u,v$ real, $\mT(F)=\mT(u)+\mT(v)$
and $F\in\PT$ if and only if $u,v\in\PT$ (as in the proof of Lemma~\ref{lem:strict}), so the counts and
the bottoms are the same over $\R$ and over $\C$.

\begin{lemma}[zero extension]\label{lem:zeroext}
For $0<T<T'$, extension by zero maps $\PT^{\mathrm{par}}$ isometrically into
$\mathcal P^{\mathrm{par}}_{T'}$ with the same values of $\mT$. Hence $n_{\mathrm{par}}$ is
nondecreasing and $\lambda_{\mathrm{par}}$ nonincreasing in $T$, with no range condition at either window.
\end{lemma}
\begin{proof}
The norm, the Tate moments, the parity and $\cF$ are unchanged, and $\mT(F)=C_0\|F\|^2+\Lev(\cF)$
depends on $F$ only through $\cF$ on $\R$ (Theorem~\ref{thm:A}); $\Ainf$ does not depend on $T$
(Proposition~\ref{prop:arch}(i)).
\end{proof}

\begin{lemma}[strict monotonicity]\label{lem:strict}
For every $T>0$ the bottoms $\lambda_{\mathrm{ev}}(T)$, $\lambda_{\mathrm{odd}}(T)$ and
$\laminfP(T)$ are attained, and all three are strictly decreasing on $(0,\infty)$. Consequently, for each
sector, $\lambda_{\mathrm{par}}(T)>0$ for $T<T^*_{\mathrm{par}}$ and $\lambda_{\mathrm{par}}(T)<0$ for
$T>T^*_{\mathrm{par}}$; $\lambda_{\mathrm{par}}(T)<0$ if and only if $n_{\mathrm{par}}(T)\ge1$;
$T^*_{\infty,\mathrm P}=\min(T^*_{\mathrm{ev}},T^*_{\mathrm{odd}})$; and $\laminfP$ is positive on
$(0,T^*_{\infty,\mathrm P})$ and negative on $(T^*_{\infty,\mathrm P},\infty)$.
\end{lemma}
\begin{proof}
\emph{Attainment.} $\mT-C_0\|\cdot\|^2=\Lev(c_\cdot)\ge0$ is a closed form whose domain embeds
compactly in $L^2(\IT)$ (proof of Proposition~\ref{prop:nodes}(i)); $\PT^{\mathrm{par}}$ and $\PT$ are
closed, $A$ and $B$ being continuous on $L^2(\IT)$. A minimizing sequence on the unit sphere has an
$L^2$-convergent subsequence whose limit has norm $1$ and lies in the subspace, and lower semicontinuity
of the closed form gives attainment. Since the symbol is real and even, $\mT(u+iv)=\mT(u)+\mT(v)$
for real $u,v$, and $u+iv\in\PT^{\mathrm{par}}$ if and only if $u,v\in\PT^{\mathrm{par}}$ ($A,B$ are real on
real functions), so the minimizer may be taken real.

\emph{Strictness.} Monotonicity is Proposition~\ref{prop:arch}(i), sector by sector: extension by zero
maps $\mathcal P^{\mathrm{par}}_{T'}$ isometrically into $\PT^{\mathrm{par}}$ for $T'<T$ with the same
values of $\mT$ (Lemma~\ref{lem:zeroext}). Suppose $T'<T$ and
$\lambda_{\mathrm{par}}(T)=\lambda_{\mathrm{par}}(T')=:\lambda$, and let $F^*$ be a real minimizer on
$\mathcal P^{\mathrm{par}}_{T'}$. Extended by zero it is a minimizer on $\PT^{\mathrm{par}}$, so
$\mT(F^*,G)=\lambda\langle F^*,G\rangle$ for every $G\in\PT^{\mathrm{par}}$ in the form domain. Take
real $G_0\in C_c^\infty(T',T)$ with $\int G_0e^{y/2}dy=\int G_0e^{-y/2}dy=0$ and $G=G_0+\epsilon JG_0$,
$\epsilon=\pm1$ the parity sign. Then $G\in\PT^{\mathrm{par}}$, $\langle F^*,G\rangle=0$ by disjoint
supports, and $\mT(F^*,G)=2\mT(F^*,G_0)$, since $\mT$ is $J$-invariant and
$JF^*=\epsilon F^*$. Polarizing the Dirichlet form of Proposition~\ref{prop:nodes}(i),
$\Lev(\cF)=\frac14\iint_{\R^2}K(|x-y|)|\widetilde F(x)-\widetilde F(y)|^2$, for real $F,G$ with
supports at positive distance gives
$\mT(F,G)=-\frac12\iint K(|x-y|)F(x)G(y)\,dx\,dy$, absolutely convergent. Hence
$\Phi(y):=\int_{I_{T'}}K(|x-y|)F^*(x)\,dx$ satisfies $\int\Phi G_0=0$ for every such $G_0$, so
$\Phi=\alpha e^{y/2}+\beta e^{-y/2}$ on $(T',T)$. For $y>T'\ge x$, Lemma~\ref{lem:K} and Tonelli give
$\Phi(y)=2\sum_{k\ge0}c_ke^{-(2k+\frac12)y}$ with $c_k=\int F^*(x)e^{(2k+\frac12)x}dx$,
$|c_k|\le\|F^*\|_1e^{(2k+\frac12)T'}$; so $e^{y/2}\Phi(y)=2\sum_kc_kw^k$, $w=e^{-2y}$, is a convergent
power series for $0<w<e^{-2T'}$, and the identity extends to all $y>T'$. As $y\to\infty$, $\Phi\to0$,
so $\alpha=0$; then $2\sum_kc_kw^k\equiv\beta$, so $c_k=0$ for $k\ge1$ (the value $c_0=\beta/2$ is not
needed). With $u=e^{2x}$, $c_k=\int f(u)u^k\,du$ for
$f(u)=\frac12F^*(\frac12\log u)u^{-3/4}\in L^1[e^{-2T'},e^{2T'}]$, so all moments of $uf(u)$ vanish and
$uf=0$ by Weierstrass, whence $F^*=0$, contradicting $\|F^*\|=1$. For $\laminfP$ use Lemma~\ref{lem:parity}.

\emph{Consequences.} If $T>T^*_{\mathrm{par}}$, pick $T'\in(T^*_{\mathrm{par}},T)$: then
$\lambda_{\mathrm{par}}(T')\le0$ and $\lambda_{\mathrm{par}}(T)<\lambda_{\mathrm{par}}(T')$. If
$T<T^*_{\mathrm{par}}$, some $T''>T$ has $\lambda_{\mathrm{par}}(T'')>0$, and monotonicity gives
$\lambda_{\mathrm{par}}(T)>0$. A negative vector forces $\lambda_{\mathrm{par}}<0$, and conversely an
attained negative bottom spans a negative line. The last two statements follow from
$\laminfP=\min(\lambda_{\mathrm{ev}},\lambda_{\mathrm{odd}})$, the sets $\{\lambda_{\mathrm{par}}>0\}$ being
intervals.
\end{proof}
The same argument, with the pole term of $\Ainf$ in place of the two multipliers, shows that $\laminf$
and the bottom of $\mT$ on $L^2(\IT)$ are strictly decreasing as well (attainment as above, the pole
term being a bounded form of rank two); this is not used below.

\begin{theorem}[the negative index of the archimedean form, $0<T\le0.90$]\label{thm:index}
In the notation of \S\ref{subsec:arch}:
\emph{(i)} $n_{\mathrm{ev}}(T)\le1$ and $n_{\mathrm{odd}}(T)\le1$ for $0<T\le0.90$;
$n_{\mathrm{ev}}(T)=n_{\mathrm{odd}}(T)=0$ and $\laminfP(T)>0$ for $T\le0.58$; $n_{\mathrm{odd}}(T)=0$
and $\lambda_{\mathrm{odd}}(T)>0$ for $T\le\log2$.
\emph{(ii)} $n_{\mathrm{ev}}(T)\ge1$ for $T\ge0.60$, and $n_{\mathrm{odd}}(T)\ge1$ for $T\ge0.81$.
\emph{(iii)} $0.58\le T^*_{\mathrm{ev}}\le0.60$, $\log2\le T^*_{\mathrm{odd}}\le0.81$,
$T^*_{\infty,\mathrm P}=T^*_{\mathrm{ev}}$, and
\[
n_-(\Ainf|\PT)=\begin{cases}0,&0<T<T^*_{\mathrm{ev}},\\ 1,&T^*_{\mathrm{ev}}<T<T^*_{\mathrm{odd}},\\
2,&T^*_{\mathrm{odd}}<T\le0.90.\end{cases}
\]
\emph{(iv)} At the windows: $n_-(\Ainf|\PT)=0$ at $T=\frac12\log3$ and $T=0.58$; $n_-(\Ainf|\PT)=1$ at
$T=0.60$ and $T=\log2$, the negative direction being even; $n_-(\Ainf|\PT)\in\{1,2\}$ at $T=0.80$,
with $n_{\mathrm{ev}}(0.80)=1$; and $n_-(\Ainf|\PT)=2$ at $T=0.81$, $0.82$ and $0.90$.

\noindent The proof includes a computer verification in interval arithmetic (Appendix~\ref{app:encl}).
\end{theorem}
\begin{proof}
(i) The upper certificates of Proposition~\ref{prop:haynsworth} (Appendix~\ref{app:encl}) give
$n_{\mathrm{par}}(T)\le r$ with $r=0$ at $T=0.58$ in both sectors and at $T=\log2$ in the odd sector,
and $r=1$ at $T=0.90$ in both sectors. Lemma~\ref{lem:zeroext} carries each bound to every smaller $T$,
and Lemma~\ref{lem:rzero} turns $r=0$ into $\lambda_{\mathrm{par}}>0$, at the certified window and
therefore, by monotonicity, below it. (ii) The explicit negative planes of
Proposition~\ref{prop:negplane} at $T=0.60$ (even) and $T=0.81$ (odd) give $n_{\mathrm{par}}\ge1$
there, and Lemma~\ref{lem:zeroext} carries each to every larger $T$, with no range condition at the
larger window. (iii) By (i), (ii) and Lemma~\ref{lem:strict}: $\lambda_{\mathrm{ev}}(0.58)>0$ gives
$T^*_{\mathrm{ev}}\ge0.58$, $\lambda_{\mathrm{ev}}(0.60)<0$ gives $T^*_{\mathrm{ev}}\le0.60$,
$\lambda_{\mathrm{odd}}(\log2)>0$ gives $T^*_{\mathrm{odd}}\ge\log2$, $\lambda_{\mathrm{odd}}(0.81)<0$
gives $T^*_{\mathrm{odd}}\le0.81$, and $T^*_{\infty,\mathrm P}=\min(T^*_{\mathrm{ev}},T^*_{\mathrm{odd}})
=T^*_{\mathrm{ev}}$ because $0.60<\log2$. By Lemma~\ref{lem:strict}, off the threshold points
$n_{\mathrm{par}}(T)\ge1$ exactly when $T>T^*_{\mathrm{par}}$, and $n_{\mathrm{par}}\le1$ on $(0,0.90]$
by (i); add over the sectors (Lemma~\ref{lem:parity}). (iv) is (i)--(iii) at the windows:
$\frac12\log3=0.5493<0.58$; at $0.60$ and $\log2$ the even count is $1$ (at least $1$ from $0.60$, at
most $1$ from $0.90$) and the odd count is $0$ (from $\log2$); at $0.80$ the even count is $1$ and the
odd count is at most $1$; at $0.81$, $0.82$, $0.90$ both counts are $1$.
\end{proof}
The values $n_{\mathrm{par}}(T^*_{\mathrm{par}})$
at the two thresholds are not decided, and nothing depends on them; nor is it decided which of the two
values $n_-$ takes at $T=0.80$. Theorem~\ref{thm:index} concerns the
archimedean form with the prime comb deleted; it says nothing about the sign of $\mathfrak a_T$ or
$W_T$.


\begin{corollary}[the archimedean threshold]\label{cor:Tarch}
\[
0.58\ \le\ T^*_{\infty,\mathrm P}\ \le\ 0.60 .
\]
The bottom $\laminfP$ is strictly decreasing, positive on $(0,T^*_{\infty,\mathrm P})$ and negative on
$(T^*_{\infty,\mathrm P},\infty)$, and $\laminfP(T)<0$ for every $T\ge0.60$. In particular
$\laminfP(\frac12\log3)>0$: $\Ainf\ge\laminfP(\frac12\log3)\,\|F\|^2$ with a positive constant, i.e.\
the archimedean form is positive definite on $\mathcal P_{\frac12\log3}$ by itself. The proof
includes a computer verification in interval arithmetic (Appendix~\ref{app:encl}).
\end{corollary}
\begin{proof}
Theorem~\ref{thm:index}(iii) gives
$T^*_{\infty,\mathrm P}=T^*_{\mathrm{ev}}\in[0.58,0.60]$, and Theorem~\ref{thm:index}(ii) gives
$n_{\mathrm{ev}}(T)\ge1$, hence $\laminfP(T)\le\lambda_{\mathrm{ev}}(T)<0$, for $T\ge0.60$
(Lemma~\ref{lem:strict}); the rest is Lemma~\ref{lem:strict}, with $\frac12\log3=0.5493<0.58$.
\end{proof}
Corollary~\ref{cor:Tarch} proves a sign, not a value. Its lower half is a proved lower bound on an
infimum, obtained from an explicit minorant of the symbol, an exact finite section and a Schur
complement over the tail controlled by the trace of Lemma~\ref{lem:KyFan}; it improves the bound
$0.74921$ of Proposition~\ref{prop:arch}(iv).

On all of $L^2(\IT)$ the pole term $2A(F)B(F)$ is present, and one explicit odd test function makes
the archimedean form negative well before $T^*_{\infty,\mathrm P}$, so that the two thresholds are distinct:
$T^*_\infty<0.58\le T^*_{\infty,\mathrm P}$.

\begin{proposition}[the archimedean complete form is negative from $T=0.4608$]\label{prop:archneg}
For $T>0$ let $F^{\mathrm{odd}}_T(t)=t$ on $\IT$ and put
\[
U(T):=\frac{2T^3}3\Big(C_0+\log\coth\frac T2+2\arctan e^{-T}\Big)+\frac{2T^3(3T+7)}9
-2\Big(4T\cosh\frac T2-8\sinh\frac T2\Big)^2 .
\]
Then $\Ainf(F^{\mathrm{odd}}_T)\le U(T)$, and $U(0.4608)\in[-7.0159575,-7.0159574]\cdot10^{-6}$. Hence
$\laminf(T)<0$ for every $T\ge0.4608$, and
\[
T^*_\infty\ \le\ 0.4608\ <\ 0.58\ \le\ T^*_{\infty,\mathrm P} .
\]
The single numerical evaluation in the proof is done in interval arithmetic.
\end{proposition}
\begin{proof}
Write $F=F^{\mathrm{odd}}_T$. It is real, odd and bounded, and $|\Fh(\tau)|=O(1/|\tau|)$, so it lies in the
form domain. By the parity paragraph after Proposition~\ref{prop:nodes}, $A(F)=-B(F)=s:=\int F\sinh\frac t2$,
so the pole term is $2A(F)B(F)=-2s^2$, and $s=2\int_0^Tt\sinh\frac t2\,dt=4T\cosh\frac T2-8\sinh\frac T2$.
For $0\le u\le2T$, $\cF(u)=\int_{u-T}^Tt(t-u)\,dt=(4T^3-6T^2u+u^3)/6$, so $\|F\|^2=\cF(0)=2T^3/3$ and
$\cF(0)-\cF(u)=u(6T^2-u^2)/6\ge0$; for $u\ge2T$, $\cF(u)=0$. On $u>0$ one has $K(u)\le\frac1u+\frac12$,
since $\sinh x\ge x$ and $\cosh x\ge1$. Against the nonnegative function $\cF(0)-\cF(u)$ this gives
\begin{align*}
\Lev(\cF)&\le\int_0^{2T}\frac{u(6T^2-u^2)}6\Big(\frac1u+\frac12\Big)du+\|F\|^2\int_{2T}^\infty K\\
&=\frac{2T^3(3T+7)}9+\frac{2T^3}3\Big(\log\coth\frac T2+2\arctan e^{-T}\Big),
\end{align*}
by Theorem~\ref{thm:S}(iii) for the tail, with $\pi-2\arctan e^T=2\arctan e^{-T}$.
Adding $C_0\|F\|^2$ and $-2s^2$ gives $\Ainf(F)\le U(T)$. The value of $U(0.4608)$ is obtained by an
outward-rounded interval evaluation at $40$ digits, with $C_0=-\gamma-\frac\pi2-3\log2-\log\pi$ and
$2\arctan e^{-T}$ enclosed by its alternating series with the remainder term; the expressions for
$\cF$, the integral over the window and $s$ are also checked symbolically. So
$\laminf(0.4608)\le\Ainf(F)/\|F\|^2<0$. The rest is Proposition~\ref{prop:arch}(i), $\laminf$ being
nonincreasing, together with Corollary~\ref{cor:Tarch}.
\end{proof}
The same evaluation gives $U(0.461)\in[-3.59383,-3.59382]\cdot10^{-5}$, $U(0.47)=-1.4048\cdot10^{-3}$
and $U(\frac12\log3)=-0.0200526$. Proposition~\ref{prop:archneg} separates the two thresholds
strictly. It is the empty-comb case of Corollary~\ref{cor:convfail}(i) below: on $[0.4608,0.58]$ the
archimedean form is positive definite on $\PT$ but not on $L^2(\IT)$.

For every $T\ge0.60$, in particular at $T=\log2$, one has $\laminfP<0$, so every proof of $\lamA\ge0$
has to exhibit $\PiT(F)\ge-\Ainf(F)>0$ on the vectors on which the archimedean block is negative; at
$\log2$, by Theorem~\ref{thm:index}, this set contains a line but no two-dimensional subspace. Beyond the
threshold the prime form is not a perturbation of a positive form but the source of its positivity.


\subsection{Bounds on the negative index for every window}\label{subsec:negindex}

Theorem~\ref{thm:index} computes the negative index for $T\le0.90$ by certified computation. The
following bounds are explicit in $T$. Throughout, $n_-$ of a closed lower-bounded form with compact
resolvent on a closed subspace is the number of its negative eigenvalues with multiplicity, i.e.\ the
largest dimension of a subspace of the form domain on which the form is negative definite, and
$\tau_0=6.28983599\ldots$ is the zero of $m$ on $(0,\infty)$.

\begin{proposition}[negative index of the archimedean form]\label{prop:NC}
For $\tau_1>0$ let $S_{T,\tau_1}=Q_TB_{\tau_1}Q_T$ be the time--band limiting operator on $L^2(\IT)$,
$\langle S_{T,\tau_1}F,F\rangle=\frac1{2\pi}\int_{|\tau|\le\tau_1}|\Fh|^2d\tau$, and, for $\tau_1>\tau_0$,
$\theta(\tau_1)=m(\tau_1)/(m(\tau_1)-C_0)\in(0,1)$. Then
\[
n_-(\Ainf|\PT)\ \le\ n_-(\mT)\ \le\ \#\{n:\lambda_n(S_{T,\tau_1})>\theta(\tau_1)\}\ \le\
\Big\lfloor\frac{2T\tau_1}{\pi\,\theta(\tau_1)}\Big\rfloor
\qquad\text{for every }\tau_1>\tau_0,
\]
$n_-(\mT)$ being the index of $\mT$ on $L^2(\IT)$. At $\tau_1=14.8815$ the right-hand side is
$\lfloor68.513\,T\rfloor$: at most $37$ at $T=\frac12\log3$, $47$ at $T=\log2$ and $54$ at $T=0.80$.
\end{proposition}
\begin{proof}
On $\PT$ the pole term vanishes and $\Ainf=\mT$, and restricting a form to a subspace does not increase its
index; this is the first inequality. For $\|F\|=1$ put $b=\langle S_{T,\tau_1}F,F\rangle$. Splitting the
archimedean integral at $\tau_1$ and using that $m$ is even and increasing on $[0,\infty)$ with $m(0)=C_0$,
$\mT(F)\ge C_0b+m(\tau_1)(1-b)$. Hence $\mT(F)<0$ forces $b>\theta(\tau_1)$. If $\mT$ were negative
definite on an $n$-dimensional subspace $V$, every unit $F\in V$ would have
$\langle S_{T,\tau_1}F,F\rangle>\theta$, so by min--max $\lambda_{n-1}(S_{T,\tau_1})>\theta$; this is the
second inequality. The third is Markov's inequality with $\operatorname{tr}S_{T,\tau_1}=2T\tau_1/\pi$
(Lemma~\ref{lem:KyFan} with $c=T\tau_1$). The numerical constant is $2\tau_1/(\pi\theta(\tau_1))$ at
$\tau_1=14.8815$, evaluated with $m=\Rel\psi(\frac14+\frac{i\tau}2)-\log\pi$.
\end{proof}

\begin{corollary}[explicit negative-index counts]\label{cor:P3}
For every $\tau_1>\tau_0$, $n_-(\Ainf|\PT)\le n_-(\mT)\le\min\{k:\
\Lambda_k(T\tau_1)\le\theta(\tau_1)\}$, and the same with $\Lambda^*_k$ in place of $\Lambda_k$
(Lemma~\ref{lem:KyFan}, Corollary~\ref{cor:P1}). With the quadrature-free $\Lambda^*$ this gives
\[
n_-(\Ainf|\PT)\le n_-(\mT)\ \le\ 5\ \ \big(T=\tfrac12\log3\big),\qquad 7\ \ (T=\log2),\qquad 8\ \ (T=0.80),
\]
against $37$, $47$ and $54$ from the trace bound of Proposition~\ref{prop:NC}.
\end{corollary}
\begin{proof}
If $\Lambda_k(T\tau_1)\le\theta(\tau_1)$ then $\lambda_k(S_{T,\tau_1})\le\theta(\tau_1)$ by
Lemma~\ref{lem:KyFan}, so $\#\{n:\lambda_n>\theta\}\le k$, and Proposition~\ref{prop:NC} applies for
that $\tau_1$. The three values come from $\tau_1=6.6159$, $6.3576$, $6.3576$, where the partial sums of
the series $\Lambda^*_k$ fall below $\theta(\tau_1)$; the omitted tails are below $10^{-170}$.
\end{proof}

\begin{remark}[the limit of this method]\label{rem:P4}
Numerically, at $T=\log2$ and the optimal $\tau_1$ the exact count
$\#\{n:\lambda_n(S_{T,\tau_1})>\theta(\tau_1)\}$ is already $5$, and a scan over $\tau_1$ does not
lower it (double precision, not certified). So no improvement of the prolate estimate can bring the
bound of Proposition~\ref{prop:NC} below $5$ there: the distance to the true value
$n_-(\Ainf|\mathcal P_{\log2})=1$ (Theorem~\ref{thm:index}) is due to the two-value minorant
$m\ge C_0\mathbf 1_{|\tau|<\tau_1}+m(\tau_1)\mathbf 1_{|\tau|\ge\tau_1}$ and to $\operatorname{codim}\PT=2$, not to
prolate theory. A multi-band minorant $m\ge\sum_ic_i\mathbf 1_{B_i}$ produces a different operator and a
different counting problem; with $28$ bands, handled exactly on a finite Legendre section of each parity
sector and with Lemma~\ref{lem:KyFan} used only through the trace of the far tail, it returns the exact
count: these are the certificates behind Theorem~\ref{thm:index} (Appendix~\ref{app:encl}).
\end{remark}

In the other direction the count admits a lower bound, from the same two ingredients as
Corollary~\ref{cor:mono}: Dirichlet modes and the concavity of $\varphi$.

\begin{theorem}[lower bound on the negative index]\label{thm:J}
\[
n_-(\mT)\ \ge\ \#\{k\ge1:\ k\pi/2T<\tau_0\}=\Big\lceil\frac{2T\tau_0}{\pi}\Big\rceil-1
\quad\text{on }L^2(\IT),\qquad
n_-(\Ainf|\PT)\ \ge\ \Big\lceil\frac{2T\tau_0}{\pi}\Big\rceil-3 .
\]
\end{theorem}
\begin{proof}
Let $D_k\subset H^1_0(\IT)$ be the span of the first $k$ Dirichlet modes $\sin(j\pi(t+T)/2T)$,
$j\le k$. For $F\in D_k$, extended by zero, Plancherel gives
$\frac1{2\pi}\int\tau^2|\Fh|^2=\|F'\|^2\le(k\pi/2T)^2\|F\|^2$. By Theorem~\ref{thm:S}(i),
$m(\tau)-C_0=\varphi(\tau^2)$ with $\varphi$ complete Bernstein, hence concave and strictly
increasing; Jensen's inequality in the variable $\tau^2$ for the probability measure
$d\mu=|\Fh|^2d\tau/(2\pi\|F\|^2)$ gives
\[
\frac{\mT(F)}{\|F\|^2}=C_0+\int\varphi(\tau^2)\,d\mu\ \le\ C_0+\varphi\Big(\int\tau^2d\mu\Big)
=C_0+\varphi\big(\|F'\|^2/\|F\|^2\big)\ \le\ C_0+\varphi\big((k\pi/2T)^2\big).
\]
Since $m(\tau_0)=0$ means $\varphi(\tau_0^2)=-C_0$, and $\varphi$ is strictly increasing,
$k\pi/2T<\tau_0$ forces $\mT<0$ on $D_k\setminus\{0\}$, i.e.\ $n_-(\mT)\ge k$; the
largest such $k$ is $\lceil2T\tau_0/\pi\rceil-1$. On $\PT$ the pole term vanishes and
$\Ainf=\mT$; $D_k\cap\PT$ has dimension at least $k-2$.
\end{proof}
The bound on $\PT$ is vacuous at $T=\frac12\log3$ and $T=\log2$: $2T\tau_0/\pi$ equals $2.1996$ and
$2.7755$ there. On $L^2(\IT)$ it reads $n_-(\mT)\ge2$ at both, so Theorem~\ref{thm:J} and
Corollary~\ref{cor:P3} give the brackets $n_-(\mT)\in[2,5]$ at $T=\frac12\log3$, $[2,7]$ at $T=\log2$ and
$[3,8]$ at $T=0.80$.

\begin{corollary}[asymptotic negative index]\label{cor:nasym}
As $T\to\infty$, $n_-(\Ainf|\PT)=\dfrac{2\tau_0}{\pi}\,T\,(1+o(1))=4.0042\,T\,(1+o(1))$, and the same holds for $n_-(\mT)$.
\end{corollary}
\begin{proof}
The lower bound is Theorem~\ref{thm:J}. For the upper bound fix $\tau_1>\tau_0$, so that
$\theta=\theta(\tau_1)\in(0,1)$ is fixed. By the Landau--Widom theorem \cite{LandauWidom1980},
$\#\{n:\lambda_n(S_{T,\tau_1})>\theta\}=\frac{2T\tau_1}\pi+O_{\tau_1}(\log T)$ as $T\to\infty$; by
Proposition~\ref{prop:NC}, $\limsup_T n_-(\mT)/T\le2\tau_1/\pi$. Let $\tau_1\downarrow\tau_0$.
\end{proof}

For $T_2>T^*_{\infty,\mathrm P}$ write $\mathrm H(T_2)$ for the statement that
$n_-(\Ainf|\PT)=1$ for every $T\in(T^*_{\infty,\mathrm P},T_2]$.

\begin{corollary}[one negative direction up to $T^*_{\mathrm{odd}}$]\label{cor:H1}
$\mathrm H(T_2)$ holds for every $T_2\in(T^*_{\infty,\mathrm P},T^*_{\mathrm{odd}})$ and for every
$T_2\in(T^*_{\infty,\mathrm P},\log2]$, and fails for every $T_2>T^*_{\mathrm{odd}}$; the case
$T_2=T^*_{\mathrm{odd}}>\log2$ is not decided. Hence the supremum of the admissible $T_2$ is
$T^*_{\mathrm{odd}}\in[\log2,0.81]$.
\end{corollary}
\begin{proof}
By Theorem~\ref{thm:index}(iii), $T^*_{\infty,\mathrm P}=T^*_{\mathrm{ev}}$, $n_-(\Ainf|\PT)=1$ on
$(T^*_{\mathrm{ev}},T^*_{\mathrm{odd}})$ and $n_-(\Ainf|\PT)=2$ on $(T^*_{\mathrm{odd}},0.90]$, which is
nonempty since $T^*_{\mathrm{odd}}\le0.81$. For $T_2<T^*_{\mathrm{odd}}$ the interval
$(T^*_{\mathrm{ev}},T_2]$ lies in the first range; for $T_2>T^*_{\mathrm{odd}}$ the interval
$(T^*_{\mathrm{ev}},T_2]$ meets the second (and for $T_2>0.90$ it contains $(T^*_{\mathrm{odd}},0.90]$).
For $T_2\le\log2$ use instead Theorem~\ref{thm:index}(i): $n_{\mathrm{odd}}=0$ on $(0,\log2]$, and
$n_{\mathrm{ev}}=1$ on $(T^*_{\mathrm{ev}},0.90]$ by Lemma~\ref{lem:strict} and Theorem~\ref{thm:index}(i).
\end{proof}

\begin{remark}
Theorem~\ref{thm:J} alone limits how far $\mathrm H(T_2)$ can reach:
$\lceil2T\tau_0/\pi\rceil-3\ge2$ as soon as $2T\tau_0/\pi>4$, so any admissible $T_2$ satisfies
$T_2\le2\pi/\tau_0=0.9989426\ldots$ Corollary~\ref{cor:H1} replaces that bound by
$\sup T_2=T^*_{\mathrm{odd}}\le0.81$: the second negative direction is odd and appears by $T=0.81$.
The window $\log2$ lies inside the admissible range, and $\frac12\log3$ lies below
$T^*_{\infty,\mathrm P}$, where $n_-=0$.
\end{remark}

\begin{remark}\label{rem:odd080}
Theorem~\ref{thm:index} leaves $n_{\mathrm{odd}}(0.80)\in\{0,1\}$, hence $n_-(\Ainf|\mathcal P_{0.80})\in\{1,2\}$;
whether $\mathrm H(0.80)$ holds, equivalently whether $T^*_{\mathrm{odd}}>0.80$, is open. The upper
certificate of Proposition~\ref{prop:haynsworth} at $T=0.80$ in the odd sector stops at $r=1$.
Numerically, Galerkin bottoms of the odd sector are $+0.0684$, $+0.0288$, $+0.0033$ at $T=0.75$, $0.78$,
$0.80$ and $-0.0092$ at $T=0.81$ (Galerkin, not certified); they suggest an odd crossing in $(0.80,0.81)$
and prove no sign at $0.80$.
\end{remark}

\begin{remark}
Theorem~\ref{thm:J}, Theorem~\ref{thm:arena} and Corollary~\ref{cor:mono} rest on one fact: $\varphi$ is
concave, so Jensen applied to $|\Fh|^2$ always produces an upper bound for the archimedean form
and never a lower bound. A lower bound would require control of the spread of $|\Fh|^2$, that is,
a quantitative statement that a function supported in $\IT$ cannot concentrate its spectrum in
$|\tau|<s$: an explicit non-asymptotic prolate concentration estimate. This is why the moment
arguments of this section bound $n_-$ from above by units and not by one.
\end{remark}

\subsection{Competing even modes}\label{subsec:arena}

No fixed finite number of archimedean modes captures the obstruction, and this is a theorem with an
explicit count. The archimedean gap along the even tower grows only logarithmically, while the prime
form is bounded by $\Aconst$ and does not decay; so the number $N^*(T)$ of even archimedean modes
whose level the prime form can reach, the competing even modes with which any truncation has to
work, is large, and grows doubly exponentially in the window.

\begin{theorem}[log-slow archimedean gaps]\label{thm:arena}
Let $\PT^{\mathrm{ev}}$ be the even part of $\PT$ and let
$\lambda^{\mathrm{ev}}_1\le\lambda^{\mathrm{ev}}_2\le\cdots$ be the eigenvalues of $\Ainf$ on it.
Then, for every $k\ge0$,
\[
\lambda^{\mathrm{ev}}_{k+1}\ \le\ C_0+\varphi\Big(\big(\tfrac{(2k+3)\pi}{2T}\big)^2\Big)
=m\Big(\frac{(2k+3)\pi}{2T}\Big)=\log k+O(1).
\]
\end{theorem}
\begin{proof}
$\Ainf$ commutes with parity, and on even functions the two Tate conditions coincide: both read
$\int F(t)\cosh\frac t2\,dt=0$. Hence $\PT^{\mathrm{ev}}$ has codimension $1$ in
$L^2_{\mathrm{ev}}(\IT)$, and the pole term vanishes on it, so $\Ainf=\mT$ there
(Proposition~\ref{prop:nodes}). Take the $k+2$ even Dirichlet modes $\cos((2j+1)\pi t/2T)$,
$j=0,\dots,k+1$, which vanish at $\pm T$ and satisfy
$\|F'\|^2\le((2k+3)\pi/2T)^2\|F\|^2$ on their span; intersecting that span with the single linear
constraint leaves a subspace of $\PT^{\mathrm{ev}}$ of dimension at least $k+1$. For $F$ in that subspace (extended by zero), Plancherel gives $\frac1{2\pi}\int\tau^2|\Fh|^2=\|F'\|^2$, and
$m-C_0=\varphi(\tau^2)$ with $\varphi$ complete Bernstein (Theorem~\ref{thm:S}(i)), hence concave and
increasing; Jensen's inequality in the variable $\tau^2$ for the probability measure
$|\Fh|^2d\tau/(2\pi\|F\|^2)$ gives
$\mT(F)/\|F\|^2\le C_0+\varphi(\|F'\|^2/\|F\|^2)\le C_0+\varphi(((2k+3)\pi/2T)^2)$, and
min--max bounds $\lambda^{\mathrm{ev}}_{k+1}$ by that number. The last equality is
$C_0+\varphi(\tau^2)=m(\tau)$ together with Lemma~\ref{lem:mbounds}(ii).
\end{proof}

\begin{corollary}[the number of competing even modes]\label{cor:arena}
Put $\tau_k=(2k+3)\pi/2T$ and let
$N^*(T)=\#\{j\ge1:\lambda^{\mathrm{ev}}_j\le\Aconst\}$ be the number of even archimedean modes
whose level the prime form, of norm at most $\Aconst$, can reach. Then
\[
N^*(T)\ \ge\ \#\{k\ge0:m(\tau_k)\le\Aconst\}\ =\ \Big\lfloor\frac{T\tau_*}{\pi}-\frac12\Big\rfloor,
\qquad \tau_*=m^{-1}(\Aconst).
\]
At $T=\log2$, where $\Aconst=2.24883$, Lemma~\ref{lem:mbounds}(ii) at the exact frequencies gives
$m(\tau_{11})\le2.2168\le\Aconst$ and $m(\tau_{12})\ge2.2596>\Aconst$, so $N^*\ge12$ and the
bound yields no more. At $T=\frac12\log3$, where $\Aconst=0.98026$, it gives
$m(\tau_1)\le0.8928\le\Aconst$ and $m(\tau_2)\ge1.1075>\Aconst$, so $N^*\ge2$. In general
$\Aconst=\sum_{n<e^{2T}}2\Lambda(n)n^{-1/2}=4e^{T}(1+o(1))$ by the prime number theorem and
partial summation, and $\tau_*=2\pi e^{\Aconst}(1+o(1))$ by Lemma~\ref{lem:mbounds}(ii), so
\[
N^*(T)\ \ge\ 2T\,e^{\Aconst}\,(1+o(1))\ =\ \exp\big(4e^{T}(1+o(1))\big):
\]
the number of competing even modes grows doubly exponentially in $T$.
\end{corollary}
\begin{proof}
By Theorem~\ref{thm:arena}, $\lambda^{\mathrm{ev}}_{k+1}\le m(\tau_k)$, so every $k$ with
$m(\tau_k)\le\Aconst$ contributes a mode. Since $m$ is increasing, $m(\tau_k)\le\Aconst$ is
$\tau_k\le\tau_*$, i.e.\ $k\le T\tau_*/\pi-\frac32$, and the number of integers $k\ge0$ with that
property is $\lfloor T\tau_*/\pi-\frac32\rfloor+1=\lfloor T\tau_*/\pi-\frac12\rfloor$. The two
window evaluations are Lemma~\ref{lem:mbounds}(ii) at $\tau_{11},\tau_{12}$ and at
$\tau_1,\tau_2$ respectively, with $\tau_k=(2k+3)\pi/2T$. The asymptotics are
$\log\Aconst$-free: $m(\tau)=\log(\tau/2\pi)+O(1/\tau)$ inverts to
$\tau_*=2\pi e^{\Aconst}(1+o(1))$, and $T\tau_*/\pi=2Te^{\Aconst}(1+o(1))$.
\end{proof}

Consequently no fixed finite truncation in the archimedean basis can decide the sign of the even
part: any argument that controls the prime form by the archimedean gap must carry at least $N^*(T)$
even modes, because up to that index the gap, growing like $\log k$, is smaller than the norm of a
perturbation that neither decays nor is compact.


\begin{remark}[no total positivity]\label{rem:PF2}
The property of $K$ that produces the upper bounds of this section also rules out the classical route to
a lower bound. By Lemma~\ref{lem:K}, $K$ is completely monotone and not a single exponential, hence
strictly log-convex on $(0,\infty)$ (Cauchy--Schwarz applied to the representing measure), and
$K(u)=u^{-1}+O(1)$ at the origin; so $u\mapsto K(|u|)$ is log-concave on no interval. By Schoenberg's
characterization of the P\'olya frequency functions of order $2$ (a nonnegative integrable $f$ is
$\mathrm{PF}_2$ if and only if it is log-concave \cite{Schoenberg1951}), $K(|\cdot|)$ is not
$\mathrm{PF}_2$, so the translation kernel $K(x-y)$ is not totally positive of order $2$, and Karlin's
variation-diminishing and oscillation theory \cite{Karlin1968} does not apply to it. This does not prove
that no such argument exists: the operator is $C_0+\Lev$, whose natural totally positive object would be
the transition density of the associated semigroup rather than $K$ itself, and we do not know whether
that density is $\mathrm{PF}_2$. Complete monotonicity thus plays two roles: it makes $\varphi$ a complete
Bernstein function, whence by Jensen the upper bounds of Corollary~\ref{cor:mono}, Theorem~\ref{thm:J}
and Theorem~\ref{thm:arena}, and it removes the total-positivity hypothesis under which nodal counting
would order the spectrum (compare Proposition~\ref{prop:nodes}).
\end{remark}

% (i.e. immediately before \section{Obstructions to general methods}).
\subsection{The negative set of the symbol}\label{subsec:band}

Fix a regime, so that $\sigma:=\sT$ and $\Aconst$ are fixed, and let $m_{\mathrm{lo}}$ be the lower bound of
Lemma~\ref{lem:mbounds}(ii); it increases from $-\infty$ to $+\infty$ on $(0,\infty)$.

\begin{proposition}[the negative set is a finite band]\label{prop:band}
Let $\tau^{\mathrm{lo}}$ be the unique positive root of $m_{\mathrm{lo}}(\tau)=\Aconst$. Then
\emph{(i)} $\sigma>0$ on $(\tau^{\mathrm{lo}},\infty)$;
\emph{(ii)} $\sigma(0)=C_0-\Aconst<0$;
\emph{(iii)} $\sigma$ is real-analytic and nonconstant on $\R$, so its zero set in $[0,\infty)$ is finite;
\emph{(iv)} $\{\sigma<0\}$ is a finite union of bounded open intervals, one of which contains $0$, and
$\tau_{\mathrm{last}}:=\sup\{\tau:\sigma(\tau)<0\}$ satisfies $\tau_{\mathrm{last}}\le\tau^{\mathrm{lo}}$;
\emph{(v)} $\tau^{\mathrm{lo}}=2\pi e^{\Aconst}\exp\big(1/\tau^{\mathrm{lo}}+1/(2(\tau^{\mathrm{lo}})^2)\big)>2\pi e^{\Aconst}$.
\end{proposition}
\begin{proof}
(i) For $\tau>\tau^{\mathrm{lo}}$, $m(\tau)\ge m_{\mathrm{lo}}(\tau)>\Aconst\ge\sum_{n\in\NT}v_n\cos(\tau\log n)$
(Lemma~\ref{lem:mbounds}(ii),(iii)). (ii) $m(0)=C_0<0$ by \eqref{eq:m}. (iii)
$\Rel\psi(\frac14+\frac{i\tau}2)=\frac12[\psi(\frac14+\frac{i\tau}2)+\psi(\frac14-\frac{i\tau}2)]$ is real-analytic in
$\tau\in\R$, the comb is a trigonometric polynomial, and $\sigma\to+\infty$; by (i) the zeros lie in the
compact interval $[0,\tau^{\mathrm{lo}}]$, and a nonconstant real-analytic function has finitely many zeros there.
(iv) follows from (i)--(iii) and the evenness of $\sigma$; (v) rearranges the equation
$m_{\mathrm{lo}}(\tau^{\mathrm{lo}})=\Aconst$.
\end{proof}

\begin{remark}\label{rem:bandvalues}
For $\NT=\{2,3\}$ ($\Aconst=2.248826$) one finds $\tau^{\mathrm{lo}}=60.5432$, and for $\NT=\{2,3,4\}$
($\Aconst=2.941974$) $\tau^{\mathrm{lo}}=120.0865$, against $2\pi e^{\Aconst}=59.543$ and $119.087$ (roots of the
explicit increasing function $m_{\mathrm{lo}}$, evaluated in double precision). So the threshold
$2\pi e^{\Aconst}$ of Proposition~\ref{prop:env}(iii) is the multiplier-free envelope zero up to the factor
$e^{1/\tau^{\mathrm{lo}}+\cdots}$.
\end{remark}

\begin{proposition}[the alignment structure at two and three contacts]\label{prop:align}
Let $\NT=\{2,3\}$, i.e.\ $T\in(\frac12\log3,\log2]$, so that $v_2=0.980258$ and $v_3=1.268568$, and let
$\tau_a=23.3409$ and $\tau_b=17.7446$ be the roots of $m_{\mathrm{lo}}=v_3$ and $m_{\mathrm{lo}}=v_2$. Every
$\tau>\tau_a$ with $\sigma(\tau)<0$ satisfies
\[
\cos(\tau\log2)>\frac{m_{\mathrm{lo}}(\tau)-v_3}{v_2}>0,\qquad\cos(\tau\log3)>\frac{m_{\mathrm{lo}}(\tau)-v_2}{v_3}>0,
\]
so that there are unique integers $k,l\ge1$ with $|\tau\log2-2\pi k|<\delta_2(\tau)$ and
$|\tau\log3-2\pi l|<\delta_3(\tau)$, where $\delta_2=\arccos\frac{m_{\mathrm{lo}}-v_3}{v_2}$ and
$\delta_3=\arccos\frac{m_{\mathrm{lo}}-v_2}{v_3}$ are $<\frac\pi2$ and decreasing; consequently
\[
|k\log3-l\log2|<\frac{\delta_2\log3+\delta_3\log2}{2\pi}<\frac{\log6}4=0.4479,\qquad\text{i.e.}\qquad
0.639<2^l/3^k<1.565 .
\]
Every component of $\{\sigma<0\}\cap(\tau_a,\tau^{\mathrm{lo}}]$ lies in one arc $\{|\tau\log2-2\pi k|<\delta_2\}$,
of length $<\pi/\log2$, and at most $6$ such arcs meet $(\tau_a,\tau^{\mathrm{lo}}]$. For $\NT=\{2,3,4\}$,
$T\in(\log2,\frac12\log5]$, the same holds for $\tau>\tau_a$ with $|\tau\log2-2\pi k|<\arccos y_0=1.1185$
($y_0=0.43702$ the positive root of $2v_4y^2+v_2y-v_4=0$), with the bracket constant
$(1.1185\log3+\frac\pi2\log2)/2\pi=0.369$ on $(\tau_b',\tau^{\mathrm{lo}}]$, where $\tau_b'=34.4909$ is the root of
$m_{\mathrm{lo}}=v_2+v_4$, and with at most $12$ arcs.
\end{proposition}
\begin{proof}
If $\sigma(\tau)<0$ then $v_2\cos(\tau\log2)>m(\tau)-v_3\cos(\tau\log3)\ge m_{\mathrm{lo}}(\tau)-v_3$, which is
positive for $\tau>\tau_a$; symmetrically for $3$, since $\tau_b<\tau_a$. The arcs of $\cos x>c>0$ are
$|x-2\pi j|<\arccos c<\frac\pi2$, and the integers are $\ge1$ because $\tau\log2>\frac\pi2$ there. Dividing the two
arc conditions by $\log2$ and $\log3$ and subtracting gives
$2\pi|k/\log2-l/\log3|<\delta_2/\log2+\delta_3/\log3$; multiply by $\log2\log3/2\pi$ and use $\delta_j<\frac\pi2$.
Consecutive arcs are $2\pi/\log2$ apart, whence the count $\lfloor(\tau^{\mathrm{lo}}-\tau_a)\log2/2\pi\rfloor+2=6$. For
$\{2,3,4\}$ the comb is $g(\tau\log2)+v_3\cos(\tau\log3)$ with $g(x)=v_2\cos x+v_4\cos2x=2v_4y^2+v_2y-v_4$,
$y=\cos x$, which is increasing for $y>-v_2/4v_4$ and positive on $[-1,1]$ exactly for $y>y_0$; for
$\tau>\tau_a$ a point with $\sigma(\tau)<0$ has $g>0$, and for $\tau>\tau_b'$ also
$v_3\cos(\tau\log3)>m_{\mathrm{lo}}-v_2-v_4>0$; the rest is as before, with the count $\lfloor10.67\rfloor+2=12$.
The numerical constants are values of explicit elementary functions, evaluated in double precision.
\end{proof}

\begin{remark}[locating the last sign change]\label{rem:taulast}
By the upper bound $m\le m_{\mathrm{up}}:=\log\frac\tau{2\pi}+\frac1\tau+\frac12\log(1+\frac1{4\tau^2})$ of
Lemma~\ref{lem:mbounds}(ii): at $\tau=45.64206$ the comb for $\{2,3\}$ equals $2.21553>m_{\mathrm{up}}=2.00492$,
and at $\tau=108.72265$ the comb for $\{2,3,4\}$ equals $2.93679>m_{\mathrm{up}}=2.86013$. So $\sigma<0$ at these
points and, with Proposition~\ref{prop:band}, $\tau_{\mathrm{last}}\in[45.642,60.544]$ for $\{2,3\}$ and
$\tau_{\mathrm{last}}\in[108.723,120.087]$ for $\{2,3,4\}$ (double-precision evaluation of explicit functions, with
margins $0.21$ and $0.077$). The measured components of $\{\sigma<0\}$ are listed in Remark~\ref{rem:bandtable}.
\end{remark}

Write $\sigma=\sigma_+-\sigma_-$ with $\sigma_\pm=\max(\pm\sigma,0)$. For a real even measurable function $h$ of at
most logarithmic growth, $h(D)$ denotes the Fourier multiplier $\widehat{h(D)F}=h\Fh$, and $Q_Th(D)Q_T$ the
corresponding form on $L^2(\IT)$, $\langle Q_Th(D)Q_TF,F\rangle=\frac1{2\pi}\int h|\Fh|^2$. Put
$\mathrm{Pos}_T:=Q_T\sigma_+(D)Q_T$, a nonnegative form on $\mathcal D_T$, and
$\mathrm{Neg}_T:=Q_T\sigma_-(D)Q_T$, a bounded operator on $L^2(\IT)$, so that $\mathfrak a_T=\mathrm{Pos}_T-\mathrm{Neg}_T$ on $\mathcal D_T$. For $s>0$ let $S_{Ts}$ be the time--band limiting
operator of Lemma~\ref{lem:KyFan} with band $|\tau|<s$.

\begin{proposition}[the negative part is trace class, and the band does not decide]\label{prop:negtrace}
\emph{(i)} $\mathrm{Neg}_T\ge0$ is trace class, with $\operatorname{tr}\mathrm{Neg}_T=\frac{2T}{2\pi}\int_\R\sigma_-(\tau)\,d\tau$.
\emph{(ii)} $\mathfrak a_T(F)\ge-\frac1{2\pi}\int_{|\tau|\le\tau_{\mathrm{last}}}\sigma_-|\Fh|^2$ for every $F\in\mathcal D_T$:
all negativity of the form is produced on the band $|\tau|\le\tau_{\mathrm{last}}$.
\emph{(iii)} $\lamA\ge-\|\mathrm{Neg}_T\|\ge-\operatorname{tr}\mathrm{Neg}_T$.
\emph{(iv)} $\lamA\ge0$ if and only if $\langle\mathrm{Neg}_TF,F\rangle\le\mathrm{Pos}_T(F)$ for all $F\in\PT\cap\mathcal D_T$.
\emph{(v)} The sign of the bottom on $\PT$ is not a function of $\sigma|_{[-\tau_{\mathrm{last}},\tau_{\mathrm{last}}]}$ alone. Let
$\lambda_0:=\|S_{T\tau_{\mathrm{last}}}\|$, $s_+:=\max_{|\tau|\le\tau_{\mathrm{last}}}\sigma_+$, and let $\mathbf 1_{\mathrm{off}}$ be the indicator
of $|\tau|>\tau_{\mathrm{last}}$. Then $\lambda_0<1$; the symbol $\sigma\mathbf 1_{|\tau|\le\tau_{\mathrm{last}}}-c\,\mathbf 1_{\mathrm{off}}$ gives a
form that is negative on $L^2(\IT)\setminus\{0\}$ for $c>s_+/(1-\lambda_0)$; and the symbol
$\sigma+c\,\mathbf 1_{\mathrm{off}}$ gives a form $\ge(c(1-\lambda_0)-\|\mathrm{Neg}_T\|)\|F\|^2>0$ on $\mathcal D_T$ for
$c>\|\mathrm{Neg}_T\|/(1-\lambda_0)$. Both symbols agree with $\sigma$ on the band.
\end{proposition}
\begin{proof}
(i) $\sigma_-$ is bounded and supported in $[-\tau_{\mathrm{last}},\tau_{\mathrm{last}}]$ (Proposition~\ref{prop:band}). So
$\mathrm{Neg}_T=(Q_T\sigma_-^{1/2}(D))(\sigma_-^{1/2}(D)Q_T)$ is the product of an operator and its adjoint, each
Hilbert--Schmidt: $\sigma_-^{1/2}(D)Q_T$ has kernel $\kappa(t-s)\mathbf 1_{\IT}(s)$ with
$\kappa=\frac1{2\pi}\int\sigma_-^{1/2}e^{-i\tau\cdot}$ in $L^2(\R)$. Its kernel on $\IT\times\IT$ is $k(t-s)$ with
$k=\frac1{2\pi}\int\sigma_-e^{-i\tau\cdot}$ continuous, and the trace of a positive trace-class operator with
continuous kernel is the integral of the diagonal, $\int_{\IT}k(0)\,dt=\frac{2T}{2\pi}\int\sigma_-$.
(ii)--(iv) are the identity $\mathfrak a_T=\mathrm{Pos}_T-\mathrm{Neg}_T$ and $\mathrm{Pos}_T\ge0$.
(v) $S:=S_{T\tau_{\mathrm{last}}}$ is compact with $0\le S\le1$, so $\lambda_0=\|S\|$ is an eigenvalue with an
eigenvector $F_0$; $\lambda_0=1$ would give $\int_{|\tau|>\tau_{\mathrm{last}}}|\Fh_0|^2=0$, impossible for the nonzero entire
function $\Fh_0$. Hence $\frac1{2\pi}\int_{\mathrm{off}}|\Fh|^2=\|F\|^2-\langle SF,F\rangle\ge(1-\lambda_0)\|F\|^2$. For the first
symbol the form is $\le s_+\langle SF,F\rangle-c(1-\lambda_0)\|F\|^2\le(s_+-c(1-\lambda_0))\|F\|^2$; for the second it is
$\mathrm{Pos}_T(F)-\langle\mathrm{Neg}_TF,F\rangle+\frac c{2\pi}\int_{\mathrm{off}}|\Fh|^2\ge(c(1-\lambda_0)-\|\mathrm{Neg}_T\|)\|F\|^2$. Both
signs hold on all of $L^2(\IT)$, hence on $\PT$. The statement concerns symbols, not $\zeta$.
\end{proof}

So the band $[0,\tau^{\mathrm{lo}}]$ carries all the negativity of $\mathfrak a_T$, it is finite and of explicit
length, and beyond $\tau_a\approx23.3$ its negative part sits at the near-solutions of $2^l\approx3^k$ in the two
regimes computed above. Positivity, however, is a global balance of $\mathrm{Pos}_T$ against $\mathrm{Neg}_T$, which the
band alone does not decide.


\section{Obstructions to general methods}\label{sec:negative}

\subsection{Mass against value}

\begin{theorem}[logarithmic exchange]\label{thm:Vp}
Let $B_*=\Aconst-C_0$ ($=7.621$ at $T\in(\frac12\log3,\log2]$).
\emph{(i)} For every feasible $\mu$ of (P), with $\varepsilon:=\mu(\{|\tau|\ge R\})$,
$\int\sT d\mu\ge-B_*+\varepsilon\,(m(R)-m(0))$, and $m(R)-m(0)=\log R+O(1)$.
\emph{(ii)} For every $\delta>0$ there is $C_1<\infty$ such that for all large $R$ and all
$\varepsilon\in[0,1]$ there is a feasible $\mu'_{\varepsilon,R}$ with
$\mu'_{\varepsilon,R}(\{|\tau|\ge R-\sqrt R\})\ge\varepsilon(1-o(1))$ and
$\int\sT d\mu'_{\varepsilon,R}\le\lamP+\delta+\varepsilon(\log R+C_1)+o(\varepsilon)$.
\emph{(iii)} Consequently no inequality $\int\sT d\mu\ge-B_*+\varepsilon\,\omega(R)$ with
$\omega(R)/\log R\to\infty$ holds on the feasible set: the exchange rate between escaped mass and
value is $\Theta(\log R)$, and converting the depth $B_*$ into a nonnegative value through escaped
mass alone requires $R\gtrsim e^{B_*/\varepsilon}$.
\end{theorem}
\begin{proof}
(i) $\sT\ge-B_*$ everywhere and $\sT\ge m(R)-\Aconst$ on $|\tau|\ge R$; split the integral.
(ii) Pick $F_0\in C_c^\infty(\IT)$ with $c_{F_0}\in\Kp$ and $\mathfrak a_T(F_0)\le\lamP+\delta$
(such $F_0$ are dense in the primitive class), $\mu_0=\mu_{F_0}$, which has a rapidly decaying
density. Set $\mu_{\varepsilon,R}=(1-\varepsilon)\mu_0+\frac\varepsilon2(\mu_0(\cdot-R)+\mu_0(\cdot+R))$;
then $\check\mu_{\varepsilon,R}=[(1-\varepsilon)+\varepsilon\cos Ru]\check\mu_0$ is still supported in
$[-2T,2T]$, so $\mu_{\varepsilon,R}$ is feasible for the relaxation without the two moment
constraints; its moment defect is $\varepsilon\,d(R)$ with $d(R)=(\int g_0\check\mu_0\cos Ru,\int g_1\check\mu_0\cos Ru)\to0$
by Riemann--Lebesgue. By Lemma~\ref{lem:W} there are $r_0>0$, $V^*<\infty$ and, for each $w$
with $|w|\le r_0$, a relaxed-feasible $\mu_w$ of value $\le V^*$ and moment $w$. Put
$\lambda=\varepsilon|d(R)|/(r_0+\varepsilon|d(R)|)$ and choose $\mu_w$ with
$w=-\frac{1-\lambda}\lambda\varepsilon d(R)$ (of norm $r_0$); then
$\mu'_{\varepsilon,R}=(1-\lambda)\mu_{\varepsilon,R}+\lambda\mu_w$ has moment $0$, is feasible for
(P), and its value exceeds that of $\mu_{\varepsilon,R}$ by at most
$\lambda(B_*+V^*)=o(\varepsilon)$. The value of $\mu_{\varepsilon,R}$ is
$(1-\varepsilon)\mathfrak a_T(F_0)+\varepsilon\int\sT(\tau\pm R)d\mu_0$, and by
Lemma~\ref{lem:mbounds}(ii) $\int m(\tau+R)d\mu_0\le\log R+\int|\tau|d\mu_0/R+O(1)$ while the prime
part is bounded by $\Aconst$. (iii) Apply (i) to $\mu'_{\varepsilon,R}$ and let $\varepsilon=1$.
\end{proof}

Theorem~\ref{thm:Vp} also limits the rate at which the lattice ceilings of Theorem~\ref{thm:R} converge: the
logarithmic rate of Corollary~\ref{cor:Rrate} can be improved on sub-families of feasible measures, but not
uniformly over the near-optimal ones.

\begin{corollary}[no uniform rate; a rate on a sub-family]\label{cor:CP}
Let $K_\beta$ be as in Theorem~\ref{thm:R}, $P=\pi M/T$, and let
$\mu\in\mathcal F_\varepsilon$ (\S\ref{subsec:conc}), so that Theorem~\ref{thm:R} reads
$\operatorname{Ceil}(M,T;\beta)\le\lamP+\varepsilon+K_\beta\,\mu(|\tau|>P/2)$.
\emph{(i)} Theorem~\ref{thm:P3}(i) gives
$\operatorname{Ceil}(M,T;\beta)\le\lamP+\varepsilon+K_\beta\,G(\varepsilon)/(m(P/2)-C_0)=\lamP+\varepsilon+O_{\beta,T}\big(G(\varepsilon)/\log M\big)$,
uniformly over $\mathcal F_\varepsilon$.
\emph{(ii)} On the sub-family of feasible measures satisfying $\int(m-C_0)^{1+\delta}d\mu\le K$ for
some $\delta>0$ (non-empty in $\mathcal F_\varepsilon$ for every $\varepsilon>0$),
\[
\operatorname{Ceil}(M,T;\beta)\ \le\ \lamP+\varepsilon+O_{\beta,T}\big(K\,(\log M)^{-1-\delta}\big).
\]
\emph{(iii)} The constant $K$ in (ii) cannot be taken uniform over $\mathcal F_\varepsilon$:
$\sup_{\mu\in\mathcal F_\varepsilon}\int(m-C_0)^{1+\delta}d\mu=\infty$ for every $\delta>0$.
\end{corollary}
\begin{proof}
(i) $\mu(|\tau|>P/2)\le G(\varepsilon)/(m(P/2)-C_0)$ by
Theorem~\ref{thm:P3}(i), and $m(P/2)-C_0=\log M+O_T(1)$ by Lemma~\ref{lem:mbounds}(ii). (ii) By Chebyshev's
inequality $\mu(|\tau|>P/2)\le K(m(P/2)-C_0)^{-1-\delta}$. The sub-family is non-empty because the
measures $\mu_{F_0}$, $F_0\in C_c^\infty(\IT)$, used in the proof of Theorem~\ref{thm:Vp}(ii) have
Schwartz densities and lie in $\mathcal F_\varepsilon$ for suitable $F_0$. (iii) Fix $\delta'\in(0,\varepsilon)$.
By Theorem~\ref{thm:Vp}(ii) with $\delta'$, for small $\eta>0$ and $\log R=(\varepsilon-\delta')/(2\eta)$ the
measure $\mu'_{\eta,R}$ lies in $\mathcal F_\varepsilon$ and has mass $\ge\eta/2$ on $|\tau|\ge R/2$, where
$m-C_0\ge\log R-O(1)$; so $\int(m-C_0)^{1+\delta}d\mu'_{\eta,R}\gg\eta^{-\delta}\to\infty$ as $\eta\to0$.
\end{proof}


\begin{remark}
Theorem~\ref{thm:Vp} excludes a rate lemma depending on $R$ alone, uniformly over the feasible
set. It does not exclude rate statements that depend on the data of a certificate class (the
lattice step $d$, the number of atoms, the $\ell^1$ size); compare \S\ref{sec:class}.
\end{remark}

\subsection{The envelope threshold}

\begin{proposition}[envelope threshold]\label{prop:env}
\emph{(i)} Every certificate $(\nu,\beta)$ satisfies $f(\tau)\ge\log\frac\tau{2\pi}-\frac1\tau-\frac1{2\tau^2}-\Aconst-\|\nu\|_{\TV}-|\beta_0|k_0(\tau)-|\beta_1|k_1(\tau)$,
so $f\ge0$ holds automatically for $\tau\ge\tau_{\mathrm{env}}:=2\pi e^{\Aconst+\|\nu\|_{\TV}+o(1)}$.
\emph{(ii)} No bound of the form $f\ge m(\tau)-E(\tau)$ with $E$ depending on $\tau$ only through
an upper bound of the prime comb can replace $\Aconst$ by a smaller constant, because
$\sup_\tau\sum_{n\in\NT}v_n\cos(\tau\log n)=\Aconst$ (Lemma~\ref{lem:mbounds}(iii)).
\emph{(iii)} $\Aconst=\sum_{n<e^{2T}}2\Lambda(n)/\sqrt n=(4+o(1))e^{T}$ by the prime number theorem,
so $\tau_{\mathrm{env}}(T)$ grows doubly exponentially in $T$. The multiplier-free part alone gives
the threshold $2\pi e^{\Aconst}$, equal to $16.7$ and $59.5$ at $T=\frac12\log3$ and $\log2$ (the last
sign change of $\sT$ is at $10.53$ and $46.13$); for a certificate of total variation $\|\nu\|_{\TV}$
the threshold is $2\pi e^{\Aconst+\|\nu\|_{\TV}}$, about $10^{40}$ for the certificate of
\S\ref{sec:cert}, which is why its tail is handled by the periodic Lemma~\ref{lem:F1} and not by this
bound.
\end{proposition}
\begin{proof}
(i) Lemma~\ref{lem:mbounds} and $|\nuh|\le\|\nu\|_{\TV}$. (ii),(iii) as stated.
\end{proof}

Thus a pointwise argument can stop at $\tau_{\mathrm{env}}(T)$ and at no earlier point, and
$2\pi e^{\Aconst}$, doubly exponential in $T$, is the shortest interval such an argument can
reach; this is the dual counterpart of \cite[Thm.~1.4]{Zhu2026}. We do not know whether every
certificate must do work on all of $[0,2\pi e^{\Aconst}]$; the periodic tail lemmas show that a
lattice certificate is checked on $[0,\tau_1+(N+1)P]$, with $\tau_1+(N+1)P\approx442$ at $T=\frac12\log3$.

\subsection{Partial prime combs}
\begin{theorem}[sub-comb negativity]\label{thm:SC}
Let $\mathcal N\subseteq\{n\ge2:\log n<2T\}$ be any set of contacts and $(v_n)_{n\in\mathcal N}$ any
nonnegative weights of total mass $A=\sum_{n\in\mathcal N}v_n$. Then the bottom of the corresponding
form on $\PT$ is at most $C_0+A+\varphi((3\pi/2T)^2)$. Since $\varphi$ is increasing with
$\varphi(0^+)=0$ and $T\mapsto\varphi((3\pi/2T)^2)$ decreases to $0$, if
\[
A<-C_0=5.3721834\ldots
\]
then the form is negative on $\PT$ for every $T>T_0(A)$, and nonpositive at $T=T_0(A)$, where $T_0(A)=\frac{3\pi}{2\sqrt{x_A}}$ and
$x_A$ is the unique solution of $\varphi(x_A)=-C_0-A$. For instance $T_0(0)=0.74921$,
$T_0(\frac12)=1.23295$, $T_0(1)=2.01848$, $T_0(2)=4.63567$, $T_0(4)=13.19461$, $T_0(5)=29.60905$.
\end{theorem}
\begin{proof}
The displayed bound is Corollary~\ref{cor:mono} with the weights $(v_n)_{n\in\mathcal N}$: the L\'evy
form $\Lev$ does not depend on the weights, its bottom on the codimension-two subspace $\PT$ is at
most $\varphi((3\pi/2T)^2)$, the comb contributes at most $A\|F\|^2$ and the constant is $C_0\|F\|^2$.
Monotonicity and the limits of $\varphi$ are immediate from the series of Theorem~\ref{thm:S}.
\end{proof}
\begin{remark}
The quantitative ancestors of Theorem~\ref{thm:SC} are the observations that the prime terms are
absent for $2T<\log2$ and that beyond that width they enter with a definite cost: Yoshida's prime-free
window \cite{Yoshida1992} and Bombieri's Theorem~12 \cite{Bombieri2000}. Theorem~\ref{thm:SC} adds
that no sub-collection of primes of total mass below $-C_0$ can carry positivity for large $T$.
\end{remark}


\begin{corollary}\label{cor:SC1}
Let the contacts be the ten primes $12049$, $124909$, $149333$, $322573$, $324011$, $373963$,
$417491$, $630229$, $782209$, $908861$, with their natural weights $v_p=2\log p/\sqrt p$ and all other
weights $0$. Then $A=0.563667\ldots$ and $T_0(A)=1.313456\ldots$. At
$T=\frac12\log908861+0.01=6.869974\ldots$, the smallest window in this family containing all ten
contacts, Theorem~\ref{thm:SC} gives
\[
\lambda_{\min}\big(A_T(v)\big|_{\PT}\big)\ \le\ -5.3721834+0.563667+2.685686=-2.122831 .
\]
\end{corollary}

Ten primes of size $10^5$--$10^6$, in a window that contains them, give a form that is negative by
$2.12$, about two fifths of the archimedean constant $|C_0|$. The complete comb at that window has
$\Aconst\sim4e^T\approx3.8\cdot10^3$, and its form is conjecturally nonnegative. Positivity, if it
holds, is thus a property of the complete initial segment $\{n<e^{2T}\}$: an argument that
establishes the inequality for a proper sub-comb and then adds contacts one at a time cannot work,
because every sub-comb of mass $<-C_0$ is negative in a large enough window. The mechanism is
explicit: each contact adds at most $v_n$ to the bound, the archimedean deficit $-C_0$ is a
constant, and the eigenvalue $\varphi((3\pi/2T)^2)$ of the window decays to $0$; only a comb whose
mass grows with $T$ keeps up. This concerns the operators $A_T(v)$, not $\zeta$. For comparison, the
window whose contact set is exactly the prime powers below $31$ (the sixteen numbers
$2,3,4,5,7,8,9,11,13,16,17,19,23,25,27,29$, among them the first ten primes) is
$T=\frac12\log31=1.71699$, where $\Aconst=16.640459$ exceeds $-C_0$ and Theorem~\ref{thm:SC} is
vacuous; the multiplier-free verification threshold there is
$2\pi e^{\Aconst}=1.06\cdot10^8$, against $16.7$ at $\frac12\log3$ and $59.5$ at $\log2$.



\subsection{A fixed window is a finite condition}\label{subsec:finite}
Fix $T$. For a nonnegative arithmetic function $a$ on $n\ge2$ write $v_n(a)=2a(n)/\sqrt n$ and
$\lambda_{\mathrm A}(T;a)$ for the bottom on $\PT$ of the form \eqref{eq:form} built with these weights
(as in Corollary~\ref{cor:SC1}). We say that \emph{Weil positivity holds at $T$ for $a$}, written
$\mathrm{WP}_T(a)$, if $\lambda_{\mathrm A}(T;a)\ge0$, and we write $\mathrm{WP}_T:=\mathrm{WP}_T(\Lambda)$, which is
$\lamA\ge0$. ($\mathrm{WP}_T$ is not the dual value $D_T$ of Definition~\ref{def:PD}.) For a real vector
$v=(v_n)_{n\in\NT}$, signed or not, $\mathfrak a^{(v)}_T$ denotes the form \eqref{eq:form} with the comb
$\sum_{n\in\NT}v_n\cos(\tau\log n)$. With $v(a):=(v_n(a))_{n\in\NT}$, the \emph{comb-level complete form} is
\[
W_T(a)(F):=\mathfrak a^{(v(a))}_T(F)+2\Rel\big(A(F)\overline{B(F)}\big),
\]
the form \eqref{eq:W} built with the weights of $a$, and $\lambda^{\mathrm W}_T(a)$ is its bottom on the unit sphere of
$L^2(\IT)$. So $W_T(\Lambda)=W_T$, and on $\PT$ the form $W_T(a)$ agrees with $\mathfrak a^{(v(a))}_T$. For
$a\ne\Lambda$, $W_T(a)$ is not the zero sum of $\zeta_a$: for the modifications of Theorem~\ref{thm:finmod}
below, $\zeta_a=\zeta E$ has exactly the zeros of $\zeta$, whose zero sum is $W_T$ (Lemma~\ref{lem:roughEF}),
while $W_T(a)(F)-W_T(F)=\sum_{n\in\NT}(v_n(\Lambda)-v_n(a))\cF(\log n)$, the contribution of $E'/E$, is in
general nonzero. Statements about $\lambda^{\mathrm W}_T(a)$ and $\mathrm{WP}_T(a)$ for $a\ne\Lambda$ are therefore
statements about forms, not about the zeros of $\zeta$ or of any $\zeta_a$.

\begin{proposition}[finite dependence]\label{prop:finite}
Put $v_T(a):=(v_n(a))_{n\in\NT}\in\R^{\NT}$, $|\NT|=\lceil e^{2T}\rceil-2$. Then:
\emph{(i)} $\mathrm{WP}_T(a)$ depends on $a$ only through $v_T(a)$;
\emph{(ii)} the set $\mathcal V_T:=\{v\in\R^{\NT}:\ \mathfrak a^{(v)}_T\ge0$ on $\PT\cap$ form domain$\}$ is closed
and convex, namely
\[
\mathcal V_T=\bigcap_{F}\Big\{v:\ \sum_{n\in\NT}v_n\,\cF(\log n)\le C_0+\Lev(\cF)\Big\},
\]
the intersection over the real unit vectors $F\in\PT$ in the form domain;
\emph{(iii)} $0\notin\mathcal V_T$ for $T>T_0(0)=0.74921$, indeed for $T\ge0.60$, and a vector $v\ge0$ with
$\sum_nv_n<-C_0-\varphi((3\pi/2T)^2)$ is not in $\mathcal V_T$ (for signed $v$ the same holds with
$\sum_n|v_n|$);
\emph{(iv)} for $a=\Lambda$ the numbers entering $\mathrm{WP}_T$ are exactly $\Lambda(n)$, $2\le n<e^{2T}$,
together with the one explicit function $m$ (equivalently $C_0$ and $K$), and nothing else.
\end{proposition}
\begin{proof}
(i) and (iv) are the definition of $\sT$. (ii) Real tests suffice: for $F=F_1+iF_2$ with $F_j$ real,
$|\Fh|^2=|\Fh_1|^2+|\Fh_2|^2+2\Iml(\Fh_1\overline{\Fh_2})$, and the last term is odd and integrates to $0$
against the even symbol, so the form splits as the sum of its values at $F_1,F_2$; and $A(F)=A(F_1)+iA(F_2)$
with $A(F_j)$ real, likewise for $B$, so $F\in\PT$ if and only if $F_1,F_2\in\PT$. For real unit $F$ in the
form domain, \eqref{eq:Aform} reads $\mathfrak a^{(v)}_T(F)=C_0+\Lev(\cF)-\sum v_n\cF(\log n)$ with
$\Lev(\cF)<\infty$, a closed half-space condition on $v$. (iii) The first statement is
Proposition~\ref{prop:arch}(ii) ($\laminfP<0$ for $T>T_0(0)$), and Corollary~\ref{cor:Tarch} for $T\ge0.60$; the second Theorem~\ref{thm:SC} with these
weights; for signed $v$ the comb term of that proof is bounded by $\sum|v_n|\,\|F\|^2$.
\end{proof}

\begin{remark}[regimes]\label{rem:regimes}
If $n_1<n_2<\cdots$ enumerate the integers $n\ge2$ with $\Lambda(n)\ne0$, then on each regime
$T\in(\frac12\log n_k,\frac12\log n_{k+1}]$ the vector $v_T(\Lambda)$, the mass $\Aconst$ and the symbol $\sT$
are constant, $\mathrm{WP}_T$ varies only through the interval $\IT$, and $\lamA$ is nonincreasing across all
regimes (Lemma~\ref{lem:monoT}). At a regime boundary one atom $v_{n_{k+1}}$ enters, and the sign of its
effect is not determined by anything proved here: $\cF(\log n)$ takes both signs on $\PT$
(Lemma~\ref{lem:cpos}). At a fixed window, $\mathrm{WP}_T$ is a closed condition on a known vector, and its
failure, if it occurs, is witnessed by a single test $F\in\PT$.
\end{remark}


\begin{lemma}[a contact of either sign]\label{lem:cpos}
Let $0<u_0<2T$. There are real $F_\pm\in\PT\cap C_c^\infty(\IT)$ with $\pm c_{F_\pm}(u_0)>0$.
\end{lemma}
\begin{proof}
Choose $0<\ell<\min(u_0,2T-u_0)$ and $J=(-T,-T+\ell)$: then $J$, $J+u_0$, $J+2u_0$ are pairwise disjoint
and $J\cup(J+u_0)\subset\IT$. For real $\phi\in C_c^\infty(J)$ put $F^\pm_\phi=\phi\pm\phi(\cdot-u_0)$. Then
$c_{F^\pm_\phi}(u_0)=\langle\phi\pm\phi_{u_0},\phi_{u_0}\pm\phi_{2u_0}\rangle=\pm\|\phi\|^2$, the three cross terms
vanishing by disjointness. The map $\phi\mapsto F^\pm_\phi$ is injective, so its image is
infinite-dimensional, and its intersection with $\PT$ has codimension at most $2$ in it and contains some
$F^\pm_\phi$ with $\phi\ne0$.
\end{proof}

\begin{corollary}[each local prime form is indefinite on $\PT$]\label{cor:Qp}
For a prime $p$ with $\log p<2T$ let $Q_p(F)=-\sum_{k\ge1,\,k\log p<2T}v_{p^k}\Rel\cF(k\log p)$ be the
$p$-part of the prime form $\PiT$. With $u_0=\log p$, $J_0:=(-T,-T+\ell)$ and
$F^\pm_\phi=\phi\pm\phi(\cdot-u_0)$ as in the proof of Lemma~\ref{lem:cpos}, one has
$Q_p(F^\pm_\phi)=\mp v_p\|\phi\|^2$ for every $\phi\in C_c^\infty(J_0)$. Hence, for every $T>\frac12\log p$,
$Q_p$ is negative definite on the infinite-dimensional subspace $\{F^+_\phi\}\cap\PT$ of
$\PT\cap C_c^\infty(\IT)$, and positive definite on $\{F^-_\phi\}\cap\PT$.
\end{corollary}
\begin{proof}
Since $\ell<u_0$, the translates $J_0+ku_0$, $k\ge0$, are pairwise disjoint. For $k\ge2$ the support
$J_0\cup(J_0+u_0)$ of $F^\pm_\phi$ and the support $(J_0+ku_0)\cup(J_0+(k+1)u_0)$ of $V_{ku_0}F^\pm_\phi$ are
therefore disjoint, and $c_{F^\pm_\phi}(k\log p)=0$. For $k=1$, $c_{F^\pm_\phi}(u_0)=\pm\|\phi\|^2$ by the proof of
Lemma~\ref{lem:cpos}, also for complex $\phi$. The subspaces are infinite-dimensional by the last sentence of
that proof.
\end{proof}

Thus no single $Q_p$ is nonnegative on $\PT$, nor is the archimedean form for $T\ge0.60$
(Corollary~\ref{cor:Tarch}); for finite sets of places see Theorem~\ref{thm:SC} and
Proposition~\ref{prop:finsupp}. Indefinite summands do not forbid a positive sum: on $\PT$ one has
$\mathfrak a_T=\mT+\PiT$ with both summands indefinite for $T\ge0.60$. The corollary only rules out
inferring the sign from a single local form.


\begin{theorem}[finite modifications]\label{thm:finmod}
Let $\psi_a(x):=\sum_{n\le x}a(n)$ (so $\psi_\Lambda$ is Chebyshev's function) and, for $\Rel s>1$,
$\zeta_a(s):=\exp\big(\sum_{n\ge2}a(n)n^{-s}/\log n\big)$.
\emph{(i) Deletion.} Let $N\ge5$ and $a:=\Lambda\mathbf 1_{n\ge N}$. For $T\le\frac12\log N$ the comb of $a$ is
empty, so $\mathrm{WP}_T(a)\iff\laminfP(T)\ge0$; hence $\mathrm{WP}_T(a)$ holds for every $T\le0.58$
and fails for every $T$ with $0.60\le T\le\frac12\log N$ (Theorem~\ref{thm:index}(i), Corollary~\ref{cor:Tarch}; for $T<\frac12\log2$
and for $T>T_0(0)$ already by Proposition~\ref{prop:arch}(iii) and (ii)); on $(0.58,0.60)$ it holds below
$T^*_{\infty,\mathrm P}$ and fails above it (Lemma~\ref{lem:strict}), the threshold not being located more
precisely here.
\emph{(ii) Amplification.} Let $n_0\ge2$, $c>0$ and $a:=\Lambda+c\,\mathbf 1_{n=n_0}$. For every $T$ with
$2T>\log n_0$ there is $c_*(T,n_0)<\infty$ such that $\mathrm{WP}_T(a)$ fails for all $c>c_*$. For $n_0=2$
this holds at every $T>\frac12\log2$, including $(\frac12\log2,\frac12\log5]$, where $\mathrm{WP}_T(\Lambda)$ holds
(Theorem~\ref{thm:log5A}).
\emph{(iii) What is preserved.} Both modifications keep $a\ge0$ and change $\psi_a$ by a constant for
$x\ge N$ (resp.\ $x\ge n_0$), so they preserve every partial-sum statement $\psi_a(x)=x+O(\varepsilon(x))$
with $\inf\varepsilon>0$; and $\zeta_a=\zeta\cdot E$ with
$E(s)=\exp\big(-\sum_{n<N}\Lambda(n)n^{-s}/\log n\big)$ (resp.\ $\exp(c\,n_0^{-s}/\log n_0)$) entire and zero-free,
so $\zeta_a$ continues meromorphically with \emph{exactly the zeros and the pole of $\zeta$}. Deletion also
preserves $a\le\Lambda$. Neither preserves the functional equation, $E(s)\ne E(1-s)$.
\end{theorem}
\begin{proof}
(i) $n\in\NT$ means $n<e^{2T}\le N$, so $a=0$ on $\NT$ and the form is the archimedean form on $\PT$, whose
bottom is $\laminfP(T)$ (\S\ref{subsec:arch}); Proposition~\ref{prop:arch}(iii) gives
$\laminfP\ge\laminf\ge0$ for $T<\frac12\log2$, and Proposition~\ref{prop:arch}(ii) gives
$\laminfP(T)\le C_0+\varphi((3\pi/2T)^2)<0$ for $T>T_0(0)$; Theorem~\ref{thm:index}(i) and Corollary~\ref{cor:Tarch} give $\laminfP>0$ on
$(0,0.58]$ and $\laminfP<0$ on $[0.60,\infty)$; $N\ge5$ makes $\frac12\log N\ge0.8047>T_0(0)$.
(ii) Take $F=F_+$ of Lemma~\ref{lem:cpos} with $u_0=\log n_0$. Being smooth, it lies in the form domain,
and $\mathfrak a^{(a)}_T(F)=\mathfrak a_T(F)-(2c/\sqrt{n_0})\,\cF(\log n_0)$, which is negative for
$c>c_*:=\sqrt{n_0}\max(\mathfrak a_T(F),0)/(2\cF(\log n_0))$. (iii) The exponent of $\zeta_a$ differs from that
of $\zeta$ by a finite Dirichlet polynomial, an entire function.
\end{proof}

\begin{remark}[resolution in the weights]\label{rem:resweights}
This remark quantifies Theorem~\ref{thm:finmod}(ii) by Corollary~\ref{cor:margin1}, first for the
primitive form. Fix $T$ with $\lamA>0$ and $n_0\in\NT$. A real unit minimizer $F_\star$ of $\mathfrak a_T$ on
$\PT$ exists: $\lamA$ is attained (Lemma~\ref{lem:compact} and the paragraph preceding
Lemma~\ref{lem:contA}), and the splitting in the proof of Proposition~\ref{prop:finite}(ii) gives a real
minimizer. Suppose $c_{F_\star}(\log n_0)\ne0$, put $s=\operatorname{sgn}c_{F_\star}(\log n_0)$ and, for
$\varepsilon>0$, $a=\Lambda+s\varepsilon\mathbf 1_{n=n_0}$; if $s=-1$ this requires $\varepsilon\le\Lambda(n_0)$, so
that $a\ge0$. In particular a contact with $\Lambda(n_0)=0$ admits only $s=+1$, and the class
$\mathcal C_{\mathrm{BT}}$ of Corollary~\ref{cor:void} admits only $s=-1$. The computation in the proof of
Theorem~\ref{thm:finmod}(ii), applied to $F_\star$ in place of $F_+$, gives
\[
\lambda_{\mathrm A}(T;a)\ \le\ \mathfrak a^{(v(a))}_T(F_\star)\ =\ \lamA-\frac{2\varepsilon}{\sqrt{n_0}}\,|c_{F_\star}(\log n_0)| .
\]
So $\mathrm{WP}_T(a)$ fails for every admissible $\varepsilon>\varepsilon_*(T,n_0):=\sqrt{n_0}\,\lamA/(2|c_{F_\star}(\log n_0)|)$.
For $T\ge1$, Corollary~\ref{cor:margin1} and $\sqrt{n_0}<e^T$ give
\[
\varepsilon_*(T,n_0)\ \le\ \frac{2\cdot10^9\,T\,e^{-14.449\,T}}{|c_{F_\star}(\log n_0)|},\qquad
\varepsilon_*(T,2)\ \le\ \frac{2\sqrt2\cdot10^9\,T\,e^{-15.449\,T}}{|c_{F_\star}(\log2)|}.
\]
For the complete form, take a real unit minimizer $F^{\mathrm W}$ of $W_T$ on $L^2(\IT)$ (attained, as in the
proof of Lemma~\ref{lem:Onat}) and let $s$ be the sign of $c_{F^{\mathrm W}}(\log n_0)$. The same computation gives
$\lambda^{\mathrm W}_T(a)\le\lamW-(2\varepsilon/\sqrt{n_0})|c_{F^{\mathrm W}}(\log n_0)|$, and the same bounds hold with
$c_{F^{\mathrm W}}$ in place of $c_{F_\star}$, since $\lamW\le\lamA$.

These are upper bounds on the threshold, not its size. Since $|c_{F_\star}|\le\|F_\star\|^2=1$, the factor
$1/|c_{F_\star}(\log n_0)|$ only enlarges the bound, and no lower bound on $\max_{n\in\NT}|c_{F_\star}(\log n)|$ is
known. Even when $|c_{F_\star}(\log n_0)|=1$ the first bound exceeds $1$ for $T\le1.51$, so the remark is vacuous
there. For $\mathcal C_{\mathrm{PNT}}(\varepsilon)$ and $\mathcal C_{\mathrm{zeros}}$, Corollary~\ref{cor:void} already
applies with no size restriction; the remark adds content only for classes closed under a change of one
weight by $\varepsilon_*(T,n_0)$, for instance inputs known only to within a given width. We do not know whether
$c_{F_\star}(\log n_0)\ne0$ for some minimizer and some contact; if this fails, every minimizer has vanishing
autocorrelation on $\log\NT$, and Lemma~\ref{lem:degen} describes that case.
\end{remark}

\begin{lemma}[the degenerate case]\label{lem:degen}
Fix $T$. For $v\in\R^{\NT}$ let
$\mu_T(v):=\inf\{\mathfrak a^{(v)}_T(F):\ F\in\PT\cap\mathcal D_T,\ \|F\|=1\}$, so that
$\mu_T(v_T(\Lambda))=\lamA$. Then $\mu_T$ is finite and concave on $\R^{\NT}$. If $F_\star$ is a real unit
minimizer at $v^0=v_T(\Lambda)$, then
$\mu_T(v)\le\mu_T(v^0)-\sum_{n\in\NT}(v_n-v^0_n)\,c_{F_\star}(\log n)$ for all $v$. Hence, if some real unit
minimizer has $c_{F_\star}(\log n)=0$ for every $n\in\NT$, then $v_T(\Lambda)$ is a global maximizer of
$\mu_T$ on $\R^{\NT}$; in particular $\lamA\ge\mu_T(0)=\laminfP(T)$.
\end{lemma}
\begin{proof}
By \eqref{eq:Aform} and the splitting in the proof of Proposition~\ref{prop:finite}(ii), $\mu_T$ is the
infimum, over real unit $F\in\PT\cap\mathcal D_T$, of the affine functions
$v\mapsto C_0+\Lev(\cF)-\sum_nv_n\cF(\log n)$; so $\mu_T$ is concave. It is bounded below by
$C_0-\sum_n|v_n|$, because $\Lev\ge0$ and $|\cF|\le1$, and above by its value at any fixed $F$. The displayed
inequality is $\mu_T(v)\le\mathfrak a^{(v)}_T(F_\star)$, written out with \eqref{eq:Aform}. The last assertion
is the case $v=0$, where the comb is empty and the form is $\mT$ on $\PT$.
\end{proof}

We do not know whether the degenerate case occurs at any $T$. If it did, the comb of $\zeta$ would be the
optimal comb for the primitive bottom at that window among all real weight vectors.


\begin{proposition}[finitely supported combs]\label{prop:finsupp}
Let $a\ge0$ have finite support $S$ and put $A_S=\sum_{n\in S}v_n(a)$. Then
$\limsup_{T\to\infty}\lambda_{\mathrm A}(T;a)\le C_0-A_S<0$; in particular $\mathrm{WP}_T(a)$ fails for all
large $T$. This applies to every truncation $\Lambda\mathbf 1_{n<N}$, and for the empty comb it recovers
a bound consistent with $\laminfP\downarrow C_0$ (Proposition~\ref{prop:arch}(ii)).
\end{proposition}
\begin{proof}
For real $g\in C_c^\infty(-1,1)$, $g\ne0$, put $G_T(t)=g(t/T)$ and $F_T=(-\partial^2+\frac14)G_T\in C_c^\infty(\IT)$.
Two integrations by parts give $A(F_T)=\int G_T(-\partial^2+\frac14)e^{t/2}dt=0$ and likewise $B(F_T)=0$, so
$F_T\in\PT$, and $\Fh_T(\tau)=(\tau^2+\frac14)T\widehat g(T\tau)$. Let $2T>\max_S\log n$, so that
$\sigma:=m-\sum_Sv_n(a)\cos(\tau\log n)$ does not depend on $T$. Substituting $s=T\tau$,
\[
\frac{\mathfrak a^{(a)}_T(F_T)}{\|F_T\|^2}=\frac{\int\sigma(s/T)(s^2/T^2+\frac14)^2|\widehat g(s)|^2ds}{\int(s^2/T^2+\frac14)^2|\widehat g(s)|^2ds}
\ \longrightarrow\ \sigma(0)=C_0-A_S\qquad(T\to\infty)
\]
by dominated convergence: for $T\ge1$, $|\sigma(s/T)|\le C(1+\log(2+|s|))$ (Lemma~\ref{lem:mbounds} and
$|\text{comb}|\le A_S$), $(s^2/T^2+\frac14)^2\le(s^2+\frac14)^2$, $\widehat g$ is Schwartz, and $\sigma$ is
continuous at $0$.
\end{proof}

\begin{corollary}[classical prime-side hypotheses decide no window]\label{cor:void}
Let $\mathcal C$ be a class of nonnegative arithmetic functions containing $\Lambda$. If $\mathcal C$ is
closed under deletion of finitely many terms, then for every $T>T_0(0)$ some $a\in\mathcal C$ fails
$\mathrm{WP}_T(a)$, namely $a=\Lambda\mathbf 1_{n\ge N}$ with $N\ge\max(5,e^{2T})$. If $\mathcal C$ is closed under
amplification of one term, then for every $T>\frac12\log2$ some $a\in\mathcal C$ fails $\mathrm{WP}_T(a)$, namely
$a=\Lambda+c\,\mathbf 1_{n=2}$ with $c>c_*(T,2)$. The following classes contain $\Lambda$ and have these closure
properties:
\begin{itemize}
\item the class $\mathcal C_{\mathrm{PNT}}(\varepsilon)=\{a\ge0:\psi_a(x)=x+O(\varepsilon(x))\}$, for any error term
$\varepsilon$ with $\inf\varepsilon>0$ that is known for $\Lambda$ (e.g.\ de la Vall\'ee Poussin's), is closed under
both operations;
\item the class $\mathcal C_{\mathrm{zeros}}=\{a\ge0:\ \zeta_a=\zeta E,\ E$ entire and zero-free$\}$ of all $a$ with
the zeros and the pole of $\zeta$, and hence with every zero-free region and every zero-density estimate
of $\zeta$, is closed under both operations;
\item the class $\mathcal C_{\mathrm{BT}}=\{a:0\le a\le\Lambda\}$, which satisfies every upper-bound sieve
inequality satisfied by $\Lambda$, is closed under deletion.
\end{itemize}
Consequently no argument that uses about the primes only membership of $\Lambda$ in such a class (for
instance a prime number theorem with error term, a zero-free region, a zero-density estimate, or an
upper-bound sieve inequality) can prove $\mathrm{WP}_T$ at a window $T>T_0(0)$, since it would equally prove
$\mathrm{WP}_T(a)$ for the member $a$ above, which is false; for $\mathcal C_{\mathrm{PNT}}(\varepsilon)$ and
$\mathcal C_{\mathrm{zeros}}$ the same holds at every window $T>\frac12\log2$.
\end{corollary}
\begin{proof}
Theorem~\ref{thm:finmod}(i),(ii) give the failing members; the closure properties are
Theorem~\ref{thm:finmod}(iii).
\end{proof}

\begin{remark}[the functional equation]\label{rem:twohalves}
A similar obstruction comes from the other side of the explicit formula. The functional equation and the
pole are shared by every Dedekind zeta function, and Epstein zeta functions of class number $>1$, which
share the functional equation, reality, the polar set and nonnegative coefficients, have zeros off the
critical line \cite{DavenportHeilbronn1936}, so that, by Weil's criterion \cite{Weil52} applied to these
functions, their Weil forms take negative values. Hence a proof of $\mathrm{WP}_T$ for $\zeta$ must use a
property of $\zeta$ that (a) is not preserved by the modified weights of Theorem~\ref{thm:finmod}, which fail
$\mathrm{WP}_T$ at the windows listed in Corollary~\ref{cor:void}, so that it depends on the actual values
$\Lambda(n)$, $n\in\NT$, and not only on their asymptotics; and (b) is not shared by all Epstein zeta
functions, so that it does not follow from the functional equation, the pole, reality and nonnegative
coefficients alone. The known proofs at individual windows are computations with the actual vector:
Theorem~\ref{thm:C} at $\frac12\log3$, \cite[Thm.~1.2]{Zhu2026} for $T\le0.8$ (the complete form), and
Theorems~\ref{thm:log5A} and \ref{thm:log5W} for $T\le\frac12\log5$, and \cite[Thm.~B]{Liu2026} for
$T\le17/16$ and \cite{Mori15} for $T\le\frac12\log15$ (the complete form; unrefereed, cited).
\end{remark}



\subsection{The sign at a fixed window is decidable}

\begin{proposition}[the sign at each window is decidable]\label{prop:decide}
Fix a computable $T>0$ with $\lamP\ne0$. Then the sign of $\lamP$ is decided by a procedure that
terminates. In the positive case it produces a finite certificate, hence a proof of $\lamA>0$; in the
negative case it produces a feasible point of \textup{(P)} of negative value, which proves $\lamP<0$ but, since
$\lamP\le\lamA$, not the sign of $\lamA$. By Theorem~\ref{thm:O} the procedure in the positive case
never certifies the value $\lamP$ itself.
\end{proposition}
\begin{proof}
All constants below ($C_0$, $\Aconst$, values of $m$, of $\gh0,\gh1$, of $k_0,k_1$, of the transforms
$\nuh_{s,p}$ of Proposition~\ref{prop:T0}, and the Tate moments of step functions) are computable reals,
enclosed to any prescribed accuracy in interval arithmetic. Run two searches in alternation.

(a) Enumerate finite rational data: rational $s_k>0$, integers $p_k\ge0$, rational $c_k$, rational
$\beta\in\Q^2$, a rational target $c>0$, and rational $\tau_1>0$. For each item, put
$\nu=\sum_kc_k\nu_{s_k,p_k}$, $f=\sT+\nuh+\beta\cdot\widehat g$, check $L(\tau_1)>c$ for the function $L$
of Lemma~\ref{lem:T} (which then gives $f>c$ on $[\tau_1,\infty)$), and check $f>c$ on $[0,\tau_1]$ by
subdividing it into cells of a fixed width and enclosing $f$ on each cell by its value at the centre
and a Lipschitz bound ($f'$ is bounded on $[0,\tau_1]$ by an explicit constant); halve the width up to a
bound depending on the item, and pass to the next item if the check fails. If $\lamP>0$, Theorem~\ref{thm:E}
gives data $(\nu^\circ,\beta^\circ)$ in the span of the exponential columns with
$\inf f^\circ>\frac23\lamP$. The map from the data to $f$ is continuous for the uniform norm on $\R$:
$|\nuh_{s,p}-\nuh_{s',p}|\le\|\nu_{s,p}-\nu_{s',p}\|_{\TV}\to0$ as $s'\to s$, and
$|\nuh|\le\|\nu\|_{\TV}$, $|\gh j|\le\gh j(0)$. So some rational item has $\inf f>\frac12\lamP$, and for
it $c=\frac13\lamP$ (or any rational below $\frac12\lamP$) passes: $L(\tau_1)>c$ for large $\tau_1$ by
Corollary~\ref{cor:T2}, and the cell check succeeds once the width is below
$(\frac12\lamP-c)/\sup_{[0,\tau_1]}|f'|$ and the enclosures are accurate enough. A passed check is a
certificate of value $c>0$, so $\lamA\ge\lamP\ge c>0$ by Proposition~\ref{prop:weak}.

(b) Feasible points of \textup{(P)} are the even probability measures $\mu$ with $\check\mu=0$ off the window and
$\mathfrak m(\mu)=0$ in the notation of Lemma~\ref{lem:W}; for real $F$, $\mu_F=|\Fh|^2d\tau/(2\pi\|F\|^2)$ has
$\mathfrak m(\mu_F)=(Q_0(F),Q_1(F))/\|F\|^2$, and finite convex combinations of such measures satisfy the window
condition. Fix four real step functions $F_1,\dots,F_4$ whose normalized moment vectors have $0$ in the interior of
their convex hull; they exist by Lemma~\ref{lem:W}, since $F\mapsto(A,B,A',B')$ is onto $\R^4$ already from step
functions. Enumerate pairs of real step functions $G_1,G_2$ with rational values on rational partitions of $\IT$ and
a rational $\theta\in[0,1]$, put $\mu_G=\theta\mu_{G_1}+(1-\theta)\mu_{G_2}$ with moment $d=\mathfrak m(\mu_G)$, and,
choose computable weights $\lambda\in[0,1)$ and convex $\lambda_1,\dots,\lambda_4$ with
$(1-\lambda)d+\lambda\sum_i\lambda_i\mathfrak m(\mu_{F_i})=0$; they exist because the hull contains a disc around
$0$, and $\lambda$ can be taken $\to0$ as $d\to0$ (as in the proof of Theorem~\ref{thm:Vp}(ii)). The measure $\mu=(1-\lambda)\mu_G+\lambda\sum_i\lambda_i\mu_{F_i}$
is feasible for \textup{(P)}, and its value $\int\sT d\mu$ is a computable real: for a step function $G$, $c_G$ is an
explicit piecewise linear function, $\int\sT d\mu_G=(C_0\|G\|^2+\Lev(c_G)-\sum_{n\in\NT}v_nc_G(\log n))/\|G\|^2$
(Theorem~\ref{thm:A}), and $\Lev(c_G)$ is the integral of the completely monotone density $K$ against $c_G(0)-c_G(u)$,
which is $O(u)$ at $0$ and constant beyond $2T$. Check $\int\sT d\mu<0$ by interval enclosure. If $\lamP<0$, the
primal is attained at some $\Phi_\star=c_{F_\star}\in\Kp$ with $\|F_\star\|=1$ (proof of Theorem~\ref{thm:O}), and
writing $F_\star=u+iv$ with $u,v$ real, $\mu_{F_\star}=\|u\|^2\mu_u+\|v\|^2\mu_v$ is of the enumerated form with
$\theta=\|u\|^2$ (the form splits over real and imaginary parts, as in the proof of Proposition~\ref{prop:finite}(ii)).
Step functions $G_j\to u$ and $G'_j\to v$ with supports in a fixed compact subset of $\IT$ and convergence in
$H^s$ for every $s<\frac12$ give $d_j\to0$, hence $\lambda_j\to0$, and the values converge, since
$\log(2+|\tau|)\le C_s(1+\tau^2)^s$ makes the form continuous for the $H^s$ norm; so some enumerated $\mu$ has
negative value, and a strict inequality is detected in finite time.

Exactly one of $\lamP>0$, $\lamP<0$ holds, and the corresponding search terminates; the other cannot
terminate, by weak duality and by the definition of $\lamP$.
\end{proof}

\subsection{Bounded remainders}
\begin{theorem}[infinite negative index]\label{thm:inertia-nogo}
Let $T>\frac12\log2$ and let $k\in L^\infty(-2T,2T)$ be real and even, with $C_k$ the operator with kernel
$k(x-y)$ on $L^2(\IT)$. For every $N\ge1$ there is an $N$-dimensional subspace of
$C_c^\infty(\IT)\cap\PT$ on which $-\langle C_kF,F\rangle+\PiT(F)<0$ for $F\ne0$. Hence
$n_-\big(-C_k+\PiT\,\big|\,\PT\cap\{\int F=0\}\big)=\infty$. Consequently, if $\mathfrak a_T=\mathfrak p+\mathfrak r$ with
$\mathfrak p\ge0$ and $\mathfrak r=-C_k+\PiT$ of this form, then $\mathfrak r$ is not bounded below by any form
with finitely many negative directions on $\PT$. By \cite[Thm.~4.7]{ConnesConsani2021} this applies to the
Sonin trace $\mathfrak p$ of Connes and Consani, whose remainder kernel is continuous.
\end{theorem}
\begin{proof}
Fix $c_0\in C_c^\infty(-1,1)$ real, $c_w=c_0(\cdot/w)$, $\beta=c_w''-\frac14c_w$. Two
integrations by parts give $\int\beta e^{\pm t/2}=0$, so translates of $\beta$ lie in $\PT$. Put
$h_j=\beta(\cdot-t_j)+\beta(\cdot-t_j-a)$, $a=\log2$, with $t_j$ chosen so that all mutual differences avoid
$\{\pm\log n\pm\{0,a\}:n\in\NT\}$ by more than $4w$. For $x=\sum c_jh_j$ the only nonzero value of $c_x(\log n)$,
$n\in\NT$, is the inner one at $\log2$, so $\PiT(x)=-v_2\|\beta\|_2^2\sum|c_j|^2=-\sqrt2\log2\,\|\beta\|_2^2\sum|c_j|^2$.
Also $|\langle C_kx,x\rangle|\le\|k\|_\infty\|x\|_1^2\le4N\|k\|_\infty\|\beta\|_1^2\sum|c_j|^2$ and
$\|\beta\|_1^2/\|\beta\|_2^2=O(w)$. For small $w$ the form is $\le-\frac12\sqrt2\log2\,\|\beta\|_2^2\sum|c_j|^2$.
Imposing $\int x=0$ costs one dimension. A bound with $k'$ negative directions fails on a $(k'+1)$-dimensional
such subspace by dimension count.
\end{proof}

\subsection{The margin decays}\label{subsec:margin}
The trivial half of Zhu's conjectural law $-\log\lamW\asymp e^{2T}$ \cite[Conj.~12.1]{Zhu2026}, namely that
the margin tends to $0$ at an explicit rate, can be proved from three facts about the zeros of $\zeta$.
They are the only numerical facts about zeros used in a proof in this paper, and nothing outside this
subsection depends on them.
\begin{itemize}
\item[(Z1)] $\rho_1=\frac12+i\gamma_1$, $\gamma_1=14.134725141734\ldots$, is a simple zero on the critical line,
and $\zeta$ has no non-trivial zero $\rho\ne\rho_1,\bar\rho_1$ with $|\Iml\rho|\le21$
\cite{RosserYoheSchoenfeld1968,PlattTrudgian2021}. Only $\gamma_1\in[14.13,14.14]$ and the gap $(\gamma_1,21]$
are used.
\item[(Z2)] Every non-trivial zero $\rho$ with $|\Iml\rho|\le H_{\mathrm{PT}}:=3\cdot10^{12}$ lies on the
critical line \cite{PlattTrudgian2021}.
\item[(Z3)] For $t\ge2$, $\big|N(t)-\frac t{2\pi}\log\frac t{2\pi e}-\frac78\big|\le0.137\log t+0.443\log\log t+1.588$,
where $N(t)$ counts the zeros with $0<\Iml\rho\le t$ with multiplicity \cite[Thm.~19]{Rosser1941}.
\end{itemize}
Besides these, we use that the non-trivial zeros lie in $0<\Rel s<1$, and Lemma~\ref{lem:roughEF}.
\begin{lemma}[a crude zero count]\label{lem:Ncount}
$N(t)\le\frac t{2\pi}\log t$ for $t\ge21$.
\end{lemma}
\begin{proof}
$\frac t{2\pi}\log\frac t{2\pi e}=\frac t{2\pi}\log t-\frac{1+\log2\pi}{2\pi}t$ and
$\frac{1+\log2\pi}{2\pi}=0.45166\ldots$, so by (Z3)
$\frac t{2\pi}\log t-N(t)\ge0.45166\,t-\frac78-1.588-0.137\log t-0.443\log\log t$, which is $6.11\ldots$ at
$t=21$ and has derivative $0.45166-0.137/t-0.443/(t\log t)>0$ for $t\ge21$.
\end{proof}

\begin{lemma}[zero sums by counting]\label{lem:zerosum}
For $p\ge2$ and $H\ge21$, with the sum over the ordinates $\gamma=\Iml\rho>0$ of the non-trivial zeros
counted with multiplicity,
\[
\Sigma_p(H):=\sum_{\gamma>H}\gamma^{-p}\ \le\ \frac p{2\pi}\,H^{1-p}\Big[\frac{\log H}{p-1}+\frac1{(p-1)^2}\Big].
\]
\end{lemma}
\begin{proof}
$\sum_{\gamma>H}\gamma^{-p}=\int_{H^+}^\infty t^{-p}dN(t)=\big[t^{-p}N(t)\big]_H^\infty+p\int_H^\infty t^{-p-1}N(t)\,dt$;
the term at $\infty$ vanishes and the one at $H$ is $-H^{-p}N(H)\le0$ by Lemma~\ref{lem:Ncount}, which
also bounds the integral by $\frac p{2\pi}\int_H^\infty t^{-p}\log t\,dt=\frac p{2\pi}H^{1-p}\big[\frac{\log H}{p-1}+\frac1{(p-1)^2}\big]$.
\end{proof}

\begin{lemma}[the test family]\label{lem:testfam}
Fix $T>0$ and an integer $k\ge5$; let $b_k=\frac k2\mathbf 1_{[-1/k,1/k]}$ and $\omega_k=b_k^{*k}$ (the
$k$-fold convolution).
\emph{(i)} $\omega_k\ge0$, $\int\omega_k=1$, $\supp\omega_k\subseteq[-1,1]$, $\omega_k\in C^{k-2}(\R)$ with
$\omega_k^{(j)}(\pm1)=0$ for $j\le k-2$, and $\|\omega_k\|_2^2\ge\frac12$.
\emph{(ii)} $\widehat\omega_k(x)=\big(\sin(x/k)/(x/k)\big)^k$; $|\widehat\omega_k(x)|\le\min\{1,(k/|x|)^k\}$ for real
$x$, and $|\widehat\omega_k(z)|\le(k/|\Rel z|)^ke^{|\Iml z|}$ for complex $z$.
\emph{(iii)} Put $G(t)=\omega_k(t/T)$, $\varpi(c)=(c^2-\frac14)(c^2+\gamma_1^2)$, $\partial=d/dt$, and
\[
F:=\varpi(\partial)G=G''''+(\gamma_1^2-\tfrac14)G''-\tfrac{\gamma_1^2}4G .
\]
Then $F$ is real, supported in $[-T,T]$, bounded, and $F\in\PT$: four integrations by parts, the
boundary terms carrying $G^{(j)}(\pm T)=0$ for $j\le3$, give $\int Fe^{ct}dt=\varpi(c)\int Ge^{ct}dt$
for every $c\in\C$, and $\varpi(\pm\frac12)=0$.
\emph{(iv)} $\Fh(\tau)=(\tau^2+\frac14)(\tau^2-\gamma_1^2)\,T\,\widehat\omega_k(T\tau)$,
$\Phi_F(s)=\varpi(s-\frac12)\Phi_G(s)$ and $\Phi_G(\frac12+\beta+i\gamma)=T\,\widehat\omega_k\big(T(\gamma-i\beta)\big)$;
in particular $\Phi_F(\rho_1)=\Phi_F(\bar\rho_1)=0$.
\emph{(v)} With $Y:=T\gamma_1/\sqrt2$ and $\varepsilon_k(Y):=\frac1\pi k^{2k}Y^{1-2k}/(2k-1)$,
$\|F\|^2\ge\frac{\gamma_1^4}{64}\,T\,\big(\frac12-\varepsilon_k(Y)\big)$ whenever $\varepsilon_k(Y)<\frac12$.
\end{lemma}
\begin{proof}
(i) $\omega_k$ is a spline of degree $k-1$ supported in $[-1,1]$ and vanishing there to order $k-1$;
$1=(\int_{-1}^1\omega_k)^2\le2\int\omega_k^2$. (ii) $\widehat b_k(x)=\sin(x/k)/(x/k)$; $|\sin y/y|\le\min(1,1/|y|)$
on $\R$, and $|\sin(a+ib)|\le\cosh b\le e^{|b|}$. (iii) $G\in C^3$ with $G'''$ absolutely continuous
and $G''''\in L^\infty$ ($k=5$), and $G^{(j)}(\pm T)=0$ for $j\le3$; $\int G''e^{ct}=c^2\int Ge^{ct}$.
(iv) is (iii) with $c=i\tau$, resp.\ $c=s-\frac12$, and $\varpi(\pm i\gamma_1)=0$. (v) On
$|\tau|\le\gamma_1/\sqrt2$ one has $(\tau^2+\frac14)^2\ge\frac1{16}$ and $(\tau^2-\gamma_1^2)^2\ge\gamma_1^4/4$,
so $\|F\|^2=\frac1{2\pi}\int|\Fh|^2\ge\frac{\gamma_1^4}{64}\,T\,\frac1{2\pi}\int_{|x|\le Y}|\widehat\omega_k|^2$, and
$\frac1{2\pi}\int_{|x|\le Y}|\widehat\omega_k|^2\ge\|\omega_k\|_2^2-\frac1\pi\int_Y^\infty(k/x)^{2k}dx\ge\frac12-\varepsilon_k(Y)$.
\end{proof}

\begin{theorem}[the margin decays]\label{thm:margin}
Let $T>0$, $k\ge5$ an integer, $p:=2k-8$, $Y:=T\gamma_1/\sqrt2$, and suppose $\varepsilon_k(Y)<\frac12$. Then
\begin{equation}\label{eq:Bmargin}
\lamW\ \le\ \lamP\ \le\ \lamA\ \le\ \mathcal B(T,k):=
\frac{2\cdot1.0012\cdot T^2(k/T)^{2k}\big[\Sigma_p(21)+e^T\,\Sigma_p(H_{\mathrm{PT}})\big]}
{(14.13^4\,T/64)\,\big(\frac12-\varepsilon_k(Y)\big)},
\end{equation}
with $\Sigma_p(H)$ bounded by Lemma~\ref{lem:zerosum}. Consequently $\limsup_{T\to\infty}\lamA\le0$, and the
same for $\lamP$ and $\lamW$.
\end{theorem}
\begin{proof}
Take $F$ of Lemma~\ref{lem:testfam}. (1) $F\in\PT$, so $\lamA\le\mathfrak a_T(F)/\|F\|^2$, and
$\lamW\le\lamP\le\lamA$ by Theorem~\ref{thm:A} and Remark~\ref{rem:Kprim}.
(2) For $|\tau|\ge21$, $(\tau^2+\frac14)^2\le1.0012\,\tau^4$ and $(\tau^2-\gamma_1^2)^2\le\tau^4$, so
$|\Fh(\tau)|^2\le1.0012\,\tau^8T^2(k/(T|\tau|))^{2k}=O(|\tau|^{8-2k})$ with $8-2k\le-2$: $F$ is in the
form domain, which is hypothesis (a) of Lemma~\ref{lem:roughEF}; hypothesis (b) is verified in step (6).
Hence $\mathfrak a_T(F)=\sum_\rho\Phi_F(\rho)\Phi_F(1-\rho)$.
(3) $\rho_1,\bar\rho_1$ contribute $0$ (Lemma~\ref{lem:testfam}(iv)); by (Z1) no other zero has $|\Iml\rho|\le21$.
(4) A zero with $21<|\gamma|\le H_{\mathrm{PT}}$ lies on the line by (Z2) and contributes
$|\Fh(\gamma)|^2\le1.0012\,T^2(k/T)^{2k}|\gamma|^{8-2k}$.
(5) A zero $\rho=\frac12+\beta+i\gamma$ with $|\gamma|>H_{\mathrm{PT}}$ has $|\beta|<\frac12$;
$|\varpi(\beta+i\gamma)|\le(\gamma^2+\beta^2+\frac14)(\gamma^2+\beta^2+\gamma_1^2)\le(1+10^{-22})\gamma^4$, $\varpi$ is even,
and $|\Phi_G(\rho)|,|\Phi_G(1-\rho)|\le T(k/(T|\gamma|))^ke^{T|\beta|}$ by Lemma~\ref{lem:testfam}(ii),(iv);
so $|\Phi_F(\rho)\Phi_F(1-\rho)|\le1.0012\,T^2(k/T)^{2k}|\gamma|^{8-2k}e^T$.
(6) Summing (4)--(5) over $\pm\gamma$ with Lemma~\ref{lem:zerosum} ($p=2k-8\ge2$) gives
$\sum_\rho|\Phi_F(\rho)\Phi_F(1-\rho)|\le2\cdot1.0012\,T^2(k/T)^{2k}\big[\Sigma_p(21)+e^T\Sigma_p(H_{\mathrm{PT}})\big]<\infty$,
which is hypothesis (b) of Lemma~\ref{lem:roughEF} and bounds $\mathfrak a_T(F)$.
(7) Divide by Lemma~\ref{lem:testfam}(v) with $\gamma_1^4\ge14.13^4$. The limit statement is
Corollary~\ref{cor:margin1}.
\end{proof}

\begin{corollary}[closed form]\label{cor:margin1}
For every $T\ge1$,
\begin{equation}\label{eq:closedmargin}
\lamW\ \le\ \lamP\ \le\ \lamA\ \le\ 4\cdot10^9\,T\,e^{-15.449\,T},\qquad 15.449<\tfrac{42}e=15.4494\ldots\footnotemark
\end{equation}
\footnotetext{The exponent $\frac{42}{e}$ belongs to the test family of Lemma~\ref{lem:testfam}: one factor removes only the first zero, so the cutoff is $21$, just below the second ordinate $\gamma_2=21.022\ldots$, and $\frac{42}e$ comes from optimizing $k$ in $(k/21T)^{2k}$. It is a property of that family, not a sharp rate for $\lamA$.}
\end{corollary}
\begin{proof}
Take $k:=\lceil21T/e\rceil\ge8$ and write $k=21T/e+\theta$, $\theta\in[0,1)$. Then
$(k/21T)^{2k}=\exp\big(2k[\log(1+\theta e/21T)-1]\big)\le\exp(-42T/e+2\theta^2e/21T)\le e^{-15.4494\,T+0.259}$.
Next $(k/T)^{2k}\Sigma_p(21)\le21^9(k/21T)^{2k}\cdot\frac1{2\pi}\frac{2k-8}{2k-9}\big[\log21+\frac1{2k-9}\big]\le21^9\cdot0.5797\,(k/21T)^{2k}$
for $k\ge8$. The bound for the $H_{\mathrm{PT}}$-term is at most $10\,e^T(21/H_{\mathrm{PT}})^{2k-9}$ times that for the first, with
$21/H_{\mathrm{PT}}=7\cdot10^{-12}$ and $2k-9\ge15.449\,T-9$, hence less than $10^{-70}$ of it for $T\ge1$.
Finally $\varepsilon_k(Y)\le0.01566$ for all $T\ge1$: for fixed $k$, $\varepsilon_k$ decreases in $T$, so it
suffices to evaluate it at the left end of each $k$-range, namely $T=1$ for $k=8$ ($\varepsilon_8=0.0060$)
and $T=(k-1)e/21$ for $k\ge9$, where the largest value is $\varepsilon_9=0.01566$ at $T=8e/21$, the values
decreasing thereafter. So $\frac12-\varepsilon_k\ge0.484$, and
$\mathcal B(T,k)\le\frac{2\cdot1.0012\cdot0.5797\cdot21^9e^{0.259}}{(14.13^4/64)\cdot0.484}\,T\,e^{-15.4494\,T}(1+10^{-70})
\le3.97\cdot10^9\,T\,e^{-15.4494\,T}$.
\end{proof}

\begin{remark}[a cited doubly exponential ceiling]\label{rem:moriceiling}
Corollary~\ref{cor:margin1} is singly exponential in $T$ and is uninformative for $T\lesssim1.5$ (its right side
exceeds $1$ there). Mori \cite{MoriUpper} states, for the complete form, $\lamW\le3.50\cdot10^{24}\mu^{9/2}\log\mu\,
e^{-4\pi\mu}$ for $50\le\mu\le10^4$ and $\lamW\le7.58\cdot10^{24}\mu^5(\log\mu)^3e^{-4\pi\mu}$ for $\mu\ge50$, with
$\mu=e^{2T}$, by an explicit prolate test vector; that is, a ceiling of order $\exp(-4\pi e^{2T})$, doubly exponential in
$T$, consistent with the envelope threshold of Proposition~\ref{prop:env}. The proof uses zero data (verification of
the zeros up to a height and a zero-free region). We have not checked it line by line; we record it as plausible,
unverified, and cite it. It bounds $\lamW\le\lamA$ and so does not replace Corollary~\ref{cor:margin1} as a bound on
$\lamA$.
\end{remark}

\begin{corollary}[the structure of the margin]\label{cor:margin2}
\emph{(i)} There is no $c>0$ with $\lamA\ge c$ for all $T$, nor for $\lamP$ or $\lamW$; for $T\ge1$ every
lower bound on any of the three lies below $4\cdot10^9Te^{-15.449T}$.
\emph{(ii)} The limits $L^{\mathrm A}:=\lim_{T\to\infty}\lamA$, $L^{\mathrm{prim}}:=\lim_{T\to\infty}\lamP$ and
$L^{\mathrm W}:=\lim_{T\to\infty}\lamW$ exist in $[-\infty,0]$, and $L^{\mathrm W}\le L^{\mathrm{prim}}\le L^{\mathrm A}$.
In particular \eqref{eq:target} is equivalent to $L^{\mathrm A}=0$.
\end{corollary}
\begin{proof}
(i) is Corollary~\ref{cor:margin1}. (ii) Monotone limits (Lemma~\ref{lem:monoT}) bounded above by
Theorem~\ref{thm:margin}; the last sentence follows from (ii) and monotonicity.
\end{proof}

% i.e. before \section{The complete form at a single window}.
\section{Inertia and remainders}\label{sec:inertia}

This section records what the inertia of the forms of this paper, that is, the number of negative
directions rather than the value of the bottom, can and cannot decide. All forms are on $L^2(\IT)$ with
its metric; $n_-(\mathfrak q|\mathcal V)$ is the largest dimension of a subspace of $\mathcal V\cap\mathcal D_T$ on which
$\mathfrak q<0$, as in Lemma~\ref{lem:parity}. Recall $\mathfrak a_T=\mT+\PiT$ with $\mT(F)=\frac1{2\pi}\int m|\Fh|^2$
and $\PiT(F)=-\sum_{n\in\NT}v_n\Rel\cF(\log n)$. Nothing in this section bears on the sign of $\mathfrak a_T$ at any
window.

\subsection{The Sonin split of the truncated form}
Connes and Consani \cite[Thm.~4.7, (83)--(84)]{ConnesConsani2021} define the Sonin projection $\mathbf S$ and a
real, even, continuous function $\varepsilon_\infty$ on $\R_+^*$ with
$\operatorname{Tr}(\mathcal U(f)\mathbf S)=W_\infty(f)+\int f(\rho^{-1})\varepsilon_\infty(\rho)\,d^*\rho$, where $\mathcal U$ is the
scaling action and $W_\infty$ the archimedean distribution. In the additive variable ($f=g*g^*$, $F(t)=g(e^t)$,
after the conjugation of \cite[App.~A]{ConnesConsani2021}) we write
\[
\mathrm{Son}(F):=\operatorname{Tr}\big(\mathcal U(g)\mathbf S\,\mathcal U(g)^*\big)\ge0,\qquad
\langle C_\varepsilon F,F\rangle:=\int_{-2T}^{2T}\cF(x)\,\varepsilon_\infty(e^{|x|})\,dx ,
\]
where $C_\varepsilon$ is the integral operator on $L^2(\IT)$ with the continuous kernel $\varepsilon_\infty(e^{|x-y|})$; it
is Hilbert--Schmidt.

\begin{theorem}[the Sonin split]\label{thm:inertia-split}
For every $T>0$ and every $F\in C_c^\infty(\IT)$, with no moment condition,
\[
\mathfrak a_T(F)=\mathrm{Son}(F)+\mathfrak r_T(F),\qquad
\mathfrak r_T(F):=-\langle C_\varepsilon F,F\rangle+\PiT(F).
\]
\end{theorem}
\begin{proof}
In the dictionary of \cite[(152)--(154)]{ConnesConsani2021}, $W_\infty(g*g^*)=\frac1{2\pi}\int h_+|\Fh|^2$ with
$h_+(\tau)=\Rel\psi(\frac14+\frac{i\tau}2)-\log\pi$, which is $m$ of \eqref{eq:symbol}; no polar term occurs, so
$W_\infty(g*g^*)=\mT(F)$, not $\Ainf(F)$. By \cite[(83)]{ConnesConsani2021} and cyclicity of the trace,
$\mT=\mathrm{Son}-C_\varepsilon$ on $C_c^\infty(\IT)$. Adding $\PiT$ gives $\mathfrak a_T=\mT+\PiT$, which is \eqref{eq:form},
since $\frac1{2\pi}\int\cos(\tau\log n)|\Fh|^2=\Rel\cF(\log n)$.
\end{proof}

The prime form is bounded, $|\PiT(F)|\le\Aconst\|F\|^2$ since $|\cF(\log n)|\le\|F\|^2$, but its operator is not compact
(Proposition~\ref{prop:shifts}).

\begin{proposition}[the prime form at the first contact]\label{thm:inertia-B2}
Let $\frac12\log2<T<\frac12\log3$, so that $\NT=\{2\}$ and $\Aconst=v_2=\sqrt2\log2$. Then the operator of $\PiT$ has
spectrum $\{-\frac{\log2}{\sqrt2},0,\frac{\log2}{\sqrt2}\}$, each point of infinite multiplicity, and
$\|\PiT\|=\Aconst/2$. For every $T>\frac12\log2$, $\PiT$ has infinite positive and infinite negative index on
$\PT\cap\{\int F=0\}$.
\end{proposition}
\begin{proof}
Put $a=\log2$. The map $F\mapsto(F(-T+y+ja))_{0\le j<n(y)}$, $y\in[0,a)$, $n(y)=\#\{j\ge0:-T+y+ja<T\}$, is unitary from
$L^2(\IT)$ onto $\int^\oplus_{[0,a)}\C^{n(y)}dy$, with $n(y)=2$ for $y<2T-a$ and $n(y)=1$ otherwise (here $a<2T<2a$).
Since $\Rel\cF(a)=\frac12\langle(V_a+V_a^*)F,F\rangle$ couples only $j$ and $j+1$ in a fibre, $\PiT$ acts on a
$2$-dimensional fibre as $-v_2\begin{psmallmatrix}0&1/2\\1/2&0\end{psmallmatrix}$, with eigenvalues $\pm v_2/2$, and as $0$ on a
$1$-dimensional one; both sets of $y$ have positive measure, whence the spectrum and the multiplicities. The
negative index is Theorem~\ref{thm:inertia-nogo} with $k=0$, and the positive index follows from the same proof
with the sign of one bump in each pair reversed (Corollary~\ref{cor:Qp}).
\end{proof}

By Theorem~\ref{thm:inertia-nogo} with $k=\varepsilon_\infty(e^{|\cdot|})$, which is bounded on $[0,2T]$ by continuity
\cite[(84)]{ConnesConsani2021}, the remainder $\mathfrak r_T=\mathfrak a_T-\mathrm{Son}$ has $n_-(\mathfrak r_T|\PT)=\infty$ for every
$T>\frac12\log2$, and no inequality $\mathfrak a_T-\mathrm{Son}\ge-(\text{a form of finite rank})+c\|\cdot\|^2$ holds on $\PT$. At
$T\le\frac12\log2$ the comb is empty and $\mathfrak r_T=-C_\varepsilon$ is compact, which is the regime of
\cite[\S6]{ConnesConsani2021}. The theorem uses neither the sign of $\mathfrak a_T$ nor the identity of the prime: the
Sonin trace, subtracted in full, leaves a remainder that no finite-rank correction controls once a contact is
present. Ceding a fraction of the Sonin trace restores a finite index:

\begin{theorem}[relative bound $0$ and finite index]\label{thm:inertia-finite}
Let $T>0$, $\eta\in(0,1]$, let $k\in L^\infty(-2T,2T)$ be real and even with $C_k$ the operator with kernel $k(x-y)$ on
$L^2(\IT)$, and $\kappa:=\|C_k\|$. Put $A':=\kappa+\Aconst$, let $s_*$ be the least $s\ge0$ with $\eta\,m(s)\ge2A'$, put
$L:=A'-\eta C_0>0$ and $c:=Ts_*$. Then on $\mathcal D_T$
\[
\eta\,\mT-C_k+\PiT\ \ge\ A'\,I-(A'+L)\,S_c ,
\]
with $S_c$ the time--band limiting operator of Lemma~\ref{lem:KyFan}. Consequently
\[
n_-\big(\eta\mT-C_k+\PiT\,\big|\,L^2(\IT)\big)\le\#\Big\{j:\lambda_j(S_c)\ge\frac{A'}{A'+L}\Big\}
\le\frac{A'+L}{A'}\cdot\frac{2Ts_*}{\pi},
\]
and the bound $\lambda_j(S_c)\le\Lambda_j(c)$ of Lemma~\ref{lem:KyFan} applies. With $k=(1-\eta)\varepsilon_\infty(e^{|\cdot|})$
one has $\eta\mT-C_k+\PiT=\mathfrak a_T-(1-\eta)\mathrm{Son}$ on $C_c^\infty(\IT)$, so
$n_-(\mathfrak a_T-(1-\eta)\mathrm{Son})<\infty$ for every $\eta>0$, while it is $\infty$ for $\eta=0$ and $T>\frac12\log2$.
Since $m(s)=\log\frac{|s|}{2\pi}+O(|s|^{-1})$ (Lemma~\ref{lem:mbounds}(ii)), $s_*=2\pi e^{2A'/\eta}(1+o(1))$ as
$\eta\downarrow0$.
\end{theorem}
\begin{proof}
By $\|C_k\|=\kappa$ and $\|\PiT\|\le\Aconst$, the form is $\ge\frac1{2\pi}\int(\eta m-A')|\Fh|^2$. By
Lemma~\ref{lem:mbounds}(i), $m$ is increasing in $|\tau|$, so $\eta m-A'\ge A'$ for $|\tau|\ge s_*$ and
$\eta m-A'\ge\eta C_0-A'=-L$ for $|\tau|<s_*$. Hence the form is $\ge A'\langle(I-S_c)F,F\rangle-L\langle S_cF,F\rangle$. On the
span of the eigenvectors of $S_c$ with eigenvalue $<A'/(A'+L)$ this is positive definite, which gives the first count;
the second follows since the eigenvalues of $S_c$ are nonnegative with sum $\operatorname{tr}S_c=2c/\pi$ (Lemma~\ref{lem:KyFan}).
\end{proof}
No rate for $n_-(\mathfrak a_T-(1-\eta)\mathrm{Son})$ as $\eta\downarrow0$ beyond the upper bound above is asserted.

\subsection{Quartets and the zero-side inertia}
For $z=\beta+i\gamma$ with $0<\beta<\frac12$ and $\gamma>0$, the four points $\frac12\pm z,\frac12\pm\bar z$ form a
\emph{quartet}. For real $F$ put $\mathsf A_z(F):=\Phi_F(\frac12+z)$, with $\Phi_F$ as in \S\ref{sec:intro}.

\begin{lemma}[signature of a quartet]\label{thm:inertia-quartet}
For real $F\in L^2(\IT)$ with $F=F_{\mathrm e}+F_{\mathrm o}$, the sum of $\Phi_F(\rho)\Phi_F(1-\rho)$ over the quartet is
\[
q_z(F)=4\Rel\big(\mathsf A_z(F_{\mathrm e})^2\big)-4\Rel\big(\mathsf A_z(F_{\mathrm o})^2\big),
\]
with no cross term between parities. For every $T>0$, $q_z$ has signature exactly $(2,2)$ on the real space $\PT$:
one negative direction in $\PT^{\mathrm{ev}}$ and one in $\PT^{\mathrm{odd}}$.
\end{lemma}
\begin{proof}
$\Phi_F(\frac12-z)=\Phi_{JF}(\frac12+z)$ and $JF=F_{\mathrm e}-F_{\mathrm o}$; for real $F$ the terms at $\bar z$ are the complex
conjugates of those at $z$. So the four terms sum to
$4\Rel\big((\mathsf A_z(F_{\mathrm e})+\mathsf A_z(F_{\mathrm o}))(\mathsf A_z(F_{\mathrm e})-\mathsf A_z(F_{\mathrm o}))\big)$. The real
quadratic form $\Rel(w^2)$ has signature $(1,1)$ in $(\Rel w,\Iml w)$, so it suffices that $\Rel\mathsf A_z,\Iml\mathsf A_z$ are
linearly independent on $\PT^{\mathrm{ev}}$ and on $\PT^{\mathrm{odd}}$. On even $F$, $\Rel\mathsf A_z(F)=\int F\cosh\beta t\cos\gamma t$,
$\Iml\mathsf A_z(F)=\int F\sinh\beta t\sin\gamma t$, and $\PT^{\mathrm{ev}}$ is the kernel of $F\mapsto\int F\cosh\frac t2$. A
dependence would give $a\cosh\beta t\cos\gamma t+b\sinh\beta t\sin\gamma t=c\cosh\frac t2$ on $(0,T)$, hence on $\R$ by
analyticity; comparing growth ($\beta<\frac12$) forces $c=0$, then $t=0$ gives $a=0$, and then $b=0$. The odd case
is the same with $\sinh\beta t\cos\gamma t$, $\cosh\beta t\sin\gamma t$ and $\sinh\frac t2$: growth gives $c=0$, and
$a\sinh\beta t\cos\gamma t+b\cosh\beta t\sin\gamma t=\frac12e^{\beta t}(a\cos\gamma t+b\sin\gamma t)+O(e^{-\beta t})$ forces $a=b=0$.
\end{proof}

\begin{proposition}[finite modifications of the zero set]\label{thm:inertia-rank}
Let $Z'$ be obtained from the multiset of nontrivial zeros of $\zeta$ by adding $m$ quartets and removing $m$ pairs
$\frac12\pm i\gamma^0_k$ of zeros on the critical line, and put $W^{Z'}_T(F)=\sum_{\rho\in Z'}\Phi_F(\rho)\Phi_F(1-\rho)$ for
real $F\in\mathcal D_T$ with $\sum_\rho|\Phi_F(\rho)\Phi_F(1-\rho)|<\infty$. Then $W^{Z'}_T-W_T$ is a real quadratic form of rank
$\le6m$, and $|n_-(W^{Z'}_T|\mathcal V)-n_-(W_T|\mathcal V)|\le6m$ for every subspace $\mathcal V$ on which both forms are defined.
\end{proposition}
\begin{proof}
By Lemma~\ref{lem:roughEF}, $W_T$ is the zero sum. Each added quartet contributes $q_z$, of rank $4$
(Lemma~\ref{thm:inertia-quartet}); each removed pair contributes $-2|\Phi_F(\frac12+i\gamma^0)|^2$, of rank $2$. A
perturbation of rank $r$ changes the negative index by at most $r$ (if $\mathcal N$ is negative for one form, then
$\mathcal N\cap\ker(\text{perturbation})$ has codimension $\le r$ in $\mathcal N$ and is negative for the other).
\end{proof}

\subsection{What these results exclude}\label{subsec:inertia-excl}
\emph{Sonin-norm bounds.} By Theorem~\ref{thm:inertia-nogo}, no bound $\mathfrak a_T\ge\mathrm{Son}-(\text{finite rank})$
holds on $\PT$ past $\frac12\log2$. By Theorem~\ref{thm:inertia-finite}, ceding a fraction $\eta>0$ of $\mathrm{Son}$
restores a finite index, but only up to $N\le\frac{A'+L}{A'}\cdot\frac{2Ts_*}\pi\asymp Te^{2A'/\eta}$ directions, and
positivity is then the positivity of a Schur complement of that size, the finite structure of
Theorem~\ref{thm:coer-schur} with a larger dimension.

\emph{Inertia alone (conditional).} By Proposition~\ref{thm:inertia-rank}, an argument that uses only the decomposition
$\mathfrak a_T=\mT-C_\varepsilon+\PiT$ together with operator properties of the three pieces that are stable under
perturbations of finite rank of the middle term (the archimedean multiplier; $\PiT$ with exactly the weights $v_n$;
a compact remainder), and inertia arithmetic (Sylvester, Haynsworth, interlacing, Lemma~\ref{lem:monoT}), applies
verbatim to the forms $W^{Z'}_T$ built from a modified zero set with $m$ added quartets. These have $2m$
negative directions of $q_z$ on $\PT$ for every $T$ (Lemma~\ref{thm:inertia-quartet}), carrying weight of order
$e^{2\beta T}$. \emph{Assume} that $n_-(W^{Z'}_T|\PT)\ge2m$ for all large $T$; this is not proved here. Then an argument
of the class just described cannot prove that $\sup_Tn_-(\mathfrak a_T|\PT)<\infty$ implies $n_-(\mathfrak a_T|\PT)=0$ for all
$T$. The exclusion concerns the class as defined: it does not cover arguments that use quantitative properties of
the compact term, such as a $T$-uniform bound on $\|C_\varepsilon\|$ or a sign.

\begin{remark}[conventions of the sources]\label{rem:inertia-conv}
The Fourier convention of \cite{ConnesConsani2021} ($e^{-2\pi ixy}$) differs from ours ($\Fh(\tau)=\int Fe^{i\tau t}$);
this is immaterial, since only $|\Fh|^2$ and even symbols enter. Their truncated form has no polar term and equals
$\mathfrak a_T$, not $W_T$; the two agree on $\PT$. Their one-directional bound for $T\le\frac12\log2$
\cite[\S5]{ConnesConsani2021} is stated in a different metric, not that of $L^2(\IT)$, and nothing in that metric
enters here.
\end{remark}

\section{The semilocal Sonin decomposition at the first prime}\label{sec:semilocal}

Theorem~\ref{thm:inertia-nogo} shows that the archimedean Sonin trace, subtracted in full, leaves a remainder of
infinite negative index as soon as the contact $n=2$ is present. Here we replace the archimedean Sonin space by
the semilocal one of Connes and Consani for $S=\{\infty,2\}$ and show that the comb of $2$ is then absorbed
exactly, with a compact remainder. The objects $\mathsf F_S$, the
quasi-inner function and the semilocal Sonin space are those of \cite{CCquasi2020,CCM2023}; what is new is listed in
Remark~\ref{rem:sl-new}. No statement of this section bears on the sign of $\mathfrak a_T$.

\subsection{Two models and their dictionary}\label{subsec:sl-dict}
The operators below live in the \emph{u-model} $\mathsf H=L^2(\R_+,du)$. Fix a level $\lambda>0$. The dictionary is:
\begin{enumerate}
\item[(D1)] $L^2(\R)^{\mathrm{ev}}\to\mathsf H$, $\xi\mapsto\sqrt2\,\xi|_{\R_+}$, is unitary; it carries the Fourier transform
$\xi\mapsto\int\xi(x)e^{-2\pi ixy}dx$ of \cite{ConnesConsani2021} to $\mathsf F_\infty$, the operator with kernel
$2\cos(2\pi ux)$, and the cut-off $\mathbf 1_{|x|\le\lambda}$ to $P=\mathbf 1_{(0,\lambda]}(u)$.
\item[(D2)] $\mathsf W:\mathsf H\to L^2(\R,dt)$, $(\mathsf W\xi)(t)=e^{t/2}\xi(e^t)$, is unitary; $P\mapsto\mathbf 1_{t\le\log\lambda}$,
dilations become translations, and $D:=-i\,d/dt$.
\item[(D3)] The scaling action is $(\mathcal U(\rho)\xi)(v)=\rho^{-1/2}\xi(\rho^{-1}v)$, $\mathcal U(f)=\int f(\rho)\mathcal U(\rho)d^*\rho$
\cite[(61)]{ConnesConsani2021}. Every functional below ($W_\infty$, the comb, and $\int f(\rho^{-1})\varepsilon(\rho)d^*\rho$ with
$\varepsilon$ inversion-symmetric) is invariant under $f\mapsto f(\cdot^{-1})$, so the opposite convention gives the same
identities.
\item[(D4)] For $F\in C_c^\infty(\IT)$ put $g_F(\rho):=F(\log\rho)$ and $f_F:=g_F*g_F^*$, so that $f_F(e^x)=\cF(x)$,
$\int f_F(\rho^{-1})\varepsilon(\rho)d^*\rho=\int_{-2T}^{2T}\cF(x)\varepsilon(e^{|x|})dx$, and $W_\infty(f_F)=\mT(F)$ (proof of
Theorem~\ref{thm:inertia-split}).
\item[(D5)] For a projection $\Pi'$ on $\mathsf H$, $\mathrm{Son}_{\Pi'}(F):=\operatorname{Tr}(\mathcal U(g_F)\Pi'\,\mathcal U(g_F)^*)
=\operatorname{Tr}(\mathcal U(f_F)\Pi')\ge0$.
\end{enumerate}
All operator statements of \S\S\ref{subsec:sl-dict}--\ref{subsec:sl-trace} are in $\mathsf H$ (equivalently $L^2(\R,dt)$);
the forms of \S\ref{subsec:sl-split} are on $L^2(\IT)$ with its own metric, through (D4).

\subsection{The semilocal Fourier transform}
\begin{definition}[\cite{CCM2023,CCquasi2020}]\label{def:FS}
Let $D_2\xi(u)=\xi(2u)$, so that $\|D_2\|=2^{-1/2}$, and $\mathsf E:=\sum_{k\ge0}D_2^k$ on $\mathsf H$, so that $\mathsf E^{-1}=1-D_2$. The semilocal Fourier transform is $\mathsf F_S:=\mathsf E\,\mathsf F_\infty\mathsf E^{-1}$
\cite[Prop.~4.1]{CCM2023}. In Mellin form ($\Rel s=\frac12$) $\mathsf F_S$ is $\mathsf F_\infty$ composed with the multiplier
$(1-2^{-s})/(1-2^{-(1-s)})$, the factor $\rho_2$ of the quasi-inner function of \cite[Lemma~3.1]{CCquasi2020}. Put
$\widehat P^S:=\mathsf F_SP\mathsf F_S$, $Q_1:=1-P$, $Q_2:=1-\widehat P^S$, $Z_S:=P\mathsf F_SP$, and let $\Pi_S$ be the orthogonal
projection onto $\operatorname{ran}Q_1\cap\operatorname{ran}Q_2$, the semilocal Sonin space \cite{CCquasi2020}. For $S=\{\infty\}$
these are $\widehat P_\infty$, $Z_\infty$ and the archimedean Sonin projection $\Pi_\infty=\mathbf S$.
\end{definition}

\begin{lemma}[u-model form; scattering phase]\label{lem:sl-phase}
In $L^2(\R,dt)$, $\mathsf F_\infty=J\,m_\infty(D)$ with $J$ the reflection and $|m_\infty|=1$, and
$\mathsf F_S=\mathsf F_\infty\mathsf V=J\,M(D)$, $M=m_\infty e^{i\varphi}$, $\mathsf V=e^{i\varphi(D)}$, with
\[
\varphi'(t)=-\sum_{k\ge1}v_{2^k}\cos(kt\log2);\qquad\text{for finite }S:\quad
\varphi_S'(t)=-\sum_{p\in S}\sum_{k\ge1}v_{p^k}\cos(kt\log p).
\]
The comb $\varphi_S'$ is untruncated; in pairings with $\cF$ only $n\in\NT$ survive, since $\supp\cF\subset[-2T,2T]$.
\end{lemma}
\begin{proof}
The identification of $\mathsf V$ with $\rho_2$ is \cite[Lemma~3.1]{CCquasi2020}; we verify the normalization. With
$q=2^{-1/2}$, $(\mathsf WD_2\mathsf W^{-1}G)(t)=e^{t/2}(2e^t)^{-1/2}G(t+\log2)$, i.e.\ $\mathsf WD_2\mathsf W^{-1}=qV_{-\log2}$, so
$\mathsf E^{-1}=1-qV_{-\log2}$, $\mathsf F_\infty\mathsf E\mathsf F_\infty=\mathsf E^*$ (since $\mathsf F_\infty$ reverses translations), and
$\mathsf V=\mathsf F_\infty\mathsf F_S=\mathsf E^*\mathsf E^{-1}=(1-qV_{\log2})^{-1}(1-qV_{-\log2})$, a quotient of conjugate multipliers. Its
phase is $\varphi(t)=2\arg(1-qe^{it\log2})=-2\sum_kq^k\sin(kt\log2)/k$, whence $\varphi'$ with $2q^k\log2=v_{2^k}$.
$J\mathsf F_\infty$ commutes with translations and is unitary, hence a multiplier, the Tate factor. For several primes
$\mathsf V_S=\prod_p\mathsf V_p$.
\end{proof}

\begin{proposition}\label{prop:sl-AS}
\emph{(i)} $\mathsf F_S$ is a unitary, self-adjoint involution, real, with $\mathsf F_S\,\mathcal U(\rho)\mathsf F_S=\mathcal U(\rho^{-1})$.
\emph{(ii)} $Z_S$ is compact and self-adjoint, $\|Z_S\|<1$, and $\pm1$ are not eigenvalues.
\emph{(iii)} For $S=\{\infty,2\}$, $Z_S$ is not Hilbert--Schmidt, and $Z_S\in\mathcal S_p$ for every $p>2$.
\end{proposition}
\begin{proof}
(i) $\mathsf F_S$ is similar to the involution $\mathsf F_\infty$, and unitary by Lemma~\ref{lem:sl-phase}; a unitary involution
is self-adjoint. $\mathsf F_\infty\mathcal U(\rho)\mathsf F_\infty=\mathcal U(\rho^{-1})$ and $\mathsf V$ commutes with $\mathcal U(\rho)$.
(ii) Compactness is \cite[Thms.~2.1, 4.4]{CCquasi2020}; directly, with $c_k(u,x)=\cos(2\pi2^kux)$ the grouped kernels
$\sum_{k\le K}(2c_k-c_{k-1})$ give Hilbert--Schmidt operators on $(0,\lambda]$ converging in norm, with
$\|PD_2^k\mathsf F_\infty\mathsf E^{-1}P\|\le(1+q)q^k$. A compact self-adjoint contraction has norm $<1$ unless $\pm1$ is an
eigenvalue. If $Z_S\xi=\pm\xi$ with $\|\xi\|=1$, then $\|P\mathsf F_S\xi\|=\|\mathsf F_S\xi\|$ forces $\mathsf F_S\xi=\pm\xi$, so
$\mathsf F_\infty\mathsf E^{-1}\xi=\pm\mathsf E^{-1}\xi$; but $\mathsf E^{-1}\xi=\xi-D_2\xi$ is a nonzero function supported in $(0,\lambda]$ (it
vanishes only if $\xi=D_2\xi$, i.e.\ $\xi=0$ by iteration), and a nonzero compactly supported function on $\R_+$ has a
cosine transform that is not compactly supported (Paley--Wiener).
(iii) If $Z_S$ had a kernel $\kappa\in L^2((0,\lambda]^2)$, test against $\chi\,c_m$ with $\chi\in C_c^\infty((\delta,\lambda-\delta)^2)$:
non-stationary phase in the lacunary sum gives $\langle\kappa,\chi c_m\rangle\to\frac12\iint\chi\ne0$ as $m\to\infty$ (the term
$k=m$ is the only non-oscillating one), while the Riemann--Lebesgue lemma gives $\to0$. For $p>2$,
$\|Pc_kP\|_{\mathcal S_p}\le\|Pc_kP\|^{1-2/p}\|Pc_kP\|_{\mathrm{HS}}^{2/p}\le(\tfrac12q^k)^{1-2/p}\lambda^{2/p}$, which is summable in $k$.
\end{proof}

\begin{lemma}[trace formula for the linear part]\label{lem:sl-linear}
Let $h$ be such that $y\mapsto h(e^y)$ is a Schwartz function, and $S$ finite. Then $\mathcal U(h)(\widehat P^S-\widehat P_\infty)$
is trace class, and for every $\lambda>0$
\[
\operatorname{Tr}\mathcal U(h)(\widehat P^S-\widehat P_\infty)=\sum_{p\in S}\sum_{k\ge1}\log p\;p^{-k/2}\big(h(p^k)+h(p^{-k})\big)
=\tfrac12\sum_{p\in S}\sum_{k\ge1}v_{p^k}\big(h(p^k)+h(p^{-k})\big).
\]
\end{lemma}
\begin{proof}
Write $\mathsf V=\sum_ju_jV_{h_j}$ with $\sum_j(1+|h_j|)|u_j|<\infty$: for $S$ finite, $h_j$ ranges over $\pm\log n$, $n$ in the
multiplicative semigroup generated by $S$, and $u_j$ is a product of geometric sequences. With $P_a$ the projection
onto $t\le a$, $\mathsf VP_a\mathsf V^*-P_a=\sum_{j,l}u_j\bar u_l(P_{a+h_j}-P_a)V_{h_j-h_l}$, and $P_{a+h}-P_a=\operatorname{sgn}(h)\mathbf 1_{I_h}$ with
$|I_h|=|h|$. By \cite[Thm.~4.5]{Simon2005}, $\|\widehat k(D)\mathbf 1_I\|_1\le C_k(1+|I|)$ for the convolution kernel $k$ of
$\mathcal U(h)$ in the $t$-variable, so the series converges in trace norm, and the trace is
$\sum_{j,l}u_j\bar u_l\,h_j\,k(h_j-h_l)$ up to sign, which equals $-\frac1{2\pi}\int\widehat k\,\varphi'$. Conjugating by
$\mathsf F_\infty$ reverses $D$, and Lemma~\ref{lem:sl-phase} gives the comb.
\end{proof}
This is the dual form of the classical phase-shift trace formula, and it is the local term at $p$ of Connes'
semilocal trace formula; it concerns the linear part $1-\widehat P$, not the Sonin projection $\Pi_S$.

\subsection{Gap and trace class}\label{subsec:sl-trace}
Let $\{\xi_n\}$ be a real orthonormal eigenbasis of $Z_S$ on $\operatorname{ran}P$, $Z_S\xi_n=\lambda_n\xi_n$, ordered by $|\lambda_n|$
nonincreasing, so $\lambda_1^2=\|Z_S\|^2<1$; put $\zeta_n:=(1-P)\mathsf F_S\xi_n/(1-\lambda_n^2)^{1/2}$ and
$\tau_n:=\lambda_n(1-\lambda_n^2)^{-1/2}$.

\begin{theorem}[operator gap]\label{thm:sl-gap}
The $\zeta_n$ are orthonormal, $\Pi_S=(1-P)-\sum_n|\zeta_n\rangle\langle\zeta_n|$, and
\[
Q_1Q_2Q_1=\Pi_S+\sum_n\lambda_n^2|\zeta_n\rangle\langle\zeta_n| .
\]
Hence $\operatorname{spec}(Q_1Q_2Q_1)\subset\{0\}\cup\{\lambda_n^2\}\cup\{1\}$ with $\sup_n\lambda_n^2=\lambda_1^2<1$, the point $1$ is isolated,
$\Pi_S$ is the Riesz projection of $Q_1Q_2Q_1$ at $1$, and $\|(Q_1Q_2Q_1)^k-\Pi_S\|=\lambda_1^{2k}$ for $k\ge1$. In the Halmos
decomposition of the pair $(Q_1,Q_2)$ the squared cosines of the angles are the $\lambda_n^2$.
\end{theorem}
\begin{proof}
Put $X:=(1-P)\mathsf F_SP$. Since $\mathsf F_S^2=1$ (Proposition~\ref{prop:sl-AS}(i)),
$X^*X=P\mathsf F_S(1-P)\mathsf F_SP=P-Z_S^2$, so $X^*X\xi_n=(1-\lambda_n^2)\xi_n$, $\|X\xi_n\|^2=1-\lambda_n^2>0$, and the
$\zeta_n=X\xi_n/\|X\xi_n\|$ are orthonormal and span $\overline{\operatorname{ran}X}$. On $\operatorname{ran}Q_1$,
$Q_1Q_2Q_1=(1-P)-XX^*=(1-P)-\sum_n(1-\lambda_n^2)|\zeta_n\rangle\langle\zeta_n|$. A vector $y\in\operatorname{ran}(1-P)$ lies in
$\operatorname{ran}Q_2$ iff $P\mathsf F_Sy=0$, i.e.\ iff $X^*y=0$, i.e.\ iff $y\perp\operatorname{ran}X$; so
$\Pi_S=(1-P)-\sum_n|\zeta_n\rangle\langle\zeta_n|$, which gives the displayed identity, and $Q_1Q_2Q_1$ vanishes on
$\operatorname{ran}P$. The spectral statements follow, with $(Q_1Q_2Q_1)^k-\Pi_S=\sum_n\lambda_n^{2k}|\zeta_n\rangle\langle\zeta_n|$. The
last sentence is the definition of the Halmos angles for the generic part of the pair.
\end{proof}
The gap is spectral, at $1$, for $Q_1Q_2Q_1$. By Proposition~\ref{prop:sl-AS}(iii), $\sum\lambda_n^2=\infty$ while
$\sum|\lambda_n|^p<\infty$ for $p>2$: the operator $Q_1Q_2Q_1-\Pi_S$ is in $\mathcal S_{p/2}$ for $p>2$, in particular
Hilbert--Schmidt, and not trace class.

\begin{lemma}[trace class]\label{lem:sl-TC}
For $g,f\in C_c^\infty(\R_+^*)$: $\mathcal U(g)Q_1Q_2$ is Hilbert--Schmidt; $\mathcal U(f)Q_1Q_2Q_1$, $\mathcal U(f)\Pi_S$ and
$\mathcal U(f)(\widehat P^S-\widehat P_\infty)$ are trace class; and $P\,\mathcal U(f)P\,Z_S^2\in\mathcal S_1$. In particular
$\sum_n\lambda_n^2\langle\xi_n,\mathcal U(f)\xi_n\rangle$ converges absolutely although $\sum_n\lambda_n^2=\infty$.
\end{lemma}
\begin{proof}
Write $a=\log\lambda$. By Lemma~\ref{lem:sl-phase} and (D2), $Q_2=\mathsf F_S(1-P)\mathsf F_S$ and
$Q_1Q_2=(1-P)JM(D)(1-P)JM(D)=J\,\mathbf 1_{t<-a}M(D)\mathbf 1_{t>a}M(D)^*J$, using $J\mathbf 1_{t>a}J=\mathbf 1_{t<-a}$ and
$M(-\tau)=\overline{M(\tau)}$. In the $t$-variable $\mathcal U(g)$ is a convolution $\widehat{\tilde g}(D)$ with a Schwartz symbol.
Split $\widehat{\tilde g}(D)\mathbf 1_{t<-a}M(D)\mathbf 1_{t>a}$ into the commutator $[\widehat{\tilde g}(D),\mathbf 1_{t<-a}]M(D)\mathbf 1_{t>a}$,
whose Hilbert--Schmidt norm is at most $(\int|k|^2|x|dx)^{1/2}$ with $k$ the kernel of $\widehat{\tilde g}(D)$, and the
Hankel block $\mathbf 1_{t<-a}(\widehat{\tilde g}M)(D)\mathbf 1_{t>a}$, whose squared Hilbert--Schmidt norm is
$\int|K_g(x)|^2(-2a-x)_+dx$ with $K_g$ the kernel of $(\widehat{\tilde g}M)(D)$; this is finite because
$\widehat{\tilde g}M$ is a Schwartz function ($m_\infty$ and its derivatives grow at most polynomially by Stirling's formula,
and $e^{i\varphi}\in C_b^\infty$). Trace class of $\mathcal U(f)Q_1Q_2Q_1$ follows by writing $f=g_1*g_2$ (Dixmier--Malliavin)
and using \cite[Lemma~D.1]{ConnesConsani2021} for the commutators; then $\mathcal U(f)\Pi_S=\mathcal U(f)Q_1Q_2Q_1\Pi_S$. The
operator $\mathcal U(f)(\widehat P^S-\widehat P_\infty)$ is Lemma~\ref{lem:sl-linear}. Finally, with $X'=\widehat P^S-\widehat P_\infty$,
$P\,\mathcal U(f)X'P=\Theta_{00}(Z_S^2-Z_\infty^2)+\Theta_{01}X'_{10}$, where $\Theta_{00}=P\mathcal U(f)P$,
$\Theta_{01}=P\,\mathcal U(f)(1-P)\in\mathcal S_1$ \cite[Thm.~4.5]{Simon2005} and $X'_{10}=(1-P)X'P$; since $Z_\infty\in\mathcal S_1$ and
the left side is trace class, so is $\Theta_{00}Z_S^2$.
\end{proof}

For $f\in C_c^\infty(\R_+^*)$ put
$\mathrm{Hal}_S(f):=\sum_n\tau_n\big(\langle\xi_n,\mathcal U(f)\zeta_n\rangle+\langle\zeta_n,\mathcal U(f)\xi_n\rangle\big)$, the Halmos cross series of
\cite[(84)]{ConnesConsani2021} in the semilocal setting, and $G:=(1-P)\mathsf F_SZ_S(1-Z_S^2)^{-1}P$, $G_0:=(1-P)\mathsf F_SP\mathsf F_SP$.
The series converges absolutely: its two halves are traces of $\Theta_{01}G$ and $\Theta_{10}G^*$ in an orthonormal
basis, with $\Theta_{01},\Theta_{10}$ trace class.

\begin{lemma}[a lacunary sum]\label{lem:sl-lac}
Let $\mathrm s(w)=\sin(2\pi\lambda w)/(2\pi w)$. Then uniformly in $t\ne0$ and $k_0\in\Z$,
\[
\sum_{k\ge k_0}\mathrm s(2^kt)=\lambda\,\#\{k\ge k_0:\ 2^k|t|\le1\}+O(1),
\]
with an implied constant depending only on $\lambda$.
\end{lemma}
\begin{proof}
For $2^k|t|\le1$ write $\mathrm s(w)=\lambda+O(\lambda^3w^2)$; the error terms sum to $O(\lambda^3\sum_{2^k|t|\le1}4^kt^2)=O(\lambda^3)$, a
geometric series dominated by its largest term. For $2^k|t|>1$, $|\mathrm s(2^kt)|\le1/(2\pi2^k|t|)$, a geometric series
dominated by its first term, $\le1/(2\pi)\cdot2$.
\end{proof}


For $u\ge\lambda>z>0$ put $a_{-1}=-1$, $a_k=1$ for $k\ge0$, and
\begin{equation}\label{eq:G0series}
G_0(u,z):=\frac12\sum_{k,l\ge-1}a_ka_l\big[\mathrm s(2^ku-2^lz)+\mathrm s(2^ku+2^lz)\big],
\end{equation}
the series that the proof of Theorem~\ref{thm:sl-L2} identifies with the kernel of $G_0$.

\begin{lemma}[the kernel of $G_0$, uniformly in the first variable]\label{lem:sl-HG}
The series \eqref{eq:G0series} converges absolutely whenever $u\notin\{2^iz:i\ge0\}$, and for every $R>1$ there is
$C=C(\lambda,R)$ such that
\[
|G_0(u,z)|\le C\Big(1+\log\frac\lambda z\Big)+\lambda\sum_{i\ge0}\log_2^+\frac1{|u-2^iz|}
\qquad(\lambda\le u\le\lambda R,\ 0<z<\lambda).
\]
Consequently
\begin{equation}\label{eq:sl-HG}
\sup_{\lambda\le u\le\lambda R}\|G_0(u,\cdot)\|_{L^2(0,\lambda)}<\infty\qquad\text{for every }R>1 .
\end{equation}
\end{lemma}
\begin{proof}
Since $|\sin x|\le|x|$, $|\mathrm s(w)|\le\lambda$, and $|\mathrm s(w)|\le1/(2\pi|w|)$. Hence, for $w\ne0$ and $k_0\ge-1$,
\begin{equation}\label{eq:sl-HG1}
\sum_{k\ge k_0}|\mathrm s(2^kw)|\le\frac{2^{-k_0}}{\pi|w|},\qquad
\sum_{k\ge-1}|\mathrm s(2^kw)|\le\lambda N(w)+\frac1\pi,
\end{equation}
where $N(w):=\#\{k\ge-1:2^k|w|\le1\}\le2+\log_2^+\frac1{|w|}$; the second because the terms with $2^k|w|>1$ form a geometric series whose first term is less than $1/(2\pi)$.
Group the series \eqref{eq:G0series} by $i=l-k\in\Z$; for fixed $i$ the index $k$ runs over $k\ge k_i:=\max(-1,-1-i)$,
and the terms are $\mathrm s(2^k(u-2^iz))$ and $\mathrm s(2^k(u+2^iz))$. Let $i_0:=\lfloor\log_2(u/z)\rfloor\ge0$, so that
$2^{i_0}z\le u<2^{i_0+1}z$ and $i_0\le\log_2(\lambda R/z)$.

\emph{The terms with the plus sign.} Here $w=u+2^iz\ge u\ge\lambda$. For $i<0$, $2^{-k_i}=2^{1+i}$ and
\eqref{eq:sl-HG1} gives a total of at most $\sum_{i<0}2^{1+i}/(\pi\lambda)=2/(\pi\lambda)$. For $0\le i\le i_0$ each
diagonal contributes at most $2/(\pi u)\le2/(\pi\lambda)$. For $i>i_0$, $w\ge2^iz>2^{i-i_0-1}u\ge2^{i-i_0-1}\lambda$, and the
total is at most $\sum_{j\ge1}2/(\pi2^{j-1}\lambda)=4/(\pi\lambda)$. Altogether at most
$\frac2{\pi\lambda}\big(4+\log_2\frac{\lambda R}{z}\big)$.

\emph{The terms with the minus sign.} Here $w_i=u-2^iz$. For $i\le i_0-1$, $2^iz\le u/2$, so $w_i\ge u/2\ge\lambda/2$,
and \eqref{eq:sl-HG1} gives at most $2^{-k_i}\cdot2/(\pi\lambda)$ per diagonal: at most $4/(\pi\lambda)$ in total over
$i<0$ and at most $4i_0/(\pi\lambda)$ over $0\le i\le i_0-1$. For $i\ge i_0+2$, $2^iz\ge2u$, so
$|w_i|\ge2^{i-1}z\ge2^{i-i_0-2}\lambda$, and the total is at most $\sum_{j\ge0}2/(\pi2^j\lambda)=4/(\pi\lambda)$. For
$i\in\{i_0,i_0+1\}$ the second bound in \eqref{eq:sl-HG1} gives at most $\lambda(2+\log_2^+\frac1{|w_i|})+\frac1\pi$ each.

Summing, and using $i_0\le\log_2(\lambda R/z)$, every term of \eqref{eq:G0series} is dominated by a convergent series
off the ridges $u=2^iz$, and
\[
|G_0(u,z)|\le\frac1{2}\Big[\frac{2}{\pi\lambda}\Big(4+\log_2\frac{\lambda R}z\Big)+\frac{8}{\pi\lambda}+\frac{4}{\pi\lambda}\log_2\frac{\lambda R}z
+4\lambda+\frac2\pi\Big]+\frac\lambda2\sum_{i\in\{i_0,i_0+1\}}\log_2^+\frac1{|u-2^iz|},
\]
which is the stated bound. For the $L^2$ norm: $\int_0^\lambda(1+\log\frac\lambda z)^2dz=5\lambda$, and the function
$f_i(z)=\log_2^+\frac1{2^i|z-u2^{-i}|}$ is supported in $|z-u2^{-i}|<2^{-i}$ with
$\|f_i\|_{L^2(\R)}^2=2^{-i}\int_{|x|<1}\log_2^2\frac1{|x|}dx=2^{2-i}/\log^22$, so by Minkowski's inequality
$\|G_0(u,\cdot)\|_{L^2(0,\lambda)}\le C\sqrt{5\lambda}+\lambda\sum_{i\ge0}2^{1-i/2}/\log2<C\sqrt{5\lambda}+10\lambda$, uniformly in
$u\in[\lambda,\lambda R]$.
\end{proof}

\begin{theorem}[structure of the Halmos density]\label{thm:sl-L2}
Let $S=\{\infty,2\}$. The functional $\mathrm{Hal}_S$ is given by a density: $\mathrm{Hal}_S(f)=\int f(\rho^{-1})d_S(\rho)\,d^*\rho$ with
$d_S(\rho^{-1})=d_S(\rho)$ and, for $\rho\ge1$, $d_S(\rho)=\rho^{1/2}\int_{\lambda/\rho}^{\lambda}G(\rho x,x)\,dx$.
For every $R>1$
\[
d_S(\rho)=\sum_{0\le j\le\log_2R}\beta_j\log^+\frac1{|\rho-2^j|}+O_R(1)\quad(1\le\rho\le R),\qquad
\beta_j=2^{j/2}(1-2^{-j})\frac{\lambda^2}{2\log2},
\]
and $d_S(\rho)\sim b\,(\rho-1)\log\frac1{\rho-1}$ as $\rho\downarrow1$, $b=\lambda^2/(2\log2)$. Hence
$d_S\in L^p_{\mathrm{loc}}(\R_+^*)$ for every $p<\infty$ and $d_S\notin L^\infty_{\mathrm{loc}}$.
\end{theorem}
\begin{proof}
The density formula is the computation of the two traces above in the kernel form of $G$ and of the convolution
$\mathcal U(f)$. Here $G=G_0+G_0R'$ below, $G_0R'$ is Hilbert--Schmidt (a bounded operator times
a Hilbert--Schmidt one), so its kernel is square integrable and contributes to $d_S$ a function in
$L^2_{\mathrm{loc}}$ after the change of variables $u=\rho x$; the explicit kernel of $G_0$ is locally integrable by
Lemma~\ref{lem:sl-HG}. So $d_S\in L^1_{\mathrm{loc}}(\R_+^*)$, and
Lemma~\ref{lem:sl-HG} is used again for the $O_R(1)$ bound on the $G_0R'$ part. Next, $G=G_0(1+R')$ with $R'=Z_S^2(1-Z_S^2)^{-1}$, which is Hilbert--Schmidt by
Proposition~\ref{prop:sl-AS}(iii) ($\sum\lambda_n^4<\infty$). Its contribution to $d_S$ is bounded on $[1,R]$: if $r'$ is the
kernel of $R'$, then $|(G_0R')(u,x)|\le\|G_0(u,\cdot)\|_{L^2(0,\lambda)}\|r'(\cdot,x)\|$, so
$|\rho^{1/2}\int_{\lambda/\rho}^\lambda(G_0R')(\rho x,x)dx|\le R^{1/2}\lambda^{1/2}\|R'\|_{\mathrm{HS}}\sup_u\|G_0(u,\cdot)\|$, which is finite by
Lemma~\ref{lem:sl-HG}. The kernel of $G_0$ on $u>\lambda\ge z>0$ is the series \eqref{eq:G0series}, absolutely
convergent off the ridges $u=2^iz$ by Lemma~\ref{lem:sl-HG}, from $\int_0^\lambda\cos(2\pi\alpha y)\cos(2\pi\beta y)dy=\frac12[\mathrm s(\alpha-\beta)+\mathrm s(\alpha+\beta)]$.
Grouping by the diagonal $i=l-k$ and applying Lemma~\ref{lem:sl-lac} to each diagonal gives
$G_0(\rho x,x)=\frac{\lambda}{2\log2}\sum_i\log^+\frac1{|\rho-2^i|}+O_R(1)$ uniformly in $x\in[\lambda/R,\lambda]$: on the diagonal $i$ the
terms $\mathrm s(2^k(\rho x-2^ix))$ with $2^k|\rho-2^i|x\le1$ contribute $\lambda$ each, and their number is
$\log_2\frac1{|\rho-2^i|x}+O(1)$. Integrating over $x\in[\lambda/\rho,\lambda]$ and multiplying by $\rho^{1/2}$ gives the
coefficient $\rho^{1/2}\lambda(1-\rho^{-1})\cdot\frac{\lambda}{2\log2}$, which at $\rho=2^j$ is $\beta_j$, and the behaviour at
$\rho\downarrow1$. The integrability statements follow because $\log^+$ is in every $L^p_{\mathrm{loc}}$, and $\beta_j>0$ for
$j\ge1$.
\end{proof}

\begin{theorem}[the semilocal trace formula]\label{thm:sl-84}
Let $S=\{\infty,2\}$, $\lambda>0$, $f\in C_c^\infty(\R_+^*)$ real with $f(\rho^{-1})=f(\rho)$, and
$\mathrm{comb}_2(f):=\sum_{k\ge1}\log2\cdot2^{-k/2}(f(2^k)+f(2^{-k}))$.
\emph{(i) Relative form.}
\[
\operatorname{Tr}\mathcal U(f)\Pi_S-\mathrm{Hal}_S(f)=\operatorname{Tr}\mathcal U(f)\Pi_\infty-\mathrm{Hal}_\infty(f)-\mathrm{comb}_2(f).
\]
\emph{(ii) Absolute form at $\lambda=1$.} For $\lambda=1$,
\[
\operatorname{Tr}\mathcal U(f)\Pi_S=W_\infty(f)-\mathrm{comb}_2(f)+\int f(\rho^{-1})\,d_S(\rho)\,d^*\rho .
\]
\end{theorem}
\begin{proof}
For such $f$, $\mathcal U(f)$ commutes with $\mathsf F\in\{\mathsf F_\infty,\mathsf F_S\}$ (Proposition~\ref{prop:sl-AS}(i)). Expanding
$\langle\xi_n,\mathsf F\mathcal U(f)\mathsf F\xi_n\rangle$ with $\mathsf F\xi_n=\lambda_n\xi_n+(1-\lambda_n^2)^{1/2}\zeta_n$ gives, for each $n$,
$\langle\zeta_n,\mathcal U(f)\zeta_n\rangle=\langle\xi_n,\mathcal U(f)\xi_n\rangle-\tau_n(c_n+c_n')$ with $c_n,c_n'$ the two cross terms
(\cite[(87)]{ConnesConsani2021}; exact, no summation). The diagonal series in $\xi$, the two cross series and the
diagonal series in $\zeta$ converge absolutely by Lemma~\ref{lem:sl-TC}, so by Theorem~\ref{thm:sl-gap}
$\operatorname{Tr}\mathcal U(f)Q_1Q_2Q_1-\operatorname{Tr}\mathcal U(f)\Pi_{\mathsf F}$ can be computed term by term; the diagonal terms $\lambda_n^2$ cancel
$n$ by $n$, which is the semilocal form of \cite[(89)--(90)]{ConnesConsani2021}, and what remains is expressed through
$\mathrm{Hal}_{\mathsf F}(f)$ and the trace of $\mathcal U(f)(1-P)$ restricted to the generic part. Subtracting the two cases
$\mathsf F=\mathsf F_S$ and $\mathsf F=\mathsf F_\infty$, the difference $Q_1Q_2^SQ_1-Q_1Q_2^\infty Q_1=-(1-P)(\widehat P^S-\widehat P_\infty)(1-P)$
has $\mathcal U(f)$-trace equal to $-\mathrm{comb}_2(f)$ by Lemma~\ref{lem:sl-linear} and the block computation of
Lemma~\ref{lem:sl-TC}. This is (i). (ii) At $\lambda=1$, \cite[Thm.~4.7, (83)--(84)]{ConnesConsani2021} gives
$\operatorname{Tr}\mathcal U(f)\Pi_\infty=W_\infty(f)+\int f(\rho^{-1})\varepsilon_\infty(\rho)d^*\rho$ with $\varepsilon_\infty$ the density of the
archimedean Halmos series, i.e.\ $\int f(\rho^{-1})\varepsilon_\infty d^*\rho=\mathrm{Hal}_\infty(f)$. Insert this into (i), and use
the density $d_S$ of $\mathrm{Hal}_S$ (Theorem~\ref{thm:sl-L2}, first sentence).
\end{proof}

\begin{remark}[the level $\lambda$]\label{rem:sl-level}
\cite[Thm.~4.7]{ConnesConsani2021} is stated for Sonin's space with cut-offs $|q|\le1$, $|p|\le1$, which under (D1) is
the level $\lambda=1$; the space at level $\lambda$ is not a dilate of it unless $\lambda=1$. The absolute form (ii) is therefore
asserted at $\lambda=1$ only; all other statements of this section hold for every $\lambda$.
\end{remark}

\subsection{The split of the truncated form}\label{subsec:sl-split}
{\sloppy For $T>0$ let $\Pi^{(2)}_T(F):=-\sum_{k:\,2^k\in\NT}v_{2^k}\Rel\cF(k\log2)$ and
$\Pi^{(\ne2)}_T:=\PiT-\Pi^{(2)}_T$, so that
$\mathfrak a_T=\mT+\Pi^{(2)}_T+\Pi^{(\ne2)}_T$, and $\Pi^{(\ne2)}_T=0$ exactly when $T\le\frac12\log3$. With (D4),
$\mathrm{comb}_2(f_F)=-\Pi^{(2)}_T(F)$, since $v_{2^k}=2\log2\cdot2^{-k/2}$ and $\cF(-x)=\overline{\cF(x)}$.\par}

\begin{corollary}[semilocal Sonin split]\label{cor:sl-split}
Let $\lambda=1$, $S=\{\infty,2\}$, $\varepsilon_S:=d_S$, and let $C_{\varepsilon_S}$ be the operator on $L^2(\IT)$ with kernel
$\varepsilon_S(e^{|x-y|})$. For every $T>0$ and $F\in C_c^\infty(\IT)$,
\[
\mT(F)+\Pi^{(2)}_T(F)=\mathrm{Son}_{\Pi_S}(F)-\langle C_{\varepsilon_S}F,F\rangle .
\]
In particular, for $T\le\frac12\log3$, $\mathfrak a_T=\mathrm{Son}_{\Pi_S}-C_{\varepsilon_S}$ on $C_c^\infty(\IT)$.
Moreover $C_{\varepsilon_S}$ is Hilbert--Schmidt, and for every $\delta>0$,
$n_-\big(\mathfrak a_T-\mathrm{Son}_{\Pi_S}+\delta\|\cdot\|^2\,\big|\,\PT\big)<\infty$ when $T\le\frac12\log3$.
\end{corollary}
\begin{proof}
The identity is Theorem~\ref{thm:sl-84}(ii) with $f=f_F$, (D4) and (D5). By Theorem~\ref{thm:sl-L2},
$\int_{\IT^2}|\varepsilon_S(e^{|x-y|})|^2\le2T\int_{-2T}^{2T}|\varepsilon_S(e^{|x|})|^2dx<\infty$ by Theorem~\ref{thm:sl-L2}
($\log^+\frac1{|e^x-2^j|}\asymp\log^+\frac1{|x-j\log2|}$ near $x=j\log2$). A compact self-adjoint operator has finitely many eigenvalues below $-\delta$.
\end{proof}
Compactness gives only the $\delta$-version: a compact self-adjoint operator may have infinitely many negative
eigenvalues, and $n_-(-C_{\varepsilon_S}|\PT)<\infty$ is not claimed. (That $n_-(\mathfrak a_T|\PT)<\infty$ is immediate from
Theorem~\ref{thm:inertia-finite} with $\eta=1$, $k=0$.) For $T>\frac12\log3$ the identity of the corollary is a
statement about $\mT+\Pi^{(2)}_T$, not about $\mathfrak a_T$.

\begin{remark}[what this does not give]\label{rem:sl-not}
(1) \emph{No sign.} For $T\le\frac12\log3$ the corollary writes $\mathfrak a_T$ as a positive trace form minus a compact form; positivity of $\mathfrak a_T$ on $\PT$ is the statement that $\mathrm{Son}_{\Pi_S}$ dominates
$C_{\varepsilon_S}$ there, which nothing here decides.
(2) \emph{Comparison with Theorem~\ref{thm:inertia-nogo}.} For $\frac12\log2<T\le\frac12\log3$,
$\mathfrak a_T-\mathrm{Son}_{\Pi_\infty}$ has infinite negative index with eigenvalues bounded away from $0$, while
$\mathfrak a_T-\mathrm{Son}_{\Pi_S}$ is compact: the contact $2$ is absorbed by the semilocal Sonin
space. For $T>\frac12\log3$ the obstruction returns with $3$: $\mathfrak a_T-\mathrm{Son}_{\Pi_S}=-C_{\varepsilon_S}+\Pi^{(\ne2)}_T$, and
the test pairs of the proof of Theorem~\ref{thm:inertia-nogo} placed at distance $\log3$ form a sequence tending weakly
to $0$, so a compact $C_{\varepsilon_S}$ contributes $o(\|x\|^2)$ on them while $\Pi^{(\ne2)}_T(x)\le-c\|x\|^2$; the negative index is
infinite. Each new contact requires enlarging $S$; Lemma~\ref{lem:sl-linear} holds for every finite $S$, but
Proposition~\ref{prop:sl-AS}(iii) and Theorem~\ref{thm:sl-L2} are proved only for $S=\{\infty,2\}$, and the case of a
set $S(T)$ growing with $T$ is open.
(3) There is no uniform-in-$T$ control of $\|C_{\varepsilon_S}\|$; the coefficients $\beta_j$ grow like $2^{j/2}$.
(4) Theorem~\ref{thm:sl-gap} is a spectral gap at $1$ for $Q_1Q_2Q_1$; a positivity argument through $\Pi_S$ would
need a gap of a different kind, which is not addressed here.
\end{remark}

\begin{remark}[attribution and status]\label{rem:sl-new}
From \cite{CCquasi2020,CCM2023,CCM2024}: $\mathsf F_S$, the quasi-inner function and the factor $\rho_2$
(Lemma~\ref{lem:sl-phase}), compactness of the Hankel block, and the semilocal Sonin space. From
\cite{ConnesConsani2021}: the archimedean identities (83)--(90) whose algebra Theorem~\ref{thm:sl-84} transports.
New here: Proposition~\ref{prop:sl-AS}(iii), Theorem~\ref{thm:sl-gap} as stated, Lemma~\ref{lem:sl-TC},
Lemma~\ref{lem:sl-HG}, Theorem~\ref{thm:sl-L2} and Theorem~\ref{thm:sl-84}; Lemma~\ref{lem:sl-linear} is not new in
substance. Status: Theorem~\ref{thm:sl-84}(ii) rests on \cite[Thm.~4.7]{ConnesConsani2021} at $\lambda=1$; everything
else, including Lemma~\ref{lem:sl-HG}, Theorem~\ref{thm:sl-L2} and Corollary~\ref{cor:sl-split}, is proved
unconditionally.
\end{remark}


\section{The complete form at a single window}\label{sec:complete}

\begin{lemma}[non-attainment for the complete form]\label{lem:Onat}
Let $T\ge\frac12\log2$ and $w\in\Wcl$. Then $\inf_\R f_w<\lamW$.
\end{lemma}
\begin{proof}
Weak duality gives $c:=\inf f_w\le\lamW$ (the paragraph after Definition~\ref{def:W}). Suppose $c=\lamW$. The
bottom of $W_T$ on the unit sphere of $L^2(\IT)$ is attained at some $F_\star$ (\cite[Thm.~3]{Bombieri2000},
or the compactness argument in the proof of Theorem~\ref{thm:O}, which uses only that the symbol tends to
$+\infty$). Then $\varrho:=\sT+\nuh_w+2\gh0-\lamW\ge0$ is continuous, and
$\frac1{2\pi}\int\varrho|\Fh_\star|^2=W_T(F_\star)-\lamW=0$ by Proposition~\ref{prop:P}, since
$\frac1{2\pi}\int\nuh_w|\Fh_\star|^2=\int c_{F_\star}d\nu_w=0$. As $\Fh_\star$ is entire and nonzero,
$\varrho\equiv0$ on $\R$. Taking inverse Fourier transforms on $0<|u|<2T$, where $\nu_w=0$ and $g_0=\cosh\frac u2$,
and where the constant contributes nothing: $2\cosh\frac u2=\frac12K(|u|)+\sum_{n\in\NT}\frac{v_n}2(\delta_{\log n}+\delta_{-\log n})$
there, impossible since $K(u)\sim1/u$ as $u\downarrow0$ while $\cosh$ is bounded.
\end{proof}

\begin{theorem}[flat-top certificates and the complete bottom]\label{thm:Al}
\emph{(a)} If $f_w\ge0$ on $\R$ for some $w\in\Wcl$, regular or not, then $W_T\ge0$ on the form domain;
hence $\lamW\ge0$ and $\lamA\ge0$.
\emph{(b)} Let $T\ge\frac12\log2$. Then there is $w\in\Wcl$ with $f_w\ge0$ on $\R$ if and only if
$\lamW>0$, and $w$ can then be taken in $\Wcl^{\mathrm{reg}}$ with $\inf f_w>0$. Likewise, if $\lamP>0$
there are $w\in\Wcl^{\mathrm{reg}}$ and $\beta\in\R^2$ with $\inf_\tau[f_w+(\beta_0-2)\gh0+\beta_1\gh1]>0$.
\end{theorem}
\begin{proof}
(a) $(\nu_w,2,0)$ is dual feasible and $W_T\ge\inf f_w$ on $L^2(\IT)$ (the paragraph after
Definition~\ref{def:W}), and $W_T=\mathfrak a_T$ on $\PT$. (b) ``$\Leftarrow$'' and the regular choice:
Theorem~\ref{thm:Zexh}(i). ``$\Rightarrow$'': $\lamW\ge\inf f_w\ge0$ by weak duality, and $\lamW>\inf f_w$ by
Lemma~\ref{lem:Onat}. The last sentence is Theorem~\ref{thm:Zexh}(ii).
\end{proof}

% (before the paragraph "\emph{The per-window converse.}")
\begin{definition}[the alignment inequality]\label{def:Al}
For $w\in\Wcl$ let $\mathrm{Al}_T(w)$ be the statement: for every $\tau\in\R$,
\[
-2\Rel\sum_{n\ge2}\frac{\Lambda(n)w(\log n)}{n^{\frac12+i\tau}}\ \ge\
-\Big[C_0+\int_0^\infty\big(1-w(u)\cos\tau u\big)K(u)\,du\Big]-2\int_\R w(u)\cosh\tfrac u2\cos(\tau u)\,du .
\]
The left side is the smoothed prime sum on the critical line, $-\sum_nv_nw(\log n)\cos(\tau\log n)$; by
\eqref{eq:fw}, $\mathrm{Al}_T(w)$ is exactly the statement $f_w\ge0$ on $\R$ of Theorem~\ref{thm:Al}.
\end{definition}

\begin{lemma}[the alignment inequality lives on a band]\label{lem:Alband}
For $w\in\Wcl$ let $\tau_{\mathrm{env}}(w)$ be the unique positive root of
$m_{\mathrm{lo}}(\tau)-\Aconst-\|\nu_w\|_{\TV}-2k_0(\tau)=0$, with $m_{\mathrm{lo}}$ as in Lemma~\ref{lem:mbounds}(ii)
and $k_0$ as in Lemma~\ref{lem:gbounds}. Then $f_w>0$ on $(\tau_{\mathrm{env}}(w),\infty)$, so $\mathrm{Al}_T(w)$ is
equivalent to its restriction to $|\tau|\le\tau_{\mathrm{env}}(w)$.
\end{lemma}
\begin{proof}
$m_{\mathrm{lo}}$ increases from $-\infty$ to $+\infty$ and $k_0$ decreases (Lemmas~\ref{lem:mbounds}, \ref{lem:gbounds}),
so the left side is strictly increasing and the root is unique. Proposition~\ref{prop:env}(i) for the feasible
triple $(\nu_w,2,0)$ gives $f_w\ge m_{\mathrm{lo}}-\Aconst-\|\nu_w\|_{\TV}-2k_0$, and $f_w$ is even.
\end{proof}

\begin{remark}[the per-window quantity]\label{rem:excess}
By Theorem~\ref{thm:Zexh}(i), $\sup_{w\in\Wcl^{\mathrm{reg}}}\inf_\tau f_w=\lamW$ for $T\ge\frac12\log2$: the largest
uniform margin with which the alignment inequality can hold is $\lamW$ itself, a number determined by
$v_T(\Lambda)$ and $m$ (Proposition~\ref{prop:finite}). The trivial envelope $-\sum_nv_n|w(\log n)|$ of the prime
sum certifies $\mathrm{Al}_T(w)$ wherever the right side of Definition~\ref{def:Al} is at most it, and by
Proposition~\ref{prop:env}(ii) the constant $\Aconst$ of that envelope cannot be lowered; on the complement the
actual phases $\tau\log p$ modulo $2\pi$ enter, and in the two regimes of Proposition~\ref{prop:align} the
negative set of the symbol itself lies at the near-solutions of $2^l\approx3^k$. For $w\in\Wcl^{\mathrm{reg}}$, by
Theorem~\ref{thm:Z}(ii), $\mathrm{Al}_T(w)$ at $\tau$ is the inequality $\sum_\rho\wh(\tau-\gamma_\rho)\ge0$.
\end{remark}

\begin{remark}[the prime side of a primitive test is the Chebyshev remainder]\label{rem:Cheb}
Let $\psi_\Lambda(x)=\sum_{n\le x}\Lambda(n)$ and $r_\Lambda(u):=\psi_\Lambda(e^u)-e^u$. For real $F\in\PT\cap\mathcal D_T$,
\begin{equation}\label{eq:Cheb}
\mathfrak a_T(F)=C_0\|F\|^2+\Lev(\cF)+2\int_0^{2T}\cF(u)e^{-u/2}\,du-2\int_{(0,2T)}e^{-u/2}\cF(u)\,dr_\Lambda(u).
\end{equation}
Indeed $\sum_{n\in\NT}v_n\cF(\log n)=2\int_{(0,2T)}e^{-u/2}\cF(u)\,d\psi_\Lambda(e^u)$, since $u\mapsto\psi_\Lambda(e^u)$
jumps by $\Lambda(n)$ at $u=\log n$ and $\cF(2T)=0$; split $d\psi_\Lambda(e^u)=e^udu+dr_\Lambda(u)$. The identity
$\int_{-2T}^{2T}\cF(u)e^{u/2}du=A(F)B(F)=0$ on $\PT$ (substitute $s=t-u$) and the evenness of $\cF$ give
$\int_0^{2T}\cF e^{u/2}du=-\int_0^{2T}\cF e^{-u/2}du$; insert into \eqref{eq:Aform}. So on primitive tests the whole
arithmetic of the comb is the last term of \eqref{eq:Cheb}, and the unconditional control of $r_\Lambda$ is the
theory of zero-free regions. If $\lamA\ge0$, then for every $\eta>\lamA$ there are tests $F\in\PT$ with
$0\le\mathfrak a_T(F)\le\eta\|F\|^2$, and $\lamA$ is at most $4\cdot10^9Te^{-15.449T}$ for $T\ge1$
(Corollary~\ref{cor:margin1}); for such $F$, \eqref{eq:Cheb} says that $2\int e^{-u/2}\cF\,dr_\Lambda$ equals
$C_0\|F\|^2+\Lev(\cF)+2\int_0^{2T}\cF e^{-u/2}du$ up to $\eta\|F\|^2$.
\end{remark}


\emph{The per-window converse.} At a single window, the converse of Theorem~\ref{thm:Al}(a) is the
implication $\mathrm{WP}_T\Rightarrow\lamW\ge0$. It reduces to the signs of two numbers. Fix $T$. We work
with real functions. This suffices for bottoms and negative indices, as in the remark after
Lemma~\ref{lem:parity}: for $F=u+iv$ with $u,v$ real, $W_T(F)=W_T(u)+W_T(v)$ and
$\mathfrak a_T(F)=\mathfrak a_T(u)+\mathfrak a_T(v)$. Let $\mathcal D_T$ be the form domain of \eqref{eq:form}
and $\mathcal D^{\mathrm{par}}_T$ its elements of parity $\mathrm{par}\in\{\mathrm{ev},\mathrm{odd}\}$. With
$\chi_{\mathrm{par}}$ as in \S\ref{subsec:conv}, put
$\ell_{\mathrm{par}}(F)=\langle F,\chi_{\mathrm{par}}\rangle$, $\varepsilon_{\mathrm{ev}}=+1$ and
$\varepsilon_{\mathrm{odd}}=-1$. Then $A=\ell_{\mathrm{ev}}+\ell_{\mathrm{odd}}$ and
$B=\ell_{\mathrm{ev}}-\ell_{\mathrm{odd}}$ on real $F$ (the parity paragraph after
Proposition~\ref{prop:nodes}). Let $\mathcal Q^{\mathrm{par}}_T:=\PT^{\mathrm{par}}\cap\mathcal D_T
=\ker(\ell_{\mathrm{par}}|\mathcal D^{\mathrm{par}}_T)$, and write $\lambda^{\mathrm{par}}_{\mathrm A}(T)$ for the
bottom of $\mathfrak a_T$ on its unit sphere. This is the bottom with the primes, not the prime-free
$\lambda_{\mathrm{par}}$ of \S\ref{subsec:arch}. Write $\lambda^{\mathrm W}_{T,\mathrm{par}}$ for the bottom of
$W_T$ on the unit sphere of $\mathcal D^{\mathrm{par}}_T$. As in Lemma~\ref{lem:parity}, $n_-$ is the largest
dimension of a subspace of the form domain on which the form is negative definite. The letters $h$,
$\mathfrak g$, $\mathfrak G$, $S^{\mathrm{Sch}}$, $\mathcal N$ are local to this section.

\begin{theorem}[the parity Schur complement of the complete form]\label{thm:parityschur}
Fix $T>0$.
\emph{(o) Parity splitting.} On real $F=F_{\mathrm e}+F_{\mathrm o}\in\mathcal D_T$,
\[
W_T(F)=\big[\mathfrak a_T(F_{\mathrm e})+2\ell_{\mathrm{ev}}(F_{\mathrm e})^2\big]
+\big[\mathfrak a_T(F_{\mathrm o})-2\ell_{\mathrm{odd}}(F_{\mathrm o})^2\big],
\]
the sum being orthogonal for $L^2$, for $\mathfrak a_T$ and for $W_T$. Hence
$\lamA=\min_{\mathrm{par}}\lambda^{\mathrm{par}}_{\mathrm A}(T)$ and $\lamW=\min_{\mathrm{par}}\lambda^{\mathrm W}_{T,\mathrm{par}}$.

\emph{(a) Non-degenerate sector.} Suppose $0$ is not in the spectrum of the self-adjoint operator of the
closed form $\mathfrak a_T|\mathcal Q^{\mathrm{par}}_T$. This holds in particular if
$\lambda^{\mathrm{par}}_{\mathrm A}(T)>0$. Then there is a unique $h_{\mathrm{par}}\in\mathcal D^{\mathrm{par}}_T$
with $\ell_{\mathrm{par}}(h_{\mathrm{par}})=1$ and $\mathfrak a_T(h_{\mathrm{par}},G)=0$ for all
$G\in\mathcal Q^{\mathrm{par}}_T$. Put $\mathfrak g_{\mathrm{par}}(T):=W_T(h_{\mathrm{par}})=\mathfrak a_T(h_{\mathrm{par}})+2\varepsilon_{\mathrm{par}}$.
Then:
\begin{itemize}
\item[(a1)] $\mathcal D^{\mathrm{par}}_T=\mathcal Q^{\mathrm{par}}_T\oplus\R h_{\mathrm{par}}$, orthogonally for
$\mathfrak a_T$ and for $W_T$, and $W_T(G+xh_{\mathrm{par}})=\mathfrak a_T(G)+\mathfrak g_{\mathrm{par}}x^2$ for
$G\in\mathcal Q^{\mathrm{par}}_T$, $x\in\R$;
\item[(a2)] $n_-(W_T|\mathcal D^{\mathrm{par}}_T)=n_-(\mathfrak a_T|\mathcal Q^{\mathrm{par}}_T)+[\mathfrak g_{\mathrm{par}}<0]$.
The radical of $W_T$ on $\mathcal D^{\mathrm{par}}_T$ is $\R h_{\mathrm{par}}$ if $\mathfrak g_{\mathrm{par}}=0$ and $0$
otherwise;
\item[(a3)] if $\lambda^{\mathrm{par}}_{\mathrm A}(T)>0$, then
$\mathfrak a_T(h_{\mathrm{par}})=\min\{\mathfrak a_T(F):F\in\mathcal D^{\mathrm{par}}_T,\ \ell_{\mathrm{par}}(F)=1\}$ and
$\mathfrak g_{\mathrm{par}}=\min\{W_T(F):F\in\mathcal D^{\mathrm{par}}_T,\ \ell_{\mathrm{par}}(F)=1\}$;
\item[(a4)] $\lambda^{\mathrm W}_{T,\mathrm{par}}\le\mathfrak g_{\mathrm{par}}/\|h_{\mathrm{par}}\|^2$ and
$\|h_{\mathrm{par}}\|\ge1/\|\chi_{\mathrm{par}}\|$. If $\lambda^{\mathrm{par}}_{\mathrm A}(T)>0$ \emph{and}
$\mathfrak g_{\mathrm{par}}>0$, then
$\lambda^{\mathrm W}_{T,\mathrm{par}}\ge\frac12\min\big(\lambda^{\mathrm{par}}_{\mathrm A}(T),\mathfrak g_{\mathrm{par}}/\|h_{\mathrm{par}}\|^2\big)$.
If in addition $2\lambda^{\mathrm W}_{T,\mathrm{par}}<\lambda^{\mathrm{par}}_{\mathrm A}(T)$, then
$\lambda^{\mathrm W}_{T,\mathrm{par}}\le\mathfrak g_{\mathrm{par}}/\|h_{\mathrm{par}}\|^2\le2\lambda^{\mathrm W}_{T,\mathrm{par}}$.
\end{itemize}

\emph{(b) Degenerate sector.} Suppose $\lambda^{\mathrm{par}}_{\mathrm A}(T)=0$, and let
$\mathcal N=\ker(\mathfrak a_T|\mathcal Q^{\mathrm{par}}_T)$. It is nonzero and finite-dimensional. For each
$F_0\in\mathcal N$ there is a unique $\alpha(F_0)\in\R$ with $\mathfrak a_T(F_0,G)=\alpha(F_0)\ell_{\mathrm{par}}(G)$
for all $G\in\mathcal D^{\mathrm{par}}_T$. If $\alpha\not\equiv0$ on $\mathcal N$, then $W_T$ is unbounded below on
$\mathcal D^{\mathrm{par}}_T$, so $\lamW<0$. If $\alpha\equiv0$ on $\mathcal N$ (the \emph{range condition}), then
$\mathcal N$ lies in the radical of $W_T$ on $\mathcal D^{\mathrm{par}}_T$ and (a) applies on the $L^2$-orthogonal
complement of $\mathcal N$ in $\mathcal D^{\mathrm{par}}_T$. With the $h_{\mathrm{par}},\mathfrak g_{\mathrm{par}}$ so obtained,
$W_T\ge0$ on $\mathcal D^{\mathrm{par}}_T$ if and only if $\mathfrak g_{\mathrm{par}}\ge0$.

\emph{(c) Assembled.} Put $\mathfrak G_T:=\mathrm{diag}(\mathfrak g_{\mathrm{ev}}(T),\mathfrak g_{\mathrm{odd}}(T))$. It is
defined whenever each sector is as in (a), or as in (b) with the range condition. If $0$ is not in the
spectrum of the operator of $\mathfrak a_T|(\PT\cap\mathcal D_T)$, then
$n_-(W_T)=n_-(\mathfrak a_T|\PT)+n_-(\mathfrak G_T)$, and the radical of $W_T$ has dimension $n_0(\mathfrak G_T)$. At every window,
without further hypotheses,
\[
\lamW>0\iff\big[\lamA>0\ \text{and}\ \mathfrak G_T\succ0\big],
\]
\[
\lamW\ge0\iff\big[\mathrm{WP}_T,\ \text{the range condition in every sector with }
\lambda^{\mathrm{par}}_{\mathrm A}(T)=0,\ \text{and}\ \mathfrak G_T\succeq0\big].
\]
In particular, if $\lamA>0$, then $\lamW\ge0$ if and only if $\mathfrak g_{\mathrm{ev}}(T)\ge0$ and
$\mathfrak g_{\mathrm{odd}}(T)\ge0$.
\end{theorem}
\begin{proof}
(o) The symbol $\sT$ is even. For real even $F_{\mathrm e}$ and real odd $F_{\mathrm o}$, $\widehat{F_{\mathrm e}}$ is even and
$\widehat{F_{\mathrm o}}$ is odd, so $\mathfrak a_T(F_{\mathrm e},F_{\mathrm o})=0$ as in the proof of
Lemma~\ref{lem:parity}. Also $\ell_{\mathrm{ev}}$ vanishes on odd functions and $\ell_{\mathrm{odd}}$ on even ones, so
$A(F)B(F)=(\ell_{\mathrm{ev}}(F_{\mathrm e})+\ell_{\mathrm{odd}}(F_{\mathrm o}))(\ell_{\mathrm{ev}}(F_{\mathrm e})-\ell_{\mathrm{odd}}(F_{\mathrm o}))
=\ell_{\mathrm{ev}}(F_{\mathrm e})^2-\ell_{\mathrm{odd}}(F_{\mathrm o})^2$, with no cross term. Polarizing gives the
orthogonality. $F\in\PT$ if and only if $\ell_{\mathrm{ev}}(F_{\mathrm e})=\ell_{\mathrm{odd}}(F_{\mathrm o})=0$, and bottoms over an
orthogonal sum take the minimum.

(a) Write $\mathcal Q=\mathcal Q^{\mathrm{par}}_T$, $\ell=\ell_{\mathrm{par}}$ and $\varepsilon=\varepsilon_{\mathrm{par}}$. On
$\mathcal D^{\mathrm{par}}_T$, $W_T(F,G)=\mathfrak a_T(F,G)+2\varepsilon\,\ell(F)\ell(G)$ by (o).

\emph{Solvability.} First, $\mathfrak a_T=\mT+\PiT$ with $\|\PiT\|\le\Aconst$, since $|\cF(\log n)|\le\|F\|^2$ for each $n$ (this is Proposition~\ref{prop:coer-prime}(i) below). The form
$\mT\ge C_0\|\cdot\|^2$ is closed, with domain $\mathcal D_T$ (Theorem~\ref{thm:A}), and this domain embeds
compactly in $L^2(\IT)$ (proof of Proposition~\ref{prop:nodes}(i)). So $\mathfrak a_T$ is closed and bounded
below by $-M$, where $M:=\Aconst-C_0$. Second, $\mathcal Q$ is closed in the form norm, because $\ell$ and the
parity projection are $L^2$-continuous. Third, $\mathcal Q$ is dense in
$H:=\ker\ell\cap L^2_{\mathrm{par}}(\IT)$. Indeed, let $x\in H$, let $d_n\in\mathcal D^{\mathrm{par}}_T$ with $d_n\to x$ in
$L^2$ (smooth functions of parity $\mathrm{par}$ are dense), and let $e\in\mathcal D^{\mathrm{par}}_T$ with $\ell(e)=1$.
Then $d_n-\ell(d_n)e\in\mathcal Q$ tends to $x$. Hence $\mathfrak a_T|\mathcal Q$ is a closed, semibounded, densely
defined form in $H$ with compact resolvent. Its spectrum is discrete, and its negative spectral subspace
$\mathcal Q_-$ is finite-dimensional and $\mathfrak a_T$-orthogonal to $\mathcal Q_+:=\mathcal Q\cap\mathcal Q_-^\perp$. If
$0$ is not in the spectrum, then $\mathfrak a_T\ge\delta\|\cdot\|^2$ on $\mathcal Q_+$ for some $\delta>0$. Hence
$\mathfrak a_T(x)\ge\frac{\delta}{\delta+M+1}\big(\mathfrak a_T(x)+(M+1)\|x\|^2\big)$ there, so $\mathfrak a_T$ is coercive
on $\mathcal Q_+$ in the form norm and the Lax--Milgram lemma applies on $\mathcal Q_+$. On $\mathcal Q_-$ the form is a
negative definite finite matrix. So for every form-bounded linear functional $f$ on $\mathcal Q$ there is a unique
$Y\in\mathcal Q$ with $\mathfrak a_T(Y,\cdot)=f$ on $\mathcal Q$. Uniqueness holds because the operator has trivial
kernel. Now fix $H_0\in\mathcal D^{\mathrm{par}}_T$ with $\ell(H_0)=1$, for instance a smooth bump of parity
$\mathrm{par}$. The functional $f=\mathfrak a_T(H_0,\cdot)$ is form-bounded, by Cauchy--Schwarz for the nonnegative
form $\mathfrak a_T+M\|\cdot\|^2$. Put $h:=H_0-Y$. Then $\ell(h)=1$ and $\mathfrak a_T(h,\cdot)=0$ on $\mathcal Q$. Two such
$h$ differ by an element of $\mathcal Q$ that is $\mathfrak a_T$-orthogonal to $\mathcal Q$, which is $0$.

(a1) Any $F\in\mathcal D^{\mathrm{par}}_T$ decomposes as $F=(F-\ell(F)h)+\ell(F)h$ with $F-\ell(F)h\in\mathcal Q$. For
$G\in\mathcal Q$, $W_T(h,G)=\mathfrak a_T(h,G)+2\varepsilon\ell(h)\ell(G)=0$. Expand, using $\ell(G)=0$ and $\ell(h)=1$.

(a2) By (a1), $\mathcal D^{\mathrm{par}}_T=(\mathcal Q_-\oplus[\R h\text{ if }\mathfrak g<0])\oplus(\mathcal Q_+\oplus[\R h\text{ if }\mathfrak g\ge0])$,
orthogonally for $W_T$. The first summand is negative definite and has the stated dimension. The second
summand carries $W_T\ge0$. A negative definite subspace meets the second summand only in $0$, so its
projection onto the first summand along the second is injective. For the radical: if
$W_T(G+xh,\cdot)=0$, then $\mathfrak a_T(G,\cdot)=0$ on $\mathcal Q$, so $G=0$, and then $\mathfrak g x=0$.

(a3) If the bottom is positive, then $\mathfrak a_T(G+h)=\mathfrak a_T(G)+\mathfrak a_T(h)\ge\mathfrak a_T(h)$ and
$W_T(G+h)=\mathfrak a_T(G)+\mathfrak g\ge\mathfrak g$.

(a4) The upper bound is the Rayleigh quotient at $h$, and $1=\ell(h)\le\|h\|\|\chi_{\mathrm{par}}\|$. Put
$\lambda:=\lambda^{\mathrm{par}}_{\mathrm A}(T)$ and assume $\lambda>0$ and $\mathfrak g>0$. With
$\mu:=\min(\lambda,\mathfrak g/\|h\|^2)$,
\[
W_T(G+xh)=\mathfrak a_T(G)+\mathfrak gx^2\ge\lambda\|G\|^2+\frac{\mathfrak g}{\|h\|^2}\|xh\|^2
\ge\mu\big(\|G\|^2+\|xh\|^2\big)\ge\tfrac12\mu\|G+xh\|^2 .
\]
If also $2\lambda^{\mathrm W}_{T,\mathrm{par}}<\lambda$, then $\frac12\mu\le\lambda^{\mathrm W}_{T,\mathrm{par}}<\frac12\lambda$ forces
$\mu=\mathfrak g/\|h\|^2$. The hypothesis $\mathfrak g>0$ is used in the second inequality and cannot be dropped.

(b) The bottom $0$ is attained (compact resolvent), so $\mathcal N\ne0$, and $\mathcal N$ is finite-dimensional. For
$F_0\in\mathcal N$ and $G\in\mathcal Q$, Cauchy--Schwarz for the nonnegative form $\mathfrak a_T|\mathcal Q$ gives
$\mathfrak a_T(F_0,G)=0$. So $\mathfrak a_T(F_0,\cdot)$ is a linear functional on $\mathcal D^{\mathrm{par}}_T$ vanishing on
$\ker\ell$, hence a multiple $\alpha(F_0)\ell$. Since $\ell(F_0)=0$, we get $W_T(F_0,H)=\alpha(F_0)\ell(H)$ and
$W_T(F_0)=0$. So $W_T(xF_0+H)=W_T(H)+2x\alpha(F_0)\ell(H)$, which is unbounded below in $x$ if
$\alpha(F_0)\ne0$ and $\ell(H)=1$. Now suppose $\alpha\equiv0$. Then $\mathcal N$ is $W_T$-orthogonal to
$\mathcal D^{\mathrm{par}}_T$. Put $\mathcal D':=\mathcal D^{\mathrm{par}}_T\cap\mathcal N^\perp$ and $\mathcal Q':=\mathcal Q\cap\mathcal N^\perp$.
Then $\mathcal Q'=\ker(\ell|\mathcal D')$, and $\ell\ne0$ on $\mathcal D'$: subtract from $H_0$ its $L^2$-projection onto
$\mathcal N\subset\mathcal Q$, which $\ell$ annihilates. On $\mathcal Q'$ the form is bounded below by the next eigenvalue,
which is $>0$. So (a) applies with $\mathcal D',\mathcal Q'$ in place of $\mathcal D^{\mathrm{par}}_T,\mathcal Q$. Moreover
$\mathcal D^{\mathrm{par}}_T=\mathcal N\oplus\mathcal D'$, $W_T$-orthogonally. Conclude by (a1).

(c) By (o), everything adds over the two sectors, and the inertia formula is (a2) in each sector. If
$\lamW>0$, then $\lamA\ge\lamW>0$ (recall $\lamW\le\lamA$), so both sectors are as in (a), and
$\mathfrak g_{\mathrm{par}}=W_T(h_{\mathrm{par}})\ge\lamW\|h_{\mathrm{par}}\|^2>0$. The converse is (a4) in each sector. If $\lamW\ge0$,
then $\lamA\ge\lamW\ge0$, which is $\mathrm{WP}_T$. In a sector with bottom $0$, (b) excludes $\alpha\not\equiv0$, and
in every sector $\mathfrak g_{\mathrm{par}}=W_T(h_{\mathrm{par}})\ge0$. Conversely, (a1) and $\mathfrak a_T\ge0$ on
$\mathcal Q^{\mathrm{par}}_T$ give $W_T\ge0$ on a sector with positive bottom when $\mathfrak g_{\mathrm{par}}\ge0$, and (b) does the
same on a sector with bottom $0$.
\end{proof}

The range condition says that the column $\rho(G):=\mathfrak a_T(H_0,G)$ of the bordered form, with
$\ell_{\mathrm{par}}(H_0)=1$, vanishes on $\ker(\mathfrak a_T|\mathcal Q^{\mathrm{par}}_T)$, since $\rho(F_0)=\alpha(F_0)$. This is the
range condition of a bordered Schur complement. Because $\mathfrak a_T|\mathcal Q^{\mathrm{par}}_T$ has a spectral gap above
$0$, it is equivalent to $\rho$ lying in the range of the square root of the operator. The theorem takes
the complement over $\PT$ and not over $L^2(\IT)$. So only $\mathfrak a_T|\PT$ has to be invertible, which is
the natural hypothesis for a converse whose premise is on the primitive side.

\begin{proposition}[$\mathfrak G_T$ is the polar capacity]\label{prop:polarcap}
Assume $\mathrm{WP}_T$, and the range condition of Theorem~\ref{thm:parityschur}(b) in every
sector with $\lambda^{\mathrm{par}}_{\mathrm A}(T)=0$. Let
$S^{\mathrm{Sch}}_T:=\mathrm{diag}\big(\mathfrak a_T(h_{\mathrm{ev}}),\mathfrak a_T(h_{\mathrm{odd}})\big)$. This is the compression of
$\mathfrak a_T$ to $\mathrm{span}(h_{\mathrm{ev}},h_{\mathrm{odd}})$, the $\mathfrak a_T$-orthogonal complement of the primitive part,
i.e.\ the Schur complement of $\mathfrak a_T$ over $\PT$. Let $D_{\mathrm{pol}}:=\mathrm{diag}(2,-2)$. Then
\[
\mathfrak G_T=S^{\mathrm{Sch}}_T+D_{\mathrm{pol}} .
\]
For $(\beta,\alpha)\in\R^2$ put
$E_T(\beta,\alpha):=\inf\{W_T(F):\ F\in\mathcal D_T\ \text{real},\ \Phi_F(0)=B(F)=\beta,\ \Phi_F(1)=A(F)=\alpha\}$. Then
\[
E_T(\beta,\alpha)=(\alpha,\beta)\,\mathfrak G^{AB}_T\,(\alpha,\beta)^{\mathsf T},\qquad
\mathfrak G^{AB}_T=\frac14\begin{pmatrix}\mathfrak g_{\mathrm{ev}}+\mathfrak g_{\mathrm{odd}}&\mathfrak g_{\mathrm{ev}}-\mathfrak g_{\mathrm{odd}}\\
\mathfrak g_{\mathrm{ev}}-\mathfrak g_{\mathrm{odd}}&\mathfrak g_{\mathrm{ev}}+\mathfrak g_{\mathrm{odd}}\end{pmatrix},
\]
and the infimum is attained. In these coordinates $D_{\mathrm{pol}}$ becomes $\big(\begin{smallmatrix}0&1\\1&0\end{smallmatrix}\big)$,
the form $2\alpha\beta=2\Phi_F(0)\Phi_F(1)$, with the one negative square of Lemma~\ref{lem:polar}. Hence
$\lamW\ge0$ if and only if $E_T\ge0$ on $\R^2$. Under $\mathrm{WP}_T$, the complete form is nonnegative
exactly when the least complete Weil energy needed to prescribe the values of $\Phi_F$ at the pole
$s\in\{0,1\}$ is nonnegative.
\end{proposition}
\begin{proof}
$\mathfrak G_T=S^{\mathrm{Sch}}_T+D_{\mathrm{pol}}$ is the definition $\mathfrak g_{\mathrm{par}}=\mathfrak a_T(h_{\mathrm{par}})+2\varepsilon_{\mathrm{par}}$. In the
degenerate case with the range condition, $\mathfrak a_T(h_{\mathrm{par}})$ does not depend on the choice of $h_{\mathrm{par}}$
modulo $\mathcal N$. For real $F$, $\Phi_F(1)=A(F)$ and $\Phi_F(0)=B(F)$ (proof of Lemma~\ref{lem:roughEF}). The
constraints are $\ell_{\mathrm{ev}}(F_{\mathrm e})=c:=\frac{\alpha+\beta}2$ and $\ell_{\mathrm{odd}}(F_{\mathrm o})=s:=\frac{\alpha-\beta}2$,
and they split over the sectors, as does $W_T$ by Theorem~\ref{thm:parityschur}(o). In each sector,
Theorem~\ref{thm:parityschur}(a1) (or (b)) gives $W_T(G+xh_{\mathrm{par}})=\mathfrak a_T(G)+\mathfrak g_{\mathrm{par}}x^2$ with $x$ the
prescribed value of $\ell_{\mathrm{par}}$ and $\mathfrak a_T(G)\ge0$ by $\mathrm{WP}_T$. So the sector minimum is
$\mathfrak g_{\mathrm{par}}x^2$, attained at $G=0$. Hence $E_T=\mathfrak g_{\mathrm{ev}}c^2+\mathfrak g_{\mathrm{odd}}s^2$, which is the displayed
matrix. The form $2c^2-2s^2=2\alpha\beta$ gives $D_{\mathrm{pol}}$ in the new coordinates. The last statement is
Theorem~\ref{thm:parityschur}(c), since $E_T\ge0$ if and only if $\mathfrak G_T\succeq0$.
\end{proof}

The converse is therefore the pair of inequalities $\mathfrak a_T(h_{\mathrm{odd}})\ge2$ and
$\mathfrak a_T(h_{\mathrm{ev}})\ge-2$. In the odd direction the primitive Schur complement must beat the negative
polar square. In the even direction the positive polar square may absorb a negative Schur value down
to $-2$. Neither is automatic; Corollary~\ref{cor:convfail} below shows that the pair is not implied by
$\mathrm{WP}_T$ at the level of combs.


\subsection{The optimal multipliers}\label{subsec:mult}
Remark~\ref{rem:Kprim} leaves open whether $\lamP=\lamW$, that is, whether the multipliers $\beta$ of
near-optimal certificates of the primitive problem tend to the polar value $(2,0)$ of
Proposition~\ref{prop:P}. The question has a variational answer. For $F\in L^2(\IT)$ put
$A'(F)=\int tF(t)e^{t/2}dt$, $B'(F)=\int tF(t)e^{-t/2}dt$ and
\[
\mathfrak m(F):=\Big(\frac1{2\pi}\int\gh0|\Fh|^2,\ \frac1{2\pi}\int\gh1|\Fh|^2\Big)
=\big(\Rel(A\bar B),\ \Rel(A'\bar B-A\bar B')\big),
\]
the two moments of Lemma~\ref{lem:W} (the second expression is the computation in the proof of
Theorem~\ref{thm:A}), and for $\beta\in\R^2$
\[
V_T(\beta):=\inf\Big\{\frac1{2\pi}\int(\sT+\beta\cdot\widehat g)|\Fh|^2\,d\tau:\ F\in\mathcal D_T,\ \|F\|=1\Big\}
=\inf_{F\in\mathcal D_T,\ \|F\|=1}\big[\mathfrak a_T(F)+\beta\cdot\mathfrak m(F)\big].
\]
By Proposition~\ref{prop:P}, $V_T(2,0)=\lamW$. For real $F=F_{\mathrm e}+F_{\mathrm o}$ the first moment is
$\Rel(A\bar B)=\ell_{\mathrm{ev}}(F_{\mathrm e})^2-\ell_{\mathrm{odd}}(F_{\mathrm o})^2$ (proof of
Theorem~\ref{thm:parityschur}(o)), and for $F=u+iv$ with $u,v$ real one has $\mathfrak m(F)=\mathfrak m(u)+\mathfrak m(v)$,
$W_T(F)=W_T(u)+W_T(v)$ and $\|F\|^2=\|u\|^2+\|v\|^2$, since $\gh0,\gh1,\sT$ are even and
$\Iml(\widehat u\,\overline{\widehat v})$ is odd.

\begin{theorem}[the optimal multipliers]\label{thm:mult}
Fix $T>0$.
\emph{(i)} $V_T$ is finite, concave and continuous on $\R^2$, $V_T(\beta)\to-\infty$ as $|\beta|\to\infty$, and
\[
\lamP=\max_{\beta\in\R^2}V_T(\beta).
\]
Let $\mathcal M_T\subset\R^2$ be the set of maximizers; it is nonempty, compact and convex.
\emph{(ii)} If $(\nu_k,\beta^{(k)})$ are certificates of values $c_k\to\lamP$, then
$\operatorname{dist}(\beta^{(k)},\mathcal M_T)\to0$. Hence every such sequence satisfies $\beta^{(k)}\to(2,0)$ if and
only if $\mathcal M_T=\{(2,0)\}$, and $\lamP=\lamW$ if and only if $(2,0)\in\mathcal M_T$.
\emph{(iii)} Let $\mathcal Z_T$ be the set of real unit vectors $F\in\mathcal D_T$ with $W_T(F)=\lamW$; it is nonempty
and compact in $L^2(\IT)$. Then $(2,0)\in\mathcal M_T$ if and only if
$0\in\operatorname{conv}\{\mathfrak m(F):F\in\mathcal Z_T\}$.
\emph{(iv)} Suppose $\lamW<\lamA$. Then every $F\in\mathcal Z_T$ of definite parity has
$\Rel(A\bar B)=\ell_{\mathrm{ev}}(F)^2>0$ if $F$ is even and $\Rel(A\bar B)=-\ell_{\mathrm{odd}}(F)^2<0$ if $F$ is odd.
Consequently, if moreover $\lambda^{\mathrm W}_{T,\mathrm{ev}}\ne\lambda^{\mathrm W}_{T,\mathrm{odd}}$, then $\lamW<\lamP$, and
the multipliers of every maximizing sequence of certificates stay at a positive distance from $(2,0)$.
\emph{(v)} If $\lamW=\lamA$, then $\lamP=\lamW$ and $(2,0)\in\mathcal M_T$.
\end{theorem}
\begin{proof}
(i) For a unit vector $F$, $|\Rel(A\bar B)|\le\|e^{t/2}\|\,\|e^{-t/2}\|=2\sinh T$ and
$|\Rel(A'\bar B-A\bar B')|\le2\|te^{t/2}\|\,\|e^{-t/2}\|\le4T\sinh T$, so $|\mathfrak m(F)|\le C_T:=2\sinh T\,(1+2T)$.
Hence $V_T(\beta)\ge C_0-\Aconst-C_T|\beta|>-\infty$, and $V_T(\beta)<\infty$ since $\mathcal D_T\ne0$. As an infimum
of affine functions of $\beta$, $V_T$ is concave; a finite concave function on $\R^2$ is continuous. For the
behaviour at infinity, Lemma~\ref{lem:W} gives $r>0$ and $V_0<\infty$ such that every point of the ball of radius
$r$ is the moment vector $\mathfrak m(F)$ of some unit $F\in\mathcal D_T$ with $\mathfrak a_T(F)\le V_0$ (the ball is covered by
convex combinations, at the level of the measures $|\Fh|^2d\tau/2\pi$, of four measures of finite value, and such a
measure is $|\Fh|^2d\tau/2\pi$ for a unit $F$ by the proof of Proposition~\ref{prop:weak}). Taking
$\mathfrak m(F)=-r\beta/|\beta|$ gives $V_T(\beta)\le V_0-r|\beta|$.
Next, $\lamP\ge V_T(\beta)$ for every $\beta$: if $\mu$ is feasible for (P) then $\int(\sT+\beta\cdot\widehat g)d\mu=\int\sT d\mu$,
and the feasible set of (P) is contained in $\{|\Fh|^2d\tau/2\pi:\|F\|=1\}$ (Proposition~\ref{prop:weak}). Conversely,
by Theorem~\ref{thm:D1} there are, for every $\varepsilon>0$, a certificate $(\nu,\beta)$ of value $c>\lamP-\varepsilon$,
and $c\le V_T(\beta)$, because for a unit $F$ one has $\frac1{2\pi}\int\nuh|\Fh|^2=\int\cF\,d\nu=0$, so that
$\frac1{2\pi}\int(\sT+\beta\cdot\widehat g)|\Fh|^2=\frac1{2\pi}\int f|\Fh|^2\ge c$. Hence $\sup V_T=\lamP$, the
supremum is attained because $V_T$ is continuous and tends to $-\infty$ at infinity, and $\mathcal M_T$, a superlevel
set of a concave function, is closed, bounded and convex.
(ii) As just shown, $c_k\le V_T(\beta^{(k)})\le\lamP$, so $V_T(\beta^{(k)})\to\lamP$. If $\operatorname{dist}(\beta^{(k)},\mathcal M_T)\ge\delta>0$
along a subsequence, that subsequence is bounded (coercivity) and has a limit point $\beta^*$ with
$V_T(\beta^*)=\lamP$ by continuity, i.e.\ $\beta^*\in\mathcal M_T$, a contradiction. The two equivalences follow, the
second one because $V_T(2,0)=\lamW\le\lamP$.
(iii) By Lemma~\ref{lem:compact} the bottom of $W_T$ on the unit sphere is attained, so $\mathcal Z_T\ne\emptyset$
(a complex minimizer has a real or an imaginary part that is a minimizer, by the additivity recalled above).
Compactness: unit vectors with $W_T(F)\le\lamW+\eta$ have $\mathfrak a_T(F)\le\lamW+\eta+4\sinh T$, so by
Lemma~\ref{lem:compact} they form a relatively compact set; $W_T$ is lower semicontinuous on $L^2(\IT)$ (a closed
form bounded below, plus a continuous term), so a limit $F$ of minimizers satisfies $\|F\|=1$ and
$W_T(F)\le\lamW$, hence is a minimizer. The map $\mathfrak m$ is continuous on $L^2(\IT)$, so
$\mathfrak m(\mathcal Z_T)$ and its convex hull are compact.
\emph{Sufficiency.} If $0=\sum_i\theta_i\mathfrak m(F_i)$ with $F_i\in\mathcal Z_T$, $\theta_i\ge0$, $\sum\theta_i=1$, then for
every $\beta'$, $V_T(\beta')\le\mathfrak a_T(F_i)+\beta'\cdot\mathfrak m(F_i)=\lamW+(\beta'-(2,0))\cdot\mathfrak m(F_i)$ for each $i$;
averaging, $V_T(\beta')\le\lamW=V_T(2,0)$, so $(2,0)\in\mathcal M_T$.
\emph{Necessity.} Suppose $0\notin\operatorname{conv}\mathfrak m(\mathcal Z_T)$. By strict separation there are $d\in\R^2$ and
$\delta>0$ with $d\cdot\mathfrak m(F)\ge\delta$ for all $F\in\mathcal Z_T$, and then also for all complex unit minimizers,
since for $F=u+iv$ one has $\mathfrak m(F)=\|u\|^2\mathfrak m(u/\|u\|)+\|v\|^2\mathfrak m(v/\|v\|)$ with $u/\|u\|,v/\|v\|\in\mathcal Z_T$
whenever nonzero. We claim that there is $\eta>0$ with $d\cdot\mathfrak m(F)\ge\delta/2$ for every unit $F\in\mathcal D_T$
with $W_T(F)\le\lamW+\eta$: otherwise there are unit $F_k$ with $W_T(F_k)\to\lamW$ and $d\cdot\mathfrak m(F_k)<\delta/2$;
by the compactness above a subsequence converges in $L^2$ to a unit minimizer $F^*$, and by continuity
$d\cdot\mathfrak m(F^*)\le\delta/2<\delta$, a contradiction. Now for $t>0$ and $\beta=(2,0)$,
$V_T(\beta+td)=\inf_F[W_T(F)+t\,d\cdot\mathfrak m(F)]$ over unit $F\in\mathcal D_T$: if $W_T(F)\le\lamW+\eta$ the bracket is
$\ge\lamW+t\delta/2$, and otherwise it is $\ge\lamW+\eta-tC_T|d|$. For $0<t<\eta/(2C_T|d|)$ this gives
$V_T(\beta+td)\ge\lamW+\min(t\delta/2,\eta/2)>\lamW=V_T(\beta)$, so $\beta\notin\mathcal M_T$.
(iv) Let $F\in\mathcal Z_T$ be even, $x=\ell_{\mathrm{ev}}(F)$. Then $A=B=x$ (the parity paragraph before
Theorem~\ref{thm:parityschur}) and $W_T(F)=\mathfrak a_T(F)+2x^2$. If $x=0$ then $F\in\PT\cap\mathcal D_T$ and
$W_T(F)=\mathfrak a_T(F)\ge\lamA>\lamW$, which is impossible; so $x\ne0$ and $\Rel(A\bar B)=x^2>0$. The odd case is the
same with $A=-B=\ell_{\mathrm{odd}}(F)$. Now let $F=F_{\mathrm e}+F_{\mathrm o}\in\mathcal Z_T$ be arbitrary. By
Theorem~\ref{thm:parityschur}(o), $W_T(F)=W_T(F_{\mathrm e})+W_T(F_{\mathrm o})$ with
$W_T(F_{\mathrm{par}})\ge\lambda^{\mathrm W}_{T,\mathrm{par}}\|F_{\mathrm{par}}\|^2\ge\lamW\|F_{\mathrm{par}}\|^2$, so each nonzero component
is a minimizer in its sector, and if the sector bottoms differ the component in the sector with the larger bottom
vanishes. Hence, when $\lambda^{\mathrm W}_{T,\mathrm{ev}}\ne\lambda^{\mathrm W}_{T,\mathrm{odd}}$, all of $\mathcal Z_T$ lies in one
sector, the first coordinates of $\mathfrak m(\mathcal Z_T)$ are all positive or all negative, $0\notin\operatorname{conv}\mathfrak m(\mathcal Z_T)$,
and (iii), (ii) give $\lamW<\lamP$ and the statement about the multipliers.
(v) $\lamW\le\lamP\le\lamA$ (Theorem~\ref{thm:A}, Remark~\ref{rem:Kprim}) and (ii).
\end{proof}

\begin{remark}[numerical]\label{rem:multnum}
At $T=\frac12\log3$ the Galerkin value $\mathfrak g_{\mathrm{ev}}\approx9\cdot10^{-8}$ of Remark~\ref{rem:polarcapdiag}
is an upper bound for $\mathfrak g_{\mathrm{ev}}$ (up to rounding and quadrature), and by
Theorem~\ref{thm:parityschur}(a4), $\lambda^{\mathrm W}_{T,\mathrm{ev}}\le\mathfrak g_{\mathrm{ev}}/\|h_{\mathrm{ev}}\|^2\le\mathfrak g_{\mathrm{ev}}\|\chi_{\mathrm{ev}}\|^2$
with $\|\chi_{\mathrm{ev}}\|^2=T+\sinh T\approx1.13$; so $\lamW\lesssim10^{-7}<8.8577\cdot10^{-5}\le\lamA$
(Theorem~\ref{thm:C}), and the hypothesis $\lamW<\lamA$ of (iv) holds there numerically. Whether the two parity
bottoms of $W_T$ differ at that window, or at any window, has not been computed; by (iii) this and the second
moment decide whether $\lamP=\lamW$.
\end{remark}

% "The converse is therefore the pair of inequalities ..." and before \section{Certified positivity ...}
\subsection{One direction at \texorpdfstring{$T=\log2$}{T = log 2}}\label{subsec:schur}

A quadratic form splits along any single direction by the classical Schur complement. We use the
direction singled out by the archimedean form: for $T^*_{\infty,\mathrm P}<T\le\log2$ the form $\mT$ has exactly one
negative direction on $\PT$, and it is even (Theorem~\ref{thm:index}(iii),(iv)). The explicit pivot $F^*$ below
approximates the bottom vector of the even sector; nothing depends on the quality of the approximation.
Throughout this subsection $W_T=\mT+\PiT$ on $\PT$ (the polar term vanishes there).

\begin{proposition}[Schur reduction]\label{prop:schur}
Let $\mathfrak w$ be a real quadratic form on a linear subspace $\mathcal Q\subseteq\PT$, $e\in\mathcal Q$ a unit vector with
$w_{11}:=\mathfrak w[e]\ne0$, $\mathcal U=e^\perp\cap\PT$, $\mathcal Q_{\mathcal U}=\mathcal Q\cap\mathcal U$, $\beta(G):=\mathfrak w(e,G)$, and let
$\mathfrak s[G]:=\mathfrak w[G]-\beta(G)^2/w_{11}$ on $\mathcal Q_{\mathcal U}$. Then for every $F=xe+G$ with $x\in\R$, $G\in\mathcal Q_{\mathcal U}$,
\[
\mathfrak w[F]=w_{11}\Big(x+\frac{\beta(G)}{w_{11}}\Big)^2+\mathfrak s[G].
\]
Consequently:
\emph{(i)} if $w_{11}>0$, then $\mathfrak w\ge0$ on $\mathcal Q$ if and only if $\mathfrak s\ge0$ on $\mathcal Q_{\mathcal U}$, and $\mathfrak w\ge c\|\cdot\|^2$
implies $\mathfrak s\ge c\|\cdot\|^2$;
\emph{(ii)} if $\mathfrak w\ge0$ on $\mathcal Q$ then $w_{11}\ge0$; and if $w_{11}=0$, then $\mathfrak w\ge0$ forces $\beta\equiv0$ on
$\mathcal Q_{\mathcal U}$, and then $\mathfrak w\ge0$ if and only if $\mathfrak w|_{\mathcal Q_{\mathcal U}}\ge0$;
\emph{(iii)} if $\mathfrak w=W_T|_{\PT}$ and $e$ lies in the domain of the self-adjoint operator $H$ of $\mathfrak a_T|_{\PT}$
(Lemma~\ref{lem:compact}), then $b:=P_{\mathcal U}He\in\mathcal U$, $\beta=\langle\cdot,b\rangle$ is bounded, and with
$L=I-w_{11}^{-1}(b\otimes e)$, i.e.\ $LF=F-w_{11}^{-1}\langle F,e\rangle b$, the map $L$ is boundedly invertible and
$\mathfrak w[L^{-1}(xe+G)]=w_{11}x^2+\mathfrak s[G]$: the reduction is a congruence, so it preserves the sign of the form
and not its eigenvalues;
\emph{(iv)} $W_T$ commutes with $J$ on $\PT$, so for an even pivot $\beta$ vanishes on the odd part of
$\mathcal Q_{\mathcal U}$ and the Schur correction lives in the even sector alone.
\end{proposition}
\begin{proof}
The identity is the completion of the square
$x^2w_{11}+2x\beta(G)=w_{11}(x+\beta(G)/w_{11})^2-\beta(G)^2/w_{11}$; it is pointwise and uses no topology.
(i) If $\mathfrak s\ge0$ both terms are $\ge0$; conversely take $x=-\beta(G)/w_{11}$, and in the coercive case
$\mathfrak s[G]=\mathfrak w[xe+G]\ge c\|xe+G\|^2\ge c\|G\|^2$. (ii) Take $G=0$, $x=1$; if $w_{11}=0$ and $\beta(G_0)\ne0$
then $\mathfrak w[xe+G_0]=2x\beta(G_0)+\mathfrak w[G_0]\to-\infty$ as $x\to\mp\infty$. (iii) $\beta(G)=\mathfrak a_T(e,G)=\langle G,He\rangle$
for $G\in\mathcal U$ in the form domain, which is $\langle G,b\rangle$; read off the block form. $L^{-1}F=F+w_{11}^{-1}\langle F,e\rangle b$
because $\langle b,e\rangle=0$. (iv) $c_{JF}=\cF$, so $\mT$ and $\PiT$ are $J$-invariant, and $J$ swaps $A$ and $B$, so
$\PT$ is $J$-invariant; an even $e$ is $\mathfrak w$-orthogonal to odd functions (Lemma~\ref{lem:parity}).
\end{proof}

\begin{theorem}[a value at $T=\log2$]\label{thm:w11}
Let $T=\log2$ and let $F^*$ be the explicit even polynomial of degree $17$ whose coefficients in the
orthonormal Legendre basis of $L^2(\IT)$ are given in the certificate data accompanying this paper. Then
\[
W_{\log2}[F^*]\ \in\ [0.0342113430,\ 0.0342182500],
\]
in particular $W_{\log2}[F^*]>0.0342$; and $|A(F^*)|,|B(F^*)|\le1.30\cdot10^{-14}$. The proof includes a computer
verification in interval arithmetic.
\end{theorem}
\begin{proof}
Everything is interval arithmetic on the L\'evy form of Theorem~\ref{thm:A},
$W_T[F]=2A(F)B(F)+C_0c(0)+\int_0^{2T}(c(0)-c(u))K(u)du+c(0)\int_{2T}^\infty K-\sum_{n\in\NT}v_nc(\log n)$,
$c=c_{F^*}$, $\NT=\{2,3\}$. The coefficient vector is read as exact decimal data. (a) $c$ is computed as an exact
polynomial of degree $\le35$ with interval coefficients, by rational arithmetic on the Legendre recurrence, the
binomial expansion of $p(x-s)$, the convolution with $p$ and term-by-term integration, with no quadrature and no
interpolation. (b) $g:=c(0)-c\ge0$ on $[0,2T]$ by Cauchy--Schwarz, which allows integrating an enclosure of $K$
against $g$ term by term. (c) On $[\delta,2T]$, $K$ is squeezed between two affine functions on each panel, since
$1/(2\sinh(u/2))$ is convex and $1/(2\cosh(u/2))$ is concave for $u/2<\operatorname{arccosh}\sqrt2=0.88137$, which covers
the range since $T<0.88137$; $\int g$ and $\int ug$ are evaluated exactly from the interval antiderivatives of $c$, so
no quadrature remainder enters. (d) On $[0,\delta]$, $\delta=10^{-6}$, use $K(u)\le u^{-1}+\frac12$ and $g(u)\le G_1u$
with $G_1$ obtained from the exact Chebyshev expansion of $g$, for an upper contribution $4.8\cdot10^{-6}$ and a lower
one $0$. (e) The remaining terms are closed forms in intervals: $C_0=-\gamma-3\log2-\frac\pi2-\log\pi$,
$\int_{2T}^\infty K=\log\coth\frac T2+\pi-2\arctan e^{T}$ (Theorem~\ref{thm:S}(iii)), and the two atoms by interval
Horner. (f) The enclosure does not assume $F^*\in\PT$: it pays for the defect,
$|A|,|B|\le1.30\cdot10^{-14}$ by exact interval integration of the polynomial against $e^{\pm t/2}$, so
$|2AB|\le3.4\cdot10^{-28}$. The total width is $6.91\cdot10^{-6}$.
\end{proof}

This is a statement about one explicit direction at $\log2$, namely the easy half, $w_{11}>0$, of
Proposition~\ref{prop:schur} at $e=F^*$; the sign of the form on $\PT$ at $\log2$ is given by
Theorem~\ref{thm:log5A}.

\begin{theorem}[positive definiteness on nested even subspaces at $T=\log2$]\label{thm:nest}
Let $T=\log2$ and let $e_1=F^*,e_3,e_5,\dots,e_{13}$ be the seven explicit even polynomials of degree $17$ given in the
certificate data accompanying this paper (approximate even eigenvectors of $\mT$ on $\PT$; their Gram matrix
differs from the identity by at most $1.07\cdot10^{-15}$ per entry, and their Tate moments are bounded as in
Theorem~\ref{thm:w11}). Then $W_{\log2}$ is positive definite on every
$E'_k=\mathrm{span}\{e_1,e_3,\dots,e_{2k-1}\}$, $k=1,\dots,7$, with certified lower bounds for the smallest
eigenvalues of its Gram-normalized matrix
\begin{gather*}
3.4211\cdot10^{-2},\quad2.1136\cdot10^{-3},\quad4.0453\cdot10^{-4},\quad9.0763\cdot10^{-5},\\
4.0951\cdot10^{-5},\quad2.2156\cdot10^{-5},\quad6.5687\cdot10^{-6};
\end{gather*}
in particular $W_{\log2}[F]\ge6.56\cdot10^{-6}\|F\|^2$ for every $F\in E'_7$, and the determinant of the
$2\times2$ block on $\mathrm{span}\{F^*,e_3\}$ is at least $2.10\cdot10^{-4}$. The same procedure certifies $E'_5$ at $T=0.62$
with smallest eigenvalue $\ge3.92\cdot10^{-4}$ and at $T=0.65$ with smallest eigenvalue $\ge1.70\cdot10^{-4}$. The proof
includes a computer verification in interval arithmetic.
\end{theorem}
\begin{proof}
Only quadratic values of $W_T$ are enclosed: the off-diagonal entries are obtained by polarization,
$W_T(e_i,e_j)=\frac14\big(W_T[e_i+e_j]-W_T[e_i-e_j]\big)$, exact because $W_T$ is a quadratic form
(Theorem~\ref{thm:A}; $A,B$ are linear, so $2AB$ is quadratic). This matters: the affine sandwich in the proof of
Theorem~\ref{thm:w11} is licensed by $c(0)-c\ge0$, which is Cauchy--Schwarz for an autocorrelation and false for a
cross-correlation. Each quadratic value is enclosed exactly as in Theorem~\ref{thm:w11}. Because the panelwise
bounds do not depend on the integrand's coefficients, the panel work is done once, as two fixed linear functionals
on the monomial coefficients. Interval inclusion monotonicity gives an interval matrix $\boldsymbol M$ containing the
true matrix, and an interval Cholesky factorization of $\boldsymbol M-\mu I$ that runs to completion with every pivot
interval strictly positive certifies $\lambda_{\min}>\mu$ for every real symmetric matrix in $\boldsymbol M$; bisection on
$\mu$ gives the stated bounds. The transfer to $\|F\|^2$ costs nothing at this precision, by the Gram bound.
\end{proof}

This is positivity on an explicit finite-dimensional subspace; it extends Theorem~\ref{thm:w11} from a line
to a $7$-plane.

\begin{proposition}[the Schur coupling is purely arithmetic]\label{prop:bprime}
In Proposition~\ref{prop:schur}(iii), if $e$ is an eigenvector of the self-adjoint operator of $\mT|_{\PT}$ lying in
the domain of $H$, then $b=P_{\mathcal U}He=P_{\mathcal U}\Pi e$, where $\Pi$ is the bounded operator of the prime form
$\PiT$ on $\PT$: the coupling between the pivot and its orthogonal complement is carried by the prime form alone.
\end{proposition}
\begin{proof}
The operator of $\mT|_{\PT}$ maps $e$ to a multiple of $e$, which $P_{\mathcal U}$ annihilates, and $H$ is the sum of
that operator and $\Pi$ on the domain of $H$, since $\PiT$ is bounded.
\end{proof}

\begin{proposition}[the prime form is a sum of compressed shifts]\label{prop:shifts}
With $V_a$ the right shift and $Q_T$ the compression to $\IT$ (\S\ref{subsec:conv}),
\[
\PiT(F)\ =\ -\sum_{n\in\NT}v_n\,\Rel\big\langle F,Q_TV_{\log n}Q_TF\big\rangle .
\]
Each $Q_TV_aQ_T$, $0<a<2T$, is isometric on $L^2(-T,T-a)$ and vanishes on $L^2(T-a,T)$, so it is a partial isometry
with infinite-dimensional defect. $\PiT$ is bounded, $\|\PiT\|\le\Aconst$, and for $T>\frac12\log2$ its
operator $\Pi$ on $L^2(\IT)$ is not compact. For any orthonormal system $(e_i)$ of $\PT$ the entries
$\pi_{ij}=\PiT(e_i,e_j)$ satisfy $|\pi_{ij}|\le\Aconst$ and $\sum_j\pi_{ij}^2\le\Aconst^2$.
\end{proposition}
\begin{proof}
$\cF(u)=\int F(t)\overline{F(t-u)}dt=\langle F,V_uF\rangle$ for $F$ extended by zero, and only the part of $V_uF$ inside
$\IT$ contributes, whence the formula. The isometry and defect statements are immediate from $(V_aG)(t)=G(t-a)$ read
on $\IT$. Boundedness is the triangle inequality with $\sum v_n=\Aconst$. Non-compactness: $\Pi=-\frac12\sum_{n\in\NT}v_n
(Q_TV_{\log n}Q_T+Q_TV_{-\log n}Q_T)$. Choose $\epsilon>0$ smaller than $\min_{n\in\NT}\log n$, than $2T-\max_{n\in\NT}\log n$
and than the smallest difference between distinct elements of $\{\log n:n\in\NT\}$, and let $(G_j)$ be an orthonormal
sequence in $L^2(-T,-T+\epsilon)$. Then $Q_TV_{-\log n}G_j=0$, and the functions $V_{\log n}G_j$, $n\in\NT$, have pairwise
disjoint supports $(-T+\log n,-T+\log n+\epsilon)\subset\IT$; so $\|\Pi G_j\|^2=\frac14\sum_{n\in\NT}v_n^2>0$ for every $j$,
while $G_j\to0$ weakly. A compact operator maps weakly null sequences to norm null ones. The row bounds are Bessel's
inequality, $\sum_j\pi_{ij}^2\le\|\Pi e_i\|^2\le\Aconst^2$, the entries being real for real $e_i$.
\end{proof}

The atoms of the prime matrix are one-sided overlaps, and their signs are nodal statements.

\begin{lemma}[one-sided overlap]\label{lem:N}
Let $F,e\in L^2(\IT)$ be real and even and $0<a<2T$. Then
\[
\Rel\big\langle Q_TV_aQ_TF,\,e\big\rangle=\int_{a-T}^{T}F(z-a)\,e(z)\,dz=:I_{F,e}(a),
\]
and consequently, for real even $e_i,e_j$,
$\PiT(e_i,e_j)=-\sum_{n\in\NT}v_n\,\frac12\big(I_{e_i,e_j}(\log n)+I_{e_j,e_i}(\log n)\big)$; for real even $F=e$ this is
$\PiT(F)=-\sum_{n\in\NT}v_nI_{F,F}(\log n)$.
\end{lemma}
\begin{proof}
For $z\in\IT$, $(Q_TV_aQ_TF)(z)=F(z-a)$ if $z-a>-T$ and $0$ otherwise, which gives the integral. By the
substitution $z\mapsto-z$ and evenness, $\langle Q_TV_{-a}Q_TF,e\rangle=\int_{-T}^{T-a}F(z+a)e(z)dz=I_{F,e}(a)$ as well,
so $\frac12\langle(Q_TV_aQ_T+Q_TV_{-a}Q_T)F,e\rangle=I_{F,e}(a)$. The polarized prime form is
$\PiT(F,e)=-\sum_nv_n\frac12\Rel\langle(Q_TV_{\log n}Q_T+Q_TV_{-\log n}Q_T)F,e\rangle$ by Proposition~\ref{prop:shifts}, which
is symmetric in $F,e$.
\end{proof}

\begin{proposition}[sign forcing]\label{prop:G}
Let $F,e\in L^2(\IT)$ be real, even and continuous, with $e>0$ on $(z_e,T)$ and $F<0$ on $(z_F,T)$ for some
$0\le z_e,z_F<T$. Then for every $a$ with $T+\max(z_F,z_e)<a<2T$ one has $I_{F,e}(a)<0$, so that the corresponding
atom $-v_nI_{F,e}(\log n)$ at $a=\log n$ is $\ge0$, and $>0$ when $\Lambda(n)\ne0$. The sign of a contact is therefore
forced once the contact lies beyond $a^*:=T+\max(z_F,z_e)$.
\end{proposition}
\begin{proof}
For $z\in(a-T,T)$ one has $z>a-T>z_e$, so $e(z)>0$; and $z-a\in(-T,T-a)$ with $T-a<-z_F$, so $z-a\in(-T,-z_F)$,
where $F<0$ by evenness. The integrand is negative on the interval of integration, which is nonempty since $a<2T$.
\end{proof}

\begin{remark}[what the numerics say at $\log2$]\label{rem:galerkin}
Galerkin (Rayleigh--Ritz) values on nested subspaces are upper bounds for the bottom of the form on $\PT$. At
$T=\log2$ they converge, in the basis size, to $\lambda_{\mathrm A}(\log2)\approx9.83\cdot10^{-8}$ (numerical, not
certified; Appendix~\ref{app:num}), far below the pivot value of Theorem~\ref{thm:w11} and the subspace bounds of
Theorem~\ref{thm:nest}. The other half of Proposition~\ref{prop:schur}, positivity of the Schur form on $\mathcal U$, is
not reachable by a norm inequality that separates the archimedean block from the prime block: numerically, even
the condition in which the prime form is dropped from the Schur form fails at $T=0.58,0.62,0.65,\log2$
(Appendix~\ref{app:diag}), so positivity on $\mathcal U$ is a cancellation between the Schur correction and the prime
form.
\end{remark}


\section{Certified positivity up to \texorpdfstring{$T=\tfrac12\log5$}{T = (1/2) log 5}}\label{sec:log5}

Put $T_5:=\frac12\log5=0.804718\ldots$. Since $\NT=\{n\ge2:\log n<2T\}$, the contact set is
$\{2,3,4\}$ for $\frac12\log4\le T\le T_5$, and $5$ enters only for $T>T_5$; at $T_5$,
$\Aconst=v_2+v_3+v_4=2.94186\ldots$. The window $T_5$ is therefore the largest one with this contact set,
and it lies slightly beyond the range $T\le0.8$ of \cite[Thm.~1.2]{Zhu2026}; the ranges $T\le1$ and $T\le17/16$ are claimed, with much
smaller constants $2^{-151}$ and $2^{-49162}$, in \cite[Thms.~A, B]{Liu2026}. The notation
$\chi_{\mathrm{par}},\ell_{\mathrm{par}},\varepsilon_{\mathrm{par}},\mathcal Q^{\mathrm{par}}_T,
\mathfrak g_{\mathrm{par}},\mathfrak G_T$ is that of \S\ref{sec:complete}, and
$\lambda^{\mathrm W}_{T,\mathrm{par}}$ is the bottom of $W_T$ on the unit sphere of
$\mathcal D^{\mathrm{par}}_T$ (Theorem~\ref{thm:parityschur}(o)).

\begin{theorem}[the primitive form at $\frac12\log5$]\label{thm:log5A}
$\lambda^{\mathrm{ev}}_{\mathrm A}(T_5)\ge2.5\cdot10^{-12}$ and
$\lambda^{\mathrm{odd}}_{\mathrm A}(T_5)\ge3.0\cdot10^{-10}$. Hence
$\lamA\ge2.5\cdot10^{-12}$ for every $0<T\le\frac12\log5$.
\end{theorem}

\noindent The proof includes a computer verification in ball arithmetic (Appendix~\ref{subsec:encllog5}).

\begin{proof}
By Theorem~\ref{thm:parityschur}(o), $\lambda_{\mathrm A}(T_5)$ is the minimum of the two
parity bottoms, which are bounded below by Proposition~\ref{prop:lowercert} with $c=0$ and the
certificate data of Appendix~\ref{subsec:encllog5}. The bound for $T<T_5$ follows from
Lemma~\ref{lem:monoT}.
\end{proof}

\begin{theorem}[the complete form at $\frac12\log5$]\label{thm:log5W}
$W_{T_5}\ge10^{-18}\|\cdot\|^2$ on $\mathcal D^{\mathrm{ev}}_{T_5}$ and $W_{T_5}\ge10^{-16}\|\cdot\|^2$ on
$\mathcal D^{\mathrm{odd}}_{T_5}$. Hence $\lamW\ge10^{-18}$ for every $0<T\le\frac12\log5$, and
$\mathfrak G_T\succ0$ there.
\end{theorem}

\noindent The proof includes a computer verification in ball arithmetic (Appendix~\ref{subsec:encllog5}).

\begin{proof}
By Theorem~\ref{thm:parityschur}(o), $\lambda^{\mathrm W}_{T_5}=\min_{\mathrm{par}}\lambda^{\mathrm W}_{T_5,\mathrm{par}}$,
and each parity is bounded below by Proposition~\ref{prop:lowercert} with $c=-\frac12\varepsilon_{\mathrm{par}}$.
Lemma~\ref{lem:monoT} gives the bound for $T<T_5$. Since $\lamW>0$, Theorem~\ref{thm:parityschur}(c)
gives $\mathfrak G_T\succ0$.
\end{proof}

\begin{corollary}[indeterminacy for $T\le\frac12\log5$]\label{cor:KEPlog5}
$\mathrm{KEP}(T)$ is indeterminate for every $0<T\le\frac12\log5$.
\end{corollary}
\begin{proof}
By Theorem~\ref{thm:log5W}, $\lamW>0$ for $T\le T_5$, so Theorem~\ref{thm:KEPdich} applies for
$\frac12\log2\le T\le T_5$. For $0<T<\frac12\log2$, note that $\mathcal W_{T'}\subseteq\Wcl$ for $T'>T$ (a weight
flat on $[-2T',2T']$ is flat on $[-2T,2T]$, and the decay condition does not involve $T$), so a weight
$w$ with $\inf f_w>0$ at $T'=\frac12\log2$ is admissible at $T$, and $f_w$ does not depend on the window
by \eqref{eq:fw}; the perturbation of Theorem~\ref{thm:KEP1}(ii) applies, its $\chi$ being supported in
$\{|u|\ge\log2\}$, outside the window.
\end{proof}

\begin{remark}[scope]\label{rem:log5scope}
The gain over \cite[Thm.~1.2]{Zhu2026} is small, from $0.8$ to $0.80472$, and it is obtained by the
window variable alone: on $[0.8,T_5]$ the contact set is still $\{2,3,4\}$. It is not the largest
certified window: \cite{Liu2026} (preprint under review) claims $T\le17/16$, and \cite{Mori15} (unrefereed,
computational part reproduced by us) gives $T\le\frac12\log15$. Our certificate is
independent, uses a different section, and gives constants larger by more than thirty orders of magnitude
on $T\le\frac12\log5$. The bound
$10^{-18}$ is consistent with Zhu's $8.9\cdot10^{-18}$ at $T=0.8$ and with the monotonicity of
$\lamW$. The complete bottom is about six orders of magnitude below the primitive one, which is the
size of the polar capacity $\mathfrak g_{\mathrm{ev}}/\|h_{\mathrm{ev}}\|^2$ of
Theorem~\ref{thm:parityschur}(a4). No certified upper bound is given; the values at which the
verification fails (Appendix~\ref{subsec:encllog5}) are not claims of negativity. Both theorems are
certificates, not a mechanism. By Theorem~\ref{thm:margin} and the doubly exponential growth of the
verification threshold (Proposition~\ref{prop:env}), the same method becomes rapidly more expensive as $T$
grows; the first window with the new contact $5$ is not treated here (it is covered by
\cite[Thm.~B]{Liu2026} and \cite{Mori15}, cited).
\end{remark}

\emph{The converse across windows.} Theorem~\ref{thm:parityschur}(c) is a statement at one window. The
following lemma lets a single window control all smaller ones.

\begin{lemma}[monotonicity of the polar values]\label{lem:gmono}
Let $0<T<T'$ and suppose $\lambda^{\mathrm{par}}_{\mathrm A}(S)>0$ for $S=T$ and $S=T'$. Then
$\mathfrak g_{\mathrm{par}}(T')\le\mathfrak g_{\mathrm{par}}(T)$.
\end{lemma}
\begin{proof}
By Theorem~\ref{thm:parityschur}(a3),
$\mathfrak g_{\mathrm{par}}(S)=\min\{W_S(F):F\in\mathcal D^{\mathrm{par}}_S,\ \ell_{\mathrm{par}}(F)=1\}$. Extension by
zero maps $\mathcal D^{\mathrm{par}}_T$ into $\mathcal D^{\mathrm{par}}_{T'}$; it preserves the parity, the value of
$\ell_{\mathrm{par}}$ (the function $\chi_{\mathrm{par}}$ on $I_{T'}$ restricts to $\chi_{\mathrm{par}}$ on $\IT$),
and the value of $W$, by the proof of Lemma~\ref{lem:monoT}. So the feasible set at $T'$ contains that at
$T$, and the minimum does not increase.
\end{proof}

\begin{proposition}[one-window converse]\label{prop:onewindow}
Suppose $\lamA>0$ for $0<T\le T_1$. Then
\emph{(i)} $\lamW\ge0$ for all $0<T\le T_1$ if and only if $\mathfrak g_{\mathrm{ev}}(T_1)\ge0$ and
$\mathfrak g_{\mathrm{odd}}(T_1)\ge0$;
\emph{(ii)} $\lamW>0$ for all $0<T\le T_1$ if and only if $\mathfrak G_{T_1}\succ0$.
In either case $\mathfrak g_{\mathrm{par}}(T)\ge\mathfrak g_{\mathrm{par}}(T_1)$ for $0<T\le T_1$.
\end{proposition}
\begin{proof}
By Lemma~\ref{lem:monoT}, the conditions on $(0,T_1]$ are equivalent to the same conditions at $T_1$. By
Theorem~\ref{thm:parityschur}(c), since $\lamA>0$ at $T_1$, $\lambda^{\mathrm W}_{T_1}\ge0$ if and only if
$\mathfrak g_{\mathrm{ev}}(T_1),\mathfrak g_{\mathrm{odd}}(T_1)\ge0$, and $\lambda^{\mathrm W}_{T_1}>0$ if and only if
$\mathfrak G_{T_1}\succ0$. Since $\lamA>0$ implies that both parity bottoms are positive
(Theorem~\ref{thm:parityschur}(o)), Lemma~\ref{lem:gmono} gives the last statement.
\end{proof}

By Theorem~\ref{thm:log5A}, Proposition~\ref{prop:onewindow} applies with $T_1=T_5$, and
Theorem~\ref{thm:log5W} is the statement that its condition (ii) holds there.

\begin{corollary}[the polar capacity beyond $\frac12\log5$]\label{cor:beyond5}
There is $\varepsilon>0$ such that $\lamW>0$, $\lamA>0$ and $\mathfrak G_T\succ0$ for every
$0<T\le\frac12\log5+\varepsilon$. In particular the polar capacity is positive definite at windows beyond
$\frac12\log5$, in the presence of the contact $5$.
\end{corollary}
\begin{proof}
$\lambda^{\mathrm W}_{T_5}\ge10^{-18}$ by Theorem~\ref{thm:log5W}, and $T\mapsto\lamW$ is continuous by
Lemma~\ref{lem:contT}, whose proof sums the comb over $n<e^{4T}$ and so covers the contact $5$ entering at $T_5$.
Hence $\lamW>0$ on $[T_5,T_5+\varepsilon]$ for some $\varepsilon>0$, and on $(0,T_5]$ by Lemma~\ref{lem:monoT}.
Theorem~\ref{thm:parityschur}(c) gives $\lamA>0$ and $\mathfrak G_T\succ0$ at each such window.
\end{proof}
\begin{remark}\label{rem:liu}
Granting \cite[Thm.~B]{Liu2026}, $\lamW\ge2^{-49162}>0$ on $(0,17/16]$, which contains the contacts
$5,7,8$, and then $\lamA>0$ there by Theorem~\ref{thm:A} and Remark~\ref{rem:Kprim}. Corollary~\ref{cor:beyond5}
is the in-paper, unconditional statement. Granting \cite{Mori15}, the same holds with $\lamW\ge2^{-2904}$ on
$(0,\frac12\log15]$, which contains the contacts $5,7,8,9,11,13$.
\end{remark}
The $\varepsilon$ that the proof of Lemma~\ref{lem:contT} yields is of no practical use: the modulus of continuity
there is logarithmic in the dilation (step (ii) of that proof), and the margin at $T_5$ is $10^{-18}$. A window
beyond $\frac12\log5$ at which $\mathfrak G_T\succ0$ is certified, with its value, remains to be produced.


\subsection{The per-window converse fails for modified weights}
Put $F^{\mathrm{odd}}_T(t)=t$ on $\IT$, as in Proposition~\ref{prop:archneg}, and use the comb-level forms of \S\ref{subsec:finite}.

\begin{corollary}[the converse fails for modified weights]\label{cor:convfail}
The proof uses Theorem~\ref{thm:index}, Corollary~\ref{cor:Tarch} and Theorem~\ref{thm:C}, whose proofs include
computer verifications, two further evaluations in interval arithmetic, and Theorem~\ref{thm:log5W}.
\emph{(i) Deletion.} Let $a=\Lambda\mathbf 1_{n\ge5}$. For every $T\in[0.4608,0.58]$, $\mathrm{WP}_T(a)$ holds
strictly while $\lambda^{\mathrm W}_T(a)<0$:
\[
\lambda_{\mathrm A}(T;a)=\laminfP(T)>0,\qquad\lambda^{\mathrm W}_T(a)=\laminf(T)<0 .
\]
On $[0.4608,\log2]$ the odd-sector implication fails. There the bottom of $\mathfrak a^{(v(a))}_T$ on $\PT^{\mathrm{odd}}$ is
$\lambda_{\mathrm{odd}}(T)>0$, while $W_T(a)$ is negative on an odd vector. On $(T^*_{\infty,\mathrm P},\log2]$,
$\mathrm{WP}_T(a)$ fails, so the full implication holds there vacuously and only its odd-sector form is refuted.
\emph{(ii) One contact.} Let $T=\frac12\log3$, so $\NT=\{2\}$, and for $\theta\in[0,1]$ let
$a_\theta:=\Lambda-(1-\theta)\Lambda(2)\mathbf 1_{n=2}$. Its comb on $\NT$ is $\theta v_2$. Then
$\lambda_{\mathrm A}(T;a_\theta)>0$ for every $\theta\in[0,1]$, and $\lambda^{\mathrm W}_T(a_\theta)<0$ for every
$\theta\in[0,0.8148]$. The set of $\theta\in[0,1]$ with $\lambda^{\mathrm W}_T(a_\theta)\ge0$ is an interval
$[\theta_*,1]$ with $0.8148<\theta_*<1$.
\emph{(iii) Consequences.} Let $\mathcal C$ be a class of nonnegative arithmetic functions containing $\Lambda$.
If $\mathcal C$ is closed under deletion of finitely many terms, it contains $\Lambda\mathbf 1_{n\ge5}$, for which
the implication $\mathrm{WP}_T(a)\Rightarrow\lambda^{\mathrm W}_T(a)\ge0$ fails at every $T\in[0.4608,0.58]$ by (i). If
$\mathcal C$ contains $a_\theta$ for some $\theta\le0.8148$ (for instance, if $\mathcal C$ is closed under lowering
one term within $[0,\Lambda(n)]$), the implication fails for that member at $T=\frac12\log3$ by (ii). The classes
$\mathcal C_{\mathrm{PNT}}(\varepsilon)$, $\mathcal C_{\mathrm{zeros}}$ and $\mathcal C_{\mathrm{BT}}$ of
Corollary~\ref{cor:void} are closed under deletion and under lowering one term. Consequently no argument
that uses about the primes only a property preserved by these modifications (for instance a prime number
theorem with error term, a zero-free region, a zero-density estimate, or an upper-bound sieve inequality)
can prove $\mathrm{WP}_T\Rightarrow\lamW\ge0$ at any $T\in[0.4608,0.58]$ or at $T=\frac12\log3$. At $T=\frac12\log3$
such an argument must distinguish $\Lambda(2)$ from every weight $\theta\Lambda(2)$ with $\theta\le0.8148$, that is,
from weights at least $18.5\%$ below it; nothing is claimed about weights above $\Lambda(2)$.
\end{corollary}
\begin{proof}
(i) For $T\le\log2<\frac12\log5$ one has $\NT\subseteq\{2,3,4\}$, where $a=0$. So the comb of $a$ is empty,
$W_T(a)=\Ainf$ and $\mathfrak a^{(v(a))}_T=\mT$ on $\PT$, as in the proof of Theorem~\ref{thm:finmod}(i). For
$T\le0.58$, $\laminfP(T)>0$ by Theorem~\ref{thm:index}(i), and for $T\ge0.4608$, $\laminf(T)<0$ by
Proposition~\ref{prop:archneg}. On $[0.4608,\log2]$, $\lambda_{\mathrm{odd}}(T)>0$ by Theorem~\ref{thm:index}(i). The
odd vector $F^{\mathrm{odd}}_{0.4608}$, extended by zero, lies in $L^2(\IT)$ with
$\Ainf(F^{\mathrm{odd}}_{0.4608})\le U(0.4608)<0$ ($\Ainf$ is $T$-free). $\mathrm{WP}_T(a)$ fails on
$(T^*_{\infty,\mathrm P},\log2]$ by Lemma~\ref{lem:strict}.

(ii) The comb: $\NT=\{n:\log n<\log3\}=\{2\}$ and $v_2(a_\theta)=\theta v_2$. So
$W_T(a_\theta)=\Ainf+\theta\PiT$, and $\mathfrak a^{(v(a_\theta))}_T=\mT+\theta\PiT$ on $\PT$.

\emph{The primitive side.} $\theta\mapsto\lambda_{\mathrm A}(T;a_\theta)=\inf_{F\in\PT,\|F\|=1}[\mT(F)+\theta\PiT(F)]$
is an infimum of affine functions of $\theta$. It is finite, since $\|\PiT\|\le v_2$, and hence concave on
$\R$. At $\theta=0$ it is $\laminfP(\frac12\log3)>0$ (Corollary~\ref{cor:Tarch}). At $\theta=1$ it is
$\lamA\ge8.8577\cdot10^{-5}$ (Theorem~\ref{thm:C}, which is stated for the actual comb $v_2$ at this window). A
concave function on $[0,1]$ is at least the smaller of its endpoint values.

\emph{The complete side.} Let
$F_5(t)=t-\frac{2634299}{500000}t^3+\frac{1919331}{250000}t^5$ on $\IT$, with $T=\frac12\log3$. As in
Proposition~\ref{prop:archneg}, $\Ainf(F_5)\le U_0$. Here $U_0$ is obtained by bounding $K\le\frac1u+\frac12$
against $c_{F_5}(0)-c_{F_5}(u)\ge0$ (Cauchy--Schwarz) on $(0,2T)$, by Theorem~\ref{thm:S}(iii) on $[2T,\infty)$,
and by $2A(F_5)B(F_5)=-2s(F_5)^2$. The autocorrelation $c_{F_5}$ is an exact polynomial in $(u,T)$ with rational
coefficients. Evaluating in outward-rounded interval arithmetic at $50$ digits with $T=\frac12\log3$ gives
\begin{gather*}
U_0\in[-0.00613778443569004,\,-0.00613778443569003],\\
\PiT(F_5)=-v_2c_{F_5}(\log2)\in[0.00753283781432786,\,0.00753283781432787],
\end{gather*}
and $U_0+0.8148\,\PiT(F_5)\in[-2.81845757,-2.81845756]\cdot10^{-8}$. Since $\PiT(F_5)>0$, we get
$W_T(a_\theta)(F_5)=\Ainf(F_5)+\theta\PiT(F_5)\le U_0+0.8148\,\PiT(F_5)<0$ for every $\theta\le0.8148$. (The bare
witness $F^{\mathrm{odd}}_T$ gives the same for $\theta\le0.4741$.) Finally, $\theta\mapsto\lambda^{\mathrm W}_T(a_\theta)$
is concave and finite on $\R$ for the same reason, hence continuous. At $\theta=1$ it is
$\lamW>0$ by Theorem~\ref{thm:log5W}. So
$\{\theta\in[0,1]:\lambda^{\mathrm W}_T(a_\theta)\ge0\}$ is a closed interval that contains $1$ and is disjoint from
$[0,0.8148]$. It has the form $[\theta_*,1]$, and $\theta_*<1$ by continuity.

(iii) By (i) and (ii), the implication is false for $\Lambda\mathbf 1_{n\ge5}$ (delete $n=2,3,4$) and for $a_\theta$.
For the closure properties: lowering one term keeps $0\le a\le\Lambda$; it changes $\psi_a$ by a constant for
$x\ge n$; and $\zeta_a=\zeta\exp\big(-(1-\theta)\Lambda(n)n^{-s}/\log n\big)$ with an entire zero-free factor, as in
Theorem~\ref{thm:finmod}(iii).
\end{proof}

As in Corollary~\ref{cor:void}, the form $W_T(a)$ is not the zero sum of $\zeta_a$ (\S\ref{subsec:finite}), so
$\lambda^{\mathrm W}_T(a)<0$ says nothing about the zeros of $\zeta_a$, which are those of $\zeta$; moreover $\zeta_a$
has no functional equation (Theorem~\ref{thm:finmod}(iii)). We have not examined a model with the
functional equation. For the empty comb, (i) is exactly the strict separation
$T^*_\infty<T^*_{\infty,\mathrm P}$, proved by Proposition~\ref{prop:archneg} with Corollary~\ref{cor:Tarch}

\section{A finite reduction}\label{sec:coer}

This section localizes the positivity question for $\mathfrak a_T$ on $\PT$ to an explicit subspace of
dimension $2k-2$, the \emph{piecewise Tate exponentials}, by proving that the archimedean form is coercive,
with explicit constants, on the complementary subspace of functions that satisfy the two Tate conditions on
each of $k$ consecutive blocks. The coercivity is a high-pass phenomenon ($m(\tau)\sim\log|\tau|$); the prime
form is handled by a norm bound. All forms are on $L^2(\IT)$ with its metric, $\mathfrak a_T=\mT+\PiT$ on
$\mathcal D_T$, and $\mT(F)=C_0\|F\|^2+\Lev(\cF)$ does not depend on $T$ except through its domain. In this
section $r=T/k$ denotes the half-length of a block.

\subsection{Blocks and the Tate-orthogonal subspace}\label{subsec:coer-blocks}

\begin{definition}\label{def:coer-V}
Let $k\ge1$ and let $B_1,\dots,B_k$ be the $k$ consecutive open intervals of length $2r$,
$r:=T/k$, that partition $\IT$ up to endpoints, $B_j=(-T+2(j-1)r,\,-T+2jr)$. Put
\[
\mathcal V_k:=\Big\{F\in L^2(\IT):\ \int_{B_j}F(t)e^{t/2}dt=\int_{B_j}F(t)e^{-t/2}dt=0\ \ (1\le j\le k)\Big\},
\qquad E_k:=\PT\ominus\mathcal V_k ,
\]
and $e_j^\pm:=e^{\pm t/2}\mathbf 1_{B_j}$. (Lemma~\ref{lem:coer-decomp} holds for any partition into
intervals; equal blocks are used from Lemma~\ref{lem:coer-lowpass} on.)
\end{definition}

\begin{lemma}[the decomposition $\PT=\mathcal V_k\oplus E_k$]\label{lem:coer-decomp}
\emph{(i)} $\mathcal V_k$ is a closed subspace of $\PT$. \emph{(ii)} The orthogonal complement of
$\mathcal V_k$ in $L^2(\IT)$ is $\operatorname{span}\{e_j^\pm\}$, of dimension $2k$.
\emph{(iii)} $E_k=\operatorname{span}\{e_j^\pm\}\cap\PT=\{\sum_jc_j^+e_j^++c_j^-e_j^-:\ A=B=0\}$ has
dimension $2k-2$, and $\PT=\mathcal V_k\oplus E_k$ (orthogonal in $L^2(\IT)$).
\emph{(iv)} $E_k\subset\mathcal D_T$, and $\PT\cap\mathcal D_T=(\mathcal V_k\cap\mathcal D_T)\oplus E_k$.
\end{lemma}
\begin{proof}
(i) The $2k$ conditions are continuous linear functionals; summing them over $j$ gives $A(F)=B(F)=0$.
(ii) $\mathcal V_k$ is the joint kernel of the functionals $F\mapsto\langle F,e_j^\pm\rangle$, and the $2k$
functions $e_j^\pm$ are linearly independent (disjoint supports; $e^{t/2},e^{-t/2}$ independent on an
interval). (iii) $E_k=\mathcal V_k^\perp\cap\PT$ since $\mathcal V_k\subset\PT$. The two functionals
$A,B$ restricted to $\operatorname{span}\{e_j^\pm\}$ are independent: on the two-dimensional space
$\operatorname{span}\{e_1^+,e_1^-\}$ their matrix is
$\begin{psmallmatrix}\int_{B_1}e^t&|B_1|\\|B_1|&\int_{B_1}e^{-t}\end{psmallmatrix}$, of determinant
$\int_{B_1}e^t\int_{B_1}e^{-t}-|B_1|^2>0$ (Cauchy--Schwarz, strict since $e^t$ is not constant). Hence
$\dim E_k=2k-2$. If $F\in\PT$ and $F=v+w$ with $v\in\mathcal V_k$, $w\in\mathcal V_k^\perp$, then
$w=F-v\in\PT$, so $w\in E_k$. (iv) $\widehat{e_j^\pm}(s)=\big[e^{(\pm\frac12+is)t}/(\pm\frac12+is)\big]_{\partial B_j}$,
so $|\widehat{e_j^\pm}(s)|\le2e^{T/2}(\frac14+s^2)^{-1/2}$ and $\int\log(2+|s|)|\widehat{e_j^\pm}|^2<\infty$.
Since $\mathcal D_T$ is a linear space containing $E_k$, the $\mathcal V_k$-component of an $F\in\PT\cap\mathcal D_T$
lies in $\mathcal D_T$.
\end{proof}

For a window $r>0$ we write $\lambda_{\mathrm A}(r)$ for the bottom of $\mathfrak a_r$ on the unit sphere of
$\mathcal P_r\cap\mathcal D_r$; it is attained (Lemma~\ref{lem:compact}) and continuous in $r$ (Lemma~\ref{lem:contA}).
For $2r\le\log2$ the contact set $\mathcal N_r$ is empty, so $\mathfrak a_r=\mT$ on $L^2(I_r)$ and $\lambda_{\mathrm A}(r)$ is a
purely archimedean quantity. We use $|\cF(u)|\le\cF(0)=\|F\|^2$ and $\cF(u)=0$ for $|u|\ge2T$.

\begin{lemma}[the reduced kernel]\label{lem:coer-K1}
Put $K_1(u):=K(u)-2e^{-u/2}=\dfrac{2e^{-5u/2}}{1-e^{-2u}}=2\sum_{j\ge1}e^{-(2j+\frac12)u}$ for $u>0$, with $K$ as in
Lemma~\ref{lem:K}. Then $0<K_1(u)\le u^{-1}$ for all $u>0$, and $K_1(u)=u^{-1}-\frac32+O(u)$ as $u\downarrow0$. Let
$F,G\in L^2(\IT)$ be real, $\supp F\subset\overline B$, $\supp G\subset\overline{B'}$, with $B,B'$ disjoint intervals and $B'$
to the right of $B$. Then
\begin{equation}\label{eq:coer-cross}
\mT(F+G)=\mT(F)+\mT(G)-\iint F(x)G(y)K(y-x)\,dx\,dy ,
\end{equation}
and if $\int F(x)e^{x/2}dx=0$ or $\int G(y)e^{-y/2}dy=0$, the kernel $K$ in \eqref{eq:coer-cross} may be replaced by $K_1$.
\end{lemma}
\begin{proof}
The three expressions for $K_1$ agree by Lemma~\ref{lem:K} ($K=2e^{-u/2}\sum_{j\ge0}e^{-2ju}$). $K_1\le u^{-1}$ is
$e^{5u/2}-e^{u/2}\ge2u$: the function $g(u)=e^{5u/2}-e^{u/2}-2u$ has $g(0)=g'(0)=0$ and
$g''=\frac{25}4e^{5u/2}-\frac14e^{u/2}>0$. The expansion follows from $\frac1{2\sinh(u/2)}=u^{-1}-\frac u{24}+O(u^3)$,
$\frac1{2\cosh(u/2)}=\frac12+O(u^2)$ and $2e^{-u/2}=2-u+O(u^2)$. For \eqref{eq:coer-cross}, apply \eqref{eq:Aform} with
$\NT=\emptyset$ to $F+G$: $c_{F+G}(u)=\cF(u)+c_G(u)+\int F(t)G(t-u)dt+\int G(t)F(t-u)dt$; for $u>0$ the first cross
integral vanishes ($t\in B$, $t-u\in B'$ is impossible), while the second equals $\iint F(x)G(y)\mathbf 1_{y-x=u}$; integrating
against $K(u)du$ over $u>0$ gives \eqref{eq:coer-cross}, the term $\|F+G\|^2=\|F\|^2+\|G\|^2$ having no cross part.
Finally $2e^{-(y-x)/2}=2e^{-y/2}e^{x/2}$, so $\iint F(x)G(y)\,2e^{-(y-x)/2}dxdy=2\int Fe^{x/2}\int Ge^{-y/2}$, which vanishes
under either Tate condition.
\end{proof}
The cross term in \eqref{eq:coer-cross} carries the factor $1$ for the unordered pair $\{B,B'\}$; the polarized form is
$\mT(F,G)=-\frac12\iint F(x)G(y)K(y-x)\,dxdy$.

\begin{lemma}[a block sector is a short window]\label{lem:coer-sector}
Let $B\subset\IT$ be an interval of length $2r$ with $2r\le\log2$. Then
\[
\inf\Big\{\mathfrak a_T(F):\ F\in\mathcal D_T,\ \supp F\subset\overline B,\ \int_BFe^{\pm t/2}=0,\ \|F\|=1\Big\}=\lambda_{\mathrm A}(r),
\]
independently of $T$ and of the position of $B$; the infimum is attained. In particular, for $F\in\mathcal V_k$ with
$r=T/k$, $2r\le\log2$, and $F_j:=F\mathbf 1_{B_j}$,
\begin{equation}\label{eq:coer-diag}
\sum_{j=1}^k\mT(F_j)\ \ge\ \lambda_{\mathrm A}(r)\,\|F\|^2 .
\end{equation}
\end{lemma}
\begin{proof}
Let $B=I_r+a$. For $F$ supported in $\overline B$ and $F_a(t):=F(t+a)$, supported in $\overline{I_r}$, one has $c_{F_a}=\cF$,
hence $\mT(F_a)=\mT(F)$ by \eqref{eq:Aform}, and $\int F_ae^{\pm t/2}=e^{\mp a/2}\int Fe^{\pm t/2}$, so the two Tate
conditions on $B$ are those of $\mathcal P_r$ for $F_a$. Since $\cF(\log n)=0$ for $\log n\ge2r$ and $\log2\ge2r$, also
$\PiT(F)=0$, so $\mathfrak a_T(F)=\mT(F)=\mathfrak a_r(F_a)$ and the infimum is $\lambda_{\mathrm A}(r)$, attained since
$\lambda_{\mathrm A}(r)$ is. \eqref{eq:coer-diag} is the sum over the blocks, each $F_j$ satisfying the Tate conditions of
$B_j$, together with $\sum_j\|F_j\|^2=\|F\|^2$.
\end{proof}


\subsection{A lower bound for the short window}\label{subsec:coer-short}

\begin{lemma}\label{lem:coer-K}
$K(u)\ge u^{-1}$ for $0<u\le2$.
\end{lemma}
\begin{proof}
By Lemma~\ref{lem:K}, $K(u)-u^{-1}=\frac{u-2\sinh(u/2)}{2u\sinh(u/2)}+\frac1{2\cosh(u/2)}$. From $\tanh x\le x$,
$2\sinh(u/2)\le u\cosh(u/2)$, so $u-2\sinh(u/2)\ge-u(\cosh(u/2)-1)=-2u\sinh^2(u/4)$, and the first term is
$\ge-\frac{\sinh^2(u/4)}{\sinh(u/2)}=-\frac12\tanh\frac u4\ge-\frac u8$. Hence $K(u)-u^{-1}\ge\frac1{2\cosh(u/2)}-\frac u8\ge0$
whenever $u\cosh(u/2)\le4$, which holds for $u\le2$ ($2\cosh1=3.09$).
\end{proof}

\begin{lemma}[the tail]\label{lem:coer-tail}
For $r>0$, $\int_{2r}^\infty K(u)\,du\ge\log\frac2r+\frac\pi2-r$.
\end{lemma}
\begin{proof}
By Theorem~\ref{thm:S}(iii) the integral is $\log\coth\frac r2+\pi-2\arctan e^r$; $\coth x\ge x^{-1}$ bounds the first
term, and $\frac{d}{dr}(\pi-2\arctan e^r)=-\frac1{\cosh r}\ge-1$ with value $\frac\pi2$ at $r=0$ bounds the second.
\end{proof}

\begin{lemma}[the truncated logarithm]\label{lem:coer-w}
Put $\mathsf w(s):=\int_0^2\frac{1-\cos sv}v\,dv=\log(2|s|)+\gamma-\operatorname{Ci}(2|s|)$ for $s\ne0$, $\mathsf w(0)=0$. Then $\mathsf w$ is
even, increasing in $|s|$, and for every $0<R\le\pi$ and all $s\in\R$, $\mathsf w(s)\ge\mathsf w(R)\min\{1,s^2/R^2\}$.
\end{lemma}
\begin{proof}
$\mathsf w'(s)=\int_0^2\sin(sv)\,dv=\frac{1-\cos2s}s\ge0$ for $s>0$. Next $\mathsf w(s)/s^2=\int_0^2v\,\frac{1-\cos sv}{(sv)^2}dv$, and
$x\mapsto\frac{1-\cos x}{x^2}=\frac12\big(\frac{\sin(x/2)}{x/2}\big)^2$ is decreasing on $(0,2\pi)$; for $0<s\le R\le\pi$ and $v\le2$
one has $sv\le2\pi$, so $\mathsf w(s)/s^2\ge\mathsf w(R)/R^2$. For $s\ge R$ use monotonicity.
\end{proof}

\begin{lemma}[a double Poincar\'e inequality]\label{lem:coer-poinc}
Let $h\in L^2(\R)$ with $\supp h\subset[-1,1]$ and $h''\in L^2(\R)$. Then $\|h''\|^2\ge\pi^2\|h'\|^2\ge\frac{\pi^4}4\|h\|^2$.
\end{lemma}
\begin{proof}
$h,h''\in L^2(\R)$ give $h\in H^2(\R)\subset C^1(\R)$, and the support condition gives $h(\pm1)=h'(\pm1)=0$. The second
inequality is the Dirichlet Poincar\'e inequality for $h\in H^1_0(-1,1)$ (first eigenvalue $(\pi/2)^2$). For the first,
$g:=h'\in H^1_0(-1,1)$ with $\int_{-1}^1g=h(1)-h(-1)=0$. Expand $g=\sum_{n\ge1}c_n\phi_n$ in the orthonormal Dirichlet
eigenbasis $\phi_n(x)=\sin\frac{n\pi(x+1)}2$ of $(-1,1)$, so that $\|g\|^2=\sum c_n^2$ and $\|g'\|^2=\sum(\frac{n\pi}2)^2c_n^2$. Since
$\int\phi_n=\frac4{n\pi}$ for odd $n$ and $0$ for even $n$, the constraint reads $c_1=-\sum_{n\ge3\ \mathrm{odd}}c_n/n$. The
even-indexed terms satisfy $(\frac{n\pi}2)^2\ge\pi^2$, and for the odd-indexed ones
\[
\sum_{n\ \mathrm{odd}}\Big(\frac{n^2\pi^2}4-\pi^2\Big)c_n^2=\frac{\pi^2}4\Big[\sum_{n\ge3\ \mathrm{odd}}(n^2-4)c_n^2-3c_1^2\Big]\ \ge\ 0,
\]
because by Cauchy--Schwarz $c_1^2\le\mathcal S\sum_{n\ge3\ \mathrm{odd}}(n^2-4)c_n^2$ with
$\mathcal S=\sum_{n\ge3\ \mathrm{odd}}\frac1{n^2(n^2-4)}=\frac14\big(\frac43-\frac{\pi^2}8\big)=0.0249<\frac13$ (partial fractions and
$\sum_{n\ \mathrm{odd}}n^{-2}=\pi^2/8$). Hence $\|g'\|^2\ge\pi^2\|g\|^2$.
\end{proof}

\begin{theorem}[the short window]\label{thm:coer-A}
Let $0<r\le1$ with $2r\le\log2$, and put $a:=r^2/4$. For every $R\in(0,\pi]$,
\begin{equation}\label{eq:coer-A}
\lambda_{\mathrm A}(r)\ \ge\ C_0+\log\frac2r+\frac\pi2-r+\mathsf w(R)\Big[1-\frac{16(R^2+a)^3}{27\pi^4R^2}\Big].
\end{equation}
In particular, with $R=2$ ($\mathsf w(2)=\log4+\gamma-\operatorname{Ci}(4)=2.1045$),
$\lambda_{\mathrm A}(r)\ge\log\frac1r-1.209-r$ for $0<r\le0.1$, so $\lambda_{\mathrm A}(r)\ge\log\frac1r-1.31$ on $(0,0.1]$.
\end{theorem}
\begin{proof}
Let $F\in\mathcal P_r\cap\mathcal D_r$, $\|F\|=1$. Since $\mathcal N_r=\emptyset$ and $\cF(u)=0$ for $u\ge2r$, \eqref{eq:Aform} gives
\[
\mathfrak a_r(F)=C_0+\int_{2r}^\infty K+\int_0^{2r}\big(1-\Rel\cF(u)\big)K(u)\,du .
\]
The tail is bounded by Lemma~\ref{lem:coer-tail}. Put $G(y):=r^{1/2}F(ry)$, so that $\supp G\subset[-1,1]$, $\|G\|=1$,
$\cF(u)=c_G(u/r)$, and $\int_0^{2r}(1-\Rel\cF)K=\int_0^2(1-\Rel c_G(v))\,rK(rv)\,dv\ge\int_0^2\frac{1-\Rel c_G(v)}v\,dv$, using
$1-\Rel c_G\ge0$ and Lemma~\ref{lem:coer-K} ($rv\le2r\le2$). By Plancherel, $1-\Rel c_G(v)=\frac1{2\pi}\int(1-\cos sv)|\widehat G(s)|^2ds$,
so by Tonelli the last integral equals $\frac1{2\pi}\int\mathsf w|\widehat G|^2$, and by Lemma~\ref{lem:coer-w}, for $0<R\le\pi$,
\[
\frac1{2\pi}\int\mathsf w|\widehat G|^2\ \ge\ \mathsf w(R)\Big[1-\frac1{2\pi}\int\Big(1-\frac{s^2}{R^2}\Big)^+|\widehat G(s)|^2ds\Big].
\]
The Tate conditions $\int Fe^{\pm t/2}=0$ read $\int Ge^{\pm ry/2}=0$, i.e.\ $\widehat G(\mp ir/2)=0$. $\widehat G$ is entire of
exponential type $\le1$ and in $L^2(\R)$, so $H(s):=\widehat G(s)/(s^2+a)$ is entire of exponential type $\le1$ and in $L^2(\R)$
($s^2+a\ge a>0$ on $\R$); by the Paley--Wiener theorem $H=\widehat h$ with $h\in L^2(\R)$, $\supp h\subset[-1,1]$, and
$G=-h''+ah$, so $h''\in L^2$. Then
$1=\|G\|^2=\frac1{2\pi}\int(s^2+a)^2|H|^2=\|h''\|^2+2a\|h'\|^2+a^2\|h\|^2\ge\|h''\|^2\ge\frac{\pi^4}4\|h\|^2$
(Lemma~\ref{lem:coer-poinc}). On the other hand $(1-\frac{s^2}{R^2})^+(s^2+a)^2\le M_R:=\frac{4(R^2+a)^3}{27R^2}$ (the
maximum of $(1-x/R^2)(x+a)^2$ on $0\le x\le R^2$ is at $x=(2R^2-a)/3$), so
$\frac1{2\pi}\int(1-s^2/R^2)^+|\widehat G|^2\le M_R\|h\|^2\le\frac{4M_R}{\pi^4}$. Collecting terms gives \eqref{eq:coer-A}. For
$R=2$ and $a\le\frac1{400}$ the bracket is $\ge0.9025$ and $\mathsf w(2)\cdot0.9025\ge1.899$; with
$C_0+\log2+\frac\pi2=-3.1082$ this gives $\log\frac1r-1.209-r$.
\end{proof}

\begin{remark}\label{rem:coer-A}
(1) The only numerical input is the value $\mathsf w(2)$, an evaluation of an elementary function
($\operatorname{Ci}(4)=-0.1410$). Optimizing $R$ in \eqref{eq:coer-A} numerically gives $R\approx2.17$ and
$\log\frac1r-1.18-r$. (2) Numerically, $\lambda_{\mathrm A}(r)\approx\log\frac1r-0.5$ for $r\in[0.02,0.1]$ (Galerkin upper
bounds, not certified); the loss of the proof is in $\int_{|s|<R}|H|^2\le\|h\|^2$ and in the constant $\pi^4/4$.
(3) Theorem~\ref{thm:coer-A} complements Lemma~\ref{lem:contA} at the left end: $\lambda_{\mathrm A}(r)\to+\infty$ as $r\to0$ at
the rate $\log\frac1r$, consistently with Corollary~\ref{cor:mono}, which gives
$\lambda_{\mathrm A}(r)\le C_0+\varphi((3\pi/2r)^2)=\log\frac1r+O(1)$ (Theorem~\ref{thm:S}(i), Lemma~\ref{lem:mbounds}(ii)).
(4) The theorem is about $\lambda_{\mathrm A}(r)$, the bottom over $\mathcal P_r$ (two Tate conditions on $F$), not about
$\lambda^{\mathrm{prim}}_r$ of Definition~\ref{def:PD}; no statement about the latter is made here.
\end{remark}

\begin{remark}[precedents]\label{rem:coer-prec}
Coercivity of the Weil form on short windows in finite codimension is classical in qualitative form: Yoshida
\cite[Lemma~3]{Yoshida1992} proves that for given $a_0,\mu>0$ there is $N$ such that the Weil form is $\ge\mu\|\phi\|^2$ for
$\phi$ supported in $[-a,a]$, $0<a\le a_0$, orthogonal to $N$ Fourier modes (with no rate in $a$); Connes and Consani
\cite[Thm.~3.6, Cor.~3.8, Rem.~3.9]{ConnesConsani2021} write the archimedean part as $-2\cdot\mathrm{Id}$ plus a
Hilbert--Schmidt operator and obtain positivity on short windows after finitely many linear conditions. The leading
term $\log(1/r)$ of the short-window bottom has been observed numerically for Suzuki's operator without primes in
\cite{Kim2026}. Finite truncations of forms with a difference kernel are also studied in \cite{CvS2025}. What is new in
Theorems~\ref{thm:coer-A} and \ref{thm:coer-B} is the rate $\log(1/r)$ with an explicit constant, under the two Tate
conditions, and the block structure.
\end{remark}


\subsection{Coercivity of the archimedean form on $\mathcal V_k$}\label{subsec:coer-B}
For $R>0$ let $P_R$ be the orthogonal projection of $L^2(\R)$ onto the band $\{|s|\le R\}$,
$\widehat{P_RF}=\Fh\mathbf 1_{|s|\le R}$.

\begin{lemma}[low-pass leakage of $\mathcal V_k$]\label{lem:coer-lowpass}
Let $F\in\mathcal V_k$ with equal blocks of half-length $r$, and $R>0$. Then
\[
\|P_RF\|^2\le\varepsilon_{\mathrm{lp}}(Rr)\,\|F\|^2,\qquad \varepsilon_{\mathrm{lp}}(\theta):=\frac{\cosh^2r}{6}\Big(\theta^2+\frac{r^2}4\Big)^2 .
\]
\end{lemma}
\begin{proof}
$G:=P_RF$ is entire. Fix a block $B$ with centre $c$ and put, on $B$,
$u(x):=G(c)\cosh\frac{x-c}2+2G'(c)\sinh\frac{x-c}2\in\operatorname{span}\{e^{x/2},e^{-x/2}\}$, so that
$u(c)=G(c)$, $u'(c)=G'(c)$, $u''=u/4$. Then $\varrho:=G-u$ satisfies $\varrho(c)=\varrho'(c)=0$ and
$\varrho''-\frac\varrho4=G''-\frac G4=:H$, whence $\varrho(x)=\int_c^x2\sinh\frac{x-t}2\,H(t)\,dt$. By Cauchy--Schwarz and
$2\sinh\frac z2\le z\cosh\frac z2\le z\cosh r$ ($0\le z\le r$),
$|\varrho(x)|^2\le\cosh^2r\,\frac{|x-c|^3}3\,\|H\|_{L^2(B)}^2$, so $\int_B\varrho^2\le\cosh^2r\,\frac{r^4}6\|H\|_{L^2(B)}^2$.
Since $F_B:=F\mathbf 1_B$ is orthogonal to $e^{\pm x/2}\mathbf 1_B$, $\langle F_B,G\rangle=\langle F_B,\varrho\rangle\le\|F_B\|\,\|\varrho\|_{L^2(B)}$.
Summing over the blocks and using Cauchy--Schwarz once more,
\begin{align*}
\|G\|^2=\langle F,G\rangle&\le\|F\|\Big(\sum_B\|\varrho\|_{L^2(B)}^2\Big)^{1/2}\le\|F\|\cosh r\,\frac{r^2}{\sqrt6}\,\|G''-\tfrac14G\|_{L^2(\R)}\\
&\le\|F\|\cosh r\,\frac{r^2}{\sqrt6}\Big(R^2+\frac14\Big)\|G\|,
\end{align*}
the first equality because $P_R$ is an orthogonal projection, the last inequality because
$\widehat G$ is supported in $|s|\le R$. Squaring gives the claim, as $r^4(R^2+\frac14)^2=((Rr)^2+\frac{r^2}4)^2$.
\end{proof}

\begin{theorem}[coercivity on $\mathcal V_k$]\label{thm:coer-B}
Let $k\ge1$, $r=T/k\le1$, $\kappa:=\cosh^2r$, $a:=r^2/4$, $\varepsilon_{\mathrm{lp}}$ as in
Lemma~\ref{lem:coer-lowpass}, and $\theta_*:=(\sqrt{6/\kappa}-a)^{1/2}$ (the root of $\varepsilon_{\mathrm{lp}}(\theta_*)=1$); $m_{\mathrm{lo}}$ is
the lower bound of Lemma~\ref{lem:mbounds}(ii).
Then for every $F\in\mathcal V_k\cap\mathcal D_T$,
\begin{equation}\label{eq:coer-B}
\mT(F)\ \ge\ \Big(\log\frac1r-C_B(r)\Big)\|F\|^2,\qquad
C_B(r):=\log\frac1r-C_0\,\varepsilon_{\mathrm{lp}}(r)-\frac{2\kappa}3\big[\Phi_r(\theta_*)-\Phi_r(r)\big],
\end{equation}
where $\Phi_r$ is the elementary primitive
\[
\Phi_r(\theta):=\Big(\frac{\theta^4}4+\frac{a\theta^2}2\Big)\big(\log\theta-\log(2\pi r)\big)-\frac{\theta^4}{16}-\frac{a\theta^2}4
-r\Big(\frac{\theta^3}3+a\theta\Big)-\frac{r^2\theta^2}4-\frac{ar^2}2\log\theta
\]
of $(\theta^3+a\theta)\,m_{\mathrm{lo}}(\theta/r)$. Moreover $C_B(r)\to\frac14-\log\frac{6^{1/4}}{2\pi}=1.6399$ as
$r\to0$, and $C_B(0.02)=1.6573$, $C_B(0.0619)=1.6957$, $C_B(0.1)=1.7331$ (evaluations of the closed form).
\end{theorem}
\begin{proof}
Let $\|F\|=1$ and $\mu:=\frac1{2\pi}|\Fh|^2ds$, a probability measure, so that $\mT(F)=\int m\,d\mu$ by
\eqref{eq:form} with $\NT=\emptyset$. By Lemma~\ref{lem:mbounds}(i) and \eqref{eq:m}, $m$ is even, real-analytic,
increasing on $[0,\infty)$ with $m(0)=C_0$; hence $m(s)=C_0+\int_0^{|s|}m'(\rho)d\rho$ with $m'\ge0$ and, by Tonelli,
\[
\int m\,d\mu=C_0+\int_0^\infty m'(\rho)\,\mu(|s|\ge\rho)\,d\rho\ \ge\ C_0+\int_0^{\rho_*}m'(\rho)\big(1-\varepsilon_{\mathrm{lp}}(\rho r)\big)d\rho,
\qquad\rho_*:=\theta_*/r,
\]
since $\mu(|s|\ge\rho)=1-\|P_{\rho}F\|^2+\mu(|s|=\rho)\ge1-\varepsilon_{\mathrm{lp}}(\rho r)$ by Lemma~\ref{lem:coer-lowpass} and the
integrand is nonnegative on $[0,\rho_*]$, where $\varepsilon_{\mathrm{lp}}(\rho r)\le1$. Integrating by parts on $[0,\rho_*]$, with
vanishing boundary terms ($1-\varepsilon_{\mathrm{lp}}(\theta_*)=0$ and $m(0)-C_0=0$),
\[
\int_0^{\rho_*}m'(\rho)(1-\varepsilon_{\mathrm{lp}}(\rho r))d\rho=\int_0^{\theta_*}\big(m(\theta/r)-C_0\big)\varepsilon_{\mathrm{lp}}'(\theta)\,d\theta,
\]
so that
\[
\int m\,d\mu\ \ge\ C_0\,\varepsilon_{\mathrm{lp}}(0)+\int_0^{\theta_*}m(\theta/r)\,\varepsilon_{\mathrm{lp}}'(\theta)\,d\theta ,
\]
using $\int_0^{\theta_*}\varepsilon_{\mathrm{lp}}'=1-\varepsilon_{\mathrm{lp}}(0)$. Here $\varepsilon_{\mathrm{lp}}'(\theta)=\frac{2\kappa}3\theta(\theta^2+a)\ge0$. By
Lemma~\ref{lem:mbounds}(i),(ii), $m\ge\max\{C_0,m_{\mathrm{lo}}\}$ on $(0,\infty)$; we use $m\ge C_0$ on $[0,r]$ (i.e.\ $\rho\le1$)
and $m\ge m_{\mathrm{lo}}$ on $[r,\theta_*]$ (note $m_{\mathrm{lo}}(1)=-3.34>C_0$, and $m_{\mathrm{lo}}$ is increasing, so no further loss occurs
there). Hence
\[
\int m\,d\mu\ \ge\ C_0\,\varepsilon_{\mathrm{lp}}(0)+C_0\big(\varepsilon_{\mathrm{lp}}(r)-\varepsilon_{\mathrm{lp}}(0)\big)+\int_r^{\theta_*}m_{\mathrm{lo}}(\theta/r)\,\varepsilon_{\mathrm{lp}}'(\theta)\,d\theta
=C_0\,\varepsilon_{\mathrm{lp}}(r)+\frac{2\kappa}3\big[\Phi_r(\theta_*)-\Phi_r(r)\big],
\]
which is \eqref{eq:coer-B}; the primitive is checked by differentiation. As $r\to0$: $\kappa\to1$, $a\to0$,
$\theta_*\to6^{1/4}$, $\varepsilon_{\mathrm{lp}}(r)\to0$, $\frac{2\kappa}3\Phi_r(r)\to0$, and
$\frac{2\kappa}3\Phi_r(\theta_*)=\log\theta_*-\log(2\pi r)-\frac14+o(1)$ (the $a$- and $r$-terms are $o(1)$),
whence $C_B(r)\to\log2\pi-\log6^{1/4}+\frac14$. The three values are evaluations of the closed form (double precision).
\end{proof}

\begin{remark}\label{rem:coer-B}
(1) Using $m\ge\max\{C_0,m_{\mathrm{lo}}\}$ with the exact crossover instead of the split at $\rho=1$ changes $C_B$ by less
than $5\cdot10^{-5}$ on $(0,0.1]$ (numerical). (2) Numerically $C_B(r)\le1.74$ on $(0,0.1]$; monotonicity of $C_B$ is not
proved and is not used, since Corollary~\ref{cor:coer-L1} needs $C_B$ only at the specific value $r=T/k$. (3) The bound
loses about $1$ against the Galerkin value of the bottom of $\mT$ on $\mathcal V_k$ (about $1.1$ against $2.1$ at
$T=\frac12\log5$, $k=13$; numerical); the loss is in Lemma~\ref{lem:coer-lowpass}. (4) The theorem uses the whole
multiplier $m$ and the Tate conditions only through the geometry of $\mathcal V_k$; it does not separate the block-diagonal
part \eqref{eq:coer-diag} from the coupling \eqref{eq:coer-cross}. Remark~\ref{rem:coer-nogo} explains why that separation
is not viable.
\end{remark}




\begin{proposition}[the norm of the prime form]\label{prop:coer-prime}
\emph{(i)} $|\PiT(F)|\le\Aconst\|F\|^2$ for all $F\in L^2(\IT)$ and all $T$.
\emph{(ii)} For $\ell>0$ let $L(T,\ell):=\lceil2T/\ell\rceil$. Then $|\Rel\cF(\ell)|\le\cos\frac{\pi}{L(T,\ell)+1}\,\|F\|^2$,
hence
\[
\|\PiT\|\ \le\ \sum_{n\in\NT}v_n\cos\frac\pi{L(T,\log n)+1}
=\begin{cases}
\Aconst/2, & T\le\log2,\\[2pt]
\dfrac{v_2}{\sqrt2}+\dfrac{v_3+v_4}2=1.6740, & \log2<T\le\tfrac12\log5 .
\end{cases}
\]
Here $\frac{v_2}{\sqrt2}+\frac{v_3+v_4}2=\log2+\frac{\log3}{\sqrt3}+\frac{\log2}2$; in the first case the bound is
$\frac{v_2}2=0.4901$ for $T\le\frac12\log3$ and $\frac{v_2+v_3}2=1.1244$ for $\frac12\log3<T\le\log2$. The bounds hold on all of $L^2(\IT)$, in particular on $\mathcal V_k$.
\end{proposition}
\begin{proof}
(i) is $|\cF(\log n)|\le\|F\|^2$. (ii) Let $V_\ell$ be the shift $F\mapsto F(\cdot-\ell)\mathbf 1_{\IT}$ on
$L^2(\IT)$, so that $\Rel\cF(\ell)=\frac12\langle(V_\ell+V_\ell^*)F,F\rangle$. The map
$F\mapsto\big(F(-T+y+j\ell)\big)_{0\le j<n(y)}$, $y\in[0,\ell)$, $n(y):=\#\{j\ge0:-T+y+j\ell<T\}\le L(T,\ell)$, is unitary
from $L^2(\IT)$ onto $\int^\oplus_{[0,\ell)}\C^{n(y)}dy$, and it carries $V_\ell$ to the fibrewise shift $j\mapsto j+1$.
Hence $V_\ell+V_\ell^*$ acts on the fibre $\C^{n}$ as the adjacency matrix of the path with $n$ vertices, whose
eigenvalues are $2\cos\frac{j\pi}{n+1}$, $1\le j\le n$; its norm is $2\cos\frac\pi{n+1}\le2\cos\frac\pi{L+1}$. For
$T\le\log2$ one has $2T/\log n\le2$ for all $n\ge2$, so $L=2$ and $\cos\frac\pi3=\frac12$. For $\log2<T\le\frac12\log5$:
$\NT=\{2,3,4\}$, $L(T,\log2)=3$ since $2<2T/\log2\le\log5/\log2=2.32$, and $L(T,\log3)=L(T,\log4)=2$ since
$2T\le\log5<2\log3$; $\cos\frac\pi4=\frac1{\sqrt2}$. The values use $v_2=\sqrt2\log2$, $v_3=2\log3/\sqrt3$, $v_4=\log2$.
\end{proof}
For $\frac12\log2<T<\frac12\log3$ the bound $\Aconst/2$ is the exact norm: on a set of fibres of positive
measure the fibre has length $2$ and the prime form acts there as $-v_2\begin{psmallmatrix}0&1/2\\1/2&0\end{psmallmatrix}$.
Once $T>\log2$ the orbit of $\log2$ has length $3$ and the path norm exceeds $1$, so the bound
$\|\PiT\|\le\Aconst/2$ is specific to $T\le\log2$.

\begin{corollary}[coercivity of the truncated form on $\mathcal V_k$]\label{cor:coer-L1}
Let $k\ge1$ with $r=T/k\le1$, and put $c(T,k):=\inf\{\mathfrak a_T(F):F\in\mathcal V_k\cap\mathcal D_T,\ \|F\|=1\}$. Then
\begin{equation}\label{eq:coer-L1}
c(T,k)\ \ge\ \log\frac kT-C_B\Big(\frac Tk\Big)-\|\PiT\| ,
\end{equation}
with $\|\PiT\|$ bounded by Proposition~\ref{prop:coer-prime}. In particular (evaluations of the closed form):
\begin{center}
\begin{tabular}{@{}lllll@{}}
\toprule
$T$ & $k$ & $\dim E_k$ & bound for $\|\PiT\|$ & lower bound for $c(T,k)$\\
\midrule
$\frac12\log3$ & $8$ & $14$ & $0.4901$ & $0.486$\\
$\log2$ & $12$ & $22$ & $1.1244$ & $0.035$\\
$0.69$ & $14$ & $26$ & $1.1244$ & $0.202$\\
$\frac12\log5$ & $25$ & $48$ & $1.6740$ & $0.094$\\
$\frac12\log5$ & $26$ & $50$ & $1.6740$ & $0.134$\\
$\frac12\log5$ & $30$ & $58$ & $1.6740$ & $0.281$\\
$\frac12\log5$ & $85$ & $168$ & $\Aconst=2.9420$ & $0.070$\\
\bottomrule
\end{tabular}
\end{center}
Moreover, for every $T>0$ there is $k_0(T)$ such that $c(T,k)>0$ for all $k\ge k_0(T)$; indeed the right-hand side
of \eqref{eq:coer-L1} tends to $+\infty$ with $k$, because $C_B(T/k)\to1.6399$ (Theorem~\ref{thm:coer-B}).
\end{corollary}
\begin{proof}
$\mathfrak a_T=\mT+\PiT$ on $\mathcal D_T$, Theorem~\ref{thm:coer-B} and Proposition~\ref{prop:coer-prime}.
\end{proof}

\begin{remark}\label{rem:coer-size}
The size of $k_0(T)$: by \eqref{eq:coer-L1} and Proposition~\ref{prop:coer-prime}(i), $c(T,k)>0$ as soon as
$\log\frac kT>C_B(T/k)+\Aconst$, so any $k$ with $k>Te^{\Aconst+C_B(T/k)}$ is admissible, and $C_B(T/k)\to1.6399$ as
$k\to\infty$ (Theorem~\ref{thm:coer-B}). Since $\Aconst\asymp e^T$ by the prime number theorem, the dimension $2k-2$ of $E_k$
obtained in this way grows doubly exponentially in $T$; on the range $T\le\frac12\log5$ it stays below $60$
(Corollary~\ref{cor:coer-L1}). The Galerkin values of $c(T,k)$ exceed the bounds by about $1$: for instance
$c(\frac12\log5,13)\approx1.2$ (numerical), where \eqref{eq:coer-L1} gives a negative number.
\end{remark}




\subsection{The finite reduction}\label{subsec:coer-schur}

\begin{theorem}[reduction to the piecewise Tate exponentials]\label{thm:coer-schur}
Fix $T>0$ and $k\ge1$ with $c:=c(T,k)>0$ (Corollary~\ref{cor:coer-L1}: every $k\ge k_0(T)$). Put
$\mathcal V:=\mathcal V_k\cap\mathcal D_T$ and $M:=\Aconst-C_0$.
\emph{(i)} $\mathfrak a_T$ is an inner product on $\mathcal V$ for which $\mathcal V$ is complete, and for every
$e\in E_k$ the functional $v\mapsto\mathfrak a_T(e,v)$ is bounded on $(\mathcal V,\mathfrak a_T)$; hence there is a unique
$\mathsf Ke\in\mathcal V$ with $\mathfrak a_T(\mathsf Ke,v)=\mathfrak a_T(e,v)$ for all $v\in\mathcal V$. The map $\mathsf K:E_k\to\mathcal V$ is linear,
and $\|\mathsf K\|:=\sup\{\|\mathsf Ke\|:\|e\|=1\}<\infty$.
\emph{(ii)} For $v\in\mathcal V$ and $e\in E_k$,
\begin{equation}\label{eq:coer-schur}
\begin{split}
\mathfrak a_T(v+e)&=\mathfrak a_T(v+\mathsf Ke)+\mathfrak s_{T,k}(e),\\
\mathfrak s_{T,k}(e)&:=\mathfrak a_T(e)-\mathfrak a_T(\mathsf Ke)=\mathfrak a_T(e)-\sup_{u\in\mathcal V}\big[2\Rel\mathfrak a_T(e,u)-\mathfrak a_T(u)\big].
\end{split}
\end{equation}
Let $S_{T,k}$ be the matrix of the real quadratic form $\mathfrak s_{T,k}$ in an orthonormal basis of $E_k$ (a
real symmetric $(2k-2)\times(2k-2)$ matrix) and $\lambda_S$ its least eigenvalue. Then
\begin{equation}\label{eq:coer-sandwich}
\frac{\min\{c,\lambda_S\}}{(1+\|\mathsf K\|)^2}\ \le\ \lamA\ \le\ \min\{c,\lambda_S\}.
\end{equation}
\emph{(iii)} Consequently $\lamA>0$ if and only if $S_{T,k}\succ0$, and $\lamA\ge0$ if and only if $S_{T,k}\succeq0$.
\end{theorem}
\begin{proof}
(i) By Lemma~\ref{lem:mbounds}, $\sigma_T\ge C_0-\Aconst$, so $\mathfrak a_T+M\|\cdot\|^2\ge0$ on $\mathcal D_T$, and
by Lemma~\ref{lem:mbounds}(ii) and $m\ge C_0$ there are constants $c_1,c_2>0$ (depending on $T$) with
$c_1\int\log(2+|s|)|\Fh|^2\le\mathfrak a_T(F)+(M+1)\|F\|^2\le c_2\int\log(2+|s|)|\Fh|^2$ on $\mathcal D_T$. On $\mathcal V$,
$\mathfrak a_T\ge c\|\cdot\|^2$ gives $\mathfrak a_T(v)\ge\frac c{c+M+1}\big(\mathfrak a_T(v)+(M+1)\|v\|^2\big)$, so
$\mathfrak a_T$ is equivalent on $\mathcal V$ to the weighted norm $\int\log(2+|s|)|\widehat v|^2$, for which
$\mathcal D_T$ is complete and $\mathcal V_k\cap\mathcal D_T$ is closed ($\mathcal V_k$ is $L^2$-closed and the weighted norm
dominates the $L^2$ norm). For $e\in E_k\subset\mathcal D_T$ (Lemma~\ref{lem:coer-decomp}(iv)), the Cauchy--Schwarz
inequality for the nonnegative form $\mathfrak a_T+M\langle\cdot,\cdot\rangle$ gives
$|\mathfrak a_T(e,v)|\le(\mathfrak a_T(e)+M\|e\|^2)^{1/2}(\mathfrak a_T(v)+M\|v\|^2)^{1/2}+M\|e\|\|v\|\le C_e\,\mathfrak a_T(v)^{1/2}$,
so $\mathsf Ke$ exists by the Riesz representation theorem; $K$ is linear, and $\|\mathsf K\|<\infty$ since $\dim E_k<\infty$.
(ii) $\mathfrak a_T(v+\mathsf Ke)=\mathfrak a_T(v)+2\Rel\mathfrak a_T(\mathsf Ke,v)+\mathfrak a_T(\mathsf Ke)=\mathfrak a_T(v)+2\Rel\mathfrak a_T(e,v)+\mathfrak a_T(\mathsf Ke)$,
which rearranges to \eqref{eq:coer-schur}; the variational expression is $\mathfrak a_T(\mathsf Ke)=\sup_u[2\Rel\mathfrak a_T(\mathsf Ke,u)-\mathfrak a_T(u)]$
(maximum at $u=\mathsf Ke$). For \eqref{eq:coer-sandwich}: every $F\in\PT\cap\mathcal D_T$ is $v+e$ (Lemma~\ref{lem:coer-decomp}(iv)),
$\|F\|^2=\|v\|^2+\|e\|^2$, and $v\mapsto v+\mathsf Ke$ is a bijection of $\mathcal V$. Taking $e=0$ gives $\lamA\le c$; taking
$v=-\mathsf Ke$ gives $\mathfrak a_T(F)=\mathfrak s_{T,k}(e)$ with $\|F\|^2\ge\|e\|^2$, so $\lamA\le\lambda_S$. Conversely, if
$\lambda_S>0$, then $\mathfrak a_T(v+e)\ge c\|v+\mathsf Ke\|^2+\lambda_S\|e\|^2$, and
$\|v\|^2+\|e\|^2\le(\|v+\mathsf Ke\|+\|\mathsf K\|\|e\|)^2+\|e\|^2\le(1+\|\mathsf K\|)^2(\|v+\mathsf Ke\|^2+\|e\|^2)$, which gives the lower bound; if
$\lambda_S\le0$ the lower bound is implied by $\lamA\le\lambda_S$. (iii) follows from \eqref{eq:coer-sandwich}
and, for the nonstrict case, from $\mathfrak a_T(v+e)\ge\mathfrak s_{T,k}(e)$ together with $v=-\mathsf Ke$.
\end{proof}



\begin{remark}[what the reduction is and is not]\label{rem:coer-honest}
(1) Theorem~\ref{thm:coer-schur} is the Schur complement of $\mathfrak a_T$ with respect to
$\PT=\mathcal V_k\oplus E_k$, in the form-domain setting required because $\mathfrak a_T$ is unbounded on $L^2$.
Its content is a localization: the sign of $\lamA$ is the sign of an explicit real symmetric matrix of size
$2k-2$, with $k=25$ at $T=\frac12\log5$, indexed by the piecewise Tate exponentials $e^{\pm t/2}\mathbf 1_{B_j}$.
(2) There is no gain in precision. By \eqref{eq:coer-sandwich}, $\lamA\le\lambda_S\le(1+\|\mathsf K\|)^2\lamA$
whenever $\lambda_S\le c$, so $S_{T,k}$ inherits the decay of Theorem~\ref{thm:margin}; at $T=\frac12\log5$ its
least eigenvalue is of the size of $\lambda_{\mathrm A}(\frac12\log5)\approx3.10\cdot10^{-12}$ (numerical). The
entries of $S_{T,k}$ are not in closed form: the correction $\mathfrak a_T(\mathsf Ke)$ is a supremum over $\mathcal V$,
so a Galerkin computation in $\mathcal V$ gives upper bounds for $S_{T,k}$, and a certified lower bound needs a
near-maximizer together with a dual estimate, which is legitimate because $c>0$ is proved.
(3) Finite reductions of the window problem in sufficient (one-directional) form exist in Legendre bases
\cite[Thms.~1.1, 1.4]{Zhu2026}, \cite{Groskin2026}; Theorem~\ref{thm:coer-schur} is an exact reduction in an
explicit basis of explicit dimension, and it does not improve any certified value.
\end{remark}

\begin{remark}[the route by operator norms]\label{rem:coer-nogo}
By Lemmas~\ref{lem:coer-sector} and \ref{lem:coer-K1}, for $F\in\mathcal V_k$ with $2r\le\log2$,
$\mT(F)=\sum_j\mT(F_j)+\mathfrak q(F)$ with $\sum_j\mT(F_j)\ge\lambda_{\mathrm A}(r)\|F\|^2$ and the coupling form
$\mathfrak q(F)=-\sum_{i<j}\iint F_i(x)F_j(y)K_1(|x-y|)\,dxdy$ (the Tate conditions on each block allow $K_1$ in place of $K$),
so that $c(T,k)\ge\lambda_{\mathrm A}(r)-\sup_{\mathcal V_k}|\mathfrak q|/\|\cdot\|^2-\|\PiT\|$. This splitting cannot replace
Theorem~\ref{thm:coer-B}:
\begin{itemize}
\item[(a)] For two adjacent blocks, the bilinear form $(F_1,F_2)\mapsto\iint F_1(x)F_2(y)K_1(|x-y|)dxdy$ on Tate-orthogonal
functions has norm exactly $\pi$. It is $\le\pi$ by $0<K_1(u)\le u^{-1}$ and Hilbert's inequality
$\iint_{(0,\infty)^2}\frac{f(x)g(y)}{x+y}dxdy\le\pi\|f\|\|g\|$ \cite[Thm.~316]{HLP}, in the distances to the common endpoint. It is
$\ge\pi$ by the sequence $f_M(x')=x'^{-1/2}\mathbf 1_{[2re^{-M},2r]}(x')$ in the distance $x'$ to the common endpoint, placed on
both blocks and projected onto the orthogonal complement of $\{e^{\pm x/2}\mathbf 1_B\}$: in logarithmic variables
$\iint f_M(x')f_M(y')/(x'+y')=\int_{[-M,0]^2}\frac{du\,dv}{2\cosh\frac{u-v}2}=\pi M-O(1)$ while $\|f_M\|^2=M$; the bounded part
$K_1-u^{-1}=-\frac32+O(u)$ contributes $O(\|f_M\|_1^2)=O(r)$; and the projection changes $f_M$ by a vector of norm $O(1)$
uniformly in $M$, which changes the pairing by $O(\sqrt M)$. Hence, with the normalization of \eqref{eq:coer-cross},
$\sup_{\mathcal V_2}|\mathfrak q|/\|\cdot\|^2=\frac\pi2$, and $\sup_{\mathcal V_k}|\mathfrak q|/\|\cdot\|^2\ge\frac\pi2$ for every $k\ge2$ and every
$r$.
\item[(b)] With \eqref{eq:coer-A} and $\|\PiT\|\le1.674$ (Proposition~\ref{prop:coer-prime}), the split bound is positive at
$T=\frac12\log5$ only if $\log\frac1r>1.21+r+\frac\pi2+1.674$, which needs $k\ge70$ ($\dim E_k\ge138$), against $k=25$ from
Theorem~\ref{thm:coer-B}. Numerically the true coupling cost at the Galerkin minimizers of $\mT$ on $\mathcal V_k$ is only
about $0.2$: the extremals of the coupling norm are not the minimizers of the form.
\end{itemize}
\end{remark}


% ---------------------------------------------------------------------
% Section: the window form as a truncated Wiener--Hopf operator
% Insert after \section{A finite reduction} (sec:coer), before sec:firstcontact.
% ---------------------------------------------------------------------
\section{The window form as a truncated Wiener--Hopf operator}\label{sec:wh}

This section studies $\mathfrak a_T$ on $\PT$ as the difference of its positive and negative parts
(Proposition~\ref{prop:negtrace}), regarded as a one-parameter family of truncated convolution operators
whose inertia changes at the eigenvalues of a trace-class pencil. The name refers only to the shape
$R_Tf(D)R_T$ of the operators; no Wiener--Hopf index is used. Every result of the section is either an exact
reformulation of the sign of $\lamA$ or a bound with explicit constants, and several are barriers: they show
that a declared class of finite or determinantal criteria reproduces $\lamA$ up to a factor of order one.

\emph{Notation.} $R_T$ is multiplication by $\mathbf 1_{\IT}$ on $L^2(\R)$ and $f(D)$ the Fourier multiplier
$\widehat{f(D)F}=f\Fh$, so that $\langle R_Tf(D)R_TF,F\rangle=\frac1{2\pi}\int f|\Fh|^2$ for $F\in L^2(\IT)$ and
$\mathfrak a_T(F)=\langle R_T\sT(D)R_TF,F\rangle$ on $\mathcal D_T$. As in Proposition~\ref{prop:negtrace},
$\sT=\sigma_+-\sigma_-$, $\mathrm{Pos}_T=R_T\sigma_+(D)R_T$ (a closed form on $\mathcal D_T$) and
$\mathrm{Neg}_T=R_T\sigma_-(D)R_T$ (nonnegative, trace class on $L^2(\IT)$). For measurable $E\subset\R$,
$S_E:=R_T\mathbf 1_E(D)R_T$ is the time--band limiting operator, $0\le S_E\le1$; for $E=[-\tau_1,\tau_1]$ it is
$S_{T,\tau_1}$ of Proposition~\ref{prop:NC}, and $\operatorname{tr}S_E=T|E|/\pi$ when $|E|<\infty$ (the kernel is
$\frac1{2\pi}\int_Ee^{i\tau(t-s)}d\tau$). If $E$ has positive measure, $S_E$ is injective and $\|S_E\|<1$ when
$|\R\setminus E|>0$: an $F$ with $\Fh=0$ a.e.\ on a set of positive measure vanishes, $\Fh$ being entire. By
Proposition~\ref{prop:band}, $U_T:=\{\sT<0\}$ is a finite union of bounded open intervals, symmetric under
$\tau\mapsto-\tau$ and containing $0$; its connected components are $E_0=(-a_0,a_0)$ and
$E_i=\Delta_i\cup(-\Delta_i)$, $1\le i<n_T$, with $\Delta_i\subset(0,\infty)$ disjoint open intervals, and
$s_i:=\sup_{E_i}\sigma_-$. Let $\mathrm{pr}_T$ be the orthogonal projection of $L^2(\IT)$ onto $\PT$; $J$
preserves $\PT$ and commutes with $\mathrm{pr}_T$ and with every $R_Tf(D)R_T$ with $f$ even. By
Lemma~\ref{lem:compact}, $\lamA$ is attained; $h^*_T$ denotes a unit minimizer. For a closed lower-bounded form
$\mathfrak q$ on a subspace, $n_-(\mathfrak q)$ is the largest dimension of a subspace of its domain on which it
is negative definite, as in Lemma~\ref{lem:parity}.

\subsection{The pencil and its inertia crossings}\label{subsec:wh-pencil}

\begin{lemma}[the positive part is coercive]\label{lem:wh-cP}
For $\delta>0$ let $c_\delta:=\min\{\sT(\tau):|\tau|\ge\tau_{\mathrm{last}}+\delta\}$ and
$\lambda_0(\delta):=\|S_{T,\tau_{\mathrm{last}}+\delta}\|$. Then $c_\delta>0$, $\lambda_0(\delta)<1$, and
$\mathrm{Pos}_T(F)\ge c_\delta(1-\lambda_0(\delta))\|F\|^2$ for $F\in\mathcal D_T$. Hence
$c_P:=\inf\{\mathrm{Pos}_T(F):F\in\PT\cap\mathcal D_T,\ \|F\|=1\}>0$, and the closed form $\mathrm{Pos}_T|_{\PT\cap\mathcal D_T}$
is the form of a self-adjoint operator $\mathfrak p\ge c_P$ on $\PT$; $\mathfrak p^{1/2}$ maps $\PT\cap\mathcal D_T$
bijectively onto $\PT$.
\end{lemma}

\begin{proof}
$\sT>0$ for $|\tau|>\tau_{\mathrm{last}}$ (Proposition~\ref{prop:band}(iv)), $\sT$ is continuous and
$\sT\ge m_{\mathrm{lo}}-\Aconst\to+\infty$ (Lemma~\ref{lem:mbounds}), so the minimum exists and is positive;
$\lambda_0(\delta)<1$ by the remark in the notation paragraph. Since
$\sigma_+\ge c_\delta\mathbf 1_{|\tau|\ge\tau_{\mathrm{last}}+\delta}$,
$\mathrm{Pos}_T(F)\ge c_\delta(\|F\|^2-\langle S_{T,\tau_{\mathrm{last}}+\delta}F,F\rangle)$. The form is closed on
$\mathcal D_T$: $\sigma_+\le\sT+\sigma_-$ with $\sigma_-$ bounded, and $\sigma_+=\sT\ge\frac12\log(2+|\tau|)$ for large
$|\tau|$, so $\mathrm{Pos}_T+\|\cdot\|^2$ is equivalent to the weighted norm $\int\log(2+|\tau|)|\Fh|^2$ that defines
$\mathcal D_T$. The last statement is the representation theorem and the spectral theorem.
\end{proof}

\begin{definition}\label{def:wh-B}
Let $\mathfrak n:=\mathrm{pr}_T\mathrm{Neg}_T\mathrm{pr}_T|_{\PT}$ (nonnegative, trace class), and
\[
B_T:=\mathfrak p^{-1/2}\,\mathfrak n\,\mathfrak p^{-1/2}\quad\text{on }\PT,\qquad
\beta_1\ge\beta_2\ge\dots\downarrow0\ \text{its eigenvalues},
\]
so that $B_T\ge0$ is trace class with $\operatorname{tr}B_T\le\operatorname{tr}\mathrm{Neg}_T/c_P$. For $\beta>0$ let
$\mathfrak w_\beta:=\mathrm{Pos}_T-\beta^{-1}\mathrm{Neg}_T$ on $\PT\cap\mathcal D_T$, the compression to $\PT$ of the
truncated convolution operator with symbol $\sigma_+-\sigma_-/\beta$; thus $\mathfrak w_1=\mathfrak a_T|_{\PT}$.
\end{definition}

\begin{theorem}[inertia crossings]\label{thm:wh-pencil}
\emph{(i)} For $u\in\PT\cap\mathcal D_T$ and $\beta>0$,
$\mathfrak w_\beta(u)=\langle(1-\beta^{-1}B_T)\mathfrak p^{1/2}u,\mathfrak p^{1/2}u\rangle$. Hence
$n_-(\mathfrak w_\beta)=\#\{j:\beta_j>\beta\}<\infty$, a nonincreasing function of $\beta$ that jumps exactly at the
$\beta_j$, by their multiplicities. In particular
\[
n_-(\mathfrak a_T|\PT)=\#\{j:\beta_j>1\},\qquad \lamA>0\iff\beta_1<1,\qquad \lamA\ge0\iff\beta_1\le1 .
\]
\emph{(ii)} $1-\beta_1=\inf\{\mathfrak a_T(u)/\mathrm{Pos}_T(u):u\in\PT\cap\mathcal D_T\setminus0\}$, and
$\beta_1=\sup\mathrm{Neg}_T(u)/\mathrm{Pos}_T(u)$ is attained at the unit vector $u_1\in\PT\cap\mathcal D_T$ proportional
to $\mathfrak p^{-1/2}v_1$, $v_1$ an eigenvector of $\beta_1$.
\emph{(iii)} If $\lamA\ge0$, then
\begin{equation}\label{eq:wh-sandwich}
\frac{\lamA}{\mathrm{Pos}_T(u_1)}\ \le\ 1-\beta_1\ \le\ \frac{\lamA}{\mathrm{Pos}_T(h^*_T)}\ \le\ \frac{\lamA}{c_P}.
\end{equation}
\emph{(iv)} $n_-(\mathfrak a_T|\PT)\le\#\{j:\nu_j(\mathfrak n)>c_P\}\le\#\{j:\nu_j(\mathrm{Neg}_T)>c_P\}$, where $\nu_j(\cdot)$ are
eigenvalues in decreasing order.
\end{theorem}

\begin{proof}
(i) $\mathrm{Pos}_T(u)=\|\mathfrak p^{1/2}u\|^2$ and $\mathrm{Neg}_T(u)=\langle\mathfrak nu,u\rangle=\langle B_T\mathfrak p^{1/2}u,\mathfrak p^{1/2}u\rangle$.
Since $\mathfrak p^{1/2}$ is a bijection of $\PT\cap\mathcal D_T$ onto $\PT$, subspaces on which $\mathfrak w_\beta$ is
negative definite correspond to subspaces on which $1-\beta^{-1}B_T$ is, and for compact self-adjoint $B_T$ the
maximal dimension of the latter is $\#\{\beta_j>\beta\}$ by min--max. For $\beta=1$: if $\beta_1<1$ then
$\mathfrak a_T(u)\ge(1-\beta_1)\mathrm{Pos}_T(u)\ge(1-\beta_1)c_P\|u\|^2$; if $\beta_1\ge1$ then
$\mathfrak a_T(\mathfrak p^{-1/2}v_1)=(1-\beta_1)\|v_1\|^2\le0$; the case $\ge$ is the same.
(ii) $\mathfrak a_T/\mathrm{Pos}_T=1-\mathrm{Neg}_T/\mathrm{Pos}_T$ and $\sup_u\mathrm{Neg}_T(u)/\mathrm{Pos}_T(u)=\sup_v\langle B_Tv,v\rangle/\|v\|^2=\beta_1$,
attained at $v_1$.
(iii) $\mathfrak a_T(u_1)=(1-\beta_1)\mathrm{Pos}_T(u_1)$ and $\mathfrak a_T(u_1)\ge\lamA$ give the left inequality.
Since $\mathrm{Neg}_T\le\beta_1\mathrm{Pos}_T$ on $\PT\cap\mathcal D_T$, $\lamA=\mathfrak a_T(h^*_T)\ge(1-\beta_1)\mathrm{Pos}_T(h^*_T)$,
which is the middle one; finally $\mathrm{Pos}_T(h^*_T)\ge c_P$ and $\lamA\ge0$.
(iv) $\mathfrak a_T(u)\ge c_P\|u\|^2-\langle\mathfrak nu,u\rangle$; if $\mathfrak a_T<0$ on an $n$-dimensional subspace, then
$\langle\mathfrak nu,u\rangle>c_P\|u\|^2$ there and $\nu_n(\mathfrak n)>c_P$ by min--max. Finally
$\nu_j(\mathfrak n)\le\nu_j(\mathrm{Neg}_T)$ because $\mathfrak n$ is a compression of $\mathrm{Neg}_T$.
\end{proof}

\begin{remark}\label{rem:wh-pencil}
Theorem~\ref{thm:wh-pencil} is Sylvester's law of inertia for the pair $(\mathrm{Pos}_T,\mathrm{Neg}_T)$. The
question whether $\lamA>0$ is the question whether the last inertia crossing $\beta_1$ of the family
$\beta\mapsto\mathfrak w_\beta$ lies below $1$: an extremal-eigenvalue question, not a counting one. By
\eqref{eq:wh-sandwich}, $1-\beta_1$ is $\lamA$ divided by a number in $[c_P,\mathrm{Pos}_T(u_1)]$; this is the
first instance of the inheritance that recurs below. Numerically, $\mathrm{Pos}_T(h^*_T)\approx0.2$ at
$T=\frac12\log3$ and $T=\frac12\log5$, $1-\beta_1$ agrees with $\lamA/\mathrm{Pos}_T(h^*_T)$ to the precision of the
computation, and the bound (iv) gives $n_-(\mathfrak a_T|\PT)\le4$ and $\le10$ at these windows (Galerkin, not
certified).
\end{remark}

\subsection{Constant weights and principal angles}\label{subsec:wh-angles}

For two bands with constant weights the pencil is a function of the principal angles between the
subspace of functions supported in $\IT$ and the band subspaces, in the sense of Halmos's theory of two
subspaces \cite{Halmos1969}. The following statements are on $L^2(\IT)$, without the Tate conditions.

\begin{theorem}[two bands with constant weights]\label{thm:wh-angles}
Let $\Omega_-,\Omega_+\subset\R$ be disjoint measurable sets with $|\Omega_-\cup\Omega_+|>0$, let $s_\pm>0$,
$\mathsf Q_\pm:=\mathbf 1_{\Omega_\pm}(D)$ and $\mathsf Q:=\mathsf Q_-+\mathsf Q_+$ on $L^2(\R)$, and put
$\mathfrak n_0(u):=s_-\|\mathsf Q_-u\|^2$, $\mathfrak p_0(u):=s_+\|\mathsf Q_+u\|^2$ for $u\in L^2(\IT)$. For
$v\in L^2(\R)\setminus0$ let $\phi(v)\in[0,\frac\pi2]$ be the angle between $v$ and $\operatorname{ran}\mathsf Q_-$,
$\cos^2\phi(v)=\|\mathsf Q_-v\|^2/\|v\|^2$. Then $u\mapsto\mathsf Qu$ is injective on $L^2(\IT)$ and, for $u$ with
$\mathsf Q_+u\ne0$,
\[
\frac{\mathfrak n_0(u)}{\mathfrak p_0(u)}=\frac{s_-}{s_+}\cot^2\phi(\mathsf Qu).
\]
Consequently, for a subspace $X\subset L^2(\IT)$ on which $\mathfrak p_0>0$,
$\sup_X\mathfrak n_0/\mathfrak p_0=\frac{s_-}{s_+}\cot^2\phi_{\min}$, with $\phi_{\min}$ the smallest angle between
$\mathsf QX$ and $\operatorname{ran}\mathsf Q_-$, and $n_-\big((\mathfrak p_0-\mathfrak n_0)|X\big)$ is the largest dimension
of a subspace $Y\subset\mathsf QX$ with $\|\mathsf Q_-v\|^2>\theta\|v\|^2$ for $v\in Y\setminus0$,
$\theta:=s_+/(s_++s_-)$.
\end{theorem}

\begin{proof}
If $\mathsf Qu=0$, then $\widehat u=0$ a.e.\ on $\Omega_-\cup\Omega_+$, so $u=0$. For $v=\mathsf Qu$ one has
$\mathsf Q_\pm v=\mathsf Q_\pm u$ and $v=\mathsf Q_-v+\mathsf Q_+v$ orthogonally, so $\|\mathsf Q_+v\|^2=\sin^2\phi(v)\|v\|^2$ and
$\mathfrak n_0(u)/\mathfrak p_0(u)=s_-\|\mathsf Q_-v\|^2/(s_+\|\mathsf Q_+v\|^2)=(s_-/s_+)\cot^2\phi(v)$. The supremum statement
follows since $\cot^2$ is decreasing on $(0,\frac\pi2]$, and the negativity condition
$s_+\|\mathsf Q_+v\|^2<s_-\|\mathsf Q_-v\|^2$ is $(s_++s_-)\|\mathsf Q_-v\|^2>s_+\|v\|^2$.
\end{proof}

\begin{corollary}[complementary bands]\label{cor:wh-nogap}
If $\Omega_+=\R\setminus\Omega_-$ with $0<|\Omega_-|<\infty$, then on $L^2(\IT)$ the operators of $\mathfrak n_0$ and
$\mathfrak p_0$ are $s_-S_{\Omega_-}$ and $s_+(1-S_{\Omega_-})$; they commute, and with $\lambda_0\ge\lambda_1\ge\dots$ the
eigenvalues of $S_{\Omega_-}$ the pencil eigenvalues are $\beta_j=\frac{s_-}{s_+}\frac{\lambda_j}{1-\lambda_j}$. So
$n_-(\mathfrak p_0-\mathfrak n_0)=\#\{j:\lambda_j>\theta\}$, and $\mathfrak p_0-\mathfrak n_0>0$ on $L^2(\IT)$ iff
$s_-\lambda_0<s_+(1-\lambda_0)$. On $\PT$ the pencil eigenvalues interlace:
$\beta_j(L^2(\IT))\ge\beta_j(\PT)\ge\beta_{j+2}(L^2(\IT))$.
\end{corollary}

\begin{proof}
$\lambda_0<1$, so $1-S_{\Omega_-}$ is invertible and the two operators have a common eigenbasis, on which the
ratio is $s_-\lambda_j/(s_+(1-\lambda_j))$. The interlacing is min--max for the ratio of forms restricted to a
subspace of codimension $2$.
\end{proof}

\begin{proposition}[the sign is not a function of the negative part]\label{prop:wh-notangles}
For $\varepsilon>0$ let $\sigma^{(\varepsilon)}:=\varepsilon\sigma_+-\sigma_-$. Then $\{\sigma^{(\varepsilon)}<0\}=U_T$,
$\sigma^{(\varepsilon)}_-=\sigma_-$, the operator $\mathrm{Neg}_T$ and all principal angles between $L^2(\IT)$ and
$\operatorname{ran}\mathbf 1_{U_T}(D)$ are the same for all $\varepsilon$, while the top pencil eigenvalue of
$\sigma^{(\varepsilon)}$ on $\PT$ is $\beta_1/\varepsilon$, and $\beta_1>0$. Hence the form of $\sigma^{(\varepsilon)}$ on
$\PT$ has a negative direction for $\varepsilon<\beta_1$ and is positive definite for $\varepsilon>\beta_1$.
\end{proposition}

\begin{proof}
$\mathrm{Pos}^{(\varepsilon)}_T=\varepsilon\mathrm{Pos}_T$ and $\mathrm{Neg}^{(\varepsilon)}_T=\mathrm{Neg}_T$, so
$B^{(\varepsilon)}_T=B_T/\varepsilon$; apply Theorem~\ref{thm:wh-pencil}(i). $\mathrm{Neg}_T$ is injective on $L^2(\IT)$
($\langle\mathrm{Neg}_Tu,u\rangle=0$ forces $\widehat u=0$ on the open set $U_T$), so $\mathfrak n\ne0$ and $\beta_1>0$.
\end{proof}

\subsection{The Kre\u{\i}n formula for a rank-one negative part}\label{subsec:wh-krein}

\begin{lemma}[Kre\u{\i}n]\label{lem:wh-krein}
Let $\mathfrak p\ge c>0$ be a self-adjoint operator on a Hilbert space $H$, let $\mathfrak n\ge0$ be bounded with
$\mathfrak ne=\nu e$, $\|e\|=1$, $\nu>0$, and $\beta_1:=\|\mathfrak p^{-1/2}\mathfrak n\mathfrak p^{-1/2}\|$.
\emph{(i)} $\beta_1\ge\nu\langle\mathfrak p^{-1}e,e\rangle$, with equality if $\mathfrak n=\nu\,e\otimes e$. In particular
$\nu\langle\mathfrak p^{-1}e,e\rangle<1$ is necessary for $\mathfrak p-\mathfrak n>0$; it is sufficient only when
$\mathfrak n$ has rank one.
\emph{(ii)} $\langle\mathfrak p^{-1}e,e\rangle=\sup_{u\ne0}|\langle u,e\rangle|^2/\mathfrak p(u)=\big(\inf_{u\perp e}\mathfrak p(e+u)\big)^{-1}$.
\emph{(iii)} If $\mathfrak n=\nu\,e\otimes e$, then $\mathfrak p-\mathfrak n>0$ iff $\nu\langle\mathfrak p^{-1}e,e\rangle<1$ iff
$\nu<\inf_{u\perp e}\mathfrak p(e+u)$.
\end{lemma}

\begin{proof}
(i) $\mathfrak n-\nu\,e\otimes e\ge0$, since it vanishes on $e$ and equals $\mathfrak n$ on $e^\perp$, which $\mathfrak n$
preserves. Hence $\mathfrak p^{-1/2}\mathfrak n\mathfrak p^{-1/2}\ge\nu\,(\mathfrak p^{-1/2}e)\otimes(\mathfrak p^{-1/2}e)$, whose norm is
$\nu\|\mathfrak p^{-1/2}e\|^2$; with equality in the rank-one case. The necessity follows from
Theorem~\ref{thm:wh-pencil}(i) in this abstract setting ($\mathfrak p-\mathfrak n>0$ iff $\beta_1<1$).
(ii) $|\langle u,e\rangle|^2=|\langle\mathfrak p^{1/2}u,\mathfrak p^{-1/2}e\rangle|^2\le\mathfrak p(u)\langle\mathfrak p^{-1}e,e\rangle$, with equality at
$u=\mathfrak p^{-1}e$; and $\{w:\langle w,e\rangle=1\}=e+e^\perp$. (iii) $\mathfrak p-\nu e\otimes e=\mathfrak p^{1/2}(1-\nu\,f\otimes f)\mathfrak p^{1/2}$ with
$f=\mathfrak p^{-1/2}e$, and $1-\nu f\otimes f>0$ iff $\nu\|f\|^2<1$; the second equivalence is (ii).
\end{proof}

\begin{corollary}[one negative interval on a constant background]\label{cor:wh-step}
Let $E=(-W,W)$, $a,b_0>0$, and $\lambda_0\ge\lambda_1\ge\dots$ the eigenvalues of $S_E=S_{T,W}$. For the symbol
$b_0\mathbf 1_{\R\setminus E}-a\mathbf 1_E$ the pencil eigenvalues on $L^2(\IT)$ are $\frac a{b_0}\frac{\lambda_j}{1-\lambda_j}$, and
the form is positive iff $a\lambda_0<b_0(1-\lambda_0)$. If $b\ge b_0$ on $\R\setminus E$ and $e$ is a unit top
eigenvector of $S_E$, then
\[
b_0(1-\lambda_0)\ \le\ \inf_{u\perp e}\ \langle R_T(b\mathbf 1_{\R\setminus E})(D)R_T(e+u),e+u\rangle\ \le\ \frac1{2\pi}\int_{\R\setminus E}b\,|\widehat e|^2 .
\]
\end{corollary}

\begin{proof}
The first part is Corollary~\ref{cor:wh-nogap}. For the second, the form of $b\mathbf 1_{\R\setminus E}$ dominates
$b_0(1-S_E)$, for which the infimum equals $b_0(1-\lambda_0)$ by Lemma~\ref{lem:wh-krein}(ii), $e$ being an
eigenvector; the right inequality is $u=0$.
\end{proof}

\begin{remark}\label{rem:wh-model}
In the model of Corollary~\ref{cor:wh-step} the sign is decided by one number per component. For $\sT$ this
mechanism is absent: $\sigma_+$ vanishes at $\partial U_T$, so the constant $b_0$ is $0$ for every component.
Numerically, the components of $U_T$ have Shannon numbers $T|\Delta_i|/\pi$ below $1$ at the windows
computed, $\|S_{U_T}\|$ stays below about $0.9$, and models of $\sT$ that are constant on each component give
values of $1-\beta_1$ of either sign and of order $10^{-2}$ to $1$ (Galerkin, not certified). So the smallness of
$1-\beta_1$ is not a plateau phenomenon of a time--band operator, and the asymptotics of Slepian, Landau and
Widom for such operators are not used.
\end{remark}

\subsection{The prolate block and coercivity on its complement}\label{subsec:wh-prolate}

\begin{definition}\label{def:wh-V}
For each component $E_i$ of $U_T$ and each parity $\pm$, let $\lambda^{(i),\pm}_0\ge\lambda^{(i),\pm}_1\ge\dots$ be
the eigenvalues of $S_{E_i}$ on the even ($+$) and odd ($-$) subspaces of $L^2(\IT)$ ($S_{E_i}$ commutes with $J$
because $E_i$ is symmetric), and $e^\pm_i$ a unit eigenvector of $\lambda^{(i),\pm}_0$. Put
\[
\mathcal E_T:=\operatorname{span}\{e^\pm_i\},\qquad \mathcal E^P_T:=\mathrm{pr}_T\mathcal E_T\subset\PT,\qquad
\mathcal X_T:=\PT\ominus\mathcal E^P_T ,
\]
and let $\mathsf M_T:=\mathfrak a_T|_{\mathcal E^P_T}$ (the \emph{prolate block}, a real symmetric form on a space of
dimension at most $2n_T$). Variant: $\mathcal E'_T:=\mathcal E_T+\operatorname{span}\{f^\pm_i\}$, with $f^\pm_i$ a unit top
eigenvector of $\mathrm{Neg}_{T,i}:=R_T(\sigma_-\mathbf 1_{E_i})(D)R_T$ in the parity sector $\pm$, and
$\mathcal E'^P_T$, $\mathcal X'_T$ accordingly.
\end{definition}

\begin{lemma}\label{lem:wh-W}
\emph{(i)} $\mathcal E_T,\mathcal E'_T,\mathcal E^P_T\subset\mathcal D_T$ and $\mathcal X_T=\PT\cap\mathcal E_T^\perp$;
$\PT\cap\mathcal D_T=(\mathcal X_T\cap\mathcal D_T)\oplus\mathcal E^P_T$, orthogonally in $L^2(\IT)$. \emph{(ii)} These spaces are
$J$-invariant.
\end{lemma}

\begin{proof}
(i) If $Ke=\mu e$ with $\mu>0$ and $K=R_Tg(D)R_T$, $g$ bounded with compact support, then $e=\mu^{-1}R_T(g\widehat e)^\vee$
is the restriction to $\IT$ of an entire function, hence smooth on $\overline{\IT}$, and integration by parts gives
$\widehat e(\tau)=O(1/|\tau|)$, so $e\in\mathcal D_T$. The eigenvalues are positive because $S_{E_i}$ and
$\mathrm{Neg}_{T,i}$ are injective. $\mathrm{pr}_Te$ differs from $e$ by a combination of $e^{\pm t/2}\mathbf 1_{\IT}\in\mathcal D_T$.
For $F\in\PT$, $\langle F,\mathrm{pr}_Tv\rangle=\langle F,v\rangle$, which gives $\mathcal X_T=\PT\cap\mathcal E_T^\perp$; the
decomposition is $F=(F-\pi F)+\pi F$, $\pi$ the orthogonal projection onto the finite-dimensional
$\mathcal E^P_T\subset\mathcal D_T$. (ii) The $e^\pm_i,f^\pm_i$ are $J$-eigenvectors and $J$ commutes with $\mathrm{pr}_T$.
\end{proof}

\begin{theorem}[the prolate block is necessary]\label{thm:wh-necessary}
$\lamA\le\lambda_{\min}(\mathsf M_T)$, $n_-(\mathsf M_T)\le n_-(\mathfrak a_T|\PT)$, and
$\sup_{\mathcal E^P_T}\mathrm{Neg}_T/\mathrm{Pos}_T\le\beta_1$. In particular, if $\mathsf M_T$ has a negative eigenvalue then
$\lamA<0$.
\end{theorem}

\begin{proof}
Rayleigh quotients on the subspace $\mathcal E^P_T\subset\PT\cap\mathcal D_T$, and Theorem~\ref{thm:wh-pencil}(ii).
\end{proof}

\begin{lemma}[the negative part on the complement]\label{lem:wh-neg}
For $F\in\mathcal X_T$,
\[
\mathrm{Neg}_T(F)\le N_{\mathrm{ub}}\|F\|^2,\qquad
N_{\mathrm{ub}}:=\min\Big\{\max_\pm\sum_is_i\lambda^{(i),\pm}_1,\ \Big\|\sum_is_i\widetilde S_{E_i}\Big\|\Big\},
\]
where $\widetilde S_{E_i}:=S_{E_i}-\sum_\pm\lambda^{(i),\pm}_0e^\pm_i\otimes e^\pm_i\ge0$.
For $F\in\mathcal X'_T$ the same holds with $s_i\lambda^{(i),\pm}_1$ replaced by the second eigenvalue of
$\mathrm{Neg}_{T,i}$ in the sector $\pm$.
\end{lemma}

\begin{proof}
The $E_i$ are disjoint, so $\mathrm{Neg}_T=\sum_i\mathrm{Neg}_{T,i}$, and $\sigma_-\mathbf 1_{E_i}\le s_i\mathbf 1_{E_i}$ gives
$\mathrm{Neg}_{T,i}\le s_iS_{E_i}$. Write $F=F_++F_-$ by parity; $F_\pm\in\mathcal X_T$ by Lemma~\ref{lem:wh-W}(ii), and
$\langle S_{E_i}F,F\rangle=\sum_\pm\langle S_{E_i}F_\pm,F_\pm\rangle$ since $S_{E_i}$ commutes with $J$. In the sector $\pm$,
$F_\pm\perp e^\pm_i$, so $\langle S_{E_i}F_\pm,F_\pm\rangle\le\lambda^{(i),\pm}_1\|F_\pm\|^2$; summing over $i$ and $\pm$ gives the
first bound. For $F\perp\mathcal E_T$, $\langle S_{E_i}F,F\rangle=\langle\widetilde S_{E_i}F,F\rangle$, and $\widetilde S_{E_i}\ge0$ since the
$e^\pm_i$ are orthonormal eigenvectors of $S_{E_i}$; this gives the second. The variant is the same argument with
$\mathrm{Neg}_{T,i}$ in place of $s_iS_{E_i}$.
\end{proof}

\begin{lemma}[the positive part on the complement, by sublevel layers]\label{lem:wh-pos}
For $s\ge0$ let $L_s:=\{\tau:\sT(\tau)\le s\}$, a bounded finite union of closed intervals containing every $E_i$,
nondecreasing in $s$, and put
\[
\widetilde S_{L_s}:=S_{L_s}-\sum_i\sum_\pm\lambda^{(i),\pm}_0\,e^\pm_i\otimes e^\pm_i,\qquad \bar\mu(s):=\|\widetilde S_{L_s}\| .
\]
Then $\widetilde S_{L_s}\ge0$, $\bar\mu$ is nondecreasing, and for $F\in\mathcal X_T\cap\mathcal D_T$ with $\|F\|=1$ and every
partition $0=s_0<s_1<\dots<s_N$,
\[
\mathrm{Pos}_T(F)=\int_0^\infty\big(1-\mu_F(L_s)\big)\,ds\ \ge\ \sum_{k=0}^{N-1}(s_{k+1}-s_k)\big(1-\bar\mu(s_{k+1})\big)^+=:P_{\mathrm{lb}},
\]
where $\mu_F(A):=\frac1{2\pi}\int_A|\Fh|^2$.
\end{lemma}

\begin{proof}
$L_s$ is bounded because $\sT\to+\infty$ (Lemma~\ref{lem:mbounds}) and has finitely many components because
$\sT$ is real-analytic and nonconstant (Proposition~\ref{prop:band}(iii)). For $\|F\|=1$, $\mu_F$ is a probability
measure, and by the layer-cake formula $\mathrm{Pos}_T(F)=\int\sigma_+d\mu_F=\int_0^\infty\mu_F(\{\sigma_+>s\})ds$ with
$\{\sigma_+>s\}=\R\setminus L_s$ for $s\ge0$. The $E_i$ are disjoint subsets of $L_0\subset L_s$, so
$S_{L_s}\ge\sum_iS_{E_i}\ge\sum_i\sum_\pm\lambda^{(i),\pm}_0e^\pm_i\otimes e^\pm_i$, whence $\widetilde S_{L_s}\ge0$; the
subtracted term does not depend on $s$, so $\widetilde S_{L_s}$ and $\bar\mu$ are nondecreasing. For $F\perp\mathcal E_T$,
$\mu_F(L_s)=\langle\widetilde S_{L_s}F,F\rangle\le\bar\mu(s)$, and on $(s_k,s_{k+1}]$,
$1-\mu_F(L_s)\ge(1-\bar\mu(s_{k+1}))^+$. Integrate.
\end{proof}

The bound of Lemma~\ref{lem:wh-pos} is not an identity at $s=0$: the $e^\pm_i$ of different components are not
mutually orthogonal, so $\widetilde S_{L_0}$ is not the compression of $S_{L_0}$ to $\mathcal E_T^\perp$.

\begin{theorem}[coercivity on the prolate complement]\label{thm:wh-LW}
For every $T>0$,
\[
\mathfrak a_T(F)\ \ge\ c_X\|F\|^2\quad(F\in\mathcal X_T\cap\mathcal D_T),\qquad c_X:=P_{\mathrm{lb}}-N_{\mathrm{ub}},
\]
and likewise on $\mathcal X'_T$ with the variant constant. Both constants are spectral data of time--band
operators built from $(T,U_T,\sT)$, and the inequality holds verbatim for any symbol whose negative set is a
finite union of bounded symmetric intervals.
\end{theorem}

\begin{proof}
$\mathfrak a_T=\mathrm{Pos}_T-\mathrm{Neg}_T$ and Lemmas~\ref{lem:wh-neg} and \ref{lem:wh-pos}.
\end{proof}

\begin{remark}\label{rem:wh-LWnum}
Numerically, $c_X\approx0.4,\ 0.3,\ 0.5$ at $T=\frac12\log3,\ \log2,\ \frac12\log5$, and $\approx0.5,\ 0.4,\ 0.5$
with $\mathcal X'_T$ (Galerkin, not certified). If $\mathcal E_T$ is replaced by the span of explicit Galerkin
approximations of the $e^\pm_i$ (polynomials), Lemmas~\ref{lem:wh-neg}--\ref{lem:wh-pos} and
Theorem~\ref{thm:wh-LW} hold for the corresponding complement with the same proofs, every constant becomes an
explicit integral, and floating-point quadrature converged to six digits gives $c_X\approx0.1$ and $0.2$ at
$T=\frac12\log3$ and $\log2$, and $\approx0.6,\ 0.5,\ 0.3$ with the variant at the three windows. These values
are numerical: no interval arithmetic was used, and $s_i$ was evaluated on a grid. Nothing in this paper
depends on them. The constant $c_X(T)$ is not known to be positive for all $T$: $\sum_is_i\lambda^{(i)}_1$ grows
with the number of components, while $P_{\mathrm{lb}}$ does not.
\end{remark}

\subsection{The exact Schur coupling over the prolate block}\label{subsec:wh-schur}

\begin{lemma}[Schur complement of a closed form]\label{lem:wh-schurabs}
Let $\mathfrak q$ be a closed real form on a Hilbert space with domain $\mathcal D$, $\mathfrak q+M\|\cdot\|^2\ge\|\cdot\|^2$ for
some $M$, and let $\mathcal D=(\mathcal X\cap\mathcal D)\oplus E$, orthogonally, with $\mathcal X$ closed and $E\subset\mathcal D$ finite-dimensional.
Suppose $c_{\mathcal X}:=\inf\{\mathfrak q(w):w\in\mathcal X\cap\mathcal D,\ \|w\|=1\}>0$. Then:
\emph{(i)} for each $e\in E$ there is a unique $\mathsf Ke\in\mathcal X\cap\mathcal D$ with $\mathfrak q(\mathsf Ke,w)=\mathfrak q(e,w)$ for all
$w\in\mathcal X\cap\mathcal D$; $\mathsf K$ is linear and bounded, and
\[
\mathfrak q(w+e)=\mathfrak q(w+\mathsf Ke)+\mathfrak s(e),\qquad
\mathfrak s(e):=\mathfrak q(e)-\mathfrak c(e),
\]
\[
\mathfrak c(e):=\mathfrak q(\mathsf Ke)=\sup_{w\in\mathcal X\cap\mathcal D}\big[2\mathfrak q(e,w)-\mathfrak q(w)\big]\ge0 ;
\]
\emph{(ii)} with $\lambda_{\mathfrak q}$ the bottom of $\mathfrak q$ and $\lambda_S:=\lambda_{\min}(\mathfrak s)$,
\[
\frac{\min\{c_{\mathcal X},\lambda_S\}}{(1+\|\mathsf K\|)^2}\ \le\ \lambda_{\mathfrak q}\ \le\ c_{\mathcal X};
\]
moreover $\lambda_{\mathfrak q}\le\lambda_S$ if $\lambda_S\ge0$, $\lambda_{\mathfrak q}<0$ if $\lambda_S<0$, and $\lambda_S\ge\lambda_{\mathfrak q}$ if
$\lambda_{\mathfrak q}\ge0$.
\end{lemma}

\begin{proof}
(i) On $\mathcal X\cap\mathcal D$, $\mathfrak q\ge c_{\mathcal X}\|\cdot\|^2$, so $\mathfrak q$ and $\mathfrak q+M\|\cdot\|^2$ define equivalent
norms there and $(\mathcal X\cap\mathcal D,\mathfrak q)$ is a Hilbert space ($\mathfrak q$ is closed and $\mathcal X$ is closed). The functional
$w\mapsto\mathfrak q(e,w)$ is bounded on it by Cauchy--Schwarz for $\mathfrak q+M\|\cdot\|^2$; the Riesz theorem gives $\mathsf Ke$,
and $\|\mathsf Ke\|$ is bounded on the finite-dimensional $E$. Expanding $\mathfrak q(w+e)$ and using $\mathfrak q(e,w)=\mathfrak q(\mathsf Ke,w)$ gives
the identity; completing the square gives the supremum.
(ii) Write $a=\|w+\mathsf Ke\|$, $k=\|\mathsf K\|$. By (i), $\mathfrak q(w+e)\ge c_{\mathcal X}a^2+\lambda_S\|e\|^2$, and by
orthogonality $\|w+e\|^2=\|w\|^2+\|e\|^2\le(a+k\|e\|)^2+\|e\|^2\le(1+k)a^2+(k(1+k)+1)\|e\|^2\le(1+k)^2(a^2+\|e\|^2)$. This
is the lower bound. For the upper bounds take $e=0$; and $w=-\mathsf Ke_S$, $e=e_S$ a unit eigenvector of
$\lambda_S$, for which $\mathfrak q(e_S-\mathsf Ke_S)=\lambda_S$ and $\|e_S-\mathsf Ke_S\|^2=1+\|\mathsf Ke_S\|^2\ge1$. Finally, if
$\lambda_{\mathfrak q}\ge0$ then $\mathfrak s(e)=\mathfrak q(e-\mathsf Ke)\ge\lambda_{\mathfrak q}(1+\|\mathsf Ke\|^2)\|e\|^2\ge\lambda_{\mathfrak q}\|e\|^2$.
\end{proof}

\begin{theorem}[the exact Schur coupling]\label{thm:wh-schur}
Let $c_X$ be the bottom of $\mathfrak a_T$ on $\mathcal X_T\cap\mathcal D_T$, and suppose $c_X>0$ (by
Theorem~\ref{thm:wh-LW}, this holds whenever $P_{\mathrm{lb}}>N_{\mathrm{ub}}$). Let $\mathsf K$, $\mathfrak c_T$, $\mathfrak s_T=\mathsf M_T-\mathfrak c_T$
be given by Lemma~\ref{lem:wh-schurabs} for $\mathfrak q=\mathfrak a_T$ on $\PT\cap\mathcal D_T=(\mathcal X_T\cap\mathcal D_T)\oplus\mathcal E^P_T$.
\emph{(i)} $\lamA>0\iff\mathfrak s_T\succ0$, and $\lamA\ge0\iff\mathfrak s_T\succeq0$.
\emph{(ii)} If $\mathsf M_T\succ0$, put $\rho_T:=\max\{\mathfrak c_T(e)/\mathsf M_T(e):e\in\mathcal E^P_T\setminus0\}$. Then
\[
\lamA>0\iff c_X>0,\ \ \mathsf M_T\succ0\ \text{ and }\ \rho_T<1 ,
\]
and the first condition is automatic when $\lamA>0$.
\emph{(iii) (Inheritance.)} If $\lamA>0$ and $\lambda_S\le c_X$, then
\[
\frac{\lamA}{\|\mathsf M_T\|}\ \le\ 1-\rho_T\ \le\ \frac{\lambda_S}{\lambda_{\min}(\mathsf M_T)}\ \le\ \frac{(1+\|\mathsf K\|)^2}{\lambda_{\min}(\mathsf M_T)}\,\lamA .
\]
\end{theorem}

\begin{proof}
Lemma~\ref{lem:wh-W}(i) gives the decomposition, and $\mathfrak a_T+M\|\cdot\|^2$ with $M=\Aconst-C_0+1$ is the closed
form of Lemma~\ref{lem:compact}. (i) If $\lamA>0$ ($\ge0$), then $\lambda_S\ge\lamA$ by the last clause of
Lemma~\ref{lem:wh-schurabs}(ii). Conversely, if $\mathfrak s_T\succ0$ the lower bound there gives $\lamA>0$, and if
$\mathfrak s_T\succeq0$ then $\mathfrak a_T(w+e)\ge c_X\|w+\mathsf Ke\|^2+\lambda_S\|e\|^2\ge0$. (ii) If $\lamA>0$ then $c_X\ge\lamA>0$, and
$\mathfrak s_T\succ0$ forces $\mathsf M_T\succeq\mathfrak s_T\succ0$; on the finite-dimensional $\mathcal E^P_T$ with $\mathsf M_T\succ0$,
$\mathsf M_T-\mathfrak c_T\succ0$ iff $\rho_T<1$. (iii) $1-\rho_T=\min_e\mathfrak s_T(e)/\mathsf M_T(e)$. At a unit eigenvector $e_S$
of $\lambda_S$ this is at most $\lambda_S/\lambda_{\min}(\mathsf M_T)$, and $\lambda_S\le(1+\|\mathsf K\|)^2\lamA$ by the lower bound of
Lemma~\ref{lem:wh-schurabs}(ii) with $\min\{c_X,\lambda_S\}=\lambda_S$. In the other direction
$\mathfrak s_T(e)\ge\lamA\|e\|^2\ge\lamA\mathsf M_T(e)/\|\mathsf M_T\|$.
\end{proof}

\begin{remark}\label{rem:wh-schur}
Theorem~\ref{thm:wh-schur} reduces the sign of $\lamA$, at any window where $c_X>0$ is established, to a
matrix of size at most $2n_T$ coupled to an operator with a margin of order one; and part (iii) shows that the
coupling constant $1-\rho_T$ carries exactly the margin $\lamA$, up to factors computable from the reduction.
Numerically, $\lambda_{\min}(\mathsf M_T)$ is much larger than $\lamA$ (by several orders of magnitude at
$T=\frac12\log5$), $\|\mathfrak c_T\|$ is of order $0.2$, and adding the index-one prolates to $\mathcal E_T$ brings
$\lambda_{\min}(\mathsf M_T)$ to within a few percent of $\lamA$ (Galerkin, not certified). So neither the block
nor a norm bound for the coupling decides the sign: the decision is in $\rho_T$.
\end{remark}

\subsection{Four points on the zero side as a form of inertia (2,2)}\label{subsec:wh-quartet}

\begin{lemma}\label{lem:wh-Q}
For real $F\in L^2(\IT)$ and $z=\delta+i\gamma$, $\delta,\gamma\in\R$, put
$\mathfrak Q_z(F):=4\Rel\big[\Phi_F(\tfrac12+z)\Phi_F(\tfrac12-z)\big]$, the sum of $\Phi_F(\rho)\Phi_F(1-\rho)$ over the four
points $\rho=\frac12\pm\delta\pm i\gamma$. Then
\[
\mathfrak Q_z(F)=4|\Phi_c(F)|^2-4|\Phi_s(F)|^2,
\]
with $\Phi_c(F):=\int F(t)\cosh(\delta t)e^{i\gamma t}dt$ and $\Phi_s(F):=\int F(t)\sinh(\delta t)e^{i\gamma t}dt$; this is
a real quadratic form of inertia $(2,2)$ on real $L^2(\IT)$ when $\delta\gamma\ne0$.
\end{lemma}

\begin{proof}
For real $F$, $\Phi_F(\frac12+z)=X_+$ and $\Phi_F(\frac12-z)=\overline{X_-}$ with $X_\pm:=\Phi_c\pm\Phi_s$, and
$X_+\overline{X_-}=|\Phi_c|^2-|\Phi_s|^2+(\Phi_s\overline{\Phi_c}-\Phi_c\overline{\Phi_s})$, the last term being purely imaginary.
The points $\frac12\pm\bar z$ contribute the complex conjugates. The four real functionals
$\Rel\Phi_c,\Iml\Phi_c,\Rel\Phi_s,\Iml\Phi_s$ are linearly independent on real $L^2(\IT)$, the functions
$\cosh\delta t\cos\gamma t$, $\cosh\delta t\sin\gamma t$, $\sinh\delta t\cos\gamma t$, $\sinh\delta t\sin\gamma t$ being linearly
independent on $\IT$ when $\delta\gamma\ne0$.
\end{proof}

\begin{proposition}[Kre\u{\i}n threshold for a quartet on a finite block]\label{prop:wh-K}
Let $X\subset\PT\cap\mathcal D_T$ be finite-dimensional with $\mathfrak a_T|_X\succ0$, let $S_X:X\to\R^2$,
$S_Xe=(\Rel\Phi_s(e),\Iml\Phi_s(e))$, and $H_+:=\mathfrak a_T|_X+4|\Phi_c|^2\succ0$. Then
\[
(\mathfrak a_T+\mathfrak Q_z)|_X\succ0\iff 4\,S_XH_+^{-1}S_X^{\mathsf T}\prec I_2 .
\]
\end{proposition}

\begin{proof}
By Lemma~\ref{lem:wh-Q}, $(\mathfrak a_T+\mathfrak Q_z)|_X=H_+-G^{\mathsf T}G$ with $G=2S_X$, and for $H_+\succ0$ this is positive
definite iff $\|GH_+^{-1/2}\|<1$, i.e.\ $GH_+^{-1}G^{\mathsf T}\prec I_2$ (Lemma~\ref{lem:wh-krein}(iii) for a perturbation of
rank at most two).
\end{proof}

\subsection{Moments of the pencil}\label{subsec:wh-moments}

\begin{lemma}[moment criterion]\label{lem:wh-mom}
Let $m_k:=\operatorname{tr}B_T^k=\sum_j\beta_j^k$ ($k\ge1$). Then $\beta_1\le m_k^{1/k}$, $m_k^{1/k}\downarrow\beta_1$, and
$\beta_1<1\iff m_k<1$ for some $k$.
\end{lemma}

\begin{proof}
$\beta_1^k\le m_k=\beta_1^k\sum_j(\beta_j/\beta_1)^k$, and the sum decreases to the multiplicity of $\beta_1$. If
$\beta_1<1$, $m_k\le\beta_1^{k-1}m_1\to0$.
\end{proof}

The next two lemmas are barriers for a declared class of arguments: criteria that are valid for \emph{all}
nonnegative trace-class operators, or for all bounded perturbations of the form. They do not concern criteria
valid only on symbols of the arithmetic class, nor certificates of positivity such as those of
\S\ref{sec:log5}.

\begin{lemma}[universal moment criteria inherit the margin]\label{lem:wh-T}
Let $\Phi:\R^k\to\R$ be such that $\Phi(m_1(B),\dots,m_k(B))<1$ implies $\|B\|<1$ for every nonnegative trace-class
operator $B$. If $\beta_1<1$, then the point $m(B_T)=(m_1(B_T),\dots,m_k(B_T))$ lies within sup-distance
$k(1-\beta_1)\le k\lamA/c_P$ of the set $\{\Phi\ge1\}$. If $\Phi$ is $L$-Lipschitz for the sup-norm, then
$1-\Phi(m(B_T))\le Lk\lamA/c_P$.
\end{lemma}

\begin{proof}
$B':=B_T+(1-\beta_1)v_1\otimes v_1$ is nonnegative and trace class with $\|B'\|=1$, so $\Phi(m(B'))\ge1$; it has the
eigenvalues of $B_T$ with one copy of $\beta_1$ replaced by $1$, so $|m_j(B')-m_j(B_T)|=1-\beta_1^j\le j(1-\beta_1)$. Conclude
with \eqref{eq:wh-sandwich}.
\end{proof}

\begin{lemma}[continuous universal criteria inherit the margin]\label{lem:wh-Tprime}
Let $\mathcal F$ be the set of forms $\mathfrak q=\mathfrak a_T|_{\PT}+\mathfrak b$ with $\mathfrak b$ a bounded symmetric form on $\PT$, and
let $\Phi:\mathcal F\to\R$ satisfy $|\Phi(\mathfrak q)-\Phi(\mathfrak q')|\le L\|\mathfrak b-\mathfrak b'\|$ and
$\Phi(\mathfrak q)>0\Rightarrow\inf_{\|u\|=1}\mathfrak q(u)>0$. If $\lamA\ge0$, then $\Phi(\mathfrak a_T|_{\PT})\le L\lamA$.
\end{lemma}

\begin{proof}
$h^*_T$ is an eigenvector of the operator of $\mathfrak a_T|_{\PT}$ (Lemma~\ref{lem:compact}), so
$\mathfrak q:=\mathfrak a_T|_{\PT}-\lamA|\langle\cdot,h^*_T\rangle|^2$ satisfies $\mathfrak q(\alpha h^*_T+u')=\mathfrak a_T(u')\ge0$ for
$u'\perp h^*_T$, and $\mathfrak q(h^*_T)=0$. Hence $\Phi(\mathfrak q)\le0$, and $\|\mathfrak b\|=\lamA$.
\end{proof}

\begin{remark}\label{rem:wh-mom}
The criterion of Lemma~\ref{lem:wh-mom} is exact and computable from the symbol, but numerically the first $k$
with $m_k<1$ is of order $20$, $10^5$ and $10^{10}$ at $T=\frac12\log3,\ \log2,\ \frac12\log5$ (Galerkin, not
certified), and at that $k$ the margin $1-m_k^{1/k}$ is $1-\beta_1$ to the precision of the computation, as
Lemma~\ref{lem:wh-T} predicts.
\end{remark}

\subsection{Determinants}\label{subsec:wh-det}

We call a pair $(\mathfrak p_\star,\mathfrak n_\star)$ an \emph{admissible splitting} of $\mathfrak a_T|_{\PT}$ if $\mathfrak p_\star$ is a closed
form on $\PT\cap\mathcal D_T$ with $\mathfrak p_\star\ge c_\star\|\cdot\|^2$, $c_\star>0$, $\mathfrak n_\star\ge0$ is bounded and trace class
on $\PT$, and $\mathfrak a_T=\mathfrak p_\star-\mathfrak n_\star$ on $\PT\cap\mathcal D_T$. Examples: $(\mathfrak p,\mathfrak n)$ of
Definition~\ref{def:wh-B}, and $(\mathfrak a^{(U_0)}_T,\mathrm{Neg}^{\mathrm{new}}_T)$ of Proposition~\ref{prop:wh-oldnew} below
whenever $\mathfrak a^{(U_0)}_T$ is coercive. Write $B_\star:=\mathfrak p_\star^{-1/2}\mathfrak n_\star\mathfrak p_\star^{-1/2}$, with eigenvalues
$\beta^\star_1\ge\beta^\star_2\ge\dots$, and $u^\star_1$ for the unit pencil extremal as in
Theorem~\ref{thm:wh-pencil}(ii). The proofs of Theorem~\ref{thm:wh-pencil} use only these properties, so its
conclusions hold for every admissible splitting.

\begin{theorem}[the Birman--Schwinger determinant of the window]\label{thm:wh-det}
For an admissible splitting, $B_\star$ is nonnegative and trace class, and
\[
D_\star(s;T):=\det(1-sB_\star)=\prod_j(1-s\beta^\star_j)
\]
is an entire function of $s$ of genus zero, real on $\R$, with $D_\star(0;T)=1$, whose zeros are the points
$1/\beta^\star_j$ with multiplicity. Moreover
\[
n_-(\mathfrak a_T|\PT)=\#\{\text{zeros of }D_\star(\cdot;T)\text{ in }(0,1)\},
\]
\[
\lamA>0\iff D_\star(s;T)\ne0\ \text{for all }s\in[0,1]\iff\beta^\star_1<1 .
\]
For $(\mathfrak p,\mathfrak n)$ of Definition~\ref{def:wh-B} we write $D_B(s;T)$.
\end{theorem}

\begin{proof}
$\mathfrak p_\star^{-1/2}$ is bounded and $\mathfrak n_\star$ is trace class, so $B_\star$ is trace class; $\sum_j\beta^\star_j<\infty$, so the
product converges locally uniformly to an entire function of genus zero, and it equals the Fredholm
determinant by Lidskii's theorem \cite[Thm.~3.7]{Simon2005}. A zero in $(0,1)$ is a $\beta^\star_j>1$, and $s=1$ is
a zero iff some $\beta^\star_j=1$. Theorem~\ref{thm:wh-pencil}(i), valid for admissible splittings, gives the rest.
\end{proof}

\begin{theorem}[determinants inherit the margin]\label{thm:wh-detinherit}
Let $(\mathfrak p_\star,\mathfrak n_\star)$ be an admissible splitting, and suppose $\lamA\ge0$. Put
$\Pi_\star:=\prod_{j\ge2}(1-\beta^\star_j)\in[0,1]$.
\emph{(i)} The first positive zero $s_1=1/\beta^\star_1$ of $D_\star(\cdot;T)$ satisfies
\[
\frac{\lamA}{\mathfrak p_\star(u^\star_1)}\ \le\ 1-\frac1{s_1}\ \le\ \frac{\lamA}{\mathfrak p_\star(h^*_T)}\ \le\ \frac{\lamA}{c_\star}.
\]
\emph{(ii)} $D_\star(1;T)=(1-\beta^\star_1)\,\Pi_\star$, so that
$\lamA\,\Pi_\star/\mathfrak p_\star(u^\star_1)\le D_\star(1;T)\le\lamA\,\Pi_\star/\mathfrak p_\star(h^*_T)\le\lamA/c_\star$.
\emph{(iii) (Leading terms.)} Let $D_0\in\R$ and $\varepsilon_0\ge0$ with $|D_\star(1;T)-D_0|\le\varepsilon_0|D_0|$. If
$\varepsilon_0<1$ and $D_0>0$, then
\[
c_\star(1-\varepsilon_0)\,D_0\ \le\ \lamA\qquad\text{and}\qquad D_0\ \le\ \frac{\lamA}{c_\star(1-\varepsilon_0)} .
\]
If $\varepsilon_0\ge1$, the hypothesis is compatible with $D_\star(1;T)=0$, i.e.\ with $\lamA=0$, so it does not
imply $\lamA>0$.
\end{theorem}

\begin{proof}
(i) is \eqref{eq:wh-sandwich} for the splitting, with $1/s_1=\beta^\star_1$. (ii) is the product formula at $s=1$
and (i); note $\Pi_\star\le1$, and $\Pi_\star\ge0$ since $\beta^\star_j\le\beta^\star_1\le1$. (iii) If $\varepsilon_0<1$ and $D_0>0$ then
$D_\star(1;T)\ge(1-\varepsilon_0)D_0>0$, and by (ii) $\lamA\ge\mathfrak p_\star(h^*_T)D_\star(1;T)/\Pi_\star\ge c_\star D_\star(1;T)$; also
$(1-\varepsilon_0)D_0\le D_\star(1;T)\le\lamA/c_\star$. The last sentence is the case $D_\star(1;T)=0$, admissible when
$\varepsilon_0\ge1$ since $|0-D_0|=|D_0|$; and $D_\star(1;T)=0$ means $\beta^\star_1=1$, i.e.\ $\lamA=0$ by Theorem~\ref{thm:wh-pencil}(i).
\end{proof}

Theorem~\ref{thm:wh-detinherit} is a barrier for the following class of arguments: establish $\lamA>0$ by
exhibiting a leading term $D_0$ for a determinant of an admissible splitting, together with a bound on the
relative error. Part (iii) shows that such an argument succeeds only if the leading term is itself a two-sided
estimate of $\lamA$ up to the explicit factors $c_\star$, $\mathfrak p_\star(u^\star_1)$, $\Pi_\star$ and $1\pm\varepsilon_0$; a leading
term with a multiplicative error of order one ($\varepsilon_0\ge1$) never decides the sign. The class does not
contain structural certificates (an identity or a dual certificate whose verification is a sign check of
quantities of order one, as in \S\ref{sec:log5}); for those, the size of $\lamA$ enters only through the
value of the bound they produce.

\begin{corollary}[determinants of the negative part alone]\label{cor:wh-negonly}
No function of $(\sigma_-,U_T)$ decides the sign of the bottom on $\PT$ within the family
$\{\varepsilon\sigma_+-\sigma_-:\varepsilon>0\}$. In particular neither $\det(1+s\,\mathrm{Neg}_T)$, nor any asymptotic formula for
it, nor the principal angles between $L^2(\IT)$ and $\operatorname{ran}\mathbf 1_{U_T}(D)$, decides it.
\end{corollary}

\begin{proof}
All these data are the same for every member of the family (Proposition~\ref{prop:wh-notangles}), while the
bottom on $\PT$ is negative for $\varepsilon<\beta_1$ and positive for $\varepsilon>\beta_1$, with $0<\beta_1<\infty$.
\end{proof}

\begin{remark}\label{rem:wh-det}
The Szeg\H{o}--Kac--Akhiezer theorem and its refinements \cite[\S10.9]{BottcherSilbermann2006} describe
$\det(1+\mathrm{Neg}_T)$, the truncation to an interval of length $2T$ of the bounded symbol $1+\sigma_-$, for a fixed
symbol as the length tends to infinity; by Corollary~\ref{cor:wh-negonly} that determinant carries no
information on the sign. The determinant $D_B(1;T)$ that does is a ratio of the determinants of
$\mathfrak a_T|\PT$ and $\mathfrak p$ in a regularized sense, and the symbols of the two blocks have disjoint supports
$\{\sT<0\}$ and $\{\sT>0\}$, so the leading terms of such theorems do not apply to it; and neither $\mathfrak a_T$
itself nor its quotient by the archimedean part has a Fredholm determinant, the symbol growing like
$\log|\tau|$. Numerically, at $T=\frac12\log5$, $\log D_B(1;T)$ is dominated by $\log(1-\beta_1)$ together with
three further eigenvalues $\beta_j$ within $10^{-4}$ of $1$ (Galerkin, not certified).
\end{remark}

\subsection{Old and new negativity, and contact thresholds}\label{subsec:wh-oldnew}

\begin{proposition}[old/new partition]\label{prop:wh-oldnew}
Let $T_0<T$, $U_0:=\{\sigma_{T_0}<0\}$,
\[
\mathrm{Neg}^{\mathrm{new}}_T:=R_T(\sigma_-\mathbf 1_{U_T\setminus U_0})(D)R_T\ \ge0\quad\text{(trace class)},\qquad
\mathfrak a^{(U_0)}_T:=\mathfrak a_T+\mathrm{Neg}^{\mathrm{new}}_T ,
\]
the form of the symbol whose only negativity is that of $\sT$ on $U_0$. On $\PT\cap\mathcal D_T$:
\emph{(i)} $\lamA\ge0\iff\mathfrak a^{(U_0)}_T\ge\mathrm{Neg}^{\mathrm{new}}_T$; so $\lamA\ge0$ implies $\mathfrak a^{(U_0)}_T\ge0$, and
$\lamA>0$ implies $\mathfrak a^{(U_0)}_T\ge\lamA\|\cdot\|^2$.
\emph{(ii)} If $\mathfrak a^{(U_0)}_T\ge c\|\cdot\|^2$ with $c>0$ on $\PT\cap\mathcal D_T$, then $(\mathfrak a^{(U_0)}_T,\mathrm{Neg}^{\mathrm{new}}_T)$ is an
admissible splitting; with $\beta_{\mathrm{new}}$ its top pencil eigenvalue and $u^*$ its unit extremal,
$\lamA>0\iff\beta_{\mathrm{new}}<1$, and when $\lamA\ge0$,
\[
\frac{\lamA}{\mathfrak a^{(U_0)}_T(u^*)}\ \le\ 1-\beta_{\mathrm{new}}\ \le\ \frac{\lamA}{\mathfrak a^{(U_0)}_T(h^*_T)} .
\]
\end{proposition}

\begin{proof}
$\sT=\sT\mathbf 1_{\R\setminus(U_T\setminus U_0)}-\sigma_-\mathbf 1_{U_T\setminus U_0}$, so $\mathfrak a_T=\mathfrak a^{(U_0)}_T-\mathrm{Neg}^{\mathrm{new}}_T$; the
added term is bounded, so $\mathfrak a^{(U_0)}_T$ is closed on $\mathcal D_T$, and its compression to $\PT$ has trace-class
negative part $\mathrm{pr}_T\mathrm{Neg}^{\mathrm{new}}_T\mathrm{pr}_T$. (i) is immediate; (ii) is Theorem~\ref{thm:wh-pencil} for this
splitting.
\end{proof}

Numerically, $\mathfrak a^{(U_0)}_T$ has a bottom of order $0.1$ on $\PT$ at the computed pairs, while
$1-\beta_{\mathrm{new}}=\lamA/\mathfrak a^{(U_0)}_T(h^*_T)$ to the precision of the computation (Galerkin, not
certified): the partition separates a coercive part from a pencil that carries the whole margin.

\begin{lemma}[a contact threshold does not change the form]\label{lem:wh-birth}
Let $n\ge2$ and $T_n:=\frac12\log n$. For every $F\in L^2(I_{T_n})$, $\frac1{2\pi}\int\cos(\tau\log n)|\Fh|^2d\tau=\Rel\cF(\log n)=0$.
Hence the form with symbol $\sigma_{T_n}-v_n\cos(\tau\log n)$, the symbol of the windows immediately to the right
of $T_n$, coincides with $\mathfrak a_{T_n}$ on $L^2(I_{T_n})$.
\end{lemma}

\begin{proof}
$\cF$ is continuous and supported in $[-2T_n,2T_n]=[-\log n,\log n]$, so $\cF(\pm\log n)=0$, and
$\widehat{\cF}=|\Fh|^2$. This is the computation in the proof of Lemma~\ref{lem:monoT}.
\end{proof}

\begin{lemma}[newborn components are sub-Shannon]\label{lem:wh-shannon}
Let $T_n=\frac12\log n$, $\sigma^{\flat}:=\sigma_{T_n}$ and $\sigma^{\sharp}:=\sigma^\flat-v_n\cos(\tau\log n)$. If
$\Delta\subset(0,\infty)$ is a connected component of $\{\sigma^\sharp<0\}$ disjoint from $\{\sigma^\flat<0\}$, then
$\cos(\tau\log n)>0$ on $\Delta$, $|\Delta|<\pi/\log n$, and $\operatorname{tr}S_{\Delta\cup(-\Delta)}=2T_n|\Delta|/\pi<1$ at the window
$T_n$.
\end{lemma}

\begin{proof}
On $\Delta$, $v_n\cos(\tau\log n)>\sigma^\flat\ge0$, so the connected set $\Delta$ lies in one interval of
$\{\cos(\tau\log n)>0\}$, of length $\pi/\log n$; and $2T_n/\log n=1$.
\end{proof}

By Lemmas~\ref{lem:wh-birth} and \ref{lem:wh-shannon}, a component of the negative set created at a contact
threshold is invisible to the form at the threshold and carries less than one prolate degree of freedom. It is
not small in the symbol, and beyond the threshold its cost is, to first order, $v_n\Rel c_{h^*_T}(\log n)$, a
quantity of the minimizer; no continuation argument with a margin of its own follows from the two lemmas.

\subsection{The negative index by parity sectors}\label{subsec:wh-index}

We use the notation of Theorem~\ref{thm:parityschur}: $\ell_{\mathrm{par}}=\langle\cdot,\chi_{\mathrm{par}}\rangle$,
$\varepsilon_{\mathrm{par}}$, $\mathcal Q^{\mathrm{par}}_T$, $h_{\mathrm{par}}$, $\mathfrak g_{\mathrm{par}}=W_T(h_{\mathrm{par}})=\mathfrak a_T(h_{\mathrm{par}})+2\varepsilon_{\mathrm{par}}$. By
Theorem~\ref{thm:parityschur}(o), on real $F=F_{\mathrm e}+F_{\mathrm o}\in\mathcal D_T$,
\begin{equation}\label{eq:wh-aW}
\mathfrak a_T(F)=W_T(F)-2\ell_{\mathrm{ev}}(F_{\mathrm e})^2+2\ell_{\mathrm{odd}}(F_{\mathrm o})^2 .
\end{equation}

\begin{proposition}[index by sectors]\label{prop:wh-index}
\emph{(i)} If $0$ is not in the spectrum of the operator of $\mathfrak a_T|\mathcal Q^{\mathrm{par}}_T$ in either sector (in
particular if $\lamA>0$), then
\[
n_-(\mathfrak a_T|\mathcal D_T)=n_-(\mathfrak a_T|\mathcal Q^{\mathrm{ev}}_T)+[\mathfrak g_{\mathrm{ev}}(T)<2]+n_-(\mathfrak a_T|\mathcal Q^{\mathrm{odd}}_T)+[\mathfrak g_{\mathrm{odd}}(T)<-2].
\]
\emph{(ii)} For $0<T\le\frac12\log5$: $\mathfrak a_T\ge0$ on $\mathcal D^{\mathrm{odd}}_T$, $n_-(\mathfrak a_T|\mathcal D^{\mathrm{ev}}_T)\le1$, and
$\mathfrak a_T\ge10^{-18}\|\cdot\|^2$ on $\{F\in\mathcal D_T:\ell_{\mathrm{ev}}(F_{\mathrm e})=0\}$, a subspace of codimension one
containing $\PT\cap\mathcal D_T$. If moreover $\mathfrak a_T(\chi_{\mathrm{ev}})<0$, then $n_-(\mathfrak a_T|\mathcal D_T)=1$, the negative
direction is even, and, if the even sector is as in (i), $\mathfrak g_{\mathrm{ev}}(T)<2$.
\emph{(iii)} $\mathfrak a_T(\chi_{\mathrm{ev}})\le\mT(\chi_{\mathrm{ev}})=C_0\|\chi_{\mathrm{ev}}\|^2+\int_0^{2T}\big(\|\chi_{\mathrm{ev}}\|^2-c_{\chi_{\mathrm{ev}}}(u)\big)K(u)\,du+\|\chi_{\mathrm{ev}}\|^2\int_{2T}^\infty K(u)\,du$,
an explicit function of $T$.
\end{proposition}

\begin{proof}
(i) By Theorem~\ref{thm:parityschur}(a1), $\mathcal D^{\mathrm{par}}_T=\mathcal Q^{\mathrm{par}}_T\oplus\R h_{\mathrm{par}}$ orthogonally for
$\mathfrak a_T$, so $n_-(\mathfrak a_T|\mathcal D^{\mathrm{par}}_T)=n_-(\mathfrak a_T|\mathcal Q^{\mathrm{par}}_T)+[\mathfrak a_T(h_{\mathrm{par}})<0]$, and
$\mathfrak a_T(h_{\mathrm{par}})=\mathfrak g_{\mathrm{par}}-2\varepsilon_{\mathrm{par}}$; the two sectors are $\mathfrak a_T$-orthogonal
(Theorem~\ref{thm:parityschur}(o)). (ii) By Theorem~\ref{thm:log5W} and Lemma~\ref{lem:monoT},
$W_T\ge10^{-18}\|\cdot\|^2$ on $\mathcal D_T$. By \eqref{eq:wh-aW}, $\mathfrak a_T(F_{\mathrm o})=W_T(F_{\mathrm o})+2\ell_{\mathrm{odd}}(F_{\mathrm o})^2\ge0$
and, if $\ell_{\mathrm{ev}}(F_{\mathrm e})=0$, $\mathfrak a_T(F)=W_T(F_{\mathrm e})+W_T(F_{\mathrm o})+2\ell_{\mathrm{odd}}(F_{\mathrm o})^2\ge10^{-18}\|F\|^2$. Since
$\ker\ell_{\mathrm{ev}}$ has codimension one in $\mathcal D^{\mathrm{ev}}_T$, $n_-(\mathfrak a_T|\mathcal D^{\mathrm{ev}}_T)\le1$ by min--max. If
$\mathfrak a_T(\chi_{\mathrm{ev}})<0$ the even index is $1$; as $\mathcal Q^{\mathrm{ev}}_T\subset\ker\ell_{\mathrm{ev}}$ has index $0$, (i) forces
$\mathfrak a_T(h_{\mathrm{ev}})<0$. (iii) $\mathfrak a_T=\mT+\PiT$ with $\PiT(\chi_{\mathrm{ev}})=-\sum v_nc_{\chi_{\mathrm{ev}}}(\log n)\le0$ because
$\chi_{\mathrm{ev}}\ge0$; and \eqref{eq:Aform} without the prime sum, with $c_{\chi_{\mathrm{ev}}}=0$ beyond $2T$, gives the formula.
\end{proof}

\begin{proposition}[the polar resolvent]\label{prop:wh-polarres}
Suppose that $\mathfrak a_T$ is invertible on $\mathcal D^{\mathrm{par}}_T$ for both parities, that the sectors are as in
Proposition~\ref{prop:wh-index}(i), and that $r_{\mathrm{par}}:=\langle\mathfrak a_T^{-1}\chi_{\mathrm{par}},\chi_{\mathrm{par}}\rangle\ne0$. Then
$h_{\mathrm{par}}=\mathfrak a_T^{-1}\chi_{\mathrm{par}}/r_{\mathrm{par}}$ and $r_{\mathrm{par}}=1/(\mathfrak g_{\mathrm{par}}-2\varepsilon_{\mathrm{par}})$. Consequently the
$2\times2$ matrix of $\mathfrak a_T^{-1}$ on the Tate functionals $(A,B)$ has eigenvalues
$2/(\mathfrak g_{\mathrm{ev}}-2)$ and $2/(\mathfrak g_{\mathrm{odd}}+2)$. In particular, when $\mathfrak g_{\mathrm{ev}}$ is small, $\mathfrak a_T^{-1}$ has an
eigenvalue $-1-\frac12\mathfrak g_{\mathrm{ev}}+O(\mathfrak g_{\mathrm{ev}}^2)$ on the span of the Tate functionals.
\end{proposition}

\begin{proof}
$g:=\mathfrak a_T^{-1}\chi_{\mathrm{par}}$ satisfies $\mathfrak a_T(g,G)=\ell_{\mathrm{par}}(G)=0$ for $G\in\mathcal Q^{\mathrm{par}}_T$ and $\ell_{\mathrm{par}}(g)=r_{\mathrm{par}}$, so
$h_{\mathrm{par}}=g/r_{\mathrm{par}}$ by the uniqueness in Theorem~\ref{thm:parityschur}(a), and
$\mathfrak a_T(h_{\mathrm{par}})=\langle\chi_{\mathrm{par}},g\rangle/r_{\mathrm{par}}^2=1/r_{\mathrm{par}}$. On real $F$, $A=\ell_{\mathrm{ev}}+\ell_{\mathrm{odd}}$ and
$B=\ell_{\mathrm{ev}}-\ell_{\mathrm{odd}}$, and the cross terms between parities vanish, so the matrix is
$\begin{pmatrix}r_{\mathrm{ev}}+r_{\mathrm{odd}}&r_{\mathrm{ev}}-r_{\mathrm{odd}}\\ r_{\mathrm{ev}}-r_{\mathrm{odd}}&r_{\mathrm{ev}}+r_{\mathrm{odd}}\end{pmatrix}$, with eigenvalues
$2r_{\mathrm{ev}}$ and $2r_{\mathrm{odd}}$.
\end{proof}

\begin{remark}\label{rem:wh-index}
Numerically, $\mT(\chi_{\mathrm{ev}})<0$ for $T\ge0.12$; this is a one-dimensional evaluation of the explicit
expression in Proposition~\ref{prop:wh-index}(iii), not certified by interval arithmetic. Granting it,
Proposition~\ref{prop:wh-index}(ii) gives $n_-(\mathfrak a_T|\mathcal D_T)=1$ on $[0.12,\frac12\log5]$ with an even negative
direction. The index is not a lever for the sign: when both sectors are nondegenerate,
Proposition~\ref{prop:wh-index}(i) shows that $n_-(\mathfrak a_T|\mathcal D_T)=1$ together with $\mathfrak g_{\mathrm{ev}}<2$ holds iff
$\mathfrak a_T$ is positive definite on $\mathcal Q^{\mathrm{ev}}_T$ and $\mathcal Q^{\mathrm{odd}}_T$ and $\mathfrak g_{\mathrm{odd}}\ge-2$; without the
sign of $\mathfrak a_T(h_{\mathrm{ev}})$, the count $n_-=1$ is compatible with a negative direction in $\mathcal Q^{\mathrm{ev}}_T$. By
Theorem~\ref{thm:parityschur}(a4), $\mathfrak g_{\mathrm{ev}}$ is comparable to $\lambda^{\mathrm W}_{T,\mathrm{ev}}\|h_{\mathrm{ev}}\|^2$ when both
are small, so Proposition~\ref{prop:wh-polarres} explains an eigenvalue near $-1$ of the inverse form on the
Tate functionals as the inverse of the negative square of the pole (Lemma~\ref{lem:polar}).
\end{remark}


\section{Sampling at the Nyquist mesh and two-sector structure}\label{sec:nyq}

The autocorrelation of $F\in L^2(\IT)$ is supported in $[-2T,2T]$, so $|\Fh|^2$ is band limited and is
determined by its values on the mesh $\frac{\pi}{2T}\Z$. This section rewrites $\mathfrak a_T$ as a sum over
that mesh. The identity is exact only for the symbol obtained by truncating the L\'evy measure of $m$ to the
window, as in Theorem~\ref{thm:A}; with $\sT$ itself there is an aliasing term, which we compute. Splitting the
mesh by parity gives two orthonormal bases of $L^2(\IT)$ in which the form is a mean of two diagonal forms,
and from this we obtain a Schur reduction with an explicit coercivity constant. Nothing in this section gives
the sign of $\mathfrak a_T$ at a new window.

Throughout, $\omega_k:=\frac{k\pi}{2T}$ $(k\in\Z)$, and
\begin{equation}\label{eq:nyq-sT}
s_T(\tau):=C_0+\int_{2T}^\infty K(u)\,du+\int_0^{2T}(1-\cos\tau u)K(u)\,du-\sum_{n\in\NT}v_n\cos(\tau\log n).
\end{equation}
By \eqref{eq:m}, $s_T=\sT+r_T$ with $r_T(\tau):=\int_{2T}^\infty\cos(\tau u)K(u)\,du$. We call $s_T$ the
\emph{truncated symbol}. We also write $L_T(\tau):=\int_0^{2T}(1-\cos\tau u)K(u)\,du\ge0$ and
\[
c_T:=C_0+\int_{2T}^\infty K(u)\,du-\Aconst,\qquad\text{so that}\qquad
s_T(\tau)=c_T+L_T(\tau)+\sum_{n\in\NT}v_n\big(1-\cos(\tau\log n)\big).
\]

\subsection{The sampling identity}

\begin{lemma}[sampling the autocorrelation]\label{lem:nyq-sampling}
Let $F\in L^2(\IT)$.
\emph{(i)} $\frac1{4T}\sum_{k\in\Z}|\Fh(\omega_k)|^2=\|F\|^2$.
\emph{(ii)} For every $u\in\R$, $\frac1{4T}\sum_k|\Fh(\omega_k)|^2e^{-i\omega_ku}=\sum_{j\in\Z}\cF(u+4Tj)$, the
series converging absolutely and uniformly in $u$; for $|u|\le2T$ the right side is $\cF(u)$, and for
$0\le u\le2T$
\[
\frac1{4T}\sum_k(-1)^k|\Fh(\omega_k)|^2e^{-i\omega_ku}=\cF(u-2T).
\]
\emph{(iii)} $\sum_k(-1)^k|\Fh(\omega_k)|^2=0$.
\emph{(iv)} If $\nu$ is a finite even signed Borel measure on $[-2T,2T]$, then
$\frac1{4T}\sum_k|\Fh(\omega_k)|^2\,\nuh(\omega_k)=\int\Rel\cF\,d\nu$, the series converging absolutely.
\end{lemma}

\begin{proof}
The functions $\epsilon_k(t)=(4T)^{-1/2}e^{-i\omega_kt}$, $k\in\Z$, form an orthonormal basis of
$L^2(-2T,2T)$, and for $F\in L^2(\IT)\subset L^2(-2T,2T)$ one has $\langle F,\epsilon_k\rangle=(4T)^{-1/2}\Fh(\omega_k)$.
(i) is Parseval's identity. (ii) Let $F^{\mathrm{per}}(x)=\sum_jF(x-4Tj)$, the $4T$-periodic extension (for each
$x$ at most one term is nonzero), and $G_u(t):=F^{\mathrm{per}}(t-u)$ on $(-2T,2T)$. Translation on the circle
$\R/4T\Z$ gives $\langle G_u,\epsilon_k\rangle=e^{i\omega_ku}\langle F,\epsilon_k\rangle$, so Parseval gives
$\langle F,G_u\rangle=\frac1{4T}\sum_k|\Fh(\omega_k)|^2e^{-i\omega_ku}$, while
$\langle F,G_u\rangle=\sum_j\int F(t)\overline{F(t-u-4Tj)}\,dt=\sum_j\cF(u+4Tj)$. Absolute and uniform convergence
follow from (i). Since $\cF$ is continuous and vanishes on $|x|\ge2T$, only $j=0$ survives when $|u|\le2T$.
Replacing $u$ by $u+2T$ and using $e^{-i\omega_k2T}=(-1)^k$ gives the second formula, since for
$0\le u\le2T$ only $j=-1$ survives. (iii) is (ii) at $u=0$ in the second formula: $\cF(-2T)=0$.
(iv) Taking real parts in (ii), $\frac1{4T}\sum_k|\Fh(\omega_k)|^2\cos(\omega_ku)=\Rel\cF(u)$ for $|u|\le2T$;
integrate against $\nu$, which is legitimate by (i) and $|\nu|(\R)<\infty$.
\end{proof}

\begin{theorem}[the Nyquist identity]\label{thm:nyq}
For every $F\in L^2(\IT)$,
\begin{equation}\label{eq:nyq}
\begin{aligned}
\mathfrak a_T(F)&=\frac1{4T}\sum_{k\in\Z}|\Fh(\omega_k)|^2\,s_T(\omega_k)\\
&=c_T\|F\|^2+\frac1{4T}\sum_k|\Fh(\omega_k)|^2L_T(\omega_k)
+\sum_{n\in\NT}v_n\big(\|F\|^2-\Rel\cF(\log n)\big),
\end{aligned}
\end{equation}
where the middle sum has nonnegative terms and equals $\int_0^{2T}(\|F\|^2-\Rel\cF(u))K(u)\,du\in[0,\infty]$.
Both sides are finite exactly when $F\in\mathcal D_T$, and then the first series converges absolutely. On the
mesh,
\begin{equation}\label{eq:nyq-r}
r_T(\omega_k)=(-1)^k\,2\sum_{m\ge0}e^{-(4m+1)T}\frac{2m+\frac12}{(2m+\frac12)^2+\omega_k^2}.
\end{equation}
\end{theorem}

\begin{proof}
For $n\in\NT$ apply Lemma~\ref{lem:nyq-sampling}(iv) to $\nu=\frac12(\delta_{\log n}+\delta_{-\log n})$ and
combine with (i): $\frac1{4T}\sum_k|\Fh(\omega_k)|^2(1-\cos(\omega_k\log n))=\|F\|^2-\Rel\cF(\log n)$. For
$0<\varepsilon<2T$ apply (iv) to the finite even measure $\nu_\varepsilon=\frac12K(|u|)\mathbf 1_{\varepsilon<|u|<2T}\,du$,
again with (i):
\[
\frac1{4T}\sum_k|\Fh(\omega_k)|^2\int_\varepsilon^{2T}(1-\cos\omega_ku)K(u)\,du
=\int_\varepsilon^{2T}\big(\|F\|^2-\Rel\cF(u)\big)K(u)\,du .
\]
As $\varepsilon\downarrow0$ each term on the left increases, because $1-\cos\ge0$ and $K>0$, and so does the
integrand on the right, because $|\cF(u)|\le\|F\|^2$. Monotone convergence on both sides gives
$\frac1{4T}\sum_k|\Fh(\omega_k)|^2L_T(\omega_k)=\int_0^{2T}(\|F\|^2-\Rel\cF)K\in[0,\infty]$. Since $\cF=0$ on
$[2T,\infty)$, $\Lev(\cF)$ is this integral plus $\|F\|^2\int_{2T}^\infty K$, and \eqref{eq:Aform} becomes the
right side of \eqref{eq:nyq}. By Theorem~\ref{thm:A} both sides are finite exactly on $\mathcal D_T$. Since
$|s_T(\omega_k)|\le|c_T|+2\Aconst+L_T(\omega_k)$, the first series then converges absolutely, and it equals the
second expression term by term. For \eqref{eq:nyq-r} use $K(u)=2\sum_me^{-b_mu}$, $b_m=2m+\frac12$
(Lemma~\ref{lem:K}), and $\int_{2T}^\infty\cos(\omega u)e^{-bu}du=e^{-2bT}\frac{b\cos2T\omega-\omega\sin2T\omega}{b^2+\omega^2}$,
with $\cos2T\omega_k=(-1)^k$, $\sin2T\omega_k=0$; the interchange is justified by $\sum_m e^{-2b_mT}/b_m<\infty$.
\end{proof}

The identity does not hold with $\sT$ in place of $s_T$: the part of the L\'evy measure outside the window is
not band limited, and its periodization folds back.

\begin{proposition}[aliasing]\label{prop:nyq-alias}
For real $F\in\mathcal D_T$,
\begin{align*}
\mathfrak a_T(F)-\frac1{4T}\sum_k|\Fh(\omega_k)|^2\sT(\omega_k)
&=\sum_{j\ge1}\int_{-2T}^{2T}K(4Tj-v)\,\cF(v)\,dv\\
&=\sum_{j\ge1}\sum_{m\ge0}2e^{-(2m+\frac12)4Tj}\,\Phi_F(2m+1)\,\Phi_F(-2m),
\end{align*}
with $\Phi_F(s)=\int F(t)e^{(s-\frac12)t}dt$. On $\PT$ the terms with $m=0$ vanish, since
$\Phi_F(1)\Phi_F(0)=A(F)B(F)$.
\end{proposition}

\begin{proof}
By Theorem~\ref{thm:nyq} the difference is $\frac1{4T}\sum_k|\Fh(\omega_k)|^2r_T(\omega_k)$, absolutely
convergent since $|r_T|\le\int_{2T}^\infty K$. By Fubini (Lemma~\ref{lem:nyq-sampling}(i) and
$K\in L^1(2T,\infty)$) and Lemma~\ref{lem:nyq-sampling}(ii) for real $F$, it equals
$\int_{2T}^\infty K(u)\sum_j\cF(u+4Tj)\,du$. For $u>2T$ the term $\cF(u+4Tj)$ is nonzero only if
$|u+4Tj|<2T$, which forces $j\le-1$; substituting $v=u+4Tj$ gives
$\sum_{j\ge1}\int_{-2T}^{2T}K(v+4Tj)\cF(v)dv$ after renaming $j\mapsto-j$, and $\cF$ is even. For the second form expand
$K(x)=2\sum_me^{-(2m+\frac12)x}$ for $x\ge2T$ and use $\int\cF(v)e^{bv}dv=\int F(t)e^{bt}dt\int F(s)e^{-bs}ds
=\Phi_F(b+\frac12)\Phi_F(\frac12-b)$ with $b=2m+\frac12$; all double series converge absolutely because
$e^{-b(4Tj-|v|)}\le e^{-b(4j-2)T}$ and $|\cF|\le\|F\|^2$.
\end{proof}

\begin{remark}\label{rem:nyq-num}
Numerically (Galerkin, not certified), on a minimizer at $T=\frac12\log3,\log2,\frac12\log5$ the aliasing term
of Proposition~\ref{prop:nyq-alias} is about $5\cdot10^{-4}$, $1\cdot10^{-4}$, $5\cdot10^{-5}$, comparable to or
far larger than $\lamA$ (Appendix~\ref{app:num}). So the samples $\sT(\omega_k)$ do not represent the form,
while the samples $s_T(\omega_k)$ represent it exactly; \eqref{eq:nyq} was checked numerically against the
$u$-side evaluation of $\mathfrak a_T$ to about $10^{-14}$ for smooth tests.
\end{remark}

\subsection{The window measure and the L\'evy--Khintchine form of the samples}
On $[0,2T)$ consider the signed measure
\[
\mu_T:=\Big(C_0+\int_{2T}^\infty K\Big)\delta_0+K(u)\,du-\sum_{n\in\NT}v_n\,\delta_{\log n},
\]
paired with $\cos\tau u$ through the renormalization $\int\cos(\tau u)K(u)du\leadsto-\int(1-\cos\tau u)K(u)du$ of
the nonintegrable density; with this convention its cosine transform is $s_T$. It is the restriction to the window
of the measure $C_0\delta_0+K\,du-\sum_nv_n\delta_{\log n}$ on $[0,\infty)$, the mass of $K$ beyond $2T$ being
moved to the atom at $0$. Moving the mass of the comb to the atom as well gives the following form.

\begin{proposition}[L\'evy--Khintchine form of the samples]\label{prop:nyq-LK}
Let $\rho_T:=K(u)\mathbf 1_{(0,2T)}(u)\,du+\sum_{n\in\NT}v_n\delta_{\log n}$, a positive measure on $(0,2T)$, and
$\psi_T(\tau):=\int(1-\cos\tau u)\,d\rho_T(u)\ge0$.
\emph{(i)} $s_T=c_T+\psi_T$, and $\mathfrak a_T(F)=c_T\|F\|^2+\frac1{4T}\sum_k|\Fh(\omega_k)|^2\psi_T(\omega_k)$.
\emph{(ii)} For every finitely supported $x\in\C^\Z$ with $\sum_kx_k=0$,
$\sum_{k,l}x_k\bar x_l\,s_T(\omega_{k-l})\le0$; and $\sum_{k,l}x_k\bar x_l\sum_nv_n\cos(\omega_{k-l}\log n)\ge0$ for every
finitely supported $x$.
\emph{(iii)} $s_T(\omega_k)\to+\infty$ as $|k|\to\infty$; more precisely
$\frac{T}{\sinh T}\int_0^{2T|\tau|}\frac{1-\cos x}x\,dx\le L_T(\tau)\le C+2\log(1+|\tau|)$.
\emph{(iv)} For $T\ge\frac12\log2$, $s_T(\omega_0)=c_T\le C_0+2.995<-2.37$.
\end{proposition}

\begin{proof}
(i) is the definition of $c_T$ and Theorem~\ref{thm:nyq}. (ii) For real $u$,
$\sum_{k,l}x_k\bar x_l\cos((\omega_k-\omega_l)u)=\frac12|\sum_kx_ke^{i\omega_ku}|^2+\frac12|\sum_kx_ke^{-i\omega_ku}|^2\ge0$,
which gives the second statement after summation against $v_n\ge0$; and
$\sum_{k,l}x_k\bar x_l(1-\cos(\omega_{k-l}u))=|\sum x_k|^2-(\text{that})\le0$ when $\sum x_k=0$, which integrates
against $\rho_T\ge0$ (the integrand is $O(u^2)$ at $0$); the constant $c_T$ contributes $c_T|\sum x_k|^2=0$.
(iii) $K(u)\ge\frac1{2\sinh(u/2)}\ge\frac{T}{u\sinh T}$ for $0<u\le2T$, because $\sinh y/y$ increases; this gives the
lower bound, and $\int_0^X\frac{1-\cos x}xdx\to\infty$. For the upper bound use $K(u)\le\frac1u+\frac12$ and
$1-\cos\tau u\le\min(\tau^2u^2/2,2)$. The comb is bounded by $\Aconst$. (iv) $\int_{2T}^\infty K$ decreases in $T$,
and by Theorem~\ref{thm:S}(iii) at $T=\frac12\log2$ it equals $\log\coth(\frac14\log2)+\pi-2\arctan\sqrt2<2.995$;
$\Aconst\ge0$ and $C_0=-5.3721\ldots$
\end{proof}

\begin{remark}\label{rem:nyq-circle}
By Lemma~\ref{lem:nyq-sampling}, \eqref{eq:nyq} says that $\mathfrak a_T$ is the compression to
$L^2(\IT)\subset L^2(\R/4T\Z)$ of the Fourier multiplier on the circle $\R/4T\Z$ with eigenvalues
$s_T(\omega_k)$. By Proposition~\ref{prop:nyq-LK}(iv) that multiplier is not nonnegative, and (ii) shows that the
sequence $(s_T(\omega_k))_k$ is conditionally negative definite on $\Z$ up to the constant $c_T$, while the sampled
comb is positive definite. Both properties concern the whole circle and carry no information on the sign of the
compression to $\IT$, which is the content of \eqref{eq:target} at the window $T$.
\end{remark}

\subsection{Two sectors}
Put $e_j(t):=(2T)^{-1/2}e^{-i\omega_{2j}t}$ and $e'_l(t):=(2T)^{-1/2}e^{-i\omega_{2l+1}t}$ on $\IT$, $j,l\in\Z$. Both
families are orthonormal bases of $L^2(\IT)$ (the second is the first multiplied by the unimodular function
$e^{-i\pi t/2T}$), and $\langle F,e_j\rangle=(2T)^{-1/2}\Fh(\omega_{2j})$, $\langle F,e'_l\rangle=(2T)^{-1/2}\Fh(\omega_{2l+1})$.
The $e_j$ are the eigenfunctions of $-i\,d/dt$ with periodic boundary condition $F(T)=F(-T)$, and the $e'_l$ those with
the antiperiodic condition $F(T)=-F(-T)$. For $F\in\mathcal D_T$ define
\[
\mathfrak s^{\mathrm{per}}_T(F):=\sum_js_T(\omega_{2j})|\langle F,e_j\rangle|^2,\qquad
\mathfrak s^{\mathrm{anti}}_T(F):=\sum_ls_T(\omega_{2l+1})|\langle F,e'_l\rangle|^2,
\]
and $h_m:=\dfrac{2(-1)^m}{\pi(2m+1)}$, $m\in\Z$.

\begin{theorem}[two sectors]\label{thm:nyq-sectors}
\emph{(i)} On $\mathcal D_T$ both sector forms are finite and $\mathfrak a_T=\frac12\big(\mathfrak s^{\mathrm{per}}_T+\mathfrak s^{\mathrm{anti}}_T\big)$.
\emph{(ii)} For $F\in\mathcal D_T$ the limit
\[
\Delta_T(F):=\lim_{\varepsilon\downarrow0}\int_\varepsilon^{2T}K(u)\,\Rel\cF(2T-u)\,du+\sum_{n\in\NT}v_n\,\Rel\cF(2T-\log n)
\]
exists, and $\mathfrak s^{\mathrm{per}}_T=\mathfrak a_T-\Delta_T$, $\mathfrak s^{\mathrm{anti}}_T=\mathfrak a_T+\Delta_T$.
\emph{(iii)} With $c_j=\langle F,e_j\rangle$ and $c'_l=\langle F,e'_l\rangle$ one has $c'=Hc$, where $H=(h_{l-j})_{l,j}$ is the
Laurent matrix with symbol $e^{i\theta/2}$, $|\theta|<\pi$. $H$ is real and unitary on $\ell^2(\Z)$, $H^2$ is the
bilateral shift, and
\[
\mathfrak a_T(F)=\tfrac12\big\langle c,(D_{\mathrm e}+H^{\mathsf T}D_{\mathrm o}H)c\big\rangle,\qquad
D_{\mathrm e}=\mathrm{diag}\big(s_T(\omega_{2j})\big),\ D_{\mathrm o}=\mathrm{diag}\big(s_T(\omega_{2l+1})\big).
\]
\end{theorem}

\begin{proof}
(i) By Parseval in each basis, $\frac1{2T}\sum_{k\text{ even}}|\Fh(\omega_k)|^2=\frac1{2T}\sum_{k\text{ odd}}|\Fh(\omega_k)|^2=\|F\|^2$.
Each sector sum is dominated by twice the absolutely convergent series of Theorem~\ref{thm:nyq}, and their mean is
that series.
(ii) Adding and subtracting the two formulas of Lemma~\ref{lem:nyq-sampling}(ii) and taking real parts gives, for
$0\le u\le2T$,
\[
\frac1{2T}\sum_{k\text{ even}}|\Fh(\omega_k)|^2\cos\omega_ku=\Rel\cF(u)+\Rel\cF(2T-u),\]\[
\frac1{2T}\sum_{k\text{ odd}}|\Fh(\omega_k)|^2\cos\omega_ku=\Rel\cF(u)-\Rel\cF(2T-u),
\]
where $\Rel\cF(u-2T)=\Rel\cF(2T-u)$ because $\cF(-x)=\overline{\cF(x)}$. The left sides have modulus at most $\|F\|^2$,
so the functions $\|F\|^2-\Rel\cF(u)\mp\Rel\cF(2T-u)$ are nonnegative. Repeating the proof of Theorem~\ref{thm:nyq}
for each parity (atoms at $\pm\log n$, then $\nu_\varepsilon$ and monotone convergence) gives
\begin{align*}
\mathfrak s^{\mathrm{per}/\mathrm{anti}}_T(F)=c_T\|F\|^2&+\int_0^{2T}\big(\|F\|^2-\Rel\cF(u)\mp\Rel\cF(2T-u)\big)K(u)\,du\\
&+\sum_{n\in\NT}v_n\big(\|F\|^2-\Rel\cF(\log n)\mp\Rel\cF(2T-\log n)\big),
\end{align*}
both integrals being finite on $\mathcal D_T$ by (i). Hence $\int_\varepsilon^{2T}K(u)\Rel\cF(2T-u)du$, one half of the
difference of two convergent integrals over $(\varepsilon,2T)$, converges as $\varepsilon\downarrow0$, and comparison
with \eqref{eq:nyq} gives the two identities.
(iii) $\langle e_j,e'_l\rangle=\frac1{2T}\int_{-T}^Te^{i(2l+1-2j)\pi t/2T}dt=\frac{\sin((l-j+\frac12)\pi)}{(l-j+\frac12)\pi}=h_{l-j}$,
and $c'_l=\sum_jc_j\langle e_j,e'_l\rangle$; $H$ is unitary as the matrix of a change of orthonormal basis. The numbers
$h_m=\frac1{2\pi}\int_{-\pi}^\pi e^{i\theta/2}e^{im\theta}d\theta$ are the Fourier coefficients of $e^{i\theta/2}$, so $H$ is
the Laurent operator with that unimodular symbol, and $H^2$ has symbol $e^{i\theta}$, the bilateral shift. Finally
$\mathfrak s^{\mathrm{per}}_T=\langle c,D_{\mathrm e}c\rangle$, $\mathfrak s^{\mathrm{anti}}_T=\langle Hc,D_{\mathrm o}Hc\rangle$, and
$H^*=H^{\mathsf T}$ since $H$ is real.
\end{proof}

The forms $\mathfrak s^{\mathrm{per}}_T$ and $\mathfrak s^{\mathrm{anti}}_T$ are the forms of $s_T(D_0)$ and
$s_T(D_\pi)$, where $D_\theta$ is the self-adjoint realization of $-i\,d/dt$ on $\IT$ with $F(T)=e^{i\theta}F(-T)$;
$\Delta_T$ is the pairing of $\cF$ with the reflection $u\mapsto2T-u$ of the positive part $\rho_T$ of the window
measure. In particular each sector is, like $\mathfrak a_T$, a functional of $\cF$ alone (a difference kernel on
$\IT$ with a reflected correction), not a multiplication operator in a variable dual to $j$.

\begin{corollary}[aliasing sandwich]\label{cor:nyq-sandwich}
Put $\sigma_{\mathrm e}:=\min_js_T(\omega_{2j})$ and $\sigma_{\mathrm o}:=\min_ls_T(\omega_{2l+1})$, which exist by
Proposition~\ref{prop:nyq-LK}(iii). For $F\in\mathcal D_T$,
$\mathfrak a_T(F)\ge\max\big(\sigma_{\mathrm e}\|F\|^2+\Delta_T(F),\ \sigma_{\mathrm o}\|F\|^2-\Delta_T(F)\big)$. In particular
$\mathfrak a_T(F)<0$ implies $\sigma_{\mathrm o}\|F\|^2<\Delta_T(F)<-\sigma_{\mathrm e}\|F\|^2$.
\end{corollary}

\begin{proof}
$\mathfrak s^{\mathrm{per}}_T(F)\ge\sigma_{\mathrm e}\|F\|^2$ and $\mathfrak s^{\mathrm{anti}}_T(F)\ge\sigma_{\mathrm o}\|F\|^2$ by Parseval;
apply Theorem~\ref{thm:nyq-sectors}(ii).
\end{proof}

\begin{proposition}[no common diagonalization]\label{prop:nyq-nopoint}
Let $S^{\mathrm{per}}$ and $S^{\mathrm{anti}}$ be the self-adjoint operators on $L^2(\IT)$ that are diagonal in
$\{e_j\}$ and $\{e'_l\}$ with eigenvalues $s_T(\omega_{2j})$ and $s_T(\omega_{2l+1})$. They have no common
eigenvector. In particular no orthonormal basis of $L^2(\IT)$ diagonalizes both sectors, and there is no
representation of $\mathfrak a_T$ as $\frac12$ times the sum of two multiplication operators in a common basis.
\end{proposition}

\begin{proof}
By Proposition~\ref{prop:nyq-LK}(iii) every eigenspace of $S^{\mathrm{per}}$ is spanned by finitely many $e_j$. Let
$v=\sum_jc_je_j\ne0$, $c$ finitely supported, satisfy $S^{\mathrm{anti}}v=\mu v$. Then
$(s_T(\omega_{2l+1})-\mu)(Hc)_l=0$ for all $l$. But $(Hc)_l=\frac{2(-1)^l}\pi R(l)$ with
$R(x)=\sum_jc_j(-1)^j/(2x-2j+1)$, a nonzero rational function (its poles $x=j-\frac12$ are distinct, with residues
$\frac12(-1)^jc_j$), so $(Hc)_l\ne0$ for all but finitely many $l$, and $s_T(\omega_{2l+1})=\mu$ for infinitely many
$l$, contradicting $s_T(\omega_k)\to\infty$.
\end{proof}

The only pointwise condition in the index $k$ is $s_T(\omega_k)\ge0$ for all $k$ (modulo the alternating sequence of
Lemma~\ref{lem:nyq-sampling}(iii) and the two Tate multipliers $\gh0,\gh1$ on $\PT$); it is sufficient for
$\mathfrak a_T\ge0$ on $\PT$, and it fails. On $L^2(\IT)$ this is Proposition~\ref{prop:nyq-LK}(iv). On $\PT$:

\begin{remark}\label{rem:nyq-pointwise}
Numerically (linear program over $k\le4000$, not certified), $\max\{t:\ s_T(\omega_k)+\beta_0\gh0(\omega_k)+
\beta_1\gh1(\omega_k)+c(-1)^k\ge t\ \forall k\}$ over $(\beta_0,\beta_1,c)\in\R^3$ is about $-0.12$, $-0.56$, $-0.95$,
$-1.3$, $-6.9$ at $T=\frac12\log3,\log2,\frac12\log5,1,2$. So no window in this list has a certificate that is
pointwise on the mesh; positivity of $\mathfrak a_T$ on $\PT$ at the first three windows
(Theorem~\ref{thm:log5A}) is a property of the compression.
\end{remark}

\subsection{The primal cone}

\begin{theorem}[the cone of admissible samples]\label{thm:nyq-cone}
Let $\mathcal K_T:=\{(|\Fh(\omega_k)|^2)_{k\in\Z}:F\in\PT\}$. Then $\mathcal K_T$ is the set of sequences $(f(\omega_k))_k$
where $f$ ranges over the entire functions of exponential type at most $2T$ that are nonnegative and integrable on $\R$
and vanish to order at least $2$ at $i/2$ and at $-i/2$. Every $a\in\mathcal K_T$ satisfies $\sum_k(-1)^ka_k=0$.
Moreover $\PT=\{(\frac14-\partial^2)G:\ G\in H^2_0(\IT)\}$, and $\Fh(\tau)=(\tau^2+\frac14)\widehat G(\tau)$.
\end{theorem}

\begin{proof}
If $F\in\PT$, $f(z):=\Fh(z)\overline{\Fh(\bar z)}$ is entire of type at most $2T$ (Paley--Wiener), equals $|\Fh|^2\ge0$
on $\R$ and is integrable there; $\Fh(-i/2)=A(F)=0$ and $\Fh(i/2)=B(F)=0$, so both factors vanish at each of
$\pm i/2$. Conversely let $f$ be as stated. Since $\log^+f\le f$, $\int\log^+f(x)\,dx/(1+x^2)<\infty$, and by
Akhiezer's theorem \cite[\S7.5]{Boas1954} $f(z)=G(z)\overline{G(\bar z)}$ with $G$ entire of exponential type at most
$T$. Then $|G|^2=f$ on $\R$, so $G\in L^2(\R)$ and $G=\Fh$ for some $F\in L^2(\IT)$ (Paley--Wiener). The order of $f$
at $i/2$ and at $-i/2$ is $\nu_G(i/2)+\nu_G(-i/2)\ge2$, where $\nu_G$ is the order of vanishing of $G$. If one of
$\nu_G(\pm i/2)$ is $0$, the other is $\ge2$; if, say, $\nu_G(i/2)=0$, multiplying $G$ by $(z-i/2)/(z+i/2)$, which is unimodular on $\R$ and
bounded at infinity, moves one zero from $-i/2$ to $i/2$ (symmetrically in the other case); the new $G$ has the same modulus on $\R$, the
same type, and vanishes at both points, so its $F$ lies in $\PT$. The alternating identity is
Lemma~\ref{lem:nyq-sampling}(iii). For the last statement, $F\in\PT$ makes $\Fh/(z^2+\frac14)$ entire of type at most
$T$ and square integrable on $\R$ together with $\tau^2$ times it, hence the transform of some $G\in H^2(\R)$
supported in $[-T,T]$, i.e.\ $G\in H^2_0(\IT)$, and $((\frac14-\partial^2)G)^\wedge=(\tau^2+\frac14)\widehat G$;
conversely two integrations by parts give $A=B=0$ for $(\frac14-\partial^2)G$.
\end{proof}

\begin{remark}\label{rem:nyq-dual}
By Theorems~\ref{thm:nyq} and \ref{thm:nyq-cone}, $\lamA\ge0$ if and only if $\sum_ka_ks_T(\omega_k)\ge0$ for every
$a\in\mathcal K_T$ with $\sum_ka_k\log(2+|k|)<\infty$, and since $\mathcal K_T$ lies in the hyperplane
$\sum(-1)^ka_k=0$, the form determines $(s_T(\omega_k))_k$ only modulo $\R(-1)^k$. A description of the dual cone by
band-limited smoothings of positive measures (a sampled form of Kre\u{\i}n's extension theorem \cite{Krein1940})
can at best hold modulo this alternating sequence: $(-1)^k$ annihilates $\mathcal K_T$ but is not the sampled
smoothing of a nonzero positive measure. We neither prove nor use such a description.
\end{remark}

\subsection{A two-sector Schur reduction}
Fix $\eta>0$. Let $N_{\mathrm e}:=\{j:s_T(\omega_{2j})<\eta\}$ and $N_{\mathrm o}:=\{l:s_T(\omega_{2l+1})<\eta\}$, finite by
Proposition~\ref{prop:nyq-LK}(iii), and
\[
\mathcal N_\eta:=\mathrm{span}\{e_j:j\in N_{\mathrm e}\}+\mathrm{span}\{e'_l:l\in N_{\mathrm o}\},\qquad
d_\eta:=\dim\mathcal N_\eta\le|N_{\mathrm e}|+|N_{\mathrm o}|,
\]
\[
s^{\mathrm e}_+:=\min_{j\notin N_{\mathrm e}}s_T(\omega_{2j}),\qquad s^{\mathrm o}_+:=\min_{l\notin N_{\mathrm o}}s_T(\omega_{2l+1}),\qquad
c_\eta:=\tfrac12\big(s^{\mathrm e}_++s^{\mathrm o}_+\big)\ \ge\ \eta .
\]
Let $\mathcal C_\eta:=\PT\cap\mathcal N_\eta^\perp$ and $N'_\eta:=\PT\ominus\mathcal C_\eta=\mathrm{pr}_T\mathcal N_\eta$, where
$\mathrm{pr}_T$ is the orthogonal projection onto $\PT$. The functions $e_j$, $e'_l$, $e^{\pm t/2}\mathbf 1_{\IT}$ have
$|\widehat{(\cdot)}(\omega_k)|=O(1/|k|)$, so they lie in $\mathcal D_T$ and $N'_\eta\subset\mathcal D_T$.

\begin{theorem}[two-sector Schur reduction]\label{thm:nyq-schur}
\emph{(i) Coercivity.} $\mathfrak a_T(F)\ge c_\eta\|F\|^2$ for every $F\in\mathcal D_T\cap\mathcal N_\eta^\perp$.
\emph{(ii) Exact reduction.} For $x\in N'_\eta$ the infimum $S_\eta(x):=\inf\{\mathfrak a_T(x+y):y\in\mathcal C_\eta\cap\mathcal D_T\}$ is
attained at a unique $y_*(x)$, depending linearly on $x$, so $S_\eta$ is a Hermitian form on the space $N'_\eta$ of
dimension at most $d_\eta$. Moreover $\lamA\ge0\iff S_\eta\succeq0$, $\lamA>0\iff S_\eta\succ0$, and if $\lamA\ge0$ then
$\lamA\le\min\operatorname{spec}S_\eta$.
\emph{(iii) A certificate without inversion.} For $F\in L^2(\IT)$ with $\sum_ks_T(\omega_k)^2|\Fh(\omega_k)|^2<\infty$ put
\[
g_F:=\tfrac12\Big[\sum_js_T(\omega_{2j})\langle F,e_j\rangle e_j+\sum_ls_T(\omega_{2l+1})\langle F,e'_l\rangle e'_l\Big]\in L^2(\IT).
\]
Then $\mathfrak a_T(G,F)=\langle G,g_F\rangle$ for all $G\in\mathcal D_T$, and for $x\in N'_\eta$
\[
S_\eta(x)\ \ge\ S^{\mathrm{low}}_\eta(x):=\mathfrak a_T(x)-\frac{\|P_{\mathcal C_\eta}g_x\|^2}{c_\eta},\qquad
\|P_{\mathcal C_\eta}g_x\|^2=\|g_x\|^2-\|P_{\mathcal N_\eta+\mathrm{span}\{e^{\pm t/2}\}}g_x\|^2 .
\]
Consequently $S^{\mathrm{low}}_\eta\succeq0$ implies $\lamA\ge0$, and $S^{\mathrm{low}}_\eta\succ0$ implies $\lamA>0$. The
entries of $S^{\mathrm{low}}_\eta$ in a basis of $N'_\eta$ are finite combinations of absolutely convergent series in the
samples $s_T(\omega_k)$ and the numbers $h_m$.
\emph{(iv) Monotonicity under truncation.} If $X_1\subset X_2\subset\cdots$ are closed subspaces of $L^2(\IT)$ and
$S^{\mathrm{low}}_{\eta,J}(x):=\mathfrak a_T(x)-\|P_{\mathcal C_\eta\cap X_J}g_x\|^2/c_\eta$, then $S^{\mathrm{low}}_{\eta,J}$ is
nonincreasing in $J$ and $S^{\mathrm{low}}_{\eta,J}\ge S^{\mathrm{low}}_\eta$ for every $J$.
\end{theorem}

\begin{proof}
(i) If $F\perp e_j$ for $j\in N_{\mathrm e}$ then $\mathfrak s^{\mathrm{per}}_T(F)=\sum_{j\notin N_{\mathrm e}}s_T(\omega_{2j})|\langle F,e_j\rangle|^2
\ge s^{\mathrm e}_+\|F\|^2$ by Parseval; likewise $\mathfrak s^{\mathrm{anti}}_T(F)\ge s^{\mathrm o}_+\|F\|^2$; average with
Theorem~\ref{thm:nyq-sectors}(i). $c_\eta\ge\eta$ because both minima are over samples $\ge\eta$.
(ii) Let $M\ge0$ with $\mathfrak a_T+M\|\cdot\|^2\ge\|\cdot\|^2$ on $\mathcal D_T$ (Lemma~\ref{lem:compact}); its square root is a
norm equivalent to the form norm. On $\mathcal C_\eta\cap\mathcal D_T$, which is closed in the form norm, (i) gives
$\mathfrak a_T\ge\frac{c_\eta}{c_\eta+M}(\mathfrak a_T+M\|\cdot\|^2)$, so $\mathfrak a_T$ is an inner product on it with an
equivalent complete norm. The functional $y\mapsto\mathfrak a_T(y,x)$ is bounded for this norm (Cauchy--Schwarz for
$\mathfrak a_T+M\|\cdot\|^2$), so by the Riesz representation theorem there is a unique $y_*(x)$, linear in $x$, with
$\mathfrak a_T(y,x+y_*(x))=0$ for all $y\in\mathcal C_\eta\cap\mathcal D_T$, and then
$\mathfrak a_T(x+y)=S_\eta(x)+\mathfrak a_T(y-y_*(x))\ge S_\eta(x)$ with $S_\eta(x)=\mathfrak a_T(x+y_*(x))$. Every
$F\in\PT\cap\mathcal D_T$ is $x+y$ with $x\in N'_\eta$, $y\in\mathcal C_\eta\cap\mathcal D_T$. If $\lamA\ge0$ then
$S_\eta(x)\ge\lamA\|x+y_*\|^2\ge\lamA\|x\|^2\ge0$. If $S_\eta\succeq0$ then $\mathfrak a_T(x+y)\ge0$. If
$S_\eta\ge s_0\|\cdot\|^2$ with $s_0>0$, let $\|y_*(x)\|\le C\|x\|$; for $y=y_*(x)+z$,
$\mathfrak a_T(x+y)=S_\eta(x)+\mathfrak a_T(z)\ge s_0\|x\|^2+c_\eta\|z\|^2$, while $\|x+y\|^2\le(1+2C^2)\|x\|^2+2\|z\|^2$, so
$\lamA>0$.
(iii) $\mathfrak a_T=\frac12(\mathfrak s^{\mathrm{per}}_T+\mathfrak s^{\mathrm{anti}}_T)$ as quadratic forms on $\mathcal D_T$, hence as
sesquilinear forms, and $\mathfrak s^{\mathrm{per}}_T(G,F)=\sum_js_T(\omega_{2j})\langle G,e_j\rangle\overline{\langle F,e_j\rangle}$
converges absolutely by Cauchy--Schwarz; this is $\langle G,g_F\rangle$, with $g_F\in L^2$ by the hypothesis. The
hypothesis holds for $e_j,e'_l,e^{\pm t/2}\mathbf 1_{\IT}$, since $|\widehat{(\cdot)}(\omega_k)|=O(1/|k|)$ and
$s_T(\omega_k)=O(\log|k|)$ (Proposition~\ref{prop:nyq-LK}(iii)), hence for every $x\in N'_\eta$. For
$y\in\mathcal C_\eta\cap\mathcal D_T$,
$\mathfrak a_T(x+y)=\mathfrak a_T(x)+2\Rel\langle y,P_{\mathcal C_\eta}g_x\rangle+\mathfrak a_T(y)
\ge\mathfrak a_T(x)-2\|y\|\|P_{\mathcal C_\eta}g_x\|+c_\eta\|y\|^2\ge S^{\mathrm{low}}_\eta(x)$. Since
$\mathcal C_\eta^\perp=\mathcal N_\eta+\mathrm{span}\{e^{\pm t/2}\mathbf 1_{\IT}\}$ is finite dimensional, the formula for
$\|P_{\mathcal C_\eta}g_x\|^2$ follows; $\langle g_x,g_y\rangle$, $\langle g_x,e_j\rangle$, $\langle g_x,e'_l\rangle$ and
$\mathfrak a_T(x,y)$ are absolutely convergent series in $s_T(\omega_k)$ and $h_m$ by Theorem~\ref{thm:nyq-sectors}(iii).
The two implications follow from (ii).
(iv) $\|P_{\mathcal C_\eta\cap X_J}g\|\le\|P_{\mathcal C_\eta\cap X_{J+1}}g\|\le\|P_{\mathcal C_\eta}g\|$.
\end{proof}

By (iv), a negative value of $S^{\mathrm{low}}_{\eta,J}$ for one $J$ shows that $S^{\mathrm{low}}_\eta$ is not
positive semidefinite, while positive values for finitely many $J$ prove nothing without a bound for the remainder.
The coercivity in (i) costs nothing: it is the two-basis structure of Theorem~\ref{thm:nyq-sectors}, in contrast with
the coercivity of the archimedean form on block-Tate-orthogonal functions used in Theorem~\ref{thm:coer-schur}.

\begin{remark}\label{rem:nyq-schur-num}
Numerically (Galerkin truncation $X_J$ of the complement with up to $800$ periodic modes; not certified, the
remainder of (iv) not being bounded), $S^{\mathrm{low}}_\eta\succ0$ with smallest eigenvalue about $0.7\lamA$ at
$T=\frac12\log3$ ($\eta=0.3$, $d_\eta=7$), about $0.9\lamA$ at $T=\log2$ ($\eta=1$, $d_\eta=29$) and about $0.95\lamA$ at
$T=\frac12\log5$ ($\eta=1.5$, $d_\eta$ between $45$ and $59$), the values decreasing in $J$ as (iv) requires; for the
form with a quartet of Section~\ref{sec:wit} at $T=\log2$, $\delta=0.01$, it is negative, about $-1.7\cdot10^{-6}$. The
smallest $\eta$ that works grows with $T$ ($0.3$, $0.6$, $1.5$). The loss $S_\eta-S^{\mathrm{low}}_\eta$ in the direction of a
minimizer is governed by the part of the minimizer outside $\mathcal N_\eta$, a quantity not tied to the size of $\lamA$;
if that part decayed like a power of the number of mesh modes, $d_\eta$ would have to grow like a negative power of
$\lamA$. This is a heuristic, not a theorem.
\end{remark}


\section{First contact}\label{sec:firstcontact}

This section studies a hypothetical first window at which positivity of the primitive form fails. We
prove that such a window carries a primitive null vector of definite parity, a finite operator identity in
which additive Poisson summation and the convolution identity $\Lambda*1=\log$ occur together
(Theorem~\ref{thm:PCfinite}), a faithful quadratic observation of a test function by the positive odd poles
of the Gamma factor \eqref{eq:Gammaenergy}, exact identities for the exterior collars produced by compressing
the prime shifts to the window \eqref{eq:collar}, an exact balance identity on the kernel \eqref{eq:Qbalance},
and a strict contraction of the Euler channel (Proposition~\ref{prop:GEgap}). The contraction does not by
itself give vanishing: on the kernel, the natural reverse inequality is equivalent to the vanishing of the
detector. What is missing is the annihilation statement of Problem~\ref{prob:annihilate}. Nothing in this
section extends the certified range.

\subsection{The first-contact reduction}
Throughout, $\mathfrak a_T(F,G)$ is the sesquilinear form of \eqref{eq:form}, $H_T$ the operator of
Lemma~\ref{lem:compact}, and $\mathrm{pr}_T$ the orthogonal projection of $L^2(\IT)$ onto $\PT$; $J$ preserves
$\PT$ and commutes with $\mathrm{pr}_T$. The orthogonal complement of $\PT$ in $L^2(\IT)$ is the
two-dimensional \emph{Tate space} $\mathrm{span}\{e^{t/2},e^{-t/2}\}$ (restricted to $\IT$).

\begin{lemma}[first-contact reduction]\label{lem:firstcontact}
If $\lamA<0$ for some $T$, then there is a window $T_*>\frac12\log5$ and a nonzero $F\in\mathcal P_{T_*}$ such
that $\lambda_{\mathrm A}(T)\ge0$ for $T\le T_*$, $\lambda_{\mathrm A}(T)<0$ for $T>T_*$, and
\[
 H_{T_*}\ge0,\qquad H_{T_*}F=0,\qquad JF=\varepsilon F,\quad \varepsilon\in\{\pm1\}.
\]
\end{lemma}
\begin{proof}
Let $T_*=\inf\{T:\lamA<0\}$. By Lemma~\ref{lem:monoT}, $\lamA<0$ for every $T>T_*$ and $\lamA\ge0$ for every
$T<T_*$; by Lemma~\ref{lem:contA}, $\lambda_{\mathrm A}(T_*)=0$. Since $\lambda_{\mathrm A}(\frac12\log5)>0$
(Theorem~\ref{thm:log5A}), monotonicity gives $T_*>\frac12\log5$. By Lemma~\ref{lem:compact} the value $0$ is
an eigenvalue of $H_{T_*}$, and since $\ker H_{T_*}$ is $J$-invariant it contains a nonzero vector of
definite parity.
\end{proof}

\begin{remark}[the complete form]\label{rem:completecontact}
Lemma~\ref{lem:contT} gives the same reduction for the complete bottom $\lamW$: if it ever becomes negative, at its
first zero it is attained by a nonzero vector in the kernel of the nonnegative complete form. That null vector need
not be primitive, and its primitive component is governed by the Schur complement including the polar block
(Theorem~\ref{thm:parityschur}), not by $H_T$. The identities below are proved for $\ker H_T$, not for that kernel;
with Lemma~\ref{lem:contA} the primitive reduction is unconditional, and the complete form is not needed below.
\end{remark}



We call a nonzero $F\in\ker H_T$ with $JF=\varepsilon F$, at a window where $H_T\ge0$, a
\emph{first-contact vector}, and $\ker H_T$ the \emph{radical}. Equivalently, $F\in\PT$ satisfies
$\mathfrak a_T(F,\mathrm{pr}_TG)=0$ for every $G$ in the form domain, i.e.\ the weak equation
\begin{equation}\label{eq:localweak}
 \mathfrak a_T(F,G)=\langle r,G\rangle\qquad(G\ \text{in the form domain of}\ L^2(\IT))
\end{equation}
holds for a unique $r$ in the Tate space. More generally, for $F$ in the domain of $H_T$ we define
$r\in(\PT)^\perp$ by $\mathfrak a_T(F,G)=\langle H_TF+r,G\rangle$.

\subsection{Operators on the line}

Let $XF(t)=tF(t)$ on the Schwartz space $\mathcal S(\R)$ (the \emph{Schwartz
core}). Let $\mathcal M=m(D)$ be the Fourier multiplier with the symbol $m$ of \eqref{eq:symbol},
$(\mathcal MF)\,\widehat{}\,=m\Fh$; we call it the completed Gamma multiplier. Let
$u_\infty(\tau)=\pi^{-i\tau}\Gamma(\frac14+\frac{i\tau}2)/\Gamma(\frac14-\frac{i\tau}2)$, a
unimodular function with $u_\infty(-\tau)=\overline{u_\infty(\tau)}$ and, by Tate's local
functional equation, $\frac{d}{d\tau}\arg u_\infty=m$. Put
\[
 \mathscr F=u_\infty(D)\,J .
\]
This is the additive Fourier transform in logarithmic coordinates: if $F(t)=e^{t/2}f(e^t)$ with $f$
an even Schwartz function on $\R$, then $\mathscr FF(t)=e^{t/2}\widehat f(e^t)$. On the Schwartz
core
\[
 \mathscr F^*=\mathscr F^{-1}=\mathscr F,\qquad
 JV_aJ=V_{-a},\qquad X+\mathscr FX\mathscr F=\mathcal M .
\]
The first identity follows from $Ju_\infty(D)J=u_\infty(D)^*$. For the third, if $U=e^{i\theta}(D)$
with $\theta$ real and of polynomial growth with its first two derivatives, then
$U^*XU=X+\theta'(D)$ (differentiate $\widehat{XG}=-i\partial_\tau\widehat G$, or see the time-delay
identity $\langle UG,XUG\rangle=\langle G,XG\rangle+\frac1{2\pi}\int\theta'|\widehat G|^2$); with
$JX=-XJ$ this gives $\mathscr FX\mathscr F=-UXU^*=-X+m(D)$.

Shifts are written with the right shift $V_a$ of \S\ref{subsec:conv}, acting on the whole line; so
$V_{-\log n}F(t)=F(t+\log n)$.

\subsection{The finite sum--product identity}

For integers $M,N\ge1$ define on the Schwartz core
\begin{align*}
 S_N&=\sum_{n\le N}n^{-1/2}V_{-\log n},&
 S'_N&=\sum_{n\le N}(\log n)n^{-1/2}V_{-\log n},\\
 L_M&=\sum_{d\le M}\Lambda(d)d^{-1/2}V_{-\log d},&
 A_M&=\mathcal M-(L_M+L_M^*),\\
 E_N&=S_N\mathscr F-JS_N,&
 H_{M,N}&=L_MS_N-S'_N,
\end{align*}
and $\{X,E\}=XE+EX$. For the full integer comb, Poisson summation formally reads
$S_\infty\mathscr F=JS_\infty$; thus $E_N$ is the finite Poisson defect, and $H_{M,N}$ is the
incomplete divisor hyperbola.

\begin{theorem}[finite Poisson--Chebyshev identity]\label{thm:PCfinite}
On the Schwartz core,
\begin{equation}\label{eq:PCfinite}
 S_NA_M\mathscr F
 =\{X,E_N\}-L_M^*E_N-JH_{M,N}-H_{M,N}\mathscr F.
\end{equation}
If $M\ge N$, then
\[
 H_{M,N}=\sum_{k>N}h_{M,N}(k)k^{-1/2}V_{-\log k},\qquad
 h_{M,N}(k)=\sum_{\substack{dm=k\\d\le M,\ m\le N}}\Lambda(d)\ge0.
\]
In particular $H_{M,N}$ has no coefficient at $k\le N$.
\end{theorem}
\begin{proof}
The commutator $[X,S_N]=-S'_N$ and the three Fourier identities above give
\[
 S_N\mathcal M\mathscr F=JS'_N+S'_N\mathscr F+\{X,E_N\}.
\]
Insert $S'_N=L_MS_N-H_{M,N}$, use $JL_M=L_M^*J$ and
$JS_N=S_N\mathscr F-E_N$, and subtract $S_N(L_M+L_M^*)\mathscr F$.
This proves \eqref{eq:PCfinite}. The coefficient of
$k^{-1/2}V_{-\log k}$ in $L_MS_N$ is the displayed divisor sum. For
$k\le N\le M$ every divisor occurs, so $\sum_{d\mid k}\Lambda(d)=\log k$ cancels the
coefficient of $S'_N$.
\end{proof}

The identity is specific to $\Lambda$: replacing $\Lambda$ by other coefficients $a(d)$ leaves, at
$k\le N$, the coefficient $\sum_{d\mid k}a(d)-\log k$. On the other hand the square of either side
of \eqref{eq:PCfinite} is invariant under $A_M\mapsto-A_M$, so squaring it cannot determine the sign
of the Weil form.

% ("The identity is specific to $\Lambda$ ... cannot determine the sign of the Weil form.")
\begin{remark}[renormalization]\label{rem:Sren}
The operators $S_N$ have no limit as $N\to\infty$ on the critical half-density. A natural renormalization is
\[
 S_N^{\rm ren}=S_N-C_N,\qquad C_Nh(t)=e^{-t/2}\int_{-\infty}^{t+\log N}e^{u/2}h(u)\,du ,
\]
defined modulo the homogeneous solution $e^{-t/2}$; $C_N$ is the continuum counterpart of the integer comb. For
smooth compactly supported primitive $h$, the difference $S_N^{\rm ren}h$ converges locally uniformly as $N\to\infty$ by
partial summation and the prime number theorem with its classical error term, and the limit extends continuously to
distributions by moving derivatives onto the test function. This weak limit does not control a Hilbert norm of the
limit, the projection $\mathrm{pr}_T$ applied to it, or its Gamma residues: an exterior exponential tail can tend to zero
in every fixed Sobolev norm while keeping a fixed meromorphic residue. None of this is used in a proof. For this reason every quadratic expression
below is formed at finite $N$, before any limit.
\end{remark}



\subsection{A faithful observation by the Gamma poles}

Fix $\delta>0$ and put, for $m\ge1$,
\[
 a_m=2m+\tfrac12,\qquad
 b_m=\sqrt{a_m}\,e^{-a_m(T+\delta)},\qquad
 \mathcal L_F(a)=\int_{-T}^T e^{-at}F(t)\,dt,\qquad
 (\mathcal OF)_m=b_m\mathcal L_F(a_m).
\]
Since $\|e^{-a_mt}\|^2_{L^2(\IT)}\le2Te^{2a_mT}$, the map $\mathcal O:L^2(\IT)\to\ell^2$ is Hilbert--Schmidt,
and
\begin{equation}\label{eq:Gammaenergy}
 \|\mathcal OF\|^2
 =\sum_{m\ge1}a_me^{-2a_m(T+\delta)}|\mathcal L_F(a_m)|^2
 =\iint F(t)\overline{F(u)}\,\bigl[-(k_\Gamma^+)'(2T+2\delta+t+u)\bigr]\,dt\,du,
\end{equation}
where
\[
 k_\Gamma^+(v)=\sum_{m\ge1}e^{-a_mv}.
\]
This is a positive Hankel form, and it is faithful: $\mathcal OF=0$ implies $F=0$. Indeed, after the
substitution $x=e^{-2t}$, vanishing of every $\mathcal L_F(a_m)$ says that $x^{1/4}F(-\frac12\log x)\,x^{-1}$
is orthogonal to every power $x^m$, $m\ge1$, on the compact interval $[e^{-2T},e^{2T}]$. Since this
interval does not contain $0$, the functions $x\,p(x)$, $p$ a polynomial, are dense in its continuous
functions; hence $F=0$.

\begin{remark}[interpretation]\label{rem:residue}
The numbers $\mathcal L_F(a_m)$ are the residues at the positive odd poles of the Gamma factor. Write
$s=\frac12+z$, $\chi(s)=\pi^{s-1/2}\Gamma(\frac{1-s}2)/\Gamma(\frac s2)$, so that
$\zeta(s)=\chi(s)\zeta(1-s)$, and $d_\Gamma=-\chi'/\chi$. Near $s_m=2m+1$ one has
$\psi(\frac{1-s}2)=2/(s-s_m)+O(1)$, hence $\operatorname{Res}_{s=s_m}d_\Gamma=1$. For a
first-contact vector $F$ of parity $\varepsilon$ let $W^0_\varepsilon(s)$ be the corrected product
defect: the function $d_\Gamma(s)\mathcal L_F(s-\frac12)$ minus the Laplace transforms of the two compactly
supported Euler histories of $F$ and minus two polar corrections supported outside $\IT$. We use only
two properties of it: its inverse transform restricts on $\IT$ to the Tate vector $r$ of
\eqref{eq:localweak}, and every term other than $d_\Gamma \mathcal L_F$ is holomorphic at each $s_m$. Hence
\begin{equation}\label{eq:residue}
 \operatorname*{Res}_{s=s_m}\bigl[\zeta(s)W^0_\varepsilon(s)\bigr]
 =\zeta(s_m)\mathcal L_F(a_m).
\end{equation}
Consequently holomorphy of $\zeta W^0_\varepsilon$ at every $s_m$, and vanishing of
\eqref{eq:Gammaenergy}, are each sufficient for $F=0$. Formula \eqref{eq:residue} places the
obstruction outside the window: corrections supported in a compact interval have entire Laplace
transforms and cannot cancel these residues. The defect $W^0_\varepsilon$ is described here only through the two properties stated, and nothing below depends on this remark.
\end{remark}

\subsection{Compressed collars and the balance on the radical}\label{subsec:collars}

For $N\ge2$ put
\[
 B_N=Q_T\sum_{2\le n\le N}\Lambda(n)n^{-1/2}V_{\log n}Q_T,\qquad
 d_{m,N}=\sum_{2\le n\le N}\frac{\Lambda(n)}{n^{a_m+1/2}},\qquad
 D_N=\operatorname{diag}(d_{m,N})_{m\ge1}.
\]
For $N\ge e^{2T}$, $B_N$ contains every shift that acts on $L^2(\IT)$, and
$\mathfrak a_T(F,G)=\mathfrak m_T(F,G)-\langle(B_N+B_N^*)F,G\rangle$, where
\[
 \mathfrak m_T(F,G)=\frac1{2\pi}\int_\R m(\tau)\,\Fh(\tau)\overline{\widehat G(\tau)}\,d\tau
\]
is the compressed Gamma form, i.e.\ the form of $\mathcal M$ on the form domain of $L^2(\IT)$.
Thus \eqref{eq:localweak} is the weak form of $(\mathcal M-B_N-B_N^*)F=r$ on $\IT$.
The exterior collar $\mathcal E_N:L^2(\IT)\to\ell^2$ is
\[
 (\mathcal E_NF)_m=b_m\sum_{2\le n\le N}\Lambda(n)n^{-(a_m+1/2)}
 \int_{\max(-T,T-\log n)}^T e^{-a_mu}F(u)\,du.
\]
The substitution $u=t-\log n$ gives the exact identity
\begin{equation}\label{eq:collar}
 \mathcal OB_N=D_N\mathcal O-\mathcal E_N.
\end{equation}
Primes with $n\ge e^{2T}$ do not contribute to $B_N$ but do contribute to $\mathcal E_N$: once
$N\ge e^{2T}$, one has $\mathcal E_{N'}-\mathcal E_N=(D_{N'}-D_N)\mathcal O$ for $N'\ge N$, so these terms
cannot be discarded.

Let $N\ge e^{2T}$, let $\mathcal O_L,D_{N,L},\mathcal E_{N,L}$ denote the truncations to $1\le m\le L$, and put
$G_L=\mathcal O_L^*\mathcal O_L$ and $\Sigma_L=G_L+JG_LJ$. The operator $\Sigma_L$ commutes with $J$ and its kernel is a finite
sum of exponentials on $\overline{\IT}$, so $\mathrm{pr}_T\Sigma_LF$ lies in the form domain and is an admissible
test. Let $F$ be a first-contact vector with Tate vector $r$. Taking $G=\Sigma_LF$ in
\eqref{eq:localweak}, using $JB_NJ=B_N^*$, equal parity of $F$ and $\Sigma_LF$, and \eqref{eq:collar}, one
obtains
\begin{equation}\label{eq:Qbalance}
 0=2c_{N,L}(F)+Q_{N,L}(F),\qquad
 c_{N,L}(F)=\langle \mathcal O_LF,D_{N,L}\mathcal O_LF\rangle,
\end{equation}
where, with no further boundary term,
\begin{align}\label{eq:Qfull}
 Q_{N,L}(F)={}&2\Rel\langle \mathcal O_LF,\mathcal O_LB_N^*F\rangle
 -2\Rel\langle \mathcal O_LF,\mathcal E_{N,L}F\rangle\notag\\
 &-\Rel\mathfrak m_T(F,\Sigma_LF)+\Rel\langle r,\Sigma_LF\rangle.
\end{align}
The backward term involves the truncated moments
\[
 (\mathcal O_LB_N^*F)_m=b_m\sum_{2\le n<e^{2T}}
 \Lambda(n)n^{a_m-1/2}\int_{-T+\log n}^T e^{-a_mu}F(u)\,du;
\]
replacing them by full moments would remove a collar.

Identity \eqref{eq:Qbalance} carries no sign information. Indeed, for $F$ in the domain of $H_T$
with $r$ defined as above, the same computation gives
\begin{equation}\label{eq:Qrecombine}
 Q_{N,L}=-2c_{N,L}-\Rel\mathfrak a_T(F,\mathrm{pr}_T\Sigma_LF),
\end{equation}
and for an eigenvector $H_TF=\lambda F$ of parity $\varepsilon$, since
$\langle F,\Sigma_LF\rangle=2\|\mathcal O_LF\|^2$,
\begin{equation}\label{eq:Qspectral}
 Q_{N,L}(F)=-2\sum_{m\le L}(d_{m,N}+\lambda)|(\mathcal O_LF)_m|^2.
\end{equation}
Thus $Q_{N,L}$ is negative on every eigenvector with $\lambda\ge0$ and $\mathcal O_LF\ne0$. In particular, the
inequality $Q_{N,L}\ge0$ on the radical is equivalent to $\mathcal O_LF=0$ there, and $Q_{N,L}$ admits no
representation as a sum of squares valid on all of $\PT$.

\subsection{The Gamma--Euler contraction and what it does not prove}

The Euler channel $D_N$ is a strict contraction, uniformly in $m$ and $N$.

\begin{proposition}[strict Gamma--Euler contraction]\label{prop:GEgap}
For all $m\ge1$ and $N\ge2$,
\[
 0<d_{m,N}\le q_0:=\frac{\log2}{4}+\frac1{16}<0.236<1.
\]
Consequently
\begin{equation}\label{eq:gapop}
 \Delta_{N,L}:=\mathcal O_L^*(I-D_{N,L})\mathcal O_L
 \ge(1-q_0)\mathcal O_L^*\mathcal O_L,
\end{equation}
and, keeping the collar,
\begin{equation}\label{eq:gapcollar}
 \langle F,\Delta_{N,L}F\rangle
 =\|\mathcal O_LF\|^2-\Rel\langle \mathcal O_LF,\mathcal O_LB_NF\rangle
  -\Rel\langle \mathcal O_LF,\mathcal E_{N,L}F\rangle.
\end{equation}
For a first-contact vector $F$,
\begin{equation}\label{eq:gapQ}
 2\langle F,\Delta_{N,L}F\rangle=2\|\mathcal O_LF\|^2+Q_{N,L}(F).
\end{equation}
\end{proposition}
\begin{proof}
Since $a_m+\frac12=2m+1\ge3$, $\Lambda(n)\le\log n$, and $x^{-3}\log x$ is decreasing on $[2,\infty)$,
\[
 d_{m,N}\le\sum_{n\ge2}\frac{\log n}{n^3}
 \le\frac{\log2}{8}+\int_2^\infty\frac{\log x}{x^3}\,dx
 =\frac{\log2}{4}+\frac1{16}.
\]
Positivity holds because the term $n=2$ is present. Equation \eqref{eq:gapop} follows coordinatewise,
and \eqref{eq:gapcollar} follows by pairing \eqref{eq:collar} with $\mathcal O_LF$. Finally, on the radical
\eqref{eq:Qbalance} gives $Q_{N,L}=-2\langle \mathcal O_LF,D_{N,L}\mathcal O_LF\rangle$.
\end{proof}

Letting $N\to\infty$, $d_m:=\lim_Nd_{m,N}=-\zeta'(s_m)/\zeta(s_m)$. Hence the residue vector
$x_m=\zeta(s_m)b_m\mathcal L_F(a_m)$ and the Chebyshev jet $j_m=-\zeta'(s_m)b_m\mathcal L_F(a_m)$ satisfy, with
$D=\operatorname{diag}(d_m)$,
\begin{equation}\label{eq:smallgain}
 j=D x,\qquad \|j\|\le q_0\|x\|.
\end{equation}
This is a quantitative identity, but it supplies no independent leverage toward vanishing.  For a
fixed $L$ injective on the radical and $F\in\ker H_T$, the following are equivalent:
\begin{equation}\label{eq:equivgate}
 \mathcal O_LF=0
 \quad\Longleftrightarrow\quad
 \langle F,\Delta_{N,L}F\rangle\le0
 \quad\Longleftrightarrow\quad
 Q_{N,L}(F)\ge0.
\end{equation}
The first equivalence follows from \eqref{eq:gapop}; the second follows from
$Q_{N,L}=-2\langle \mathcal O_LF,D_{N,L}\mathcal O_LF\rangle$ and strict positivity of $D_{N,L}$.
Likewise, because $j=Dx$ with every diagonal entry of $D$ in $(0,q_0]$,
\[
 \|x\|^2\le\Rel\langle x,j\rangle\quad\Longleftrightarrow\quad x=0.
\]
Thus neither the finite return inequality nor its infinite-channel version is an intermediate
estimate: each is the desired conclusion in another notation.  Cauchy--Schwarz and $q_0<1$ merely
verify this equivalence.

\begin{problem}[Gamma-moment annihilation at first contact]\label{prob:annihilate}
Starting from the finite Poisson--Chebyshev identity \eqref{eq:PCfinite}, prove directly that for every primitive first-contact
vector $F$ there is a finite $L$, injective on the radical, for which
\begin{equation}\label{eq:annihilate}
 \mathcal O_LF=0.
\end{equation}
The proof must use all three pieces absent from the formal radical equation: (i) minimality of the
contact window, not merely $H_TF=0$; (ii) an oriented causal or Poisson--Tate identity which changes
under reversal of the Weil sign; and (iii) $\Lambda*1=\log$ in the step which forces vanishing or
fixes the sign, not only in the preliminary contraction estimate.  It must retain both Fourier
orientations, the continuum polar term, the full exterior collars and the crossed Tate endpoint
before taking a limit.
\end{problem}

Condition (i) is a falsification test: by \eqref{eq:Qspectral}, reasoning which also applied to
positive eigenvalues would give the wrong conclusion.  Condition (ii) excludes norm-squaring of
\eqref{eq:PCfinite}, which is invariant under $A_M\mapsto-A_M$.  Condition (iii) retains the genuine
arithmetic advantage of this route: changing one von Mangoldt weight creates an interior divisor
defect, while Poisson and a functional equation alone also occur for Epstein/Davenport--Heilbronn
examples without an Euler product.

A finite $L$ with $\mathcal O_L$ injective on $\ker H_T$ exists at each fixed window. The kernel is
finite-dimensional because $H_T$ has compact resolvent; the subspaces $\ker H_T\cap\ker \mathcal O_L$ decrease
with $L$ and have zero intersection by faithfulness of $\mathcal O$; a decreasing chain of subspaces of a
finite-dimensional space stabilizes, so one of them is $\{0\}$. The argument gives no bound on $L$
that is uniform in $T$.

\begin{proposition}[annihilation is window positivity]\label{prop:annihilate-equiv}
Call a \emph{first-contact vector} a nonzero $F\in\ker H_{T_*}$ of definite parity at a window $T_*$ as in
Lemma~\ref{lem:firstcontact}, that is, with $\lambda_{\mathrm A}(T)\ge0$ for $T\le T_*$ and $\lambda_{\mathrm A}(T)<0$ for
$T>T_*$ (the minimality of the window is part of the definition here, as in Problem~\ref{prob:annihilate}(i), unlike the
wider use of the term before \eqref{eq:localweak}). The following are equivalent:
\emph{(a)} \eqref{eq:target};
\emph{(b)} there is no first-contact vector at any window;
\emph{(c)} for every first-contact vector $F$ there is a finite $L$ with $\mathcal O_L$ injective on $\ker H_{T_*}$ and
$\mathcal O_LF=0$, the conclusion of Problem~\ref{prob:annihilate};
\emph{(d)} every first-contact vector satisfies $\mathcal OF=0$.
\end{proposition}
\begin{proof}
(a)$\Rightarrow$(b): under \eqref{eq:target} no window has $\lambda_{\mathrm A}<0$ beyond it. (b)$\Rightarrow$(a) is
Lemma~\ref{lem:firstcontact}. (b)$\Rightarrow$(c) and (b)$\Rightarrow$(d) are vacuous. (c)$\Rightarrow$(b): $\mathcal O_LF=0$
with $\mathcal O_L$ injective on $\ker H_{T_*}\ni F$ forces $F=0$. (d)$\Rightarrow$(b): $\mathcal O$ is faithful.
\end{proof}
Problem~\ref{prob:annihilate} is therefore not a weaker statement than \eqref{eq:target} but \eqref{eq:target} itself,
in the notation of this section: it asks for a proof of \eqref{eq:target} that proceeds through the Gamma moments
of a hypothetical first-contact vector and uses the three ingredients (i)--(iii). No intermediate estimate
toward it is known beyond the identities of this section.

% "A finite $L$ with $\mathcal O_L$ injective on $\ker H_T$ exists at each fixed window. ... uniform in $T$."
\begin{remark}[on moment annihilation]\label{rem:returnremarks}
The following shortcuts do not solve Problem~\ref{prob:annihilate}. The weak renormalized limit of
Remark~\ref{rem:Sren} cannot be squared, and it does not make the extraction of residues continuous. The identity
$Q_{N,L}=-2c_{N,L}$ on the radical cannot be combined with an inequality $Q_{N,L}\ge0$, since that inequality is
equivalent to the desired conclusion by \eqref{eq:equivgate}, and \eqref{eq:Qspectral} excludes a proof of it valid on all
of $\PT$. The variance of a sum labelled by divisor paths is not the norm of the coherent sum over $k$, which contains
Gram cross terms between different $k$. Additive phases diagonalize the sum over $k$ by Parseval but do not factor
through $de=k$, whereas multiplicative characters factor through $de=k$ but replace the Gamma--Poisson operator by a
twisted one, different from the $H_T$ whose kernel is under study. Unitarity of $\mathscr F$ gives only the contractive
direction, and reversing it without a new oriented identity for the compressed polar and Tate terms amounts to
assuming \eqref{eq:annihilate}.
\end{remark}


\section{Witnesses}\label{sec:wit}

A failure of \eqref{eq:target} is witnessed by a single test function. This section describes the set of such
functions, shows that it can only grow with $T$ and that a witness never changes once it exists, and compares it
with a family of model forms obtained by adding finitely many quartets to the zero side. The model forms share every
structural property used below and do have witnesses; this delimits what a proof of \eqref{eq:target} by structural
arguments can use. We end with the coherence of the prime comb on the mesh of \S\ref{sec:nyq}. Nothing in this section
decides the sign of $\lamA$ at a window $T>\frac12\log5$.

Since $\sT$ is even and real, $\mathfrak a_T(F)=\mathfrak a_T(\Rel F)+\mathfrak a_T(\Iml F)$, and $F\in\PT$ if and only if
$\Rel F,\Iml F\in\PT$ (proof of Proposition~\ref{prop:finite}(ii)); real test functions therefore suffice for every
statement about signs, and $H_T$ has real eigenfunctions.

\subsection{The witness set}

\begin{definition}\label{def:wit}
$\mathscr W_T:=\{F\in\PT\cap\mathcal D_T:\ \mathfrak a_T(F)<0\}$ is the set of \emph{witnesses} at $T$, and
$\mathscr T:=\{T>0:\ \mathscr W_T\ne\emptyset\}$.
\end{definition}

\begin{lemma}[structure of the witness set]\label{lem:wit-basic}
\emph{(i)} $\mathscr W_T\ne\emptyset$ $\iff$ $\lamA<0$ $\iff$ every minimizer of $\mathfrak a_T$ on the unit sphere of
$\PT\cap\mathcal D_T$ lies in $\mathscr W_T$.
\emph{(ii)} $\mathscr W_T$ is a cone ($cF\in\mathscr W_T$ for $F\in\mathscr W_T$, $c\in\C\setminus\{0\}$), open in the form norm,
and $J$-invariant; if $F\in\mathscr W_T$ then its even part or its odd part lies in $\mathscr W_T$.
\emph{(iii)} If $F\in\mathscr W_T$, then $\Lev(\cF)<(\Aconst-C_0)\|F\|^2$, and for every $R>0$ with $m(R)>C_0$
\[
\frac1{2\pi\|F\|^2}\int_{|\tau|>R}|\Fh|^2\,d\tau\ <\ \frac{\Aconst-C_0}{m(R)-C_0}.
\]
\emph{(iv)} If $F\in\mathscr W_T$ and $\mT(F)\ge0$, then $\sum_{n\in\NT}v_n\Rel\cF(\log n)>\mT(F)\ge0$; so a witness either
makes the archimedean form negative or has an autocorrelation positively aligned with the comb.
\emph{(v)} $\lamA\ge C_0-\Aconst$.
\end{lemma}

\begin{proof}
(i) The bottom is attained (Lemma~\ref{lem:compact}). (ii) Homogeneity of degree two; continuity of $\mathfrak a_T$
in the form norm; $J$ preserves $\PT$ and $\mathfrak a_T(JF)=\mathfrak a_T(F)$ because $\sT$ is even and $A(JF)=B(F)$. If
$F_\pm=\frac12(F\pm JF)$, then $A(F_\pm)=\frac12(A(F)\pm B(F))=0$ and likewise $B(F_\pm)=0$; and
$\mathfrak a_T(F_+,F_-)=\mathfrak a_T(JF_+,JF_-)=-\mathfrak a_T(F_+,F_-)$, so $\mathfrak a_T(F)=\mathfrak a_T(F_+)+\mathfrak a_T(F_-)$ and a
negative sum has a negative term. (iii) By \eqref{eq:Aform} and $|\cF|\le\|F\|^2$,
$0>\mathfrak a_T(F)\ge C_0\|F\|^2+\Lev(\cF)-\Aconst\|F\|^2$. Moreover
$\Lev(\cF)=\frac1{2\pi}\int(m-C_0)|\Fh|^2$ by Theorem~\ref{thm:A} and \eqref{eq:m}, and $m-C_0\ge0$ increases in
$|\tau|$ (Lemma~\ref{lem:mbounds}(i)), so $\Lev(\cF)\ge(m(R)-C_0)\frac1{2\pi}\int_{|\tau|>R}|\Fh|^2$. (iv) is
$\mathfrak a_T=\mT+\PiT$. (v) is (iii) read as a lower bound.
\end{proof}

\begin{theorem}[structure of the set of windows with witnesses]\label{thm:wit-Tc}
Let $T_c:=\inf\mathscr T\in(0,\infty]$.
\emph{(i)} Extension by zero $\iota:L^2(\IT)\to L^2(I_{T'})$, $T<T'$, maps $\mathscr W_T$ into $\mathscr W_{T'}$ with
$\mathfrak a_{T'}(\iota F)=\mathfrak a_T(F)$.
\emph{(ii)} $\mathscr T=(T_c,\infty)$, and $T_c>\frac12\log5$.
\emph{(iii)} If $T_c<\infty$, then $\lambda_{\mathrm A}(T_c)=0$, every minimizer $h$ at $T_c$ satisfies
$\mathfrak a_{T_c}(h,G)=0$ for all $G\in\mathcal P_{T_c}\cap\mathcal D_{T_c}$, and for $T>T_c$ no witness at $T$ is supported in
$[-T_c,T_c]$.
\emph{(iv)} The limit $L:=\lim_{T\to\infty}\lamA$ exists in $[-\infty,0]$, and the following are equivalent:
\eqref{eq:target}; $T_c=\infty$; $L=0$; $\lambda_{\mathrm A}(T_k)\ge0$ for some sequence $T_k\to\infty$;
$\lamA\ge-\varepsilon(T)$ for all $T$, for some function $\varepsilon(T)\to0$. In particular \eqref{eq:target} holds as
soon as $\lambda_{\mathrm A}(\frac12\log n_k)\ge0$ along one sequence $n_k\to\infty$.
\emph{(v)} Either $0\le\lamA\le4\cdot10^9Te^{-15.449T}$ for every $T\ge1$, or there are $c>0$ and $T_0$ with
$\lamA\le-c$ for every $T\ge T_0$.
\end{theorem}

\begin{proof}
(i) is the computation in the proof of Lemma~\ref{lem:monoT}. (ii) By (i), $\mathscr T$ is an up-set; it is open
because $\lamA$ is continuous (Lemma~\ref{lem:contA}) and $\mathscr T=\{\lamA<0\}$ by
Lemma~\ref{lem:wit-basic}(i). $\lambda_{\mathrm A}(\frac12\log5)>0$ (Theorem~\ref{thm:log5A}), so by continuity
$T_c>\frac12\log5$. (iii) $\lamA\ge0$ for $T<T_c$ and $<0$ for $T>T_c$, so $\lambda_{\mathrm A}(T_c)=0$ by continuity.
$\mathfrak a_{T_c}\ge0$ on $\mathcal P_{T_c}\cap\mathcal D_{T_c}$ and $\mathfrak a_{T_c}(h)=0$, so $h$ lies in the radical by the
Cauchy--Schwarz inequality for nonnegative forms. If $F\in\mathscr W_T$ were supported in $[-T_c,T_c]$, its restriction
would be a witness at $T_c$, by (i) read backwards. (iv) Monotone limit (Lemma~\ref{lem:monoT}), $L\le0$ by
Corollary~\ref{cor:margin2}. If $\lamA\ge-\varepsilon(T)$ with $\varepsilon\to0$, then for every $T$,
$\lamA\ge\lim_{T'\to\infty}\lambda_{\mathrm A}(T')\ge0$; the other implications are monotonicity. (v) If $L=0$, use
monotonicity and Corollary~\ref{cor:margin1}; if $L<0$, take $T_0>T_c$ and $c=-\lambda_{\mathrm A}(T_0)$.
\end{proof}

\begin{lemma}[freezing]\label{lem:wit-freeze}
Let $F_0\in\mathscr W_{T_0}$ and $T\ge T_0$. Then $\iota F_0\in\mathscr W_T$, and the following do not depend on $T$:
$\mathfrak a_T(\iota F_0)=\mathfrak a_{T_0}(F_0)$, the autocorrelation $c_{\iota F_0}=c_{F_0}$, the spectral measure
$|\widehat{\iota F_0}|^2d\tau$, the parity, $\mT(\iota F_0)$, and the comb term
$\sum_{n\in\mathcal N_T}v_n\Rel c_{\iota F_0}(\log n)=\sum_{n\in\mathcal N_{T_0}}v_n\Rel c_{F_0}(\log n)$; every contact $n$ with
$\log n\ge2T_0$ satisfies $c_{\iota F_0}(\log n)=0$.
\end{lemma}

\begin{proof}
Extension by zero preserves $F$, hence $\Fh$, $\cF$, $A$, $B$ and $J$-parity, and $\cF=0$ on $|u|\ge2T_0$; apply
\eqref{eq:Aform}.
\end{proof}

\begin{theorem}[asymptotic conditions are fixed-window conditions]\label{thm:wit-asym}
Let $\mathscr C_T\subset\PT\cap\mathcal D_T$, $T>0$, satisfy $\mathscr W_T\subset\mathscr C_T$ for every $T$ (a family of
necessary conditions), and put $\mathscr C^\infty_{T_0}:=\{F\in\mathcal P_{T_0}\cap\mathcal D_{T_0}:\ \iota F\in\mathscr C_T\ \text{for all}\
T\ge T_0\}$. Then $\mathscr W_{T_0}\subset\mathscr C^\infty_{T_0}$ for every $T_0$. Consequently, if $\mathscr C_T=\emptyset$ for all
$T\ge T_1$, then $\mathscr C^\infty_{T_0}=\emptyset$, hence $\mathscr W_{T_0}=\emptyset$, for every $T_0>0$: a proof that a
family of necessary conditions becomes contradictory for large $T$ is, at each fixed $T_0$, a proof that the
fixed-window conditions $\mathscr C^\infty_{T_0}$ are contradictory.
\end{theorem}

\begin{proof}
Lemma~\ref{lem:wit-freeze} gives $\iota\mathscr W_{T_0}\subset\mathscr W_T\subset\mathscr C_T$ for $T\ge T_0$. If
$\mathscr C_T=\emptyset$ for $T\ge T_1$, no $F$ has $\iota F\in\mathscr C_{\max(T_0,T_1)}$, so $\mathscr C^\infty_{T_0}=\emptyset$.
\end{proof}

By Lemma~\ref{lem:wit-freeze}, a frozen witness does not interact with the contacts that enter after $T_0$ and
none of its components grows with $T$. Hence an argument of the form ``as $T\to\infty$ the new primes cannot sustain the
coupling a witness needs'', or any limit law for $\sT$ and its eigenfunctions, says nothing about frozen witnesses
beyond what the corresponding statement at the window $T_0$ says.

\subsection{Quartet forms and model forms}
For $0<\delta<\frac12$, $\gamma>0$ and real $F\in L^2(\IT)$ put
\[
\mathfrak q_{\delta,\gamma}(F):=\int_{-2T}^{2T}\cF(u)\,4\cos(\gamma u)\cosh(\delta u)\,du,\]\[
\Phi_{\mathrm c}(F):=\int F(t)\cosh(\delta t)e^{i\gamma t}dt,\quad \Phi_{\mathrm s}(F):=\int F(t)\sinh(\delta t)e^{i\gamma t}dt,
\]
extended to complex $F$ by $\mathfrak q(F)=\mathfrak q(\Rel F)+\mathfrak q(\Iml F)$ (its kernel $4\cos(\gamma(t-s))\cosh(\delta(t-s))$ is real
and symmetric), and $S_\delta(T):=\int_{-T}^T\sinh^2(\delta t)dt=\frac{\sinh2\delta T}{2\delta}-T$.

\begin{lemma}[the node lemma]\label{lem:wit-node}
For real $F\in L^2(\IT)$:
\emph{(i)} $\mathfrak q_{\delta,\gamma}(F)=\sum_{\rho}\Phi_F(\rho)\Phi_F(1-\rho)$, the sum over the four points
$\rho\in\{\frac12\pm\delta\pm i\gamma\}$, and $\mathfrak q_{\delta,\gamma}(F)=4|\Phi_{\mathrm c}(F)|^2-4|\Phi_{\mathrm s}(F)|^2$.
\emph{(ii)} $-4S_\delta(T)\|F\|^2\le\mathfrak q_{\delta,\gamma}(F)\le4\big(S_\delta(T)+2T\big)\|F\|^2$.
\emph{(iii)} $\Phi_{\mathrm s}(F)=-i\delta\Fh'(\gamma)+E_{\mathrm s}$ and $\Phi_{\mathrm c}(F)=\Fh(\gamma)+E_{\mathrm c}$, with
$|E_{\mathrm s}|\le\frac{\delta^3}6\cosh(\delta T)\big(\frac{2T^7}7\big)^{1/2}\|F\|$ and
$|E_{\mathrm c}|\le\frac{\delta^2}2\cosh(\delta T)\big(\frac{2T^5}5\big)^{1/2}\|F\|$.
In particular, if $\Fh(\gamma)\ne0$ and $\Fh'(\gamma)=0$, then $\mathfrak q_{\delta,\gamma}(F)=4|\Fh(\gamma)|^2+O(\delta^2)>0$ for small
$\delta$: such a test sees the quartet as a double point $\frac12\pm i\gamma$ of the zero side. A negative value requires
$|\Phi_{\mathrm s}|>|\Phi_{\mathrm c}|$, i.e., to first order in $\delta$, a node of $\Fh$ near $\gamma$ with nonzero slope.
\end{lemma}

\begin{proof}
(i) With $\Psi(z)=\int F(t)e^{zt}dt=\Phi_F(\frac12+z)$, $\int\cF(u)e^{zu}du=\Psi(z)\Psi(-z)$ (Fubini), and
$\sum_{z\in\{\pm\delta\pm i\gamma\}}e^{zu}=4\cos(\gamma u)\cosh(\delta u)$; this is the first formula. For $z=\delta+i\gamma$ the four
terms pair to $4\Rel[\Psi(z)\Psi(-z)]$ because $\Psi(\bar z)=\overline{\Psi(z)}$; and $\Psi(z)=\Phi_{\mathrm c}+\Phi_{\mathrm s}$,
$\Psi(-z)=\overline{\Phi_{\mathrm c}-\Phi_{\mathrm s}}$, so
$\Psi(z)\Psi(-z)=|\Phi_{\mathrm c}|^2-|\Phi_{\mathrm s}|^2+(\Phi_{\mathrm s}\overline{\Phi_{\mathrm c}}-\Phi_{\mathrm c}\overline{\Phi_{\mathrm s}})$,
whose last term is purely imaginary. (ii) Cauchy--Schwarz on $\IT$: $|\Phi_{\mathrm s}|^2\le\|F\|^2S_\delta(T)$ and
$|\Phi_{\mathrm c}|^2\le\|F\|^2\int_{-T}^T\cosh^2(\delta t)dt=\|F\|^2(S_\delta(T)+2T)$. (iii) $\int F(t)te^{i\gamma t}dt=-i\Fh'(\gamma)$;
$|\sinh x-x|\le\frac{|x|^3}6\cosh x$ and $|\cosh x-1|\le\frac{x^2}2\cosh x$; Cauchy--Schwarz.
\end{proof}

\begin{definition}\label{def:wit-model}
A \emph{quartet configuration} is a finite list $\mathbf q=((\delta_j,\gamma_j,m_j))_{j\le J}$ with $0<\delta_j<\frac12$,
$\gamma_j>0$, $m_j\in\N$. The \emph{model form} is $\mathfrak a^{\mathbf q}_T:=\mathfrak a_T+\sum_jm_j\,\mathfrak q_{\delta_j,\gamma_j}$ on
$\mathcal D_T$, with bottom $\lambda^{\mathbf q}_{\mathrm A}(T)$ on the unit sphere of $\PT\cap\mathcal D_T$ and witness set
$\mathscr W^{\mathbf q}_T$.
\end{definition}

By Lemma~\ref{lem:roughEF} and Lemma~\ref{lem:wit-node}(i), for real $F\in\PT$ satisfying (a), (b) of
Lemma~\ref{lem:roughEF}, $\mathfrak a^{\mathbf q}_T(F)=\sum_\rho\Phi_F(\rho)\Phi_F(1-\rho)$, the sum over the zeros of the entire
function $\xi(s)\prod_j\big[((s-\frac12)^2-z_j^2)((s-\frac12)^2-\bar z_j^2)\big]^{m_j}$, $z_j=\delta_j+i\gamma_j$, which satisfies the
same functional equation $s\mapsto1-s$ as the completed zeta function
$\xi(s)=\frac12s(s-1)\pi^{-s/2}\Gamma(\frac s2)\zeta(s)$ and has, besides the zeros of $\xi$, the quartets
$\frac12\pm\delta_j\pm i\gamma_j$ with multiplicity $m_j$. Nothing below uses this reading; the model forms are defined by
Definition~\ref{def:wit-model}.

\begin{theorem}[model forms]\label{thm:wit-model}
Let $\mathbf q$ be a quartet configuration.
\emph{(i) Shared structure.} $\mathfrak a^{\mathbf q}_T-\mathfrak a_T$ is a bounded form on $L^2(\IT)$ with a real, even, translation
kernel; so $\mathfrak a^{\mathbf q}_T$ is closed on $\mathcal D_T$, bounded below, $J$-invariant, its bottom on $\PT$ is attained, and
Lemma~\ref{lem:wit-basic}(i),(ii), Theorem~\ref{thm:wit-Tc}(i), Lemma~\ref{lem:wit-freeze} and Theorem~\ref{thm:wit-asym}
hold verbatim for $\mathfrak a^{\mathbf q}_T$, $\lambda^{\mathbf q}_{\mathrm A}$, $\mathscr W^{\mathbf q}_T$.
\emph{(ii) Weyl bounds.} $\lamA-4\sum_jm_jS_{\delta_j}(T)\le\lambda^{\mathbf q}_{\mathrm A}(T)\le\lamA+4\sum_jm_j\big(S_{\delta_j}(T)+2T\big)$.
\emph{(iii) Witnesses at every window.} For every $T>0$, $\delta\in(0,\frac12)$ and $\gamma>0$ there is $m_0<\infty$ such that the
configuration $((\delta,\gamma,m))$ has a witness at $T$, hence at every $T'\ge T$, for all $m\ge m_0$.
\emph{(iv) A single quartet.} If $\mathbf q=((\delta,\gamma,1))$ and $\mathscr W^{\mathbf q}_T\ne\emptyset$, then
$4T\big(g(\delta T)-1\big)>\lamA$, where $g(x)=\frac{\sinh2x}{2x}$. In particular the single quartet has no witness at
$T=\frac12\log3$ when $\delta\le0.013$.
\emph{(v) The degenerate quartet.} The form $\mathfrak a_T+\int\cF(u)4\cos(\gamma u)\,du=\mathfrak a_T+4|\Fh(\gamma)|^2$ (the
limit $\delta\to0$, a double point $\frac12\pm i\gamma$) has bottom $\ge\lamA$ on $\PT$.
\end{theorem}

\begin{proof}
(i) By Lemma~\ref{lem:wit-node}(ii) and the splitting into real and imaginary parts, $|\mathfrak q_{\delta,\gamma}(F)|\le
4(S_\delta(T)+2T)\|F\|^2$; a bounded perturbation preserves closedness, the domain and lower semiboundedness, and the
compactness argument of Lemma~\ref{lem:compact} applies with $\Aconst$ replaced by $\Aconst+4\sum_jm_j(S_{\delta_j}(T)+2T)$.
The kernel is even, so $J$-invariance holds. Under extension by zero $c_{\iota F}=\cF$ and the integral defining
$\mathfrak q$ runs over the support of $\cF$, so $\mathfrak q(\iota F)$ does not depend on the window; the proofs cited use
only these properties. (ii) Lemma~\ref{lem:wit-node}(ii) and the variational characterization of the bottom.
(iii) The six real functions $e^{\pm t/2}$, $\cosh(\delta t)\cos\gamma t$, $\cosh(\delta t)\sin\gamma t$, $\sinh(\delta t)\cos\gamma t$,
$\sinh(\delta t)\sin\gamma t$ are linear combinations of the six distinct exponentials $e^{\pm t/2}$, $e^{(\pm\delta\pm i\gamma)t}$ by an
invertible matrix, hence linearly independent on $\IT$. So the real-linear map
$F\mapsto(A,B,\Rel\Phi_{\mathrm c},\Iml\Phi_{\mathrm c},\Rel\Phi_{\mathrm s},\Iml\Phi_{\mathrm s})$ from real $C_c^\infty(\IT)$ to $\R^6$ is onto:
otherwise a nontrivial real combination of the six functions would be orthogonal to $C_c^\infty(\IT)$. Take a real
$F\in C_c^\infty(\IT)$ with $A=B=0$, $\Phi_{\mathrm c}(F)=0$ and $\Phi_{\mathrm s}(F)=1$. Then $F\in\PT\cap\mathcal D_T$ and, by
Lemma~\ref{lem:wit-node}(i), $\mathfrak a_T(F)+m\,\mathfrak q_{\delta,\gamma}(F)=\mathfrak a_T(F)-4m<0$ for $m>m_0:=\max(\mathfrak a_T(F),0)/4$.
Then use (i) and Theorem~\ref{thm:wit-Tc}(i). (iv) By (ii), $0>\lambda^{\mathbf q}_{\mathrm A}(T)\ge\lamA-4S_\delta(T)$ and
$4S_\delta(T)=4T(g(\delta T)-1)$. At $T=\frac12\log3$, $\lamA\ge8.8577\cdot10^{-5}$ (Theorem~\ref{thm:C}), while for
$\delta\le0.013$, with $y=2\delta T\le0.01429$, $4T(g(\delta T)-1)=4T(\frac{\sinh y}y-1)\le4T\frac{y^2}6\cosh y<7.5\cdot10^{-5}$.
(v) $\mathfrak q$ with $\delta=0$ is $4|\Phi_{\mathrm c}|^2=4|\Fh(\gamma)|^2\ge0$ by the computation in Lemma~\ref{lem:wit-node}(i).
\end{proof}

\begin{corollary}[consistency of the structural conditions]\label{cor:wit-consistent}
The conditions of Lemma~\ref{lem:wit-basic}(i),(ii) (with $\mathfrak a^{\mathbf q}_T$ in place of $\mathfrak a_T$),
Theorem~\ref{thm:wit-Tc}(i) and Lemma~\ref{lem:wit-freeze} are satisfied by the witnesses of every model form, and for
every $T>0$ there are model forms with witnesses at $T$. Hence an argument whose only inputs about $\mathfrak a_T$ are
properties it shares with all model forms (closedness, the domain, lower semiboundedness, $J$-invariance, invariance
under extension by zero) cannot show $\mathscr W_T=\emptyset$; it would show
the same for a model form with witnesses. An exclusion of witnesses for $\mathfrak a_T$ must use a property of its
arithmetic part, for instance that the measure $\sum_nv_n\delta_{\log n}$ is atomic and positive, whereas the model forms
add the density $-8\cos(\gamma u)\cosh(\delta u)\,du$ on $(0,2T)$.
\end{corollary}

\begin{proof}
Theorem~\ref{thm:wit-model}(i),(iii).
\end{proof}

\begin{remark}\label{rem:wit-model-num}
Numerically (Galerkin, not certified), the single quartet at $\gamma=\gamma_1=14.1347\ldots$ has
$\lambda^{\mathbf q}_{\mathrm A}(\log2)\approx-1.6\cdot10^{-6}$ for $\delta=0.01$ and $\approx-1.4\cdot10^{-3}$ for $\delta=0.1$, with a
first window carrying a witness near $T\approx0.57$ for $\delta=0.1$; in these computations the minimizer of the model form
has a node near $\gamma_1$, as Lemma~\ref{lem:wit-node}(iii) predicts. With the values of Appendix~\ref{app:num}, the
threshold of Theorem~\ref{thm:wit-model}(iv), $4T(g(\delta T)-1)=\lamA$, is $\delta\approx0.05$, $3\cdot10^{-4}$,
$1.5\cdot10^{-6}$ at $T=\frac12\log3,\log2,\frac12\log5$.
\end{remark}

\subsection{The spectral-derivative family and the inequality \texorpdfstring{$(E_\gamma)$}{(E gamma)}}
Fix a real even $\eta\in C_c^\infty(-1,1)$, $\eta\ne0$. For $\gamma>0$ and $T>0$ put
\[
g_T(t):=t\,\eta(t/T)\sin(\gamma t),\qquad F_T:=\big(\tfrac14-\partial^2\big)g_T,
\]
a real even function in $C_c^\infty(\IT)$ with $\Fh_T(\tau)=(\tau^2+\frac14)\widehat g_T(\tau)$; two integrations by parts give
$A(F_T)=B(F_T)=0$, so $F_T\in\PT$. $\Fh_T$ has a node at $\pm\gamma$ with slope of order $T^2$: it is, up to the factor
$\tau^2+\frac14$, the derivative in $\gamma$ of the transform of a window. Put
\[
\mathfrak m_T(\gamma):=\frac{\mT(F_T)}{\|F_T\|^2},\qquad \mathfrak S_T(\gamma):=\frac1{\|F_T\|^2}\sum_{n\in\NT}v_n\,c_{F_T}(\log n),\qquad
\text{so}\quad\frac{\mathfrak a_T(F_T)}{\|F_T\|^2}=\mathfrak m_T(\gamma)-\mathfrak S_T(\gamma).
\]

\begin{lemma}[concentration]\label{lem:wit-conc}
Let $p$ be continuous and even on $\R$ with $|p(\tau)|\le C(1+|\tau|)^M$. Then
$\frac1{2\pi\|F_T\|^2}\int p\,|\Fh_T|^2d\tau\to p(\gamma)$ as $T\to\infty$. In particular $\mathfrak m_T(\gamma)\to m(\gamma)$ and
$\|F_T'\|/\|F_T\|\to\gamma$.
\end{lemma}

\begin{proof}
With $h_1(s)=s\eta(s)$ and $h_T(t)=t\eta(t/T)$ one has $\widehat h_T(\xi)=T^2\widehat h_1(T\xi)$ and
$\widehat g_T(\tau)=\frac1{2i}\big[\widehat h_T(\tau+\gamma)-\widehat h_T(\tau-\gamma)\big]$, with $\widehat h_1$ a Schwartz function, not
identically $0$. Put $P(\tau)=p(\tau)(\tau^2+\frac14)^2$. The two diagonal terms of $\int P|\widehat g_T|^2$ are
$\frac{T^3}4\int P(\pm\gamma+x/T)|\widehat h_1(x)|^2dx=\frac{T^3}4P(\gamma)\|\widehat h_1\|_2^2(1+o(1))$ by dominated convergence
($|P(\pm\gamma+x/T)|\le C'(1+\gamma+|x|)^{M+4}$ for $T\ge1$; $P$ is even). In the cross term,
$\max(|\tau+\gamma|,|\tau-\gamma|)\ge\gamma$ gives $|\widehat h_1(T(\tau+\gamma))\widehat h_1(T(\tau-\gamma))|\le C_N(1+T\gamma)^{-N}
\big[(1+T|\tau+\gamma|)^{-N}+(1+T|\tau-\gamma|)^{-N}\big]$, so the cross term is $O_N(T^4(1+T\gamma)^{-N})=o(T^3)$. Divide by the
same expression with $p\equiv1$. For the last two statements take $p=m$ (Lemma~\ref{lem:mbounds}(ii)) and $p=\tau^2$.
\end{proof}

\begin{proposition}[the inequality $(E_\gamma)$]\label{prop:wit-Egamma}
If \eqref{eq:target} holds, then for every $\gamma>0$
\[
(E_\gamma)\qquad\limsup_{T\to\infty}\mathfrak S_T(\gamma)\le m(\gamma).
\]
\end{proposition}

\begin{proof}
$\mathfrak a_T(F_T)\ge0$ gives $\mathfrak S_T(\gamma)\le\mathfrak m_T(\gamma)$; apply Lemma~\ref{lem:wit-conc}.
\end{proof}

$(E_\gamma)$ is a statement about the primes alone: the comb, tested against the autocorrelation of the
spectral-derivative direction at $\gamma$, does not exceed the archimedean level $m(\gamma)$ in the limit. It is the
restriction of \eqref{eq:target} to a one-parameter family; by Lemma~\ref{lem:wit-freeze} a witness of fixed support,
which is not of the form $F_T$, is not seen by it. For a model form with a quartet $(\delta,\gamma)$ the left side of
$(E_\gamma)$ acquires the term $-\mathfrak q_{\delta,\gamma}(F_T)/\|F_T\|^2$; numerically (without the comb, $\delta=0.1$,
$\gamma=\gamma_1$) it is about $1.3$, $205$, $1.6\cdot10^3$ at $T=5,20,30$, consistent with growth of order $e^{2\delta T}/T$.

\subsection{Coherence of the comb on the mesh}
Let $\varrho_T:=\sum_{n\in\NT}v_n\delta_{\log n}-2e^{u/2}\mathbf 1_{(0,2T)}(u)\,du$, a signed measure on $(0,2T)$, and define the
\emph{coherence}
\[
\mathcal C_T(\omega):=-\frac1{4T}\int\cos(\omega u)\,d\varrho_T(u)
=\frac1{4T}\Big[\Rel\frac{2\big(e^{(\frac12+i\omega)2T}-1\big)}{\frac12+i\omega}-\sum_{n\in\NT}v_n\cos(\omega\log n)\Big].
\]
The normalization is such that adding to $-\varrho_T$ the density $4\cos(\gamma u)$ on $(0,2T)$, which is half the
degenerate quartet of Theorem~\ref{thm:wit-model}(v), adds $\frac1T\int_0^{2T}\cos^2(\gamma u)du=1+O(1/(\gamma T))$ to
$\mathcal C_T(\gamma)$. Write $\mathfrak b_T(F):=C_0\|F\|^2+\Lev(\cF)-2\int_0^{2T}\Rel\cF(u)e^{u/2}du$, so that
$\mathfrak a_T(F)=\mathfrak b_T(F)-\int\Rel\cF\,d\varrho_T$ by \eqref{eq:Aform}.

\begin{proposition}[coherence on the mesh]\label{prop:wit-coh}
\emph{(i)} For $F\in L^2(\IT)$,
\[
\int\Rel\cF\,d\varrho_T=-\sum_k|\Fh(\omega_k)|^2\,\mathcal C_T(\omega_k),
\]
the series converging absolutely. Hence every $F\in\mathscr W_T$ satisfies $\sum_k|\Fh(\omega_k)|^2\,\mathcal C_T(\omega_k)<-\mathfrak b_T(F)$.
\emph{(ii) Comparison of a quartet with its degenerate limit.} Let $\mathcal C^{(\delta)}_T$ be defined as $\mathcal C_T$ with
$\varrho_T$ replaced by $\varrho^{(\delta)}_T:=\varrho_T-8\cos(\gamma u)\cosh(\delta u)\mathbf 1_{(0,2T)}du$, so that
$\mathfrak a_T+\mathfrak q_{\delta,\gamma}=\mathfrak b_T-\int\Rel\cF\,d\varrho^{(\delta)}_T$, and $\mathcal C^{(0)}_T$ likewise with $\delta=0$. Then
\[
\sup_\omega\big|\mathcal C^{(\delta)}_T(\omega)-\mathcal C^{(0)}_T(\omega)\big|\le4\big(g(\delta T)-1\big)\le\tfrac23(2\delta T)^2\cosh(2\delta T),
\]
while the form with $\delta=0$ has bottom $\ge\lamA$ on $\PT$ (Theorem~\ref{thm:wit-model}(v)), which is $>0$ for
$T\le\frac12\log5$.
\end{proposition}

\begin{proof}
(i) Apply Lemma~\ref{lem:nyq-sampling}(iv) to the finite even measure $\frac12(\varrho_T+\check\varrho_T)$ on
$[-2T,2T]$, $\check\varrho_T$ the reflection; its transform is $\int\cos(\omega u)d\varrho_T=-4T\,\mathcal C_T(\omega)$, and
$|\mathcal C_T|$ is bounded. The second statement is $\mathfrak a_T(F)<0$. (ii) The difference of the two measures is
$8\cos(\gamma u)(\cosh\delta u-1)du$ on $(0,2T)$, so the difference of coherences is at most
$\frac1{4T}\int_0^{2T}8(\cosh\delta u-1)du=\frac2T\big(\frac{\sinh2\delta T}\delta-2T\big)=4(g(\delta T)-1)$ in modulus; and
$\frac{\sinh y}y-1=\sum_{k\ge1}\frac{y^{2k}}{(2k+1)!}\le\frac{y^2}6\cosh y$. Positivity for $T\le\frac12\log5$ is
Theorem~\ref{thm:log5A}.
\end{proof}

So no condition on the coherence that is Lipschitz in the supremum norm with constant $L$ can separate the quartet
model from its degenerate limit by more than $\frac23L(2\delta T)^2\cosh(2\delta T)$; numerically, at $T=\log2$ the single
quartet with $\delta=0.01$ has a witness (Remark~\ref{rem:wit-model-num}) while $4(g(\delta T)-1)\approx1.3\cdot10^{-4}$.

The comb term of $(E_\gamma)$ is a smoothed, twisted Chebyshev sum. With $\psi(y)=\sum_{n\le y}\Lambda(n)$:

\begin{proposition}[Chebyshev form]\label{prop:wit-cheb}
Let $F\in\PT\cap H^1_0(\IT)$ be real with $\|F\|=1$, and $w:=\cF$. Then
\[
\sum_{n\in\NT}v_nw(\log n)=2\int_0^{2T}w(u)e^{u/2}du+\mathcal F_T(F),\qquad
\mathcal F_T(F):=2\int_{[1,e^{2T})}y^{-1/2}w(\log y)\,d\big(\psi(y)-y\big),
\]
with $\big|2\int_0^{2T}w(u)e^{u/2}du\big|\le4$ and
$|\mathcal F_T(F)|\le2+(2\|F'\|+1)\int_1^{e^{2T}}|\psi(y)-y|\,y^{-3/2}dy$.
\end{proposition}

\begin{proof}
$\sum_{n<e^{2T}}2\Lambda(n)n^{-1/2}w(\log n)=2\int_{[1,e^{2T})}y^{-1/2}w(\log y)d\psi(y)$, and the part $d y$ gives
$2\int_0^{2T}w(u)e^{u/2}du$ ($y=e^u$). Since $F\in\PT$, $\int_{-2T}^{2T}\cF\cosh\frac u2\,du=\Rel(A\bar B)=0$
(Proposition~\ref{prop:P}), and $\cF$ is even, so $\int_0^{2T}we^{u/2}=-\int_0^{2T}we^{-u/2}$, of modulus at most
$\int_0^\infty e^{-u/2}=2$. For $\mathcal F_T$ integrate by parts: $w(2T)=0$, $w(0)=1$, $\psi(1)-1=-1$ give the boundary
term $2$; and $\big|\frac{d}{dy}\big(y^{-1/2}w(\log y)\big)\big|\le y^{-3/2}(|w'|+\frac12)$ with $|w'(u)|=|\int F(t)F'(t-u)dt|\le\|F'\|$.
\end{proof}

\begin{remark}\label{rem:wit-cheb}
For $F=F_T/\|F_T\|$, Lemma~\ref{lem:wit-conc} gives $\|F'\|\to\gamma$, so $(E_\gamma)$ holds if
$\limsup_T\mathcal F_T(F_T/\|F_T\|)\le m(\gamma)-4$, and implies $\limsup_T\mathcal F_T(F_T/\|F_T\|)\le m(\gamma)+4$. The prime number
theorem with the Vinogradov--Korobov error term \cite[\S8.5]{IK} bounds the integral in Proposition~\ref{prop:wit-cheb} by
$e^{T-c(2T)^{3/5}(\log2T)^{-1/5}}$ up to a constant, and a bound $|\psi(y)-y|\le Cy^\theta$ with $\theta>\frac12$ would give
$O(e^{(2\theta-1)T})$; so the absolute-value bound does not give $(E_\gamma)$ from any such estimate, and $(E_\gamma)$, if
proved from the primes, needs cancellation in the twisted sum $\mathcal F_T$ at the fixed frequency $\gamma$.
\end{remark}

\begin{proposition}[sensitivity to the weights]\label{prop:wit-sens}
For weight vectors $v,v'\in\R^{\NT}$ let $\lambda_{\mathrm A}(T;v)$ be the bottom of $\mathfrak a^{(v)}_T$ on $\PT$
(\S\ref{subsec:finite}). Then $|\lambda_{\mathrm A}(T;v)-\lambda_{\mathrm A}(T;v')|\le\sum_n|v_n-v'_n|$. Let $h$ be a real unit
minimizer at $T$ for the weights $v_n=2\Lambda(n)/\sqrt n$ and $\kappa_T:=\sum_{n\in\NT}v_n\,c_h(\log n)$. Then for
$\varepsilon>0$, $\lambda_{\mathrm A}(T;(1+\varepsilon)v)\le\lamA-\varepsilon\kappa_T$, while
$\sup_\omega|\mathcal C_T^{(1+\varepsilon)v}(\omega)-\mathcal C_T(\omega)|\le\varepsilon\Aconst/(4T)$. Consequently, if $\kappa_T>0$, the
weights $(1+\varepsilon)v$ with $\varepsilon>\lamA/\kappa_T$ violate positivity at $T$; for $T\ge1$ this holds for all
$\varepsilon>4\cdot10^9Te^{-15.449T}/\kappa_T$ (Corollary~\ref{cor:margin1}).
\end{proposition}

\begin{proof}
The two forms differ by $\sum_n(v'_n-v_n)\Rel\cF(\log n)$, bounded by $\sum|v_n-v'_n|\,\|F\|^2$. Rayleigh with $h$:
$\mathfrak a^{((1+\varepsilon)v)}_T(h)=\lamA-\varepsilon\kappa_T$. The coherence changes by $\frac\varepsilon{4T}\sum v_n\cos(\omega\log n)$.
\end{proof}

\begin{remark}\label{rem:wit-sens}
Numerically $\kappa_{\log2}\approx0.09$, so at $T=\log2$ a uniform relative change of the weights of about $10^{-6}$
destroys positivity while moving the coherence by less than $10^{-5}$ in the supremum norm (Galerkin, not certified).
Thus a sufficient condition for \eqref{eq:target} at a window that is stated through a functional of the weights,
continuous in the weights, must be sharp at the scale of $\lamA/\kappa_T$; compare Corollary~\ref{cor:void}, which gives
the corresponding statement, proved and without numerics, for modifications preserving classical prime-side
information.
\end{remark}


\section{Combs, translates of one profile, and the polar-free form}\label{sec:combs}

This section treats the prime comb as one member of a family of positive measures and asks which members
are positive at every window. It then reduces \eqref{eq:target} to statements about translates of a single
test function. Subsection~\ref{subsec:polarfree} removes the polar term exactly. Subsection~\ref{subsec:poscombs}
shows that combs positive at every window exist, continuous and atomic, and identifies the density threshold.
Subsection~\ref{subsec:reson} proves that combs supported on $\log\N$ cannot exceed density $1$, and that $\Lambda$
is positive on every test space on which near-coincidences of integers are invisible.
Subsections~\ref{subsec:translates} and \ref{subsec:oneprofile} prove a two-sided band for the local Chebyshev
fluctuation, a large-deviation bound for random combs, and the equivalence of \eqref{eq:target} with the positive
definiteness of one explicit function of one real variable. None of these results gives the sign of $\lamA$ at a
window $T>\frac12\log5$.

\subsection{The polar-free form}\label{subsec:polarfree}
For $F\in L^2(\IT)$ put
\[
\mathfrak e(F):=\frac1{2\pi}\int_\R\frac{|\Fh(\tau)|^2}{\frac14+\tau^2}\,d\tau ,\qquad
\mathcal F_T(F):=\sum_{n\in\NT}v_n\Rel\cF(\log n)-2\int_0^{2T}\Rel\cF(u)\,e^{u/2}\,du .
\]
For real $F\in\PT\cap H^1_0(\IT)$ with $\|F\|=1$ this $\mathcal F_T$ is the Chebyshev remainder of
Proposition~\ref{prop:wit-cheb}.

\begin{lemma}[polar elimination]\label{lem:polarelim}
For every $F\in\PT$,
\[
2\int_0^{2T}\Rel\cF(u)\,e^{u/2}\,du\;=\;-2\int_0^{2T}\Rel\cF(u)\,e^{-u/2}\,du\;=\;-\mathfrak e(F),
\]
and if $F=(\frac14-\partial^2)G$ with $G\in H^2_0(\IT)$ (Theorem~\ref{thm:nyq-cone}), then
$\mathfrak e(F)=\|G'\|^2+\frac14\|G\|^2$. In particular $0\le\mathfrak e(F)\le4\|F\|^2$.
\end{lemma}

\begin{proof}
Substituting $s=t-u$ gives $\int\cF(u)e^{u/2}du=A(F)\overline{B(F)}$ and
$\int\cF(u)e^{-u/2}du=B(F)\overline{A(F)}$, which vanish on $\PT$. Since $\cF(-u)=\overline{\cF(u)}$, $\Rel\cF$ is even, so taking real parts gives
$\int_0^{2T}\Rel\cF(u)(e^{u/2}+e^{-u/2})\,du=0$, the first equality. Next
$2\int_0^{2T}\Rel\cF\,e^{-u/2}du=\int_\R\cF(u)e^{-|u|/2}du$ (the imaginary part of $\cF$ is odd), and the transform of
$e^{-|u|/2}$ is $(\frac14+\tau^2)^{-1}$, so Plancherel gives $\mathfrak e(F)$. Finally
$|\Fh|^2/(\frac14+\tau^2)=(\tau^2+\frac14)|\widehat G|^2$, and $(\frac14+\tau^2)^{-1}\le4$.
\end{proof}

\begin{proposition}[the polar-free form]\label{prop:sharp}
Put $m^\sharp(\tau):=m(\tau)+\frac1{\frac14+\tau^2}$ and $\mathfrak a^\sharp(F):=\frac1{2\pi}\int m^\sharp|\Fh|^2$.
\emph{(i)} $m^\sharp(\tau)=\Rel\psi(\frac54+\frac{i\tau}2)-\log\pi=C_0+4+\sum_{k\ge1}\frac{\tau^2}{a_k(\tau^2+4a_k^2)}$ with
$a_k=k+\frac14$; it is even, increasing in $|\tau|$, and $m^\sharp(0)=C_0+4=-1.3721834\ldots$.
\emph{(ii)} For every $F\in\PT\cap\mathcal D_T$, $\mathfrak a_T(F)=\mathfrak a^\sharp(F)-\mathcal F_T(F)$. Hence $\lamA\ge0$
if and only if $\mathcal F_T(F)\le\mathfrak a^\sharp(F)$ for every $F\in\PT\cap\mathcal D_T$.
\emph{(iii)} The bottom $\lambda^\sharp(T):=\inf\{\mathfrak a^\sharp(F):F\in\PT\cap\mathcal D_T,\ \|F\|=1\}$ is nonincreasing in $T$,
$\lambda^\sharp(T)\ge C_0+4$ for every $T$, and $\lambda^\sharp(T)\to C_0+4$ as $T\to\infty$.
\end{proposition}

\begin{proof}
(i) $\psi(z+1)=\psi(z)+1/z$ with $z=\frac14+\frac{i\tau}2$, and $\Rel\frac1z=\frac1{\frac14+\tau^2}$. In Theorem~\ref{thm:S}(i) the
summand $k=0$ is $\frac{4\tau^2}{\tau^2+1/4}=4-\frac1{\frac14+\tau^2}$, which gives the series; each summand with $k\ge1$ is
nonnegative and increasing in $|\tau|$. (ii) By \eqref{eq:Aform} and Lemma~\ref{lem:polarelim},
$\sum_{n\in\NT}v_n\Rel\cF(\log n)=\mathcal F_T(F)-\mathfrak e(F)$, so
$\mathfrak a_T(F)=\mT(F)+\mathfrak e(F)-\mathcal F_T(F)$. (iii) Monotonicity: extension by zero maps $\PT$ into
$\mathcal P_{T'}$ for $T<T'$ and does not change $\mathfrak a^\sharp$. The lower bound is $m^\sharp\ge m^\sharp(0)$. For the
limit take the tests $F_T=(\frac14-\partial^2)G_T$ of the proof of Proposition~\ref{prop:finsupp}; the computation there,
with the continuous function $m^\sharp$ of logarithmic growth in place of $\sigma$, gives
$\mathfrak a^\sharp(F_T)/\|F_T\|^2\to m^\sharp(0)$.
\end{proof}

By Proposition~\ref{prop:finite}(ii), positivity at $T$ is the inequality
$\sum_nv_n\Rel\cF(\log n)\le\mT(F)$ on $\PT$. Its right side has bottom $\laminfP(T)\downarrow C_0$
(Proposition~\ref{prop:arch}). Proposition~\ref{prop:sharp} moves the main term of the comb to the archimedean side, where the
Tate conditions turn it into the positive energy $\mathfrak e(F)$. The fluctuation $\mathcal F_T$ must then be paid by
$\mathfrak a^\sharp$, whose bottom tends to $C_0+4$, not to $C_0$.

\begin{proposition}[the exact form of $(E_\gamma)$]\label{prop:Egamma-exact}
Let $F_T$ be the spectral-derivative family of \S\ref{sec:wit} at $\gamma>0$. Then $(E_\gamma)$ of
Proposition~\ref{prop:wit-Egamma} holds if and only if
\[
\limsup_{T\to\infty}\frac{\mathcal F_T(F_T)}{\|F_T\|^2}\ \le\ m^\sharp(\gamma).
\]
\end{proposition}

\begin{proof}
$F_T\in\PT$, so $\mathfrak S_T(\gamma)=(\mathcal F_T(F_T)-\mathfrak e(F_T))/\|F_T\|^2$ by Lemma~\ref{lem:polarelim}. Lemma~\ref{lem:wit-conc}
with the bounded continuous even function $p=(\frac14+\tau^2)^{-1}$ gives $\mathfrak e(F_T)/\|F_T\|^2\to(\frac14+\gamma^2)^{-1}$.
\end{proof}

This replaces the two-sided slack $\pm4$ of Remark~\ref{rem:wit-cheb} by an exact threshold. The threshold is the
boundary value of the logarithmic derivative whose smoothed partial sums $\mathcal F_T(F_T)/\|F_T\|^2$ are.

\begin{proposition}[the threshold is a boundary value]\label{prop:boundary}
Let $\gamma\in\R$ with $\zeta(\frac12+i\gamma)\ne0$ and $s=\frac12+i\gamma$. Then
\[
m(\gamma)=-2\Rel\frac{\zeta'}{\zeta}(s),\qquad m^\sharp(\gamma)=\Rel\Big[-2\frac{\zeta'}{\zeta}(s)-\frac2{s-1}\Big].
\]
In particular $C_0=-2\,\zeta'(\frac12)/\zeta(\frac12)$.
\end{proposition}

\begin{proof}
Let $\xi(s)=\frac12s(s-1)\pi^{-s/2}\Gamma(\frac s2)\zeta(s)$. From $\xi(1-s)=\xi(s)$ and $\xi(\bar s)=\overline{\xi(s)}$, the function
$t\mapsto\xi(\frac12+it)$ is real, so $\frac{d}{dt}\xi(\frac12+it)=i\xi'(\frac12+it)$ is real and $\frac{\xi'}{\xi}(\frac12+it)$ is purely
imaginary where $\xi(\frac12+it)\ne0$. Now $\frac{\xi'}{\xi}(s)=\frac1s+\frac1{s-1}-\frac12\log\pi+\frac12\psi(\frac s2)+\frac{\zeta'}{\zeta}(s)$, and on
$\Rel s=\frac12$ one has $\Rel(\frac1s+\frac1{s-1})=0$ and $\Rel\frac12\psi(\frac s2)-\frac12\log\pi=\frac12m(\gamma)$. This gives the first
identity. The second follows from $\Rel\frac{-2}{s-1}=\frac1{\frac14+\gamma^2}$. At $\gamma=0$, $\zeta(\frac12)\ne0$ and $m(0)=C_0$.
\end{proof}

\begin{remark}[criticality of packet limits]\label{rem:critical}
For the normalized autocorrelations $w_T=c_{F_T}/\|F_T\|^2$, Lemma~\ref{lem:wit-conc} with $p=\cos(\tau u)$ gives $w_T(u)\to\cos(\gamma u)$
locally uniformly. So $\mathcal F_T(F_T)/\|F_T\|^2$ is a smoothed partial sum of $\Rel\sum_n2\Lambda(n)n^{-s}$ at
$s=\frac12+i\gamma$ minus its polar part, and Proposition~\ref{prop:boundary} identifies the threshold of
Proposition~\ref{prop:Egamma-exact} with the boundary value of the corresponding function on the critical line. If these
smoothed sums converge to that boundary value, then $(E_\gamma)$ holds with equality, at every $\gamma$ and at the
centre: the concentrating packets carry no slack, and an inequality $(E_\gamma)$ weakened by a constant can only be
proved on them by proving the convergence itself. We do not prove the convergence. The proposition is the exact part
of this picture.
\end{remark}

\subsection{Positive combs}\label{subsec:poscombs}

\begin{definition}\label{def:comb}
A \emph{comb} is a positive Borel measure $\nu$ on $(0,\infty)$ that is finite on bounded sets. Its window form is
\[
\mathfrak a^\nu_T(F):=\mT(F)-\int_{(0,2T)}\Rel\cF\,d\nu,\qquad F\in\PT\cap\mathcal D_T,
\]
with bottom $b_T(\nu)$ on the unit sphere. Put $\mathfrak C^+:=\{\nu:\ b_T(\nu)\ge0\text{ for every }T>0\}$, and let
$\nu_\Lambda:=\sum_{n\ge2}v_n\delta_{\log n}$. Then $\mathfrak a^{\nu_\Lambda}_T=\mathfrak a_T$, and \eqref{eq:target} is the
statement $\nu_\Lambda\in\mathfrak C^+$. For $\kappa\ge0$ let $\rho_\kappa:=2\kappa e^{u/2}du$ on $(0,\infty)$ and
$g_\kappa(\tau):=m(\tau)+\kappa(\frac14+\tau^2)^{-1}$.
\end{definition}

Since $\sum_{n<x}2\Lambda(n)n^{-1/2}$ has the main term $\int_1^x2y^{-1/2}dy$, the main term of $\nu_\Lambda$ is $\rho_1$, and
$g_1=m^\sharp$. Let $\tau_0>0$ be the zero of $m$ on $(0,\infty)$; it exists and is unique because $m$ is increasing there,
$m(0)=C_0<0$ and $m\to\infty$ (Lemma~\ref{lem:mbounds}).

\begin{proposition}[structure of $\mathfrak C^+$]\label{prop:Cplus}
\emph{(i)} $\nu\mapsto b_T(\nu)$ is concave, so $\mathfrak C^+$ is convex.
\emph{(ii)} $\mathfrak C^++\rho_\kappa\subset\mathfrak C^+$ for every $\kappa\ge0$.
\emph{(iii)} There is a constant $c_1>0$ such that every $\nu\in\mathfrak C^+$ satisfies $\nu((0,2T))\ge c_1T^2$ for all large $T$.
\end{proposition}

\begin{proof}
(i) $b_T$ is an infimum of functions affine in $\nu$. (ii) By Lemma~\ref{lem:polarelim},
$\int_{(0,2T)}\Rel\cF\,d\rho_\kappa=-\kappa\,\mathfrak e(F)\le0$ on $\PT$. (iii) Fix a real even $g\ge0$ in $C_c^\infty(-1,1)$, $g\ne0$,
put $G_T(t)=T^{-1/2}g(t/T)$ and $F_T=(\frac14-\partial^2)G_T\in\PT$. Then
$c_{F_T}=\frac1{16}c_{G_T}-\frac12c_{G_T}''+c_{G_T}''''$ with $c_{G_T}(u)=c_g(u/T)\ge0$, so
$c_{F_T}\ge-\frac12T^{-2}\|c_g''\|_\infty-T^{-4}\|c_g''''\|_\infty$ and
$\int_{(0,2T)}c_{F_T}\,d\nu\ge-C_gT^{-2}\nu((0,2T))$ for $T\ge1$. As in Proposition~\ref{prop:finsupp},
$\mT(F_T)/\|F_T\|^2\to C_0$ and $\|F_T\|^2\to\frac1{16}\|g\|^2$. If $\nu((0,2T))\le c_1T^2$ along a sequence $T\to\infty$ with
$16C_gc_1/\|g\|^2<|C_0|$, then $\mathfrak a^\nu_T(F_T)<0$ along it.
\end{proof}

\begin{theorem}[main-term combs]\label{thm:kappa}
\emph{(i)} For every $T$ and $F\in\PT\cap\mathcal D_T$, $\mathfrak a^{\rho_\kappa}_T(F)=\frac1{2\pi}\int g_\kappa|\Fh|^2$.
\emph{(ii)} $\rho_\kappa\in\mathfrak C^+$ if and only if $g_\kappa\ge0$ on $\R$, if and only if
\[
\kappa\ \ge\ \kappa_\star:=\sup_{|\tau|<\tau_0}\big(-m(\tau)\big)\big(\tfrac14+\tau^2\big),\qquad
\tfrac14|C_0|\le\kappa_\star\le|C_0|\big(\tfrac14+\tau_0^2\big).
\]
\emph{(iii)} For $\kappa>\kappa_\star$, $\eta_\kappa:=\inf_\R g_\kappa>0$ and $b_T(\rho_\kappa)\ge\eta_\kappa$ for every $T$.
\emph{(iv)} $\rho_1\notin\mathfrak C^+$, since $1<\frac14|C_0|=1.3430\ldots$; indeed $g_1(0)=m^\sharp(0)=C_0+4<0$.
\end{theorem}

\begin{proof}
(i) is Lemma~\ref{lem:polarelim}. (ii) If $g_\kappa\ge0$, (i) gives $b_T\ge0$. If $g_\kappa(\gamma_1)<0$ for some $\gamma_1$, then by continuity and
evenness $g_\kappa(\gamma)<0$ for some $\gamma>0$; the spectral-derivative family $F_T$ at $\gamma$ lies in $\PT$, and
Lemma~\ref{lem:wit-conc} with $p=g_\kappa$ (continuous, even, $|g_\kappa|\le|m|+4\kappa$) gives
$\mathfrak a^{\rho_\kappa}_T(F_T)/\|F_T\|^2\to g_\kappa(\gamma)<0$. For $|\tau|\ge\tau_0$, $g_\kappa\ge m\ge0$; for $|\tau|<\tau_0$,
$g_\kappa(\tau)\ge0$ iff $\kappa\ge-m(\tau)(\frac14+\tau^2)$. The bounds on $\kappa_\star$ are the value at $\tau=0$ and
$-m\le|C_0|$. (iii) On $|\tau|<\tau_0$, $g_\kappa=(\kappa+m(\frac14+\tau^2))/(\frac14+\tau^2)\ge(\kappa-\kappa_\star)/(\frac14+\tau_0^2)$;
on $|\tau|\ge\tau_0$, $g_\kappa$ is continuous, at least $\kappa(\frac14+\tau^2)^{-1}>0$, and tends to $\infty$. Then
$\frac1{2\pi}\int g_\kappa|\Fh|^2\ge\eta_\kappa\|F\|^2$. (iv) is (ii) and Proposition~\ref{prop:sharp}(i).
\end{proof}

\begin{remark}\label{rem:beurling}
For $x_0>1$ the comb $\rho_\kappa|_{(\log x_0,\infty)}$ is the comb of the continuous Beurling system with
$\psi_B(x)=\kappa(x-x_0)_+$ \cite{DiamondZhang2016}. Removing $\rho_\kappa|_{(0,\log x_0)}$ changes the form by at most
$2\kappa\int_0^{\log x_0}e^{u/2}du\,\|F\|^2=4\kappa(\sqrt{x_0}-1)\|F\|^2$, so for $\kappa>\kappa_\star$ and $4\kappa(\sqrt{x_0}-1)<\eta_\kappa$
this comb lies in $\mathfrak C^+$ with $b_T\ge\eta_\kappa-4\kappa(\sqrt{x_0}-1)>0$ for every $T$. Positivity at every window is
therefore not special by itself; it is forced here by a smooth density $\kappa>\kappa_\star>1$. By
Proposition~\ref{prop:sharp}, the comb $\nu_\Lambda$ has density $1$, and its fluctuation must cover the symbol
$m^\sharp=g_1$, which is negative near $0$.
\end{remark}

Atomic combs can be positive at every window too. For $a<b$, $h>0$ with $b-a\in h\N$ and $\kappa>0$ let
\[
\tau_h[a,b]:=h\Big(\tfrac12\delta_a+\sum_{0<j<(b-a)/h}\delta_{a+jh}+\tfrac12\delta_b\Big)\cdot2\kappa e^{u/2},
\]
the trapezoid discretization of $\rho_\kappa|_{[a,b]}$, and $\tau_h[a,a]:=0$. Write $w=\frac12+i\tau$.

\begin{lemma}[exact trapezoid error]\label{lem:trap}
$\varepsilon_h(\tau;a,b):=\int e^{i\tau u}d\tau_h[a,b]-\int_a^b2\kappa e^{wu}du=2\kappa(e^{wb}-e^{wa})\big[\frac h2\coth\frac{hw}2-\frac1w\big]$. Moreover
\emph{(i)} $|\varepsilon_h|\le2\kappa(e^{b/2}+e^{a/2})(4+\frac{h^2}{24})$ for all $\tau$, and
\emph{(ii)} $|\varepsilon_h|\le\frac29\kappa(e^{b/2}+e^{a/2})h^2|w|$ for $|w|\le\pi/h$.
\end{lemma}

\begin{proof}
With $M=(b-a)/h$, $\sum_{j=0}^{M-1}e^{w(a+jh)}=(e^{wb}-e^{wa})/(e^{wh}-1)$ and $\frac1{e^x-1}+\frac12=\frac12\coth\frac x2$ give the
first term; the integral is $(e^{wb}-e^{wa})/w$. (i) For $\Rel z>0$, $|\coth z|^2=\frac{\cosh2\Rel z+\cos2\Iml z}{\cosh2\Rel z-\cos2\Iml z}\le\coth^2\Rel z$;
here $\Rel\frac{hw}2=\frac h4$ and $\coth x\le\frac1x+\frac x3$, so $|\frac h2\coth\frac{hw}2|\le2+\frac{h^2}{24}$, and $|\frac1w|\le2$.
(ii) With $z=\frac{hw}2$, $|z|\le\frac\pi2$: $\frac h2\coth\frac{hw}2-\frac1w=(z\coth z-1)/w$ and
$z\coth z-1=\sum_{k\ge1}\frac{2z^2}{z^2+\pi^2k^2}$ with $|z^2+\pi^2k^2|\ge\frac34\pi^2k^2$, so $|z\coth z-1|\le\frac49|z|^2$.
\end{proof}

\begin{theorem}[atomic combs positive at every window]\label{thm:atomic}
Let $\kappa>\kappa_\star$ and $\eta=\eta_\kappa$. Choose $h_j=2^{-n_j}$, $j\ge0$, with $h_0\le1$, $h_{j+1}\le h_j$, and
\[
\textup{(H1)}\ \ \tfrac{4\pi}9\kappa e^{j+1}h_j\le\eta\,2^{-j-4},\qquad
\textup{(H2)}\ \ \log\tfrac1{4h_j}-1\ge30\kappa e^{j+1}+\eta .
\]
Then $\nu:=\sum_{j\ge0}\tau_{h_j}[2j,2j+2]$, restricted to $(0,\infty)$, is a purely atomic comb with $b_T(\nu)\ge\frac\eta2$ for
every $T>0$. In particular $\mathfrak C^+$ contains purely atomic combs.
\end{theorem}

\begin{proof}
Let $\bar\nu$ be the same sum on $[0,\infty)$. For a positive measure $\mu$ on $[0,2T)$ and $F\in\PT\cap\mathcal D_T$, Fubini and
$\Rel\cF(u)=\frac1{2\pi}\int\cos(\tau u)|\Fh|^2d\tau$ give, by Theorem~\ref{thm:kappa}(i),
\begin{equation}\label{eq:star-comb}
\mT(F)-\int_{[0,2T)}\Rel\cF\,d\mu=\frac1{2\pi}\int\big[g_\kappa-\Rel E_\mu\big]|\Fh|^2d\tau,\qquad
E_\mu(\tau):=\int_{[0,2T)}e^{i\tau u}d\mu-\int_0^{2T}2\kappa e^{wu}du .
\end{equation}
Removing the atom of $\bar\nu$ at $0$ adds $\bar\nu(\{0\})\|F\|^2\ge0$, so it suffices to show $g_\kappa-\Rel E\ge\frac\eta2$ for
$E:=E_{\bar\nu|[0,2T)}$. Fix $T$, let $J=\lfloor T\rfloor$, and let $b'$ be the largest point of $2J+h_J\Z$ in $[2J,2T)$. On $[0,2T)$,
$\bar\nu$ is $\sum_{j<J}\tau_{h_j}[2j,2j+2]+\tau_{h_J}[2J,b']$ plus a half atom of weight $\kappa h_Je^{b'/2}$ at $b'$. Hence
$E=\sum_{j<J}\varepsilon_{h_j}(\tau;2j,2j+2)+\varepsilon_{h_J}(\tau;2J,b')+R$ with
$|R|\le\kappa h_Je^{T}+2\kappa e^T(2T-b')\le3\kappa h_Je^{J+1}\le\frac{27}{4\pi}\eta2^{-J-4}\le\frac\eta7$ by (H1).

\emph{Low band}, $|w|\le\pi/h_0$. Then $|w|\le\pi/h_j$ for all $j$, and Lemma~\ref{lem:trap}(ii) bounds each block, including the
partial one, by $\frac29\kappa\cdot2e^{j+1}h_j^2\frac\pi{h_j}=\frac{4\pi}9\kappa e^{j+1}h_j\le\eta2^{-j-4}$. So $|E|\le\frac\eta8+\frac\eta7$
and $g_\kappa-\Rel E\ge\eta-0.27\eta$.

\emph{Band $i\ge1$}, $\pi/h_{i-1}<|w|\le\pi/h_i$. The blocks $j\ge i$ satisfy $|w|\le\pi/h_i\le\pi/h_j$, so as above they contribute at
most $\sum_{j\ge i}\eta2^{-j-4}\le\frac\eta{16}$. The blocks $j<i$, including the partial block if $J<i$, contribute by
Lemma~\ref{lem:trap}(i) at most $\sum_{j<i}8.1(1+e^{-1})\kappa e^{j+1}\le8.1(1+e^{-1})\frac e{e-1}\kappa e^i\le17.6\kappa e^i$.
Since $h_{i-1}\le1<\pi$, $|\tau|\ge|w|-\frac12\ge\pi/(2h_{i-1})$, and Lemma~\ref{lem:mbounds}(ii) gives
$m(\tau)\ge\log\frac1{4h_{i-1}}-\frac{2h_{i-1}}\pi-\frac{2h_{i-1}^2}{\pi^2}\ge\log\frac1{4h_{i-1}}-1\ge30\kappa e^i+\eta$ by (H2). As
$g_\kappa\ge m$, we get $g_\kappa-\Rel E\ge30\kappa e^i+\eta-17.6\kappa e^i-\frac\eta{16}-\frac\eta7\ge\frac\eta2$.

By (H2), $h_i\to0$, so the bands cover $\R$. Then \eqref{eq:star-comb} gives the claim.
\end{proof}

\begin{remark}\label{rem:atomic}
The atoms of $\nu$ sit at $x=e^u$ with $u$ in dyadic lattices; they are not integers. The weights $2\kappa h_je^{u/2}$ are arbitrary
positive reals, not of the form $2\Lambda(n)n^{-1/2}$, and the spacing is doubly exponentially small,
$h_j\approx\exp(-30\kappa e^{j+1})$. Two questions are open: whether a Beurling generalized prime system realizes such a comb,
and whether a spacing this small is necessary for an atomic comb to be positive at every window. The comparison with
$\nu_\Lambda$ that is proved is the density one: $\nu_\Lambda$ has density $1<\kappa_\star$. Theorem~\ref{thm:ceiling} below shows that
on $\log\N$ no larger density is possible.
\end{remark}

\subsection{Resonators and the integer density ceiling}\label{subsec:reson}
Fix a real $\chi\in C_c^\infty(-1,1)$, $\chi\ne0$. For $0<\varepsilon\le1$ put $\chi_\varepsilon(t)=\chi(t/\varepsilon)$ and
$\psi_\varepsilon:=(\frac14-\partial^2)\chi_\varepsilon\in\mathcal P_\varepsilon$, so $\widehat{\psi_\varepsilon}(\tau)=(\tau^2+\frac14)\varepsilon\widehat\chi(\varepsilon\tau)$.
For $M\ge1$, $b\in\R^M$ and $x_m:=\log m-\frac12\log M$ put
\[
F_b:=\sum_{m\le M}b_m\,V_{x_m}\psi_\varepsilon,
\]
the \emph{resonator} with coefficient vector $b$. For a real arithmetic function $a$ on $n\ge2$ write
$\lambda_{\mathrm A}(T;a)$ for the bottom on $\PT$ of the form with the weights $v_n(a)=2a(n)/\sqrt n$, as in
\S\ref{subsec:finite}, and $H(y):=\sum_{k\le y}1/k$.

\begin{lemma}[resonators]\label{lem:reson}
Let $\varepsilon\le1/(8M)$ and $T\ge\frac12\log M+\varepsilon$.
\emph{(i)} $F_b\in\PT\cap C_c^\infty(\IT)$ and $\|F_b\|^2=\|\psi_\varepsilon\|^2\|b\|^2$.
\emph{(ii)} For every $n\ge2$, $c_{F_b}(\log n)=\|\psi_\varepsilon\|^2\sum_{k\le M/n}b_{nk}b_k$.
\emph{(iii)} $\mT(F_b)\le(\log\frac1\varepsilon+C_\chi)\|F_b\|^2$, with $C_\chi$ depending only on $\chi$.
Consequently, for every real $a$ and $b\ne0$,
\[
\lambda_{\mathrm A}(T;a)\ \le\ \log\tfrac1\varepsilon+C_\chi-\frac1{\|b\|^2}\sum_{n\ge2}v_n(a)\sum_{k\le M/n}b_{nk}b_k .
\]
\end{lemma}

\begin{proof}
\emph{Separation.} Let $m,k\le M$ and $n\ge1$ with $nk\ne m$. If $nk>m$ then $nk/m\ge1+\frac1m$; if $nk<m$ then $m/(nk)\ge1+\frac1{nk}$
with $nk<M$. In both cases $|\log\frac{nk}m|\ge\log(1+\frac1M)\ge\frac1{2M}\ge4\varepsilon$.
(i) $A(V_x\psi)=e^{x/2}A(\psi)$ and $B(V_x\psi)=e^{-x/2}B(\psi)$; $A(\psi_\varepsilon)=B(\psi_\varepsilon)=0$ by two integrations by parts. The
supports lie in $(-\frac12\log M-\varepsilon,\frac12\log M+\varepsilon)\subset\IT$ and are disjoint by the separation with $n=1$.
(ii) $c_{F_b}(u)=\sum_{m,k}b_mb_kc_{\psi_\varepsilon}(u-\log\frac mk)$ with $\supp c_{\psi_\varepsilon}\subset(-2\varepsilon,2\varepsilon)$; at $u=\log n$ the
separation leaves only $m=nk$, where the term is $c_{\psi_\varepsilon}(0)=\|\psi_\varepsilon\|^2$.
(iii) By Lemma~\ref{lem:mbounds}, $m(\tau)\le\log(1+|\tau|)$ for all $\tau$ (for $|\tau|\ge1$ the upper bound there is below $\log|\tau|$, and
$m(\tau)\le m(1)<0$ for $|\tau|<1$). For $\varepsilon\le1$, $\log(1+|\tau|)\le\log\frac1\varepsilon+\log(1+\varepsilon|\tau|)\le\log\frac1\varepsilon+\frac12\log2+\ell(\varepsilon\tau)$
with $\ell(x)=\frac12\log(1+x^2)$. Since $|\Fh_b|^2=|\widehat{\psi_\varepsilon}|^2|\sum_mb_me^{i\tau x_m}|^2$,
\[
\mT(F_b)\le(\log\tfrac1\varepsilon+\tfrac12\log2)\|F_b\|^2+\sum_{m,k}b_mb_kJ(x_m-x_k),\qquad
J(u):=\frac1{2\pi\varepsilon^3}\int G_\varepsilon(x)e^{ixu/\varepsilon}dx,
\]
where $G_\varepsilon(x)=\ell(x)(x^2+\frac{\varepsilon^2}4)^2|\widehat\chi(x)|^2$ (substitute $\tau=x/\varepsilon$). The family $\{G_\varepsilon\}_{\varepsilon\le1}$ is bounded in
$W^{2,1}(\R)$; let $M_0,M_2$ bound $\|G_\varepsilon\|_1,\|G_\varepsilon''\|_1$ and $D_\chi:=\int x^4|\widehat\chi|^2$, so that
$\|\psi_\varepsilon\|^2\ge D_\chi/(2\pi\varepsilon^3)$. Then $J(0)\le(M_0/D_\chi)\|\psi_\varepsilon\|^2$ and, by two integrations by parts,
$|J(u)|\le(M_2/D_\chi)(\varepsilon/u)^2\|\psi_\varepsilon\|^2$. For $m\ne k$, $|x_m-x_k|=|\log\frac mk|\ge|m-k|/M$, so $(\varepsilon/u)^2\le\frac1{64|m-k|^2}$,
and the Schur test gives $\sum_{m\ne k}|b_mb_k||m-k|^{-2}\le2\zeta(2)\|b\|^2$. So (iii) holds with
$C_\chi=\frac12\log2+(M_0+M_2\pi^2/192)/D_\chi$. The final inequality is (i)--(iii) inserted in \eqref{eq:form}.
\end{proof}

\begin{theorem}[the integer density ceiling]\label{thm:ceiling}
Let $a$ be real with $\lambda_{\mathrm A}(T;a)\ge0$ for every $T$. Then for every $M\ge1$
\[
\mathsf G_M(a):=\frac2{H(M)}\sum_{2\le n\le M}\frac{a(n)}nH\Big(\frac Mn\Big)\ \le\ \log(8M)+C_\chi .
\]
If moreover $a\ge0$ and $S(x):=\sum_{2\le n\le x}a(n)/n$, then $\liminf_{x\to\infty}S(x)/\log x\le1$. In particular
$\sum_{n\le x}a(n)\sim\kappa x$ forces $\kappa\le1$, and $\sum_na(n)n^{-\sigma}\sim r/(\sigma-1)$ as $\sigma\downarrow1$ forces $r\le1$.
\end{theorem}

\begin{proof}
Apply Lemma~\ref{lem:reson} with $\varepsilon=1/(8M)$, $T=\frac12\log M+\varepsilon$ and $b_m=m^{-1/2}$: $\sum_kb_{nk}b_k=n^{-1/2}H(M/n)$ and
$\|b\|^2=H(M)$, so $0\le\log(8M)+C_\chi-\mathsf G_M(a)$. For $a\ge0$, $\sum_{n\le M}\frac{a(n)}nH(\frac Mn)=\sum_{k\le M}\frac1kS(\frac Mk)$,
and since $S$ is nondecreasing, $\frac1kS(\frac Mk)\ge\int_k^{k+1}S(\frac My)\frac{dy}y$, so this sum is at least
$\int_1^MS(x)\frac{dx}x$. If $S(x)\ge r'\log x-C$ for all $x\ge1$ with $r'>1$, then
$\mathsf G_M(a)\ge\frac{2}{\log M+1}(\frac{r'}2\log^2M-C\log M)=r'\log M-O(1)$, which contradicts the first bound for large $M$.
The density statement follows by partial summation. For the last, the Hardy--Littlewood--Karamata theorem
\cite[Thm.~I.15.1]{Korevaar2004} applied to the nondecreasing function $t\mapsto S(e^t)$ gives $S(x)\sim r\log x$.
\end{proof}

For $a=\Lambda$, $S(x)=\log x+O(1)$ by Mertens' theorem, so the prime comb sits exactly at the ceiling. Combined with
Theorem~\ref{thm:kappa}, the main-term symbol $g_r$ of a nonnegative comb on $\log\N$ that lies in $\mathfrak C^+$ satisfies
$g_r(0)=C_0+4r\le C_0+4<0$.

\begin{remark}[resolution against density]\label{rem:resolution}
The proof of Lemma~\ref{lem:reson} uses that the numbers $\log\frac mk$, $m,k\le M$, are $\frac1{2M}$ apart, so profiles of width
$\varepsilon\approx\frac1M$ see every atom $\log n$ exactly and buy an archimedean energy of only $\log M$. Against
$\rho_\kappa$ the comb part of a resonator is $-\kappa\,\mathfrak e(F_b)$ and does not grow with $M$. Against the lattice combs of
Theorem~\ref{thm:atomic} the admissible width is the spacing $h_j$, whose archimedean energy $\log\frac1{h_j}$ dominates. So
Theorem~\ref{thm:ceiling} expresses that $\log\N$ is too coarse to carry a density above $1$. It does not show by what mechanism
$\nu_\Lambda$ is positive.
\end{remark}

On resonators the near-coincidences $m\approx kn$, $m\ne kn$, are invisible, and there $\Lambda$ is positive with a margin.

\begin{proposition}[positivity on resonators]\label{prop:nearfree}
There are $K_0$ and $c_0>0$, depending only on $\chi$, with the following property. Let $K\ge K_0$ and
$\varepsilon_M:=1/(8KM(2+\log M))$. Then for every $M\ge1$, every $T\ge\frac12\log M+\varepsilon_M$ and every $b\in\R^M$, the
resonator $F_b$ built with $\varepsilon=\varepsilon_M$ satisfies $\mathfrak a_T(F_b)\ge c_0\|F_b\|^2$.
\end{proposition}

\begin{proof}
Write $\psi=\psi_{\varepsilon_M}$ and $\varepsilon=\varepsilon_M$. \emph{Comb.} By Lemma~\ref{lem:reson}(ii), $\sum_nv_nc_{F_b}(\log n)=\|\psi\|^2\cdot\frac12b^\top\mathbb Db$
with $\mathbb D_{mk}=v_{m/k}$ if $k\mid m$, $k<m$, $\mathbb D_{mk}=v_{k/m}$ if $m\mid k$, $m<k$, and $0$ otherwise. $\mathbb D$ has nonnegative entries.
The Schur test with weights $w_k=k^{-1/2}$ gives, for each row $m$,
$m^{1/2}\sum_k\mathbb D_{mk}k^{-1/2}=2\sum_{d\mid m}\Lambda(d)+2\sum_{j\le M/m}\frac{\Lambda(j)}j\le2\log m+2(\log\frac Mm+E)$, where
$E:=\sup_{x\ge1}(\sum_{j\le x}\Lambda(j)/j-\log x)<\infty$ by Mertens' theorem \cite[Thm.~2.7]{MontgomeryVaughan2007}. The identity
$\sum_{d\mid m}\Lambda(d)=\log m$ enters here. Hence the comb term is at most $(\log M+E)\|F_b\|^2$.
\emph{Archimedean diagonal.} By Lemma~\ref{lem:mbounds}(ii), $m(\tau)\ge\log\frac{|\tau|}{2\pi}-4$ for every $\tau\ne0$ (for $|\tau|<1$ the right side
is below $C_0$). Substituting $\tau=x/\varepsilon$, $\mT(\psi)\ge(\log\frac1\varepsilon-C'_\chi)\|\psi\|^2$ with
$C'_\chi=\log2\pi+4+\sup_{\varepsilon\le1}\int(\log|x|)^-(x^2+\frac{\varepsilon^2}4)^2|\widehat\chi|^2/D_\chi$.
\emph{Archimedean cross terms.} For $u_0\ge4\varepsilon$, \eqref{eq:m} gives
$\frac1{2\pi}\int m|\widehat\psi|^2\cos(\tau u_0)d\tau=-\frac12\int_0^\infty c_\psi(u_0-u)K(u)\,du$, because $c_\psi$ vanishes at $u_0$ and at
$u_0+u$; only $u\in(\frac{u_0}2,2u_0)$ contributes, and $K(u)\le\frac1u+\frac12$, so the modulus is at most
$\frac12\|\psi\|_1^2(\frac2{u_0}+\frac12)$. As $\|\psi\|_1\le\varepsilon^{-1}L_\chi$ with $L_\chi=\|\chi''\|_1+\frac14\|\chi\|_1$, this is at most
$C_1\varepsilon(\frac1{u_0}+1)\|\psi\|^2$ with $C_1=2\pi L_\chi^2/D_\chi$. With $u_0=|\log\frac mk|\ge|m-k|/M$ and the Schur test,
(note $u_0\ge\frac1M>4\varepsilon$) the cross terms total at most $C_1\varepsilon M(3+2\log M)\|b\|^2\|\psi\|^2\le\frac{C_1}{4K}\|F_b\|^2$.
\emph{Collecting}, $\mathfrak a_T(F_b)/\|F_b\|^2\ge\log\frac1{\varepsilon}-\log M-C'_\chi-\frac{C_1}{4K}-E\ge\log(16K)-C'_\chi-\frac{C_1}{4K}-E$, which is
at least $1$ for $K\ge K_0$.
\end{proof}

\begin{remark}\label{rem:boxes}
On a multiplicative box $S=\{\prod_ip_i^{k_i}:0\le k_i\le K_i\}$ the matrix $\mathbb D$ restricted to $S$ is a Kronecker sum of one-prime
Toeplitz matrices with symbols $\delta_p(\theta)=2\log p\,\Rel\frac x{1-x}$, $x=p^{-1/2}e^{i\theta}$, the logarithmic derivatives of the
local Euler factors. Their ranges are $[-\frac{2\log p}{\sqrt p+1},\frac{2\log p}{\sqrt p-1}]$, so $\|\mathbb D|_S\|\le\sum_i\frac{2\log p_i}{\sqrt{p_i}-1}$
independently of the exponents. Proposition~\ref{prop:nearfree} shows that tests on which no near-coincidence is visible are
decided by the archimedean energy at the resolution of the test against such local bounds. They carry no information on the
sign beyond that, and this is a barrier for arguments that use only such tests (\S\ref{subsec:bar-arith}).
\end{remark}

\subsection{Translates of one profile}\label{subsec:translates}
For a compactly supported $F$ in the form domain, $\mathfrak a_T(F)$ does not depend on $T$ once $\supp F\subset\IT$
(Lemma~\ref{lem:wit-freeze}); we write $\mathfrak a(F)$, $\mathfrak a^\nu(F)$ for this common value and use the same letters for the
associated real bilinear forms.

\begin{lemma}[translates]\label{lem:transl}
Let $F_0\in\mathcal P_{T_0}$ be real, $\|F_0\|=1$, and $c_0:=c_{F_0}$.
\emph{(i)} $V_aF_0\in\PT$ whenever $a+I_{T_0}\subset\IT$; no correction of the Tate moments is needed.
\emph{(ii)} For real $s_k$ and $a_k$, $c_{\sum_ks_kV_{a_k}F_0}(u)=\sum_{k,l}s_ks_l\,c_0(u-(a_k-a_l))$.
\emph{(iii)} $\int c_0(u-D)e^{\pm u/2}du=0$ for every $D\in\R$.
\emph{(iv)} For $D>2T_0$, $\int c_0(v)K(D-v)\,dv=2\sum_{k\ge1}e^{-(2k+\frac12)D}\int c_0(v)e^{(2k+\frac12)v}dv$, and
\[
\Big|\int c_0(v)K(D-v)\,dv\Big|\ \le\ \frac{4T_0\,e^{-5(D-2T_0)/2}}{1-e^{-2(D-2T_0)}} .
\]
\end{lemma}

\begin{proof}
(i) $A(V_aF_0)=e^{a/2}A(F_0)$, $B(V_aF_0)=e^{-a/2}B(F_0)$. (ii) Bilinearity. (iii) $\int c_0(v)e^{v/2}dv=A(F_0)B(F_0)=0$, likewise with
$e^{-v/2}$. (iv) Expand $K$ by Lemma~\ref{lem:K}; the term $k=0$ is $2e^{-D/2}\int c_0e^{v/2}=0$ by (iii); and
$\|c_0\|_1\le\|F_0\|_1^2\le2T_0$.
\end{proof}

\begin{theorem}[the two-sided band]\label{thm:band}
Let $\nu$ be a comb, $F_0$ as in Lemma~\ref{lem:transl}, $a_0:=\mathfrak a^\nu(F_0)$, $D>2T_0$, and
\[
X_\nu(D):=\int c_0(u-D)\,d\nu(u)+\int c_0(v)K(D-v)\,dv .
\]
For every $s\in\R$ and $T\ge T_0+D/2$, $\mathfrak a^\nu_T(V_{-D/2}F_0+sV_{D/2}F_0)=(1+s^2)a_0-sX_\nu(D)$. Hence $\mathfrak a^\nu_T\ge0$ on this
two-dimensional subspace of $\PT$ if and only if $|X_\nu(D)|\le2a_0$. In particular, if \eqref{eq:target} holds, then for every such
$F_0$ and every $D>2T_0$
\[
\big|\Xi_{F_0}(D)\big|\ \le\ 2\,\mathfrak a(F_0)+\frac{4T_0\,e^{-5(D-2T_0)/2}}{1-e^{-2(D-2T_0)}},
\]
where
\[
\Xi_{F_0}(D):=\sum_{n\ge2}v_nc_0(\log n-D)=2\int_{[1,\infty)}y^{-1/2}c_0(\log y-D)\,d\big(\psi(y)-y\big).
\]
\end{theorem}

\begin{proof}
By Lemma~\ref{lem:transl}(ii) the autocorrelation is $(1+s^2)c_0+s[c_0(\cdot-D)+c_0(\cdot+D)]$. By \eqref{eq:m}, the archimedean cross
term is $s\big[2C_0c_0(D)+\int_0^\infty(2c_0(D)-c_0(D-u)-c_0(D+u))K(u)du\big]=-s\int c_0(v)K(D-v)dv$, since $c_0$ vanishes at $D$ and at
$D+u$. On $(0,2T)$ the comb sees $(1+s^2)c_0$ on $(0,2T_0)$ and $sc_0(\cdot-D)$ on $(D-2T_0,D+2T_0)\subset(0,2T)$. This gives the
identity; the matrix $\big(\begin{smallmatrix}a_0&-X/2\\-X/2&a_0\end{smallmatrix}\big)$ is positive semidefinite iff $|X|\le2a_0$. For
$\nu=\nu_\Lambda$ use Lemma~\ref{lem:transl}(iv); the integral form of $\Xi_{F_0}$ holds because the main term
$2\int c_0(\log y-D)y^{-1/2}dy=2\int c_0(u-D)e^{u/2}du$ vanishes by Lemma~\ref{lem:transl}(iii).
\end{proof}

So the smoothed Chebyshev fluctuation $\Xi_{F_0}(D)$ at scale $e^D$, with its main term removed exactly by the Tate conditions, must
lie in the fixed band $[-2\mathfrak a(F_0),2\mathfrak a(F_0)]$ at every scale, up to $O(e^{-5D/2})$. If $2T_0<\log2$, then
$\mathfrak a(F_0)=\mT(F_0)$. The band is the $2\times2$ case of Theorem~\ref{thm:oneprofile} below.

The band fails for random combs with overwhelming probability. Let $\kappa>0$ and let $\Pi$ be a Poisson point process on
$(\log2,\infty)$ with intensity $\kappa e^u/u\,du$; put $\nu_\Pi:=\sum_{x\in\Pi}2xe^{-x/2}\delta_x$, whose mean measure is
$\rho_\kappa|_{(\log2,\infty)}$.

\begin{theorem}[random combs]\label{thm:poisson}
For every $c>0$ there is $T(c,\kappa)$ such that for $T\ge T(c,\kappa)$ the probability that $\mathfrak a^{\nu_\Pi}_T\ge0$ on $\PT$ is at
most $\exp(-cT\log T)$.
\end{theorem}

\begin{proof}
Fix $T_0<\min(\frac1{4c},\frac12\log2)$ and a real unit $F_0\in\mathcal P_{T_0}$. By Corollary~\ref{cor:Tarch},
$a_0:=\mT(F_0)\ge\laminfP(T_0)>0$; since $\nu_\Pi$ has no mass on $(0,2T_0]$, $a_0=\mathfrak a^{\nu_\Pi}(F_0)$. Let $D_k:=1+4T_0k$ for
$1\le k\le K:=\lfloor(2T-4T_0-1)/(4T_0)\rfloor$, so that $T\ge T_0+D_k/2$ and the intervals $I_k:=(D_k-2T_0,D_k+2T_0)\subset(\log2,\infty)$ are
disjoint. By Theorem~\ref{thm:band} the event in question lies in $\bigcap_k\{|X_{\nu_\Pi}(D_k)|\le2a_0\}$. Here
$X_{\nu_\Pi}(D_k)=R_k+d_k$ with $d_k$ deterministic and $R_k=\sum_{x\in\Pi\cap I_k}f_k(x)$, $f_k(x)=2xe^{-x/2}c_0(x-D_k)$; the $R_k$
are independent \cite[\S2.2]{Kingman1993}. By Esseen's inequality \cite[Ch.~III, Lemma~3]{Petrov1975},
$\sup_y\mathbb P(|R_k-y|\le2a_0)\le C_E2a_0\int_{|t|\le1/(2a_0)}|\phi_k(t)|dt$, where by Campbell's theorem
\cite[\S3.2]{Kingman1993} $|\phi_k(t)|=\exp(-\int_{I_k}(1-\cos tf_k(x))\kappa e^x\frac{dx}x)$. Since $\|c_0\|_\infty\le1$,
$|tf_k|\le(D_k+2T_0)e^{-(D_k-2T_0)/2}/a_0\le\pi$ for $D_k\ge D_*(a_0,T_0)$, and then $1-\cos y\ge\frac2{\pi^2}y^2$ gives
$|\phi_k(t)|\le\exp(-\frac2{\pi^2}t^2\sigma_k^2)$ with
$\sigma_k^2=4\kappa\int_{I_k}xc_0(x-D_k)^2dx\ge4\kappa(D_k-2T_0)\|c_0\|_2^2$. Hence
$\mathbb P(|X_{\nu_\Pi}(D_k)|\le2a_0)\le C_2(D_k-2T_0)^{-1/2}$ with $C_2$ depending on $a_0,\kappa,\|c_0\|_2$. By independence, and
$D_k-2T_0\ge4T_0k$, the probability is at most $\prod_{k_0\le k\le K}\min(1,C_2(4T_0k)^{-1/2})$, whose logarithm is
$-\frac12K\log K+O(K)$. As $K\sim T/(2T_0)$, this is $-\frac{T\log T}{4T_0}(1+o(1))\le-cT\log T$ for large $T$.
\end{proof}

The lattice combs of Theorem~\ref{thm:atomic} satisfy the band trivially, since their $|X(D)|$ is of the size of the quadrature
error. A Poisson comb of any density, including $\kappa>\kappa_\star$, violates it at a fixed scale with probability
$1-O(D^{-1/2})$, independently at disjoint scales.

\subsection{One profile and the integer lattice}\label{subsec:oneprofile}
Fix $a>0$ and a real $b$ with $|b|>\frac12$. Put $g(t)=e^{bt}\mathbf 1_{[0,a]}(t)$, $c_g=g*\tilde g$ with $\tilde g(t)=g(-t)$,
$G_0:=c_g*c_g$, $T_0:=2a$, and $F_0:=(\frac14-\partial^2)G_0$. Then $G_0$ is even, supported in $[-T_0,T_0]$, of class $C^2$ with
$G_0,G_0'$ vanishing at $\pm T_0$, so $G_0\in H^2_0(I_{T_0})$ and $F_0\in\mathcal P_{T_0}$ (Theorem~\ref{thm:nyq-cone}). Moreover
\[
\Fh_0(\tau)=(\tau^2+\tfrac14)|\widehat g(\tau)|^4>0\quad(\tau\in\R),\qquad \widehat g(z)=\frac{e^{(b+iz)a}-1}{b+iz},
\]
and the zeros of $\widehat g(z)\widehat g(-z)$ are the points $\pm(\frac{2\pi k}a+ib)$, $k\ne0$, on the lines $|\Iml z|=|b|$. We call $F_0$ the
\emph{profile}, write $c_0=c_{F_0}$, so that $\widehat{c_0}=\Fh_0^2$, and define for $d\in\R$
\[
\mathsf B(d):=\mathfrak a(F_0,V_dF_0)=\frac1{2\pi}\int m\,\Fh_0^2\cos(\tau d)\,d\tau-\frac12\sum_{n\ge2}v_n\big[c_0(\log n-d)+c_0(\log n+d)\big],
\]
a finite sum for each $d$. $\mathsf B$ is real, even and continuous, and for real $s_k$, $\mathfrak a(\sum_ks_kV_{a_k}F_0)=\sum_{k,l}s_ks_l\mathsf B(a_k-a_l)$.

\begin{theorem}[one profile suffices]\label{thm:oneprofile}
\eqref{eq:target} holds if and only if $\mathsf B$ is a positive definite function on $\R$.
\end{theorem}

\begin{proof}
($\Rightarrow$) Given $a_1,\dots,a_K$, take $T$ with all $a_k+I_{T_0}\subset\IT$; by Lemma~\ref{lem:transl}(i) the function
$\sum_ks_kV_{a_k}F_0$ lies in $\PT$, so $\sum_{k,l}s_ks_l\mathsf B(a_k-a_l)\ge0$ for real $s$; for complex $s$ split into real and imaginary
parts, the matrix being real symmetric.

($\Leftarrow$) By Bochner's theorem $\mathsf B(d)=\int e^{i\tau d}d\mu(\tau)$ with $\mu$ a finite positive measure.
\emph{Step 1.} For even $\varphi\in C_c^\infty(\R)$, Fubini gives
$\int\mathsf B\varphi=\frac1{2\pi}\int m\Fh_0^2\widehat\varphi-\sum_nv_n(c_0*\varphi)(\log n)$, a finite sum, and $\int\mathsf B\varphi=\int\widehat\varphi\,d\mu$.
\emph{Step 2.} Let $G\in\PT\cap C_c^\infty(\IT)$ be real and $\Psi:=\widehat{c_G}/\Fh_0^2=|\widehat G|^2/\Fh_0^2\ge0$. As entire functions,
$\widehat{c_G}(z)=\widehat G(z)\widehat G(-z)$ vanishes to order $\ge2$ at $\pm\frac i2$ (because $\widehat G(\pm\frac i2)=0$), and
$\Fh_0(z)^2=(z^2+\frac14)^2\widehat g(z)^4\widehat g(-z)^4$ vanishes to order exactly $2$ there and otherwise only on $|\Iml z|=|b|$. So $\Psi$
is holomorphic in $|\Iml z|<|b|$. For $|\Iml z|\le\beta<|b|$ one has $|e^{(b+iz)a}-1|\ge|e^{(b-\Iml z)a}-1|\ge c_\beta>0$, hence
$|\widehat g(\pm z)|^{-1}\le C(1+|z|)$, while $\widehat G(x+iy)=O_N(e^{T|y|}(1+|x|)^{-N})$. Thus $\Psi$ decreases rapidly on horizontal
lines in the closed substrip, and a contour shift gives $|\check\Psi(u)|\le C_\beta e^{-\beta|u|}$; also $c_0*\check\Psi=c_G$, since both
sides have transform $\widehat{c_G}$. Fix $\beta\in(\frac12,|b|)$.
\emph{Step 3.} Let $\chi\in C_c^\infty$ be even, $0\le\chi\le1$, $\chi=1$ on $[-1,1]$, and $\varphi_j:=\check\Psi\,\chi(\cdot/j)$, which is even.
Then $\varphi_j\to\check\Psi$ in $L^1$, so $\widehat{\varphi_j}\to\Psi$ uniformly and boundedly; hence $\int\widehat{\varphi_j}d\mu\to\int\Psi d\mu$
and $\int m\Fh_0^2\widehat{\varphi_j}\to\int m\Fh_0^2\Psi=\int m|\widehat G|^2$, as $m\Fh_0^2\in L^1$. Moreover
$|(c_0*\varphi_j)(\log n)|\le\|c_0\|_1\sup_{|w-\log n|\le2T_0}|\check\Psi(w)|\le C'n^{-\beta}$ and $\sum_nv_nn^{-\beta}=\sum_n2\Lambda(n)n^{-1/2-\beta}<\infty$, so by
dominated convergence $\sum_nv_n(c_0*\varphi_j)(\log n)\to\sum_nv_nc_G(\log n)=\sum_{n\in\NT}v_nc_G(\log n)$. Step~1 gives
$\mathfrak a_T(G)=\int\Psi\,d\mu\ge0$.
\emph{Step 4.} Every $F\in\PT\cap\mathcal D_T$ is $(\frac14-\partial^2)H$ with $H\in H^2_0(\IT)$, and the form norm of $F$ is
$\big(\frac1{2\pi}\int(1+\log(2+|\tau|))(\tau^2+\frac14)^2|\widehat H|^2\big)^{1/2}$. Contracting $H$ towards the centre of $\IT$ and
mollifying gives $H_k\in C_c^\infty(\IT)$ converging to $H$ in this weighted norm, and $(\frac14-\partial^2)H_k\in\PT\cap C_c^\infty(\IT)$
automatically. (Dilating $F$ itself would not preserve the Tate conditions.) The form is continuous in the form norm, and real
and imaginary parts are treated separately (proof of Proposition~\ref{prop:finite}(ii)). Hence $\lamA\ge0$ for every $T$.
\end{proof}

The condition $|b|>\frac12$ is used in Step~3 alone: the prime sum converges because the profile is analytic in a strip wider
than the polar points $\pm\frac i2$. Theorem~\ref{thm:oneprofile} is a reformulation, of the same nature as the Kre\u{\i}n
extension problem of \S\ref{subsec:KEP}. Its new feature is that one fixed explicit function of one real variable, with no window
parameter, carries the whole statement \eqref{eq:target}. Theorem~\ref{thm:band} is its $2\times2$ case. Restricting the points to
logarithms of integers loses nothing.

\begin{proposition}[the integer-lattice matrices]\label{prop:AM}
Assume $2T_0<\log2$. For $M\ge2$ let $\mathbb A_M:=[\mathsf B(\log m-\log k)]_{m,k\le M}$ and $P_s(\tau):=\sum_{m\le M}s_mm^{i\tau}$.
\emph{(i)} \eqref{eq:target} holds if and only if $\mathbb A_M\succeq0$ for every $M$.
\emph{(ii)} For $s\in\R^M$,
\[
s^\top\mathbb A_Ms=\frac1{2\pi}\int\Fh_0(\tau)^2|P_s(\tau)|^2\Big[m(\tau)-\sum_{2\le n<Me^{2T_0}}v_n\cos(\tau\log n)\Big]d\tau .
\]
\emph{(iii)} $\mathbb A_M=\mathrm{Arch}_M-\mathrm{Div}_M-\mathrm{Near}_M$, where $\mathrm{Arch}_M$ has entries
$\frac1{2\pi}\int m\Fh_0^2\cos(\tau\log\frac mk)\,d\tau$; $\mathrm{Div}_M=\frac12(\mathbb V_M+\mathbb V_M^\top)$ with $(\mathbb V_M)_{mk}=v_{m/k}$ if $k\mid m$,
$k<m$, and $0$ otherwise, the Dirichlet convolution by $v$; and, for $m>k$,
\[
\mathrm{Near}_M(m,k)=\frac12\sum_{n\ge2,\ kn\ne m}v_n\,c_0\Big(\log\frac{kn}m\Big),
\]
a smoothed sum of $\Lambda(n)n^{-1/2}$ over the integers $n\ne\frac mk$ in the multiplicative interval
$(\frac mke^{-2T_0},\frac mke^{2T_0})$. The diagonals of $\mathrm{Div}_M$ and $\mathrm{Near}_M$ vanish, and
$(\mathrm{Div}_M+\mathrm{Near}_M)_{mk}=\frac12\Xi_{F_0}(\log\frac mk)$ for $m>k$.
\end{proposition}

\begin{proof}
(i) ($\Rightarrow$) is Theorem~\ref{thm:oneprofile}. ($\Leftarrow$) Given $a_1,\dots,a_K$ and real $s$, let $m_j:=\lfloor e^{C+a_j}\rfloor$. For large
$C$ the $m_j$ are distinct and $\log m_j-\log m_l\to a_j-a_l$, so by continuity of $\mathsf B$,
$\sum s_js_l\mathsf B(a_j-a_l)=\lim_C\sum s_js_l\mathsf B(\log m_j-\log m_l)\ge0$. Then apply Theorem~\ref{thm:oneprofile}.
(ii) $\frac1{2\pi}\int\Fh_0^2|P_s|^2\cos(\tau\log n)d\tau=\sum_{m,k}s_ms_k\frac12[c_0(\log\frac mk+\log n)+c_0(\log\frac mk-\log n)]$, and the
terms with $n\ge Me^{2T_0}$ vanish because $|\log n\pm\log\frac mk|\ge2T_0$. (iii) For $m>k$, $c_0(\log n+\log\frac mk)=0$ as
$\log n\ge\log2>2T_0$; the term $n=m/k$ has $c_0(0)=1$, and the remaining terms need $|\log\frac{kn}m|<2T_0$. On the diagonal
$c_0(\log n)=0$ for all $n\ge2$.
\end{proof}

So \eqref{eq:target} is the positive semidefiniteness of $\mathrm{Arch}_M-\mathrm{Div}_M-\mathrm{Near}_M$ for all $M$. Here
$\mathrm{Arch}_M$ is explicit, and $\mathrm{Div}_M$ is the Dirichlet convolution by the local Euler-factor weights. $\mathrm{Near}_M$
is the local Chebyshev fluctuation $\Xi_{F_0}$ at the multiplicative scales $\frac mk$, with the exact coincidences $n=\frac mk$
removed. Proposition~\ref{prop:nearfree} shows that tests on which $\mathrm{Near}$ is invisible are positive; for the fixed profile
$F_0$, near-coincidences $|\log\frac{kn}m|<2T_0$ occur for all large $M$. The sign therefore has to be decided on test vectors that
see $\mathrm{Near}_M$ (\S\ref{sec:open}).


\section{Barriers for classes of arguments}\label{sec:barriers}

Section~\ref{sec:negative} limits certificates and prime-side hypotheses. This section does the same for five
further kinds of argument: linear bounds in a fixed frame, arguments that use the family of windows,
arguments that use observable patterns of one window, local or finite-rank criteria on the symbol, and arguments
from arithmetic structure, orders and sums of squares.
Each subsection declares a class, proves a theorem about it, and states what the class does not exclude.
All the barriers rest on the following principle.

\begin{lemma}[transfer]\label{lem:transfer}
Let $\mathfrak X$ be a set of closed semibounded forms, each on a closed subspace of some $L^2(I_{T'})$, and let
$\mathsf P$ be a property defined on $\mathfrak X$. Suppose $q_0,q_1\in\mathfrak X$ both have $\mathsf P$, $q_0$ is the
form of $\zeta$ in question, and the bottom of $q_1$ is negative. Then there is no valid deduction of the form
``$q\in\mathfrak X$ has $\mathsf P$, hence the bottom of $q$ is nonnegative''.
\end{lemma}
\begin{proof}
Such a deduction applies to $q_1$ and yields a false conclusion.
\end{proof}

A barrier is an instance of Lemma~\ref{lem:transfer} in which $q_1$, the \emph{falsifier}, is given explicitly. It
excludes the deduction, not the conclusion. Where the negativity of $q_1$ is only computed, the barrier is stated
as numerical.

\subsection{Linear bounds in a fixed frame}\label{subsec:bar-lin}

The first class consists of lower bounds for $\mathfrak a_T$ on $\PT$ that are obtained by bounding the prime
form $\PiT$ term by term in norm, and of dual certificates whose exterior measure has a sign.

\begin{lemma}[the one-half lemma]\label{lem:half}
Let $T<a<2T$. For every $F\in L^2(\IT)$, $|\cF(a)|\le\frac12\|F\|^2$. The constant is optimal on $\PT$:
$\sup\{|\cF(a)|:F\in\PT,\ \|F\|=1\}=\frac12$.
\end{lemma}
\begin{proof}
$\cF(a)=\int_{J_1}F(t)\overline{F(t-a)}\,dt$ with $J_1=[a-T,T]$, and $t-a$ then runs over $J_2=[-T,T-a]$. Since
$a>T$, $J_1\cap J_2=\emptyset$. By Cauchy--Schwarz and the arithmetic--geometric mean inequality,
$|\cF(a)|\le\|F\mathbf 1_{J_1}\|\,\|F\mathbf 1_{J_2}\|\le\frac12(\|F\mathbf 1_{J_1}\|^2+\|F\mathbf 1_{J_2}\|^2)\le\frac12\|F\|^2$.
For optimality let $0<\delta<2T-a$ and $h_\delta=(2\delta)^{-1/2}(\mathbf 1_{[-T,-T+\delta]}+\mathbf 1_{[a-T,a-T+\delta]})$.
The two intervals are disjoint, so $\|h_\delta\|=1$, and $c_{h_\delta}(a)=\frac12$: the only overlap of
$h_\delta$ with $h_\delta(\cdot-a)$ is the second interval against the translate of the first. Let $p_\delta$ be
the orthogonal projection of $h_\delta$ onto $\mathrm{span}\{e^{t/2},e^{-t/2}\}\subset L^2(\IT)$ and
$G_\delta=h_\delta-p_\delta\in\PT$. Since $|A(h_\delta)|,|B(h_\delta)|\le e^{T/2}\|h_\delta\|_{L^1}\le e^{T/2}(2\delta)^{1/2}$
and the Gram matrix of $e^{\pm t/2}$ on $\IT$ is invertible, $\|p_\delta\|=O(\delta^{1/2})$. The map
$(F,G)\mapsto\int F(t)\overline{G(t-a)}dt$ is sesquilinear and bounded by $\|F\|\|G\|$, so
$|c_{G_\delta}(a)-\frac12|\le2\|p_\delta\|+\|p_\delta\|^2$ and $\|G_\delta\|\to1$. Hence
$|c_{G_\delta}(a)|/\|G_\delta\|^2\to\frac12$ as $\delta\to0$.
\end{proof}

\begin{corollary}[the direct bound]\label{cor:directbound}
For every $T>0$,
\[
\lamA\ \ge\ \laminfP(T)-\frac12\sum_{n\in\NT,\ \log n>T}v_n-\sum_{n\in\NT,\ \log n\le T}v_n .
\]
In particular $\lamA\ge\laminfP(T)-\frac12v_2=\laminfP(T)-0.49013\ldots$ for $\frac12\log2<T\le\frac12\log3$.
For a single contact ($\frac12\log2<T\le\frac12\log3$) the constant $\frac12$ in a bound
$|\Rel\cF(\log2)|\le\kappa\|F\|^2$ on $\PT$ cannot be lowered.
\end{corollary}
\begin{proof}
On $\PT$, $\mathfrak a_T=\mT+\PiT$ and $\mT\ge\laminfP(T)\|\cdot\|^2$. Bound $|\Rel\cF(\log n)|$ by
Lemma~\ref{lem:half} when $\log n>T$ and by $\|F\|^2$ otherwise. For $\frac12\log2<T\le\frac12\log3$ the only
contact is $n=2$, and $T<\log2<2T$; $v_2=\sqrt2\log2$. Optimality is the second statement of Lemma~\ref{lem:half}.
\end{proof}

\begin{remark}\label{rem:directbound}
The bound of Corollary~\ref{cor:directbound} is positive only where $\laminfP>0.49$. Numerically (Galerkin, not
certified) this holds up to $T\approx0.367$, below $\frac12\log3$; the extremal functions of
Lemma~\ref{lem:half} concentrate at the edges of the window, where the minimizer of $\mathfrak a_T$ has little mass.
Replacing $\mT$ by any other positive frame form and bounding the remainder by its operator norm does not help:
the bound is then at most the one above, because the prime form enters exactly and with the extremal constant
$\frac12$.
\end{remark}

We turn to certificates of the complete form, $\beta=(2,0)$ (Proposition~\ref{prop:P}). Write
$f=\sT+\nuh+2\gh0$ for the function of a certificate $(\nu,2,0)$.

\begin{proposition}[accounting at the origin]\label{prop:acct}
For every finite even measure $\nu$ on $\{|u|\ge2T\}$,
\[
f(0)=-\Delta(T)+\nu(\R),\qquad \Delta(T):=-C_0+\Aconst-8\sinh T .
\]
Hence a certificate $(\nu,2,0)$ of value $c$ has $\nu(\R)\ge c+\Delta(T)$. One has
$\Delta(\frac12\log3)=1.733639\ldots$ and $\Delta(\frac12\log5)=1.158739\ldots$.
\end{proposition}
\begin{proof}
$m(0)=C_0$ by \eqref{eq:m}, so $\sT(0)=C_0-\Aconst$; $\gh0(0)=2\int_0^{2T}\cosh\frac u2\,du=4\sinh T$; and
$\nuh(0)=\nu(\R)$. The values follow from $\Aconst=v_2$, resp.\ $v_2+v_3+v_4$, and $C_0=-5.3721834\ldots$.
\end{proof}

The positive mass that Proposition~\ref{prop:acct} requires cannot be supplied by a positive measure.

\begin{theorem}[positive exterior measures]\label{thm:posext}
Let $\mu$ be an even probability measure on $\R$ with $\int\log(2+|\tau|)\,d\mu<\infty$ whose inverse transform
$\Phi=\check\mu$ satisfies $\Phi\le0$ on $\{|u|\ge2T\}$. Then for every finite positive even measure $\nu$ on
$\{|u|\ge2T\}$,
\begin{equation}\label{eq:posext}
\begin{split}
\inf_{\tau\in\R}\big[\sT+\nuh+2\gh0\big](\tau)\ \le\ \kappa_T(\Phi):={}&C_0+\int_0^\infty\big(1-\Phi(u)\big)K(u)\,du\\
&-\sum_{n\in\NT}v_n\Phi(\log n)+2\int_{-2T}^{2T}\Phi(u)\cosh\tfrac u2\,du .
\end{split}
\end{equation}
Admissible $\Phi$ are given by $\Phi(u)=(e^{-\alpha u^2}-b\,e^{-\beta u^2})/(1-b)$ with $\alpha>\beta>0$ and
$e^{-4T^2(\alpha-\beta)}\le b\le\sqrt{\beta/\alpha}$.
\end{theorem}
\begin{proof}
The infimum of $f$ is at most its $\mu$-mean. By Tonelli and \eqref{eq:m},
$\int m\,d\mu=C_0+\int_0^\infty(1-\Phi)K$; $\int\cos(\tau\log n)\,d\mu=\Phi(\log n)$;
$\int\gh0\,d\mu=\int g_0\Phi=\int_{-2T}^{2T}\Phi\cosh\frac u2$; and $\int\nuh\,d\mu=\int\Phi\,d\nu\le0$ by Fubini,
since $\nu\ge0$ lives where $\Phi\le0$. For the family: $\Phi(0)=1$; the transform of $e^{-\alpha u^2}$ is
$\sqrt{\pi/\alpha}\,e^{-\tau^2/4\alpha}$, and $\sqrt{\pi/\alpha}\,e^{-\tau^2/4\alpha}\ge b\sqrt{\pi/\beta}\,e^{-\tau^2/4\beta}$
for all $\tau$ because $\alpha>\beta$ and $b\le\sqrt{\beta/\alpha}$, so $\Phi$ is the inverse transform of a
probability measure with Gaussian tails; and $\Phi(u)\le0$ iff $e^{-(\alpha-\beta)u^2}\le b$, which holds for
$|u|\ge2T$ by the lower bound on $b$.
\end{proof}

\begin{remark}\label{rem:posext}
Numerically (the integrals in \eqref{eq:posext} are elementary and were evaluated to $25$ digits in floating
point, not enclosed), $b=\sqrt{\beta/\alpha}$ and $(\alpha,\beta)=(2.312,1.926)$ give
$\kappa_T(\Phi)=-0.1377$ at $T=\frac12\log3$, and $(\alpha,\beta)=(1.068,0.890)$ give $-0.1284$ at $T=\frac12\log5$.
So, subject to these evaluations, no certificate of the complete form with $\nu\ge0$ and value $\ge0$ exists at these
windows, although $\lamW>0$ there (Theorem~\ref{thm:log5W}): by Proposition~\ref{prop:acct} the exterior measure must carry net
positive mass, and by Theorem~\ref{thm:posext} it cannot be positive. The mechanism is visible in the proof: the
mass that repairs $\tau=0$ produces, through atoms near $|u|=2T$, a negative contribution on the first negative lobe
of $2\gh0$, at $\tau\in(\pi/2T,\pi/T)$. The certificate of Theorem~\ref{thm:C} has total variation $90.75$ against net
mass $\nu(\R)\approx1.73$.
\end{remark}

The prime side contributes no sign of its own to a bound that treats the comb term by term.

\begin{proposition}[the comb is translation blind]\label{prop:primeblind}
Let $P(\tau)=\sum_{n\in\NT}v_n\cos(\tau\log n)$. For every $\varepsilon>0$ the set $\{\tau:P(\tau)>\Aconst-\varepsilon\}$ is
relatively dense in $\R$: there is $L_\varepsilon>0$ such that every interval of length $L_\varepsilon$ meets it.
Consequently, if $\sT\ge m-E$ on $\R$ for some function $E$, then every interval of length $L_\varepsilon$ contains a
point where $E>\Aconst-\varepsilon$: a lower bound of the window symbol that treats the comb uniformly in $\tau$ costs
the full mass $\Aconst$.
\end{proposition}
\begin{proof}
$P$ is a trigonometric polynomial, hence almost periodic, and $\sup P=\Aconst$ by Lemma~\ref{lem:mbounds}(iii). Pick
$\tau_*$ with $P(\tau_*)>\Aconst-\varepsilon/2$. The $\varepsilon/2$-almost periods $t$ of $P$
($\sup_\tau|P(\tau+t)-P(\tau)|<\varepsilon/2$) form a relatively dense set with some inclusion length $L_\varepsilon$
\cite[Ch.~VI]{Katznelson}, and $P(\tau_*+t)>\Aconst-\varepsilon$ for each of them. At such points $E\ge P$.
\end{proof}

\begin{remark}[what this class does not exclude]\label{rem:bar-lin}
Lemma~\ref{lem:half}, Corollary~\ref{cor:directbound}, Theorem~\ref{thm:posext}, Proposition~\ref{prop:primeblind},
Lemma~\ref{lem:Rrig} with Corollary~\ref{cor:nontermwise} (no flat-top kernel is positive), Theorem~\ref{thm:N}
(small lattices) and the ceiling of Proposition~\ref{prop:TK} exclude, respectively: norm bounds of the prime
form; certificates of the complete form with a positive exterior measure; term-by-term bounds of the comb;
termwise positivity on the zero side; and certificates of small lattice order. They do not exclude
\emph{signed} exterior measures, which is where the certificate of Theorem~\ref{thm:C} lives, nor primitive
certificates with $\nu\ge0$ and free $\beta$: the weak-duality argument above requires $\int\gh j\,d\mu=0$ in that
case, which the Gaussian family does not satisfy, and we have only a numerical indication (a discretized linear
program) that such certificates also fail. By Theorem~\ref{thm:R} the lattice classes are not lossy in the limit.
\end{remark}

\subsection{Families of windows}\label{subsec:bar-fam}

The second class consists of arguments that use the window bottoms themselves as data: the values $\lambda(T')$ for
$T'$ in a set $J$, together with the qualitative properties of $T\mapsto\lambda(T)$. For a nonnegative arithmetic
function $a$ we use the notation $\lambda_{\mathrm A}(T;a)$ of \S\ref{subsec:finite}.

\begin{lemma}[transport between nested windows]\label{lem:transport}
Let $0<T<T'$. \emph{(i)} Extension by zero maps $\PT$ isometrically into $\mathcal P_{T'}$ with
$\mathfrak a_{T'}(F)=\mathfrak a_T(F)$; there is no boundary term and no term from the contacts in
$\mathcal N_{T'}\setminus\NT$. \emph{(ii)} If $(\nu',\beta)$ is a certificate of value $c$ at the window $T'$, then
$\nu_T:=\nu'-\frac12\sum_{n\in\mathcal N_{T'}\setminus\NT}v_n(\delta_{\log n}+\delta_{-\log n})$ is supported in
$\{|u|\ge2T\}$ and $(\nu_T,\beta)$ is a certificate of value $c$ at the window $T$, with the same function $f$.
\end{lemma}
\begin{proof}
(i) is the proof of Lemma~\ref{lem:monoT}. (ii) $\sigma_{T'}=\sT-\sum_{n\in\mathcal N_{T'}\setminus\NT}v_n\cos(\tau\log n)$,
and the transform of the added atoms is $-\sum v_n\cos(\tau\log n)$; the atoms lie at $|u|=\log n\ge2T$. For a
flat-top weight $w\equiv1$ on $[-2T',2T']$ this is the identity $\nu_w^{(T)}=\nu_w^{(T')}-\frac12\sum v_n(\delta_{\pm\log n})$
read off from \eqref{eq:nuw}.
\end{proof}

So the only comparison that nesting provides is $\lamA\ge\lambda_{\mathrm A}(T')$, and inherited certificates carry
the negative atoms $-\frac12v_n$ at the new contacts.

\begin{theorem}[bounded families decide nothing]\label{thm:famblind}
Let $J\subset(0,T_1]$ with $T_1<\infty$. Let $\mathsf P$ be any property of a function $\lambda:(0,\infty)\to\R$ that
depends only on the restriction $\lambda|_J$, on monotonicity and on continuity, and let $\mathsf C$ be any class of
nonnegative arithmetic functions containing $\Lambda$ and closed under amplification of one term
(Theorem~\ref{thm:finmod}(ii)). No deduction that uses about the windows only that $T\mapsto\lambda_{\mathrm A}(T;a)$
has $\mathsf P$, and about the weights only that $a\in\mathsf C$, proves $\lambda_{\mathrm A}(T;a)\ge0$ for every $T$.
\end{theorem}
\begin{proof}
Let $n_0>e^{2T_1}$ and $a=\Lambda+c\,\mathbf 1_{n=n_0}$. For $T'\le T_1$ the contact $n_0$ is absent, so
$\lambda_{\mathrm A}(T';a)=\lambda_{\mathrm A}(T')$ on $J$. The proofs of Lemmas~\ref{lem:monoT} and \ref{lem:contA}
use only the nonnegativity and finiteness of the weights on each window, so $T\mapsto\lambda_{\mathrm A}(T;a)$ is
nonincreasing and continuous; hence it has $\mathsf P$, and $a\in\mathsf C$. By Theorem~\ref{thm:finmod}(ii), for any
$T>\frac12\log n_0$ and $c>c_*(T,n_0)$ one has $\lambda_{\mathrm A}(T;a)<0$. Apply Lemma~\ref{lem:transfer}.
\end{proof}

The amplified weights do not satisfy Riemann's functional equation. A falsifier that does is obtained from a
self-dual measure with a hole.

\begin{lemma}[self-dual measures with a hole]\label{lem:selfdual}
Let $N_0\ge2$, let $M$ be a non-square integer with $\sqrt M>4N_0+4$, and $R=\lfloor N_0\sqrt M\rfloor$. There is a
real, even, nonzero $c:\Z/M\Z\to\R$ with $\mathcal F_Mc=c$ for the unitary discrete Fourier transform and $c(r)=0$
for $|r|\le R$. The measure $\nu_c=\sum_{n\in\Z}c(n)\delta_{n/\sqrt M}$ satisfies $\widehat{\nu_c}=\nu_c$ (in the
normalization $\widehat\mu(\xi)=\int e^{-2\pi ix\xi}d\mu$), vanishes on $(-N_0,N_0)$, and $\sum_{r\bmod M}c(r)=0$.
\end{lemma}
\begin{proof}
Poisson summation over the classes $n\equiv r\pmod M$ gives $\widehat{\nu_c}=\sum_k(\mathcal F_Mc)(k)\delta_{k/\sqrt M}$.
The eigenvalue $1$ of $\mathcal F_M$ has multiplicity $\lfloor M/4\rfloor+1$ \cite{McClellanParks1972}; its
eigenvectors are even (odd vectors belong to $\pm i$), and since $\mathcal F_M$ is real symmetric on even vectors
they can be taken real. The conditions $c(0)=\dots=c(R)=0$ are $R+1$ linear conditions on even vectors, and
$R+1<\lfloor M/4\rfloor+1$ because $M/4>(N_0+1)\sqrt M$. The points $n/\sqrt M\in(-N_0,N_0)$ have $|n|\le R<M/2$, so
their residues are among the $|r|\le R$. Finally $\sum_rc(r)=\sqrt M(\mathcal F_Mc)(0)=\sqrt Mc(0)=0$.
\end{proof}

With $c$ as in Lemma~\ref{lem:selfdual} put $g(s)=\sum_{m>R}c(m)(m/\sqrt M)^{-s}$ and $f_t=\zeta+t\,g$ ($t\in\R$).
Riemann's proof of the functional equation, applied to the self-dual measure $\sum_n\delta_n+t\nu_c$, shows that
$\pi^{-s/2}\Gamma(\frac s2)f_t(s)$ is invariant under $s\mapsto1-s$ and has the poles of the corresponding function
of $\zeta$ \cite{Bochner1951,KahaneMandelbrojt1958}. For $\Rel s$ large, $-f_t'/f_t=\sum_\lambda b_t(\lambda)\lambda^{-s}$
over the multiplicative semigroup generated by $\N$ and $\{m/\sqrt M\}_{m>R}$. Let $\lambda_1=m_1/\sqrt M$ with $m_1>R$
minimal such that $c(m_1)\ne0$; $\lambda_1>N_0$ is irrational. Since $\log f_t=\log\zeta+\log(1+tg/\zeta)$ and every
frequency of $g/\zeta$ is at least $\lambda_1$, one has $b_t(\lambda)=\Lambda(\lambda)$ for $\lambda<\lambda_1$ (and
$0$ off $\N$), and $b_t(\lambda_1)=t\,c(m_1)\log\lambda_1$. We call \emph{window forms of $f_t$} the forms
\eqref{eq:form} and \eqref{eq:W} built with the weights $v_\lambda=2b_t(\lambda)/\sqrt\lambda$ at the points
$\log\lambda<2T$, and write $\lambda_{\mathrm A}(T;f_t)$ for the primitive bottom.

\begin{proposition}[the falsifier $f_t$]\label{prop:ft}
\emph{(i)} For $T\le\frac12\log\lambda_1$ the window forms of $f_t$ coincide with those of $\zeta$; in particular
$\lambda_{\mathrm A}(T;f_t)=\lamA$. \emph{(ii)} Let $a=\log\lambda_1$ and $0<\varepsilon<\frac12\log2$ smaller than
half the distance from $a$ to every other point $\log\lambda$ with $b_t(\lambda)\ne0$ for some $t$ and
$\lambda<\lambda_1^2$. For $|t|$ large the complete window form of $f_t$ at $T=\frac a2+\varepsilon$ takes a negative
value on $L^2(\IT)$.
\end{proposition}
\begin{proof}
(i) is the statement about $b_t$. (ii) The frequencies below $\lambda_1^2$ other than $\lambda_1$ are integers, the
numbers $m/\sqrt M$ with $m>m_1$, and products $n\,m/\sqrt M\ge2\lambda_1$; these form a discrete set independent of
$t$, so $\varepsilon$ exists. Let $0\le\phi_0\in C_c^\infty(-\varepsilon,\varepsilon)$, $\phi_0\ne0$, and
$\phi=\phi_0+\eta V_a\phi_0$. The complete form $W$ is invariant under translation ($|\widehat\phi|$, $A\bar B$ and
$c_\phi$ are), so $W(V_a\phi_0)=W(\phi_0)=:A_0$, which involves no contact since $\supp c_{\phi_0}\subset(-2\varepsilon,2\varepsilon)$.
In the cross term the only contact is $\lambda_1$, where $c_{V_a\phi_0,\phi_0}(a)=\|\phi_0\|^2$; hence
$W(\phi)=(1+\eta^2)A_0+2\eta B_t$ with $B_t=B_0-\frac12 v_{\lambda_1}\|\phi_0\|^2$, $B_0$ independent of $t$, and
$v_{\lambda_1}=2t\,c(m_1)\log\lambda_1/\sqrt{\lambda_1}$. If $A_0\le0$ a suitable $\eta$ gives a negative value as
soon as $B_t\ne0$; if $A_0>0$, $\min_\eta W(\phi)=A_0-B_t^2/A_0<0$ once $|B_t|>A_0$. After translating
$\phi$ to be centered, $\phi\in L^2(\IT)$ with $T=\frac a2+\varepsilon$.
\end{proof}

\begin{proposition}[the falsifier $f_t$ on the primitive subspace]\label{prop:ftprim}
In the situation of Proposition~\ref{prop:ft}(ii), for $|t|$ large the primitive window form of $f_t$ at
$T=\frac a2+\varepsilon$ takes a negative value on $\PT$, so $\lambda_{\mathrm A}(T;f_t)<0$. Hence, for every
$N_0$ and $M$ as in Lemma~\ref{lem:selfdual}, the falsifier $f_t$ has a primitive witness at a window
$T=\frac12\log\lambda_1+\varepsilon$, with $\frac12\log\lambda_1>\frac12\log N_0$.
\end{proposition}
\begin{proof}
Choose a real $\phi_0\in C_c^\infty(-\varepsilon,\varepsilon)$, $\phi_0\ne0$, with
$\int\phi_0(t)e^{t/2}dt=\int\phi_0(t)e^{-t/2}dt=0$; this is possible because the two functionals are linearly
independent on $C_c^\infty(-\varepsilon,\varepsilon)$ (a nontrivial combination $c_1e^{t/2}+c_2e^{-t/2}$ is real-analytic
and cannot annihilate $C_c^\infty$), so their common kernel has codimension two. Put $\phi=\phi_0+\eta V_a\phi_0$
as in the proof of Proposition~\ref{prop:ft}(ii); that proof uses the nonnegativity of $\phi_0$ nowhere: the cross
term $c_{V_a\phi_0,\phi_0}(a)=\|\phi_0\|^2$, the choice of $\varepsilon$ and the translation invariance of the complete
form hold for every $\phi_0$. Since $\int(V_a\phi_0)(t)e^{\pm t/2}dt=e^{\pm a/2}\int\phi_0e^{\pm t/2}=0$, both Tate
moments of $\phi$ vanish, and a translation by $s$ multiplies $A$ by $e^{s/2}$ and $B$ by $e^{-s/2}$, so they
still vanish after $\phi$ is centred in $\IT$, $T=\frac a2+\varepsilon$. Thus the centred $\phi$ lies in $\PT$,
where the primitive form \eqref{eq:form} and the complete form \eqref{eq:W} of $f_t$ agree, since
$2\Rel(A\bar B)=0$. The computation of Proposition~\ref{prop:ft}(ii) gives $W(\phi)=(1+\eta^2)A_0+2\eta B_t$ with
$A_0,B_0$ independent of $t$ and $B_t=B_0-\frac12v_{\lambda_1}\|\phi_0\|^2$, $v_{\lambda_1}=2t\,c(m_1)\log\lambda_1/\sqrt{\lambda_1}$,
$c(m_1)\ne0$; explicitly, $|B_t|>|A_0|$ as soon as
$|t|>\sqrt{\lambda_1}\,(|A_0|+|B_0|)/(|c(m_1)|\log\lambda_1\,\|\phi_0\|^2)$, and then the choice of $\eta$ made there
($\eta=-B_t/A_0$ if $A_0>0$; $\eta$ small with $\eta B_t<0$ if $A_0\le0$) gives $W(\phi)<0$. Hence
$\mathfrak a_T(\phi)<0$ for a nonzero $\phi\in\PT$; nothing about the zeros of $f_t$ is used, the argument being a
computation with the comb at the single contact $\lambda_1$.
\end{proof}

\begin{theorem}[bounded families and the functional equation]\label{thm:famFE}
Let $J\subset(0,T_1]$ with $T_1<\infty$. The conclusion of Theorem~\ref{thm:famblind}
holds for every class of deductions that uses, besides $\lambda|_J$, monotonicity and
continuity, only that the weights are the coefficients of $-f'/f$ for a Dirichlet series $f$ over a discrete
semigroup with the functional equation, pole and archimedean factor of $\zeta$.
\end{theorem}
\begin{proof}
Take $N_0\ge e^{2T_1}$, $M$ as in Lemma~\ref{lem:selfdual}, and $f_t$ with $|t|$ large. By
Proposition~\ref{prop:ft}(i), $\lambda_{\mathrm A}(\cdot;f_t)=\lambda_{\mathrm A}(\cdot)$ on $J$, since
$\frac12\log\lambda_1>\frac12\log N_0\ge T_1$; monotonicity and
continuity hold with the same proofs (the weights are finite on each window, of either sign, and
Lemma~\ref{lem:contA} uses only their total absolute mass); the properties of $f_t$ are those listed; and
$\lambda_{\mathrm A}(T;f_t)<0$ at $T=\frac12\log\lambda_1+\varepsilon$ by Proposition~\ref{prop:ftprim}. Apply
Lemma~\ref{lem:transfer}.
\end{proof}

\begin{remark}[status and scope]\label{rem:bar-fam}
Both theorems are unconditional. For Theorem~\ref{thm:famFE} the negativity of the \emph{complete} window form of
$f_t$ (Proposition~\ref{prop:ft}(ii)) would not suffice, since negativity on $L^2(\IT)$ does not imply negativity on
the codimension-two subspace $\PT$; Proposition~\ref{prop:ftprim} supplies the primitive witness directly, with a
primitive $\phi_0$. Neither theorem excludes arguments that use $\lambda$ on an unbounded set of windows. By
Lemma~\ref{lem:monoT}, \eqref{eq:target} is the single statement $\lim_{T\to\infty}\lamA\ge0$; and two-sided
asymptotic laws of the form $\lamA=L(T)(1+o(1))$ with an explicit $L>0$, which are strictly stronger than
\eqref{eq:target}, are neither shared by $f_t$ (which differs from $\zeta$ beyond $\frac12\log\lambda_1$) nor by the
falsifiers of the next subsection, which become negative. This class is open.
\end{remark}

\subsection{Patterns of window observables}\label{subsec:bar-pat}

The third class consists of arguments that read the sign from observable features of single windows: the value
of the bottom below a threshold, the parity of the minimizer, the location of the peak of $|\widehat{h^*}|^2$,
eigenvalue counts, or a law fitted to such data. The falsifiers are window forms obtained by adding finitely many
zeros off the critical line to $\xi(s)=\frac12s(s-1)\pi^{-s/2}\Gamma(\frac s2)\zeta(s)$.

\begin{lemma}[the quartet form]\label{lem:quartet}
Let $z=\delta+i\gamma$ with $\delta>0$, $\gamma>0$. For real $F\in L^2(\IT)$ put
$\Phi_c(F)=\int F(t)\cosh(\delta t)e^{i\gamma t}dt$ and $\Phi_s(F)=\int F(t)\sinh(\delta t)e^{i\gamma t}dt$. Then
the four points $\frac12\pm z$, $\frac12\pm\bar z$ contribute to $\sum_\rho\Phi_F(\rho)\Phi_F(1-\rho)$ exactly
\[
Q_z(F):=\int_{-2T}^{2T}\cF(u)\,4\cosh(\delta u)\cos(\gamma u)\,du=4|\Phi_c(F)|^2-4|\Phi_s(F)|^2 ,
\]
a bounded form on $L^2(\IT)$ of inertia $(2,2)$. (The second expression is Lemma~\ref{lem:wh-Q}; we record the first, which is the form used below.)
\end{lemma}
\begin{proof}
$\Phi_F(\frac12+w)\Phi_F(\frac12-w)=\iint F(t)F(s)e^{w(t-s)}dt\,ds=\int\cF(u)e^{wu}du$. The four points pair as
$\rho,1-\rho$, each pair counted twice, and $\sum_{w\in\{\pm z,\pm\bar z\}}e^{wu}=4\cosh(\delta u)\cos(\gamma u)$; this is
the first expression. For the second, with $P_\pm=\int Fe^{\pm\delta t}\cos\gamma t$ and $R_\pm=\int Fe^{\pm\delta t}\sin\gamma t$,
$\Rel[\Phi_F(\frac12+z)\Phi_F(\frac12-z)]=P_+P_-+R_+R_-=\frac14[(P_++P_-)^2+(R_++R_-)^2-(P_+-P_-)^2-(R_+-R_-)^2]$, and
$P_++P_-+i(R_++R_-)=2\Phi_c(F)$, $P_+-P_-+i(R_+-R_-)=2\Phi_s(F)$; the sum is $4\Rel[\cdot]$. Boundedness is
Cauchy--Schwarz. The four functionals $\Rel\Phi_c,\Iml\Phi_c,\Rel\Phi_s,\Iml\Phi_s$ are integrals against
$\cosh(\delta t)\cos\gamma t,\dots,\sinh(\delta t)\sin\gamma t$, which are linearly independent on $\IT$ (they are
independent combinations of the distinct exponentials $e^{(\pm\delta\pm i\gamma)t}$); so the inertia is $(2,2)$.
\end{proof}

\begin{definition}\label{def:quartetmodel}
For $z=\delta+i\gamma$ with $0<\delta<\frac12$, $\gamma>0$, and $k\in\N$ put
$\xi_{z,k}(s)=\xi(s)\,\big[((s-\frac12)^2-z^2)((s-\frac12)^2-\bar z^2)\big]^k$. Its \emph{primitive window form} is
$\mathfrak a^{z,k}_T:=\mathfrak a_T+kQ_z$ on $\PT$, with bottom $\lambda_{\mathrm A}(T;\xi_{z,k})$. Equivalently it is
the form \eqref{eq:form} with symbol $\sT+k\widehat{\beta_z}$, $\beta_z(u)=4\cosh(\delta u)\cos(\gamma u)\mathbf 1_{|u|<2T}$.
\end{definition}

\begin{proposition}[properties shared with $\zeta$]\label{prop:quartetmodel}
\emph{(i)} $\xi_{z,k}$ is entire of order $1$, real on $\R$, invariant under $s\mapsto1-s$; its zeros are those of
$\xi$ together with the four points $\frac12\pm z,\frac12\pm\bar z$, each with additional multiplicity $k$, all in
$0<\Rel s<1$. \emph{(ii)} For real $F\in\PT\cap C_c^\infty(\IT)$, the sum over the zeros of $\xi_{z,k}$, with
multiplicity, of $\Phi_F(\rho)\Phi_F(1-\rho)$ converges absolutely and equals $\mathfrak a^{z,k}_T(F)$: the window
form of $\xi_{z,k}$ satisfies the explicit formula with the archimedean term $m$, no polar term on $\PT$, and the
arithmetic side $\sum_nv_n\cF(\log n)-k\int\cF\beta_z$. \emph{(iii)} $\mathfrak a^{z,k}_T$ is closed and bounded below on
$\PT\cap\mathcal D_T$ with compact resolvent, and $T\mapsto\lambda_{\mathrm A}(T;\xi_{z,k})$ is nonincreasing and
continuous.
\end{proposition}
\begin{proof}
(i) The factor is a real polynomial invariant under $s\mapsto1-s$, with the stated roots. (ii) Lemma~\ref{lem:roughEF}
gives the sum over the zeros of $\xi$ (hypothesis (b) holds since $\Phi_F$ is rapidly decreasing on the strip), and
Lemma~\ref{lem:quartet} gives the quartet, counted $k$ times. (iii) $Q_z$ is bounded, so the form domain, closedness
and lower bound are those of $\mathfrak a_T$, and a sublevel set of $\mathfrak a^{z,k}_T$ lies in a sublevel set of
$\mathfrak a_T$, which is relatively compact (Lemma~\ref{lem:compact}). Extension by zero leaves $Q_z(F)$ unchanged,
which gives monotonicity as in Lemma~\ref{lem:monoT}. In the proof of Lemma~\ref{lem:contA} the added term
changes by $O(|r-1|)\|F\|^2$ under the dilation $U_r$ and by $O(\|G\|)\|F\|$ under the Tate correction, uniformly
on $[T/2,2T]$; the proof goes through.
\end{proof}

The models are negative at every window once the multiplicity is large: by Theorem~\ref{thm:wit-model}(iii) there is a real $F_0\in\PT\cap C_c^\infty(\IT)$ with $\Phi_c(F_0)=0$ and $\Phi_s(F_0)\ne0$, so that $\lambda_{\mathrm A}(T;\xi_{z,k})\le(\mathfrak a_T(F_0)-4k|\Phi_s(F_0)|^2)/\|F_0\|^2\to-\infty$ as $k\to\infty$.


\begin{remark}[numerical facts for one quartet]\label{rem:quartetnum}
Take $k=1$ and $\gamma=\gamma_1=14.1347\ldots$. The following values are Galerkin values (Legendre sections of
$\PT$, $128$-bit ball arithmetic for the linear algebra, Gauss quadrature not enclosed); a negative Galerkin
value is the Rayleigh quotient of an explicit polynomial, so only the quadrature is uncertified.
(a) $\lambda_{\mathrm A}(\log2;\xi_{z,1})\approx-1.52\cdot10^{-3}$ for $\delta=0.1$ and $\approx-2.85\cdot10^{-2}$ for
$\delta=0.3$, with two negative eigenvalues in each case, while $\lambda_{\mathrm A}(\log2)\approx9.83\cdot10^{-8}$;
also $\lambda_{\mathrm A}(\frac12\log5;\xi_{z,1})\approx-4.97\cdot10^{-3}$ for $\delta=0.1$.
(b) For $\delta=10^{-3}$ and $T\in\{\frac12\log3,\,0.61,\,0.67\}$, and for $\delta=10^{-6}$ at $T=\frac12\log5$, the
bottom of the model agrees with $\lamA$ to within $3\%$, the minimizers have the same parity, the peak of
$|\widehat{h^*}|^2$ is at the same frequency to two decimals, and the counts of near-critical eigenvalues of the
Wiener--Hopf pencil coincide. Below its first negative window the model therefore reproduces every observable
of the class; at $\delta=0.1$ it changes sign between $T=0.55$ and $T=0.69$, keeping parity and counts.
(c) As $\delta\to0$ the model tends to $\xi$ with a triple zero at $\frac12\pm i\gamma_1$, whose form
$\mathfrak a_T+4|\Fh(\gamma_1)|^2$ is $\ge\mathfrak a_T$; the sign change for small $\delta$ is the term
$-4\delta^2|\int F\,t\,e^{i\gamma_1t}dt|^2$ of $Q_z$ at a minimizer with $\widehat{h^*}(\gamma_1)\approx0$.
\end{remark}

\begin{corollary}[observable patterns]\label{cor:pattern}
Let $\mathsf O$ be any list of window observables that is computed from a window form alone (for example the
bottom up to a relative error, the parity of the minimizer, the peak frequency of $|\widehat{h^*}|^2$, eigenvalue
counts, or the parameters of a law fitted to these on a finite set of windows). If the quartet model $\xi_{z,1}$
has the same values of $\mathsf O$ as $\zeta$ on the windows used and a negative bottom at a window, then no
deduction from $\mathsf O$ proves \eqref{eq:target}. Numerically (Remark~\ref{rem:quartetnum}) this is the case for
the observables listed there; for the quartet models with $k$ large the negative bottom is proved
(Theorem~\ref{thm:wit-model}(iii)), and the shared properties of Proposition~\ref{prop:quartetmodel} are proved.
\end{corollary}
\begin{proof}
Lemma~\ref{lem:transfer}.
\end{proof}

\begin{remark}[the arithmetic-side counterpart and what is not excluded]\label{rem:bar-pat}
Theorem~\ref{thm:finmod} and Corollary~\ref{cor:void} are the arithmetic-side members of the same family: deletion
and amplification keep the zeros and the pole and change the comb; the quartet models keep the comb and change
the zeros. Every falsifier of this subsection differs from $\zeta$ in one respect that no listed observable sees:
the arithmetic side of $\xi_{z,k}$ has the absolutely continuous part $-k\beta_z$ besides the atoms at $\log n$, and
$\Lambda+c\mathbf 1_{n=n_0}$ breaks the functional equation. Arguments that use the atomic structure of the arithmetic
side together with the exact values $\Lambda(n)$, or two-sided laws on unbounded ranges (Remark~\ref{rem:bar-fam}),
are not excluded.
\end{remark}

\subsection{Local and finite-rank criteria}\label{subsec:bar-loc}

The fourth class consists of criteria computed from the symbol near its negative set
$U_T:=\{\tau:\sT(\tau)<0\}$, or from the compression of $\mathfrak a_T$ to a finite-dimensional subspace. Since $m$
is increasing with $m(0)=C_0<0$ and $m(\tau)\to\infty$, $U_T$ is a bounded open set containing $0$.

\begin{lemma}[compressions are necessary conditions]\label{lem:compress}
Let $X\subset\PT\cap\mathcal D_T$ be a subspace of finite dimension $d$ and $H_X$ the matrix of $\mathfrak a_T$ on an
orthonormal basis of $X$. Then $\lamA\le\lambda_{\min}(H_X)$ and $n_-(H_X)\le n_-(\mathfrak a_T|_{\PT})$.
\end{lemma}
\begin{proof}
The first is the Rayleigh principle; the second is the min--max principle (Cauchy interlacing for the compression).
\end{proof}

The exact converse, a Schur reduction onto such a compression with a coercive complement, is
Theorem~\ref{thm:wh-schur}; its margin is the full bottom. No condition on the symbol near $U_T$ suffices.

\begin{theorem}[no local criterion]\label{thm:nolocal}
Let $T>0$ and $\eta>0$, and put $\sigma^{(\eta)}:=\min(\sT,\eta)$. Then $\sigma^{(\eta)}$ is continuous and even, it
coincides with $\sT$ on the open neighbourhood $\{\sT<\eta\}$ of $\overline{U_T}$, it has the same negative set
$U_T$, and $\sigma^{(\eta)}>0$ off $\overline{U_T}$. There is $\eta_0(T)>0$ such that for $\eta<\eta_0(T)$ the form
$F\mapsto\frac1{2\pi}\int\sigma^{(\eta)}|\Fh|^2$ is negative somewhere on $\PT$. Consequently no criterion that depends
on $\sT$ only through its restriction to a set of the form $\{\sT<\eta\}$, $\eta>0$, decides the sign of $\lamA$ at
any window; in particular no criterion that is a function of the restrictions of $\sT$ to the components of $U_T$
and their neighbourhoods does.
\end{theorem}
\begin{proof}
The first assertions are immediate from the continuity of $\sT$. Fix $F\in\PT$, $\|F\|=1$. Its transform is a
nonzero entire function, so it does not vanish identically on the open set $U_T$, and
$-\kappa:=\frac1{2\pi}\int_{U_T}\sT|\Fh|^2<0$. Since $\sigma^{(\eta)}\le\eta$ off $U_T$,
$\frac1{2\pi}\int\sigma^{(\eta)}|\Fh|^2\le-\kappa+\eta$, which is negative for $\eta<\eta_0:=\kappa$. A criterion of the
stated kind returns the same answer for $\sT$ and $\sigma^{(\eta)}$; at the windows $T\le\frac12\log5$ the first form is
positive on $\PT$ (Theorem~\ref{thm:log5A}) and the second is not, and at every window a criterion that proves
positivity for $\sT$ proves it for $\sigma^{(\eta)}$ (Lemma~\ref{lem:transfer}).
\end{proof}

The sign is carried by the positive part of the symbol on all of $\R$ jointly with the negative set. What
positivity at a larger window does impose on the new contacts is an averaged inequality on $U_{T_0}$, never a
pointwise one. For $T_0<T$ put $S_{T_0,T}(\tau)=\sum_{2T_0\le\log n<2T}v_n\cos(\tau\log n)$, so that
$\sT=\sigma_{T_0}-S_{T_0,T}$, and let $\mathsf S_{U}=Q_T\mathbf 1_{U}(D)Q_T$ be the time--band limiting operator of a
bounded open set $U$ ($\mathbf 1_U(D)$ the Fourier multiplier), a compact positive contraction on $L^2(\IT)$. For
$0<\eta<1$ let $E_\eta(U)$ be the span of its eigenvectors with eigenvalue $\ge1-\eta$.

\begin{theorem}[what positivity imposes on new contacts]\label{thm:X}
Let $0<T_0<T$, $U_0=U_{T_0}$, and for $F\in\PT\cap\mathcal D_T$ put $\ell_T(F)=\frac1{2\pi}\int_{\R\setminus U_0}\sT|\Fh|^2$.
\emph{(i)} For every such $F$,
\[
\frac1{2\pi}\int_{U_0}S_{T_0,T}|\Fh|^2=\frac1{2\pi}\int_{U_0}\sigma_{T_0}|\Fh|^2+\ell_T(F)-\mathfrak a_T(F).
\]
\emph{(ii)} If $\lamA\ge0$, then for every $F\in\PT\cap\mathcal D_T$,
$\frac1{2\pi}\int_{U_0}S_{T_0,T}|\Fh|^2\le-\frac1{2\pi}\int_{U_0}|\sigma_{T_0}||\Fh|^2+\ell_T(F)$.
\emph{(iii)} For $F\in E_\eta(U_0)$, $\frac1{2\pi}\int_{\R\setminus U_0}|\Fh|^2\le\eta\|F\|^2$, and
$\ell_T(F)\ge-\frac1{2\pi}\int_{U_T\setminus U_0}|\sT||\Fh|^2$. The subspace $E_\eta(U_0)\cap\PT$ has codimension at most
$2$ in $E_\eta(U_0)$.
\emph{(iv)} The pointwise form of (ii) is false: $S_{T_0,T}(0)>0>\sigma_{T_0}(0)$ whenever
$\mathcal N_T\ne\mathcal N_{T_0}$.
\end{theorem}
\begin{proof}
(i) On $U_0$, $\sT=\sigma_{T_0}-S_{T_0,T}$; split $\mathfrak a_T(F)=\frac1{2\pi}\int_{U_0}\sT|\Fh|^2+\ell_T(F)$. (ii) is
(i) with $\mathfrak a_T(F)\ge0$ and $\sigma_{T_0}<0$ on $U_0$. (iii) $\langle\mathsf S_{U_0}F,F\rangle=\frac1{2\pi}\int_{U_0}|\Fh|^2$
and $\mathsf S_{U_0}\ge1-\eta$ on $E_\eta(U_0)$; for the lower bound, $\sT\ge0$ off $U_T$. $\PT$ has codimension $2$.
(iv) $S_{T_0,T}(0)=\sum v_n>0$ over the new contacts, and $\sigma_{T_0}(0)=C_0-\mathsf A_{T_0}<0$.
\end{proof}

\begin{remark}\label{rem:X}
By the Landau--Widom asymptotics \cite{LandauWidom1980}, $\dim E_\eta(U_0)=\frac T\pi|U_0|+O(\log T)$ as $T\to\infty$
with $U_0$ and $\eta$ fixed, so Theorem~\ref{thm:X}(ii) is a family of about $T|U_0|/\pi$ linear inequalities on the
values $\Lambda(n)$, $e^{2T_0}\le n<e^{2T}$, one for each frequency packet of width $\pi/T$ in $U_0$. Two remarks limit
its use. It is a necessary condition only: Lemma~\ref{lem:compress} and Theorem~\ref{thm:nolocal} show that the
converse requires control of the form off $E_\eta(U_0)$. And numerically (Galerkin, not certified) at the windows
$T\le\frac12\log7$ the space $E_{0.1}(U_0)\cap\PT$ is $\{0\}$, so the inequalities are void there; with $\eta=\frac12$ the
compression of $\mathfrak a_T$ to $E_\eta(U_0)\cap\PT$ has bottom between $6\cdot10^{-4}$ and $0.13$, many orders above
$\lamA$, and the quartet model with $\delta=0.1$ satisfies the same inequalities at $\log2$ and $\frac12\log5$ while
its bottom is negative.
\end{remark}

Finally, an evolution law in $T$ does not avoid the minimizer.

\begin{lemma}[the flow of a determinant]\label{lem:flowdet}
Let $T\mapsto A(T)$ be a $C^1$ family of real symmetric positive definite $d\times d$ matrices, with eigenvalues
$\mu_1(T)\le\dots\le\mu_d(T)$ and orthonormal eigenvectors $\varphi_j(T)$. Then
$\frac d{dT}\log\det A(T)=\sum_j\langle A'(T)\varphi_j,\varphi_j\rangle/\mu_j$, and where $\mu_1$ is simple,
$\mu_1'=\langle A'\varphi_1,\varphi_1\rangle$.
\end{lemma}
\begin{proof}
Jacobi's formula $(\log\det A)'=\operatorname{tr}(A^{-1}A')$, evaluated in the eigenbasis; the second statement is the
Hellmann--Feynman formula, obtained by differentiating $\langle A\varphi_1,\varphi_1\rangle$ with
$\langle\varphi_1',\varphi_1\rangle=0$.
\end{proof}

\begin{remark}[what this class does not exclude]\label{rem:bar-loc}
Applied to Galerkin sections of $\mathfrak a_T|_{\PT}$ (the windows being identified by a dilation),
Lemma~\ref{lem:flowdet} shows that the coefficient of any flow for $\log\det$ near a window where the bottom is small
is dominated by the term $\langle A'h^*,h^*\rangle/\lamA$, a functional of the minimizer. The same holds for the
quartet models: numerically, the difference between $\zeta$ and $\xi_{z,1}$ for small $\delta$ is the inequality
$\lamA<4\delta^2|\int h^*(t)\,t\,e^{i\gamma_1t}dt|^2$ between two functionals of $h^*$, and finite-rank compressions
built from the components of $U_T$ satisfy the criterion of Proposition~\ref{prop:wh-K} for $\delta$ up to about ten
times the threshold at which the full form changes sign. The class does not exclude criteria that use the positive
part of $\sT$ on all of $\R$ together with the negative set, such as the exact reductions of
Theorems~\ref{thm:coer-schur} and \ref{thm:wh-schur}; these carry the whole margin $\lamA$.
\end{remark}


\subsection{Arithmetic structure, orders and squares}\label{subsec:bar-arith}

The fifth class consists of arguments that use a structural property of the prime comb (an Euler product,
statistics in progressions, a relative bound, an operator identity, an order, a sum of squares) rather than the
vector $(\Lambda(n))_{n<e^{2T}}$ itself. We use the notation of \S\S\ref{subsec:finite} and \ref{sec:combs}.

\begin{proposition}[summable combs]\label{prop:summable}
Let $a\ge0$ with $\mathsf A(a):=\sum_{n\ge2}v_n(a)<\infty$. Then $\limsup_{T\to\infty}\lambda_{\mathrm A}(T;a)\le C_0-\mathsf A(a)<0$.
\end{proposition}

\begin{proof}
Take the tests $F_T$ of the proof of Proposition~\ref{prop:finsupp}. There $\mT(F_T)/\|F_T\|^2\to C_0$, and for each fixed $n$,
$c_{F_T}(\log n)/\|F_T\|^2\to1$ (the same computation with $p=\cos(\tau\log n)$). Since $|c_{F_T}|\le\|F_T\|^2$ and
$\sum_nv_n(a)<\infty$, dominated convergence gives $\sum_{n\in\NT}v_n(a)c_{F_T}(\log n)/\|F_T\|^2\to\mathsf A(a)$.
\end{proof}

\begin{corollary}[exact small values and Euler products]\label{cor:wtail}
\emph{(i)} Let $n_0\ge2$ and let $\mathcal C$ be a class of nonnegative arithmetic functions that contains $\Lambda$ and every $a$ with
$a(n)=\Lambda(n)$ for $n<n_0$ and $0\le a\le\Lambda$. Then $\Lambda\mathbf 1_{n<n_0}\in\mathcal C$ fails $\mathrm{WP}_T$ for all large $T$.
So no argument that uses about the primes only upper bounds of sieve type, even together with the exact values
$\Lambda(n)$, $n<n_0$, proves $\mathrm{WP}_T$ for all $T$.
\emph{(ii)} For a finite set $S$ of primes, $a_S:=\Lambda\cdot\mathbf 1\{n\text{ is a power of a prime in }S\}$ satisfies
$\sum_na_S(n)n^{-s}=-\zeta_S'/\zeta_S(s)$ with the Euler product $\zeta_S(s)=\prod_{p\in S}(1-p^{-s})^{-1}$, and
$\mathsf A(a_S)=\sum_{p\in S}\frac{2\log p}{\sqrt p-1}<\infty$. Hence $a_S$ fails $\mathrm{WP}_T$ for all large $T$, and no argument
that uses only the fact that the comb is the logarithmic derivative of an Euler product proves $\mathrm{WP}_T$ for all $T$.
The same holds for real powers $\zeta^{1+\varepsilon}$, whose comb is $(1+\varepsilon)\Lambda$, whenever $\varepsilon\kappa_T>\lamA$ in the
notation of Proposition~\ref{prop:wit-sens}.
\end{corollary}

\begin{proof}
(i) and the first part of (ii) are Proposition~\ref{prop:summable}; $\sum_{k\ge1}2\log p\cdot p^{-k/2}=2\log p/(\sqrt p-1)$. The last
statement is Proposition~\ref{prop:wit-sens}.
\end{proof}

Proposition~\ref{prop:finsupp} already covers the truncations; Corollary~\ref{cor:wtail}(i) adds that keeping the exact small
values does not help, and (ii) that multiplicativity alone does not either. The failing members preserve the Euler product,
the logarithm and M\"obius inversion; what they destroy is the identity $\sum_{d\mid n}a(d)=\log n$ for $n<e^{2T}$, which
determines $a=\Lambda$ on $\NT$.

\begin{proposition}[progressions]\label{prop:prog}
Let $\mathcal C$ be a class of nonnegative arithmetic functions containing $\Lambda$, defined by conditions on the sums
$\psi_a(x;q,c)=\sum_{n\le x,\,n\equiv c\,(q)}a(n)$ each of which is preserved when every $\psi_a(x;q,c)$ is changed by at most a
constant independent of $x$, $q$ and $c$. Then $\mathcal C$ is closed under deletion of finitely many terms and under
amplification of one term, so Corollary~\ref{cor:void} applies: no argument that uses only membership of $\Lambda$ in $\mathcal C$
proves $\mathrm{WP}_T$ at any $T>T_0(0)$.
\end{proposition}

\begin{proof}
Deleting $\Lambda(n)$ for $n<N$ changes every $\psi_a(x;q,c)$ by at most $\psi(N)$; amplifying one term by $c$ changes it by at most $c$.
\end{proof}

The hypothesis covers, with constants depending on the modification, asymptotic formulas with errors $\gg1$ (Siegel--Walfisz),
averaged forms $\sum_{q\le Q}\max_c|E_a(x;q,c)|\ll R(x)$ with $Q=o(R(x))$ (Bombieri--Vinogradov), mean squares
$\sum_{q\le Q}\sum_c|E_a(x;q,c)|^2\ll xQ\log x$ (Barban--Davenport--Halberstam, by Cauchy--Schwarz), and Brun--Titchmarsh in the
asymptotic range. In the integer-lattice matrices of Proposition~\ref{prop:AM}, such statements enter only through
$\mathrm{Near}_M$ correlated with residue classes of the test vector. On twisted vectors $s_m=\chi(m)u_m$, $\mathrm{Div}_M$ becomes the
Dirichlet convolution by $\chi(n)v_n$, but $\mathrm{Arch}_M$ stays the archimedean part of $\zeta$ and $\mathrm{Near}_M$ carries the
phases $\chi(m)\overline{\chi(k)}$ with $m\ne kn$. So the window form of $L(s,\chi)$ does not appear exactly in $\mathbb A_M$.

\begin{proposition}[relative bounds]\label{prop:relative}
Let $0\le\theta<1$, $\kappa\ge|C_0+4|$, $C\ge0$, and $T$ with $(1-\theta)\lambda^\sharp(T)<\theta\kappa+C$; by
Proposition~\ref{prop:sharp}(iii) this holds for all large $T$. Then
$\mathfrak q':=\theta(\mathfrak a^\sharp+\kappa\|\cdot\|^2)+C\|\cdot\|^2$ satisfies $|\mathfrak q'|\le\theta(\mathfrak a^\sharp+\kappa\|\cdot\|^2)+C\|\cdot\|^2$ on
$\PT\cap\mathcal D_T$, and $\mathfrak a^\sharp-\mathfrak q'$ is negative somewhere on $\PT$. Hence no deduction of the form ``$\mathcal F_T$, or
any form substituted for it, satisfies a relative bound with constants $(\theta,\kappa,C)$, hence $\mathfrak a^\sharp-\mathcal F_T\ge0$'' is
valid at such $T$.
\end{proposition}

\begin{proof}
$\mathfrak a^\sharp+\kappa\|\cdot\|^2\ge0$ by Proposition~\ref{prop:sharp}(i), so $\mathfrak q'\ge0$ and the bound holds. The bottom of
$\mathfrak a^\sharp-\mathfrak q'=(1-\theta)\mathfrak a^\sharp-(\theta\kappa+C)\|\cdot\|^2$ is $(1-\theta)\lambda^\sharp(T)-\theta\kappa-C<0$. Apply
Lemma~\ref{lem:transfer}.
\end{proof}

This is the form-boundedness (KLMN) class in polar-free coordinates. It includes bounds for the Chebyshev fluctuation transported
through the M\"obius resolvent, since the conclusion depends only on the constants. Even $\theta=0$ fails: the archimedean
side alone is negative.

\begin{proposition}[operator identities are linear in the autocorrelation]\label{prop:opcalc}
Let $A_i,B_i$ ($i\le r$) be finite complex Borel measures on $\R$, acting by convolution. For $F\in L^2(\IT)$,
\[
\sum_{i\le r}\Rel\langle A_i*F,B_i*F\rangle=\int_{(-2T,2T)}\Rel\cF\,d\nu,\qquad \nu:=\Rel\sum_i(\overline{A_i}*\check B_i),
\]
where $\check B(E)=B(-E)$ and $\overline{A}$ is the conjugate measure; in particular every identity among such operators is an identity
of measures on $(-2T,2T)$, and for $A_i=B_i$ the measure $\overline{A_i}*\check A_i$ is positive definite.
\end{proposition}

\begin{proof}
$\langle A*F,B*F\rangle=\frac1{2\pi}\int\widehat A\,\overline{\widehat B}\,|\Fh|^2d\tau$ by Plancherel, and $\cF(u)=\frac1{2\pi}\int|\Fh|^2e^{-i\tau u}d\tau$
vanishes for $|u|\ge2T$; the real part symmetrizes.
\end{proof}

So an argument written in the language of shift operators (Euler products, M\"obius inversion, logarithmic derivatives,
Gram identities among $U_\mu F$, $U_{\log}F$) uses two tools only: linear identities between measures on the window, which carry
the whole margin, and the positivity $\|A*F\|^2\ge0$, used through Cauchy--Schwarz or as a sum of squares. The second use is
treated next.

\begin{theorem}[arithmetic squares]\label{thm:sos}
Call an \emph{arithmetic square decomposition at $T$} a finite family of Dirichlet polynomials $D_i(\tau)=\sum_{n\le N}x^{(i)}_nn^{i\tau}$
and a constant $\mu_0\in\R$ with
\[
\sum_i|D_i(\tau)|^2=\mu_0-\sum_{n\in\NT}v_n\cos(\tau\log n)\quad(\tau\in\R),\qquad \mT-\mu_0\|\cdot\|^2\ge0\ \text{on }\PT .
\]
Such a decomposition proves $\lamA\ge0$, since then $\mathfrak a_T(F)=\sum_i\frac1{2\pi}\int|D_i|^2|\Fh|^2+\mT(F)-\mu_0\|F\|^2$. No such
decomposition exists at any $T\ge\frac12\log3$.
\end{theorem}

\begin{proof}
The first condition and Lemma~\ref{lem:mbounds}(iii) give $\mu_0\ge\sup_\tau\sum_{n\in\NT}v_n\cos(\tau\log n)=\Aconst$; the second gives
$\laminfP(T)\ge\mu_0$. For $T\ge\frac12\log3$, $2\in\NT$, so $\Aconst\ge v_2=0.9802\ldots$, while
$\laminfP(T)\le\laminfP(\frac12\log3)\le C_0+\varphi((3\pi/\log3)^2)=1.291-v_2=0.311$ by Proposition~\ref{prop:arch}(ii) and the value in
Corollary~\ref{cor:mono}.
\end{proof}

Identity of almost periodic symbols is the same as equality of the atomic parts of the corresponding measures. A
decomposition in which the comb is paid by squares of Dirichlet polynomials and the archimedean form pays a constant
therefore does not exist once a contact is present (and $T\ge\frac12\log3$; the interval $(\frac12\log2,\frac12\log3)$ is not
decided here). If a free positive definite part is allowed in addition, the decompositions are exactly the certificates of
value $\ge0$ of Definition~\ref{def:PD}, whose content is the bottom itself (Theorem~\ref{thm:D1}).

Let $\mathcal K^*_T$ be the set of finite measures $\nu$ on $[0,2T)$ with $\int\Rel\cF\,d\nu\ge0$ for every $F\in\PT$; it
defines an order on such measures. Let $\mathsf Z:=\sum_{1\le n<e^{2T}}n^{-1/2}\delta_{\log n}$, $\mathsf M:=\delta_0+e^{u/2}du$ on $[0,2T)$, and
compute convolutions modulo measures supported in $[2T,\infty)$. Then $\mathsf M^{-1}=\delta_0-e^{-u/2}du$, the M\"obius comb
$\mathsf Z^{-1}=\sum_{n<e^{2T}}\mu(n)n^{-1/2}\delta_{\log n}$ satisfies $\mathsf Z^{-1}=\mathsf M^{-1}-\rho*\mathsf Z^{-1}$ with
$\rho:=\mathsf M^{-1}*(\mathsf Z-\mathsf M)$, which determines $\mathsf Z^{-1}$ recursively since $\rho$ has no atom at $0$.

\begin{proposition}[the M\"obius resolvent is not monotone]\label{prop:cone}
Let $T>\frac12\log2$. \emph{(i)} $\delta_0\in\mathcal K^*_T$ and $e^{-u/2}du\in\mathcal K^*_T$. \emph{(ii)} For $0<a<2T$, $\pm\delta_a\notin\mathcal K^*_T$.
\emph{(iii)} $\rho=-e^{-u/2}du$ on $(0,\log2)$, $\rho$ has an atom of mass $2^{-1/2}$ at $\log2$, and $\pm\rho\notin\mathcal K^*_T$. Hence neither
$\nu\mapsto\rho*\nu$ nor $\nu\mapsto-\rho*\nu$ maps $\mathcal K^*_T$ into itself, and the iteration for $\mathsf Z^{-1}$ is not monotone for this
order.
\end{proposition}

\begin{proof}
(i) $c_F(0)=\|F\|^2$, and $\int_0^{2T}e^{-u/2}\Rel\cF=\frac12\mathfrak e(F)\ge0$ (Lemma~\ref{lem:polarelim}). (ii) Take a real unit $\varphi$ with $A(\varphi)=B(\varphi)=0$,
supported in an interval of length $w<\min(a,2T-a)$, and $F_\pm=\varphi\pm V_a\varphi$ placed inside $\IT$; both lie in $\PT$
(Lemma~\ref{lem:transl}(i)) and $c_{F_\pm}(a)=\pm1$. (iii) On $(0,\log2)$, $\mathsf Z-\mathsf M=-e^{u/2}du$ and
$\rho=-e^{u/2}+\int_0^ue^{-(u-v)/2}e^{v/2}dv=-e^{-u/2}$; the atom of $\mathsf Z$ at $\log2$ passes to $\rho$, and $e^{-u/2}*(\mathsf Z-\mathsf M)$ has no
atoms and a locally bounded density. For $F\in\PT$ with support length below $\frac12\log2$, $\int\Rel\cF\,d\rho=-\frac12\mathfrak e(F)<0$. For
$F=\varphi+V_{\log2}\varphi$ as in (ii) with $w\to0$, $\int\Rel\cF\,d(-\rho)=-2^{-1/2}+O(w)<0$. Finally $\pm\rho*\delta_0=\pm\rho$.
\end{proof}

The obstruction is local: the first integer atom against the continuous density. A monotone iteration would give one-sided
bounds for M\"obius sums at resolution $O(1)$ correlated with every primitive autocorrelation.

\begin{remark}[exterior certificates far from the window]\label{rem:farext}
If $(\nu,\beta)$ is a certificate of value $c$ at $T$ with $\supp\nu\subset\{|u|\ge U\}$, $U\ge2T$, then weak duality with
$\mu=|\widehat G|^2d\tau/2\pi$ gives $c\|G\|^2\le\frac1{2\pi}\int\sT|\widehat G|^2$ for every $G\in L^2(I_{U/2})$ with $\int\gh j|\widehat G|^2=0$, $j=0,1$.
This is the form at the window $U/2$ of the comb truncated to $\NT$, on the subset cut out by the two window-$T$ moment conditions.
The dilated tests of Proposition~\ref{prop:finsupp} drive this form to $C_0-\Aconst<0$ but do not satisfy the two conditions,
and we have not shown that the constrained value is negative. What is proved is that exterior mass far from the window can only
certify what the truncated comb certifies at the larger window.
\end{remark}

\begin{remark}[the functional equation on a family of symbols]\label{rem:fe}
For $\sigma>1$ let $D_\sigma(\tau):=2\Rel\frac{\xi'}{\xi}(\sigma+i\tau)=m_\sigma(\tau)+p_\sigma(\tau)-\sum_n2\Lambda(n)n^{-\sigma}\cos(\tau\log n)$, with
$m_\sigma=\Rel\psi(\frac{\sigma+i\tau}2)-\log\pi$ and $p_\sigma=2\Rel[\frac1{\sigma+i\tau}+\frac1{\sigma-1+i\tau}]$. At $\sigma=\frac12$ these become $m$, $0$
and $v_n$, the formal symbol of $\mathsf B$. For $\sigma>1$, $D_\sigma\ge0$ is classical (pair the factors of the Hadamard product of $\xi$
under $\rho\mapsto1-\bar\rho$; \cite[Ch.~12]{Davenport2000}). The functional equation and reality give
$\frac{\xi'}{\xi}(1-s)=-\frac{\xi'}{\xi}(s)$ and hence $D_{1-\sigma}=-D_\sigma$ exactly. So every combination symmetric under $\sigma\mapsto1-\sigma$
vanishes identically, and no positivity can be obtained by symmetrizing this family. At $\sigma=\frac12$ the antisymmetry forces
$D_{1/2}=0$ off the ordinates of zeros, the pointwise form of Proposition~\ref{prop:boundary}. A transport from $\sigma>1$ to
$\sigma=\frac12$ that uses only the location of the zeros is an argument from membership in $\mathcal C_{\mathrm{zeros}}$ and is
excluded by Corollary~\ref{cor:void}.
\end{remark}

\begin{remark}[what this class does not exclude]\label{rem:bar-arith}
Propositions~\ref{prop:summable}--\ref{prop:cone}, Theorem~\ref{thm:sos} and Proposition~\ref{prop:nearfree} exclude, respectively:
arguments from summable or Euler-product combs; from progression statistics; from relative form bounds; from monotone M\"obius
iterations; from arithmetic squares; and from tests that do not see near-coincidences of integers. They do not exclude an exact
inequality $\mathcal F_T\le\mathfrak a^\sharp$ on $\PT$ proved from the vector $(\Lambda(n))_{n<e^{2T}}$ jointly with the archimedean form,
for instance the positive definiteness of $\mathsf B$ (Theorem~\ref{thm:oneprofile}) on test vectors that see $\mathrm{Near}_M$.
\end{remark}


\begin{remark}[a Davenport--Heilbronn falsifier, proved]\label{rem:FDH}
Let $F_{\mathrm{DH}}=\tfrac{1-i\kappa}2L(s,\chi)+\tfrac{1+i\kappa}2L(s,\bar\chi)$, $\chi$ the odd character modulo $5$ with
$\chi(2)=i$, $\kappa=0.2840790\ldots$, so that $a_n\in\{1,\kappa,-\kappa,-1,0\}$ is real and periodic modulo $5$,
$\Lambda_{\mathrm{DH}}(s)=(5/\pi)^{s/2}\Gamma(\tfrac{s+1}2)F_{\mathrm{DH}}(s)=\Lambda_{\mathrm{DH}}(1-s)$ (verified to
$2.5\cdot10^{-30}$), with no Euler product and no pole, and with a zero at $\rho=0.808517\ldots+85.699348\ldots i$ off the critical line
(\cite{DavenportHeilbronn1936} for the existence of such zeros; the value is Spira's \cite{Spira1994}, verified here to $10^{-35}$). Its window form (archimedean constant $\log(5/\pi)+\mathrm{Re}\,\psi(\tfrac34)$, L\'evy kernel of
$\Gamma_\R(s+1)$, comb $2b_n/\sqrt n$ with $-F'/F=\sum b_nn^{-s}$, same $\PT$) is negative on $\PT$ for every
$T\ge\frac95$: an explicit odd witness of degree $\le259$ in Legendre polynomials with dyadic coefficients gives the Rayleigh
quotient $[-7.05197470947908\cdot10^{-9}\pm4\cdot10^{-25}]$ at $T=\frac95$ (ball arithmetic, 7000 bits; full sector
$-2.3799\cdot10^{-7}$; $T=2$ odd sector $-0.75194$), and negativity propagates to every $T\ge\frac95$ by Lemma~\ref{lem:monoT}, whose proof uses only the structure of the form and applies verbatim to the DH form. This reproduces, in sector,
window and order of magnitude, the value $-6.169\cdot10^{-9}$ reported by Ziad \cite{ZiadDH}; the $T=2$ witness concentrates $|\hat F|$
near $\xi\approx84.5$ (numerical observation), next to the off-line zero, which is why low-degree test families do not detect it. The same pipeline applied to
$L(s,\chi_{-4})$ shows no negative eigenvalue above $10^{-13}$. Read as a falsifier in the sense of Lemma~\ref{lem:transfer}: a
method using only the functional equation, a gamma factor and a Dirichlet series with bounded real coefficients, and not the
multiplicativity of the coefficients, cannot prove window positivity for any $T\ge\frac95$ (phase-200 memos \texttt{DH\_TRANSFER\_001}, \texttt{AUDIT\_200\_CERTDH\_001},
\texttt{CERT\_DH\_001}). No lower bound for the threshold $T_{\mathrm{DH}}$ is certified (numerically in $[1.6,1.8]$).
\end{remark}

\section{The shape of a proof}\label{sec:shape}

This section collects the barriers of \S\ref{sec:negative}, \S\ref{sec:barriers} and \S\ref{sec:wh} into one statement about
the form that any valid proof of \eqref{eq:target} must have. Nothing here is a new route and nothing moves $\lamA$. Every
clause is either an instance of Lemma~\ref{lem:transfer} with an explicit falsifier, or a constraint on the precision or the
exactness of a proof; every clause excludes a kind of deduction, never the conclusion. Each statement carries one of the labels
\emph{proved}, \emph{cited}, \emph{conditional}, \emph{argued}, \emph{open} or \emph{numerical}; a clause that is only argued or
only numerical is stated as such and is not used in the theorem.

\subsection{Form objects, modifications and classes}

\begin{definition}[form objects]\label{def:shape-obj}
A \emph{comb} is a function $a:\N_{\ge2}\to\R_{\ge0}$, finite on every window, with weights $v_n(a)=2a(n)/\sqrt n$; its window
form $\mathfrak a^{(a)}_T$, its bottom $\lambda_{\mathrm A}(T;a)$ on $\PT$ and the statement $\mathrm{WP}_T(a)$ are those of
\S\ref{subsec:finite}, and $a=\Lambda$ gives $\lamA$ and \eqref{eq:target}. A \emph{form object} is a family $q=(q_T)_{T>0}$ of
closed semibounded forms, $q_T$ defined on $\PT\cap\mathcal D_T$ or on a closed subspace of $L^2(\IT)$, with bottoms
$\lambda(T;q)$. Every comb is a form object through $q_T=\mathfrak a^{(a)}_T|_{\PT}$, and $q_0:=(\mathfrak a_T|_{\PT})_T$ is the
object of $\zeta$. The zero-side objects of \S\ref{subsec:bar-fam} and \S\ref{subsec:bar-pat} are form objects that are not combs:
the quartet models $\mathfrak a_T+kQ_z$ of Definition~\ref{def:quartetmodel}, and the deformations $f_t$ of
Proposition~\ref{prop:ft}, whose weights have both signs.
\end{definition}

\begin{definition}[modifications]\label{def:shape-ops}
The following maps on form objects are used below.
\begin{enumerate}
\item[$\mathrm O_1$] finite deletion, $a\mapsto a\mathbf 1_{n\ge N}$;
\item[$\mathrm O_2$] one-term amplification, $a\mapsto a+c\mathbf 1_{n=n_0}$;
\item[$\mathrm O_3$] tail deletion, $a\mapsto a\mathbf 1_{n<n_0}$;
\item[$\mathrm O_4$] uniform rescaling, $a\mapsto ca$ with $c>0$, and its small form $\mathrm O_4^\varepsilon$ with $c\in(1-\varepsilon,1+\varepsilon)$;
\item[$\mathrm O_5$] a change of one weight by $\varepsilon$;
\item[$\mathrm O_6$] Euler-preserving prime deletion, $a\mapsto a_P:=\Lambda\mathbf 1\{n\text{ is a power of a prime }p\ge P\}$;
\item[$\mathrm O_7$] multiplication by quartets, $\xi\mapsto\xi_{z,k}$ (Definition~\ref{def:quartetmodel});
\item[$\mathrm O_8$] self-dual deformation, $\zeta\mapsto f_t=\zeta+tg$ (Lemma~\ref{lem:selfdual});
\item[$\mathrm O_9$] a functional-equation-preserving twist $\xi\mapsto\xi R$ by a rational function (Proposition~\ref{prop:shape-twist});
\item[$\mathrm O_{10}$] Euler removal: passage to an Epstein zeta function of class number greater than one (Remark~\ref{rem:twohalves});
\item[$\mathrm O_{11}$] Beurling repositioning: $\N$ replaced by a generalized-integer system with a positive prime measure;
\item[$\mathrm O_{12}$] field change, $\xi\mapsto\xi_K$ for a number field $K$;
\item[$\mathrm O_{13}$] restriction of the family $T\mapsto\lambda(T;q)$ to a bounded set of windows $J\subset(0,T_1]$;
\item[$\mathrm O_{14}$] symbol truncation, $\sT\mapsto\sigma^{(\eta)}=\min(\sT,\eta)$ (Theorem~\ref{thm:nolocal});
\item[$\mathrm O_{15}$] Poisson randomization of the positions $\{\log n\}$ at the intensity of the primes (Theorem~\ref{thm:poisson}).
\end{enumerate}
Two further modifications produce no falsifier and enter only the structural part: the reflection $\sigma\mapsto1-\sigma$ of the
family $D_\sigma$ of Remark~\ref{rem:fe}, and a bounded perturbation $q_T\mapsto q_T+\mathfrak b$ with $\|\mathfrak b\|\le\lamA$
(Lemmas~\ref{lem:wh-T}, \ref{lem:wh-Tprime}).
\end{definition}

\begin{definition}[deduction classes]\label{def:shape-class}
A \emph{deduction class} $\mathcal A$ is specified by a set $\mathsf P$ of properties of form objects; a deduction in $\mathcal A$ is a
valid chain ``$q$ has $\mathsf P$, hence $\lambda(T;q)\ge0$'', for every $T$ or at a fixed $T$. The class is \emph{closed under} a
modification $\mathrm O$, for the parameters named in the corresponding lemma, if $q$ having $\mathsf P$ implies that $\mathrm O(q)$
has $\mathsf P$. The classes already named in this paper are: identities linear in the autocorrelation with Bochner or Cauchy--Schwarz
positivity (Proposition~\ref{prop:opcalc}); norm, moment, determinant, local, family, pattern and coherence functionals
(\S\ref{sec:wh}, \S\ref{sec:barriers}, Proposition~\ref{prop:wit-sens}); the zero-information class $\mathcal C_{\mathrm{zeros}}$ and
the classes $\mathcal C_{\mathrm{PNT}}(\varepsilon)$, $\mathcal C_{\mathrm{BT}}$ of Corollary~\ref{cor:void}; the progression classes of
Proposition~\ref{prop:prog}; the dual, lattice, flat-top and square certificates (Definition~\ref{def:PD}, Theorems~\ref{thm:Z},
\ref{thm:sos}); the symmetrizations of $D_\sigma$ (Remark~\ref{rem:fe}); and the axioms A1--A6 and the field-uniform properties of
\cite{Paper46}.
\end{definition}

\subsection{Transfer lemmas}

Each lemma names the modification, the class it closes, the falsifier, the windows at which the falsifier is negative, and the label.
The first group is proved in this paper.

\begin{lemma}[$\mathrm O_1$, finite deletion; proved]\label{lem:shape-del}
If $\mathsf P$ is preserved by $a\mapsto a\mathbf 1_{n\ge N}$, as it is for $\mathcal C_{\mathrm{PNT}}(\varepsilon)$,
$\mathcal C_{\mathrm{zeros}}$ and $\mathcal C_{\mathrm{BT}}$, then $\Lambda\mathbf 1_{n\ge N}$ with $N\ge\max(5,e^{2T})$ has $\mathsf P$
and fails $\mathrm{WP}_T$ for every $T\in[0.60,\frac12\log N]$, and for every $T>T_0(0)=0.74921$ (Theorem~\ref{thm:finmod}(i),
Corollary~\ref{cor:void}, Proposition~\ref{prop:arch}(ii)).
\end{lemma}

\begin{lemma}[$\mathrm O_2$, amplification; proved]\label{lem:shape-amp}
If $\mathsf P$ is preserved by $a\mapsto a+c\mathbf 1_{n=2}$, then $\Lambda+c\mathbf 1_{n=2}$ with $c>c_*(T,2)$ has $\mathsf P$ and fails
$\mathrm{WP}_T$ at every $T>\frac12\log2$, including the windows $(\frac12\log2,\frac12\log5]$ at which $\mathrm{WP}_T(\Lambda)$ holds
(Theorem~\ref{thm:finmod}(ii)).
\end{lemma}

\begin{lemma}[$\mathrm O_3$, tail deletion; proved]\label{lem:shape-tail}
If $\mathsf P$ holds for every $a$ with $a=\Lambda$ on $n<n_0$ and $0\le a\le\Lambda$ (the exact small values together with every upper-bound
sieve inequality), then $\Lambda\mathbf 1_{n<n_0}$ has $\mathsf P$ and fails $\mathrm{WP}_T$ for all large $T$; the dilated tests of
Proposition~\ref{prop:finsupp} have limit value $C_0-\sum_{n<n_0}v_n<0$ (Corollary~\ref{cor:wtail}(i)). The first failing window is not located.
\end{lemma}

\begin{lemma}[$\mathrm O_4$, rescaling; proved for $c$ large and small, conditional for $\mathrm O_4^\varepsilon$]\label{lem:shape-scale}
If $\mathsf P$ is preserved by $a\mapsto ca$ for every $c>0$ (an Euler product $\zeta^c$, the support, the relative sizes of the weights),
then $c\Lambda$ fails $\mathrm{WP}_T$ at every $T>\frac12\log2$ for $c$ large, and at every $T\ge0.60$ for $c$ small. If $\mathsf P$ is
preserved only by $c\in(1-\varepsilon,1+\varepsilon)$, then $(1+\operatorname{sgn}(\kappa_T)\varepsilon)\Lambda$ fails $\mathrm{WP}_T$ as soon as
$\varepsilon>\lamA/|\kappa_T|$, \emph{provided} $\kappa_T\ne0$, where $\kappa_T=\sum_nv_nc_h(\log n)$ for a real unit minimizer $h$
(Proposition~\ref{prop:wit-sens}); whether $\kappa_T\ne0$ at some window is not known, and $\kappa_{\log2}\approx0.09$ is a numerical value
(Remark~\ref{rem:wit-sens}).
\end{lemma}
\begin{proof}
For $c$ large take $F_+=\phi+\phi(\cdot-\log2)$ of Lemma~\ref{lem:cpos} with $u_0=\log2$ and $\ell<\log\frac32$. The autocorrelation
$c_{F_+}$ is supported in $\{|u|<\ell\}\cup\{|u\mp\log2|<\ell\}$, which contains no $\log n$ with $n\ge3$, and $c_{F_+}(\log2)=\|\phi\|^2>0$.
Hence $\mathfrak a^{(c\Lambda)}_T(F_+)=\mT(F_+)-cv_2\|\phi\|^2\to-\infty$ as $c\to\infty$. For $c$ small, the function $\mu_T$ of
Lemma~\ref{lem:degen} is finite and concave on $\R^{\NT}$, hence continuous, and $\mu_T(0)=\laminfP(T)<0$ for $T\ge0.60$
(Corollary~\ref{cor:Tarch}); so $\lambda_{\mathrm A}(T;c\Lambda)=\mu_T(cv)<0$ for $c$ small. The small form is Proposition~\ref{prop:wit-sens}.
\end{proof}

\begin{lemma}[$\mathrm O_5$, one weight; conditional]\label{lem:shape-weight}
Let $F_\star$ be a real unit minimizer at $T$ and $n_0\in\NT$ with $c_{F_\star}(\log n_0)\ne0$. If $\mathsf P$ is preserved by changing the
weight at $n_0$ by $\varepsilon$, then $\Lambda+s\varepsilon\mathbf 1_{n=n_0}$ with $s=\operatorname{sgn}c_{F_\star}(\log n_0)$ fails $\mathrm{WP}_T$
for every admissible $\varepsilon>\varepsilon_*(T,n_0)=\sqrt{n_0}\,\lamA/(2|c_{F_\star}(\log n_0)|)$, where for $s=-1$ admissibility means
$\varepsilon\le\Lambda(n_0)$ (Remark~\ref{rem:resweights}). If every real minimizer has $c_{F_\star}=0$ on $\log\NT$, the clause is empty
and Lemma~\ref{lem:degen} describes that case; whether this happens at any window is not known. The numerical scale $10^{-5}$ at $T=\log2$
(Remark~\ref{rem:wit-sens}) is not certified.
\end{lemma}

\begin{proposition}[$\mathrm O_6$, Euler-preserving prime deletion; proved]\label{prop:shape-eulerdel}
Let $P\ge3$ and $a_P:=\Lambda\mathbf 1\{n\text{ is a power of a prime }p\ge P\}$. Then $a_P\ge0$ has the generating function
$\zeta_{a_P}(s)=\prod_{p\ge P}(1-p^{-s})^{-1}=\zeta(s)\prod_{p<P}(1-p^{-s})$, an Euler product over $\N$ with nonnegative integral
Dirichlet coefficients; $\zeta_{a_P}$ has exactly the zeros of $\zeta$ in $0<\Rel s<1$ and the pole at $s=1$; $a_P(p^k)=\log p$ on its
support; $a_P\le\Lambda$; and $a_P$ fails $\mathrm{WP}_T$ for every $T\in[0.60,\frac12\log P]$. Of the properties listed in
Definition~\ref{def:shape-class}, $a_P$ breaks only the functional equation.
\end{proposition}
\begin{proof}
$\sum_na_P(n)n^{-s}/\log n=\sum_{p\ge P}\sum_k p^{-ks}/k=-\sum_{p\ge P}\log(1-p^{-s})$, which gives the Euler product; the coefficients of
$\prod_{p\ge P}(1-p^{-s})^{-1}$ count factorizations, so they are nonnegative integers. The finite product $\prod_{p<P}(1-p^{-s})$ is
entire and vanishes only on $\Rel s=0$. The comb of $a_P$ is empty on $\NT$ when $e^{2T}\le P$, so $\mathrm{WP}_T(a_P)$ is
$\laminfP(T)\ge0$, false for $T\ge0.60$ (Corollary~\ref{cor:Tarch}). Finally $\zeta_{a_P}(1-s)\ne\zeta_{a_P}(s)\,\Gamma$-factor, since
$\prod_{p<P}(1-p^{-s})$ is not invariant under $s\mapsto1-s$.
\end{proof}

\begin{lemma}[$\mathrm O_7$, quartets; proved]\label{lem:shape-quartet}
The models $\xi_{z,k}$ keep the comb $(\Lambda(n))_n$ exactly, the functional equation, reality and the pole, and share with $\zeta$ the
properties of Proposition~\ref{prop:quartetmodel}; for $k\ge k_0(T,\delta,\gamma)$ they have a witness at $T$, hence at every $T'\ge T$
(Theorem~\ref{thm:wit-model}(iii)). Their arithmetic side differs from that of $\zeta$ only by the absolutely continuous part $-k\beta_z$.
\end{lemma}

\begin{lemma}[$\mathrm O_8$, self-dual deformation; proved]\label{lem:shape-ft}
$f_t=\zeta+tg$ keeps every $\Lambda(n)$ with $n<\lambda_1$ and Riemann's functional equation with the archimedean factor and pole of
$\zeta$, and for $|t|$ large has a primitive witness at the window $\frac12\log\lambda_1+\varepsilon$ (Propositions~\ref{prop:ft},
\ref{prop:ftprim}). Its frequencies $m/\sqrt M$ are not integers: it breaks the Euler product over $\N$.
\end{lemma}

\begin{proposition}[$\mathrm O_9$, a twist that keeps the line data; proved]\label{prop:shape-twist}
Fix $\gamma_0>0$ with $\zeta(\frac12+i\gamma_0)\ne0$ and $0<\delta<\frac18$, and put
\[
R(s)=\prod_{\pm}\frac{(s-\frac12-\delta\mp i\gamma_0)(s-\frac12+\delta\mp i\gamma_0)}{(s-\frac12-2\delta\mp i\gamma_0)(s-\frac12+2\delta\mp i\gamma_0)},
\qquad\tilde\xi:=\xi R,\quad\tilde\zeta:=\zeta R,
\]
and $\ell(\tau):=\log|\xi(\frac12+i\tau)|$, $\tilde\ell(\tau):=\log|\tilde\xi(\frac12+i\tau)|$. Then:
\emph{(i)} $\tilde\xi(1-s)=\tilde\xi(s)$ and $\tilde\xi$ is real on $\R$; $\tilde\xi$ is meromorphic with four poles, at
$\frac12\pm2\delta\pm i\gamma_0$, inside the strip.
\emph{(ii)} On the line, $R(\frac12+i\tau)=\prod_\pm\frac{(\tau\mp\gamma_0)^2+\delta^2}{(\tau\mp\gamma_0)^2+4\delta^2}\in[\frac1{16},1]$; hence
$|\tilde\zeta|\le|\zeta|$ on the line, and, off zeros and poles, $\Rel\frac{\tilde\zeta'}{\tilde\zeta}(\frac12+i\tau)=-\frac12m(\tau)$,
the same identity as for $\zeta$.
\emph{(iii)} $\tilde\ell=\ell+\log R(\frac12+i\,\cdot)$ with $\int_\R|\log R(\tfrac12+i\tau)|\,d\tau=4\pi\delta$.
\emph{(iv)} $\tilde\ell''(\gamma_0)=\ell''(\gamma_0)+\frac3{2\delta^2}+O(\delta^2\gamma_0^{-4})$.
Consequently, let $\mathcal M$ be the class of arguments whose inputs about $\zeta$ on the line are: the functional equation with
$\Gamma_\R$ and conductor $1$, reality, the identity in (ii), and finitely many inequalities of the forms
$\int_I\varphi\,\ell\,d\tau<B$ or $\int_I|\zeta(\frac12+i\tau)|^{2k}d\tau<B_k$ with fixed bounded $I$, bounded $\varphi$ and strict
inequality. No argument in $\mathcal M$ bounds $\ell''$ from above on an interval containing $\gamma_0$. The class does not contain
arguments that use the holomorphy of $\xi$ in the strip together with its order-one growth.
\end{proposition}
\begin{proof}
(i) The zero and pole sets of $R$ are invariant under $s\mapsto1-s$ and $s\mapsto\bar s$, and a direct check of the eight linear factors
gives $R(1-s)=R(s)$ and $\overline{R(\bar s)}=R(s)$; $\xi$ has both symmetries. (ii) On $s=\frac12+i\tau$ each factor has modulus
$((\tau\mp\gamma_0)^2+\delta^2)/((\tau\mp\gamma_0)^2+4\delta^2)\in[\frac14,1]$ and the product is real positive. A function real on the
line has purely imaginary logarithmic derivative there, so $\Rel\frac{\tilde\xi'}{\tilde\xi}(\frac12+i\tau)=0$; since
$\Rel[\frac1s+\frac1{s-1}]=0$ on the line and $\Rel[\frac12\psi(\frac s2)-\frac12\log\pi]=\frac12m(\tau)$ by \eqref{eq:m}, this is the
identity stated, for $\tilde\zeta$ exactly as for $\zeta$. (iii) $\int_\R[\log(x^2+a^2)-\log(x^2+b^2)]dx=2\pi(a-b)$ for $a,b>0$, applied
with $(a,b)=(\delta,2\delta)$ to each of the two factors; $\log R\le0$ on the line. (iv) $\frac{d^2}{dx^2}\log(x^2+a^2)|_{x=0}=2/a^2$, so
the factor centred at $\gamma_0$ contributes $2/\delta^2-2/(4\delta^2)=3/(2\delta^2)$ to the second derivative at $\tau=\gamma_0$, and the
factor centred at $-\gamma_0$, evaluated at distance $2\gamma_0$, contributes $O(\delta^2\gamma_0^{-4})$. For the consequence: suppose an
argument in $\mathcal M$ derived $\ell''\le K$ on an interval $J\ni\gamma_0$. By (i), (ii) the first three inputs hold for $\tilde\zeta$ for
every $\delta$; by (ii) and (iii) each inequality of the last kind holds for $\tilde\zeta$ once $\delta$ is small (the moments of
$|\tilde\zeta|^{2k}$ do not exceed those of $|\zeta|^{2k}$, and $\int_I\varphi\tilde\ell$ differs from $\int_I\varphi\ell$ by at most
$4\pi\delta\|\varphi\|_\infty$). So the argument gives $\tilde\ell''\le K$ on $J$, contradicting (iv) for $\delta$ small. The last
sentence is (i): $\tilde\xi$ is not entire.
\end{proof}

No negative value of a window form of $\tilde\zeta$ on $\PT$ has been computed; Proposition~\ref{prop:shape-twist} excludes a class of
pointwise arguments about the curvature of $\ell$ on the line, and is not a falsifier in the sense of Lemma~\ref{lem:transfer}. An
$m$-fold twist with a witness at every window, in the manner of Theorem~\ref{thm:wit-model}(iii), would make it one; the form
$Q_{\delta+i\gamma_0}-Q_{2\delta+i\gamma_0}$ of Lemma~\ref{lem:quartet} is the natural candidate and has not been examined.

The second group of transfer lemmas rests on inputs that are cited, argued or open.

\begin{lemma}[$\mathrm O_{10}$, Euler removal; proved]\label{lem:shape-epstein}
Let $Q=x^2+5y^2$ (discriminant $-20$, class number $2$), $Z_Q(s)=\sum_{(x,y)\ne0}Q(x,y)^{-s}$ its Epstein zeta function and
$\Xi_Q(s)=\tfrac12\,20^{s/2}\Gamma_\R(s)\Gamma_\R(s+1)Z_Q(s)=(\sqrt{20}/2\pi)^s\Gamma(s)Z_Q(s)=\Xi_Q(1-s)$ its completion.
$\Xi_Q$ shares with $\xi$ reality, a functional equation (degree-two archimedean factor $\Gamma_\R(s)\Gamma_\R(s+1)$, conductor
$20$), a simple pole at $s=1$ (residue $\pi/\sqrt5$) and nonnegative Dirichlet coefficients $r_Q=1*\chi_{-20}+\chi_{-4}*\chi_5$,
and has zeros in $\Rel s>1$ \cite{DavenportHeilbronn1936}. Write $-Z_Q'/Z_Q=\sum\Lambda_Q(n)n^{-s}$ on a line $\Rel s=c$ to the
right of all zeros; $\Lambda_Q$ has both signs (no Euler product; $\Lambda_Q(2)=\Lambda_Q(3)=\Lambda_Q(7)=0$,
$\Lambda_Q(36)=-4\log6$ is the first negative value). The complete window form built from $\Lambda_Q$ and the archimedean factor
of $\Xi_Q$ (constant $C_0^Q=\log20+\psi(\tfrac14)+\psi(\tfrac34)-2\log\pi=-4.60704$, L\'evy kernel $1/\sinh(u/2)$, the same polar
term and primitive subspace $\PT$ as for $\zeta$) is negative on $\PT$ for every $T\ge1.59$: explicit even witnesses certified by
interval arithmetic give $\mathfrak a_T^Q(F)/\|F\|^2\in[-0.0599398587,-0.0599398456]$ at $T=159/100$ (degree $56$) and
$[-0.0619395254,-0.0619395173]$ at $T=8/5$ (degree $48$), and negativity propagates to larger windows by Lemma~\ref{lem:monoT},
whose proof does not use $\Lambda\ge0$. In these windows every contact $n\in\{4,5,6,9,14,16,21\}$ has $\Lambda_Q(n)\ge0$, so
nonnegativity of the comb is not the separating property. Formally $Z_Q=2\sqrt{\zeta_K L_\psi}\cosh y$ with $K=\Q(\sqrt{-5})$,
$L_\psi=L(s,\chi_{-4})L(s,\chi_5)$ (same archimedean factor and conductor), $y=\tfrac12\log(L_\psi/\zeta_K)$, whence
$\Lambda_Q=\tfrac12(\Lambda_K+\Lambda_\psi)+\Lambda_{\cosh}$ with $\Lambda_{\cosh}(n)=\log n\sum_{j\ge1}\kappa_j b^{*2j}(n)$
($\kappa_1=\tfrac12$, $b$ the coefficients of $y$; finite sums). On the $T=8/5$ witness the two Euler-product forms are positive
($+5.049$, $+1.785$) and the quadratic defect ($-3.479$, of which $99.8\%$ is the $j=1$ term) decides the sign. Hence every class of
combs closed under $\Lambda\mapsto\tfrac12(\Lambda_1+\Lambda_2)+\Lambda_{\cosh}$ contains a falsifier of the window positivity.
(Phase-200 memos \texttt{LEMMA\_E\_EPSTEIN\_001}, \texttt{COSH\_DEFECT\_001}, audit \texttt{AUDIT\_200\_EPSTEIN\_001};
Remark~\ref{rem:twohalves}; axiom list of \cite[Thms.~18.10, 18.16]{Paper46}.)
\end{lemma}

\begin{remark}[additivity of the window form under products; proved]\label{rem:shape-additive}
For an $L$-function with $\Xi_L(s)=q^{s/2}\prod_j\Gamma_\R(s+\mu_j)L(s)$ and $-L'/L=\sum\Lambda_L(n)n^{-s}$, the window form
of Lemma~\ref{lem:shape-epstein} has archimedean part $m^L(\tau)=\log q+\sum_j\big[\Rel\psi(\frac{\frac12+\mu_j}2+\frac{i\tau}2)-\log\pi\big]$
and weights $v^L_n=2\Lambda_L(n)/\sqrt n$; on $\PT$ (which depends only on $T$) the polar term vanishes. All three data are additive under
$L=L_1L_2$ (logarithmic derivatives add, conductors multiply), so $\mathfrak a_T^{L_1L_2}(F)=\mathfrak a_T^{L_1}(F)+\mathfrak a_T^{L_2}(F)$ for
$F\in\PT$. Check: $\Gamma_\R(s)$ and $\Gamma_\R(s+1)$ give densities $\frac1{2\sinh(u/2)}\pm\frac1{2\cosh(u/2)}$, summing to $1/\sinh(u/2)$, and
$C_0+\log20+\psi(\frac34)-\log\pi=C_0^Q$. Hence, for a fundamental $D<0$ with $h(D)=1$, $\mathfrak a_T^{\zeta}\ge0$ and $\mathfrak a_T^{L(\cdot,\chi_D)}\ge0$
on $\PT$ imply $\mathfrak a_T^{\zeta_{\Q(\sqrt D)}}\ge0$ on $\PT$; the converse does not follow, since $\mathfrak a_T^{L(\cdot,\chi_D)}$ has no
controlled sign.
\end{remark}

\begin{lemma}[$\mathrm O_{11}$, Beurling systems; argued]\label{lem:shape-beurling}
There are generalized-integer systems with nonnegative prime measure, an Euler product and a counting function $N(x)=\rho x+O(x^\theta)$
whose zeta function has zeros off the line $\Rel s=\frac12$ \cite{DiamondZhang2016}. Such a system keeps every property of $\Lambda$ that is
stated inside the semigroup (nonnegativity, the Euler product, the identity $\Lambda*1=\log$, the sieve axioms) and breaks only the
integrality of the frequencies. It has no functional equation with $\Gamma_\R$, and no window form of it has been shown negative: the
step from an off-line zero to a negative window bottom would require an explicit formula for the system with a test of compact support,
which is not written. The lemma is therefore an argued clause and is not used in Theorem~\ref{thm:shape}(A).
\end{lemma}

\begin{lemma}[$\mathrm O_{12}$, field change; transfer principle proved, falsifier open]\label{lem:shape-field}
A derivation of $\lamA\ge0$ that invokes only properties shared by the completed zeta functions $\xi_K$ of all number fields $K$ is a
derivation of the corresponding conclusion for every $\xi_K$. No window $T$ and field $K$ with a negative primitive bottom for $\xi_K$
is exhibited in this paper; the clause therefore says only that a proof must name the property that separates $\zeta$ from every
$\xi_K$, or else be held to the standard of proving the sign for all of them (the principle of \cite[Thm.~18.1]{Paper46}, whose own
proof passes through a statement about $\zeta_K$ that is not used here).
\end{lemma}

\begin{lemma}[$\mathrm O_{13}$, bounded families; proved]\label{lem:shape-fam}
If $\mathsf P$ depends on $T\mapsto\lambda(T;q)$ only through $J\subset(0,T_1]$, monotonicity and continuity, and on the weights only
through membership in a class closed under $\mathrm O_2$, then $\Lambda+c\mathbf 1_{n=n_0}$ with $n_0>e^{2T_1}$ has $\mathsf P$ and fails
$\mathrm{WP}_T$ for $T>\frac12\log n_0$ and $c>c_*(T,n_0)$; if instead the weights are only known to come from a self-dual Dirichlet
series with the functional equation of $\zeta$, $f_t$ does the same (Theorems~\ref{thm:famblind}, \ref{thm:famFE}). Not excluded: the use
of $\lambda$ on an unbounded set of windows, which enters only through $\lim_T\lamA\ge0$, and two-sided laws $\lamA=L(T)(1+o(1))$
with $L>0$ (Remark~\ref{rem:bar-fam}).
\end{lemma}

\begin{lemma}[$\mathrm O_{14}$, local criteria; proved]\label{lem:shape-local}
No criterion that depends on $\sT$ only through its restriction to $\{\sT<\eta\}$ decides the sign at any window: $\sigma^{(\eta)}$ has the
same negative set and a negative form for $\eta<\eta_0(T)$ (Theorem~\ref{thm:nolocal}).
\end{lemma}

\begin{lemma}[$\mathrm O_{15}$, random positions; proved]\label{lem:shape-poisson}
For $T\ge T(c,\kappa)$ the Poisson comb of intensity $\kappa$ fails $\mathrm{WP}_T$ with probability at least $1-e^{-cT\log T}$; in
particular some realization fails it, so no deduction that uses the positions $\{\log n\}$ only through their intensity and moments proves
$\mathrm{WP}_T$ at such $T$ (Theorem~\ref{thm:poisson}).
\end{lemma}

\begin{lemma}[derived classes; proved]\label{lem:shape-classes}
\emph{(i)} Classes defined by statements on $\psi_a(x;q,c)$ that are stable under $O(1)$ changes are closed under $\mathrm O_1$ and
$\mathrm O_2$ (Proposition~\ref{prop:prog}); classes $0\le a\le\Lambda$ under $\mathrm O_1$. \emph{(ii)} $\mathcal C_{\mathrm{zeros}}$ is
closed under $\mathrm O_1$ and $\mathrm O_2$, since $\zeta_a=\zeta E$ with $E$ entire and zero-free (Corollary~\ref{cor:void}); and no finite
symmetric multiset of ordinates carries a flat-top certificate (Theorem~\ref{thm:finZ}). \emph{(iii)} Deductions from window observables
(the bottom up to a relative error, the parity of the minimizer, the peak of $|\widehat{h^*}|^2$, eigenvalue counts, laws fitted to these)
are excluded numerically by $\xi_{z,1}$ and, for the shared properties of Proposition~\ref{prop:quartetmodel}, by $\xi_{z,k}$ with $k$ large
(Corollary~\ref{cor:pattern}).
\end{lemma}

\subsection{Structural lemmas}

These constrain every proof, in whatever class, and produce no falsifier.

\begin{lemma}[margin and universal criteria; proved, inputs cited]\label{lem:shape-margin}
\emph{(i)} $\lamW\le\lamP\le\lamA\le4\cdot10^9Te^{-15.449T}$ for $T\ge1$ (Corollary~\ref{cor:margin1}; the zero data are those of
\cite{PlattTrudgian2021}). Any lower bound that proves \eqref{eq:target} at $T$ has this absolute precision; relative to the comb mass
$\Aconst\asymp e^T$ it is $e^{-16.4T}$. \emph{(ii)} A criterion valid for every nonnegative trace-class operator, or for every bounded
perturbation of the form, has margin at most $k\lamA/c_P$ (moments), $L\lamA$ (Lipschitz functionals), and a leading term of the window
determinant with relative error $\varepsilon_0\ge1$ never decides (Lemmas~\ref{lem:wh-T}, \ref{lem:wh-Tprime},
Theorem~\ref{thm:wh-detinherit}(iii)). The modified objects in (ii) have bottom $0$, not a negative bottom: these are precision
constraints, not instances of Lemma~\ref{lem:transfer}.
\end{lemma}

\begin{proposition}[criticality]\label{prop:shape-crit}
\emph{(a) The centre (proved).} $m(0)=C_0=-2\zeta'(\frac12)/\zeta(\frac12)$ (Proposition~\ref{prop:boundary}), and the zero sum behind it
vanishes termwise: for a pair $\rho$, $1-\bar\rho$ one has $\Rel[\frac1{\frac12-\rho}+\frac1{\frac12-(1-\bar\rho)}]=0$. The identity holds
for every zero configuration symmetric under $\rho\mapsto1-\bar\rho$ and carries no slack.
\emph{(b) The packets (conditional).} For every $\gamma>0$ the threshold of $(E_\gamma)$ is exactly $m^\sharp(\gamma)$
(Proposition~\ref{prop:Egamma-exact}). If, in addition, the Riesz means of the twisted sums $\mathcal F_T(F_T)/\|F_T\|^2$ converge
whenever they are bounded, then $(E_\gamma)$ with any additive constant $C\ge0$ is equivalent to $(E_\gamma)$ itself: there is no slack at
any packet frequency. The additional hypothesis is a statement about the divergent alternative that is not proved here.
\emph{(c) The integer resolution (proved; a Mertens-type estimate cited).} Let $P_T(\tau)=\sum_{n\in\NT}v_n\cos(\tau\log n)$. Then
\[
\lim_{Y\to\infty}\frac1{2Y}\int_{-Y}^{Y}P_T(\tau)^2d\tau=\frac12\sum_{n\in\NT}v_n^2=4T^2\,(1+o(1))\qquad(T\to\infty),
\]
so the root mean square of the comb is $2T(1+o(1))$, while the archimedean symbol at the frequency $e^{2T}$ of the integer resolution is
$m(e^{2T})=2T+O(1)$ (Lemma~\ref{lem:mbounds}(ii)): comb and symbol are of the same size at that resolution.
\emph{(d) The density ceiling (proved).} A comb $a\ge0$ with $\lambda_{\mathrm A}(T;a)\ge0$ for every $T$ has
$\liminf S(x)/\log x\le1$, and $\Lambda$ sits at the ceiling (Theorem~\ref{thm:ceiling}); the main-term comb $\rho_\kappa$ is positive at
every window if and only if $\kappa\ge\kappa_\star$, where $\kappa_\star\ge\frac14|C_0|=1.343>1$ (Theorem~\ref{thm:kappa}). So the main
term of $\Lambda$ is below the positivity threshold of main terms.
\emph{(e) No asymptotic slack (proved).} \eqref{eq:target} holds as soon as $\lamA\ge-\varepsilon(T)$ for all $T$ with
$\varepsilon(T)\to0$, or $\lamA\ge0$ along one sequence $T_k\to\infty$ (Theorem~\ref{thm:wit-Tc}(iv)); there is no weaker asymptotic
sub-statement in $T$.
\end{proposition}
\begin{proof}
(a) Proposition~\ref{prop:boundary} and the displayed computation, in which $(\frac12-\beta)+(\beta-\frac12)=0$ for $\rho=\beta+i\gamma$.
(b) Proposition~\ref{prop:Egamma-exact}: under the hypothesis, a bounded $\limsup$ is a limit, and a limit that is at most
$m^\sharp(\gamma)+C$ for the packets $F_T$ at every scale is at most $m^\sharp(\gamma)$, since the limit does not depend on $C$.
(c) The frequencies $\log n$, $n\in\NT$, are distinct, so the Ces\`aro mean of $P_T^2$ is $\frac12\sum v_n^2=2\sum_{n<e^{2T}}\Lambda(n)^2/n$.
Since $\Lambda(n)^2=\Lambda(n)\log n$ on primes and $\sum_{p^k,k\ge2}\Lambda(p^k)\log p^k/p^k<\infty$, one has
$\sum_{n\le x}\Lambda(n)^2/n=\sum_{n\le x}\Lambda(n)\log n/n+O(1)=\frac12\log^2x+O(\log x)$ \cite[Thm.~2.7 and partial summation]{MontgomeryVaughan2007},
and $x=e^{2T}$ gives $2\cdot\frac12(2T)^2=4T^2$. (d) and (e) are the cited theorems.
\end{proof}

\begin{lemma}[antisymmetry; proved]\label{lem:shape-antisym}
$D_{1-\sigma}=-D_\sigma$ exactly, every combination of the family $D_\sigma$ symmetric under $\sigma\mapsto1-\sigma$ vanishes, and
$D_{1/2}\equiv0$ off the ordinates (Remark~\ref{rem:fe}). No positivity is obtained by symmetrizing the family.
\end{lemma}

\begin{lemma}[operator calculus and squares; proved]\label{lem:shape-lin}
\emph{(i)} Every expression $\sum_i\Rel\langle A_i*F,B_i*F\rangle$ with finite measures $A_i,B_i$ equals $\int_{(-2T,2T)}\Rel\cF\,d\nu$
(Proposition~\ref{prop:opcalc}): an identity in shift-multiplier calculus is an identity of measures on the window and carries the whole
margin, and the only further tool of that calculus is the positivity $\|A*F\|^2\ge0$. The statement covers convolution operators only.
\emph{(ii)} No arithmetic square decomposition exists at any $T\ge\frac12\log3$; with a free positive definite part the decompositions are
exactly the certificates of Definition~\ref{def:PD}, whose value is $\lamP$ (Theorems~\ref{thm:sos}, \ref{thm:D1}). The interval
$(\frac12\log2,\frac12\log3)$ is not decided.
\end{lemma}

\begin{lemma}[orders, relative bounds, exterior measures, termwise bounds; proved]\label{lem:shape-misc}
The M\"obius resolvent is not monotone for the order of $\mathcal K^*_T$ (Proposition~\ref{prop:cone}); no relative bound with constants
$(\theta,\kappa,C)$ decides for large $T$, not even $\theta=0$ (Proposition~\ref{prop:relative}); no certificate of the complete form with
a positive exterior measure exists at the windows where $\kappa_T(\Phi)<0$ (Theorem~\ref{thm:posext}; the two values in
Remark~\ref{rem:posext} are floating evaluations); and a bound that treats the comb term by term costs its full mass
(Proposition~\ref{prop:primeblind}).
\end{lemma}

\begin{lemma}[flows and new contacts; identity proved, readings argued]\label{lem:shape-flow}
The coefficient of any $C^1$ flow of $\log\det$ on Galerkin sections near a window with small bottom is dominated by
$\langle A'h^*,h^*\rangle/\lamA$, a functional of the minimizer (Lemma~\ref{lem:flowdet}, Remark~\ref{rem:bar-loc}); positivity at $T>T_0$
imposes on the new contacts an averaged inequality on $U_{T_0}$, never a pointwise one (Theorem~\ref{thm:X}), and the converse of that
inequality fails for the quartet model, numerically (Remark~\ref{rem:X}).
\end{lemma}

\begin{proposition}[the three inputs enter jointly]\label{prop:shape-inputs}
Consider the three properties: (E) the Euler product over $\N$, equivalently $\sum_{d\mid n}a(d)=\log n$ for $n\in\NT$; (F) the functional
equation with $\Gamma(s/2)$, conductor $1$, reality and the pole; (I) the integrality of the frequencies, that is, support on $\log\N$.
\emph{(i) (proved)} A deduction that uses (E) and (I) with nonnegativity, the local values $a(p^k)=\log p$ and the location of the zeros
and the pole, but not (F), proves nothing: $a_P$ has all of these (Proposition~\ref{prop:shape-eulerdel}).
\emph{(ii) (cited)} A deduction that uses (F) and (I) with nonnegative coefficients, but not (E), proves nothing for the complete form:
$\Xi_Q$ has all of these (Lemma~\ref{lem:shape-epstein}).
\emph{(iii) (argued)} A deduction that uses (E) and (F)-free semigroup properties, but not (I), is matched by the Beurling systems
of Lemma~\ref{lem:shape-beurling}.
\emph{(iv) (proved)} The modifications $\mathrm O_7$ and $\mathrm O_8$ keep (F): the quartet models keep (E) as a comb identity but add an
absolutely continuous arithmetic part, and $f_t$ breaks (I). Hence a property set that survives (i), (ii) and (iv) contains (E) and (F)
jointly, with (I); and (E) with (F) at conductor $1$ determines $\zeta$ up to a constant, by Hamburger's theorem \cite[\S2.13]{Titchmarsh1986}.
\end{proposition}
\begin{proof}
(i), (ii), (iv) are the lemmas cited; the quartet models have logarithmic derivative $\xi'/\xi$ plus a rational function, which is not a
Dirichlet series over $\N$, and the frequencies $m/\sqrt M$ of $f_t$ are irrational. The last sentence is Hamburger's theorem: a Dirichlet
series over $\N$ of finite order with the functional equation of $\zeta$ is a constant multiple of $\zeta$.
\end{proof}

\begin{lemma}[no weaker substatement; proved]\label{lem:shape-restate}
Each of the following is an exact restatement of \eqref{eq:target}, not a weaker statement: a flat-top certificate $f_w\ge0$ at every
window (Theorem~\ref{thm:Z}: $f_w=\sum_\rho\wh(\tau-\gamma_\rho)$, so its content is the explicit formula on the tests $w\cos(\tau\,\cdot)$);
the indeterminacy of the extension problem at every window off $\{\lamW=0\}$ (Theorem~\ref{thm:KEPdich}); the positive definiteness of
the one-profile function $\mathsf B$ (Theorem~\ref{thm:oneprofile}); the positivity of $\mathrm{Arch}_M-\mathrm{Div}_M-\mathrm{Near}_M$ for
every $M$ (Proposition~\ref{prop:AM}); the finite matrices $S_{T,k(T)}\succeq0$ (Theorem~\ref{thm:coer-schur}); and the one-sided
asymptotic $\lamA\ge-\varepsilon(T)$ (Proposition~\ref{prop:shape-crit}(e)). A proof of any of them is a proof of \eqref{eq:target} and has
the shape described below.
\end{lemma}

\subsection{Estimates and exact identities}\label{subsec:shape-est}

The modifications of Definition~\ref{def:shape-ops} act on the object. This subsection constrains the two remaining ways in which a sign
step can use the vector $(\Lambda(n))_{n<e^{2T}}$: through an inequality on prime sums with a remainder, and through an exact identity.
For a comb $a$ and $F\in\PT\cap\mathcal D_T$ put
\[
\mathcal F^{(a)}_T(F):=\sum_{n\in\NT}v_n(a)\Rel\cF(\log n)-2\int_0^{2T}\Rel\cF(u)\,e^{u/2}\,du ,
\]
so that $\mathcal F^{(\Lambda)}_T=\mathcal F_T$ (Proposition~\ref{prop:wit-cheb}) and, by Lemma~\ref{lem:polarelim} exactly as in
Proposition~\ref{prop:sharp}(ii), $\mathfrak a^{(a)}_T(F)=\mathfrak a^\sharp(F)-\mathcal F^{(a)}_T(F)$; the polar part does not depend on $a$.
An \emph{estimate at $T$} is a pair $(E_T,R_T)$, $E_T$ a real function on the unit vectors of $\PT\cap\mathcal D_T$ that does not depend
on the comb and $R_T\ge0$; a comb $a$ \emph{satisfies} it if $|\mathcal F^{(a)}_T(F)-E_T(F)|\le R_T$ for every such $F$. A deduction
\emph{decides by the estimate} if it concludes $\lambda_{\mathrm A}(T;a)\ge0$ from ``$a$ satisfies $(E_T,R_T)$'' together with
$\mu_T\ge R_T$, where $\mu_T:=\inf_{\|F\|=1}\big(\mathfrak a^\sharp(F)-E_T(F)\big)$. The prime number theorem with a remainder, sieve
bounds, mean-value theorems over progressions and large-sieve inequalities enter a proof of \eqref{eq:target} in this form when they
are used to replace the twisted sum $\mathcal F_T$ by a main term.

\begin{lemma}[estimates with a remainder; (i), (ii) proved, (iii) conditional, (iv) argued]\label{lem:shape-est}
\emph{(i) Pinching.} If $\Lambda$ satisfies $(E_T,R_T)$, then $\mu_T-R_T\le\lamA\le\mu_T+R_T$. Hence a deduction that decides by the
estimate at $T\ge1$ has $0\le\mu_T-R_T\le4\cdot10^9Te^{-15.449T}$: its main term and its remainder agree to that absolute precision.
\emph{(ii) Exact reductions inherit the margin.} Let $\mathfrak q=\mathfrak a_T$ on $\PT$, decomposed as in Lemma~\ref{lem:wh-schurabs}
with $E$ finite-dimensional and $c_{\mathcal X}>0$ (for instance the reductions of Theorems~\ref{thm:coer-schur} and \ref{thm:wh-schur}).
If $\lamA\ge0$ and $\lambda_S<c_{\mathcal X}$, then $\lamA\le\lambda_S\le(1+\|\mathsf K\|)^2\lamA$. A step that decides the sign of
the reduced form $\mathfrak s$ by an estimate is therefore subject to (i) with $\lamA$ replaced by $(1+\|\mathsf K\|)^2\lamA$.
\emph{(iii) Slack.} Suppose $\lamA\ge0$, let $\kappa_T$ be as in Proposition~\ref{prop:wit-sens}, and suppose $\Lambda$ satisfies the
estimate with slack $s_T$, that is $|\mathcal F_T(F)-E_T(F)|\le R_T-s_T$ for every unit $F$. If $\kappa_T\ne0$ and
$s_T>\Aconst\lamA/|\kappa_T|$, $\lamA<\frac12|\kappa_T|$, then there is a comb $(1+\delta)\Lambda$, $0<|\delta|<\frac12$, that satisfies
$(E_T,R_T)$ and has $\lambda_{\mathrm A}(T;(1+\delta)\Lambda)<0$; so no deduction that decides by this estimate is valid at $T$. For
$T\ge1$ the threshold is at most $4\cdot10^9Te^{-15.449T}\Aconst/|\kappa_T|$, with $\Aconst=(4+o(1))e^T$.
\emph{(iv) Reading.} The classical estimates are proved for classes of sequences and with remainders of size at least
$e^{-cT}$ relative to the comb mass for a fixed $c$ at best: with the Vinogradov--Korobov error term the twisted integral of
Proposition~\ref{prop:wit-cheb} is bounded by $e^{T-c(2T)^{3/5}(\log2T)^{-1/5}}$, and a bound $|\psi(y)-y|\le Cy^\theta$ gives
$O(e^{(2\theta-1)T})$ (Remark~\ref{rem:wit-cheb}). By (i) such a remainder is matched by the main term to precision
$4\cdot10^9Te^{-15.449T}$ only if the main term is itself an exact evaluation of the sum it replaces, and by (iii) a stated remainder that
exceeds the true error by more than the threshold admits a falsifier. We read this as: an estimate with a remainder cannot carry the
sign step of a proof of \eqref{eq:target}. The reading is argued, not proved: an estimate is a deduction class in the sense of
Definition~\ref{def:shape-class} only once its property set is named, and (iii) needs $\kappa_T\ne0$.
\end{lemma}

\begin{proof}
(i) For unit $F$, $\mathfrak a_T(F)=\mathfrak a^\sharp(F)-\mathcal F_T(F)$ lies within $R_T$ of $\mathfrak a^\sharp(F)-E_T(F)$; take
infima. The second sentence combines $\mu_T\ge R_T$ with Corollary~\ref{cor:margin1}. (ii) Lemma~\ref{lem:wh-schurabs}(ii): if
$\lambda_{\mathfrak q}\ge0$ then $\lambda_S\ge\lambda_{\mathfrak q}$, and the lower bound with $\min\{c_{\mathcal X},\lambda_S\}=\lambda_S$
gives $\lambda_S\le(1+\|\mathsf K\|)^2\lambda_{\mathfrak q}$. (iii) Put $\varepsilon:=s_T/\Aconst$ if $s_T/\Aconst<\frac12$ and otherwise any
$\varepsilon\in(\lamA/|\kappa_T|,\frac12)$; in both cases $\lamA/|\kappa_T|<\varepsilon<\frac12$ and $\varepsilon\Aconst\le s_T$. Let
$\delta:=\varepsilon\,\mathrm{sgn}\,\kappa_T$. Then $\mathcal F^{((1+\delta)\Lambda)}_T(F)-\mathcal F_T(F)=\delta\sum_nv_n\Rel\cF(\log n)$ has
modulus at most $\varepsilon\Aconst\le s_T$ for unit $F$, since $|\cF|\le\|F\|^2$; so $(1+\delta)\Lambda$ satisfies $(E_T,R_T)$. With $h$ the
minimizer of Proposition~\ref{prop:wit-sens}, $\lambda_{\mathrm A}(T;(1+\delta)\Lambda)\le\mathfrak a^{((1+\delta)v)}_T(h)=\lamA-\varepsilon|\kappa_T|<0$.
Lemma~\ref{lem:transfer} applies. The bound for $T\ge1$ is Corollary~\ref{cor:margin1}, and $\Aconst=(4+o(1))e^T$ is
Proposition~\ref{prop:env}(iii). (iv) is a reading of (i)--(iii) and of
Remark~\ref{rem:wit-cheb}.
\end{proof}

\begin{remark}[the size of the margin and three exact formulations]\label{rem:shape-est}
(a) Only the single-exponential bound of Corollary~\ref{cor:margin1} is proved. If the conjectural law $-\log\lamP\asymp e^{2T}$
\cite{Zhu2026} holds, the precision in Lemma~\ref{lem:shape-est}(i) is $\exp(-e^{(2-\varepsilon)T})$ for large $T$, against remainders of
size $e^{-cT}$. This is not used anywhere.
(b) Three ways of selecting frequencies were examined as candidate estimates, and each turned out to be exact or false. At one fixed
frequency $\gamma$ the statement is $(E_\gamma)$, whose threshold is exactly $m^\sharp(\gamma)$ (Proposition~\ref{prop:Egamma-exact}, proved);
the absence of any additive slack is conditional (Proposition~\ref{prop:shape-crit}(b)). At the Nyquist mesh $\omega_k=\pi k/2T$, the pointwise form, a sign for $\sT(\omega_k)$ at every mesh point without correlation
with the admissible samples of $\PT$, fails at $\omega_0=0$ at every window, since $\sT(0)=C_0-\Aconst<0$ (Proposition~\ref{prop:band}(ii)); the correlated form is $(N)$ of \S\ref{sec:open}, an exact restatement.
On the family $s_w$ of Proposition~\ref{prop:shape-family} below, the comb collapses exactly through $\sum_{d\mid N}\Lambda(d)=\log N$ and
no twisted prime sum remains (Proposition~\ref{prop:shape-family}(i), proved). None of the three leaves room for a remainder.
\end{remark}

\begin{proposition}[the exterior field of a primitive function; proved]\label{prop:shape-ext}
Let $h\in\PT$ be real, put $H(x):=x^{1/2}h(\log x)$ for $x>0$ and $K_1(u):=K(u)-2e^{-u/2}$ (Lemma~\ref{lem:K}). For $t>T$ let
\[
E_h(t):=-\frac12\int_{\IT}K(t-s)h(s)\,ds-\frac12\sum_{n\in\NT}v_n\,h(t-\log n),
\]
the restriction to $t>T$ of the distribution whose Fourier transform is $\sT\Fh$ with $F=h$. Then, for almost every $t>T$, with $\int_{-T}^{a}:=0$ for $a<-T$,
\[
E_h(t)=e^{t/2}\int_{-T}^{t-2T}h(s)e^{-s/2}ds\;-\;e^{-t/2}\int_{[1,e^{2T})}H(e^t/u)\,d\big(\psi(u)-u\big)\;-\;\frac12\int_{\IT}K_1(t-s)h(s)\,ds .
\]
The first term is supported in $(T,3T)$ and is the field of the main-term comb $e^{u/2}du$ on $(0,2T)$: at $t=T+\epsilon$ it is
$e^{t/2}$ times the partial Tate moment of $h$ on the opposite collar $(-T,-T+\epsilon)$. The second is the Chebyshev fluctuation tested
by $h$ read backwards, and the third is $\sum_{k\ge1}e^{-(2k+\frac12)t}\int h(s)e^{(2k+\frac12)s}ds$. For $t<-T$ the same holds with $h(s)$
replaced by $h(-s)$.
\end{proposition}

\begin{proof}
For $t>T$ only $h(t-\log n)$ and the kernel term survive in $\sT(D)h$, since $h$ vanishes near $t$. The kernel term is
$-\frac12\int2e^{-(t-s)/2}h(s)ds-\frac12\int K_1(t-s)h(s)ds=-e^{-t/2}A(h)-\frac12\int K_1(t-s)h(s)ds$, and $A(h)=0$. The comb term is
$-\int_{[1,e^{2T})}u^{-1/2}h(t-\log u)\,d\psi(u)$, and $u^{-1/2}h(t-\log u)=e^{-t/2}H(e^t/u)$. Split $d\psi=d(\psi-u)+du$; with $u=e^r$,
$s=t-r$, the $du$ part is $e^{t/2}\int_{t-2T}^Te^{-s/2}h(s)ds=e^{t/2}\big(B(h)-\int_{-T}^{t-2T}e^{-s/2}h(s)ds\big)$, and $B(h)=0$. For
$t\ge3T$ both prime-side terms vanish, the first being $e^{t/2}B(h)$. The expansion of $K_1$ is Lemma~\ref{lem:K}.
\end{proof}

\begin{proposition}[the integer family; (i)--(iii) proved, the growth of $\zeta$ cited]\label{prop:shape-family}
Let $F_0$, $T_0$, $c_0$ be as in \S\ref{subsec:oneprofile} with $2T_0<\log2$, let $M\ge2$, $\eta\in(0,1)$, and let $w\in C^\infty[0,\infty)$
with $\supp w\subseteq[0,1]$ and $w\equiv c$ on $[0,\eta]$. Put $s_w:=(m^{-1/2}w(m/M))_{m\le M}$ and
$\zeta_{w,M}(s):=\sum_{m\le M}m^{-s}w(m/M)$.
\emph{(i) Collapse.} With $D_w(N):=\sum_{d\mid N,\,d\ge2}2\Lambda(d)\,w\big(\frac N{dM}\big)$ and
$E_w(N):=\sum_{m\le M}m^{-1/2}w(\frac mM)c_0(\log\frac mN)$,
\[
s_w^\top\mathbb A_Ms_w=s_w^\top\mathrm{Arch}_Ms_w-\sum_{2\le N<Me^{2T_0}}N^{-1/2}D_w(N)E_w(N),
\]
\[
D_w(N)=2c\log N-\sum_{\substack{d\mid N\\2\le d<N/(\eta M)}}2\Lambda(d)\Big(c-w\big(\tfrac N{dM}\big)\Big).
\]
So on this family the comb enters only through $\sum_{d\mid N}\Lambda(d)=\log N$ for $N<Me^{2T_0}$ and through the values $\Lambda(d)$,
$d<e^{2T_0}/\eta$.
\emph{(ii) Spectral form.} With $R_X(s):=-\frac{\zeta'}{\zeta}(s)-\sum_{n<X}\Lambda(n)n^{-s}$ and $X'=Me^{2T_0}$,
\[
s_w^\top\mathbb A_Ms_w=\frac1\pi\int_\R\Fh_0(\tau)^2\,|\zeta_{w,M}(\tfrac12+i\tau)|^2\,\Rel R_{X'}(\tfrac12+i\tau)\,d\tau ,
\]
where $\Rel R_{X'}(\frac12+i\tau)=\frac12\sigma_{T'}(\tau)$, $e^{2T'}=X'$, extends continuously across the ordinates.
\emph{(iii) Mellin form of the test.} Let $\mathcal W(z):=\int_0^\infty w(x)x^{z-1}dx$ for $\Rel z>0$. Then
$\mathcal W(z)=-z^{-1}\int_0^\infty w'(x)x^z\,dx$ is meromorphic with one simple pole, at $z=0$, of residue $c$, decreases rapidly on
vertical lines, and for every $A>0$ and $s=\frac12+i\tau$,
\[
\zeta_{w,M}(s)=c\,\zeta(s)+M^{1-s}\mathcal W(1-s)+\frac1{2\pi i}\int_{\Rel z=-A}\zeta(s+z)M^z\mathcal W(z)\,dz .
\]
\emph{(iv) The continuous term vanishes to second order.} For $u>0$ let
\[
\Psi_w(u):=\int_1^\infty w(y/M)\int_\R e^{v/2}c_0(v)\,w\big(uye^v/M\big)\,dv\,dy ,
\]
so that $n^{1/2}\Psi_w(n)$ is the integral, over $x>0$ and $y\ge1$, of $(xy)^{-1/2}w(\frac xM)w(\frac yM)c_0(\log\frac x{ny})$, the continuous
counterpart of the atom-$n$ part of $s_w^\top\mathrm{Div}_Ms_w+s_w^\top\mathrm{Near}_Ms_w$. The inner integral equals
$\int e^{v/2}c_0(v)\big[w(ae^v)-w(a)-aw'(a)v\big]dv$ with $a=uy/M$; it vanishes when $ae^{2T_0}\le\eta$, and its modulus is at most
$\frac12\sup|(x\partial_x)^2w|\int v^2e^{v/2}|c_0(v)|dv$. Moreover $\int_0^\infty\Psi_w(u)\,du=0$ and
\[
\sum_{n\ge2}2\Lambda(n)\Psi_w(n)=2\int_{[1,\infty)}\Psi_w(u)\,d\big(\psi(u)-u\big)-2\int_0^1\Psi_w(u)\,du .
\]
\end{proposition}

\begin{proof}
(i) By Proposition~\ref{prop:AM}(ii) and the evenness of $c_0$, the comb part of $s^\top\mathbb A_Ms$ is
$\sum_n v_n\sum_{m,k}s_ms_kc_0(\log m-\log k-\log n)$. Put $N=kn$: $v_ns_k=2\Lambda(n)N^{-1/2}w(\frac N{nM})$, which gives
$N^{-1/2}D_w(N)E_w(N)$, and $E_w(N)=0$ unless $N<Me^{2T_0}$ since $\supp c_0\subseteq[-2T_0,2T_0]$. If $d\ge N/(\eta M)$ then
$w(\frac N{dM})=c$, and $\sum_{d\mid N}\Lambda(d)=\log N$ gives the second formula; $d<N/(\eta M)<e^{2T_0}/\eta$.
(ii) Proposition~\ref{prop:AM}(ii) with $|P_{s_w}(\tau)|=|\zeta_{w,M}(\frac12+i\tau)|$ ($w$ real), Proposition~\ref{prop:boundary}, and
$\sum_{n<X'}v_n\cos(\tau\log n)=2\Rel\sum_{n<X'}\Lambda(n)n^{-1/2-i\tau}$. At an ordinate the pole of $-\zeta'/\zeta$ on the line is
$i\,m_\rho/(\tau-\gamma)$, purely imaginary.
(iii) $w'$ is smooth with compact support in $[\eta,1]$, which gives the continuation by one integration by parts, the residue
$-\int w'=c$, and rapid decay by further integrations by parts. Mellin inversion gives $w(n/M)=\frac1{2\pi i}\int_{(b)}\mathcal W(z)(n/M)^{-z}dz$
for $b>0$; summing over $n$ with $b>\frac12$ gives $\frac1{2\pi i}\int_{(b)}\zeta(s+z)M^z\mathcal W(z)dz$. Moving the line to $\Rel z=-A$
crosses the poles $z=1-s$ (residue $M^{1-s}\mathcal W(1-s)$) and $z=0$ (residue $c\zeta(s)$); the horizontal segments vanish because
$\zeta$ grows at most polynomially in vertical strips \cite[\S5.1]{Titchmarsh1986}.
(iv) Substituting $x=nye^v$ gives the description of $n^{1/2}\Psi_w(n)$. Write $g(v)=w(ae^v)$, so $g'(0)=aw'(a)$ and
$g''(v)=((x\partial_x)^2w)(ae^v)$. Since $\Fh_0(\pm\frac i2)=0$ ($F_0\in\mathcal P_{T_0}$), the entire function
$\widehat{c_0}(z)=\Fh_0(z)\Fh_0(-z)$ has a double zero at $z=-\frac i2$, that is $\int e^{v/2}c_0=\int ve^{v/2}c_0=0$; this removes $g(0)$ and
$g'(0)v$, and Taylor's formula bounds the rest. If $ae^{2T_0}\le\eta$ then $g\equiv c$ on $\supp c_0$. For the integral of $\Psi_w$ substitute
$z=uye^v/M$: $\int_0^\infty\Psi_w=\int_0^\infty w\cdot\int_1^\infty w(y/M)\frac{M\,dy}y\cdot\int e^{-v/2}c_0(v)dv=0$, by the other Tate
condition; all integrals converge because $w$ is bounded with support in $[0,1]$ and $y\ge1$. Finally $\Psi_w$ is bounded and vanishes for
$u\ge Me^{2T_0}$, $\psi$ has no mass on $[0,2)$, and $\sum_n2\Lambda(n)\Psi_w(n)=2\int_{[1,\infty)}\Psi_w\,d\psi$; split $d\psi=d(\psi-u)+du$ and use
$\int_1^\infty\Psi_w=-\int_0^1\Psi_w$.
\end{proof}

\begin{proposition}[the window phase; proved]\label{prop:shape-phase}
Let $X=e^{2T}$, $\vartheta(\tau):=\Iml\log\Gamma(\frac14+\frac{i\tau}2)-\frac\tau2\log\pi$ (continuous branch, $\vartheta(0)=0$), and
\[
\Theta_T(\tau):=\vartheta(\tau)-\sum_{n\in\NT}\frac{\Lambda(n)}{\sqrt n}\,\frac{\sin(\tau\log n)}{\log n},
\]
\[
Z_T(\tau):=e^{i\vartheta(\tau)}\zeta(\tfrac12+i\tau)\exp\Big(-\sum_{n\in\NT}\frac{\Lambda(n)}{\sqrt n}\,\frac{\cos(\tau\log n)}{\log n}\Big).
\]
Then $\sT=2\Theta_T'$, so $\mathfrak a_T(F)=\frac1\pi\int\Theta_T'|\Fh|^2$; $Z_T$ is real on $\R$ and vanishes exactly at the ordinates of the zeros on
the line; and off these ordinates $R_X(\frac12+i\tau)=\Theta_T'(\tau)+i\,Z_T'(\tau)/Z_T(\tau)$, with $R_X$ as in
Proposition~\ref{prop:shape-family}(ii).
\end{proposition}

\begin{proof}
$\vartheta'=\frac12\Rel\psi(\frac14+\frac{i\tau}2)-\frac12\log\pi=\frac12m$, and the derivative of the sum is $\frac12\sum v_n\cos(\tau\log n)$. The
function $e^{i\vartheta}\zeta(\frac12+i\tau)$ is real (Hardy's function, \cite[\S4.17]{Titchmarsh1986}), hence so is $Z_T$. With
$Z_{<X}(s):=\exp\sum_{n<X}\Lambda(n)n^{-s}/\log n$, entire and zero-free, one has $Z_T(\tau)=e^{i\Theta_T(\tau)}(\zeta/Z_{<X})(\frac12+i\tau)$, and
$-\frac{d}{ds}\log(\zeta/Z_{<X})=R_X$; so $\frac{d}{d\tau}\log Z_T=i\Theta_T'-iR_X(\frac12+i\tau)$, which is real.
\end{proof}

\begin{lemma}[identities valid for every comb; proved]\label{lem:shape-ident}
Let $\mathsf I$ be a property that every real weight vector $a:\N_{\ge2}\to\R$ has, and let $\mathsf P'$ be a property set. If the class with property set $\mathsf P'$
is closed under a modification $\mathrm O\in\{\mathrm O_1,\dots,\mathrm O_6\}$ (in the ranges of Lemmas~\ref{lem:shape-del}--\ref{lem:shape-scale},
\ref{lem:shape-weight}, \ref{prop:shape-eulerdel}), then so is the class with property set $\mathsf P'\cup\{\mathsf I\}$, and
Theorem~\ref{thm:shape}(A), (A$'$) apply to it unchanged. The identities of Propositions~\ref{prop:shape-ext}, \ref{prop:shape-family}(i),(ii),(iv)
and \ref{prop:shape-phase} hold, with the same proofs, for every real weight vector $a$ in place of $\Lambda$ (with $\psi_a(u)=\sum_{n\le u}a(n)$, and in
\ref{prop:shape-family}(i) with $\sum_{d\mid N}a(d)$ in place of $\log N$); those of Propositions~\ref{prop:shape-ext}, \ref{prop:shape-family}(iv)
and \ref{prop:shape-phase}(first assertion) hold for every finite measure on $(0,2T)$ in place of the comb, hence survive $\mathrm O_{11}$ and
$\mathrm O_{15}$; Proposition~\ref{prop:shape-family}(iii) does not involve the comb.
\end{lemma}

\begin{proof}
$\mathrm O_1$--$\mathrm O_6$ map weight vectors to weight vectors, so $\mathrm O(q)$ has $\mathsf I$ whenever $q$ is a comb; closure is then a property of
$\mathsf P'$ alone. The proofs of the propositions use the comb only through $d\psi$, written as $d(\psi-u)+du$, and never its values;
Proposition~\ref{prop:shape-family}(i) uses $\sum_{d\mid N}\Lambda(d)=\log N$ only to rename $\sum_{d\mid N}a(d)$.
\end{proof}

\begin{remark}[what the identities do and do not give]\label{rem:shape-ident}
By Lemma~\ref{lem:shape-ident} the identities above are coordinates, not mechanisms: a deduction in which one of them is load-bearing
must add a property of $\zeta$ that the modifications destroy, and the sign lies in that property. For the three objects it is, respectively,
the sign of the Schur complement onto the collar in the decomposition of Lemma~\ref{lem:wh-schurabs} (the exterior field is the coupling
functional there; argued, since the collar space is infinite-dimensional and the lemma is stated for finite-dimensional $E$); the inequality
$s_w^\top\mathbb A_Ms_w\ge0$ on the family, which follows from \eqref{eq:target} by Proposition~\ref{prop:AM}(i) and would be equivalent to it
if the minimizing vectors of $\mathbb A_M$ lay in the closed span of the family for every $M$ (observed numerically, not proved); and nothing,
for Proposition~\ref{prop:shape-phase}, which rewrites \eqref{eq:symbol}. The statement that every expression for $\sT$ obtained by the
functional equation, the Euler product, Mellin inversion, Poisson summation or Riemann--Siegel identities is affine in $(\Lambda(n))_{n<e^{2T}}$
is a tautology, since $\sT$ is affine in this vector by definition (Proposition~\ref{prop:finite}(iv)); it is not a barrier, and it does not bear
on Corollary~\ref{cor:shape}(a), which concerns the use of the vector in the deduction, not the symbol. Proposition~\ref{prop:shape-family}(iii)
does not show that the family is small at the low ordinates: at bounded height the term $M^{1/2-i\tau}\mathcal W(\frac12-i\tau)$ has size
$M^{1/2}|\mathcal W(\frac12-i\tau)|$ and dominates. The observation that the minimizing vectors of $\mathbb A_M$ have Dirichlet polynomials
that are small at the low ordinates therefore remains numerical; if it holds, (iii) reads it as a property of the cutoff,
$\mathcal W(\frac12-i\gamma)\approx0$. Read through the Poisson summation formula on $\Z^2$, with the boundary terms at $y=1$ kept, the
nonzero dual frequencies of the atom-$n$ sum are again near-coincidences $\ell_2\approx-\ell_1n$; we do not use this.
\end{remark}

\subsection{The theorem}

\begin{theorem}[the shape of a proof]\label{thm:shape}
Let $\mathcal A$ be a deduction class with property set $\mathsf P$.
\begin{enumerate}
\item[(A)] \textup{(transfer, proved)} If $\mathcal A$ is closed under one of $\mathrm O_1$, $\mathrm O_2$, $\mathrm O_3$, $\mathrm O_4$ (for
all $c>0$), $\mathrm O_6$, $\mathrm O_7$, $\mathrm O_8$, $\mathrm O_{13}$, $\mathrm O_{14}$, $\mathrm O_{15}$, in the ranges of
Lemmas~\ref{lem:shape-del}--\ref{lem:shape-scale}, \ref{prop:shape-eulerdel}--\ref{lem:shape-ft}, \ref{lem:shape-fam}--\ref{lem:shape-poisson},
or if $\mathcal A$ is one of the classes of Lemma~\ref{lem:shape-classes}(i),(ii), then $\mathcal A$ contains no valid deduction of
$\mathrm{WP}_T(\Lambda)$ at the windows named in the corresponding lemma; in particular none proving \eqref{eq:target}.
\item[(A$'$)] \textup{(transfer with a label)} The same conclusion holds for $\mathrm O_4^\varepsilon$ and $\mathrm O_5$ under the
hypotheses $\kappa_T\ne0$, respectively $c_{F_\star}(\log n_0)\ne0$ (conditional); for $\mathrm O_{10}$ for the complete form (cited) and,
for the primitive bottom, argued; for $\mathrm O_{11}$, argued; for $\mathrm O_{12}$ the falsifier is open; and for $\mathrm O_9$ the exclusion
concerns pointwise bounds on the curvature of $\ell$, not the window form.
\item[(B)] \textup{(structure)} Every valid deduction of \eqref{eq:target}, in whatever class, satisfies
Lemmas~\ref{lem:shape-margin}--\ref{lem:shape-restate}: its margin at $T\ge1$ is at most $4\cdot10^9Te^{-15.449T}$; it is exact at the
centre, at the integer resolution and at the density ceiling, and, under the hypothesis of Proposition~\ref{prop:shape-crit}(b), at every
packet frequency; it is not a symmetrization of $D_\sigma$; written in shift-multiplier calculus it is a measure identity on $(-2T,2T)$
plus Bochner positivity, and it is no arithmetic square decomposition for $T\ge\frac12\log3$; it uses the Euler product over $\N$ and the
functional equation at conductor $1$ jointly; and each statement of Lemma~\ref{lem:shape-restate} is \eqref{eq:target} itself.
\item[(B$'$)] \textup{(no estimate)} A deduction that decides the sign at $T\ge1$ by an estimate with a remainder, in the sense of
\S\ref{subsec:shape-est}, has main term and remainder equal to within $4\cdot10^9Te^{-15.449T}$, and every exact reduction by
Lemma~\ref{lem:wh-schurabs} with $\lambda_S<c_{\mathcal X}$ passes this precision to the reduced form up to the factor $(1+\|\mathsf K\|)^2$ (proved); if $\kappa_T\ne0$, an
estimate whose slack exceeds $\Aconst\lamA/|\kappa_T|$ admits the falsifier $(1+\delta)\Lambda$ (conditional). That no estimate with a
remainder carries the sign step is argued (Lemma~\ref{lem:shape-est}).
\item[(B$''$)] \textup{(no comb-generic identity, proved)} Adding to a property set an identity valid for every weight vector does not remove
closure under $\mathrm O_1$--$\mathrm O_6$; so (A) and (A$'$) apply unchanged, and the exact identities of
Propositions~\ref{prop:shape-ext}--\ref{prop:shape-phase} carry no sign by themselves (Lemma~\ref{lem:shape-ident}).
\item[(C)] \textup{(positive form)} Hence every valid proof of \eqref{eq:target} uses, at a load-bearing line, a property of $\zeta$ that
is destroyed by each modification named in (A), and, with the labels stated there, by those named in (A$'$). Corollary~\ref{cor:shape}
lists these properties.
\end{enumerate}
\end{theorem}
\begin{proof}
(A) is Lemma~\ref{lem:transfer}: each cited lemma exhibits $\mathrm O(q_0)$ with property $\mathsf P$ and a negative bottom at the named
windows, so a deduction in $\mathcal A$ would apply to $\mathrm O(q_0)$ and yield a false conclusion there. (A$'$) is the same argument,
with the falsifier's negativity carrying the label of its lemma. (B) is the conjunction of the structural lemmas. (B$'$) is Lemma~\ref{lem:shape-est}(i)--(iii) with its reading (iv),
and (B$''$) is Lemma~\ref{lem:shape-ident}. (C) is the contrapositive of (A): a proof whose property set were invariant under one of the named modifications would lie in a class closed under
it.
\end{proof}

\subsection{Corollaries}

\begin{corollary}[what a proof must be]\label{cor:shape}
A valid proof of \eqref{eq:target}:
\begin{enumerate}
\item[(a)] uses the vector $(\Lambda(n))_{n<e^{2T}}$ exactly at every $T$, not through membership in a class closed under any modification
of (A) (Lemmas~\ref{lem:shape-del}--\ref{lem:shape-scale}, \ref{prop:shape-eulerdel}), and non-linearly: a linear use is a measure identity
on the window (Lemma~\ref{lem:shape-lin}(i)) and so a certificate of class (D), whose content is $\lamP$ itself
(Lemma~\ref{lem:shape-lin}(ii), Lemma~\ref{lem:shape-restate}); and it uses a property not shared by the zero-side modifications
(Lemmas~\ref{lem:shape-quartet}, \ref{lem:shape-ft}): the quartet models keep the comb exactly, so the property is not a property of the
comb alone;
\item[(b)] uses the Euler product over $\N$, the functional equation at conductor $1$ and the integrality of the frequencies jointly
(Proposition~\ref{prop:shape-inputs}; the third input argued), hence is a computation with the data of $\zeta$ at every window, uniform in
$T$, and not a class argument over a family of objects;
\item[(c)] decides a correlation. If $F_\star$ is a real unit minimizer at $T$, then by Lemma~\ref{lem:degen}
$\lambda_{\mathrm A}(T;v')\le\lamA-\sum_n(v'_n-v_n)c_{F_\star}(\log n)$ for every weight vector $v'$: the vector
$(-c_{F_\star}(\log n))_n$ is a supergradient of the concave bottom, and the one-sided derivative in a direction is the minimum of this
expression over the minimizers. No bounded functional of the weights or of the symbol (Lemma~\ref{lem:shape-weight},
Lemma~\ref{lem:shape-margin}(ii), Remark~\ref{rem:wit-sens}), no local criterion (Lemma~\ref{lem:shape-local}) and no termwise bound
(Lemma~\ref{lem:shape-misc}) decides it; numerically, a relative change of $10^{-5}$ in the weights changes the sign at $T=\log2$ while
every coherence functional moves by less (Remark~\ref{rem:wit-sens});
\item[(d)] has no slack: Lemma~\ref{lem:shape-margin}(i) and Proposition~\ref{prop:shape-crit};
\item[(e)] does not decide the sign by replacing a prime sum with a main term and a remainder, unless the main term is exact to the
precision of Lemma~\ref{lem:shape-est}(i) (proved) and the remainder has no slack beyond Lemma~\ref{lem:shape-est}(iii) (conditional);
the general form of this item is argued;
\item[(f)] uses an exact identity, if at all, only together with a property of $\zeta$ that the modifications of (A) destroy, and the sign
lies in that property (Lemma~\ref{lem:shape-ident}, Remark~\ref{rem:shape-ident}).
\end{enumerate}
\end{corollary}
\begin{proof}
Each item restates the lemmas it cites; in (c), Lemma~\ref{lem:degen} gives the inequality for every $v'$, which is the definition of a
supergradient of $\mu_T$ at $v_T(\Lambda)$, and a concave function has one-sided directional derivatives equal to the infimum of its
supergradients in that direction.
\end{proof}

\begin{remark}[what the theorem does not say]\label{rem:shape-scope}
(i) It does not say that \eqref{eq:target} is unprovable: every clause excludes a deduction, and the proofs of $\mathrm{WP}_T$ for
$T\le\frac12\log5$ (Theorems~\ref{thm:C}, \ref{thm:log5A}, \ref{thm:log5W}), with the cited ranges of \cite{Zhu2026}, \cite{Liu2026} and \cite{Mori15},
are computations with the exact vector, inside the region of Problem~\ref{prob:region} below. (ii) No equivalence between
\eqref{eq:target} and a statement about another object is used; the equivalences of Lemma~\ref{lem:shape-restate} are internal
restatements, used only to say that none of them is weaker than \eqref{eq:target}. (iii) The clauses that are not proved are listed in
Theorem~\ref{thm:shape}(A$'$) and in Proposition~\ref{prop:shape-crit}(b); what would upgrade each is stated with it: a proof that
$\kappa_T\ne0$ or that the degenerate case of Lemma~\ref{lem:degen} does not occur; a primitive witness for $\Xi_Q$; an explicit formula
with compactly supported tests for one Beurling system; a window and a field with a negative primitive bottom; an $m$-fold twist with
a window witness; for Theorem~\ref{thm:shape}(B$'$), a proof that $\kappa_T\ne0$ and a property set that makes an estimate a deduction class. (iv) The theorem does not rank what remains; it describes the region in which every proof must lie, and that region
is stated as Problem~\ref{prob:region}.
\end{remark}


\subsection{Geometry of decisive premises}\label{subsec:shape-geom}

This subsection records tools on the set of weight vectors at which the window form is nonnegative, and on what a
premise about the comb can and cannot contribute to the sign. Throughout, $T$ is fixed, $v^0:=v_T(\Lambda)$ with
$v_n(a)=2a(n)/\sqrt n$ (\S\ref{subsec:finite}), $\mu_T$ is the concave bottom of Lemma~\ref{lem:degen}, so $\mu_T(v^0)=\lamA$,
and all statements are in the coordinates $v\in\R^{\NT}$; by Proposition~\ref{prop:finite}(i) nothing else of a comb enters.
For combs $a,b$ put $d_v(a,b):=\sum_{n\in\NT}|v_n(a)-v_n(b)|=2\sum_{n\in\NT}|a(n)-b(n)|/\sqrt n$. By Proposition~\ref{prop:wit-sens},
$|\lambda_{\mathrm A}(T;a)-\lambda_{\mathrm A}(T;b)|\le d_v(a,b)$.

A \emph{premise set} at $T$ is a set $S\subset\R^{\NT}$ with $v^0\in S$; it is \emph{decisive} if $\mu_T\ge0$ on $S$. A deduction
that uses about the comb only the membership of its weight vector in $S$ and concludes $\mathrm{WP}_T$ is valid only if $S$ is decisive
(Lemma~\ref{lem:transfer}). Put $K_T:=\{v:\mu_T(v)\ge0\}$; this is the set $\mathcal V_T$ of Proposition~\ref{prop:finite}(ii), closed and
convex, and $S$ is decisive if and only if $S\subset K_T$.

\begin{corollary}[the supporting half-space; (i) proved, (ii) conditional]\label{cor:geo-half}
\emph{(i)} For every real unit minimizer $F_\star$ of $\mathfrak a_T$ on $\PT$, with $w(F_\star):=(c_{F_\star}(\log n))_{n\in\NT}$,
\[
K_T\ \subset\ H_T(F_\star):=\Big\{v:\ \sum_{n\in\NT}(v_n-v^0_n)\,c_{F_\star}(\log n)\le\lamA\Big\},
\]
so every decisive premise set lies in $\bigcap_{F_\star}H_T(F_\star)$, the intersection over all real unit minimizers.
\emph{(ii)} If $w(F_\star)\ne0$, every $v$ in a decisive premise set satisfies
$\langle v-v^0,w(F_\star)\rangle/\|w(F_\star)\|_2\le\lamA/\|w(F_\star)\|_2$, which for $T\ge1$ is at most
$4\cdot10^9Te^{-15.449T}/\|w(F_\star)\|_2$: a decisive premise set may have codimension one, but in the direction $w(F_\star)$ it reaches
beyond $v^0$ by at most the margin.
\end{corollary}
\begin{proof}
(i) Lemma~\ref{lem:degen} gives $\mu_T(v)\le\lamA-\langle v-v^0,w(F_\star)\rangle$; on $K_T$ the left side is $\ge0$. Real unit minimizers
exist (Remark~\ref{rem:resweights}). (ii) is (i) divided by $\|w(F_\star)\|_2$, with Corollary~\ref{cor:margin1}.
\end{proof}

Whether $w(F_\star)\ne0$ for some minimizer is not known; if $w(F_\star)=0$ for every real unit minimizer, $H_T(F_\star)=\R^{\NT}$,
part (ii) is empty, and Lemma~\ref{lem:degen} describes that case.

\paragraph{Premises robust in total variation.}
The next statement concerns premises on combs in the sense of Definition~\ref{def:comb}. A set $\Pi$ of combs is \emph{robust of
radius $r>0$ at $T$} if it contains every comb $\nu'$ with $\nu'=\nu_\Lambda$ off $(0,2T)$ and $|\nu'-\nu_\Lambda|((0,2T))\le r$. An
estimate of \S\ref{subsec:shape-est} stated on combs with slack $s$ (the stated remainder exceeding the true error by $s$ for every unit
$F$) is robust of radius $s$, because $|c_F|\le\|F\|^2$; so is any finite list of premises $X_j(\nu)\in[\alpha_j,\beta_j]$ with $X_j$
$L_j$-Lipschitz in total variation and slacks $s_j$, with $r=\min_js_j/L_j$.

\begin{theorem}[total-variation precision; proved]\label{thm:geo-tv}
If $\Pi$ is robust of radius $r$ at $T$ and $b_T(\nu)\ge0$ for every $\nu\in\Pi$, then $r\le\lamA$; for $T\ge1$, $r\le4\cdot10^9Te^{-15.449T}$.
\end{theorem}
\begin{proof}
Let $\eta>0$. Choose a unit $F\in\PT\cap\mathcal D_T$ with $\mathfrak a_T(F)<\lamA+\eta$, and then, since $u\mapsto c_F(u)$ is continuous with
$c_F(0)=1$, some $u\in(0,2T)$ with $\Rel c_F(u)\ge1-\eta$. For $m>0$ the comb $\nu'=\nu_\Lambda+m\delta_u$ has
$b_T(\nu')\le\mathfrak a^{\nu'}_T(F)=\mathfrak a_T(F)-m\Rel c_F(u)<\lamA+\eta-m(1-\eta)$, and $|\nu'-\nu_\Lambda|((0,2T))=m$. If $r>\lamA$, take
$m\in(\lamA,r]$ and $\eta$ with $m(1-\eta)>\lamA+\eta$: then $\nu'\in\Pi$ and $b_T(\nu')<0$, a contradiction. The bound for $T\ge1$ is
Corollary~\ref{cor:margin1}.
\end{proof}

The falsifying atom sits at a point $u$ near $0$, in general not of the form $\log n$. A single exact support premise (for instance
$\nu((0,\log2))=0$, or $\operatorname{supp}\nu\subset\log\N$) removes it, so Theorem~\ref{thm:geo-tv} does not apply to classes that mix
estimates with an exact support axiom. The estimate notion of \S\ref{subsec:shape-est} is stated on weight vectors $v\in\R^{\NT}$,
not on measures; for it the relevant statement is the following.

\begin{proposition}[integer-supported premises; (i), (iii) proved, (ii) conditional]\label{prop:geo-int}
Suppose $2T>\log2$. For $\eta>0$ put
\[
\hat\kappa_T(\eta):=\sup\Big\{\sum_{n\in\NT}v^0_n\,|c_F(\log n)|:\ F\in\PT\cap\mathcal D_T\ \text{real},\ \|F\|=1,\ \mathfrak a_T(F)<\lamA+\eta\Big\}.
\]
\emph{(i)} $\hat\kappa_T(\eta)>0$, and if a decisive premise set contains every $v$ with $|v_n-v^0_n|\le r\,v^0_n$ for all $n$
(relative slack $r\in(0,1]$), then $r\,\hat\kappa_T(\eta)\le\lamA+\eta$ for every $\eta>0$.
\emph{(ii)} If $h$ is a real unit minimizer, the same premise set satisfies $r\sum_nv^0_n|c_h(\log n)|\le\lamA$. This gives a
quantitative bound as soon as $c_h(\log n)\ne0$ for some $n\in\NT$, a weaker hypothesis than $\kappa_T\ne0$ in
Lemma~\ref{lem:shape-est}(iii).
\emph{(iii)} $\hat\kappa_T(\eta)\ge\sup\{\mT(F)\}-\lamA-\eta$, the supremum over the same set of $F$.
\end{proposition}
\begin{proof}
(i) For such $F$ let $s_n=\operatorname{sgn}c_F(\log n)$ and $v_n=v^0_n(1+rs_n)\ge0$. Then $\mu_T(v)\le\mathfrak a^{(v)}_T(F)=\mathfrak a_T(F)-r\sum_nv^0_n|c_F(\log n)|$;
decisiveness gives $r\sum_nv^0_n|c_F(\log n)|\le\lamA+\eta$, and the supremum over $F$ gives the inequality. Positivity: the set of
admissible $F$, together with its positive multiples, is open in the form domain of $\PT\cap\mathcal D_T$, and $F\mapsto c_F(\log2)$ is a
continuous quadratic form there; a quadratic form vanishing on a nonempty open cone vanishes identically, while
Lemma~\ref{lem:cpos} gives $F\in\PT\cap C_c^\infty(\IT)$ with $c_F(\log2)>0$. Since $v^0_2=\log2>0$, $\hat\kappa_T(\eta)>0$. (ii) is the same
computation with $F=h$, $\mathfrak a_T(h)=\lamA$. (iii) $\sum_nv^0_n|c_F(\log n)|\ge\sum_nv^0_nc_F(\log n)=\mT(F)-\mathfrak a_T(F)$.
\end{proof}

The qualitative bound of (i) is proved for every window. A quantitative lower bound $\hat\kappa_T(\eta)\ge k_0(T)>0$, with which (i) becomes
$r\le(\lamA+\eta)/k_0(T)$, is open; by (iii) it would follow from a lower bound for $\mT$ at near-minimizers. At a fixed window this is a
finite certified computation; uniformly in $T$ it is a statement on the structure of minimizers.

\paragraph{Functional-equation defects of Euler modifications.}
For a finite set $P$ of primes and real power series $R_p(x)=\sum_{k\ge0}r_{p,k}x^k$ with $r_{p,0}=1$ and
$\sum_k|r_{p,k}|p^{-k/2}<\infty$, let $a$ be the comb with $\zeta_a=\zeta R$, $R(s):=\prod_{p\in P}R_p(p^{-s})$ (notation of
Theorem~\ref{thm:finmod}); writing $\log R_p(x)=\sum_k\ell_{p,k}x^k$, $a(p^k)=\Lambda(p^k)+k\log p\,\ell_{p,k}$ and $a=\Lambda$ off the
powers of $P$. The completed function $\xi_a:=\xi R$ satisfies $\xi_a(s)=\xi_a(1-s)$ if and only if $R(s)=R(1-s)$, and, the coefficients
being real, if and only if $R(\frac12+it)\in\R$ for all $t$. Its \emph{Bohr defect} is
$\Delta(a):=\lim_{X\to\infty}\frac1{2X}\int_{-X}^X|\Iml R(\tfrac12+it)|^2dt$.

\begin{proposition}[exact Bohr defect; proved]\label{prop:geo-bohr}
$\Delta(a)=\frac12\big(\prod_{p\in P}\sum_{k\ge0}r_{p,k}^2p^{-k}-1\big)$. In particular $\Delta(a)=0$ if and only if $a=\Lambda$, and every
factor is $\ge1$, so primes added to $P$ can only increase the defect. For the Euler deformations
$R_p(x)=(1-x)/(1-\alpha_px)$, $|\alpha_p|<\sqrt p$, for which $a(p^k)=\alpha_p^k\log p$,
\[
\Delta(a)=\frac12\Big(\prod_{p\in P}\Big[1+\frac{(\alpha_p-1)^2}{p-\alpha_p^2}\Big]-1\Big)
=\sum_{p\in P}\frac{(\alpha_p-1)^2}{2(p-1)}+O\big(|\alpha-1|^3\big).
\]
\end{proposition}
\begin{proof}
$|\Iml R|^2=\frac12(|R|^2-\Rel R^2)$. On the critical line $R=\prod_pR_p(p^{-1/2}u_p)$ with $u_p=p^{-it}$, an absolutely convergent
series in the frequencies $\sum_pk_p\log p$. Since the $\log p$ are linearly independent over $\Q$, the Bohr mean of such a series is
its constant term, which is the Haar mean over the torus $(u_p)_{p\in P}$, and the mean of a product of functions of distinct $u_p$ is
the product of the means. $R_p(p^{-1/2}u)^2$ has constant term $r_{p,0}^2=1$, and $|R_p(p^{-1/2}u)|^2$ has mean $\sum_kr_{p,k}^2p^{-k}$ by
Parseval. For the deformations, $r_{p,k}=(\alpha_p-1)\alpha_p^{k-1}$ for $k\ge1$, and the geometric sum gives the factor.
\end{proof}

The closed form was also compared with direct Bohr averages for one and two primes (several $\alpha_p$ of both signs), with agreement
to ten digits; the proof does not use this.

\begin{proposition}[defect premises are blind at first order; reduction, conditional on $w(F_\star)\ne0$]\label{prop:geo-blind}
Let $(a_s)_{0\le s<s_0}$ be combs with $a_0=\Lambda$ and $s\mapsto v_T(a_s)$ differentiable at $0$ with derivative $d$, let $\Delta\ge0$ be any
functional with $\Delta(\Lambda)=0$ and $\Delta(a_s)\to0$ as $s\downarrow0$, and suppose $\langle d,w(F_\star)\rangle>0$ for a real unit
minimizer $F_\star$. For $\varepsilon>0$ put $s(\varepsilon):=\sup\{s:\Delta(a_{s'})\le\varepsilon\text{ for all }s'\le s\}>0$. If
$\{v_T(a):\Delta(a)\le\varepsilon\}$ is decisive, then $\lamA\ge\sup_{s<s(\varepsilon)}\langle v_T(a_s)-v^0,w(F_\star)\rangle
=s(\varepsilon)\langle d,w(F_\star)\rangle(1+o(1))>0$ as $\varepsilon\to0$.
For the Euler deformations $\alpha_p=1+t_p$, $p<e^{2T}$, put $W_p:=\sum_{k\ge1,\,p^k\in\NT}2k\log p\,p^{-k/2}c_{F_\star}(k\log p)$; then to
leading order in $\varepsilon$ the decisive premise $\Delta\le\varepsilon$ forces $\lamA\ge\big(2\varepsilon\sum_p(p-1)W_p^2\big)^{1/2}$.
\end{proposition}
\begin{proof}
For $a$ in the premise set, Lemma~\ref{lem:degen} and decisiveness give $0\le\mu_T(v_T(a))\le\lamA-\langle v_T(a)-v^0,w(F_\star)\rangle$.
Apply this to $a_s$, $s<s(\varepsilon)$, and expand $v_T(a_s)=v^0+sd+o(s)$. For the deformations,
$v_{p^k}(a)-v^0_{p^k}=2(\alpha_p^k-1)\log p\,p^{-k/2}=2kt_p\log p\,p^{-k/2}+O(t^2)$, so the pairing is $\sum_pt_pW_p+O(|t|^2)$; by
Proposition~\ref{prop:geo-bohr} the premise is the ellipsoid $\sum_pt_p^2/(2(p-1))\le\varepsilon$ to leading order, on which the maximum of
$\sum_pt_pW_p$ is $(2\varepsilon\sum_p(p-1)W_p^2)^{1/2}$ by Cauchy--Schwarz.
\end{proof}

So a premise of the form ``the defect of the functional equation is at most $\varepsilon$'', with any defect continuous along one curve
of positive pairing, decides $T$ only together with a strictly stronger, quantitative form of $\mathrm{WP}_T(\Lambda)$; it supplies no
sign of its own. The hypothesis is that of Corollary~\ref{cor:geo-half}(ii); in the degenerate case of Lemma~\ref{lem:degen} the
statement is empty. Continuity is all that is used: a premise that escapes it must be bounded below on the side of positive pairing
arbitrarily near $\Lambda$, that is, it must contain a half-space condition in the direction $w(F_\star)$ itself.

\begin{theorem}[splitting and continuation; proved]\label{thm:geo-split}
Let $\mathcal E$ be a class of combs containing $\Lambda$, $\Delta:\mathcal E\to[0,\infty]$ any functional with $\Delta(\Lambda)=0$,
$\mu:=\lamA$ and $C>0$. Consider
\[
(\mathrm I_{T;C})\qquad \lambda_{\mathrm A}(T;a)\ge-C\,\Delta(a)\quad\text{for all }a\in\mathcal E.
\]
\emph{(i) Splitting.} $(\mathrm I_{T;C})$ implies $\mu\ge0$ and $\Delta(a)\ge C^{-1}\big(\langle v_T(a)-v^0,w(F_\star)\rangle-\mu\big)$ on $\mathcal E$ for
every real unit minimizer $F_\star$; conversely $\mu\ge0$ together with $\Delta(a)\ge C^{-1}(d_v(a,\Lambda)-\mu)_+$ on $\mathcal E$ implies
$(\mathrm I_{T;C})$. The sign in $(\mathrm I_{T;C})$ is the sign at $\Lambda$; the rest is a rigidity property of $\Delta$ that carries no sign.
\emph{(ii) Continuation.} Let $(a_s)_{s\in[0,1]}\subset\mathcal E$ with $s\mapsto v_T(a_s)$ continuous, $a_1=\Lambda$. Information of the form
$\lambda_{\mathrm A}(T;a_s)\ge\varphi(s)$ for $s\le s_0$ yields $\mu\ge\max_{s\le s_0}[\varphi(s)-d_v(a_s,\Lambda)]$, and no more among
$1$-Lipschitz functions of $v$: the function $g(v)=\max_{s\le s_0}[\varphi(s)-\|v-v_T(a_s)\|_1]$ is $1$-Lipschitz, $\ge\varphi$ on the data,
and equals that bound at $v^0$. In particular positivity at $a_0$ transfers to $\Lambda$ only if $\lambda_{\mathrm A}(T;a_0)>d_v(a_0,\Lambda)$.
\emph{(iii)} On the segment $v_T(a_s)=(1-s)v_T(a_0)+sv^0$, concavity gives $\lambda_{\mathrm A}(T;a_s)\ge(1-s)\lambda_{\mathrm A}(T;a_0)+s\mu$, so interior
values bound $\mu$ only from above. \emph{(iv)} If $(\mathrm I_{T;C})$ holds on $\{a_s:s<1\}$ and $\Delta(a_s)\to0$, then $\mu\ge0$ by the Lipschitz
bound, without using $a_0$.
\end{theorem}
\begin{proof}
(i) Evaluate at $a=\Lambda$; combine $(\mathrm I_{T;C})$ with $\lambda_{\mathrm A}(T;a)\le\mu-\langle v_T(a)-v^0,w(F_\star)\rangle$ (Lemma~\ref{lem:degen}).
Conversely $\lambda_{\mathrm A}(T;a)\ge\mu-d_v(a,\Lambda)\ge-C\Delta(a)$. (ii) $\mu\ge\lambda_{\mathrm A}(T;a_s)-d_v(a_s,\Lambda)\ge\varphi(s)-d_v(a_s,\Lambda)$;
the properties of $g$ are immediate. (iii) is concavity of $\mu_T$ (Lemma~\ref{lem:degen}). (iv) $\mu\ge\lambda_{\mathrm A}(T;a_s)-d_v(a_s,\Lambda)\ge
-C\Delta(a_s)-d_v(a_s,\Lambda)$, and $d_v(a_s,\Lambda)\to0$ by continuity of the path.
\end{proof}

Optimality of the bound in (ii) among concave $1$-Lipschitz functions is argued, not proved ($g$ is not concave). For finitely many
deformed primes and the defect $\sup_{|t|\le\tau}|\log(\xi_a/\xi_a(1-\cdot))(\frac34+it)|$ of the deformations
$\zeta_a=\prod_p(1-p^{-s})^{-\theta_p}$, the rigidity in (i) holds with some finite $C$, by injectivity of $\theta\mapsto\log(\xi_a(s)/\xi_a(1-s))$
(its coefficient at the frequency $\log p$ on $\Rel s=\frac34$ is $-(1-\theta_p)p^{-3/4}$) and finite dimension; so on that class
$(\mathrm I_{T;C})$ holds exactly when $\mathrm{WP}_T$ holds, in particular for $T\le\frac12\log5$ by Theorem~\ref{thm:log5A}.
With infinitely many deformed primes the rigidity is open.

\paragraph{The negative index along the window.}
Let $n_-(T)$ be the number of negative eigenvalues, with multiplicity, of the self-adjoint operator of $\mathfrak a_T$ on $\PT$
(lower bounded, with compact resolvent, Lemma~\ref{lem:compact}), and $z(T)$ the dimension of its kernel.

\begin{theorem}[monotone index; proved]\label{thm:geo-index}
\emph{(W1)} $T\mapsto n_-(T)$ and $T\mapsto n_-(T)+z(T)$ are nondecreasing. Hence the set $\{T:\lamA<0\}$ is an interval $(T_*,\infty)$
(possibly empty), and every crossing of an eigenvalue through $0$ as $T$ increases is downward: the spectral flow on $[T_1,T_2]$ is
$n_-(T_2)-n_-(T_1)\ge0$. The same holds for $\mathfrak a^{(v)}_T$ for every weight vector $v$ (with $v_n$ given for all $n\ge2$).
\emph{(W1$'$)} There is $\varepsilon>0$ with $n_-(T)=z(T)=0$ for $0<T<\frac12\log5+\varepsilon$.
\emph{(W2)} Let $v(s)=(1-s)v^1+sv^0$ with $\mu_T(v^1)>0$. Then $s\mapsto\mu_T(v(s))$ is concave on $[0,1]$, and: if $\mu_T(v(s'))=0$ for some
$s'<1$, then $\lamA\le-\mu_T(v^1)(1-s')/s'<0$; if $\lamA\ge0$, then $\mu_T(v(s))>0$ for all $s<1$; if $\lamA<0$, there is exactly one
$s_*\in(0,1)$ with $\mu_T(v(s_*))=0$, and $\mu_T(v(s))<0$ on $(s_*,1]$. So the absence of a kernel before the endpoint is equivalent to
$\mathrm{WP}_T$.
\end{theorem}
\begin{proof}
(W1) By min--max, $n_-(T)$ (resp.\ $n_-(T)+z(T)$) is the largest dimension of a subspace of the form domain in $\PT$ on which the
form is negative (resp.\ nonpositive). Extension by zero maps $\PT$ into $\mathcal P_{T'}$, $T<T'$, isometrically with the same form
values (proof of Lemma~\ref{lem:monoT}, which uses only that $c_F$ vanishes on $|u|\ge2T$ and so holds for every $v$); a subspace of
either kind at $T$ is one at $T'$. Monotonicity of $\lamA$ is Lemma~\ref{lem:monoT}. (W1$'$) Theorem~\ref{thm:log5A} gives
$\lamA\ge2.5\cdot10^{-12}$ at $T=\frac12\log5$, and Lemma~\ref{lem:contA} extends $\lamA>0$ to a neighbourhood. (W2) Concavity is
Lemma~\ref{lem:degen}; slopes of a concave function decrease, which gives the first claim; the second is
$\mu_T(v(s))\ge(1-s)\mu_T(v^1)+s\lamA$; the third is continuity and concavity.
\end{proof}

The $\varepsilon$ of (W1$'$) is not explicit; on $(\frac12\log5,\frac12\log7)$ no contact enters ($\Lambda(6)=0$), so only the growth of the
window acts there.

\begin{remark}[codimension of a crossing; (a), (b) proved in finite dimension, (c) argued]\label{rem:geo-codim}
(W3) (a) In a one-parameter real symmetric family, a simple eigenvalue reaching $0$ is one real condition, so a crossing has
codimension one; the von Neumann--Wigner count (codimension two for real symmetric families) concerns the coincidence of two
eigenvalues and only predicts that a first crossing, if any, has a one-dimensional kernel, of definite parity since $J$ commutes
with the form. (b) The candidate extra linear conditions on a kernel vector $h$ that could raise the codimension (the Tate moments,
already in $\PT$; complementary slackness of a dual certificate, which constrains the certificate rather than $h$; the vanishing
$c_h(\log n)=0$ of Lemma~\ref{lem:degen}, which is not forced) give none. (c) A genericity count says nothing about the single
family $T\mapsto\mathfrak a_T$, along which a touching is one-sided by (W1). All of (W1)--(W3) hold verbatim for weight vectors that
do cross $0$ (finite truncations, Proposition~\ref{prop:finsupp}; deletions, Theorem~\ref{thm:finmod}), so no structure of this kind
excludes a crossing without reading the values $v_n$: at a simple bottom the only such datum is the derivative
$-\sum_n\dot v_n\,c_{h}(\log n)$ of Lemma~\ref{lem:degen}.
\end{remark}

\paragraph{Parity invariance of distributional defects.}
Keep the finite Euler modifications $R=\prod_{p\in P}R_p(p^{-s})$ above, with real $R_p$ convergent on the closed discs used below. For
$Q\subset P$ let $\sigma_Q$ replace $R_p(x)$ by $R_p(-x)$ for $p\in Q$; on $d:=a-\Lambda$ it acts by $d(p^j)\mapsto(-1)^jd(p^j)$ for $p\in Q$.
A \emph{distributional defect} is a functional of the joint Bohr distribution of the process
$t\mapsto(R(\sigma_i+it),R(1-\sigma_i+it))_{i\le m}$ on finitely many vertical lines; this includes the Bohr defect, $L^q$ means, suprema
over $t$, the closure of the value set, extrema, lagged correlations and counts of equioscillation, and $\sup_t|\log(\xi_a/\xi_a(1-\cdot))|$
on a line.

\begin{theorem}[parity invariance; proved]\label{thm:geo-parity}
Every distributional defect $D$ satisfies $D(a)=D(\sigma_Qa)$. Consequently, if a premise set $S$ of such modifications is closed under
$\sigma_P$ and is decisive, then writing $d=d_{\mathrm e}+d_{\mathrm o}$ for the parts of $v_T(a)-v^0$ on even and odd prime-power levels,
$\langle d_{\mathrm e},w(F_\star)\rangle+|\langle d_{\mathrm o},w(F_\star)\rangle|\le\lamA$ for every $a\in S$ and every real unit minimizer
$F_\star$: on odd levels (in particular on every prime) a premise of the form $D\le\varepsilon$ on a $\sigma_P$-closed class is symmetric.
\end{theorem}
\begin{proof}
On each line the vector $(R(\sigma_i+it),R(1-\sigma_i+it))_i$ is a continuous function of $u=(p^{-it})_{p\in P}$ on the finite torus, and
$u_p\mapsto-u_p$ ($p\in Q$) turns it into the corresponding vector for $\sigma_Qa$; the sign is real, so it acts compatibly at $s$ and
$1-s$. By Kronecker's theorem the Bohr distribution, and the closure of the set of values, are the image of Haar measure (resp.\ of the
torus), and $u_p\mapsto-u_p$ preserves Haar measure. The second statement is Corollary~\ref{cor:geo-half}(i) applied to $a$ and $\sigma_Pa$,
whose odd-level parts are opposite.
\end{proof}

The class of all finite real modifications is $\sigma_P$-closed. The class with nonnegative Dirichlet coefficients is not
($R_p=1+2x$ goes to $1-2x$); there the theorem applies to those single-level moves near $\Lambda$ that $\sigma$ preserves. Not covered
are functionals that are not distributional (evaluations at a fixed $t$, for instance on the real axis, where $u=1$ is moved to $-1$)
and functionals that couple $R$ to $\zeta$ itself. Order constraints such as $0\le b_n\le1$ do break the symmetry; near $\Lambda$ they act
as a polyhedral cone of linear conditions on the comb and reduce a decision to finitely many window forms of the same kind.

\begin{remark}[a classification of uniform premises; argued, partly conditional]\label{rem:geo-trich}
Let a deduction at $T\ge0.60$ use a premise set $S$ given by a finite conjunction of properties of the types met in
\S\ref{sec:shape} (comb-generic identities, cones on $\log\zeta_a$, coefficientwise integrality, Euler products, local functional
equations, sparse-stable asymptotics, statements on zeros), with a conclusion that is either independent of the comb or linear in it.
Then at least one of the following holds: (i) $S$ contains a falsifier of \S\ref{sec:shape} with negative bottom at $T$; (ii) the
conclusion is linear and its value at $\Lambda$ is at most $\lamP$ (Proposition~\ref{prop:weak}, Theorem~\ref{thm:D1}), which is a
reduction to the dual program, not an obstruction; (iii) $S$ fixes $v^0$, and the deduction's statement is $\mathrm{WP}_T$ itself. The
reduction to (i) uses the falsifiers with their labels in Theorem~\ref{thm:shape}(A),(A$'$): it is conditional where those are
(support or prefix-value premises on fewer than all window primes, through $\mathrm O_4^\varepsilon$ and $\mathrm O_5$), argued where
they are argued ($\mathrm O_{10}$ for the primitive bottom, $\mathrm O_{11}$), and argued for cones combined with a functional equation
without an Euler product; the list of property types is itself argued. Corollary~\ref{cor:geo-half} constrains what lies outside it.
\end{remark}

\section{What a sign-carrying object must satisfy}\label{sec:spec}

\begin{theorem}[no imported positivity; (i)--(ii) proved, (iii) proved if $\kappa_T\ne0$, (iv) argued]\label{thm:noimport}
Let $(g,a)\mapsto O_{g,a}$ be a construction defined by one formula for every weight datum $a$ of the degree-one theory
$\mathfrak T_1$ (functional equation, $\Gamma$-factor, bounded real coefficients, residue $\ge0$, no Euler product; see
Lemma~\ref{lem:transfer}), let $H(O)\Rightarrow P(O)\ge0$ be a positivity theorem, and suppose
$\mathfrak a_T^a(g)\ge\Phi(P(O_{g,a}))$ with $\Phi$ nondecreasing and $\Phi\ge0$ on $[0,\infty)$.
(i) If this bridge is a theorem of $\mathfrak T_1$, then for $T\ge\frac95$ the hypothesis $H(O_{g,a})$ fails for
$a=F_{\mathrm{DH}}$ at the witness of Remark~\ref{rem:FDH}; likewise in degree two, with Lemma~\ref{lem:shape-epstein}, for
$T\ge1.59$. (ii) If $H$ does not depend on $a$, no such deduction proves \eqref{eq:target}. (iii) If $H$ is affine in $a$,
its margin at $\Lambda$ on a minimizer is at most a constant times $\lamA$. (iv) Otherwise, verifying $H$ for $\zeta$ uses a
premise outside $\mathfrak T_1$ and is at least as strong as \eqref{eq:target}. Parts (i)--(ii) are instances of
Lemma~\ref{lem:transfer}, (iii) of Theorem~\ref{thm:shape}(A$'$). Consequently the sign cannot be imported from a positivity
theorem whose construction is uniform in the comb: it must be derived from a construction that uses the exact
multiplicativity of $\Lambda$ itself (phase-201 memo \texttt{CPOS\_CLASS\_001}).
\end{theorem}


Theorem~\ref{thm:shape} excludes classes of deductions by falsifiers. This section adds six further exclusions, obtained by testing
specific constructions (pair counts over $\N\times\N$, cones carried by the logarithmic derivative, coefficient integrality, additive
prime-pair statements, divisor gauges) against the window form, and collects them into one specification. Each clause carries its own
label; the specification as a whole is argued, since clauses (5) and (6) are not proved. No clause moves $\lamA$, and no object meeting
all six clauses is known.

Throughout, $N:=e^{2T}$, $\mathfrak a_T$ is the window form on $\PT\cap\mathcal D_T$, $m$ is the archimedean symbol \eqref{eq:m}, and a
\emph{comb} is as in Definition~\ref{def:shape-obj}.

\subsection{Preliminary facts}

\begin{lemma}[window collapse; proved]\label{lem:spec-collapse}
For $F\in\PT$ the autocorrelation $\cF$ is supported in $[-2T,2T]$, so $\mathfrak a_T(F)$ depends on the comb only through
$(\Lambda(n))_{n\in\NT}$. Consequently any identity $\Lambda(n)=\Phi(n)$ valid for $n<N$, however derived, yields the same value of
$\mathfrak a_T(F)$; it can change a lower bound only through what it places outside the window or through an inequality applied to its
pieces.
\end{lemma}
\begin{proof}
$\cF(u)=\int F(t)\overline{F(t-u)}\,dt$ vanishes for $|u|\ge2T$ when $\supp F\subset\IT$, and the prime form is
$\PiT(F)=-\sum_{n\in\NT}v_n\Rel\cF(\log n)$.
\end{proof}

\begin{theorem}[no atomic realization; proved]\label{thm:spec-noatom}
There is no measure $\rho$ supported on $\log\Q_{>0}$ with $\int\Rel\cF\,d\rho=\mathfrak a_T(F)$ for all $F\in\PT\cap\mathcal D_T$. In
particular no weighted pair count $\sum_{m,n}K(m,n)(mn)^{-1/2}\Rel\cF(\log(n/m))$, and no difference of two such counts, equals
$\mathfrak a_T$ on $\PT$.
\end{theorem}
\begin{proof}
Every pair count regroups by the reduced ratio $n/m=a/b$ into a measure on $\log\Q_{>0}$. If $\rho$ represented $\mathfrak a_T$, then
$\rho$ minus the representing distribution of $\mathfrak a_T$ would annihilate every even test function on $(-2T,2T)$ orthogonal to
$\cosh\frac u2$ (these are limits of autocorrelations of elements of $\PT$), hence would equal $c\cosh\frac u2$ there. The archimedean
part of $\mathfrak a_T$ has an absolutely continuous component with density proportional to $K(|u|)$, singular at $u=0$ (Theorem~\ref{thm:A}),
which is not a multiple of $\cosh\frac u2$ plus atoms.
\end{proof}

Hence an exact identity between a pair count and $\mathfrak a_T$ is necessarily a limit identity: off-diagonal atoms must converge
weakly to the archimedean density while the diagonal reproduces the comb. That limit is the split of $\mathfrak a_T$ into archimedean
and prime parts, and a positivity statement about it is \eqref{eq:target} itself. In the same setting, for $\sigma>1$,
\begin{equation}\label{eq:spec-dlog}
\sum_{n\ge2}2\Lambda(n)n^{-\sigma}F(\log n)=-\frac1{2\pi}\int_\R\Fh(t)\,\partial_\sigma\log|\zeta(\sigma+it)|^2\,dt
\end{equation}
for even $F\in C^2_c(\R)$ (proved, by Parseval for Dirichlet series): $\Lambda$ appears with exact weights only through the logarithmic
derivative of the pair-count density $|\zeta(\sigma+it)|^2$, which is non-linear in the count; continuing \eqref{eq:spec-dlog} to
$\sigma=\frac12$ is the explicit formula.

\begin{theorem}[the divisor gauge; proved]\label{thm:spec-gauge}
Let $\varphi:\N\to\R$ be finitely supported with $\varphi(d)=1$ for $d<N$, and $Q_\varphi(s):=-\sum_d\varphi(d)\mu(d)(\log d)d^{-s}$. For
$F\in\PT\cap\mathcal D_T$ with $\int|\Fh|^2(1+|\tau|)^{1/6+\varepsilon}<\infty$,
\[
\mathfrak a_T(F)=\frac1{2\pi}\int_\R|\Fh(\tau)|^2\,s_\varphi(\tau)\,d\tau,\qquad
s_\varphi(\tau):=m(\tau)-2\Rel\big[\zeta(\tfrac12+i\tau)\,Q_\varphi(\tfrac12+i\tau)\big].
\]
\end{theorem}
\begin{proof}
The comb $a_\varphi(n)=-\sum_{d\mid n}\varphi(d)\mu(d)\log d$ has Dirichlet series $\zeta Q_\varphi$ and equals $\Lambda$ on $n<N$, so by
Lemma~\ref{lem:spec-collapse} it may replace $\Lambda$. The measure $\sum_na_\varphi(n)n^{-1/2}\delta_{\log n}-Q_\varphi(1)e^{u/2}du$ on
$(0,\infty)$ has cosine transform $\Rel[(\zeta Q_\varphi)(\frac12+i\tau)]+\frac12Q_\varphi(1)(\frac14+\tau^2)^{-1}$, since
$\zeta Q_\varphi-Q_\varphi(1)/(s-1)$ is entire (Abel summation and the bound $|\zeta(\frac12+i\tau)|\ll(1+|\tau|)^{1/6+\varepsilon}$,
\cite[Ch.~V]{Titchmarsh1986}). The polar part contributes $Q_\varphi(1)(\frac14+\tau^2)^{-1}$ on $\PT$ by Lemma~\ref{lem:polarelim}, and
the two terms in $(\frac14+\tau^2)^{-1}$ cancel.
\end{proof}

\begin{theorem}[gauge invariance and the common factor; proved]\label{thm:spec-carrier}
(i) For two admissible gauges $\varphi,\varphi'$ the forms $\int|\Fh|^2s_\varphi$ and $\int|\Fh|^2s_{\varphi'}$ coincide on $\PT$; only
the pointwise bound $\lamA\ge\inf_\tau s_\varphi(\tau)$ depends on the gauge. (ii) Every comb $b$ on the window has a symbol
$m-2\Rel[\zeta(\frac12+i\tau)Q(\frac12+i\tau)]$ with the same factor $\zeta(\frac12+i\tau)$ and a Dirichlet polynomial $Q$ determined on
$d<N$ by $\mu*b$. (iii) For the comb $a_P$ of Proposition~\ref{prop:shape-eulerdel}, $\zeta_{a_P}=\zeta\prod_{p<P}(1-p^{-s})$ and the
finite Euler factor is absorbed into $Q$; so $a_P$ and $\Lambda$ have symbols with the same factor $\zeta(\frac12+i\tau)$.
\end{theorem}
\begin{proof}
(i) Both equal $\mathfrak a_T(F)$ by Theorem~\ref{thm:spec-gauge}. (ii) Möbius inversion on $n<N$ and the proof of
Theorem~\ref{thm:spec-gauge} verbatim. (iii) $\zeta_{a_P}Q^{(P)}=\zeta\cdot(\prod_{p<P}(1-p^{-s})Q^{(P)})$, and the second factor is a
Dirichlet polynomial.
\end{proof}

\begin{proposition}[absolutely convergent gauges do not certify; argued]\label{thm:spec-K2}
Let $T>\frac12\log2$ and let $\varphi$ satisfy $\varphi(d)=1$ for $d<N$ and $\sum_d|\varphi(d)|(\log d)d^{-1/2}<\infty$; this includes
every finitely supported $\varphi$. Then $\inf_\tau s_\varphi(\tau)=-\infty$ (argued; the missing step is stated after the proof). If so, the bound
$\lamA\ge\inf s_\varphi$ is void for every such gauge at every $T$, and, conditionally on that step, a gauge that certifies must satisfy $|\zeta Q_\varphi(\sigma+it)|\ll_\delta\log|t|$ on $\Rel s\ge\frac12+\delta$
for every $\delta>0$, i.e.\ must be a mollifier with $\sum_d|\varphi(d)|(\log d)d^{-1/2}=\infty$.
\end{proposition}
\begin{proof}[Argument]
Put $G=\zeta Q_\varphi$. If $\inf s_\varphi>-\infty$, then $\Rel G(\frac12+it)\le\frac12m(t)+C\le\frac12\log(|t|+2)+C'$, since $m(\tau)=\log(|\tau|/2\pi)+O(1/|\tau|)$.
On $\Rel s\ge2$, $G$ is bounded, and $Q_\varphi$ is bounded and continuous on $\Rel s\ge\frac12$. The Phragmén--Lindelöf principle for
subharmonic functions of finite order on the strip $\frac12\le\sigma\le2$, applied to $\Rel(G-Q_\varphi(1)/(s-1))-\frac12\log|s+2|$, gives
$\Rel G\le\frac12\log|t|+O(1)$ there, and the Borel--Carathéodory inequality gives $|G(\sigma+it)|\ll_\delta\log|t|$ for
$\sigma\ge\frac12+\delta$. Since $a_\varphi(2)=\log2\ne0$, $Q_\varphi\not\equiv0$; fix real $\sigma\in(\frac12,1)$ with $Q_\varphi(\sigma)\ne0$.
The proof of $\zeta(\sigma+it)=\Omega(\exp((\log t)^{1-\sigma-\varepsilon}))$ \cite[Thm.~8.12]{Titchmarsh1986} selects $t$ by Dirichlet's
theorem so that $\|t\log p/2\pi\|$ is small for all $p\le X$, with $X\to\infty$. For such $t$, $d^{-it}$ is close to $1$ for every $d\le X$
(if $\|t\log p/2\pi\|\le1/q$ for $p\le X$ then $|d^{-it}-1|\le2\pi\log d/(q\log2)$), so the part of $Q_\varphi$ supported on $d\le X$ tends to its value at $\sigma$, while the tail is bounded by
$\sum_{d>X}|\varphi(d)|(\log d)d^{-\sigma}\to0$. Thus $|G(\sigma+it)|\ge\frac12|Q_\varphi(\sigma)|\exp((\log t)^{1-\sigma-\varepsilon})$
along a sequence $t\to\infty$, a contradiction.
\end{proof}
The gap is in the last step: Theorem~8.12 of \cite{Titchmarsh1986} asserts large values of $|\zeta(\sigma+it)|$ along some sequence,
and its proof produces them in a range attached to an aligned point, not at a point where $p^{-it}\approx1$ holds simultaneously for all
$p\le X(t)$ with $X(t)\to\infty$. What is needed is a joint $\Omega$-result for $\big(\zeta(\sigma+it),(p^{-it})_{p\le X(t)}\big)$; it
is not proved here. The Phragmén--Lindelöf and Borel--Carathéodory steps are complete. Theorem~\ref{thm:spec-carrier} does not depend on
this proposition.

\begin{proposition}[the integer comb; proved]\label{prop:spec-intcomb}
The comb $a\equiv1$, with measure $\nu_\Z=\sum_{n\ge2}2n^{-1/2}\delta_{\log n}$, has on $\PT\cap\mathcal D_T$ the symbol
$g_\Z(\tau)=m(\tau)+2-2\Rel\zeta(\frac12+i\tau)$, and $\nu_\Z\notin\mathfrak C^+$. Numerically $g_\Z(0)=C_0+2-2\zeta(\frac12)=-0.45147\ldots$,
$g_\Z(1)=-0.313$, $g_\Z(2)=-0.043$, $g_\Z(3)=+0.190$, with sign change near $\tau\approx2.16$.
\end{proposition}
\begin{proof}
The symbol is Theorem~\ref{thm:spec-gauge} with $Q\equiv1$ and the atom at $n=1$ removed. Since $g_\Z$ is continuous, even, of polynomial
growth and $g_\Z(0)<0$ (with $\zeta(\frac12)=-1.4603545\ldots$ and $C_0=-5.3721834\ldots$), Lemma~\ref{lem:wit-conc} with $\gamma=1$ and
the packet argument of Theorem~\ref{thm:kappa}(ii) give windows with negative bottom. The values at $\tau=1,2,3$ are numerical and are not
used.
\end{proof}

The integer comb sits at the density ceiling of Theorem~\ref{thm:ceiling} and is still negative, though less so than its main term
$\rho_1$ ($g_1(0)=C_0+4=-1.37$). Writing the finite Ramanujan expansion of $\Lambda$ on $n<N$, its $q=1$ component is $\nu_\Z$ and the
components $q\ge2$ are pole-free; so a proof must draw positive contribution at small $\tau$ from the pole-free part. That statement,
read as a Toeplitz inequality on $\PT$, is \eqref{eq:target} again, and the split depends on the gauge.

\subsection{Identities on the strip}\label{subsec:strip}

This subsection records exact identities that express $\mathfrak a_T$ through holomorphic data on the strip $\Rel s>\frac12$, and
the limits of the methods built on them. None carries a sign. Throughout, $F\in C_c^\infty(\IT)$, $\Phi=\Phi_F$, so that
$\Phi(\frac12+i\tau)=\Fh(\tau)$, $A(F)=\Phi(1)$ and $B(F)=\Phi(0)$. For a finite set $S$ of primes put
$\Gamma_\R(s)=\pi^{-s/2}\Gamma(\frac s2)$, $\xi_S(s)=\Gamma_\R(s)\prod_{p\in S}(1-p^{-s})^{-1}$, $u_S=\log|\xi_S|$,
$d_S(s)=u_S(s)-u_S(1-\bar s)$ (the \emph{semilocal functional-equation defect}), $\rho_S(s)=\xi_S(s)/\xi_S(1-s)$, and
$S_N=\{p<N\}$. For a domain $\Omega$ and $f$ holomorphic on it, $\mathcal E_\Omega(f)=\iint_\Omega|f'|^2\,dA$.

\begin{lemma}[the window symbol as a normal derivative; proved]\label{lem:strip-L1}
$2\partial_\sigma u_S(\frac12+i\tau)=m(\tau)-\sum_{p\in S,\,k\ge1}v_{p^k}\cos(k\tau\log p)$. Hence for $S\supseteq S_N$ and
$F\in C_c^\infty(\IT)$, $\mathfrak a_T(F)=\frac1{2\pi}\int_\R2\partial_\sigma u_S(\frac12+i\tau)|\Fh(\tau)|^2\,d\tau$.
\end{lemma}
\begin{proof}
$\partial_\sigma\log|\Gamma_\R(s)|=\frac12\Rel\psi(\frac s2)-\frac12\log\pi$, and $-\log|1-p^{-s}|=\Rel\sum_kp^{-ks}/k$; differentiate
at $\sigma=\frac12$. Atoms with $k\log p\ge2T$ pair with $\cF$ outside its support (Lemma~\ref{lem:spec-collapse}).
\end{proof}

\begin{theorem}[Green identity for the functional-equation defect; proved]\label{thm:strip-L3}
Let $S\supseteq S_N$ and $\Omega=\{\frac12<\sigma<2\}$. Then
\[
2\pi\,\mathfrak a_T(F)=4\iint_\Omega d_S|\Phi'|^2\,dA-2\pi(1+|S|)|A(F)|^2-2\pi\sum_{p\in S}\sum_{k\ne0}\Big|\Phi\big(1+\tfrac{2\pi ik}{\log p}\big)\Big|^2-B_2(F),
\]
with $B_2(F)=\int_{\sigma=2}(d_S\,\partial_\sigma|\Phi|^2-|\Phi|^2\partial_\sigma d_S)\,d\tau$.
\end{theorem}
\begin{proof}
On $\sigma=\frac12$, $d_S=0$ and $\partial_\sigma d_S=2\partial_\sigma u_S$. In $\Omega$, $u_S(s)$ is harmonic and $u_S(1-\bar s)$ is
harmonic except at the reflected poles of $\Gamma_\R$ (at $s=1$; the next, $s=3$, lies outside) and of the local factors (at
$s=1+2\pi ik/\log p$), where $\Delta d_S=2\pi\delta$. Apply Green's second identity to $|\Phi|^2$ (with $\Delta|\Phi|^2=4|\Phi'|^2$) and
$d_S$ on $\Omega$; the rapid decay of $\Phi$ in $\tau$ removes the horizontal edges. Conclude with Lemma~\ref{lem:strip-L1}.
\end{proof}
The sources at $s=1$ are killed exactly by $A(F)=0$, so on $\PT$ the Tate moment appears as the Green residue of the reflected pole
($B(F)$ appears in the same way from the strip $0<\sigma<\frac12$). The prime sources are not killed by $\PT$; by Parseval on the
lattice $\log p\cdot\Z$,
$\sum_{k\in\Z}|\Phi(1+2\pi ik/\log p)|^2=\log p\int_0^{\log p}\big|\sum_{j\in\Z}g(t+j\log p)\big|^2dt$ with $g(t)=F(t)e^{t/2}$.
The identity holds verbatim with Beurling primes in place of $S$.

\begin{theorem}[a sourceless energy identity; proved]\label{thm:strip-R1}
Let $S\supseteq S_N$, $\sigma_1>\frac12$ and $\Omega_1=\{\frac12<\sigma<\sigma_1\}$. Then
\[
\mathfrak a_T(F)=\frac1\pi\big[\mathcal E_{\Omega_1}(\Phi)-\mathcal E_{\Omega_1}(\rho_S\Phi)\big]
+\frac1{4\pi}\int_{\sigma=\sigma_1}\partial_\sigma\big(|\rho_S\Phi|^2-|\Phi|^2\big)\,d\tau .
\]
\end{theorem}
\begin{proof}
For $f$ holomorphic on $\overline\Omega_1$ with rapid decay in $\tau$, $\int_{\sigma=1/2}\partial_\sigma|f|^2=\int_{\sigma=\sigma_1}\partial_\sigma|f|^2-4\mathcal E_{\Omega_1}(f)$.
Apply it to $f=\rho_S\Phi$ and $f=\Phi$ and subtract; $\rho_S$ is holomorphic on $\Rel s>0$ with polynomial growth in $\tau$ on
vertical strips. On the line $|\rho_S|=e^{d_S}=1$, so $\partial_\sigma(|\rho_S\Phi|^2-|\Phi|^2)=|\Fh|^2\cdot2\partial_\sigma d_S=4|\Fh|^2\partial_\sigma u_S$.
Conclude with Lemma~\ref{lem:strip-L1}.
\end{proof}

\begin{remark}
A numerical evaluation at $T=0.6$, $S=\{2,3\}$, $F(t)=(1-(t/T)^2)^4$, $\sigma_1=0.85$ gives both sides equal to $-0.479431449594\ldots$,
with agreement $1.7\cdot10^{-13}$ (numerical). The interior terms of Theorem~\ref{thm:strip-R1} are functionals of
$\rho_S=\Gamma_\R(s)\Gamma_\R(1-s)^{-1}\prod_{p\in S}(1-p^{s-1})(1-p^{-s})^{-1}$, a product of local factors; the identity holds verbatim for
Beurling systems and for the comb $a_P$ of Proposition~\ref{prop:shape-eulerdel}.
\end{remark}

\paragraph{Gauges.} Let $I_1,I_2$ be holomorphic near $\overline\Omega_1$, unimodular on $\sigma=\frac12$, of polynomial growth in $|\tau|$
uniformly on $\overline\Omega_1$, with $q=I_2/I_1$ holomorphic, and put $\psi_q(\tau)=\partial_\sigma\log|q|(\frac12+i\tau)$.

\begin{theorem}[gauged identity; proved]\label{thm:strip-T1}
For every $\sigma_1>\frac12$,
$\frac1{2\pi}\int\psi_q|\Fh|^2d\tau=\frac1\pi[\mathcal E_{\Omega_1}(I_1\Phi)-\mathcal E_{\Omega_1}(I_2\Phi)]
+\frac1{4\pi}\int_{\sigma=\sigma_1}\partial_\sigma(|I_2\Phi|^2-|I_1\Phi|^2)d\tau$.
The left side equals $\mathfrak a_T(F)$ for all $F\in C_c^\infty(\IT)$ iff the Fourier transform of the distribution $\psi_q-\sT$
vanishes on $(-2T,2T)$; on $\PT$ one may add $\beta_0\gh0+\beta_1\gh1$. Only $q$ matters.
\end{theorem}
\begin{proof}
As for Theorem~\ref{thm:strip-R1}, with $f=I_2\Phi,I_1\Phi$ and $|q|=1$ on the line. Autocorrelations of $C_c^\infty(\IT)$ span, by
polarization, a dense subspace of the even test functions on $(-2T,2T)$.
\end{proof}
The admissible $\psi_q$ are thus $\sT+\nuh+\beta\cdot\widehat g$ with $\nu$ an even tempered distribution supported in
$\{|u|\ge2T\}$. This freedom strictly contains that of Definition~\ref{def:PD}, where $\nu$ is a finite measure; at the level of
values the two agree to the extent stated in Theorem~\ref{thm:strip-T4}.

\begin{corollary}[an entire gauge; proved]\label{cor:strip-T1p}
$q_T(s)=\Gamma_\R(s)\Gamma_\R(1-s)^{-1}\exp\big(\sum_{n<N}\frac{\Lambda(n)}{\log n}(n^{-s}-n^{s-1})\big)$ satisfies $\psi_{q_T}=\sT$
exactly, so $\mathfrak a_T(F)=\frac1\pi[\mathcal E_{\Omega_1}(\Phi)-\mathcal E_{\Omega_1}(q_T\Phi)]+\frac1{4\pi}\int_{\sigma=\sigma_1}\partial_\sigma(|q_T\Phi|^2-|\Phi|^2)$
for every $\sigma_1>\frac12$, with no set $S$.
\end{corollary}
\begin{proof}
On $s=\frac12+i\tau$, $n^{-s}-n^{s-1}=-2in^{-1/2}\sin(\tau\log n)$, so $|q_T|=1$ there, and
$\partial_\sigma\Rel(n^{-s}-n^{s-1})=-2n^{-1/2}\log n\cos(\tau\log n)$; also $\partial_\sigma\log|\Gamma_\R(s)/\Gamma_\R(1-s)|=m(\tau)$.
$q_T$ is holomorphic on $\Rel s>0$ with polynomial growth on vertical strips.
\end{proof}
Corollary~\ref{cor:strip-T1p} is a holomorphic rewriting of the symbol and carries no information beyond $\sT$
(Lemma~\ref{lem:spec-collapse}).

\begin{theorem}[contractive gauges; proved for $T\le\frac12\log5$]\label{thm:strip-T3}
If $q$ is admissible on $\PT$ and $|q|\le1$ on a one-sided neighbourhood of the line, then $\psi_q\le0$ and $\mathfrak a_T\le0$ on
$\PT\cap C_c^\infty(\IT)$. Hence no admissible gauge is contractive near the line for any $T\le\frac12\log5$
(Theorem~\ref{thm:log5A}); for larger $T$ the same holds conditionally on $\lamA>0$.
\end{theorem}
\begin{proof}
$\log|q|\le0$ near the line and $=0$ on it give $\psi_q\le0$; apply Theorem~\ref{thm:strip-T1}.
\end{proof}

\begin{theorem}[gauge methods reach the certificate values; proved in the declared class]\label{thm:strip-T4}
Consider the class of methods that use a sourceless identity of the form of Theorem~\ref{thm:strip-T1}, with gauges independent of
$F$, and conclude from a pointwise bound $\psi_q\ge c$.
\emph{(i)} If $\log|q|\ge c(\sigma-\frac12)$ on a one-sided neighbourhood of the line, then $\psi_q\ge c$ and $c\le\lamA$; if the
freedom $\nu$ of $q$ is a finite measure, $(\nu,\beta)$ is a certificate in the sense of Definition~\ref{def:PD} and $c\le\lamP$.
\emph{(ii)} Conversely, for every finitely atomic certificate $f=\sT+\sum_j2a_j\cos(u_j\tau)+\beta\cdot\widehat g\ge c'$ and every
$c<c'$ there is an admissible gauge $q=e^{\ell}$ with $\log|q|\ge c(\sigma-\frac12)$ near the line. By Theorem~\ref{thm:D1}, the
supremum of the values of the method over gauges with finite-measure freedom is $\lamP$.
Outside this class (gauges depending on $F$, identities with sources, bounds on the bulk energy not passing through a pointwise bound on
$\psi_q$) the conclusion is argued only.
\end{theorem}
\begin{proof}
(i) By Theorem~\ref{thm:strip-T1}, $\mathfrak a_T(F)=\frac1{2\pi}\int\psi_q|\Fh|^2\ge c\|F\|^2$ on a dense subset of $\PT$; the
second statement is Proposition~\ref{prop:weak}. (ii) Put $G(\tau)=\int_0^\tau f$ and $\ell(s)=iG(-i(s-\frac12))$. Since $m$ is
analytic on $|\Iml\tau|<\frac12$ and the cosines and $\gh j$ are entire, $\ell$ is holomorphic on $\frac12\le\sigma<1$, $\Rel\ell=0$
and $\partial_\sigma\Rel\ell=f$ on the line, and $\Rel\ell=(\sigma-\frac12)f+O((\sigma-\frac12)^3)$ with $f''$ bounded, uniformly in
$\tau$; $|e^{\ell}|$ grows polynomially since $\int_0^\tau m\ll|\tau|\log|\tau|$ enters only through $\Rel\ell$ off the line.
\end{proof}
Expansive gauges are therefore a holomorphic coordinate on the space of dual certificates: they supply the pair $(\nu,\beta)$ of
\S\ref{sec:duality} and nothing else.

\begin{proposition}[no Hermite--Biehler structure; proved]\label{prop:strip-R4}
For every finite $S$, and for every Beurling system with finitely many primes in place of $S$, $d_S(\sigma)<0$ on the real segment
$\frac12<\sigma<1$, while $d_S(\sigma+i\tau)>0$ for fixed $\sigma\in(\frac12,1)$ and $|\tau|$ large. So $|\xi_S(s)|>|\xi_S(1-\bar s)|$ fails
on $\frac12<\sigma<1$ in both directions; $\rho_S$ is neither contractive nor expansive there.
\end{proposition}
\begin{proof}
$\partial_\sigma\log\Gamma_\R(\sigma)=\frac12(\psi(\frac\sigma2)-\log\pi)<0$ on $(0,1)$, so the archimedean part of $d_S$ is negative on
$(\frac12,1)$; $\log|1-p^{\sigma-1}|-\log|1-p^{-\sigma}|<0$ since $\sigma-1>-\sigma$, for every $p>1$. At large $|\tau|$ the
archimedean part is $(\sigma-\frac12)\log(|\tau|/2\pi)+O(1)$, while the finitely many local parts are bounded.
\end{proof}

\begin{proposition}[exact band formula; proved]\label{prop:strip-Q2}
For every $p>1$ and $\sigma\ne1$,
$d_p(s):=\log|1-p^{s-1}|-\log|1-p^{-s}|=(\sigma-1)_+\log p+\sum_{k\ge1}\frac{p^{-k\sigma}-p^{-k|\sigma-1|}}{k}\cos(k\tau\log p)$, and
$d_S=d_\infty+\sum_{p\in S}d_p$. Against $|\Phi(\sigma+i\cdot)|^2$, whose Fourier transform is supported in $[-2T,2T]$, only the terms
with $p^k<N$ enter, with weight $1/k=\Lambda(p^k)/\log p^k$.
\end{proposition}
\begin{proof}
Expand both logarithms for $\sigma<1$; for $\sigma>1$ write $\log|1-p^{s-1}|=(\sigma-1)\log p+\log|1-p^{1-s}|$. Both series converge
absolutely. The support statement is Lemma~\ref{lem:spec-collapse} on the line $\Rel s=\sigma$.
\end{proof}
\begin{remark}
Expanding $\log\xi_S$ through sums over $S$-smooth integers produces sum sets of all $S$-smooth numbers; the logarithm resums them into
prime powers, so the area term of Theorem~\ref{thm:strip-L3} sees only $\Lambda(n)$, $n<N$, with an exact cutoff. The formula is
additive over places and holds verbatim for Beurling systems and for $a_P$.
\end{remark}

\begin{proposition}[place-splitting of the singular data; proved]\label{prop:strip-S1}
Read through \eqref{eq:m}, $\sT(\tau)=C_0+\int_0^\infty(1-\cos\tau u)K(u)\,du-\sum_{n\in\NT}v_n\cos(\tau\log n)$; its transform in
$u$ has archimedean singular data on $i\pi\Z$ (the poles of $K$) and prime data at the atoms $\pm\log n$. The two sets are disjoint,
and the data split additively over places.
\end{proposition}
\begin{proof}
Linearity of the transform, the partial fractions of $K$, and the definition of $\sT$.
\end{proof}

\begin{proposition}[a frozen comb; closed form proved]\label{prop:strip-S2}
Fix the contacts $\{2,3,4\}$ of the range $\frac12\log3<T\le\frac12\log5$ and let $T\to\infty$ with no further contacts. Then the
bottom of the frozen form tends to $\min_\tau s^{(3)}=s^{(3)}(0)=\psi(\frac14)-\log\pi-(\sqrt2\log2+\frac2{\sqrt3}\log3+\log2)$,
with $\psi(\frac14)=-\gamma-\frac\pi2-3\log2$. Numerically $s^{(3)}(0)=-8.314157\ldots$ (numerical).
\end{proposition}
\begin{proof}
$\Rel\psi(x+iy)$ is increasing in $|y|$ (series formula), so $s^{(3)}$ is minimal at $0$; as $T\to\infty$ the two linear constraints
defining $\PT$ do not prevent concentration of $|\Fh|^2$ at $\tau=0$ (min-max with a codimension-two constraint).
\end{proof}

\begin{remark}[a ceiling for curvature methods]
Write $\mathfrak a_T=\mathcal E_T+\int W_TF^2$, with $\mathcal E_T$ a censored jump Dirichlet form (jump kernel $K$ plus
$\frac12\sum v_n\delta_{\pm\log n}$) and $W_T$ a killing potential (proved). Any lower bound for $\lamA$ obtained by adding a
curvature-dimension or Poincaré bound for $\mathcal E_T$ to the infimum of the potential term is at most $g_T+\operatorname{ess\,inf}W_T$,
$g_T$ the spectral gap. At $T=\frac12\log3$ the constants are $g_T\le1.941$ and $\operatorname{ess\,inf}W_T=-2.534$
(numerical), so every such bound is negative.
\end{remark}

\begin{remark}[real stability]
For the Gram matrix of $\mathfrak a_T$ on a box of dilates, polarized to a multi-affine polynomial, a strictly positive-definite
$2\times2$ principal face already prevents real stability (Rayleigh criterion of \cite{Branden2007}); the face on $\{1,2\}$ is strictly
positive definite at $T=\log2$ (Theorem~\ref{thm:band}). So real-stable and hyperbolic-polynomial methods, which certify Lorentzian
signatures, do not apply to this Gram symbol.
\end{remark}

\subsection{The specification}

For a comb $a\ge0$ write $f_a=\exp\big(\sum_na(n)(\log n)^{-1}n^{-s}\big)=\sum_nb_n n^{-s}$, so that $L\cdot b=b*a$ with $L(n)=\log n$, and
$a=\Lambda$ iff $b\equiv1$.

\begin{proposition}[what a sign-carrying object must satisfy]\label{prop:spec}
An argument proving $\mathrm{WP}_T(\Lambda)$ at every $T$ (or at the windows named below) must use an object that simultaneously:
\begin{enumerate}
\item[(1)] is not linear in $\Lambda$ and not comb-generic: an identity valid for every weight vector, or for every Beurling system whose
zeta function reaches the line, carries no sign (Theorem~\ref{thm:shape}(B$''$), proved). In particular the divisor gauge
(Theorem~\ref{thm:spec-gauge}) and finite Ramanujan and Vaughan-type expansions are identities of this kind (proved, by
Lemma~\ref{lem:spec-collapse}); the exclusion of the class of positive-definite, hyperbola-method and Dirichlet-algebra arguments is
\emph{conditional} on an itemised check that each of its ingredients holds verbatim for Beurling systems. Cf.\ Problem~\ref{prob:region}(3),(11).
\item[(2)] is not a cone condition on $\log\zeta$ (proved). If every arithmetic premise has the form ``$\log\zeta\in K$'' or
``$-\zeta'/\zeta\in K'$'' on $\Rel s>1$ with $K,K'$ cones, the deduction class is closed under $\mathrm O_4$, since $\log\zeta^c=c\log\zeta$
and the comb of $\zeta^c$ is $c\Lambda$; by Lemma~\ref{lem:shape-scale} it proves nothing at any $T>\frac12\log2$ and any $T\ge0.60$.
Cf.\ Problem~\ref{prob:region}(1).
\item[(3)] is not coefficientwise integrality (proved). Let $\mathcal C_{01}$, $\mathcal C_\Z$, $\mathcal C_1$ be the combs with
$b\in\{0,1\}^\N$, $b\in\Z_{\ge0}^\N$, $b\equiv1$. Then $\mathcal C_1=\{\Lambda\}$; $\{c>0:c\Lambda\in\mathcal C_{01}\}=\{1\}$ and
$\{c>0:c\Lambda\in\mathcal C_\Z\}=\N$; and for $a\in\mathcal C_{01}$, $0\le a(n)\le\log n$ with $a(p)\in\{0,\log p\}$. A deduction from
membership in $\mathcal C_\Z$ proves nothing at $T>\frac12\log2$ (the proof of Lemma~\ref{lem:shape-scale} with $c\in\N$); from
$\mathcal C_{01}$ nothing at $T\ge0.60$ ($0\in\mathcal C_{01}$ and Corollary~\ref{cor:Tarch}; also $a_P\in\mathcal C_{01}$,
Proposition~\ref{prop:shape-eulerdel}); and $\mathcal C_1$ is identification. Cf.\ Problem~\ref{prob:region}(1).
\item[(4)] is not sparse-stable (proved). Call a property sparse-stable if it holds for $\Lambda+e$ whenever
$\#(\supp e\cap[1,X])=O(\log X)$ and $|e|$ is bounded. Asymptotic statements on $\sum_{n\le X}a(n)a(n+h)$ with main term $\gg X^\theta$,
uniform in $h$ (the Hardy--Littlewood pair and $k$-tuple conjectures, Bombieri--Vinogradov, sieve lower bounds), are sparse-stable;
$a_P$ differs from $\Lambda$ by such an $e$, has the same zeros, and fails $\mathrm{WP}_T$ on $[0.60,\frac12\log P]$
(Proposition~\ref{prop:shape-eulerdel}). So no sparse-stable property, even together with any statement about the zeros, yields
$\mathrm{WP}_T(\Lambda)$ at any $T\ge0.60$. Cf.\ Problem~\ref{prob:region}(1),(2).
\item[(5)] couples $(\Lambda(n))_{n<N}$ with the functional equation by a route other than the factor $\zeta(\frac12+i\tau)$
(\emph{open}; no such object is known). Failure for Beurling systems is necessary but not sufficient: the input must also separate
$\Lambda$ from $a_P$ on $\NT$ (proved, from Proposition~\ref{prop:shape-eulerdel}); that the coupling of prime-pair data with the
functional equation reduces to the polar and archimedean terms of the window form is \emph{argued}. The factor $\zeta(\frac12+i\tau)$
does not serve: it is common to every window comb, including $a_P$ (Theorem~\ref{thm:spec-carrier}, proved), and the only gauges that
could turn it into a certificate are mollifiers (Proposition~\ref{thm:spec-K2}, argued; conditional on a joint $\Omega$-result). That the functional
equation at conductor~$1$ together with a Dirichlet series over $\N$ determines $\zeta$ is Hamburger's theorem \cite{Hamburger1921};
it is a rigidity statement, not an inequality. The identities of \S\ref{subsec:strip} (Theorems~\ref{thm:strip-L3}, \ref{thm:strip-R1} and \ref{thm:strip-T1}) meet the wording of this clause: they couple $\Lambda(n)$, $n<N$, with the functional equation through the local factors and the reflection $s\mapsto1-\bar s$, not through $\zeta(\frac12+i\tau)$. But they are comb-generic (additive or multiplicative over places, valid verbatim for Beurling systems and for $a_P$), so they fail clause~(1), and their gauge freedom returns exactly the dual certificates (Theorem~\ref{thm:strip-T4}). Cf.\ Problem~\ref{prob:region}(2).
\item[(6)] crosses the logarithmic derivative without losing positivity (\emph{argued}). The positivities of $f$ that are carried to
$-f'/f$ in the Dirichlet algebra are exactly the infinitely divisible ones: $-f'/f$ has nonnegative coefficients iff $\log f$ does
(proved, since the coefficients differ by the factor $\log n$); this cone is a multiplicative semigroup, closed under $f\mapsto f^c$,
and not convex, while convex cones such as nonnegative coefficients, Bochner positivity or $|\zeta|^2$ are not carried (proved for the
listed cones; that no other cone crosses, beyond these and the real-zero cone, is argued). By \eqref{eq:spec-dlog}, pair counts give
$\Lambda$ only through $\partial_\sigma\log$ (proved). From the Kronecker--Hurwitz and Eichler--Selberg class-number family, linear
operations, convolutions and multiplication by $\C(s)$ never produce $-\zeta'/\zeta$ in the same-argument part, since
$-\zeta'/\zeta\notin\C(s)(\zeta)$ by the hypertranscendence of $\zeta$ \cite{Holder1887,Ostrowski1920} (proved there); for the family with
its shifted factors $\zeta(s-1)$ this is argued. Cf.\ Problem~\ref{prob:region}(3).
\end{enumerate}
\end{proposition}
\begin{proof}
(1) is Theorem~\ref{thm:shape}(B$''$) with Lemma~\ref{lem:spec-collapse} and Theorems~\ref{thm:spec-gauge}--\ref{thm:spec-carrier}, which
hold verbatim with $\mu_B,\zeta_B$ for a Beurling system whose zeta function reaches the line. (2) is Lemma~\ref{lem:shape-scale} applied
to the class described. (3): in $L\cdot b=b*a$ every term is nonnegative and the term $d=n$ is $a(n)b_1=a(n)$, so $a(n)\le b_n\log n$;
at $n=p$, $b_p\log p=a(p)$, which gives the bounds and the two orbit intersections ($b_p=c$; and $d_k(n)\in\N$ for $c=k\in\N$). For
$c=k$ large, the test function $F_+$ in the proof of Lemma~\ref{lem:shape-scale} sees only the contact $2$, so the proof applies with
$c$ restricted to $\N$. (4): for $a_P$, $e=-\Lambda\mathbf 1\{p<P\}$ is supported on at most $\pi(P)\log_2X$ integers up to $X$ with
$|e|\le\log P$, so the correlation changes by $O(\pi(P)\log P(\log X)^2)=o(X^\theta)$; take $P>N$ and apply
Proposition~\ref{prop:shape-eulerdel}. (5) and (6) are as stated, with the labels given there.
\end{proof}

\begin{remark}[scope]\label{rem:spec-scope}
Clauses (2)--(4) are instances of Theorem~\ref{thm:shape}(A) through $\mathrm O_4$, its integer sub-orbit, and $\mathrm O_6$; they use only
the proved parts of Lemma~\ref{lem:shape-scale}, never the conditional small form $\mathrm O_4^\varepsilon$. Clauses (2) and (3) name
the hypotheses an object must break: homogeneity in $\log\zeta$, and monotonicity of integrality premises (down-closed or closed along
$c\in\N$) unless they identify $\Lambda$. The first exact relation in $(1,\Lambda)$ that is not homogeneous is of degree two (Selberg's
$\Lambda L+\Lambda*\Lambda=\mu*L^2$). The open clause (5) is the one not covered by any falsifier of this paper.
\end{remark}

\section{The exterior form and the leakage lemma}\label{sec:ext}

This section rewrites the complete form on summation-map images as the complete Weil form of an aliasing tail lying
outside the window, states a uniform coercivity inequality for that exterior form as an open target, proves that the
target gives $\lamW>0$ for every $T$ with an explicit margin, and collects the structural identities that are proved about
the exterior form. Throughout, $\lambda:=e^T$, $N:=\lambda^2=e^{2T}$, and $\mu(\cdot)$ is the M\"obius function. For
functions on the additive line we use, in this section only, $\mathcal Ff(\xi)=\int_\R f(x)e^{-2\pi ix\xi}dx$. Each
statement carries its label; ``given'' names the cited input.

\subsection{The summation map, the aliasing tail and the gauge}\label{subsec:ext-defs}

\begin{definition}\label{def:ext}
A real function $f$ on $\R$ is \emph{admissible at $T$} if it is even, vanishes outside $[-\lambda,\lambda]$, is $C^2$ on
$[-\lambda,\lambda]$, and $f(0)=\mathcal Ff(0)=f(\lambda)=0$. Its \emph{leakage} is
$\mathrm{leak}(f)=\int_{|\xi|>\lambda}|\mathcal Ff|^2d\xi$. The \emph{summation map} is
\[
\Phi_f(x)=e^{x/2}\sum_{n\ge1}f(ne^x)\quad(x\in\R),\qquad Ef:=\Phi_f|_{\IT}.
\]
Since $ne^x\ge\lambda$ for $x\ge T$, $\Phi_f=0$ on $[T,\infty)$. The \emph{aliasing tail} is $\Psi_f:=\Phi_f|_{(-\infty,-T)}$;
by Poisson summation ($f$ even, $f(0)=\mathcal Ff(0)=0$), $\Phi_f(x)=e^{-x/2}\sum_{k\ge1}\mathcal Ff(ke^{-x})$, so for $x<-T$
only the out-of-band part $r_f:=\mathcal Ff|_{(\lambda,\infty)}$ enters, and in the reflected variable $y=-x>T$
\[
\check\Psi_f(y):=\Psi_f(-y)=(Zr_f)(e^y),\qquad (Zr)(w)=w^{1/2}\sum_{n\ge1}r(nw),
\]
\[
(Z_{\log}r)(w)=w^{1/2}\sum_{n\ge2}(\log n)\,r(nw)\qquad(w>\lambda).
\]
The \emph{gauge} is $\mathcal G_T=\{\varphi\in C^2(\R)\ \text{even}:\ \supp\varphi\subset[-1/\lambda,1/\lambda],\
\varphi(0)=\mathcal F\varphi(0)=0\}$, $\overline{\mathcal G}_T$ is the closure of $\{\mathcal F\varphi|_{(\lambda,\infty)}:\varphi\in\mathcal G_T\}$
in $L^2(\lambda,\infty)$, and for $g$ in the \emph{core} (real $g\in C^2(\overline{\IT})$ vanishing to second order at $\pm T$)
\[
\ell(g):=\inf\{\mathrm{leak}(f):\ f\ \text{admissible at }T,\ Ef=g\}.
\]
For a real function $\Psi$ on $\R$ with $|\Psi(y)|\le Ce^{-\frac32|y|}$ for large $|y|$ and $\int\log(2+|\tau|)|\widehat\Psi|^2<\infty$,
the \emph{complete Weil form on the line} is
\[
\mathcal W[\Psi]=\frac1{2\pi}\int_\R m(\tau)|\widehat\Psi(\tau)|^2d\tau-\sum_{n\ge2}v_n\,c_\Psi(\log n)+2A(\Psi)B(\Psi),
\]
with $m$, $v_n=2\Lambda(n)/\sqrt n$, $c_\Psi$, $A$, $B$ as in \eqref{eq:symbol}--\eqref{eq:W}; the comb runs over \emph{all}
$n\ge2$. For $\Psi$ supported in $\IT$ this is $W_T(\Psi)$. $\mathcal W$ is invariant under translations and under
$t\mapsto-t$.
\end{definition}

\begin{lemma}[M\"obius bijection; proved]\label{lem:ext-mob}
$Ef$ depends only on $f|_{[1/\lambda,\lambda]}$, and $f|_{[1/\lambda,\lambda]}\mapsto Ef$ is a bijection
$L^2([1/\lambda,\lambda],du)\to L^2(\IT,dx)$ with inverse $f(u)=\sum_{m\le\lambda/u}\mu(m)(mu)^{-1/2}g(\log(mu))$ and
$\|E\|,\|E^{-1}\|\le\sum_{m\le N}m^{-1/2}\le2\sqrt N$. Every $g$ in the core is $Ef$ for some admissible $f$, and two
admissible representatives of $g$ differ by an element of $\mathcal G_T$.
\end{lemma}
\begin{proof}
With $u=e^x$ and $H(u)=u^{-1/2}g(\log u)=\sum_{n\le\lambda/u}f(nu)$, only $nu\ge u\ge1/\lambda$ occurs;
$\sum_m\mu(m)H(mu)=\sum_kf(ku)\sum_{m\mid k}\mu(m)=f(u)$; $\|g\|_{L^2(dx)}=\|H\|_{L^2(du)}$ and $\|H(m\,\cdot)\|=m^{-1/2}\|H\|$.
For $g$ in the core the terms of the inverse switch on at $u=\lambda/m$, where $g$ vanishes to second order, so
$f\in C^2(1/\lambda,\lambda]$ with $f(\lambda)=\lambda^{-1/2}g(T)=0$; the values on $(-1/\lambda,1/\lambda)$ are free and are
chosen in $C^2$ to meet $f(0)=\int f=0$ (the gauge is infinite dimensional).
\end{proof}

\begin{theorem}[{half-line identity; proved, given \cite[Lemmas~4.2--4.4, Thm.~4.5]{Guo2026}}]\label{thm:ext-half}
For $f$ admissible at $T$, $|\Psi_f(y)|\le Ce^{3y/2}$ and
\[
W_T(Ef)=\mathcal W[\Psi_f]=\mathcal W[\check\Psi_f]
=\frac1{2\pi}\int_\R m\,|\widehat{\check\Psi_f}|^2-2\big\langle\check\Psi_f,(Z_{\log}r_f)(e^{\cdot})\big\rangle_{L^2(T,\infty)}
+2A(\check\Psi_f)B(\check\Psi_f).
\]
\end{theorem}
\begin{proof}
By \cite[Lemma~4.4]{Guo2026}, $\Phi_f$ lies in the radical of the full Weil distribution: $\mathcal W(\Phi_f,\Phi)=0$ for the
relevant test functions, with absolutely convergent integrals. By \cite[Thm.~4.5]{Guo2026}, read term by term on $\IT$,
$W_T(Ef)=-\mathcal W(Ef,\Psi_f)$, the cross term of the full form between the window piece and the tail (extended by
zero). Pairing the radical relation with $\Psi_f$ gives $\mathcal W(\Psi_f,Ef)+\mathcal W[\Psi_f]=0$; the splitting is
legitimate because both pieces converge absolutely ($Ef$ and $\Psi_f$ are bounded near $-T$ and
$\int_0^1\!\int_0^1\frac{ds\,dt}{s+t}<\infty$ controls the L\'evy cross term). Symmetry of the bilinear form for real
functions gives the first equality; reflection gives the second. For the last, expand the comb on $\check\Psi_f$: for $y>T$,
$\check\Psi_f(y+\log q)=\sum_l l^{-1/2}\tilde r(y+\log lq)$ with $\tilde r(x)=e^{x/2}r_f(e^x)$, so
$\sum_q\frac{v_q}2c(\log q)=\sum_{m,n}(mn)^{-1/2}\big(\sum_{q\mid n}\Lambda(q)\big)I(m,n)$ with
$I(m,n)=\int_T^\infty\tilde r(x+\log m)\tilde r(x+\log n)dx$, and $\sum_{q\mid n}\Lambda(q)=\log n$ (all sums absolutely
convergent since $r_f=O(\xi^{-2})$).
\end{proof}

\begin{remark}\label{rem:ext-killing}
(i) The archimedean term is $C_0\|\check\Psi\|^2+\Lev(c_{\check\Psi})$ (Theorem~\ref{thm:A}, whose proof does not use the
support). Written as a double integral over the half-line only, it acquires the killing term
$\int_{y>T}|\check\Psi(y)|^2\kappa(y)dy$, $\kappa(y)=\int_{y-T}^\infty\frac12K(t)\,dt$, exactly as for the window form in the
proof of Theorem~\ref{thm:A}; the tail has a jump at $T$ in general, and $\kappa(y)\sim\frac12\log\frac1{y-T}$ is integrable
against it. (ii) After Theorem~\ref{thm:ext-half} the prime comb does not appear: it enters only through
$\Lambda*1=\log$, that is, through the coefficients $a_n\equiv1$ of the summation map. (iii) The zero sum moves to the tail:
for every nontrivial zero $\rho$, $\Phi_{Ef}(\rho)=\zeta(\rho)\int_0^\infty f(u)u^{\rho-1}du-\Phi_{\Psi_f}(\rho)
=-\Phi_{\Psi_f}(\rho)$, so by Lemma~\ref{lem:roughEF} $W_T(Ef)=\sum_\rho\Phi_{\Psi_f}(\rho)\Phi_{\Psi_f}(1-\rho)$ (proved).
\end{remark}

\begin{proposition}[{the gauge is the radical; proved, given \cite[Lemma~4.4]{Guo2026}}]\label{prop:ext-gauge}
For $\varphi\in\mathcal G_T$, $\Phi_\varphi$ is supported in $(-\infty,-T]$ and lies in the radical of $\mathcal W$;
$\mathcal W[\Psi_f+\Phi_\varphi]=\mathcal W[\Psi_f]$, $r_{f+\varphi}=r_f+\mathcal F\varphi|_{(\lambda,\infty)}$, and for $g$ in the core
\[
\ell(g)=2\,\mathrm{dist}^2_{L^2(\lambda,\infty)}(r_f,\overline{\mathcal G}_T)\qquad(Ef=g).
\]
\end{proposition}
\begin{proof}
$ne^x>1/\lambda$ for $x>-T$, so $\Phi_\varphi=0$ there; $\varphi$ is admissible, so Guo's radical lemma applies, and both extra
terms in $\mathcal W[\Psi_f+\Phi_\varphi]$ are pairings against the radical. The leakage of an even function is twice the
integral over $(\lambda,\infty)$; Lemma~\ref{lem:ext-mob} identifies the representatives.
\end{proof}

\begin{lemma}[the gauge-minimised leakage is polynomially bounded; proved]\label{lem:ext-ellbound}
There is an absolute constant $C_1$ such that $\ell(g)\le C_1N^2\|g\|^2$ for every $g$ in the core and $N\ge5$.
\end{lemma}
\begin{proof}
Let $f_1=E^{-1}g$ on $[1/\lambda,\lambda]$, extended evenly; $\|f_1\|^2\le8N\|g\|^2$ by Lemma~\ref{lem:ext-mob}. Patch $f_1$ in
$C^2$ to zero across $\pm1/\lambda$ at width $\varepsilon\to0$, and add $b(x)=-2\lambda I_1\beta(2\lambda x)$, with $\beta\in
C_c^2(-1,1)$ even, $\beta(0)=0$, $\int\beta=1$, and $I_1$ the integral of the patched $f_1$; then $f=f_1+\text{patch}+b$ is
admissible, $|I_1|\le(2\lambda)^{1/2}\|f_1\|+o(1)$, $\|b\|^2\le4N\|\beta\|^2\|f_1\|^2+o(1)$, and $\ell(g)\le\mathrm{leak}(f)\le\|f\|^2
\le16N(1+4N\|\beta\|^2)\|g\|^2$.
\end{proof}

\subsection{The leakage lemma}\label{subsec:ext-lemma}

\begin{problem}[leakage coercivity; open]\label{prob:ext-L}
Find $\kappa>0$ and $A_0\ge0$ such that for every $T$ with $N\ge5$ and every $g$ in the core
\begin{equation}\label{eq:ext-L}
W_T(g)\ \ge\ \kappa N^{-A_0}\,\ell(g).
\end{equation}
\end{problem}

The inequality must be stated with the gauge-minimised $\ell$: with $\mathrm{leak}(f)$ in its place it is false, since any
nonzero $\varphi\in\mathcal G_T$ has $E\varphi=0$ and $\mathrm{leak}(\varphi)>0$ (proved).

\begin{proposition}[consolidated form; proved]\label{prop:ext-consol}
By Theorem~\ref{thm:ext-half} and Proposition~\ref{prop:ext-gauge}, \eqref{eq:ext-L} for real $g$ is equivalent to: for all
$N\ge5$ and all $r=\mathcal Ff|_{(\lambda,\infty)}$ with $f$ admissible at $T$,
\begin{equation}\label{eq:ext-starstar}
\begin{split}
\frac1{2\pi}\int_\R m\,|\widehat{\check\Psi}|^2+2A(\check\Psi)B(\check\Psi)&-2\big\langle\check\Psi,(Z_{\log}r)(e^\cdot)\big\rangle_{L^2(T,\infty)}\\
&\ge\ 2\kappa N^{-A_0}\,\mathrm{dist}^2(r,\overline{\mathcal G}_T),
\end{split}
\end{equation}
where $\check\Psi=(Zr)(e^\cdot)\mathbf 1_{(T,\infty)}$.
Complex $g$ follow from the real case with $\kappa/2$. Equivalent spellings, all proved: (a) with $\tilde a(\tau)=\Rel\psi(\frac54+
\frac{i\tau}2)-\log\pi$ and $C(u)=\int\check\Psi(x)\check\Psi(x+u)dx$, the left side equals
$\frac1{2\pi}\int\tilde a|\widehat{\check\Psi}|^2-2\int_1^\infty t^{-1/2}C(\log t)\,d(\psi_\Lambda(t)-t)$, with $\psi_\Lambda(t)=\sum_{n\le t}\Lambda(n)$, because $2AB-2\int_1^\infty t^{-1/2}C(\log t)dt=\int e^{-|u|/2}C(u)du=\frac1{2\pi}\int\frac{|\widehat{\check\Psi}|^2}{\frac14+\tau^2}$;
(b) in the prolate basis $\psi_k$ of $[-\lambda,\lambda]$ with concentrations $\chi_k=\chi_k(2\pi N)$, writing $f=\sum a_k\psi_k$ on the
constraint space, $D=\mathrm{diag}(1-\chi_k)$ and $\mathfrak M_{jk}=W_T(E\psi_j,E\psi_k)$, it reads
$a^{\mathsf T}\mathfrak Ma\ge\kappa N^{-A_0}\inf_{b}(a+b)^{\mathsf T}D(a+b)$, $b$ ranging over the gauge coefficients.
\end{proposition}

\begin{proposition}[the universal form; proved]\label{prop:ext-univ}
Put $\varrho(X)=\lambda r(\lambda X)$ on $(1,\infty)$ and $\check\Psi_\varrho(y)=e^{y/2}\sum_{n\ge1}\varrho(ne^y)$ for $y>0$, $0$ for $y\le0$.
Then $\mathcal W[\check\Psi_f]=\lambda^{-1}\mathfrak Q[\varrho]$ with $\mathfrak Q[\varrho]:=\mathcal W[\check\Psi_\varrho]$, and
$\mathrm{dist}^2(r,\overline{\mathcal G}_T)=\lambda^{-1}\mathrm{dist}^2(\varrho,\overline{\mathcal G}_1)$, where $\overline{\mathcal G}_1$ is the
closure of $\{\mathcal F\varphi|_{(1,\infty)}:\varphi$ even $C^2$, $\supp\varphi\subset[-1,1]$, $\varphi(0)=\mathcal F\varphi(0)=0\}$. Neither
$\mathfrak Q$ nor $\overline{\mathcal G}_1$ depends on $T$; all dependence on $T$ sits in the data class
$\mathrm{PW}_N=\{\mathcal Fh|_{(1,\infty)}:h$ even $C^2$, $\supp h\subset[-N,N]$, $h(0)=\mathcal Fh(0)=h(N)=0\}$.
\end{proposition}
\begin{proof}
$\check\Psi_f(T+y)=\lambda^{-1/2}\check\Psi_\varrho(y)$, translation invariance of $\mathcal W$, $\|\varrho\|^2=\lambda\|r\|^2$, and
$h(t)=f(t/\lambda)$, $\mathcal Fh(X)=\lambda\mathcal Ff(\lambda X)$.
\end{proof}

\begin{theorem}[the leakage lemma gives every window; proved, conditional on \eqref{eq:ext-L}]\label{thm:ext-margin}
If \eqref{eq:ext-L} holds, then for every $T$ with $N\ge5$ and every $g$ in the core
\[
W_T(g)\ \ge\ \frac{\kappa}{4}\,N^{-A_0-1}\big(1-\chi_0(2\pi N)\big)\,\|g\|^2,
\]
where $\chi_0(c)<1$ is the largest concentration of the time--band limiting operator with $c=2\pi\cdot\lambda\cdot\lambda$ and
$1-\chi_0(c)\sim4\sqrt\pi\,c^{1/2}e^{-2c}$ \cite{SlepianPollak1961,Fuchs1964}. Since $C_c^\infty(\IT)$ lies in the core and is
dense in $\mathcal D_T$ for the norm $(\int\log(2+|\tau|)|\Fh|^2)^{1/2}$, in which $W_T$ is continuous, the same bound holds for
$\lamW$; in particular $\lamW>0$ for every $T>0$ (for $N<5$ by \cite{Zhu2026} and Theorem~\ref{thm:log5W}).
\end{theorem}
\begin{proof}
Let $f$ be any admissible representative of $g$. Since $\supp f\subset[-\lambda,\lambda]$, $\mathrm{leak}(f)\ge(1-\chi_0)\|f\|^2$; since
$f$ is even and $\|g\|\le2\sqrt N\|f|_{[1/\lambda,\lambda]}\|$ (Lemma~\ref{lem:ext-mob}), $\|f\|^2\ge\|g\|^2/(2N)$. Taking the
infimum over representatives, $\ell(g)\ge(1-\chi_0)\|g\|^2/(4N)$; insert in \eqref{eq:ext-L}. Density: dilate $F\in\mathcal D_T$
towards the origin and mollify; $|m(\tau)|\le\log(2+|\tau|)+O(1)$, and the comb and $A,B$ are continuous on $L^2(\IT)$.
\end{proof}

\begin{proposition}[base case; proved, given \cite{Mori15}]\label{prop:ext-base}
\eqref{eq:ext-L} holds for $5\le N\le15$ with $A_0=2$ and $\kappa=2^{-2904}/C_1$.
\end{proposition}
\begin{proof}
\cite{Mori15} gives $W_T(g)\ge2^{-2904}\|g\|^2$ for $C^2$ test functions at every $T\le\frac12\log15$; combine with
Lemma~\ref{lem:ext-ellbound}.
\end{proof}

\begin{lemma}[high-frequency coercivity; proved, comb-generic]\label{lem:ext-hf}
Let $C_\infty$ satisfy $|m(\tau)-\log(2+|\tau|)|\le C_\infty$, put $K_T=C_\infty+\Aconst$ and $\Omega_*=e^{K_T+\lambda+1}-2$, and
let $P_<$ be the Fourier projection of $L^2(\IT)$ onto $|\tau|<\Omega_*$. Then $W_T(g)\ge\|g\|^2-(K_T+\lambda+1)\|P_<g\|^2$.
Hence \eqref{eq:ext-L} holds with $A_0=2$, $\kappa=1/(2C_1)$ on $\{g:\|P_<g\|^2\le\|g\|^2/(2(K_T+\lambda+1))\}$.
\end{lemma}
\begin{proof}
$\sT\ge\log(2+|\tau|)-K_T$, which is $\ge\lambda+1$ for $|\tau|\ge\Omega_*$; the rank-two form $2\Rel(A\bar B)$ has eigenvalues
$2T\pm2\sinh T\ge-\lambda$. So $W_T(g)\ge(\lambda+1)\|P_>g\|^2-K_T\|P_<g\|^2-\lambda\|g\|^2$. Then Lemma~\ref{lem:ext-ellbound}.
\end{proof}

The complement of that subspace has dimension of order $T\,\Omega_*$, superpolynomial in $N$, and contains the whole
prolate range $k\lesssim4N$; Lemma~\ref{lem:ext-hf} says nothing there.

\subsection{Structure of the exterior form}\label{subsec:ext-struct}

In the universal frame put, for $s=\frac12+i\tau$ and $\zeta_X(w)=\sum_{n\le X}n^{-w}$,
\[
\mathsf F(s)=\int_1^\infty\varrho(X)X^{-s}\zeta_X(1-s)\,dX,\qquad
\mathsf F_{\log}(s)=\int_1^\infty\varrho(X)X^{-s}\sum_{n\le X}n^{s-1}\log n\,dX,
\]
so that $\mathsf F(s)=\int_0^\infty\check\Psi_\varrho(y)e^{-(s-\frac12)y}dy$, $A(\check\Psi_\varrho)=\mathsf F(0)=\int\varrho H$ with
$H(X)=\sum_{j\le X}j^{-1}$, and $B(\check\Psi_\varrho)=\mathsf F(1)=\int\varrho(X)\lfloor X\rfloor X^{-1}dX$. For pairs $U=(U,U_{\log})$ write
$Q(U,V)=\frac1{2\pi}\int[m\Rel(\bar UV)-\Rel(\bar UV_{\log}+\bar VU_{\log})]d\tau$. Data are $\varrho\in C^1(1,\infty)$ with
$\varrho=O(X^{-2})$; $\varrho_0=\int_1^\infty\varrho$, $\varrho_1=\int_1^\infty\varrho\log X$, $R(s)=\int_1^\infty\varrho X^{-s}dX$.

\begin{theorem}[critical-line form; proved]\label{thm:ext-line}
$\mathfrak Q[\varrho]=Q(\mathsf F,\mathsf F)+2\mathsf F(0)\mathsf F(1)$, and on the line $m(\tau)=-2\Rel(\zeta'/\zeta)(\frac12+i\tau)$ wherever
$\zeta(\frac12+i\tau)\ne0$, the right side extending continuously across such points.
\end{theorem}
\begin{proof}
$\int_0^\infty\check\Psi_\varrho e^{-(s-\frac12)y}dy=\sum_nn^{s-1}\int_n^\infty\varrho X^{-s}dX$, absolutely convergent; then
Theorem~\ref{thm:ext-half} in Plancherel form. The symbol identity follows from $\xi(s)=\xi(1-s)$ on the line, where
$\Rel(\frac1s+\frac1{s-1})=0$.
\end{proof}

\begin{theorem}[Euler--Maclaurin--Poisson partition; proved, an identity]\label{thm:ext-EM}
With $\zeta_X(1-s)=\zeta(1-s)+X^s/s+e_s(X)$ (Euler--Maclaurin, \cite[\S3.5]{Titchmarsh1986}), split
$\mathsf F=F_0+G$, $F_0=\zeta(1-s)R(s)$, $G=\varrho_0/s+\int\varrho X^{-s}e_s(X)dX$, and likewise $\mathsf F_{\log}$ with
$F_0^{\log}=-\zeta'(1-s)R(s)$. Then $m|F_0|^2-2\Rel(\bar F_0F_0^{\log})=0$ pointwise on the line, so
\[
\mathfrak Q=2Q(F_0,G)+Q(G,G)+2\mathsf F(0)\mathsf F(1)=-Q(G,G)+2Q(\mathsf F,G)+2\mathsf F(0)\mathsf F(1),
\]
the second form being the reflection identity $Q(F_0,G)=Q(\mathsf F,G)-Q(G,G)$. $F_0$ is null for $Q$ but not for $\mathfrak Q$:
the pole term contains $F_0(1)=\zeta(0)R(1)$.
\end{theorem}

\begin{proposition}[the pole as a residue; proved]\label{prop:ext-residue}
For $c'\in(-\frac12,0)$,
\[
2Q(F_0,G)+2F_0(1)\mathsf F(0)=2J_{c'},\quad
J_{c'}=\frac1{2\pi i}\int_{(c')}R(1-s)\big[-\chi(s)\zeta'(1-s)\mathsf F(s)-\zeta(s)\mathsf F_{\log}(s)\big]ds,
\]
$\chi(s)=\zeta(s)/\zeta(1-s)$. The shift crosses only $s=0$: $\zeta(s)(\zeta'/\zeta)(1-s)=\chi(s)\zeta'(1-s)$ has no poles at
$1-\rho$, and its only pole in $c'\le\Rel s\le\frac12$ is at $s=0$ ($\chi(s)\sim s/2$, $\zeta'(1-s)\sim-s^{-2}$), with residue
$\frac12R(1)\mathsf F(0)$. Hence $\mathfrak Q=Q(G,G)+2G(1)\mathsf F(0)+2J_{c'}$; the pole couples the $Q$-null vector through one
rank-one term only.
\end{proposition}

\begin{proposition}[a pointwise identity; proved]\label{prop:ext-pointwise}
For $\zeta(\frac12+i\tau)\ne0$ put $\Phi(s,w)=\int_1^\infty\varrho X^{-s}\zeta_X(w)dX$ and $h_\tau(w)=\Phi(s,w)/\zeta(w)$ near
$w_0=1-s$. Then $m|\mathsf F|^2-2\Rel(\bar{\mathsf F}\mathsf F_{\log})=|\zeta(\frac12+i\tau)|^2\,\partial_u|h_\tau|^2(w_0)$, $w=u+iv$; the
right side equals $\partial_u|\Phi|^2+m|\Phi|^2$ at $w_0$ and is continuous across zeros on the line. $Q$-null vectors are the data
whose quotient $h_\tau$ does not depend on $w$; $F_0$ is one.
\end{proposition}

\begin{proposition}[dual structure and the polar mode; proved]\label{prop:ext-dual}
Let
\[\varrho^\vee(\xi)=2\int_1^\infty\varrho(X)\cos(2\pi\xi X)\,dX,\]
the transform of the even extension, and let $\mathsf F^\vee$ be
$\mathsf F$ computed for the data $\varrho^\vee|_{(1,\infty)}$.
(a) On the line, $G(s)=\varrho_0/s-\overline{\mathsf F^\vee(s)}$, and the $k$-th
Poisson harmonic of $G$ is the transform of $u_k(v)=k^{-1/2}u_1(v+\log k)\mathbf 1_{v>0}$: the harmonics are dilates of one
function with coefficients $\equiv1$. (b) The polar part $G_{\mathrm{pol}}=\varrho_0/s$ is the transform of
$\varrho_0e^{v/2}\mathbf 1_{v<0}$, and $Q(G_{\mathrm{pol}},G_{\mathrm{pol}})=c_0\varrho_0^2-2\varrho_0\varrho_1$ with
$c_0=\psi(\frac32)-\log\pi=-1.1082399\ldots$, using
\[\frac1{2\pi}\int m(\tau)(\tfrac14+\tau^2)^{-1}d\tau=c_0-2;\]
this $2\times2$ form
in $(\varrho_0,\varrho_1)$ has determinant $-1$.
\end{proposition}

\begin{theorem}[Poisson mirror; proved for $\varrho\in C_c^\infty(1,\infty)$, conditional for Paley--Wiener tails]\label{thm:ext-mirror}
Put $P(y)=e^{y/2}\sum_{k\ge1}\varrho^\vee(ke^y)$ on $\R$, $G_h=G-\varrho_0/s$, and
$\mathfrak Q_{\mathrm{line}}[\varrho^\vee]=Q(\mathsf F^\vee,\mathsf F^\vee)$. Then:
(i) $P(y)=e^{-y/2}\sum_{m\ge1}\varrho(me^{-y})-\varrho_0e^{y/2}$ for all $y$; on $y<0$ the dual function is the reflected tail of
$\varrho$ minus the polar mode;
(ii) with $M(z)=(2\pi)^{-z}\Gamma(z)\cos\frac{\pi z}2$, $\zeta(1-z)=2M(z)\zeta(z)$ gives $M'/M(z)=-(\zeta'/\zeta)(1-z)-(\zeta'/\zeta)(z)$,
which on $\Rel z=\frac12$ is real and equals $m$;
(iii)
\[
Q(G,G)=-\mathfrak Q_{\mathrm{line}}[\varrho^\vee]+2\mathfrak B[\varrho]-R(1)G_h(0)+c_0\varrho_0^2-2\varrho_0\varrho_1,\qquad
\mathfrak B[\varrho]=\frac1{2\pi}\int m\,\Rel(\mathsf F\,\overline{G_h})\,d\tau,
\]
where $G_h(0)$ is the Laplace transform of the time function of $G_h$ at $s=0$ (convergent). $\mathfrak B$ couples functions with
time supports on opposite sides of $0$, so only the off-diagonal L\'evy kernel enters it.
\end{theorem}
\begin{proof}
(i) Poisson summation for the even extension, which vanishes at $0$. (ii) Logarithmic differentiation of the functional
equation; on the line $1-z=\bar z$. (iii) The log weight of $G$ splits as $\log(\xi X)-\log k$ under $\xi=ke^v$, which writes
$Q(G_h,G_h)+\mathfrak Q_{\mathrm{line}}[\varrho^\vee]$ as twice $\frac1{2\pi}\int m|\mathsf F^\vee|^2$ minus the pairing of the dual
function with its log-weighted version; by (ii) and M\"untz's formula \cite[\S2.11]{Titchmarsh1986}, the latter is the
multiplier $m$ applied to $P$ up to $-\frac12R(1)e^{-y/2}$; by (i) the half of $P$ on $y<0$ is the reflected tail minus the
polar mode. The pairing of the polar mode with $G_h$ cancels against $2Q(G_{\mathrm{pol}},G_h)$, and
Proposition~\ref{prop:ext-dual}(b) adds $Q(G_{\mathrm{pol}},G_{\mathrm{pol}})$. For Paley--Wiener tails ($\varrho(1)\ne0$,
$\varrho^\vee=O(1/\xi)$) the M\"untz step is not verified.
\end{proof}

By Theorem~\ref{thm:ext-EM}, the defect part of $\mathfrak Q$ is $-Q(G,G)$; by Theorem~\ref{thm:ext-mirror} its positivity is the
exterior positivity of the Poisson-dual data plus interface, residue and polar terms, a statement of the same type as
\eqref{eq:ext-starstar}. The identity is exact only for the weight $1$ on the comb: replacing $\Lambda$ by $c\Lambda$ adds
$2(1-c)\frac1{2\pi}\int m|\mathsf F^\vee|^2$.

\begin{proposition}[off-line contour form; proved]\label{prop:ext-contour}
For $f$ admissible, $\Phi_{\Psi_f}$ is holomorphic on $\Rel s>-1$, $O_\sigma(1/|\Iml s|)$ there, and for every $\sigma\in(1,2)$
\[
W_T(Ef)=\frac1{2\pi}\int m(\tau)|\Phi_{\Psi_f}(\tfrac12+i\tau)|^2d\tau-\frac1{\pi i}\int_{(\sigma)}\Big(-\frac{\zeta'}\zeta\Big)(s)
\Phi_{\Psi_f}(s)\Phi_{\Psi_f}(1-s)\,ds+2\Phi_{\Psi_f}(0)\Phi_{\Psi_f}(1).
\]
\end{proposition}
\begin{proof}
$c_{\Psi}(a)=\frac1{2\pi i}\int_{(\sigma)}\Phi_\Psi(s)\Phi_\Psi(1-s)e^{-(s-\frac12)a}ds$ for $-1<\sigma<2$; sum against $v_q/2$ for
$\sigma>1$, where $\sum\Lambda(q)q^{-\sigma}<\infty$ and $\sum_q\Lambda(q)q^{-s}=-\zeta'/\zeta(s)$.
\end{proof}

\begin{proposition}[pair expansion and comb collapse; proved]\label{prop:ext-pairs}
With $\tilde r(x)=e^{x/2}r(e^x)$ and $r_n=\tilde r\,\mathbf 1_{(T+\log n,\infty)}$, $\check\Psi=\sum_nn^{-1/2}r_n(\cdot+\log n)$, so
$\mathcal W[\check\Psi]=\sum_{m,n}(mn)^{-1/2}\mathcal W(r_m(\cdot+\log m),r_n(\cdot+\log n))$, the archimedean part read as a
polarised L\'evy form. Grouping $(m,n)=(j\alpha,j\beta)$ with $(\alpha,\beta)=1$ introduces the harmonic cut
$H(\min(e^{y-T}/\alpha,e^{y'-T}/\beta))$, and the prime part of class $(\alpha,\beta)$ is
$-\int\tilde r(y)\tilde r(y+\log\frac\beta\alpha)G_{\alpha\beta}(e^{y-T}/\alpha)dy$ with $G_{\alpha\beta}(X)=\sum_{j\le X}(2\log j+\log\alpha\beta)/j$.
The diagonal class $(1,1)$ is indefinite on $C_c^\infty(T,\infty)$: for a unit bump at depth $D$, its archimedean and polar part
grows like $D$ and its prime part like $-D^2$.
\end{proposition}

\begin{proposition}[the deep quasimode; proved, given the prime number theorem]\label{prop:ext-deep}
Let $\phi\in C_c^\infty(0,1)$, $B_\phi=\int\phi(s)e^{s/2}ds$, and $\varrho_D(X)=X^{-1/2}\phi(\log X-D)$, so $\int\varrho_D=e^{D/2}B_\phi$. Then,
as $D\to\infty$,
\[
\mathfrak Q[\varrho_D]=\Big(2+\gamma-\log4\pi\Big)B_\phi^2\,e^{D}+o(e^D),\qquad
2+\gamma-\log4\pi=2\frac{\xi'}{\xi}(1)=0.0461914\ldots
\]
The constant is the sum of $\psi(\frac14)-\log\pi=-5.3721834\ldots$ (constant of $m$), $\log2+\frac\pi2=2.2639435\ldots$ (L\'evy energy
of the profile $e^{-x/2}$ with its jump), $2$ (pole against the main term of the comb, residue $1$) and $2\gamma=1.1544313\ldots$
(Mertens' $\sum_{n\le Y}\Lambda(n)/n=\log Y-\gamma+o(1)$ \cite{MontgomeryVaughan2007}); the first three give
$c_0=\psi(\frac32)-\log\pi<0$.
\end{proposition}
\begin{proof}
Poisson in the dilation index gives $\check\Psi=\varrho_0e^{-y/2}\mathbf 1_{(0,D)}$ plus a term superpolynomially small on $(0,D-C)$,
bounded on $(D-C,D+1)$ and zero beyond. On the main term, the archimedean form with the pole is
$c_0\varrho_0^2+o(\varrho_0^2)$ (via $\tilde a$ of Proposition~\ref{prop:ext-consol}(a)), and the Stieltjes term contributes
$2\gamma\varrho_0^2+o(\varrho_0^2)$ by Mertens' formula and dominated convergence; the cross terms are $O(e^{D/2})$ by
Cauchy--Schwarz, and $\psi_\Lambda(Y)-Y=o(Y)$ controls the remaining pairing.
\end{proof}

\begin{remark}\label{rem:ext-deep}
The same constant is $\sum_\rho\frac1{\rho(1-\rho)}=2\sum_\rho\frac1\rho=2+\gamma-\log4\pi$ \cite[Ch.~12]{Davenport2000}: by
Lemma~\ref{lem:roughEF} applied to the profile $\varrho_0e^{-y/2}\mathbf 1_{(0,D)}$, the zero sum gives
$\varrho_0^2\sum_\rho\frac1{\rho(1-\rho)}+o(\varrho_0^2)$ with no input from the prime number theorem. The deep quasimode
is not in $\mathrm{PW}_N$ modulo $\overline{\mathcal G}_1$ (a bump at depth $D$ comes from a function of support $\asymp e^{-D}$, a gauge
direction), so Proposition~\ref{prop:ext-deep} is a test value for arguments valid on all of $C_c^\infty$, not information on
\eqref{eq:ext-starstar}.
\end{remark}

\begin{proposition}[minimal-leakage admissible prolate combination; proved, given \cite{SlepianPollak1961,Fuchs1964}]\label{prop:ext-K1b}
Among even $f\in\mathrm{span}\{\psi_0,\psi_2,\psi_4,\dots\}$ with $\mathcal Ff(0)=0$, the leakage quotient $\mathrm{leak}(f)/\|f\|^2$ is
minimised at $f=\psi_0-\sqrt2\,\psi_2\,(1+O(c^{-1}))$, with minimum $\frac23(1-\chi_2)(1+O(c^{-2}))=\frac{64}3c^2(1-\chi_0)(1+o(1))$,
$c=2\pi N$. The condition $f(0)=0$ costs nothing at the level of the form (it is a gauge move), but leakage is not invariant under
gauge moves; this is why \eqref{eq:ext-L} uses $\ell$.
\end{proposition}

\begin{proposition}[{the window flow; proved, given \cite[Lemma~4.4]{Guo2026}}]\label{prop:ext-flow}
In the dilated frame $\phi(x)=\lambda^{1/2}f(\lambda x)$, $G_\phi(y)=e^{y/2}\sum_n\phi(ne^y)$ ($G_\phi=0$ on $[0,\infty)$), one has
$Ef=G_\phi(\cdot-T)$ on $\IT$ and, for fixed $\phi$ and $b=-2T$,
\[
\frac d{dT}W_T(Ef)=-2\,G_\phi(b)\,(W^-G_\phi)(b),\qquad G_\phi(b)=N^{1/2}\sum_{k\ge1}\mathcal F\phi(kN),
\]
\[
(W^-G)(b)=\int_{x>b}G(x)k(x-b)dx-\int_{y<b}G(y)k(b-y)dy+\sum_{q\ge2}\tfrac{v_q}2\big[G(b+\log q)-G(b-\log q)\big],
\]
$k(t)=\frac12K(t)-2\cosh\frac t2$: the window form changes only through the band-edge alias $G_\phi(-2T)$ paired with an odd Weil
distribution. Integrating in $T$ from the base case of Proposition~\ref{prop:ext-base} yields \eqref{eq:ext-L} under an additional
inequality at the minimiser of the pencil $W_T-\kappa N^{-A_0}\ell$ (a virial bound of relative precision $O(1/N)$); this
implication is conditional on existence, simplicity and continuity of that minimiser, and the additional inequality is open.
\end{proposition}

\subsection{Positive regions of the exterior form and the M\"obius lattice}\label{subsec:ext-pos}

This subsection evaluates the universal form $\mathfrak Q$ of Proposition~\ref{prop:ext-univ} on three explicit families of
data (chirps in one integer cell, multi-cell chirps of prolate phase, single prolate tails) and records the exact
M\"obius-lattice structure of the exterior datum of an admissible function. None of these results gives a sign of
$\mathfrak Q$ on $\mathrm{PW}_N$ modulo $\overline{\mathcal G}_1$; they are tools, and each carries its own label.

\paragraph{Conventions.} For complex data $\varrho$ on $(1,\infty)$ put $\mathfrak Q[\varrho]:=\mathfrak Q[\Rel\varrho]+\mathfrak Q[\Iml\varrho]$
(the Hermitian extension; the pole term becomes $2\Rel(A\bar B)$). For $\varrho$ on $(1,\infty)$ write
$u(x)=e^{x/2}\varrho(e^x)$, so $\|u\|_{L^2(dx)}=\|\varrho\|_{L^2(dt)}$, and $u_n(y)=u(y+\log n)\mathbf 1_{y>0}$, so that
$\check\Psi_\varrho=\sum_{n\ge1}n^{-1/2}u_n$. Since $m$ is even, $\frac1{2\pi}\int m|\widehat u|^2$ does not depend on the sign
convention of the Fourier transform; we write it $\langle m(D)u,u\rangle$. The archimedean part of $\mathfrak Q$ is the
full-line form of the zero-extended $\check\Psi_\varrho$: by Lemma~\ref{lem:K}, $\langle m(D)\Psi,\Psi\rangle
=C_0\|\Psi\|^2+\frac14\iint|\Psi(x)-\Psi(x')|^2K(|x-x'|)\,dx\,dx'$. Put
\[
H(t)=\sum_{n\le t}\frac1n,\qquad H_1(t)=\sum_{n\le t}\frac{\log n}n,\qquad H_J=H(J),
\]
and for $\nu>0$, $t\ge1$ and $N\ge1$
\[
w(t,\nu)=H(t)\log\frac{\nu t}{2\pi}-2H_1(t),\qquad
\mathfrak w_N(t)=H(t)\log\frac{Nt^2}{\sqrt{t^2-1}}-2H_1(t)\quad(t>1).
\]
$w$ is the \emph{phase-space weight}; $\mathfrak w_N(t)=w(t,\phi_c'(t))$ for the \emph{prolate phase}
$\phi_c(t)=c\sqrt{t^2-1}$, $c=2\pi N$, whose Mellin frequency is $\tau(t)=t\phi_c'(t)=ct^2/\sqrt{t^2-1}$. For a comb weight
$c'>0$, $\mathfrak Q^{c'}$ denotes $\mathfrak Q$ with every $v_n$ replaced by $c'v_n$ (archimedean part and pole unchanged).

\begin{lemma}[the weight lemma; proved]\label{lem:ext-weight}
Put $F(J)=H_J\log J-2H_1(J)$ for integers $J\ge1$. Then $0\le F(J)\le\gamma\log J+1$. Consequently:
\emph{(i)} $w(t,\nu)\ge H(t)\log\frac\nu{2\pi}$ for all $t\ge1$, $\nu>0$, and on $[J,J+1)$,
$w(t,\nu)=H_J\log\frac\nu{2\pi}+F(J)+H_J\log\frac tJ$;
\emph{(ii)} for $N\ge4$ and all $t>1$, $\mathfrak w_N(t)\ge\log2N$, with equality if and only if $t=\sqrt2$.
\end{lemma}
\begin{proof}
$F(1)=0$ and $F(J+1)-F(J)=H_J\log(1+\frac1J)-\frac{\log(J+1)}{J+1}\ge0$, since $H_J\ge\log(J+1)$ and
$\log(1+\frac1J)\ge\frac1{J+1}$. For the upper bound use $H_J\le\log J+\gamma+\frac1{2J}$ and $H_1(J)\ge\frac12\log^2J+\gamma_1$, where
$\gamma_1=-0.0728\ldots$ is the first Stieltjes constant: $H_1(J)-\frac12\log^2J$ is nonincreasing for $J\ge3$ (its increment is
$f(J+1)-\int_J^{J+1}f\le0$ with $f(x)=\frac{\log x}x$ decreasing on $[e,\infty)$) and tends to $\gamma_1$, and $J=1,2$ are checked
directly; so $F(J)\le\gamma\log J+\frac{\log J}{2J}-2\gamma_1\le\gamma\log J+1$. (i) On $[J,J+1)$, $H(t)=H_J$ and $H_1(t)=H_1(J)$,
which gives the identity, and $F(J)\ge0$, $\log\frac tJ\ge0$. (ii) On $(1,2)$, $H=1$, $H_1=0$ and $\frac{t^4}{t^2-1}\ge4$ is
$(t^2-2)^2\ge0$. On $t\ge2$, $\frac{t^2}{\sqrt{t^2-1}}\ge t$, so by (i) $\mathfrak w_N(t)\ge H(t)\log N\ge\frac32\log N\ge\log2N$ for
$N\ge4$.
\end{proof}

In the variable $t=\sec\theta$ the prolate density $dt/(t\sqrt{t^2-1})$ is $d\theta$, and on $(1,2)$
$\mathfrak w_N(\sec\theta)=\log2N-\log\sin2\theta$; this is the form in which the weight enters Theorem~\ref{thm:ext-layer}.

\subsubsection*{Chirps in one integer cell}

\begin{lemma}[disjoint dilates and the exact comb; proved]\label{lem:ext-cell}
Let $J\ge2$, $\delta\in(0,\frac14)$ and $\supp\varrho\subset[J+\delta,J+1-\delta]$, $u\in H^1$. Then:
\emph{(i)} $u_n=u(\cdot+\log n)$ is supported in $(0,\infty)$ for $n\le J$ and $u_n=0$ for $n\ge J+1$, so
$\check\Psi_\varrho=\sum_{n\le J}n^{-1/2}u_n$ with no truncation inside any support; for $1\le m<n\le J$ the supports are disjoint,
with gap $g_{m,m+1}\ge2\delta/J$ and $g_{m,m+k}\ge k/(2J)$ for $k\ge2$.
\emph{(ii)} If the prime comb is replaced by any comb $\sum_qw(q)q^{-1/2}(\delta_{\log q}+\delta_{-\log q})$, its contribution to
$\mathfrak Q[\varrho]$ is exactly $-2\|\varrho\|^2\sum_{n\le J}(w*1)(n)/n$; for $w=\Lambda$ it is $-2H_1(J)\|\varrho\|^2$, independent of
the phase of $\varrho$.
\emph{(iii)} With $A=J\int\varrho(t)t^{-1}dt$, $B=H_J\int\varrho(t)dt$,
\[
\mathfrak Q[\varrho]=H_J\langle m(D)u,u\rangle-2H_1(J)\|\varrho\|^2+\mathfrak X[\varrho]+2\Rel(A\bar B),\quad
\mathfrak X[\varrho]=-\tfrac12\sum_{\substack{m\ne n\le J}}(mn)^{-1/2}\iint u_m(y)\overline{u_n(y')}K(|y-y'|)\,dy\,dy'.
\]
\end{lemma}
\begin{proof}
(i) $(J+\delta)/n>1$ for $n\le J$ and $(J+1-\delta)/n<1$ for $n\ge J+1$. The gap between the supports of $u_m$ and $u_n$ is
$\log\frac{J+\delta}m-\log\frac{J+1-\delta}n$. For $n=m+k$, $\log\frac nm\ge\frac kJ$ and $\log\frac{J+1-\delta}{J+\delta}\le\frac{1-2\delta}J$,
so the gap is $\ge\frac{k-1+2\delta}J\ge\frac k{2J}$ for $k\ge2$. For $k=1$ the worst case is $m=J-1$, where
$J(J+\delta)-(J-1)(J+1-\delta)=2\delta J+1-\delta$; with $\log(1+x)\ge x/(1+x)$ the gap is
$\ge\frac{2\delta J+1-\delta}{J(J+\delta)}\ge\frac{2\delta}J$, the last step being $1-\delta\ge2\delta^2$.
(ii) $\check\Psi(y+\log q)=\sum_ll^{-1/2}u(y+\log lq)$, and by (i) $\int u_m\bar u_n=\delta_{mn}\|u\|^2$ for $m,n\le J$, so
$\langle\check\Psi,\check\Psi(\cdot+\log q)\rangle=q^{-1/2}H(J/q)\|u\|^2$; summing,
$-2\sum_qw(q)q^{-1}H(J/q)\|u\|^2=-2\|u\|^2\sum_{n\le J}n^{-1}\sum_{q\mid n}w(q)$, and $\Lambda*1=\log$.
(iii) Diagonal terms: $m(D)$ commutes with translations, $\sum_{n\le J}\frac1n=H_J$. Off-diagonal terms: the supports are
disjoint, so the $C_0$ part and the terms $u_m(x)\bar u_n(x)$ of the polarised L\'evy form vanish, leaving the stated kernel,
absolutely convergent since $K$ is bounded on $[g_{mn},\infty)$. Pole: $\sum_{n\le J}n^{-1/2}\int u_ne^{-y/2}dy=J\int\varrho\,dt/t$,
$\sum_{n\le J}n^{-1}\int ue^{x/2}dx=H_J\int\varrho\,dt$.
\end{proof}

\begin{lemma}[oscillatory log-symbol; proved, given Stirling's remainder]\label{lem:ext-osc}
Let $u=\mathcal Ae^{i\Phi}$ with $\mathcal A\in C_c^2(\R)$, $\Phi\in C^3$ real, $\Phi'\ge\Omega\ge4\pi$ on $\supp\mathcal A$, and
\[
K_0=\|\mathcal A'\|_2+\|\mathcal A\Phi''/\Phi'\|_2,\qquad K_1=\|\mathcal A'\|_1+2\|\mathcal A\Phi''/\Phi'\|_1,
\]
\[
K_2=\|\mathcal A''\|_1+6\|\mathcal A'\Phi''/\Phi'\|_1+2\|\mathcal A\Phi'''/\Phi'\|_1+12\|\mathcal A\Phi''^2/\Phi'^2\|_1 .
\]
With $C_\Theta=\sup_{|\tau|\ge1}\tau^2|m(\tau)-\log(|\tau|/2\pi)|<\infty$,
\[
\Big|\langle m(D)u,u\rangle-\int\log\frac{\Phi'(x)}{2\pi}|\mathcal A(x)|^2dx\Big|\le\frac{2\log2}\Omega K_0\|\mathcal A\|_2
+\frac{4C_\Theta\|\mathcal A\|_2^2}{\Omega^2}+\frac{4\log2\,K_1K_2}{\pi\Omega^2}+\frac{K_2^2(7+2\log(1+\Omega))}{2\Omega^3}.
\]
\end{lemma}
\begin{proof}[Sketch]
Write the frequency variable so that $\widehat u$ concentrates at $\tau=\Phi'$. (a) For $\tau\le\Omega/2$, put $\sigma=\Omega/2-\tau$; on
$\supp\mathcal A$, $\Phi'-\tau\ge\Phi'/2$ and $\ge\Omega/2+\sigma$. Two integrations by parts with $Lf=(f/(\Phi'-\tau))'$, using
$L^2\mathcal A=\mathcal A''/g^2-3\mathcal A'\Phi''/g^3-\mathcal A\Phi'''/g^3+3\mathcal A\Phi''^2/g^4$ ($g=\Phi'-\tau$) and
$g^{-j}\le(2/\Phi')^{j-2}(\Omega/2+\sigma)^{-2}$, give $|\widehat u(\tau)|\le K_2(\Omega/2+\sigma)^{-2}$; one integration gives
$K_1/(\Omega/2+\sigma)$. (b) With $\tilde D=\max(D,\Omega/2)$, the difference $\langle m(D)u,u\rangle-\langle\log(\tilde D/2\pi)u,u\rangle$
is $\le4C_\Theta\|\mathcal A\|^2/\Omega^2$ on $\tau>\Omega/2$ (Stirling, \cite[\S5.11]{DLMF}) plus, by (a), at most
$K_2^2(7+2\log(1+\Omega))/(2\Omega^3)$ on $\tau\le\Omega/2$. (c) With $M$ the multiplication by $\max(\Phi',\Omega)$ and
$\log X-\log Y=\int_0^\infty(X+s)^{-1}(X-Y)(Y+s)^{-1}ds$, put $v_s=(M+s)^{-1}u=\mathcal A_se^{i\Phi}$; then
$(\tilde D-M)v_s=-i\mathcal A_s'e^{i\Phi}+(\Omega/2-D)_+v_s$. The first piece is $\le K_0\|\mathcal A\|\int_0^\infty\frac{ds}{(\Omega/2+s)(\Omega+s)}
=\frac{2\log2}\Omega K_0\|\mathcal A\|$ because $(\Phi'+s)^2\ge\Phi'(\Omega+s)$; the second is supported in $\tau\le\Omega/2$ and by
(a) applied to $\mathcal A_s$ and $\mathcal A$ contributes at most $4\log2\,K_1K_2/(\pi\Omega^2)$.
\end{proof}

\begin{theorem}[single-cell chirps; proved]\label{thm:ext-cell}
Let $J\ge2$, $\delta\in(0,\frac14)$, $\varrho=ae^{i\phi}$ with $a\in C_c^2((J+\delta,J+1-\delta))$, $\phi\in C^3$ real, $\phi'\ge\nu_0$ on
$\supp a$ and $J\nu_0\ge4\pi$. Put $\Omega=(J+\delta)\nu_0$, $\mathcal A(x)=e^{x/2}a(e^x)$, $\Phi(x)=\phi(e^x)$, with $K_0,K_1,K_2$ as in
Lemma~\ref{lem:ext-osc}; in the variable $t$, $K_0=\|\frac12a+ta'\|+\|a\,t(\log t\phi')'\|$ and $K_1\le2J^{-1/2}K_0$. Then
\[
\mathfrak Q[\varrho]=\int w(t,\phi'(t))\,|a(t)|^2\,dt+E,\qquad |E|\le E_{\mathrm{arch}}+E_{\mathrm{cross}}+E_{\mathrm{pole}},
\]
\[
E_{\mathrm{arch}}=H_J\Big[\frac{2\log2}\Omega K_0\|a\|+\frac{4C_\Theta\|a\|^2}{\Omega^2}+\frac{4\log2\,K_1K_2}{\pi\Omega^2}+\frac{K_2^2(7+2\log(1+\Omega))}{2\Omega^3}\Big],
\]
\[
E_{\mathrm{cross}}=\frac{2H_J}{\nu_0}\Big[K_0\|a\|\frac{\frac1{2\delta}+2\log J}J+\|a\|^2\Big(\frac3{4\delta^2}+8\Big)\Big],\qquad
E_{\mathrm{pole}}=\frac{2\sqrt2\,H_J(K_0+\frac12\|a\|)^2}{(J\nu_0)^2}.
\]
No monotonicity of $t\phi'$ is needed.
\end{theorem}
\begin{proof}
By Lemma~\ref{lem:ext-cell}(iii) it remains to treat the three terms. \emph{Diagonal:} Lemma~\ref{lem:ext-osc} with $\Phi'=t\phi'\ge\Omega$
gives $H_J\langle m(D)u,u\rangle=H_J\int\log\frac{t\phi'}{2\pi}|a|^2dt+O(E_{\mathrm{arch}})$, and $H_J\log\frac{t\phi'}{2\pi}-2H_1(J)=w(t,\phi')$
on the cell; this is the \emph{key cancellation}: the position part $H_J\log t\approx\log^2J$ of the archimedean diagonal, coming
from the Mellin frequency $t\phi'$, is cancelled by the comb $-2H_1(J)\approx-\log^2J$ down to $F(J)=O(\log J)$ (Lemma~\ref{lem:ext-weight}),
exactly because $\Lambda*1=\log$ and each dilate is counted once with the harmonic cut. \emph{Cross classes:} for $m\ne n$
integrate by parts once in $y$ against $e^{i\Phi(y+\log m)}$; with $K(s)/2\le1/s$ and $|K'(s)|/2\le3/s^2$, $\|\mathcal A\|_1\le J^{-1/2}\|a\|$,
$\|\mathcal A'\|_1+\|\mathcal A\Phi''/\Phi'\|_1\le J^{-1/2}K_0$, and the gaps of Lemma~\ref{lem:ext-cell}(i),
$\sum_{m<n}(mn)^{-1/2}g_{mn}^{-1}\le H_JJ(\frac1{2\delta}+2\log J)$ and
$\sum_{m<n}(mn)^{-1/2}g_{mn}^{-2}\le H_JJ^2(\frac1{4\delta^2}+4(\frac{\pi^2}6-1))$, which gives $E_{\mathrm{cross}}$. \emph{Pole:} one
integration by parts in $x$ gives $|\int\varrho\,dt/t|\le(K_0+\frac12\|a\|)/(J\Omega)$ and
$|\int\varrho\,dt|\le\sqrt{(J+1)/J}\,(K_0+\frac12\|a\|)/\Omega$.
\end{proof}

\begin{corollary}[a positive region; proved]\label{cor:ext-cellpos}
Under the hypotheses of Theorem~\ref{thm:ext-cell},
$\mathfrak Q[\varrho]\ge\big[H_J\log\frac{\nu_0}{2\pi}+F(J)+H_J\log(1+\frac\delta J)\big]\|a\|^2-|E|$. For $a(t)=b(t-J)$ with $b\in
C_c^2(\delta,1-\delta)$ fixed and $\phi'\equiv\nu_0$, $|E|\le C(b,\delta)H_J\log(J+1)/\nu_0$, so $\mathfrak Q[\varrho]>0$ as soon as
$\nu_0\log\frac{\nu_0}{2\pi}\ge C'(b,\delta)\log(J+2)$: the frequency threshold is logarithmic in the cell index, whereas
Lemma~\ref{lem:ext-hf}, being comb-generic, needs a threshold exponential in $\sqrt J$.
\end{corollary}

\begin{corollary}[the comb weight is detected; proved]\label{cor:ext-cfals}
Theorem~\ref{thm:ext-cell} holds for $\mathfrak Q^{c'}$ with $w$ replaced by $w^{c'}(t,\nu)=w(t,\nu)-2(c'-1)H_1(J)$ (Lemma~\ref{lem:ext-cell}(ii)
with $w=c'\Lambda$; the error bounds are unchanged).
\emph{(i)} If $c'\le1$, then $w^{c'}\ge w$, and Corollary~\ref{cor:ext-cellpos} holds verbatim for $\mathfrak Q^{c'}$. This is a statement
about the cone of high-frequency single-cell chirps only; it says nothing about $\mathfrak Q^{c'}$ on other data.
\emph{(ii)} Let $c'>1$, $b\in C_c^2(\delta,1-\delta)$ fixed, $\varrho_J(t)=b(t-J)e^{i\nu_Jt}$ with $\nu_J=2\pi J^{(c'-1)/2}$. Then there is
$J_0(c',b,\delta)$ such that $\mathfrak Q^{c'}[\varrho_J]<0$ for every $J\ge J_0$. Since $\check\Psi_{\varrho_J}$ is supported in
$(0,\log(J+1))$ and $\mathcal W$ is translation invariant, this is also a negative direction of the complete window form
with comb $c'v_n$ at $T_0=\frac12\log(J+1)$.
\emph{(iii)} For fixed $J$ and $c'>1$, positivity on the cell still holds for $\nu_0\ge2\pi J^{c'-1}e^{O(1)}$: the separation of
$c'v_n$ from $v_n$ is the uniformity in the number $J$ of active dilates at slowly growing frequency, not a property of one
chirp family at fixed $J$.
\end{corollary}
\begin{proof}
(ii) On $[J,J+1)$, by Lemma~\ref{lem:ext-weight}(i),
$w^{c'}(t,\nu_J)=\frac{c'-1}2H_J\log J+F(J)+H_J\log\frac tJ-2(c'-1)H_1(J)$. With $H_J\le\log J+1$, $F(J)\le\gamma\log J+1$,
$\log\frac tJ\le\log2$ and $H_1(J)\ge\frac12\log^2J+\gamma_1$, this is $\le-\frac{c'-1}2\log^2J+O_{c'}(\log J)$, while
$|E|\le C(b,\delta)H_J\log(J+1)/\nu_J\to0$. Hence $\mathfrak Q^{c'}[\varrho_J]\le(-\frac{c'-1}2\log^2J+O(\log J))\|b\|^2+o(1)<0$ for large
$J$. (iii) $w^{c'}\ge H_J\log\frac{\nu_0}{2\pi}-2(c'-1)H_1(J)\ge H_J\log\frac{\nu_0}{2\pi}-(c'-1)H_J\log J-O_{c'}(\log J)$.
\end{proof}

\subsubsection*{Multi-cell chirps of prolate phase}

For $N\ge4$, $c=2\pi N$, $\eta\in(0,1)$ and $X\ge2$, the class $\mathcal L_{\eta,X}$ consists of $\varrho=ae^{i\phi_c}$ with
$a\in C_c^2((1+\eta,X))$ and $a(k)=a'(k)=0$ at every integer $k\in(1+\eta,X)$ (no truncation jumps); the data are not required to
lie in $\mathrm{PW}_N$, the edge value $f(\lambda)$ and the gauge play no role. Now the dilates overlap:
$\langle u_m,u_n\rangle=\int_{s>1}(mn)^{1/2}s\,a(ms)\bar a(ns)e^{i\Theta_{mn}(y)}\frac{ds}s$, $s=e^y$, with relative phase
$\Theta_{mn}(y)=\phi_c(ms)-\phi_c(ns)$.

\begin{lemma}[stationary points of the relative phase; proved]\label{lem:ext-phase}
Let $\varpi(t)=t(t^2-2)/(t^2-1)^{3/2}$, so that $\tau'(t)=c\,\varpi(t)$ and $\Theta_{mn}''=c[ms\,\varpi(ms)-ns\,\varpi(ns)]$.
\emph{(i)} $\varpi<0$ on $(1,\sqrt2)$, $\varpi$ is increasing on $(\sqrt2,\infty)$, $\varpi(2)=0.7698\ldots$, $\varpi\to1$.
\emph{(ii)} For $2\le m<n$ there is no stationary point: $\Theta_{mn}'=\tau(ms)-\tau(ns)$ has modulus $\ge0.77\,cs(n-m)$.
\emph{(iii)} For $m=1$ and $n\ge2$, $\Theta_{1n}$ is strictly concave with $|\Theta_{1n}''|\ge0.4\,cn$, and has exactly one stationary point,
at $s_n=\sqrt{1+n^{-2}}$; it pairs $t=s_n$ (dilate $1$) with $t=ns_n=\sqrt{n^2+1}$ (dilate $n$), at the common Mellin frequency
$\tau_n=c(n^2+1)/n$, and $\Theta_{1n}''(y_n)=-c(n^4-1)(n^2+1)/n^3$.
\emph{(iv)} At the stationary point the log-symbol weight of the pair after the comb is
$\frac12\log\frac{\tau_n^2}{(2\pi)^2}-\log n=\log\big(N(1+n^{-2})\big)$.
\end{lemma}
\begin{proof}
(i) The logarithmic derivative of $\varpi$ has numerator $t^2+2>0$ on $(\sqrt2,\infty)$. (ii) Mean value theorem, $ms\ge2$ and
$\tau'\ge0.77c$ on $[2,\infty)$. (iii) For $s<\sqrt2$, $s\varpi(s)<0$; for $s\ge\sqrt2$, $t\varpi(t)$ is increasing; in both cases
$\Theta_{1n}''<0$, with the stated lower bound from $ns\ge2$. $\tau(s)=\tau(ns)$ is $n^2s^2-1=n^4(s^2-1)$, i.e. $s^2=1+n^{-2}$; then
$\tau(s_n)=c(1+n^{-2})n$, $s_n\tau'(s_n)=-c(n^4-1)/n$ and $\sqrt{n^2+1}\,\tau'(\sqrt{n^2+1})=c(n^4-1)/n^3$. (iv) Direct.
\end{proof}

\begin{theorem}[multi-cell chirps; proved as an asymptotic statement]\label{thm:ext-multi}
For $\varrho=ae^{i\phi_c}\in\mathcal L_{\eta,X}$,
\[
\mathfrak Q[\varrho]=\int_1^\infty\mathfrak w_N(t)|a(t)|^2dt+\mathcal R_{\mathrm{res}}+E,\qquad
|\mathcal R_{\mathrm{res}}|\le25.3\,c^{-1/2}H(X)\,(\beta_\eta\|a\|+2L_XK_a)\|a\|,
\]
where $\mathcal R_{\mathrm{res}}$ is the sum of the $(1,n)$ pairs, $K_a=\|a\|+\|ta'\|$, $\beta_\eta=\max(1,(2\eta+\eta^2)^{-1})$,
$L_X=\log(NT_\eta X)+2\log X$, $T_\eta=(1+\eta)(2\eta+\eta^2)^{-1/2}$, and $|E|\le C(\eta,X,a)\,c^{-1}$ with an explicit constant
polynomial in $X$. For a fixed profile $a$ (independent of $N$) and fixed $\eta$, $X$:
$\mathfrak Q[\varrho]\ge(\log2N-R(N))\|a\|^2$ with $R(N)=O_{\eta,X,a}(N^{-1/2}\log N)$, hence $\mathfrak Q[\varrho]>0$ for
$N\ge N_0(\eta,X,a)$.
\end{theorem}
\begin{proof}[Sketch]
\emph{Pair expansion.} $\check\Psi=\sum_{n<X}n^{-1/2}u_n$ and
$\mathfrak Q[\varrho]=\sum_{m,n}(mn)^{-1/2}[\langle m(D)u_m,u_n\rangle-(\log m+\log n)\langle u_m,u_n\rangle]+2\Rel(A\bar B)$ with
$A=\int\varrho\lfloor t\rfloor dt/t$, $B=\int\varrho H(t)dt$ (Theorem~\ref{thm:ext-half} and $\Lambda*1=\log$). Under the vanishing at
integers each truncated dilate $u_m$ is $C^1$ with $L^1$ second derivative, so Lemma~\ref{lem:ext-osc} applies to it.
\emph{Diagonal.} $\sum_mm^{-1}\mathbf 1_{t>m}=H(t)$ and $\sum_mm^{-1}\log m\,\mathbf 1_{t>m}=H_1(t)$: the harmonic cut reproduces the
weight $H(t)\log\frac{\tau(t)}{2\pi}-2H_1(t)=\mathfrak w_N(t)$ exactly, the cancellation of Theorem~\ref{thm:ext-cell} holding cell by
cell; the error is $O(K_a\|a\|/c)$ since $\sum_m1/(m\Omega_m)\le1.15/c$. \emph{Off-diagonal.} Replace $m(D)$ by $\log(\tau(ms)/2\pi)$
on each dilate (operator form of Lemma~\ref{lem:ext-osc}, error $O(\sqrt X K_a\|a\|/c)$); the pair $\{m,n\}$ becomes
$2(mn)^{-1/2}\Rel\int G_{mn}e^{i\Theta_{mn}}dy$ with
$G_{mn}=[\frac12\log\frac{\tau(ms)\tau(ns)}{(2\pi)^2}-\log m-\log n]\mathcal A_m\bar{\mathcal A}_n$, $\mathcal A_m=(ms)^{1/2}a(ms)$.
For $2\le m<n$ one integration by parts with Lemma~\ref{lem:ext-phase}(ii) gives $O(c^{-1})$ per pair, summable to $O(X\,L_X/c)$.
For $(1,n)$, the second-derivative van der Corput lemma \cite[Ch.~VIII, \S1.2]{Stein1993} with $|\Theta_{1n}''|\ge0.4cn$
(Lemma~\ref{lem:ext-phase}(iii)) gives $8(0.4cn)^{-1/2}\|G_{1n}'\|_1$, whence the bound on $\mathcal R_{\mathrm{res}}$. \emph{Pole:}
one integration by parts in $t$, $O(c^{-2})$. Finally $\int\mathfrak w_N|a|^2\ge\log2N\,\|a\|^2$ by Lemma~\ref{lem:ext-weight}(ii).
\end{proof}

The bounds are not uniform in $X$ ($\mathcal R_{\mathrm{res}}$ grows like $H(X)\log X$ and $E$ like $X\log X$), so the theorem
says nothing for data spread over all cells, and its constants are far from effective at $N=15$. The $(1,n)$ pairs are the only
$c^{-1/2}$ cross-dilate effects; their smallness is a regularity property of the amplitude on the Fresnel scale $(cn^3)^{-1/2}$.
The analogue of Corollary~\ref{cor:ext-cfals}(ii) holds inside $\mathcal L_{\eta,X}$ for $N\ge N_0(c')$, by taking a profile supported
in a single cell.

\subsubsection*{The caustic layer of a single prolate tail}

Let $\psi_k$ ($k$ even) be the prolate function of bandwidth $c=2\pi N$ on $[-1,1]$, normalised in $L^2(\R)$, with concentration
$\chi_k=\chi_k(c)$ and differential eigenvalue $\vartheta_k=c(2k+1)+O(k^2)$; then $\mathcal F[\psi_k(\cdot/N)\mathbf 1_{[-N,N]}]$ restricted to
$(1,\infty)$ is a constant multiple of $\psi_k|_{(1,\infty)}$, the one-sided exterior mass is $\frac12(1-\chi_k)$, and we take $\varrho_k=\psi_k|_{(1,\infty)}$ (the constraints defining $\mathrm{PW}_N$ are not imposed). The layer is $1<t<2$, where $H=1$ and $H_1=0$.

\begin{lemma}[the uniform prolate envelope; proved]\label{lem:ext-envelope}
Fix $c=2\pi N$ and an even $k$ with differential eigenvalue $\vartheta=\vartheta_k$, and assume the hypothesis
\[
(H_\nu)\qquad \nu:=\vartheta_k/c^2\le\tfrac14,\qquad c\ge200.
\]
(By the cited asymptotic $\vartheta_k=c(2k+1)+O(k^2)$, $(H_\nu)$ holds for every fixed $k$ once $c\ge c_0(k)$, and uniformly for
$k\le k_0(c)\asymp c/8$.) Write $p=t^2-1$, $q=c^2t^2-\vartheta$, $P=pq$ on $t>1$, so the prolate equation is $(p\psi')'+q\psi=0$.
Put
\[
a=P^{1/4}\psi,\quad b=P^{-1/4}p\psi',\quad G=a^2+b^2=\sqrt P\,\psi^2+\frac{(p\psi')^2}{\sqrt P},\quad E=\frac G{\sqrt P}=\psi^2+\frac{(p\psi')^2}P,
\]
and $a=\sqrt G\cos\phi$, $b=-\sqrt G\sin\phi$, $\Omega=\sqrt{q/p}$, $\varepsilon=P'/(4P)$, $\rho=\varepsilon/\Omega$. Then:
\emph{(a) (adiabatic invariant; exact identities)}
\[
(\log G)'=2\varepsilon\cos2\phi,\qquad \phi'=\Omega-\varepsilon\sin2\phi,\qquad E'=-\frac{(p\psi')^2P'}{P^2}\le0,
\]
so $E$ is nonincreasing with $E(1^+)=\psi(1)^2$, hence $|\psi(t)|\le|\psi(1)|$ for every $t\ge1$, and $\psi(t)^2=(G(t)/\sqrt{P(t)})\cos^2\phi(t)$.
\emph{(b) (explicit envelope)} $G_\infty:=\lim_{t\to\infty}G(t)$ exists, and for every $t>1$, with
$\sigma(t):=(1.45/\sqrt{t^2-1}+2.77)/c$,
\[
G(t)=G_\infty\,e^{x(t)},\qquad |x(t)|\le\sigma(t).
\]
Consequently, with $C_c^2:=G_\infty/c$: for $z:=c\sqrt{t^2-1}\ge2$ and $c\ge200$,
\[
|\psi(t)|\le1.57\,C_c\,(t^2(t^2-1))^{-1/4}\qquad\text{(Langer envelope)},
\]
and $\psi(1)^2\le35\,G_\infty$, whence for $1\le t\le2$,
\[
|\psi(t)|\le10.3\,C_c\sqrt c\,(1+z)^{-1/2}\qquad\text{(uniform Bessel-zone bound).}
\]
\emph{(c) (normalisation)} $C_c^2=\frac2\pi(1-\chi_k)(1+e_1)$ with $e_1$ as in part (d) below; in particular
$C_c^2=\frac2\pi(1-\chi_k)(1+O(\nu+c^{-1}\log c))$.
\emph{(d) (layer diagonal with explicit error)} Under $(H_\nu)$,
\[
\int_1^2\mathfrak w_N|\psi_k|^2\,dt=\tfrac12(1-\chi_k)\,D_{\mathrm{lay}}(N)\,(1+\eta),\qquad
|\eta|\le\frac{21.4\log c+1.54\,\vartheta_k/c+523}{c\,(1-|e_1|)},\quad
|e_1|\le0.77\nu+\frac{5.14\log(c/2)+106.2}c.
\]
\end{lemma}
\begin{proof}[Sketch]
(a) Differentiating $a,b$ along the ODE gives $a'=\varepsilon a+\Omega b$, $b'=-\varepsilon b-\Omega a$, from which the two
identities for $(\log G)'$ and $\phi'$ follow by the polar substitution; $E'$ is a direct computation, and $P'>0$ since $q>0$
under $\nu<1$, so $E$ is nonincreasing; $(p\psi')^2/P=p\psi'^2/q\to0$ at $t=1^+$ gives $E(1^+)=\psi(1)^2$.
(b) From $(\sin2\phi)'=2\cos2\phi\,(\Omega-\varepsilon\sin2\phi)$ one eliminates $\cos2\phi$ from the identity for $(\log G)'$,
integrates by parts on $[t,\infty)$ (using $\rho\to0$ at infinity) and bounds the resulting total variation and remainder
integral using elementary monotonicity properties of $t\,P^{-1/2}$ and of $c^2t\sqrt p\,q^{-3/2}$, together with the integrals
$\int_t^\infty sp^{-3/2}ds=p(t)^{-1/2}$, $\int_1^\infty ds/(s\sqrt p)=\pi/2$, $\int_1^\infty\sqrt p\,s^{-3}ds=\pi/4$; summing the
resulting bounds gives $|x(t)|\le\sigma(t)$ for $\gamma:=(1-\nu)^{-1/2}\le2/\sqrt3$. The Langer bound follows from
$\psi^2\le E=G/\sqrt P$ and $\sqrt\gamma\,e^{\sigma(t)/2}\le1.57$ for $z\ge2$, $c\ge200$. The constant $35$ follows from a
second application of the adiabatic estimate to the arc-length variable $\zeta=\int_1^t\Omega$, comparing $E$ at $\zeta=0$ and
at a fixed $\zeta_1=0.8$ via $|\psi_\zeta|\le\zeta|\psi(1)|$ and the bound $\zeta A_\zeta/A\le1.16$ on $\zeta\le1$ ($A=\sqrt P$),
giving $\psi(1)^2\le35\,G_\infty$; the Bessel-zone bound follows from (2.3) and this constant.
(c),(d) Integrating $\psi^2=(G/\sqrt P)\cos^2\phi$ against $W\in\{1,\mathfrak w_N\}$, writing $\cos^2\phi=\frac12+\frac12\cos2\phi$ and
using the identity of (a) once more (an integration by parts, "oscillation averaging") turns $\int W\cos2\phi/\sqrt P$ into a
boundary term plus an $O(c^{-1}\log c)$ remainder; comparing the resulting two mass identities ($W=1$ over $(1,\infty)$ gives
$\tfrac12(1-\chi_k)$, by the definition of $\chi_k$ as the $L^2(\R)$ concentration and $\varrho_k=\psi_k|_{(1,\infty)}$; $W=\mathfrak w_N$
over $(1,2)$ gives the layer integral) to the common model $\frac{G_\infty}{2c}\int W\,d\theta$ ($\theta=\operatorname{arcsec}t$)
and adding the Bessel-zone contribution (bounded via $\psi(1)^2\le35G_\infty$) and the $e^x-1$ error from (b) gives the stated
$e_1,\eta$ after collecting all five error sources (the $\nu$-correction $(1-\nu)^{-1/2}-1\le0.77\nu$, the missing arc
$\theta<\theta_B\le2/c$, the envelope error $e^x-1$, the oscillation remainder, and the Bessel zone).
\end{proof}

\begin{theorem}[the caustic layer; (i)--(iii) proved under the hypothesis $(H_\nu)$ of Lemma~\ref{lem:ext-envelope}, (iv) open]\label{thm:ext-layer}
\emph{(i) (exact; proved)} $\frac2\pi\int_1^2\mathfrak w_N(t)\frac{dt}{t\sqrt{t^2-1}}=D_{\mathrm{lay}}(N):=\frac23\log4N-\frac{\mathrm{Cl}_2(2\pi/3)}{2\pi}$,
with $\mathrm{Cl}_2(2\pi/3)=\sum_{k\ge1}\sin(2\pi k/3)/k^2=0.6766277\ldots$, so $\mathrm{Cl}_2(2\pi/3)/(2\pi)=0.107689\ldots$.
\emph{(ii) (local form at $t=1$; proved)} $t=1$ is a regular singular point of the prolate equation
$((t^2-1)\psi')'+(c^2t^2-\vartheta_k)\psi=0$ with indicial exponents $(0,0)$; the regular solution is locally
$J_0(\sqrt{2(c^2-\vartheta_k)(t-1)})$, so the transition zone has width $t-1\asymp c^{-2}$ and there is no Airy point on $t\ge1$.
\emph{(iii) (layer diagonal; proved under $(H_\nu)$: $\vartheta_k\le c^2/4$ and $c\ge200$)} By
Lemma~\ref{lem:ext-envelope}(a)--(b) the adiabatic invariant $G=\sqrt P\psi^2+(p\psi')^2/\sqrt P$ is, up to the explicit factor
$e^{x(t)}$, $|x|\le\sigma(t)$, equal to its limit $G_\infty=cC_c^2$ for every $t>1$; no asymptotic expansion of $\psi_k$ and no
citation of a uniform Bessel--Langer expansion is used. Then, under $(H_\nu)$, for every even $k$ with $\vartheta_k\le c^2/4$,
\[
\int_1^2\mathfrak w_N|\varrho_k|^2dt=\tfrac12(1-\chi_k)\,D_{\mathrm{lay}}(N)\,(1+\eta),\qquad
|\eta|\le\frac{21.4\log c+1.54\,\vartheta_k/c+523}{c\,(1-|e_1|)}\quad(\text{Lemma~\ref{lem:ext-envelope}(d)}),
\]
that is, $(1-\chi_k)\big[\frac23\log4N-0.107689\big](1+\eta)$ when both sides $\pm t>1$ are counted. For fixed $k$ and
$N\to\infty$ this recovers $(1+o(1))$ unconditionally; the explicit bound on $\eta$ is non-vacuous (i.e.\ $<1$) only for
$c\gtrsim820$ ($N\gtrsim130$), so any use of this diagonal at $N\le130$ remains outside the proved range and is
NUM\'ERICO, not PROVED, there.
\emph{(iv) (resonances)} Under the same expansion, at $s_n$ and $T_n=\sqrt{n^2+1}$ the amplitudes satisfy exactly
$s_n^{1/2}|\mathcal A(s_n)|\cdot T_n^{1/2}|\mathcal A(T_n)|=C_c^2$ for every $n\ge2$ ($\mathcal A=C_c(t^2(t^2-1))^{-1/4}$), and the
stationary-phase leading term of the pair $(1,n)$, per unit one-sided mass, is
\[
\widehat R_n(k)=\frac4\pi N^{-1/2}\log\big(N(1+n^{-2})\big)\,w_n\cos\big(\Psi_n-\tfrac\pi4\big),\qquad w_n=\frac n{(n^2+1)\sqrt{n^2-1}},
\]
$\Psi_n=-c\frac{n^2-1}n+\frac{2k+1}2(\frac\pi2-2\arctan\frac1n)+O(k^2/c)$; $w_2=0.23094\ldots$, $w_3=0.10607\ldots$,
$\sum_{n\ge2}w_n=0.61738\ldots$. The remainder of this stationary-phase step is not controlled (not proved). At leading order
$\sum_n|\widehat R_n|\le\frac4\pi N^{-1/2}\sum_nw_n\log(N(1+n^{-2}))<D_{\mathrm{lay}}(N)$ for every $N\ge4$, and the bracket
$D_{\mathrm{lay}}+\sum_n\widehat R_n$ grows like $\frac23\log4N$.
\end{theorem}
\begin{proof}[Sketch]
(i) With $t=\sec\theta$, $\theta\in(0,\pi/3)$: $dt/(t\sqrt{t^2-1})=d\theta$ and $\mathfrak w_N=\log2N-\log\sin2\theta$; then
$\int_0^{\pi/3}(-\log\sin2\theta)d\theta=-\frac12\int_0^{2\pi/3}\log\sin\varphi\,d\varphi=\frac\pi3\log2-\frac14\mathrm{Cl}_2(\frac{2\pi}3)$, using
$\int_0^x\log\sin\varphi\,d\varphi=-x\log2-\frac12\mathrm{Cl}_2(2x)$ and $\mathrm{Cl}_2(\frac{4\pi}3)=-\mathrm{Cl}_2(\frac{2\pi}3)$; multiply by $\frac2\pi$.
The $(t-1)^{-1/2}$ growth of the envelope is exactly the Jacobian $d\theta$; nothing diverges. (ii) With $s=t-1$ the equation is
$2s\psi''+2\psi'+(c^2-\vartheta_k)\psi=O(s)$, and $J_0(\sqrt{2\kappa s})$ solves $sw''+w'+\frac\kappa2w=0$. (iii) This is exactly
Lemma~\ref{lem:ext-envelope}(d), with the oscillation-averaging and Bessel-zone errors bounded there rather than
discarded as $o(1)$; combined with (i) this gives the stated $\eta$. (iv) $s_n^{1/2}|\mathcal A(s_n)|=C_cn^{1/2}$ and $T_n^{1/2}|\mathcal A(T_n)|=C_cn^{-1/2}$
since $s_n^2-1=n^{-2}$, $T_n^2-1=n^2$; then stationary phase \cite[Ch.~VIII]{Stein1993} with the Hessian of
Lemma~\ref{lem:ext-phase}(iii), the weight of Lemma~\ref{lem:ext-phase}(iv), $\sqrt{2\pi/c}=N^{-1/2}$, and
$\operatorname{arcsec}s_n=\arctan\frac1n$, $\operatorname{arcsec}T_n=\arctan n$. The leading-order comparison is
arithmetic on these constants.
\end{proof}

With comb weight $c'$ the layer diagonal is unchanged ($H_1=0$ on $(1,2)$) and the resonant weight becomes
$\log N+(1-c')\log n+\log(1+n^{-2})$ (algebra; proved); the extra term is summable against $w_n$. So for one prolate tail the layer
does not distinguish $c'v_n$ from $v_n$ at leading order; the dependence on the comb sits in the cells $t\ge2$
(Corollary~\ref{cor:ext-cfals}). Theorem~\ref{thm:ext-layer} concerns the $m=1$ line of the exterior datum only: the exterior
datum of an admissible function carries further lines, described next, to which single-tail estimates do not apply.

\subsubsection*{The M\"obius lattice of the exterior datum}

For $g$ in the core put $h(v)=v^{-1/2}g(\log v)$ on $[1/\lambda,\lambda]$, $\mathsf p(v)=h(|v|)\mathbf 1_{[1/\lambda,\lambda)}(|v|)$ (even) and
$\sigma_g=\mathcal F\mathsf p$. Since $g$ and $g'$ vanish at $\pm T$, $\mathsf p$ is $C^1$ and piecewise $C^2$, so $\sigma_g(x)=O(\min(1,x^{-2}))$.

\begin{theorem}[M\"obius-lattice structure; proved]\label{thm:ext-lattice}
Let $f$ be admissible at $T$ with $Ef=g$ in the core. Then
\begin{equation}\label{eq:ext-lattice}
f(u)=\sum_{1\le m<N}\mu(m)\,\mathsf p(mu)+\varphi_0(u),\qquad \supp\varphi_0\subset[-1/\lambda,1/\lambda],
\end{equation}
\begin{equation}\label{eq:ext-lattice-r}
r_f(\omega)=\sum_{1\le m<N}\frac{\mu(m)}m\,\sigma_g\Big(\frac\omega m\Big)+\mathcal F\varphi_0(\omega)\qquad(\omega>\lambda).
\end{equation}
At $u=\lambda/m$ ($m<N$ squarefree) the jet of $f$ jumps by $\mu(m)$ times the jet of $h$ at $\lambda$ dilated by $m$, so the line
$m$ of $r_f$ carries the phase $e^{\pm2\pi i\lambda\omega/m}$ (frequency $c/m$ in the universal frame); all lower cuts and
$\mathcal F\varphi_0$ carry the single phase $e^{\pm2\pi i\omega/\lambda}$.
\end{theorem}
\begin{proof}
For $1/\lambda\le u\le\lambda$, $\mathsf p(mu)\ne0$ only if $mu<\lambda$ (and $mu\ge u\ge1/\lambda$ is automatic), so
$\sum_m\mu(m)\mathsf p(mu)=\sum_{m<\lambda/u}\mu(m)(mu)^{-1/2}g(\log mu)=f(u)$ by Lemma~\ref{lem:ext-mob}; $m<\lambda/u\le N$. For
$|u|>\lambda$ both sides vanish, and for $m\ge N$, $\mathsf p(mu)=0$ when $|u|\ge1/\lambda$. So $f-\sum_{m<N}\mu(m)\mathsf p(m\cdot)$
vanishes off $[-1/\lambda,1/\lambda]$. \eqref{eq:ext-lattice-r} is $\mathcal F[\mathsf p(m\cdot)](\omega)=m^{-1}\sigma_g(\omega/m)$. The jets: the
$m$-th term switches off at $mu=\lambda$.
\end{proof}

Line $m\ge2$ samples $\sigma_g$ at $\omega/m\in(\lambda/m,\infty)$, below the threshold $\lambda$: it reads the bulk of $\mathcal F\mathsf p$,
not its tail, which is why these lines are of order one and are absent from every single prolate tail.

\begin{proposition}[the grouped identity; proved]\label{prop:ext-grouped}
For $y>T$, with $\check\Psi_0=(Z\mathcal F\varphi_0)(e^\cdot)$ and its logarithmic analogue $\check\Psi_{0,\log}$,
\[
\check\Psi_f(y)=e^{y/2}\sum_{q\in\Q_{>0}}\mathsf c_q\,\sigma_g(qe^y)+\check\Psi_0(y),\qquad
(Z_{\log}r_f)(e^y)=e^{y/2}\sum_{q\in\Q_{>0}}\mathsf d_q\,\sigma_g(qe^y)+\check\Psi_{0,\log}(y),
\]
where, for $q=b/a$ in lowest terms and $\widetilde M_a(Y)=\sum_{k<Y,(k,a)=1}\mu(k)/k$,
\[
\mathsf c_{b/a}=\sum_{k<N/a}\frac{\mu(ak)}{ak}=\frac{\mu(a)}a\,\widetilde M_a\Big(\frac Na\Big),\qquad
\mathsf d_{b/a}=\sum_{k<N/a}\frac{\mu(ak)\log(bk)}{ak}=\log b\cdot\mathsf c_{b/a}+\frac{\mu(a)}a\sum_{\substack{k<N/a\\(k,a)=1}}\frac{\mu(k)\log k}k .
\]
In $\check\Psi_f(y)=e^{y/2}\sum_{n'}r_f(n'e^y)$, line $m$ in the summand $n'=1$ and line $nm$ in the summand $n'=n$ are multiples of the same function $e^{y/2}\sigma_g(e^y/m)$, with product coefficient $\mu(m)\mu(nm)/(nm^2)$, which is $\mu(m)^2\mu(n)/(nm^2)$ if $(m,n)=1$ and
$0$ otherwise. Moreover $|\widetilde M_1(Y)|\le1$ for all $Y$, and for every prime $P$,
$\sum_{k<Y,P\nmid k}\mu(k)/k=\sum_{j\ge0}P^{-j}\widetilde M_1(Y/P^j)$.
\end{proposition}
\begin{proof}
Insert \eqref{eq:ext-lattice-r} into $(Zr_f)(e^y)=e^{y/2}\sum_nr_f(ne^y)$ (here $ne^y>\lambda$ for all $n\ge1$, so no truncation
enters); the double sum converges absolutely because $\sigma_g=O(\min(1,x^{-2}))$. Write $(m,n)=(ak,bk)$ with $(a,b)=1$; the only
constraint is $ak<N$, and $\mu(ak)=\mu(a)\mu(k)$ if $(k,a)=1$, $0$ otherwise. For $Z_{\log}$ use $\log n=\log b+\log k$. Coincident
lines: $\frac{\mu(nm)}{nm}\sigma_g(\frac{ne^y}{nm})=\frac{\mu(nm)}{nm}\sigma_g(\frac{e^y}m)$. Bound: $\sum_{k\le Y}\mu(k)\lfloor Y/k\rfloor=1$
gives $Y\widetilde M_1(Y)=1+\sum_{k\le Y}\mu(k)\{Y/k\}$, and the last sum has modulus $\le\{Y\}+\lfloor Y\rfloor-1=Y-1$ (apply at $Y-0$ for the
strict inequality). The last identity is $\mu\mathbf 1_{P\nmid\cdot}=\mu*\mathbf 1_{\{P^j\}}$.
\end{proof}

\begin{proposition}[exact reduction on the lattice class; proved]\label{prop:ext-latred}
Let $f,g$ be as in Theorem~\ref{thm:ext-lattice}. Put $\Upsilon_m=(Z\sigma^{(m)})(e^\cdot)\mathbf 1_{(T,\infty)}$ with
$\sigma^{(m)}=m^{-1}\sigma_g(\cdot/m)|_{(\lambda,\infty)}$ for squarefree $m<N$, $\Upsilon_0=\check\Psi_0\mathbf 1_{(T,\infty)}$, let $G$ be the
Gram matrix of the polarised form $\mathcal W$ on $\{\Upsilon_m\}\cup\{\Upsilon_0\}$, and $v=(\mu(m))_m\oplus1$. Then
\[
W_T(g)=\mathcal W[\check\Psi_f]=v^{\mathsf T}Gv,
\]
and for this $g$ the inequality \eqref{eq:ext-starstar} is equivalent to $v^{\mathsf T}Gv\ge2\kappa N^{-A_0}\mathrm{dist}^2(r_f,\overline{\mathcal G}_T)$.
Equivalently, with $s_q=e^{y/2}\sigma_g(qe^y)\mathbf 1_{y>T}$ and $\mathcal B_\infty$ the archimedean-plus-pole bilinear form,
$\mathcal W[\check\Psi_f]=\sum_{q,q'}\big[\mathsf c_q\mathsf c_{q'}\mathcal B_\infty(s_q,s_{q'})-(\mathsf c_q\mathsf d_{q'}+\mathsf d_q\mathsf c_{q'})\langle s_q,s_{q'}\rangle\big]$
plus the terms involving $\varphi_0$.
\end{proposition}
\begin{proof}
$\check\Psi_f=\sum_m\mu(m)\Upsilon_m+\Upsilon_0$ by \eqref{eq:ext-lattice-r} and linearity of $Z$; apply Theorem~\ref{thm:ext-half} and
bilinearity; the grouped form is Proposition~\ref{prop:ext-grouped} inserted in Theorem~\ref{thm:ext-half}.
\end{proof}

\begin{remark}[three mechanisms excluded on the lattice class; proved]\label{rem:ext-latno}
\emph{(a)} The archimedean Gram matrix $(\langle m(D)s_q,s_{q'}\rangle)$ is not positive semidefinite in general: $m(0)=C_0<0$, and
only the L\'evy part is positive. \emph{(b)} Positivity of $G$ is \eqref{eq:ext-starstar} on the span of the $\Upsilon_m$, a stronger
statement than needed, not an independent input. \emph{(c)} The size bounds and limits of the class coefficients do not
detect the comb: for the system with the Euler factor at $P$ deleted, Theorem~\ref{thm:ext-lattice} and
Proposition~\ref{prop:ext-grouped} hold verbatim with $\mu$ replaced by $\mu\mathbf 1_{P\nmid\cdot}$, and the integer-class coefficient
$\sum_{k<Y,P\nmid k}\mu(k)/k$ is bounded by $(1-P^{-1})^{-1}$ and tends to $0$ (by the last identity of
Proposition~\ref{prop:ext-grouped} and $\widetilde M_1(Y)\to0$ \cite{MontgomeryVaughan2007}); for the system with Dirichlet series
$\zeta^{c'}$, $0<c'\ne1$, the M\"obius inverse has $\mu_{c'}(p^2)=c'(c'-1)/2\ne0$, so non-squarefree lines appear and the coincident
pair $(1,4)$ acquires the sign of $c'(c'-1)$. Hence any argument using only $|\mathsf c_q|\lesssim1/a$ or $\mathsf c_b\to0$ applies equally
to these systems; the sign has to come from the non-integer classes $a\ge2$ and their coupling to $\varphi_0$, through the shape of
$\sigma_g$. Whether $v^{\mathsf T}Gv\ge2\kappa N^{-A_0}\mathrm{dist}^2$ holds is open.
\end{remark}

\begin{remark}[numerical observations at $N=15$; numerical, not used in any proof]\label{rem:ext-latnum}
On the bottom eigenvector $g^*$ of the certified matrices of \cite{Mori15} ($T=\frac12\log15$, even sector), evaluated in
ball arithmetic: the constant part $C_0\|g^*\|^2$, the L\'evy part, the pole and the prime part of $W_T(g^*)$ are $-5.372$, $+3.916$, $+1.519$, $-0.063$ (they cancel to the bottom value, of order $10^{-70}$); $f^*=E^{-1}g^*=1.106\,\psi_0-1.822\,\psi_4+O(10^{-5})$ in the prolate basis of $[-\lambda,\lambda]$; the ratio
$W_T(g^*)/\ell(g^*)$ lies in $[0.56,1.09]$ and $W_T(g^*)/(1-\chi_4)=4.08$. On the exterior tail of $\psi_4$ alone the exterior form
equals $0.97$--$0.98$ times $\frac12(1-\chi_4)\int_1^\infty\mathfrak w_N\,\frac{2}{\pi}\frac{dt}{t\sqrt{t^2-1}}$ and is positive for
$N=10,15,20,30$. On the smooth low-mode part of $g^*$ the $(1,n)$ cross-dilate pairs for $n=2,3,5,6,7$ have the sign of $\mu(n)$,
while the pair $(1,4)$ ($\mu(4)=0$) is not small; Theorem~\ref{thm:ext-layer} bounds the $(1,2)$ pair of a single prolate tail by
$0.054$ of the diagonal, against $-0.48$ measured on $g^*$, so the observed resonant deficit is not a single-tail effect.
\end{remark}

\subsection{What a proof of the leakage lemma cannot be}\label{subsec:ext-constraints}

The following are consequences of results proved above or earlier in the paper; each excludes a class of arguments for
\eqref{eq:ext-L}.
\begin{enumerate}
\item \emph{Comb-robust premises on the window.} A premise about $W_T$ that remains true when the comb is moved by total
variation $\delta$ gives at most $\delta\le\lamW$ (Theorem~\ref{thm:geo-tv}). On the exterior form the comb is not a parameter
(Theorem~\ref{thm:ext-half}).
\item \emph{Identities plus Cauchy--Schwarz, Bochner or Schur bounds.} Every identity above (half-line identity, M\"obius
bijection, $\Lambda*1=\log$, pole--Chebyshev cancellation) holds for the comb $c\Lambda$, for prime deletions and for
Beurling systems with their own summation map; any argument built from them and such bounds is comb-generic.
\item \emph{Symbol calculus.} The on-line symbol of the full form, $m+2\Rel(\zeta'/\zeta)$, vanishes off the zeros
(Theorem~\ref{thm:ext-line}); all of $W_T(Ef)$ is the truncation of the summation at $ne^x=\lambda$.
\item \emph{Sums of squares by ratio class, per-harmonic signs, Hilbert kernels.} The diagonal ratio class is indefinite
(Proposition~\ref{prop:ext-pairs}); the harmonics are dilates of one function with forced coefficients
(Proposition~\ref{prop:ext-dual}); the archimedean pairing of dilates is a positive-definite Toeplitz kernel in $\log(n/m)$, not
an antisymmetric one.
\item \emph{Ingredients.} A proof must use together $a_n\equiv1$ on all of $\N$, residue $1$ at the pole, $\Gamma_\R$ with
conductor $1$, and the Paley--Wiener restriction of the data; without the last, \eqref{eq:ext-starstar} with $\kappa=0$ on all
half-line data is the positivity of the complete form on every window. It must fail for $c\Lambda$ with $c>1$ and for the
other falsifiers of \S\ref{sec:shape}.
\item \emph{Scale.} The upper side is linear in the leakage amplitude on prolate blocks \cite[Thm.~1.6]{Guo2026}; a lower bound
must be quadratic, with polynomial loss, against the distance to the gauge, never against the leakage of a fixed vector.
\item \emph{Flows in $T$.} By Proposition~\ref{prop:ext-flow}, a continuation argument is a statement about the band-edge alias
and closes only with an independent bound at the minimiser.
\item \emph{The defect half.} By Theorem~\ref{thm:ext-mirror}, positivity of $-Q(G,G)$ is exterior positivity for the dual data;
it cannot be counted as a reduction.
\item \emph{The pole.} It couples only through rank-one residues (Proposition~\ref{prop:ext-residue}); the remaining couplings
$J_{c'}$ are full bilinear forms without sign.
\end{enumerate}

\section{Open problems}\label{sec:open}

The main problem is \eqref{eq:target} for $T>\frac12\log5$. We state it in the forms in which the results of
this paper reduce it, then list secondary problems.

\begin{enumerate}
\item[(L2)] \emph{The finite matrix.} For every $T>\frac12\log5$ prove $S_{T,k(T)}\succeq0$, where $k(T)\ge k_0(T)$ is
the number of blocks for which the archimedean form is coercive on $\mathcal V_k$ (Theorem~\ref{thm:coer-B}) and
$S_{T,k}$ is the Schur complement on the $(2k-2)$-dimensional space of piecewise Tate exponentials
(Theorem~\ref{thm:coer-schur}). By that theorem this is equivalent to \eqref{eq:target} at $T$, and by
Lemma~\ref{lem:firstcontact} to the absence of a first-contact window. By \eqref{eq:coer-sandwich} its least
eigenvalue satisfies $\lamA\le\lambda_S$ and, whenever $\lambda_S<c(T,k)$, $\lambda_S\le(1+\|\mathsf K\|)^2\lamA$; so for
$T\ge1$ Corollary~\ref{cor:margin1} gives $\lambda_S\le(1+\|\mathsf K\|)^2\cdot4\cdot10^9Te^{-15.449T}$ (no bound on
$\|\mathsf K\|$ uniform in $T$ is proved). The barriers of \S\ref{sec:barriers} and
Theorem~\ref{thm:wh-detinherit} exclude, each within the class of arguments stated there, norm bounds,
determinants, moments and a local analysis of the symbol; they do not exclude an argument outside those
classes. A first test case is
the window $T=\frac12\log5+\epsilon$ at which the contact $5$ enters.
\item[(N)] \emph{The signed twisted Chebyshev bound at Nyquist points.} By Proposition~\ref{prop:wit-coh}, a
primitive $F$ is a witness exactly when $\sum_k|\Fh(\omega_k)|^2\,\mathcal C_T(\omega_k)<-\mathfrak b_T(F)$,
$\omega_k=\pi k/2T$, where $\mathcal C_T$ is the coherence of the prime comb against its polar part. Prove, for
every $T$, a pointwise bound with sign for $\mathcal C_T(\omega_k)$ at the Nyquist points, correlated with the
admissible samples $(|\Fh(\omega_k)|^2)_k$ of $\PT$, strong enough to exclude this inequality. By
Remark~\ref{rem:wit-cheb} and Proposition~\ref{prop:wit-sens}, absolute-value bounds for $\psi(y)-y$ with any
exponent $\theta>\frac12$ and mean-square (large-sieve) bounds do not suffice, and no bounded functional of the
weights does; what is needed is cancellation of exponent $\frac12$ in the twisted sum, with sign, at about
$e^{2T}$ frequencies. A first case is the inequality $(E_\gamma)$ of Proposition~\ref{prop:wit-Egamma} for one fixed
$\gamma$, which is necessary for \eqref{eq:target}. In polar-free form (Proposition~\ref{prop:sharp}) the problem is the
inequality $\mathcal F_T\le\mathfrak a^\sharp$ on $\PT$, whose archimedean side has bottom $\lambda^\sharp(T)\downarrow C_0+4$; for the
packets the threshold is exactly $m^\sharp(\gamma)$ (Proposition~\ref{prop:Egamma-exact}).
\item[(S)] \emph{The near-coincidence matrix.} Fix the profile $F_0$ of \S\ref{subsec:oneprofile} with $2T_0<\log2$. Prove that
$\mathrm{Arch}_M-\mathrm{Div}_M-\mathrm{Near}_M\succeq0$ for every $M$; by Proposition~\ref{prop:AM} and Theorem~\ref{thm:oneprofile}
this is equivalent to \eqref{eq:target}. The three parts are known explicitly: $\mathrm{Arch}_M$ is the archimedean convolution
in the real lag $\log\frac mk$; $\mathrm{Div}_M$ is the Dirichlet convolution by $v_n=2\Lambda(n)n^{-1/2}$, whose restriction to
multiplicative boxes is bounded by the local Euler factors (Remark~\ref{rem:boxes}); and for $m>k$,
$\mathrm{Div}_M+\mathrm{Near}_M=\frac12\Xi_{F_0}(\log\frac mk)$, the smoothed Chebyshev fluctuation at the scale $\frac mk$ with its
main term removed exactly. On test spaces where near-coincidences $m\approx kn$, $m\ne kn$, are invisible, positivity holds
unconditionally (Proposition~\ref{prop:nearfree}). Each $2\times2$ section is the band of Theorem~\ref{thm:band}, and classical
progression statistics enter only as class properties excluded by Proposition~\ref{prop:prog}. What remains is a statement about
$\mathrm{Near}_M$, the part of the Toeplitz-in-$\log\frac mk$ matrix $[\Xi_{F_0}(\log\frac mk)]$ coming from integers that are close
to $\frac mk$ but not equal to it, tested jointly with $\mathrm{Arch}_M$ on vectors that see these near-coincidences.
\item[(R)] \emph{The admissible region.} By Theorem~\ref{thm:shape} and Corollary~\ref{cor:shape}, every valid proof of
\eqref{eq:target} lies in the following region, and nothing in this paper excludes an argument with all of these properties.
\begin{problem}[the admissible region]\label{prob:region}
Find an argument that proves $\lamA\ge0$ for every $T>\frac12\log5$ and that: (1) uses the exact values $\Lambda(n)$, $n<e^{2T}$,
at every $T$, and not membership in a class closed under any modification of Theorem~\ref{thm:shape}(A); (2) uses the Euler product
over $\N$ and the functional equation with $\Gamma(s/2)$ at conductor $1$ jointly; (3) is non-linear in $\Lambda$, that is, is not a
certificate, lattice certificate, flat-top certificate or operator identity with Cauchy--Schwarz, and uses arithmetic squares only with a
free positive definite part; (4) does not reduce to zero-free regions, zero densities or finitely many zeros, and, if it passes through
pointwise statements about the curvature of $\log|\xi|$ on the line, uses holomorphy and order-one growth jointly or the Euler product on
the line (Proposition~\ref{prop:shape-twist}); (5) is not a bounded functional of the weights or of the symbol, not local on the
negative set, not a leading term of the window determinant, not a pattern of window observables, and not a law on a bounded range of
windows; (6) is sensitive to changes of the weights at the scale $\lamA/\kappa_T$ and has precision at most $4\cdot10^9Te^{-15.449T}$
at $T\ge1$; (7) is exact at the centre and at the integer resolution, wins from the fluctuation of $\Lambda$ and not from its density;
(8) uses the positive part of $\sT$ on all of $\R$ jointly with the negative set, that is, correlation inside $\PT$ at all frequencies
up to $e^{2T}$; (9) is not a symmetrization of the family $D_\sigma$; (10) uses a property beyond the axioms A1--A6 of \cite{Paper46},
and is uniform over number fields only if it names the property separating $\zeta$ from every $\xi_K$; (11) is specific to $\zeta$ by
construction at the sign step: that step is not an inequality on prime sums with a remainder (Theorem~\ref{thm:shape}(B$'$)) and not an exact
identity valid for every weight vector (Theorem~\ref{thm:shape}(B$''$)); exact identities enter only as bookkeeping beside a property that
separates $\zeta$ from the modifications.
\end{problem}
Objects that this paper does not exclude and that lie in the region: two-sided laws $\lamA=L(T)(1+o(1))$ on unbounded ranges
(Remark~\ref{rem:bar-fam}); signed exterior measures and primitive certificates with $\nu\ge0$ and free $\beta$ (Remark~\ref{rem:bar-lin});
arithmetic squares on $(\frac12\log2,\frac12\log3)$ (Theorem~\ref{thm:sos}); the exact Schur reductions of Theorems~\ref{thm:coer-schur}
and \ref{thm:wh-schur}, which carry the whole margin; the problems (L2), (N) and (S) above, each an exact restatement of
\eqref{eq:target}. A certificate (class (D)) is admissible only as the final object of a computation that already carries the margin.
The identities of Propositions~\ref{prop:shape-ext}--\ref{prop:shape-phase} are coordinates in which such a computation may be written; by
Lemma~\ref{lem:shape-ident} none of them is an object of the region by itself.
\end{enumerate}

Secondary problems.
\begin{enumerate}
\item \emph{Gamma-moment annihilation at first contact} (Problem~\ref{prob:annihilate}). By
Proposition~\ref{prop:annihilate-equiv} its conclusion is equivalent to \eqref{eq:target}, and the return
inequalities are equivalent to it by \eqref{eq:equivgate}. It is a prescribed route to (L2), not a weaker
statement, and no intermediate estimate toward it is known.
\item \emph{Class-dependent value bounds} (Problem~\ref{prob:A}). The ceilings of Theorem~\ref{thm:R} exceed
$\lamP$ by at most $O(1/\log M)$ (Corollary~\ref{cor:Rrate}); a polynomial rate would follow from a Fourier tail
bound for the primal minimizer beyond its form domain, which is not known. What remains is the other direction, an
explicit deficit $D(T,d,U)$ with $c(T,d,U)\ge\lamP-D$ for the best lattice certificate, first for one contact and
then for two; no such lower bound is known at any $M$, and by Theorem~\ref{thm:Vp} the deficit must depend on
the class and not on a radius alone.
\item \emph{The polar capacity beyond $\frac12\log5$.} $\mathfrak G_T\succ0$ holds on an interval
$(\frac12\log5,\frac12\log5+\varepsilon]$ by continuity (Corollary~\ref{cor:beyond5}), with an $\varepsilon$ of no
practical size. What remains is a certified window $T>\frac12\log5$ with $\mathfrak G_T\succ0$, that is, with
$\lamW>0$ in the presence of the contact $5$; by Proposition~\ref{prop:onewindow} and Theorem~\ref{thm:parityschur}
this is the complete-form counterpart of (L2) at that window.
\item \emph{A certified lower bound for $S_{T,k}$} at a window where $\lamA$ is not already certified
(Remark~\ref{rem:coer-honest}(2)), and a certified version of the two-sector certificate of
Theorem~\ref{thm:nyq-schur} beyond $\frac12\log5$.
\item \emph{An asymptotic law with sign.} A two-sided asymptotic $\lamA=L(T)(1+o(1))$ with explicit $L>0$ is
stronger than \eqref{eq:target} and is not excluded by any barrier of \S\ref{sec:barriers}
(Remark~\ref{rem:bar-fam}). What is missing is its lower half, $\lamA\ge L(T)(1-o(1))$: by
Theorem~\ref{thm:wit-Tc}(iv) already the one-sided statement $\lamA\ge-\varepsilon(T)$ with $\varepsilon(T)\to0$ is
equivalent to \eqref{eq:target}, and beyond $\frac12\log5$ no lower bound better than the trivial
$\lamA\ge C_0-\Aconst$ (Lemma~\ref{lem:wit-basic}(v)) is known at any window. The
upper half is known only in the weak form of Corollary~\ref{cor:margin1}; the conjectural law
$-\log\lamP\asymp e^{2T}$ mentioned after Problem~\ref{prob:A} is an order-of-magnitude statement for the
logarithm, not a two-sided asymptotic with sign.
\item Whether the optimal multipliers of the primitive problem converge to $\beta=(2,0)$, i.e.\ whether
$\lamP=\lamW$ at a given window. By Theorem~\ref{thm:mult} this holds if and only if the moment vectors of the
minimizers of $W_T$ have $0$ in their convex hull, and it fails whenever $\lamW<\lamA$ and the two parity bottoms
of $W_T$ differ. At $T=\frac12\log3$ the Galerkin values indicate $\lamW<\lamA$ (Remark~\ref{rem:multnum},
numerical); whether the parity bottoms coincide there, or at any window, has not been examined.
\item \emph{Atomic positive combs.} Whether a Beurling generalized prime system (positive prime measure, prime powers
included) has a comb in $\mathfrak C^+$, and whether the doubly exponential spacing of Theorem~\ref{thm:atomic} is necessary
for an atomic comb to be positive at every window (Remark~\ref{rem:atomic}).
\item \emph{Arithmetic squares on the first window with a contact.} Theorem~\ref{thm:sos} leaves
$\frac12\log2<T<\frac12\log3$ open.
\end{enumerate}

\appendix

\section{Computer verifications}\label{app:encl}

This appendix describes the computer verifications used in Theorems~\ref{thm:C}, \ref{thm:N},
\ref{thm:Np}, \ref{thm:index}, \ref{thm:log5A}, \ref{thm:log5W}, Proposition~\ref{prop:archneg} and
Corollary~\ref{cor:convfail}. All of them are carried out in outward-rounded interval arithmetic
(\texttt{mpmath}), except those of Theorems~\ref{thm:log5A} and \ref{thm:log5W}, which use ball arithmetic
(\S\ref{subsec:encllog5}). The data (atoms, coefficient vectors, test kernels and interval endpoints) and the
scripts accompany the paper; each verification runs in at most a few minutes on one core with Python~3,
\texttt{mpmath}~1.3.0 and \texttt{scipy}~1.11.4, and the ball-arithmetic certificates use
\texttt{python-flint}~0.9.0 (Arb) and take about seven minutes per parity.

\subsection{The time--band limiting operator}
\begin{lemma}[a Ky Fan ceiling for the time--band limiting operator]\label{lem:KyFan}
Let $S_c=Q_TB_sQ_T$ on $L^2(\IT)$, where $Q_T$ is restriction to $\IT$ and $B_s$ the band projection
to $|\tau|<s$, and put $c=Ts$. Let $j_n$ be the spherical Bessel functions and
$G_k(x)=\sum_{j<k}(2j+1)j_j(x)^2\in[0,1]$. Then, in the orthonormal Legendre basis
$e_j(t)=\sqrt{(2j+1)/2T}\,P_j(t/T)$,
\[
\langle S_ce_j,e_j\rangle=\frac{2(2j+1)}\pi\int_0^cj_j(x)^2dx,\qquad
\sum_{j\ge0}\langle S_ce_j,e_j\rangle=\frac{2c}\pi=\operatorname{tr}S_c ,
\]
and for every $k\ge0$
\[
\lambda_k(S_c)\ \le\ \Lambda_k(c):=\frac2\pi\int_0^c\big(1-G_k(x)\big)\,dx .
\]
In particular $\Lambda_0(c)=2c/\pi$ is the trace and $\Lambda_k(c)=O(c^{2k+1})$ as $c\to0$.
\end{lemma}
\begin{proof}
From $\int_{-1}^1P_j(u)e^{-ixu}du=2(-i)^jj_j(x)$ one gets
$\widehat{e_j}(\tau)=\sqrt{(2j+1)/2T}\,\cdot2T(-i)^jj_j(\tau T)$, so
$|\widehat{e_j}(\tau)|^2=2T(2j+1)j_j(\tau T)^2$ and Plancherel gives the displayed diagonal entry.
The addition theorem $\sum_{j\ge0}(2j+1)j_j(x)^2=1$ gives the trace identity, so the basis loses no
trace. By the Ky Fan maximum principle $\sum_{j<k}\lambda_j(S_c)\ge\sum_{j<k}\langle
S_ce_j,e_j\rangle$; since $S_c=(B_sQ_T)^*(B_sQ_T)\ge0$ its eigenvalues are nonnegative and
non-increasing, whence
$\lambda_k\le\sum_{n\ge k}\lambda_n=\operatorname{tr}S_c-\sum_{j<k}\lambda_j\le\sum_{j\ge k}\langle
S_ce_j,e_j\rangle=\Lambda_k(c)$.
\end{proof}

\begin{corollary}[quadrature-free form]\label{cor:P1}
$|j_n(x)|\le x^n/(2n+1)!!$ for $x\ge0$, hence
$\lambda_k(S_c)\le\Lambda^*_k(c):=\frac2\pi\sum_{j\ge k}c^{2j+1}/\big((2j+1)!!\big)^2$, a convergent
series of rationals in $c$: the bound is evaluated by arithmetic alone, with no quadrature.
\end{corollary}


\subsection{The interval enclosures of Theorems~\ref{thm:N} and \ref{thm:Np}}
\emph{The window $T=\frac12\log3$.} All three integrals are computed in interval arithmetic at $25$ digits
(\texttt{mpmath}), the only non-interval inputs being the rational $\tau_0$ and the decimal $q$, which
are data of the bound and not approximations. The archimedean multiplier is summed from
Theorem~\ref{thm:S}(i), $m(\tau)=C_0+\sum_{k\ge0}\tau^2/(a_k(\tau^2+4a_k^2))$, with
$C_0=\psi(\frac14)-\log\pi=-\gamma-\frac\pi2-3\log2-\log\pi$ by the Gauss digamma theorem, the first
$100$ terms in interval arithmetic and the tail bracketed by the convex midpoint rule: with $A=K-\frac14$
and $t(a)=\tau^2/(\tau^2a+4a^3)$, the function $t$ is positive, decreasing and convex on $(0,\infty)$
(because $2u'^2-u''u=2\tau^4+24\tau^2a^2+192a^4>0$ for $u=\tau^2a+4a^3$), the points $a_k=k+\frac14$
are the midpoints of unit cells tiling $[A,\infty)$, so
$\sum_{k\ge K}t(a_k)\in[\Lambda-E,\Lambda]$ with $\Lambda=\frac12\log(1+\tau^2/4A^2)=\int_A^\infty t$
and $E=\frac1{24}|t'(A)|=\frac1{24}\tau^2(\tau^2+12A^2)/(\tau^2A+4A^3)^2$, at most $2.4\cdot10^{-7}$
on the range used. The multipliers are evaluated from the closed form
$\int_0^{2T}e^{su}\cos\tau u\,du=[e^{2Ts}(s\cos2T\tau+\tau\sin2T\tau)-s]/(s^2+\tau^2)$ at
$s=\pm\frac12$ and its $s$-derivative, without quadrature. The integrals over $\mathcal I_0$ are composite
Simpson with the classical remainder $Ph^4C_4/180$, the fourth derivatives being bounded by
$|m^{(j)}|\le2^{-j}j!\,\zeta(j+1,\frac14)$, $|(v_2\cos(\tau\log2))^{(j)}|\le v_2(\log2)^j$,
$|\gh0^{(j)}|\le2\cosh T(2T)^{j+1}/(j+1)$, $|\gh1^{(j)}|\le2\sinh T(2T)^{j+2}/(j+2)$ and the Bernstein
bound $|\vartheta^{(j)}|\le(d(M-1))^j(\sum_k|q_k|)^2$, combined by the Leibniz rule; the resolutions are
chosen so that the remainder is $7\cdot10^{-5}$ for $a$, which needs only $10^{-4}$, and below
$10^{-8}$ for $b_0,b_1$, which are multiplied by $B$. The denominator $\int_{\mathcal I_0}\vartheta=P\|q\|^2$ is
exact. An independent evaluation by adaptive quadrature and a library digamma, outside the interval
arithmetic, lands inside all three enclosures. Steps $d$ with $2T/d\notin\N$ are not covered by
Theorem~\ref{thm:N}: each such $d$ needs its own enclosure with $M'=\lceil2T/d\rceil$.

\emph{The windows $T=0.68$ and $T=\log2$.} The archimedean multiplier, the constant $C_0$, the closed forms of
$\gh0,\gh1$ and the exact denominator $P\|q\|^2$ are as in Theorem~\ref{thm:N}; nonnegativity of
$\vartheta$ is structural, since it is evaluated as a sum of two squares, so no property of the
particular $q$ is used. Two things change. First, Simpson's rule is replaced by composite
Gauss--Legendre of order $p$ on panels of width $H$, with the classical remainder
$|E_{\mathrm{panel}}|\le H^{2p+1}(p!)^4\big/\big[(2p+1)((2p)!)^3\big]\cdot\sup|f^{(2p)}|$ and the
derivative bounds of Theorem~\ref{thm:N} extended from order $4$ to order $2p$, the archimedean one
being $|m^{(j)}|\le2^{-j}j!\,\zeta(j+1,\frac14)\le2^{-j}j!\,(4^{j+1}+\zeta(j+1))$ for $j\ge1$ and
$|m|\le\max(|C_0|,m(\tau_{\max}))$ for $j=0$; the nodes are Newton-refined at $60$ digits, widened
before entering interval arithmetic, and guarded by the archived residuals
$\max_i|P_p(x_i)|$ and $|\sum_iw_i-2|$, both below $10^{-60}$ in every run; this guard ensures that the
computed rule is the one to which the remainder formula applies. Second, $q$ is optimised: rather than the bottom eigenvector of
$C_M(\tau_0,\beta^\star)$, which makes $a$ small and leaves $b$ where it falls, one minimises
$a+B(|b_0|+|b_1|)$ over the span of the three or four lowest eigenvectors, which drives
$|b_0|+|b_1|$ to the $10^{-16}$--$10^{-18}$ level while $a$ stays within a fraction of a percent of
the double-precision ceiling. Consequently the $\beta$-allowance $(|b_0|+|b_1|)B$ is below
$1.3\cdot10^{-11}$ in every row, four orders below the smallest $|a|$, and the binding error is the
quadrature on $a$, between $0\%$ and $6.7\%$ of $|a|$. The bound survives any $B\le4.7\cdot10^9$,
a factor $2.5\cdot10^5$ above the constants of Step~1, so only the qualitative content of
Theorem~\ref{thm:Bbeta} --- that some finite a priori box exists --- is used. At two values of $M$
the minimax over $\tau_0$ has two active period placements, and no single pair $(\tau_0,q)$ can then
drive $b$ to zero; a scan over the best period placements supplies one that closes, at the cost of
recovering only a few percent of the double-precision ceiling there. As in Theorem~\ref{thm:N},
steps $d$ with $2T/d\notin\N$ are not covered: each needs its own enclosure with
$M'=\lceil2T/d\rceil$.

\subsection{Verification on the bounded part}

\begin{proposition}[second-order cover]\label{prop:TAY}
Write $f=m+q+r$ with $q=-\sum_{n\in\NT}v_n\cos(\tau\log n)+\nuh$ ($\nu$ atomic) and
$r=\beta_0\gh0+\beta_1\gh1$, and put
$M_2=\sum_{n\in\NT}v_n\log^2n+2\sum_j|a_j|u_j^2+|\beta_0|H_0+|\beta_1|H_1$ with
$H_0=\int_{-2T}^{2T}u^2\cosh\frac u2\,du$, $H_1=\int_{-2T}^{2T}|u|^3\sinh\frac{|u|}2\,du$. Then
$\sup|q''+r''|\le M_2$, and on every cell $[x_0-\frac w2,x_0+\frac w2]\subset[0,\infty)$
\[
\min_{\text{cell}}f\ \ge\ m\!\big(x_0-\tfrac w2\big)+(q+r)(x_0)-|(q+r)'(x_0)|\,\tfrac w2-M_2\tfrac{w^2}8 .
\]
The same bound applies to each row $\nuh+Y_k$ of Lemmas~\ref{lem:F1}, \ref{lem:hybrid} with
$m(x+kP)$ and $-|\beta_i|k_i(x+kP)$ evaluated at the left endpoint. A cell is certified when the
right side exceeds the target and is bisected otherwise; near a contact of $f$ the admissible width
is $\approx\sqrt{8c/M_2}$, so the number of evaluations grows like $c^{-1/2}$ rather than the
$c^{-1}$ of a Lipschitz cover.
\end{proposition}
\begin{proof}
Taylor's formula with integral remainder for $q+r\in C^2$, and monotonicity of $m$ and of
$-|\beta_i|k_i$.
\end{proof}



\subsection{Grids and the band-limited minimum}
The following lemma quantifies how far the minimum over a grid can be from the true infimum: on a
$\delta$-grid a band-limited function can hide a dip of size $\asymp U^2\delta^2$ between two adjacent
nodes.

\begin{lemma}[Bernstein grid bound]\label{lem:bern}
Let $f:\R\to\R$ be bounded and band-limited to $[-U,U]$, i.e.\ its distributional Fourier transform
is supported in $[-U,U]$. Then for every $\delta>0$ and every $x_0\in\R$,
\[
\inf_{\tau\in\R}f(\tau)\ \ge\ \inf_{k\in\Z}f(x_0+k\delta)\ -\ \tfrac18U^2\|f\|_\infty\delta^2 .
\]
\end{lemma}
\begin{proof}
By Bernstein's inequality $\|f''\|_\infty\le U^2\|f\|_\infty=:M$. Put $g=f-\inf f\ge0$; $g$ is
bounded with $\|g''\|_\infty\le M$. For every $x$ and $t$, $0\le g(x+t)\le g(x)+g'(x)t+\frac M2t^2$,
and a nonnegative quadratic in $t$ has nonpositive discriminant, so $g'(x)^2\le2Mg(x)$. Fix
$\eta>0$ and $x$ with $g(x)\le\eta$; let $x_0+k\delta$ be a nearest grid point, so
$|x-(x_0+k\delta)|\le\delta/2$. Taylor's formula with integral remainder gives
\[
f(x_0+k\delta)\ \le\ f(x)+|f'(x)|\tfrac\delta2+\tfrac M2\Big(\tfrac\delta2\Big)^2
\ \le\ \inf f+\eta+\sqrt{2M\eta}\,\tfrac\delta2+\tfrac18M\delta^2 .
\]
Letting $\eta\downarrow0$ yields $\inf_kf(x_0+k\delta)\le\inf_\R f+\frac18M\delta^2$.
\end{proof}

For the certificate function $f=\sT+\nuh+\beta\cdot\widehat g$ the archimedean term is not
band-limited, and the verification of Theorem~\ref{thm:C} therefore uses no grid: the exhaustive
second-order cover of Proposition~\ref{prop:TAY} is the same Taylor estimate with the exact constant
$M_2$ in place of the Bernstein bound $U^2\|f\|_\infty$, and it certifies the infimum over $\R$, not over
a finite set of nodes. For a compactly supported flat-top weight $w$ with $\supp w\subset[-U,U]$,
every finite sum $\sum_{\gamma\in Z}\wh(\tau-\gamma)$ is bounded and band-limited to $[-U,U]$, and
Lemma~\ref{lem:bern} turns a grid evaluation of it into a two-sided statement.

\subsection{The certificates of Theorem~\ref{thm:index}}\label{subsec:enclindex}
In this subsection $\mathrm{par}\in\{\mathrm{ev},\mathrm{odd}\}$, $\chi_{\mathrm{ev}}(t)=\cosh\frac t2$ and
$\chi_{\mathrm{odd}}(t)=\sinh\frac t2$ on $\IT$, and $L^2_{\mathrm{par}}(\IT)$ is the space of real functions
of that parity, so that a real $F\in L^2_{\mathrm{par}}(\IT)$ lies in $\PT$ if and only if
$\langle F,\chi_{\mathrm{par}}\rangle=0$ (\S\ref{subsec:arch}). The basis $e_j$, the operators $S_c$ and the
tails $\Lambda_k$ are those of Lemma~\ref{lem:KyFan}. Every other letter is local to this subsection.

\begin{lemma}[bordered inertia]\label{lem:border}
Let $q$ be a real quadratic form on a real vector space $V$, $f$ a linear functional on
$V$ and $v_0\in V$ with $f(v_0)=1$ and $q(\cdot,v_0)$ finite. On $V\oplus\R$ put
$b(v,\xi)=q[v]+2\xi f(v)$. Then $n_-(b)=n_-(q|_{\ker f})+1$, $n_-$ being the largest dimension of a
subspace on which the form is negative definite.
\end{lemma}
\begin{proof}
Write $v=w+tv_0$ with $w\in\ker f$. Then $b=q[w]+2tq(w,v_0)+t^2q[v_0]+2\xi t$, and the invertible
substitution $\eta=\xi+q(w,v_0)$ gives $b=q[w]+(t^2q[v_0]+2t\eta)$: a $b$-orthogonal sum of
$q|_{\ker f}$ and a binary form of determinant $-1$, whose inertia is $(1,1)$.
\end{proof}

\begin{lemma}[multi-band minorant]\label{lem:multiband}
Let $0=\tau^{(0)}<\dots<\tau^{(N)}=\tau_N$, let $m^-$ be a lower bound for $m$ at the edges,
and let $M_a\ge\sup m''$ on the band $[\tau^{(a)},\tau^{(a+1)}]$, of width $\Delta_a$. Define the even
function $\underline m$ by $\underline m=C_0$ on $[0,\tau^{(1)}]$; on the band $a\ge1$, the chord through
$(\tau^{(a)},m^-(\tau^{(a)}))$ and $(\tau^{(a+1)},m^-(\tau^{(a+1)}))$ lowered by $\max(M_a,0)\Delta_a^2/8$;
and $\underline m=\underline m_\infty:=m^-(\tau_N)$ for $|\tau|\ge\tau_N$. Then $\underline m\le m$ on
$\R$, and $\mT\ge\underline{\mathfrak m}$ as forms, where
$\underline{\mathfrak m}(F)=\frac1{2\pi}\int\underline m|\Fh|^2$. If $m^-$ increases along the edges, then also
$\underline m\le\underline m_\infty$.
\end{lemma}
\begin{proof}
On $[0,\tau^{(1)}]$ and beyond $\tau_N$ use that $m$ is increasing with $m(0)=C_0$
(Lemma~\ref{lem:mbounds}(i)). On a band $a\ge1$ let $L_a$ be the chord through the exact values and
$h=m-L_a$: $h$ vanishes at the ends and $h''\le M_a$, so $h\ge0$ if $M_a\le0$, and otherwise
$h+\frac12M_a(\tau-\tau^{(a)})(\tau^{(a+1)}-\tau)$ is concave and vanishes at the ends, whence
$h\ge-M_a\Delta_a^2/8$; lowering the end values to $m^-$ lowers the chord. Every end value of a
lowered chord is at most $m^-$ at that edge, hence at most $m^-(\tau_N)$ when $m^-$ increases along the
edges, and the chords are affine. Finally $\mT-\underline{\mathfrak m}=\frac1{2\pi}\int(m-\underline m)|\Fh|^2\ge0$.
\end{proof}
In the verification, $m^-$ is the lower end of an interval enclosure of $C_0+\sum_ks/(a_k(a_k^2+s))$, $s=\tau^2/4$
(Theorem~\ref{thm:S}(i)), summed to $3000$ terms with the tail of the decreasing summand bracketed by the
integral test; its width is at most $7\cdot10^{-9}$, while consecutive edge values of $m$ differ by at
least $0.03$, so $m^-$ increases along the edges. The bound $M_a$ sums, over $k<3000$, the termwise
majorant $a_k(a_k^2-2s)/(2(a_k^2+s)^3)$ of the summand $a_k(a_k^2-3s)/(2(a_k^2+s)^3)$ of $m''$, which is
decreasing-then-increasing in $s$ and is therefore bounded by its values at the ends of an $s$-interval
containing the band, plus the positive tail $\sum_{k\ge K_0}(2a_k^3)^{-1}\le(4(K_0-\frac34)^2)^{-1}$,
$K_0=3000$. There are $28$ bands: band $0$ is $[0,0.15]$, twenty fixed edges reach $8$, and eight
geometric edges run from $9$ to $\tau_N=c_*/T$, $c_*=16$ (to six decimals); the largest offset is $0.05791$.

\begin{proposition}[the upper certificate]\label{prop:haynsworth}
Fix $\mathrm{par}$, integers $k_1<k_2$ ($22$ and $36$ in the verification) and $\underline m$ as in
Lemma~\ref{lem:multiband}, with $\underline m\le\underline m_\infty$, and put $c_*=T\tau_N$,
$\omega=\underline m_\infty-\min\underline m$ and $\mathcal M=Q_T(\underline m_\infty-\underline m)(D)Q_T$,
so that $\underline{\mathfrak m}(F)=\underline m_\infty\|F\|^2-\langle\mathcal MF,F\rangle$ and
$0\le\mathcal M\le\omega S_{c_*}$. Let $E$ be the span of the $e_j$, $j<k_1$, of parity $\mathrm{par}$,
$Y$ its orthogonal complement in $L^2_{\mathrm{par}}(\IT)$, $P_Y$ the orthogonal projection onto $Y$,
and $\mathfrak b(F,\xi)=\underline{\mathfrak m}(F)+2\xi\langle F,\chi_{\mathrm{par}}\rangle$ on
$L^2_{\mathrm{par}}(\IT)\oplus\R$.
\emph{(a)} On $Y$, $\mathfrak b\ge q_Y:=\underline m_\infty-\omega\Lambda_{k_1}(c_*)$.
\emph{(b)} Suppose $q_Y>0$, and let $\mathsf R$ be the symmetric matrix on $E\oplus\R$ with entries
\[
\mathsf R_{ij}=\underline{\mathfrak m}(e_i,e_j)-\frac{\langle P_Y\mathcal Me_i,P_Y\mathcal Me_j\rangle}{q_Y},\quad
\mathsf R_{i\xi}=\langle e_i,\chi_{\mathrm{par}}\rangle+\frac{\langle P_Y\mathcal Me_i,P_Y\chi_{\mathrm{par}}\rangle}{q_Y},\quad
\mathsf R_{\xi\xi}=-\frac{\|P_Y\chi_{\mathrm{par}}\|^2}{q_Y}.
\]
Then the Schur complement $\mathfrak s$ of $\mathfrak b|_Y$ in $\mathfrak b$ satisfies $\mathfrak s\ge\mathsf R$
and $n_-(\mathfrak b)=n_-(\mathfrak s)$.
\emph{(c)} If $\mathsf R$ is positive definite on a subspace of codimension $r+1$, then
$n_{\mathrm{par}}(T)\le r$.
\end{proposition}
\begin{proof}
(a) $\mathfrak b|_Y=\underline m_\infty-P_Y\mathcal MP_Y$, and
$\|P_Y\mathcal MP_Y\|\le\omega\|P_YS_{c_*}P_Y\|\le\omega\operatorname{tr}(P^\perp S_{c_*}P^\perp)
=\omega\sum_{j\ge k_1}\langle S_{c_*}e_j,e_j\rangle=\omega\Lambda_{k_1}(c_*)$, where $P^\perp$ removes every
$e_j$, $j<k_1$, of either parity: $Y$ lies in its range, the norm of a positive operator is at most its
trace, and the last equality is the trace identity of Lemma~\ref{lem:KyFan}. (b) For $i<k_1$ and $y\in Y$,
$\mathfrak b(e_i,y)=\underline{\mathfrak m}(e_i,y)=-\langle P_Y\mathcal Me_i,y\rangle$, since $\langle e_i,y\rangle=0$,
and the border pairs with $y$ through $+\langle P_Y\chi_{\mathrm{par}},y\rangle$. The congruence
$(x,y)\mapsto(x,y+(\mathfrak b|_Y)^{-1}\mathfrak b_{YX}x)$ gives $\mathfrak b\cong\mathfrak s\oplus\mathfrak b|_Y$
with $\mathfrak b|_Y>0$, so $n_-(\mathfrak b)=n_-(\mathfrak s)$; and
$\mathfrak b_{XY}(\mathfrak b|_Y)^{-1}\mathfrak b_{YX}$ is at most $q_Y^{-1}$ times the Gram matrix of the
$Y$-columns $-P_Y\mathcal Me_i$ and $+P_Y\chi_{\mathrm{par}}$, which is $\mathfrak s\ge\mathsf R$.
(c) Since $\mT\ge\underline{\mathfrak m}$ and $\PT^{\mathrm{par}}=\ker\langle\cdot,\chi_{\mathrm{par}}\rangle$,
Lemma~\ref{lem:border} gives
$n_{\mathrm{par}}(T)\le n_-(\underline{\mathfrak m}|_{\PT^{\mathrm{par}}})=n_-(\mathfrak b)-1=n_-(\mathfrak s)-1\le n_-(\mathsf R)-1\le r$.
\end{proof}
\emph{The computation.} For $i-l$ even, $\widehat{e_i}\overline{\widehat{e_l}}=2T\sqrt{(2i+1)(2l+1)}\,
(-1)^{(l-i)/2}j_ij_l(\tau T)$ (proof of Lemma~\ref{lem:KyFan}), so with $\underline m=\alpha_a+\beta_a|\tau|$
on the band $a$ and $I_p^{(il)}(c)=\int_0^cx^pj_ij_l\,dx$,
\[
\underline{\mathfrak m}(e_i,e_l)=\underline m_\infty\delta_{il}+\frac2\pi\sqrt{(2i+1)(2l+1)}\,(-1)^{\frac{l-i}2}
\Big[\sum_a\big(\alpha_a\,\Delta I^{(il)}_0+\tfrac{\beta_a}T\,\Delta I^{(il)}_1\big)-\underline m_\infty I^{(il)}_0(c_*)\Big],
\]
$\Delta$ denoting the increment over the band, and the entries with $i-l$ odd vanish. The moments
$I_p$ come from the power series of $j_ij_l$ with exact rational coefficients, the alternating tail being
bounded by its first omitted term once the index exceeds $2c$; $\langle e_l,\chi_{\mathrm{par}}\rangle$ comes
from the positive series of the modified spherical Bessel function $i_l(T/2)$ (valid for $T<1$), and
$\|\chi_{\mathrm{par}}\|^2=\sinh T\pm T$. The inner products over $Y$ are exact sums over $k_1\le l<k_2$ of
the given parity ($\mathcal M$ preserves parity) plus boxes from
$\|P^\perp_{k_2}\mathcal Me_i\|^2\le\rho\,\mathcal M_{ii}$, $\rho=\omega\Lambda_{k_2}(c_*)$, and
Cauchy--Schwarz. In the verification $\Lambda_{22}(16)\le7.2151\cdot10^{-6}$ and $\rho\le4.02\cdot10^{-20}$, so
$q_Y\ge\underline m_\infty-5\cdot10^{-5}>0$. For (c), the eigenvectors of the midpoint of $\mathsf R$,
read as decimal data, give a matrix $U_r$ omitting the $r+1$ lowest; interval Cholesky of
$U_r^{\mathsf T}\mathsf RU_r-\mu I$ over the whole interval box of $\mathsf R$ succeeds for a certified
$\mu>0$, which forces $U_r$ to be injective and $\mathsf R$ to be positive definite on its range.
Everything is interval arithmetic at $40$ digits, with no quadrature; the entry widths of $\mathsf R$ are
at most $1.7\cdot10^{-13}$. The first eigenvalue of $\mathsf R$ is always the hyperbolic direction of
the border (Lemma~\ref{lem:border}); $r$ counts the others. The verified cases are:
\begin{center}\small
\begin{tabular}{@{}lcccc@{}}
\toprule
$T$ & sector & $\underline m_\infty$ & certified & margin $\mu$\\
\midrule
$0.58$ & even & $1.479384$ & $n_{\mathrm{ev}}\le0$ & $0.015998$\\
$0.58$ & odd & $1.479384$ & $n_{\mathrm{odd}}\le0$ & $0.047516$\\
$\log2$ & odd & $1.301146$ & $n_{\mathrm{odd}}\le0$ & $0.041093$\\
$0.90$ & even & $1.039940$ & $n_{\mathrm{ev}}\le1$ & $0.185541$\\
$0.90$ & odd & $1.039940$ & $n_{\mathrm{odd}}\le1$ & $0.144840$\\
\bottomrule
\end{tabular}
\end{center}
The cases $\log2$ (even, $\le1$, $\mu=0.420225$) and $0.80$ (even and odd, $\le1$,
$\mu=0.298000$, $0.161624$) were also verified; they follow from the table by
Lemma~\ref{lem:zeroext}.

\begin{proposition}[an explicit negative plane]\label{prop:negplane}
Let $T<\operatorname{arccosh}\sqrt2=0.88137$ and let $F,H\in L^2_{\mathrm{par}}(\IT)$ be
explicit real polynomials, given by their Legendre coefficients, such that $\langle\cdot,\chi_{\mathrm{par}}\rangle$
does not vanish on $\mathrm{span}\{F,H\}$. If the matrix
\[
B_3=\begin{pmatrix}\mT[F]&\mT(F,H)&\langle F,\chi_{\mathrm{par}}\rangle\\
\mT(F,H)&\mT[H]&\langle H,\chi_{\mathrm{par}}\rangle\\
\langle F,\chi_{\mathrm{par}}\rangle&\langle H,\chi_{\mathrm{par}}\rangle&0\end{pmatrix}
\]
is negative definite on a two-dimensional subspace, then $\mathrm{span}\{F,H\}\cap\PT^{\mathrm{par}}$
contains an $F_0\ne0$ with $\mT[F_0]<0$, and $n_{\mathrm{par}}(T)\ge1$.
\end{proposition}
\begin{proof}
$B_3$ is the bordered form of $\mT|_{\mathrm{span}\{F,H\}}$ with border $\langle\cdot,\chi_{\mathrm{par}}\rangle$,
so Lemma~\ref{lem:border} in dimension two gives
$n_-(\mT|_{\mathrm{span}\{F,H\}\cap\PT})=n_-(B_3)-1\ge1$.
\end{proof}
In the verification $F$ is of degree $<18$, $H$ is the Legendre truncation of
$\chi_{\mathrm{par}}$, and $\langle H,\chi_{\mathrm{par}}\rangle\in[0.072,1.69]$. The quadratic values
$\mT[G]=C_0c_G(0)+\int_0^{2T}(c_G(0)-c_G)K+c_G(0)\int_{2T}^\infty K$ (Theorem~\ref{thm:A} with
$\NT=\emptyset$), for $G=F,H,F\pm H$, are enclosed as follows. (a) $c_G$ is computed as an exact polynomial with interval
coefficients, by rational arithmetic on the Legendre recurrence, with no quadrature. (b) $g:=c_G(0)-c_G\ge0$ on
$[0,2T]$ by Cauchy--Schwarz, which allows integrating an enclosure of $K$ against $g$ term by term. (c) On
$[\delta,2T]$, $K$ is squeezed between two affine functions on each panel, since $1/(2\sinh(u/2))$ is convex
and $1/(2\cosh(u/2))$ is concave for $u/2<\operatorname{arccosh}\sqrt2$; this is the range condition
$T<0.88137$; $\int g$ and $\int ug$ are evaluated exactly from the interval antiderivatives. (d) On
$[0,\delta]$, $\delta=10^{-6}$, use $K(u)\le u^{-1}+\frac12$ and $g(u)\le G_1u$ with $G_1$ from the exact
Chebyshev expansion of $g$. (e) The remaining terms are closed forms in intervals:
$C_0=-\gamma-3\log2-\frac\pi2-\log\pi$ and $\int_{2T}^\infty K=\log\coth\frac T2+\pi-2\arctan e^{T}$. The
off-diagonal entry is obtained by polarisation, $\mT(F,H)=\frac14(\mT[F+H]-\mT[F-H])$, and the
pairings with $\chi_{\mathrm{par}}$ from exact moment series.

\begin{lemma}[rounding in the polarisation]\label{lem:polround}
If $F\pm H$ are formed in double precision, the value of $\mT(F,H)$ computed from
the rounded combinations differs from the true one by at most $10^{-10}$, for $0.6\le T\le0.9$ and the
vectors used.
\end{lemma}
\begin{proof}
The rounded combinations are $F\pm H+\delta_\pm$ with Legendre coefficients $|\delta_{\pm,k}|\le3\cdot10^{-16}$
(all coefficients of $F\pm H$ are at most $2.2$ in modulus), so on $\IT$ one has
$\|\delta\|_\infty\le3\cdot10^{-14}$ and $\|\delta'\|_\infty\le4\cdot10^{-12}$ for degree $<18$. The
Dirichlet form of Proposition~\ref{prop:nodes}(i) with $K(u)\le u^{-1}+\frac12$ gives
$\Lev(c_\delta)\le\|\delta'\|_\infty^2T^2(2T+2T^2)+\frac12\|\delta\|_\infty^2\int_{\IT}\kappa\le10^{-22}$,
and then
$|\mT[v+\delta]-\mT[v]|\le|C_0|(2\|v\|\|\delta\|+\|\delta\|^2)+2\sqrt{\Lev(c_v)\Lev(c_\delta)}+\Lev(c_\delta)\le2\cdot10^{-10}$
with $\Lev(c_v)\le50$; divide by four.
\end{proof}
Two further effects are covered by the same inflation: the stored interval ends are rounded to double
precision, and the value of Euler's constant entering $C_0$ is correct only to about $10^{-51}$. The
verification therefore rebuilds $B_3$ from the stored enclosures, inflates every entry by $\pm10^{-9}$, certifies negative definiteness on
a two-plane over the inflated box by interval Cholesky with a certified margin, and re-imposes
$T<0.88137$. It certifies $n_{\mathrm{ev}}\ge1$ at $T=0.60$ (margin $0.006627$) and $n_{\mathrm{odd}}\ge1$
at $T=0.81$ (margin $0.005124$), which is what Theorem~\ref{thm:index} uses, and also $n_{\mathrm{ev}}\ge1$
at $\log2$ and $0.80$ and $n_{\mathrm{odd}}\ge1$ at $0.82$ (margins $0.147170$, $0.290056$, $0.012368$),
which follow from the first two by Lemma~\ref{lem:zeroext}. A run at $T=0.90$ lies outside the range
condition $T<0.88137$ and is not used.


\begin{lemma}[a certificate with $r=0$ is strict]\label{lem:rzero}
If Proposition~\ref{prop:haynsworth}(c) holds with $r=0$ at $T$, then
$\lambda_{\mathrm{par}}(T)>0$.
\end{lemma}
\begin{proof}
Then $\underline{\mathfrak m}\ge0$ on $\PT^{\mathrm{par}}$. At the attained minimiser $F^*$ of
Lemma~\ref{lem:strict}, $\mT[F^*]-\underline{\mathfrak m}[F^*]=\frac1{2\pi}\int(m-\underline m)|\Fh^*|^2>0$,
because $m-\underline m=m-C_0=\varphi(\tau^2)>0$ on $0<|\tau|<\tau^{(1)}$ (Theorem~\ref{thm:S}(i)) and
$\Fh^*$ is entire and not identically zero.
\end{proof}


\subsection{The certificates at \texorpdfstring{$T=\frac12\log5$}{T = (1/2) log 5}}\label{subsec:encllog5}
The notation $\chi_{\mathrm{par}}$, $e_j$, $S_c$, $\Lambda_k$ is that of \S\ref{subsec:enclindex}, and
$\ell_{\mathrm{par}}$, $\varepsilon_{\mathrm{par}}$ are as in \S\ref{sec:complete}. We work with real functions,
which suffices as in \S\ref{sec:complete}. Unlike Proposition~\ref{prop:haynsworth}, the primes are kept:
the lower form is a frequency truncation of the full symbol $\sT$, and the operator subtracted from a
multiple of the identity is indefinite.

\emph{The lower form.} Fix $\tau_N>0$ and a number $s_\infty\le m(\tau_N)-\Aconst$. Since $m$ is increasing
with $m(0)=C_0$ and the comb is at least $-\Aconst$ (Lemma~\ref{lem:mbounds}(i),(iii)), the even function
$\underline s$ equal to $\sT$ on $|\tau|\le\tau_N$ and to $s_\infty$ on $|\tau|>\tau_N$ satisfies
$\underline s\le\sT$. Write $g=s_\infty-\underline s=(s_\infty-\sT)\mathbf 1_{|\tau|<\tau_N}=g_+-g_-$ with
$g_\pm\ge0$, and $\mathcal G_{(\pm)}=Q_Tg_{(\pm)}(D)Q_T$. Then
\[
\mathfrak a_T(F)\ \ge\ s_\infty\|F\|^2-\langle\mathcal GF,F\rangle,\qquad
0\le\mathcal G_\pm\le\omega_\pm S_{c_*},\quad c_*=T\tau_N,
\]
with $\omega_+=s_\infty-C_0+\Aconst$ and $\omega_-=m(\tau_N)+\Aconst-s_\infty$. For $\mu\in\R$ and
$c\in\{0,\frac12,-\frac12\}$ put, on $L^2_{\mathrm{par}}(\IT)\oplus\R$,
\[
\mathfrak b_{\mu,c}(F,\xi)=(s_\infty-\mu)\|F\|^2-\langle\mathcal GF,F\rangle+2\xi\,\ell_{\mathrm{par}}(F)+c\,\xi^2 .
\]

\begin{proposition}[the lower certificate]\label{prop:lowercert}
Let $k_1<k_2$, let $E$ be the span of the $e_j$, $j<k_1$, of parity $\mathrm{par}$, $Y$ its orthogonal
complement in $L^2_{\mathrm{par}}(\IT)$, and $P_Y$ the orthogonal projection onto $Y$.
\emph{(a)} On $Y$, $\mathfrak b_{\mu,c}\ge q_Y:=s_\infty-\mu-\omega_+\Lambda_{k_1}(c_*)$.
\emph{(b)} Suppose $q_Y>0$, and let $\mathsf R$ be the symmetric matrix on $E\oplus\R$ with
\[
\begin{gathered}
\mathsf R_{ij}=(s_\infty-\mu)\delta_{ij}-\langle\mathcal Ge_i,e_j\rangle
-\frac{\langle P_Y\mathcal Ge_i,P_Y\mathcal Ge_j\rangle}{q_Y},\\
\mathsf R_{i\xi}=\langle e_i,\chi_{\mathrm{par}}\rangle+\frac{\langle P_Y\mathcal Ge_i,P_Y\chi_{\mathrm{par}}\rangle}{q_Y},\qquad
\mathsf R_{\xi\xi}=c-\frac{\|P_Y\chi_{\mathrm{par}}\|^2}{q_Y}.
\end{gathered}
\]
Then $n_-(\mathfrak b_{\mu,c})\le n_-(\mathsf R)$, and the same holds for any symmetric $\mathsf R'\le\mathsf R$.
\emph{(c)} If $n_-(\mathsf R')\le1$ with $c=0$, then $\mathfrak a_T\ge\mu$ on the unit sphere of
$\PT^{\mathrm{par}}$. If $n_-(\mathsf R')=0$ with $c=\frac12$ and $\mathrm{par}=\mathrm{odd}$, or $n_-(\mathsf R')\le1$ with
$c=-\frac12$ and $\mathrm{par}=\mathrm{ev}$, then $W_T\ge\mu\|\cdot\|^2$ on $\mathcal D^{\mathrm{par}}_T$.
\end{proposition}
\begin{proof}\sloppy
(a) As in Proposition~\ref{prop:haynsworth}(a): $-P_Y\mathcal GP_Y\ge-P_Y\mathcal G_+P_Y$ and
$\|P_Y\mathcal G_+P_Y\|\le\omega_+\operatorname{tr}(P^\perp S_{c_*}P^\perp)=\omega_+\Lambda_{k_1}(c_*)$.
(b) For $i<k_1$ and $y\in Y$, $\mathfrak b_{\mu,c}(e_i,y)=-\langle P_Y\mathcal Ge_i,y\rangle$ and the border pairs
with $y$ through $+\langle P_Y\chi_{\mathrm{par}},y\rangle$. The congruence of the proof of
Proposition~\ref{prop:haynsworth}(b) gives $\mathfrak b_{\mu,c}\cong\mathfrak s\oplus\mathfrak b_{\mu,c}|_Y$ with
$\mathfrak b_{\mu,c}|_Y\ge q_Y>0$, so $n_-(\mathfrak b_{\mu,c})=n_-(\mathfrak s)$, and the subtracted term of the Schur
complement $\mathfrak s$ is at most $q_Y^{-1}$ times the Gram matrix of the $Y$-columns
$-P_Y\mathcal Ge_i$ and $P_Y\chi_{\mathrm{par}}$; hence $\mathfrak s\ge\mathsf R\ge\mathsf R'$.
(c) Put $\mathfrak q_\mu(F)=\mathfrak b_{\mu,0}(F,0)$, so that $\mathfrak a_T-\mu\|\cdot\|^2\ge\mathfrak q_\mu$. For
$c=0$, Lemma~\ref{lem:border} gives $n_-(\mathfrak q_\mu|_{\ker\ell_{\mathrm{par}}})=n_-(\mathfrak b_{\mu,0})-1=0$. For
$c=\frac12$, $\mathfrak b_{\mu,1/2}=\mathfrak q_\mu-2\ell_{\mathrm{par}}^2+\frac12(\xi+2\ell_{\mathrm{par}})^2$, so
$\mathfrak q_\mu-2\ell_{\mathrm{odd}}^2\ge0$, and $W_T=\mathfrak a_T-2\ell_{\mathrm{odd}}^2$ on odd $F$
(Theorem~\ref{thm:parityschur}(o)). For $c=-\frac12$,
$\mathfrak b_{\mu,-1/2}=\mathfrak q_\mu+2\ell_{\mathrm{par}}^2-\frac12(\xi-2\ell_{\mathrm{par}})^2$; the substitution
$\eta=\xi-2\ell_{\mathrm{par}}(F)$ is a congruence, so
$n_-(\mathfrak q_\mu+2\ell_{\mathrm{ev}}^2)=n_-(\mathfrak b_{\mu,-1/2})-1=0$, and $W_T=\mathfrak a_T+2\ell_{\mathrm{ev}}^2$ on even
$F$. The statements pass from $L^2_{\mathrm{par}}(\IT)$ to the form domain because $\mathfrak q_\mu$ is bounded.
\end{proof}

\emph{The computation.} For $i-l$ even,
\[
\langle\mathcal Ge_i,e_l\rangle=\frac{2T}\pi\sqrt{(2i+1)(2l+1)}\,(-1)^{(l-i)/2}\int_0^{\tau_N}g(\tau)j_i(\tau T)j_l(\tau T)\,d\tau
\]
(proof of Lemma~\ref{lem:KyFan}), and the entries with $i-l$ odd vanish. The integral is computed by composite
Gauss--Legendre quadrature with $46$ nodes on panels of length $\frac12$, with the remainder bound for
functions analytic in a Bernstein ellipse \cite[Thm.~19.3]{Trefethen2013}, $\rho=3.2$. The ellipse lies
within distance $0.439<\frac12$ of the panel, while the poles of $\psi(\frac14\pm\frac{i\tau}2)$ lie at
$|\Iml\tau|\ge\frac12$. The modulus on the ellipse is bounded with $|j_n(z)|\le e^{|\Iml z|}$,
$|\cos(z\log n)|\le\cosh(\Iml z\log n)$ and, for $\psi$, the recurrence $\psi(w)=\psi(w+8)-\sum_{k<8}(w+k)^{-1}$
together with Binet's bound $|\psi(u)|\le\log|u|+\frac\pi2+\frac1{2\Rel u}+\frac1{12(\Rel u)^2}$ for
$\Rel u\ge8$. The resulting remainder, $1.5\cdot10^{-43}$ per entry (below $2\cdot10^{-42}$ under either
point-count convention for the rule), is added as a ball to every entry. The inner products over $Y$ are
exact sums over $k_1\le l<k_2$ of the given parity plus tail boxes: by Cauchy--Schwarz for the positive
operators $\mathcal G_\pm$, $\|P^\perp_{k_2}\mathcal G_\pm e\|^2\le\|P^\perp_{k_2}\mathcal G_\pm P^\perp_{k_2}\|
\langle\mathcal G_\pm e,e\rangle\le\omega_\pm^2\Lambda_{k_2}(c_*)$, so
$\|P^\perp_{k_2}\mathcal Ge\|\le(\omega_++\omega_-)\Lambda_{k_2}(c_*)^{1/2}$, and
$\|P^\perp_{k_2}\chi_{\mathrm{par}}\|^2=\sinh T\pm T-\sum_{l<k_2}\langle e_l,\chi_{\mathrm{par}}\rangle^2$ with
$\langle e_l,\chi_{\mathrm{par}}\rangle=\sqrt{(2l+1)/2T}\,2T\,i_l(T/2)$. The tails $\Lambda_k$ are bounded by
$\Lambda^*_k$ of Corollary~\ref{cor:P1}. Everything is midpoint--radius ball arithmetic at $640$ bits
(Arb \cite{Johansson2017}, through \texttt{python-flint}), including the digamma and Bessel values and the
Gauss--Legendre nodes and weights; every Bessel ball is asserted to have radius below $10^{-60}$, and
comparisons are ball predicates.

\emph{Inertia.} In $\mathsf R$ the term $\|P_Y\chi_{\mathrm{par}}\|^2/q_Y$ is of order $10^{-190}$; it is replaced
by an exact upper bound $\delta'$, which gives $\mathsf R'\le\mathsf R$. When $n_-(\mathsf R')=0$ is to be shown,
interval Cholesky of $\mathsf R'$ is applied. When $n_-(\mathsf R')\le1$ is to be shown, a double-precision
Householder matrix $H$ that sends the midpoint of the column $(\mathsf R'_{i\xi})_i$ to a multiple of $e_0$ is
used as an exact matrix; $\|H^{\mathsf T}H-I\|_\infty<1$ is verified in balls, so $\mathrm{diag}(H,1)$ is a
congruence. In the new coordinates the $2\times2$ block $B$ on $(e_0,\xi)$ has $\det B<0$, certified, so its
inertia is $(1,1)$, and by Haynsworth's inertia additivity \cite{Haynsworth1968}
$n_-(\mathsf R')=1+n_-(S)$ for the Schur complement $S$ of $B$, whose positivity is verified by interval
Cholesky. The only non-rigorous objects are $H$ and the midpoints, which enter only through this exact
congruence.

\emph{Data.} $T=\frac12\log5$, $\tau_N=240$, $c_*=193.13$, $k_1=280$, $k_2=330$ ($\dim E=140$ in each parity).
From $m(\tau_N)=3.64276\ldots$: $s_\infty=0.700787608\ldots$, $\omega_+=9.0150$, $\omega_-=5.8839$,
$\Lambda^*_{k_1}(c_*)\le4.4\cdot10^{-20}$, $\Lambda^*_{k_2}(c_*)\le3.9\cdot10^{-70}$, $q_Y\ge0.70079$, tail
box $(\omega_++\omega_-)\Lambda^*_{k_2}(c_*)^{1/2}\le2.9\cdot10^{-34}$; the entry radii of $\mathsf R$ are at
most $4.4\cdot10^{-41}$. The verified cases are:
\begin{center}\small
\begin{tabular}{@{}llccc@{}}
\toprule
form & sector & certified $\mu$ & minimal pivot at that $\mu$ & fails at\\
\midrule
$\mathfrak a_T|\PT$ ($c=0$) & even & $2.5\cdot10^{-12}$ & $8.47\cdot10^{-3}$ & $3\cdot10^{-12}$\\
$\mathfrak a_T|\PT$ ($c=0$) & odd & $3.0\cdot10^{-10}$ & $2.03\cdot10^{-2}$ & $4\cdot10^{-10}$\\
$W_T$ ($c=-\frac12$) & even & $10^{-18}$ & $4.47\cdot10^{-4}$ & $10^{-16}$\\
$W_T$ ($c=\frac12$) & odd & $10^{-16}$ & $8.14\cdot10^{-14}$ & $10^{-14}$\\
\bottomrule
\end{tabular}
\end{center}
The radius of the smallest pivot in the last row is below $10^{-24}$. A failed Cholesky is not evidence of
negativity. A double-precision eigenvalue computation of the midpoint of $S$ at $\mu=0$ gives
$2.53\cdot10^{-12}$ (even) and $3.62\cdot10^{-10}$ (odd) for the primitive form, so the certified bounds are
within a small factor of the bottom of the lower form. As a check, the same procedure at
$T=\frac12\log3$ ($\tau_N=40$, $k_1=50$) certifies $\lambda^{\mathrm{ev}}_{\mathrm A}(T)\ge6.3\cdot10^{-4}$,
consistent with Theorem~\ref{thm:C}. Each run takes about seven minutes per parity on one core.


\section{Numerical values}\label{app:num}

The values in this appendix are computed in double precision and are not certified; no proof in the paper
depends on them. Galerkin values are upper bounds for the bottoms; the three values of $\lamA$ below are
converged in the basis size.

\begin{center}\small
\begin{tabular}{@{}lll@{}}
\toprule
quantity & value & certified bound in the paper\\
\midrule
$\lambda_{\mathrm A}(\frac12\log3)$ & $9.448\cdot10^{-4}$ & $\ge8.8577\cdot10^{-5}$ (Theorem~\ref{thm:C}), $\ge2.5\cdot10^{-12}$ (Theorem~\ref{thm:log5A})\\
$\lambda_{\mathrm A}(\log2)$ & $9.83\cdot10^{-8}$ & $\ge2.5\cdot10^{-12}$ (Theorem~\ref{thm:log5A})\\
$\lambda_{\mathrm A}(\frac12\log5)$ & $3.10\cdot10^{-12}$ & $\ge2.5\cdot10^{-12}$ (Theorem~\ref{thm:log5A})\\
\bottomrule
\end{tabular}
\end{center}

\begin{itemize}
\item Test-kernel ceilings at $T=\frac12\log3$: $\operatorname{Ceil}(8)=-5.4\cdot10^{-3}$,
$\operatorname{Ceil}(12)=7.6\cdot10^{-5}$, $\operatorname{Ceil}(24)=5.9\cdot10^{-4}$; the certificate of
Theorem~\ref{thm:C} realizes about $15\%$ of its class ceiling.
\item Minimized over $k$, the right side of \eqref{eq:Bmargin} exceeds $1$ for $T<1.430$; the bound of
Theorem~\ref{thm:margin} is asymptotic and vacuous at the windows where the sign is known.
\end{itemize}


% (before \begin{thebibliography}).
\section{Numerical observations}\label{app:diag}

The values in this appendix are computed in double precision (Galerkin values in Legendre bases, quadratures,
bisection on grids). They are numerical and not certified, and no proof in the paper depends on them. A Galerkin
bottom is an upper bound for the true bottom; the reference values of $\lamA$ are those of Appendix~\ref{app:num}.

\subsection{The archimedean form}
\begin{itemize}
\item Galerkin values give $\laminf(T)<0$ already at $T=0.3716$ and $\laminfP(T)<0$ at $T=0.5958$, so numerically
$T^*_\infty\le0.372$ and $T^*_{\infty,\mathrm P}\le0.596$, the second inside the certified bracket $[0.58,0.60]$ of
Corollary~\ref{cor:Tarch}.
\item The Galerkin bottoms of $\mT$ on $\PT^{\mathrm{odd}}$ are about $+0.07$, $+0.03$, $+0.003$ at $T=0.75,0.78,0.80$ and
about $-0.009$ at $T=0.81$, which places the odd crossing $T^*_{\mathrm{odd}}$ of Theorem~\ref{thm:index} in $(0.80,0.81)$;
the sign at $0.80$ is not decided by these values.
\item $\laminfP(\frac12\log3)\approx+0.083$; with the primes the bottom drops to $\lambda_{\mathrm A}(\frac12\log3)\approx9.4\cdot10^{-4}$.
\item At $T=\frac12\log3$ a quadrature of the exact form gives $\Ainf(F^{\mathrm{odd}}_T)\approx-0.0269$, against the
proved upper bound $-0.0201$ of Proposition~\ref{prop:archneg}.
\item At $T=\log2$ the Galerkin bottom of $\mT$ on $\PT$ is about $-0.154$, and on that direction the prime form
supplies about $+0.188$; the nodes of the bottom vector lie at $\pm0.3152$ (one nodal pair, as forced by
Proposition~\ref{prop:nodes}).
\end{itemize}

\subsection{The single direction at \texorpdfstring{$\log2$}{log 2}}
\begin{itemize}
\item At $T=0.58,0.62,0.65,\log2$ the pivot value of \S\ref{subsec:schur} is $w_{11}\approx0.092,0.063,0.042,0.034$ (the
last one certified by Theorem~\ref{thm:w11}), far above $\lamA$ at these windows.
\item At the same four windows the crude sufficient condition $\lambda_2-\|\PiT|_{\mathcal U}\|-c^2/w_{11}>0$, $c=\|b\|$, for
the Schur form of Proposition~\ref{prop:schur} fails by about $1.1$, $1.2$, $1.3$, $1.5$ ($\lambda_j$ the archimedean levels on
$\mathcal U$).
\item By Proposition~\ref{prop:schur}(iv) the Schur correction acts only in the even sector, where the first available
archimedean level is the third; the ground state of the Schur form has about $99\%$ of its mass on the third
archimedean eigenvector at $\log2$. The condition $\lambda_3-c^2/w_{11}>0$, in which the prime form is dropped, fails by
about $0.06$, $0.1$, $0.15$, $0.3$ at $T=0.58,0.62,0.65,\log2$.
\item In the even archimedean basis at $T=\log2$, the entry $\pi_{13}$ of the prime matrix that makes the $2$-plane of
Theorem~\ref{thm:nest} positive is a cancellation between the two atoms: $\pi_{13}\approx+0.187\ (\text{atom }2)-0.141\
(\text{atom }3)=+0.046$; with the atom at $\log2$ alone the $2\times2$ determinant would be negative ($\approx-0.03$).
\end{itemize}

\begin{remark}\label{rem:nodal}
At $T=\log2$, with $F^*=e_1$ and $e_3$ as in Theorem~\ref{thm:nest}, the overlap $I(a)=I_{F^*,e_3}(a)$ of
Lemma~\ref{lem:N} vanishes at $a\approx0.213$, $0.609$ and $1.023$, and the two contacts straddle the third zero:
$\log2=0.693<1.023<1.099=\log3$. That straddle is the cancellation above; the two atoms computed from
Lemma~\ref{lem:N} agree with the matrix entries to three digits. The forced-sign threshold of Proposition~\ref{prop:G}
is $a^*=T+z_{e_3}\approx1.17$, which $\log3$ misses by about $7\%$: one more contact, or a window a few per cent wider,
and both atoms would be forced to the same sign. The same straddle holds at $T=0.62$ and $T=0.65$. The node of $F^*$
is forced (Proposition~\ref{prop:nodes}(ii)).
\end{remark}

\begin{remark}\label{rem:signrule}
The nodes decide the large entries of the prime matrix and only those. Refine $(\log n-T,T)$ by the nodes of $e_j$
and by those of $e_i$ shifted by $\log n$; on each cell the integrand of Lemma~\ref{lem:N} has a constant sign, readable
from the nodes alone, and predict the sign of the $n$-th atom of $\PiT(e_i,e_j)$ from the longest cell
(Proposition~\ref{prop:G} is the case of a single cell, where the prediction is a theorem). On the $8\times8$ even block
at $T=\log2$ the rule predicts every entry with $|I|\ge0.10$, which together carry about $90\%$ of $\sum|I|$, and fails
only on small, nearly cancelling entries. In the same computation the $k$-th even eigenvector of $\mT$ on $\PT$ has
exactly $k$ nodes in $(0,T)$ for $k\le8$, with strict interlacing of consecutive modes; beyond the base case of
Proposition~\ref{prop:nodes}(ii)--(iii) this Sturm pattern is not proved.
\end{remark}

\subsection{The negative set of the symbol}
\begin{remark}\label{rem:bandtable}
The negative set of $\sT$ in the two regimes of Proposition~\ref{prop:align}, by bisection on a grid of $4\cdot10^5$
points. For $\NT=\{2,3\}$ its components in $\tau\ge0$ are $(0,2.278)$, $(4.816,7.272)$, $(9.666,11.966)$,
$(16.251,18.494)$, $(27.558,28.938)$, $(45.149,46.129)$, so $\tau_{\mathrm{last}}\approx46.13$ (consistent with
Remark~\ref{rem:taulast} and Proposition~\ref{prop:env}(iii)); for $\NT=\{2,3,4\}$ the same six within $0.6$, then
$(62.616,63.777)$ and $(108.495,108.950)$, so $\tau_{\mathrm{last}}\approx108.95$. The components beyond $\tau_a$ sit at the pairs
$(k,l)=(3,5)$, $(5,8)$ for $\{2,3\}$, with $2^l/3^k=1.185$, $1.053$, and additionally $(7,11)$, $(12,19)$ for $\{2,3,4\}$,
with $2^l/3^k=0.936$, $0.9865$; the fractions $8/5$ and $19/12$ carrying the last components are convergents of
$\log3/\log2$. The traces of Proposition~\ref{prop:negtrace}(i) are $\operatorname{tr}\mathrm{Neg}_T\approx5.2$ at $T=\log2$ (with
$\sup\sigma_-\approx7.6$) and $\approx7.0$ at $T=\frac12\log5$; the bound $\lamA\ge-\operatorname{tr}\mathrm{Neg}_T$ is useless for the sign.
The test function of Theorem~\ref{thm:margin}, evaluated by quadrature of the symbol (no zeros enter), has Rayleigh
quotient smaller than the proved bound $\mathcal B(T,k)$ by six to ten orders of magnitude for
$\frac12\log3\le T\le1.5$.
\end{remark}

\subsection{The polar capacity}
\begin{remark}\label{rem:polarcapdiag}
Legendre--Galerkin values of the constants $\mathfrak g_{\mathrm{par}}$ of Theorem~\ref{thm:parityschur}. A Galerkin
$\mathfrak g_{\mathrm{par}}$ is a minimum over a subspace, hence an upper bound for the true constant: a negative value is
evidence of a true negative, up to rounding and quadrature error, and a positive value is no evidence. At
$T=\frac12\log3$ with the primes deleted, $\mathfrak g_{\mathrm{ev}}\approx-0.45$ and $\mathfrak g_{\mathrm{odd}}\approx-20$, with primitive
sector bottoms $\approx+0.08$ and $+0.38$, so both channels fail. For $\zeta$ at the same window
$\mathfrak g_{\mathrm{ev}}\approx9\cdot10^{-8}$ and $\mathfrak g_{\mathrm{odd}}\approx1\cdot10^{-3}$, stable over $16$, $24$ and $32$ modes, and
$\mathrm{diag}(\mathfrak a_T(h_{\mathrm{ev}}),\mathfrak a_T(h_{\mathrm{odd}}))\approx\mathrm{diag}(-2,+2)$: the primitive Schur complement nearly
cancels the polar block. In the family $a_\theta$ of Corollary~\ref{cor:convfail}(ii) the complete form fails for $\theta$
below a value within about $10^{-4}$ of $1$, and its even channel already fails at $\theta=1.001$. With the primes
deleted, the odd and the even channel of the complete form cross zero at $T\approx0.37$ and $T\approx0.44$. The positivity
of $\mathfrak g_{\mathrm{ev}},\mathfrak g_{\mathrm{odd}}$ for $\zeta$ at $T\le\frac12\log5$ rests on Theorem~\ref{thm:log5W} and
Theorem~\ref{thm:parityschur}(c), not on these numbers.
\end{remark}

\subsection{Flat-top weights and lattice classes}
The flat-top identity is checked numerically in Remark~\ref{rem:Zcheck}, and the comb-vanishing class in Remark~\ref{rem:comblp}.

\begin{itemize}
\item The test-kernel ceilings of Proposition~\ref{prop:TK} for the archimedean form alone approach $\laminfP$ from
below at $T=\frac12\log3$ and $T=\log2$ with a deficit of about $0.3\,M^{-1.2}$ over $8\le M\le64$. This concerns the
archimedean program, where the polar block is absent and there is no edge; it is not comparable with the edge rate of
Theorem~\ref{thm:latedge}.
\item Measured Galerkin norms of the prime form on $\PT$ are about $0.5$ on the one-contact range and $0.7$--$0.9$ on
the two-contact range, against the proved bounds of Proposition~\ref{prop:coer-prime}; the jump is the entry of the
contact $3$.
\end{itemize}


\begin{thebibliography}{99}
\bibitem{Branden2007} P.~Br\"and\'en, Polynomials with the half-plane property and matroid theory, \emph{Adv. Math.} \textbf{216} (2007), 302--320.
\bibitem{Boas1954} R.~P.~Boas, \emph{Entire Functions}, Academic Press, New York, 1954.

\bibitem{BoasKac1945} R.~P.~Boas and M.~Kac, \emph{Inequalities for Fourier transforms of positive
functions}, Duke Math.~J. \textbf{11} (1945), 189--206.

\bibitem{Bochner1951} S.~Bochner, \emph{Some properties of modular relations}, Ann.\ of Math.\ (2) \textbf{53} (1951), 332--363.

\bibitem{Bombieri2000} E.~Bombieri, \emph{Remarks on Weil's quadratic functional in the theory of
prime numbers, I}, Atti Accad.~Naz.~Lincei Cl.~Sci.~Fis.~Mat.~Natur.~Rend.~Lincei (9) Mat.~Appl.
\textbf{11} (2000), no.~3, 183--233.

\bibitem{BottcherSilbermann2006} A.~B\"ottcher and B.~Silbermann, \emph{Analysis of Toeplitz Operators}, 2nd ed., Springer, Berlin, 2006.

\bibitem{CCM2023} A.~Connes, C.~Consani and H.~Moscovici, \emph{Zeta zeros and prolate wave operators}, arXiv:2310.18423 (2023).

\bibitem{CCM2024} A.~Connes, C.~Consani and H.~Moscovici, \emph{On $q$-series and the moment problem associated to local factors}, arXiv:2403.01247 (2024).

\bibitem{CCquasi2020} A.~Connes and C.~Consani, \emph{Quasi-inner functions and local factors}, arXiv:2008.10974 (2020).

\bibitem{CohnElkies2003} H.~Cohn and N.~Elkies, \emph{New upper bounds on sphere packings I},
Ann.~of Math. \textbf{157} (2003), 689--714.

\bibitem{ConnesConsani2021} A.~Connes and C.~Consani, \emph{Weil positivity and trace formula, the archimedean place}, Selecta Math.\ (N.S.) \textbf{27} (2021), no.~4, Paper No.~77; arXiv:2006.13771.

\bibitem{CvS2025} A.~Connes and W.~D.~van Suijlekom, \emph{Quadratic forms, real zeros and echoes of the spectral action}, arXiv:2511.23257 (2025).

\bibitem{DavenportHeilbronn1936} H.~Davenport and H.~Heilbronn, \emph{On the zeros of certain
Dirichlet series}, J.~London Math.~Soc. \textbf{11} (1936), 181--185 and 307--312.

\bibitem{Davenport2000} H.~Davenport, \emph{Multiplicative Number Theory}, 3rd ed., Grad.\ Texts in Math.\ \textbf{74}, Springer, New York, 2000.

\bibitem{DiamondZhang2016} H.~G.~Diamond and W.-B.~Zhang, \emph{Beurling Generalized Numbers}, Math.\ Surveys Monogr.\ \textbf{213}, Amer.~Math.~Soc., Providence, RI, 2016.

\bibitem{DLMF} NIST Digital Library of Mathematical Functions, \url{https://dlmf.nist.gov}, \S5.9.

\bibitem{Groskin2026} A.~Groskin, \emph{A finite Guinand--Weil dictionary and archimedean tail
order for the truncated Weil quadratic form}, arXiv:2607.02828 (2026).

\bibitem{Halmos1969} P.~R.~Halmos, \emph{Two subspaces}, Trans.\ Amer.\ Math.\ Soc.\ \textbf{144} (1969), 381--389.

\bibitem{Haynsworth1968} E.~V.~Haynsworth, \emph{Determination of the inertia of a partitioned Hermitian
matrix}, Linear Algebra Appl.\ \textbf{1} (1968), 73--81.

\bibitem{Hamburger1921} H.~Hamburger, \emph{\"Uber die Riemannsche Funktionalgleichung der $\zeta$-Funktion}, Math.\ Z.\ \textbf{10} (1921), 240--254.

\bibitem{Holder1887} O.~H\"older, \emph{\"Uber die Eigenschaft der Gammafunction, keiner algebraischen Differentialgleichung zu gen\"ugen}, Math.\ Ann.\ \textbf{28} (1887), 1--13.

\bibitem{HLP} G.~H.~Hardy, J.~E.~Littlewood and G.~P\'olya, \emph{Inequalities}, 2nd ed., Cambridge Univ.~Press, 1952.

\bibitem{IK} H.~Iwaniec and E.~Kowalski, \emph{Analytic Number Theory}, Amer.~Math.~Soc.
Colloq.~Publ.~53, 2004.

\bibitem{Johansson2017} F.~Johansson, \emph{Arb: efficient arbitrary-precision midpoint-radius interval
arithmetic}, IEEE Trans.\ Comput.\ \textbf{66} (2017), 1281--1292.

\bibitem{KahaneMandelbrojt1958} J.-P.~Kahane and S.~Mandelbrojt, \emph{Sur l'\'equation fonctionnelle de Riemann et la formule sommatoire de Poisson}, Ann.\ Sci.\ \'Ecole Norm.\ Sup.\ (3) \textbf{75} (1958), no.~1, 57--80. DOI 10.24033/asens.1066.

\bibitem{Karlin1968} S.~Karlin, \emph{Total Positivity, Vol.~I}, Stanford Univ.~Press, 1968.

\bibitem{Katznelson} Y.~Katznelson, \emph{An Introduction to Harmonic Analysis}, 3rd ed.,
Cambridge Univ.~Press, 2004.

\bibitem{Kim2026} T.~Kim, Y.~Hong, M.~Kim, S.~Choi, J.~Jang and M.~Kim, \emph{A numerical realization of Suzuki's Weil-quadratic-form operator: the archimedean spectral law, its universality, and an operator form of Weil's positivity criterion}, arXiv:2607.24830 (2026).

\bibitem{Kingman1993} J.~F.~C.~Kingman, \emph{Poisson Processes}, Oxford Univ.~Press, Oxford, 1993.

\bibitem{Korevaar2004} J.~Korevaar, \emph{Tauberian Theory: A Century of Developments}, Grundlehren Math.\ Wiss.\ \textbf{329}, Springer, Berlin, 2004.

\bibitem{Krein1940} M.~G.~Kre\u{\i}n, \emph{Sur le probl\`eme du prolongement des fonctions
hermitiennes positives et continues}, Dokl.~Akad.~Nauk SSSR \textbf{26} (1940), 17--22.

\bibitem{LandauWidom1980} H.~J.~Landau and H.~Widom, \emph{Eigenvalue distribution of time and frequency
limiting}, J.~Math.\ Anal.\ Appl.\ \textbf{77} (1980), 469--481.

\bibitem{Liu2026} V.~Liu, \emph{Certified Weil positivity beyond the unit window: source-exact
block-Schur and tail-compensation bounds for the Riemann zeta function}, manuscript dated 14 Sep 2026,
alphaXiv; under review at Math.\ Comp.\ (unrefereed; AI tools declared by the author).

\bibitem{Mori15} R.~Mori, \emph{Weil positivity on the window of width $\log 15$}, Zenodo (2026),
doi:\nolinkurl{10.5281/zenodo.23281236}; certificate code doi:\nolinkurl{10.5281/zenodo.23092422}, data
doi:\nolinkurl{10.5281/zenodo.23281176} (unrefereed; AI tools declared by the author).

\bibitem{MoriUpper} R.~Mori, \emph{Improved upper bounds for windowed Weil forms with exponential factor
$e^{-4\pi\mu}$}, Zenodo (2026), v3 doi:\nolinkurl{10.5281/zenodo.23363767}; concept doi:\nolinkurl{10.5281/zenodo.23175060}
(unrefereed).

\bibitem{Spira1994} R.~Spira, \emph{Some zeros of the Titchmarsh counterexample}, Math.\ Comp.\ 63 (1994), 747--748.
\bibitem{ZiadDH} Ziad, \emph{Localized constrained Weil forms: near-degeneracy and a Davenport--Heilbronn first
kernel}, Zenodo (2026), doi:\nolinkurl{10.5281/zenodo.23145454} (unrefereed).

\bibitem{McClellanParks1972} J.~H.~McClellan and T.~W.~Parks, \emph{Eigenvalue and eigenvector decomposition of the discrete Fourier transform}, IEEE Trans.\ Audio Electroacoust.\ \textbf{20} (1972), 66--74.

\bibitem{MontgomeryVaughan2007} H.~L.~Montgomery and R.~C.~Vaughan, \emph{Multiplicative Number Theory I. Classical Theory}, Cambridge Univ.~Press, Cambridge, 2007.

\bibitem{Ostrowski1920} A.~Ostrowski, \emph{\"Uber Dirichletsche Reihen und algebraische Differentialgleichungen}, Math.\ Z.\ \textbf{8} (1920), 241--298.

\bibitem{Paper46} D.~A.~Trejo Pizzo, \emph{Row D: the sign}, preprint, 2026.

\bibitem{Paper45} D.~A.~Trejo Pizzo, \emph{An explicit formula and a comparison theorem for the canonical trace on a correspondence
algebra of the idele class group}, preprint, 2026.

\bibitem{Petrov1975} V.~V.~Petrov, \emph{Sums of Independent Random Variables}, Springer, Berlin, 1975.

\bibitem{PlattTrudgian2021} D.~J.~Platt and T.~S.~Trudgian, \emph{The Riemann hypothesis is true up to
$3\cdot10^{12}$}, Bull.\ London Math.\ Soc.\ \textbf{53} (2021), 792--797.

\bibitem{Rockafellar} R.~T.~Rockafellar, \emph{Convex Analysis}, Princeton Univ.~Press, 1970.

\bibitem{Rosser1941} J.~B.~Rosser, \emph{Explicit bounds for some functions of prime numbers}, Amer.\ J.\
Math.\ \textbf{63} (1941), 211--232.

\bibitem{RosserYoheSchoenfeld1968} J.~B.~Rosser, J.~M.~Yohe and L.~Schoenfeld, \emph{Rigorous computation and
the zeros of the Riemann zeta-function}, Proc.\ IFIP Congress 1968, 70--76.

\bibitem{Sasvari1994} Z.~Sasv\'ari, \emph{Positive Definite and Definitizable Functions},
Akademie Verlag, Berlin, 1994.

\bibitem{Schoenberg1951} I.~J.~Schoenberg, \emph{On P\'olya frequency functions.~I. The
totally positive functions and their Laplace transforms}, J.~Analyse Math. \textbf{1} (1951), 331--374.

\bibitem{Simon2005} B.~Simon, \emph{Trace Ideals and Their Applications}, 2nd ed., Math.\ Surveys Monogr.\ \textbf{120}, Amer.~Math.~Soc., 2005.

\bibitem{SongVondracek2003} R.~Song and Z.~Vondra\v{c}ek, \emph{Potential theory of subordinate
killed Brownian motion in a domain}, Probab.~Theory Related Fields \textbf{125} (2003), 578--592.

\bibitem{SSV2012} R.~L.~Schilling, R.~Song and Z.~Vondra\v{c}ek, \emph{Bernstein Functions:
Theory and Applications}, 2nd ed., de Gruyter, 2012.

\bibitem{Suzuki2023jlms} M.~Suzuki, \emph{Aspects of the screw function corresponding to the
Riemann zeta function}, J.~Lond.~Math.~Soc.~(2) \textbf{107} (2023), 1846--1888.

\bibitem{Suzuki2023screwline} M.~Suzuki, \emph{The screw line of the Riemann zeta-function and
its applications}, arXiv:2209.04658.

\bibitem{Suzuki2025hilbert} M.~Suzuki, \emph{On the Hilbert space derived from the Weil
distribution}, arXiv:2301.00421.

\bibitem{Trefethen2013} L.~N.~Trefethen, \emph{Approximation Theory and Approximation Practice},
SIAM, Philadelphia, 2013.

\bibitem{Titchmarsh1986} E.~C.~Titchmarsh, \emph{The Theory of the Riemann Zeta-Function}, 2nd ed., revised by D.~R.~Heath-Brown,
Oxford Univ.\ Press, Oxford, 1986.

\bibitem{Weil52} A.~Weil, \emph{Sur les ``formules explicites'' de la th\'eorie des nombres
premiers}, Comm.~S\'em.~Math.~Univ.~Lund (1952), 252--265.

\bibitem{Yoshida1992} H.~Yoshida, \emph{On Hermitian forms attached to zeta functions}, in:
\emph{Zeta Functions in Geometry}, Adv.~Stud.~Pure Math.~\textbf{21}, Kinokuniya, Tokyo, 1992, 281--325.

\bibitem{Guo2026} Guo~Chen, \emph{Prolate leakage and near-zero spectral clusters of the semilocal Weil operator},
Zenodo (2026), doi:\nolinkurl{10.5281/zenodo.23297284} (unrefereed).

\bibitem{Fuchs1964} W.~H.~J.~Fuchs, \emph{On the eigenvalues of an integral equation arising in the theory of band-limited
signals}, J.\ Math.\ Anal.\ Appl.\ \textbf{9} (1964), 317--330.

\bibitem{Stein1993} E.~M.~Stein, \emph{Harmonic Analysis: Real-Variable Methods, Orthogonality, and Oscillatory Integrals},
Princeton Univ.~Press, Princeton, 1993.

\bibitem{Olver1974} F.~W.~J.~Olver, \emph{Asymptotics and Special Functions}, Academic Press, New York, 1974.

\bibitem{Dunster1986} T.~M.~Dunster, \emph{Uniform asymptotic expansions for prolate spheroidal functions with large
parameters}, SIAM J.\ Math.\ Anal.\ \textbf{17} (1986), 1495--1524.

\bibitem{OsipovRokhlinXiao2013} A.~Osipov, V.~Rokhlin and H.~Xiao, \emph{Prolate Spheroidal Wave Functions of Order Zero},
Springer, New York, 2013.

\bibitem{SlepianPollak1961} D.~Slepian and H.~O.~Pollak, \emph{Prolate spheroidal wave functions, Fourier analysis and
uncertainty~I}, Bell System Tech.\ J.\ \textbf{40} (1961), 43--63.

\bibitem{Zhu2026} X.~Zhu, \emph{Weil positivity in compact windows: a finite reduction, certified
two-sided bounds, and a Landau--Widom decay law}, arXiv:2608.24827 (2026).

\end{thebibliography}

\end{document}

